\unitopener{1}{Fundamentals \& Data Representation}
\lessonheading{01}{What Is a Computer?}
\begin{enumerate}[label=\textbf{Q\arabic*.}, leftmargin=*]
\item \textit{[Original]} Which description best captures the basic role of
  a computer?
  \begin{enumerate}[label=\Alph*.]
    \item It accepts data, processes it under instructions, produces output,
      and can store data or results.
    \item It decides whether every input is truthful before processing.
    \item It is a display device that presents information but cannot store.
    \item It stores data but cannot process instructions.
  \end{enumerate}
\item \textit{[Original]} A clerk enters three raw marks into a spreadsheet.
  Before any total or average is calculated, the marks are best described as:
  \begin{enumerate}[label=\Alph*.]
    \item Information
    \item Data
    \item Output
    \item Storage
  \end{enumerate}
\item \textit{[Original]} Which operation is processing in a marks
  spreadsheet?
  \begin{enumerate}[label=\Alph*.]
    \item Typing the marks
    \item Displaying the average
    \item Applying a formula to calculate the average
    \item Saving the workbook
  \end{enumerate}
\item \textit{[Original]} In IPO, the letter P stands for:
  \begin{enumerate}[label=\Alph*.]
    \item Printing
    \item Processing
    \item Programming
    \item Peripheral
  \end{enumerate}
\item \textit{[Original]} Which sequence gives the IPO cycle in its usual
  order?
  \begin{enumerate}[label=\Alph*.]
    \item Output, input, processing
    \item Input, processing, output
    \item Processing, output, input
    \item Storage, input, output
  \end{enumerate}
\item \textit{[Original]} When storage is added to IPO, the four-part
  shorthand is:
  \begin{enumerate}[label=\Alph*.]
    \item IPOS
    \item PIOS
    \item OIPS
    \item SIPO
  \end{enumerate}
\item \textit{[Original]} A saved spreadsheet is an example of which
  function?
  \begin{enumerate}[label=\Alph*.]
    \item Input
    \item Processing
    \item Output
    \item Storage
  \end{enumerate}
\item \textit{[Original]} Which statement about computer accuracy is
  \textbf{NOT} correct?
  \begin{enumerate}[label=\Alph*.]
    \item Correct instructions can produce consistent calculations.
    \item A wrong formula can produce a wrong result.
    \item High processing speed proves that input data are true.
    \item Incorrect input can lead to an incorrect result.
  \end{enumerate}
\item \textit{[Original]} The phrase ``Garbage In, Garbage Out'' warns that:
  \begin{enumerate}[label=\Alph*.]
    \item Poor input can lead to poor output.
    \item Output is always more reliable than input.
    \item Storage automatically repairs incorrect data.
    \item A computer cannot process numerical data.
  \end{enumerate}
\item \textit{[Original]} Which characteristic refers to performing many
  operations quickly?
  \begin{enumerate}[label=\Alph*.]
    \item Diligence
    \item Speed
    \item Versatility
    \item Storage
  \end{enumerate}
\item \textit{[Original]} A computer repeats a long calculation without
  tiring in the human sense. This characteristic is:
  \begin{enumerate}[label=\Alph*.]
    \item Diligence
    \item Accuracy
    \item Versatility
    \item Input
  \end{enumerate}
\item \textit{[Original]} Running a word processor, spreadsheet, and browser
  on the same computer illustrates:
  \begin{enumerate}[label=\Alph*.]
    \item Diligence
    \item Versatility
    \item Storage
    \item Output
  \end{enumerate}
\item \textit{[Original]} A programmed task continues automatically after a
  user starts it and supplies the required input. This is:
  \begin{enumerate}[label=\Alph*.]
    \item Automation
    \item Data
    \item Output
    \item GIGO
  \end{enumerate}
\item \textit{[Original]} Evaluate the statements:
  \begin{enumerate}[label=\Roman*.]
    \item Data can be raw facts or symbols.
    \item Information is data organized or processed for a purpose.
    \item The same value can be data in one task and information in another.
  \end{enumerate}
  \begin{enumerate}[label=\Alph*.]
    \item I only
    \item I and II only
    \item II and III only
    \item I, II, and III
  \end{enumerate}
\item \textit{[Original]} Which is an example of computer use in
  healthcare?
  \begin{enumerate}[label=\Alph*.]
    \item Maintaining patient records and scheduling appointments
    \item Replacing every clinical decision with an automatic guarantee
    \item Turning a monitor into secondary storage
    \item Making input data truthful by processing them quickly
  \end{enumerate}
\item \textit{[Original]} Which statement is \textbf{NOT} a limitation of a
  computer as described in basic computer fundamentals?
  \begin{enumerate}[label=\Alph*.]
    \item It does not independently determine whether input is truthful.
    \item Its output may need interpretation in context.
    \item Incorrect instructions may produce incorrect results.
    \item It performs calculations quickly when given suitable instructions.
  \end{enumerate}
\item \textit{[Original]} A displayed average is output in the current task.
  If the value is typed into a new program tomorrow, it is then acting as:
  \begin{enumerate}[label=\Alph*.]
    \item Input
    \item Processing
    \item Storage
    \item Accuracy
  \end{enumerate}
\item \textit{[Original]} A user enters expense amounts and a formula totals
  them. Which pair is matched correctly?
  \begin{enumerate}[label=\Alph*.]
    \item Amounts -- output; formula -- input
    \item Amounts -- input; formula -- processing
    \item Amounts -- storage; formula -- output
    \item Amounts -- processing; formula -- storage
  \end{enumerate}
\item \textit{[Original]} Which statement about a computer's output is
  most accurate?
  \begin{enumerate}[label=\Alph*.]
    \item Output is guaranteed to be fair if the computer is fast.
    \item Output is valid even when the input is incomplete.
    \item Output depends on the input and instructions used.
    \item Output can only be shown on a monitor.
  \end{enumerate}
\item \textit{[Original]} A user types 18 instead of 81, but the spreadsheet
  correctly calculates the average from the entered values. What best
  explains the misleading result?
  \begin{enumerate}[label=\Alph*.]
    \item The formula is processing the supplied data; the input is wrong.
    \item The monitor has converted output into input.
    \item Storage has increased the marks.
    \item The computer's speed has validated the wrong entry.
  \end{enumerate}

\item \textit{[JKSSB FAA-pattern]} Which of the following best distinguishes data from information?
  \begin{enumerate}[label=\Alph*.]
    \item Data are raw facts; information is processed data with meaning.
    \item Information is raw; data is always meaningful.
    \item Data and information are identical in every task.
    \item Information cannot be stored, but data can.
  \end{enumerate}
\item \textit{[JKPSI-pattern]} A computer is described as ``diligent'' because it:
  \begin{enumerate}[label=\Alph*.]
    \item can repeat long calculations without tiring in the human sense.
    \item understands human emotions during long tasks.
    \item refuses to repeat any instruction twice.
    \item slows down in order to stay accurate over time.
  \end{enumerate}
\item \textit{[SSC CGL-pattern]} Which of the following is NOT a characteristic of a computer?
  \begin{enumerate}[label=\Alph*.]
    \item Speed
    \item Accuracy
    \item Diligence
    \item Intuition
  \end{enumerate}
\item \textit{[JKSSB-pattern]} \textbf{(MSQ: select ALL correct options.)} Which of the following are basic functions of a computer?
  \begin{enumerate}[label=\Alph*.]
    \item Accepting data as input
    \item Processing data under instructions
    \item Producing results as output
    \item Feeling human emotions about the result
  \end{enumerate}
\item \textit{[SSC CGL-pattern]} \textbf{(MSQ: select ALL correct options.)} Which statements about the IPO cycle are correct?
  \begin{enumerate}[label=\Alph*.]
    \item Input is accepted before processing begins.
    \item Output depends on the input and the instructions used.
    \item Storage is the letter S in the IPO abbreviation.
    \item Processing converts data into information.
  \end{enumerate}
\item \textit{[Original]} In the IPO cycle, which stage converts processed data into information for a user?
  \begin{enumerate}[label=\Alph*.]
\item Input
   \item Processing
   \item Output
   \item Storage
  \end{enumerate}

\item \textit{[Original]} Which property describes a computer's ability to repeat a task for long periods without tiring?
  \begin{enumerate}[label=\Alph*.]
\item Diligence
   \item Versatility
   \item Accuracy
   \item Automation
  \end{enumerate}

\item \textit{[Original]} A list of unprocessed daily sales entries is best described as:
  \begin{enumerate}[label=\Alph*.]
\item Information
   \item Data
   \item Knowledge
   \item Output
  \end{enumerate}

\item \textit{[Original]} Which statement is a limitation of a computer?
  \begin{enumerate}[label=\Alph*.]
\item It can process data rapidly.
   \item It can store large amounts of data.
   \item It does not independently decide its goals without instructions.
   \item It can perform repetitive work consistently.
  \end{enumerate}

\item \textit{[Original, MSQ drill]} \textbf{(MSQ: select ALL correct options.)} Which statements about the IPO cycle are correct?
  \begin{enumerate}[label=\Alph*.]
\item Input provides data or instructions.
   \item Processing applies operations to data.
   \item Output presents a result.
   \item Storage is the same stage as input.
  \end{enumerate}

\end{enumerate}
\lessonheading{02}{Types of Computers}
\begin{enumerate}[label=\textbf{Q\arabic*.}, leftmargin=*]
\item \textit{[PYQ-adapted: JKSSB Computer Instructor/Operator 2025 Q5;
paraphrased from the official paper]} Which pair correctly contrasts the
characteristic workload of a mainframe and a supercomputer?
  \begin{enumerate}[label=\Alph*.]
    \item Mainframes emphasize many concurrent transactions; supercomputers
      emphasize very large, high-performance calculations.
    \item Mainframes are analog; supercomputers are embedded controllers.
    \item Mainframes are personal-use computers; supercomputers are only
      mobile devices.
    \item Both labels refer only to physical dimensions.
  \end{enumerate}
\item \textit{[Original]} A computer represents a measured quantity as a
  continuously varying value. By signal type, it is:
  \begin{enumerate}[label=\Alph*.]
    \item Digital
    \item Analog
    \item Mainframe
    \item Embedded
  \end{enumerate}
\item \textit{[Original]} A modern laptop ordinarily represents data using
  discrete binary states. By signal representation it is:
  \begin{enumerate}[label=\Alph*.]
    \item Analog
    \item Digital
    \item Hybrid only
    \item Minicomputer
  \end{enumerate}
\item \textit{[Original]} A system combines a continuous sensor measurement
  with digital processing and display. The complete arrangement is best
  called:
  \begin{enumerate}[label=\Alph*.]
    \item Analog
    \item Mainframe
    \item Hybrid
    \item Supercomputer
  \end{enumerate}
\item \textit{[Original]} Which computer class is traditionally associated
  with personal use and a microprocessor?
  \begin{enumerate}[label=\Alph*.]
    \item Microcomputer
    \item Mainframe
    \item Minicomputer
    \item Supercomputer
  \end{enumerate}
\item \textit{[Original]} A question describes a historical mid-range
  multi-user class. Which term is the expected match?
  \begin{enumerate}[label=\Alph*.]
    \item Supercomputer
    \item Minicomputer
    \item Microcomputer
    \item Analog computer
  \end{enumerate}
\item \textit{[Original]} Which class is most closely associated with
  reliable, high-volume transaction processing for many users?
  \begin{enumerate}[label=\Alph*.]
    \item Mainframe
    \item Supercomputer
    \item Embedded controller
    \item Analog computer
  \end{enumerate}
\item \textit{[Original]} A research facility runs very large parallel
  scientific calculations. The most appropriate class is:
  \begin{enumerate}[label=\Alph*.]
    \item Microcomputer
    \item Mainframe
    \item Supercomputer
    \item Minicomputer
  \end{enumerate}
\item \textit{[Original]} Which statement is \textbf{NOT} correct?
  \begin{enumerate}[label=\Alph*.]
    \item A laptop can be both digital and a microcomputer.
    \item A desktop label describes a form factor.
    \item Mainframe and supercomputer are distinguished only by cabinet size.
    \item An embedded computer is built into a larger device.
  \end{enumerate}
\item \textit{[Original]} A smartphone is best classified as:
  \begin{enumerate}[label=\Alph*.]
    \item A digital microcomputer and a mobile form factor
    \item An analog mainframe
    \item A supercomputer because it is portable
    \item A minicomputer by signal representation
  \end{enumerate}
\item \textit{[Original]} A washing machine contains a controller designed
  for appliance functions. It is:
  \begin{enumerate}[label=\Alph*.]
    \item An embedded, special-purpose computer
    \item A general-purpose mainframe
    \item A supercomputer by purpose
    \item An analog microcomputer by form
  \end{enumerate}
\item \textit{[Original]} A desktop capable of running office, browsing, and
  media programs is best described by purpose as:
  \begin{enumerate}[label=\Alph*.]
    \item Special-purpose
    \item General-purpose
    \item Embedded
    \item Analog
  \end{enumerate}
\item \textit{[Original]} Which pairing is matched correctly?
  \begin{enumerate}[label=\Alph*.]
    \item Analog -- discrete binary values only
    \item Digital -- continuously varying representation only
    \item Hybrid -- combines analog and digital characteristics
    \item Mainframe -- a signal representation
  \end{enumerate}
\item \textit{[Original]} Evaluate the statements:
  \begin{enumerate}[label=\Roman*.]
    \item A supercomputer is associated with high-performance computation.
    \item A mainframe is associated with high-volume transaction workloads.
    \item Physical size alone is sufficient to distinguish the two.
  \end{enumerate}
  \begin{enumerate}[label=\Alph*.]
    \item I only
    \item I and II only
    \item II and III only
    \item I, II, and III
  \end{enumerate}
\item \textit{[Original]} A laptop receives a photograph originally captured
  from a continuously varying physical scene. Which classification is
  appropriate for the laptop itself?
  \begin{enumerate}[label=\Alph*.]
    \item Analog, because every photographed scene is continuous
    \item Digital, because it represents and processes stored data discretely
    \item Mainframe, because it stores images
    \item Hybrid, solely because it displays colour
  \end{enumerate}
\item \textit{[Original]} Which clue most strongly indicates a mainframe
  workload?
  \begin{enumerate}[label=\Alph*.]
    \item One extremely large weather simulation
    \item Many concurrent users completing transactions
    \item A controller embedded in a washing machine
    \item A single user's laptop form factor
  \end{enumerate}
\item \textit{[Original]} A large weather simulation emphasizes massive
  computation rather than transaction volume. Which class is the best fit?
  \begin{enumerate}[label=\Alph*.]
    \item Mainframe
    \item Supercomputer
    \item Minicomputer
    \item Embedded computer
  \end{enumerate}
\item \textit{[Original]} A digital medical monitor uses a sensor to
  measure a continuous physical signal and then processes the reading
  digitally. Why can ``hybrid'' be appropriate?
  \begin{enumerate}[label=\Alph*.]
    \item The complete system combines analog measurement with digital
      processing.
    \item Every digital screen is a mainframe.
    \item The sensor converts the monitor into a supercomputer.
    \item Hybrid means a computer has two users.
  \end{enumerate}
\item \textit{[Original]} Which term describes a computer built into a larger
  product to support that product's function?
  \begin{enumerate}[label=\Alph*.]
    \item Embedded computer
    \item Supercomputer
    \item Analog computer
    \item Minicomputer
  \end{enumerate}
\item \textit{[Original]} A prompt asks whether a machine is analog or
  digital. Which clue should be examined first?
  \begin{enumerate}[label=\Alph*.]
    \item The machine's cabinet size
    \item The way it represents or processes signals
    \item Whether it is a laptop or desktop
    \item The number of people using the room
  \end{enumerate}

\item \textit{[JKSSB FAA-pattern]} Which class of computer is used for weather forecasting and large scientific simulations?
  \begin{enumerate}[label=\Alph*.]
    \item Supercomputer
    \item Palmtop
    \item Embedded controller
    \item Digital diary
  \end{enumerate}
\item \textit{[JKPSI-pattern]} A laptop is best classified as which type of computer?
  \begin{enumerate}[label=\Alph*.]
    \item A portable microcomputer
    \item A mainframe computer
    \item A supercomputer
    \item An analog computer
  \end{enumerate}
\item \textit{[SSC CGL-pattern]} Which of the following is an Indian-built supercomputer series?
  \begin{enumerate}[label=\Alph*.]
    \item PARAM
    \item ENIAC
    \item Abacus
    \item Pager
  \end{enumerate}
\item \textit{[JKSSB-pattern]} \textbf{(MSQ: select ALL correct options.)} Which of the following are microcomputers?
  \begin{enumerate}[label=\Alph*.]
    \item Desktop computer
    \item Laptop computer
    \item Tablet computer
    \item Mainframe computer
  \end{enumerate}
\item \textit{[SSC CGL-pattern]} \textbf{(MSQ: select ALL correct options.)} Which pairs are correctly matched?
  \begin{enumerate}[label=\Alph*.]
    \item Analog computer -- works on continuous data
    \item Digital computer -- works on discrete data
    \item Hybrid computer -- combines analog and digital features
    \item Supercomputer -- meant for routine office typing
  \end{enumerate}
\item \textit{[Original]} A computer that represents values using continuously varying physical quantities is:
  \begin{enumerate}[label=\Alph*.]
\item Digital
   \item Analog
   \item Hybrid
   \item Microcomputer
  \end{enumerate}

\item \textit{[Original]} A system combining analog measurement with digital control is a:
  \begin{enumerate}[label=\Alph*.]
\item Hybrid computer
   \item Mainframe
   \item Minicomputer
   \item General-purpose digital computer
  \end{enumerate}

\item \textit{[Original]} Which class is designed to support very large numbers of users and high-volume transaction processing?
  \begin{enumerate}[label=\Alph*.]
\item Mainframe
   \item Tablet computer
   \item Embedded controller
   \item Analog computer
  \end{enumerate}

\item \textit{[Original]} Which classification is based primarily on the kind of data represented rather than physical size?
  \begin{enumerate}[label=\Alph*.]
\item Analog, digital, hybrid
   \item Micro, mini, mainframe
   \item Laptop, desktop, tablet
   \item Client, server, workstation
  \end{enumerate}

\item \textit{[Original, MSQ drill]} \textbf{(MSQ: select ALL correct options.)} Which pairings are valid?
  \begin{enumerate}[label=\Alph*.]
\item Analog computer --- continuous quantities
   \item Digital computer --- discrete values
   \item Hybrid computer --- combination of analog and digital features
   \item Mainframe --- a type of analog representation
  \end{enumerate}

\end{enumerate}
\lessonheading{03}{Bits, Bytes \& Units}
\begin{enumerate}[label=\textbf{Q\arabic*.}, leftmargin=*]
\item \textit{[PYQ-adapted: JKSSB Website Operator 2026 Q72; official
provisional answer key; paraphrased]} In the usual hierarchy of digital data
units, what is the smallest unit?
  \begin{enumerate}[label=\Alph*.]
    \item Block
    \item Bit
    \item Byte
    \item Word
  \end{enumerate}
\item \textit{[Original]} How many bits are in one nibble?
  \begin{enumerate}[label=\Alph*.]
    \item 2
    \item 4
    \item 8
    \item 16
  \end{enumerate}
\item \textit{[Original]} How many bits are in one byte?
  \begin{enumerate}[label=\Alph*.]
    \item 4
    \item 6
    \item 8
    \item 10
  \end{enumerate}
\item \textit{[Original]} One byte is equal to:
  \begin{enumerate}[label=\Alph*.]
    \item Two nibbles
    \item Two bits
    \item Four bytes
    \item Eight nibbles
  \end{enumerate}
\item \textit{[Original]} What does lowercase \texttt{b} ordinarily denote in
  data-rate notation?
  \begin{enumerate}[label=\Alph*.]
    \item Byte
    \item Bit
    \item Binary file
    \item Buffer
  \end{enumerate}
\item \textit{[Original]} Which notation is a rate rather than a storage
  amount?
  \begin{enumerate}[label=\Alph*.]
    \item 20 MB
    \item 20 GB
    \item 20 Mbps
    \item 20 bytes
  \end{enumerate}
\item \textit{[Original]} Under the lesson's 1024 convention, 1 KB is:
  \begin{enumerate}[label=\Alph*.]
    \item 100 bytes
    \item 1000 bytes
    \item 1024 bytes
    \item 2048 bytes
  \end{enumerate}
\item \textit{[Original]} How many bytes are in 3 KB under the 1024
  convention?
  \begin{enumerate}[label=\Alph*.]
    \item 3072
    \item 3000
    \item 24
    \item 1027
  \end{enumerate}
\item \textit{[Original]} Under the same convention, 2 MB equals:
  \begin{enumerate}[label=\Alph*.]
    \item 512 KB
    \item 1024 KB
    \item 2048 KB
    \item 2048 bytes
  \end{enumerate}
\item \textit{[Original]} At matching decimal prefixes, 8 Mbps is equal to
  which ideal rate before protocol overhead?
  \begin{enumerate}[label=\Alph*.]
    \item 1 MB/s
    \item 8 MB/s
    \item 64 MB/s
    \item 0.125 MB/s
  \end{enumerate}
\item \textit{[Original]} A 64-bit processor word contains how many bytes?
  \begin{enumerate}[label=\Alph*.]
    \item 4
    \item 8
    \item 16
    \item 64
  \end{enumerate}
\item \textit{[Original]} Which statement about word size is correct?
  \begin{enumerate}[label=\Alph*.]
    \item It is always exactly 8 bits.
    \item It is the same as a nibble on every processor.
    \item It depends on the processor architecture.
    \item It always means one character.
  \end{enumerate}
\item \textit{[Original]} Which statement is \textbf{NOT} correct?
  \begin{enumerate}[label=\Alph*.]
    \item Lowercase \texttt{b} commonly denotes bit.
    \item Uppercase \texttt{B} commonly denotes byte.
    \item A byte contains 8 bits.
    \item A word has one universal size on all processors.
  \end{enumerate}
\item \textit{[Original]} A file is 4 MB. Which interpretation is correct?
  \begin{enumerate}[label=\Alph*.]
    \item It is a storage amount measured in bytes.
    \item It is a data rate measured per second.
    \item It is 4 megabits per second.
    \item It specifies the processor word size.
  \end{enumerate}
\item \textit{[Original]} Evaluate the statements:
  \begin{enumerate}[label=\Roman*.]
    \item A nibble is 4 bits.
    \item A byte is 8 bits.
    \item A byte is two nibbles.
  \end{enumerate}
  \begin{enumerate}[label=\Alph*.]
    \item I only
    \item I and II only
    \item II and III only
    \item I, II, and III
  \end{enumerate}
\item \textit{[Original]} A connection is rated at 24 Mbps. Ignoring
  overhead, what is its ideal byte rate using matching decimal prefixes?
  \begin{enumerate}[label=\Alph*.]
    \item 3 MB/s
    \item 24 MB/s
    \item 192 MB/s
    \item 0.3 MB/s
  \end{enumerate}
\item \textit{[Original]} Which statement correctly distinguishes capacity
  from transfer rate?
  \begin{enumerate}[label=\Alph*.]
    \item Capacity is how much data can be held; rate is how much moves per
      time.
    \item Capacity is always measured per second; rate never includes time.
    \item Both terms mean processor word size.
    \item A larger capacity guarantees a faster transfer rate.
  \end{enumerate}
\item \textit{[Original]} The 1024 ladder used in this lesson gives
  $1$ GB as:
  \begin{enumerate}[label=\Alph*.]
    \item 1024 MB
    \item 1000 MB
    \item 8 MB
    \item 1024 bytes
  \end{enumerate}
\item \textit{[Original]} Why must a question's stated unit convention be
  checked before a capacity conversion?
  \begin{enumerate}[label=\Alph*.]
    \item Decimal prefixes use powers of 1000 while the traditional binary
      ladder uses powers of 1024.
    \item Bytes have different sizes in every file.
    \item A bit is sometimes 4 bytes.
    \item Capacity always depends on network speed.
  \end{enumerate}
\item \textit{[Original]} A character is stored using UTF-8. Which assumption
  is unsafe?
  \begin{enumerate}[label=\Alph*.]
    \item The encoded size may vary.
    \item Every character must occupy exactly one byte.
    \item Encoding represents text as bytes.
    \item Character storage depends on the encoding.
  \end{enumerate}

\item \textit{[JKSSB FAA-pattern]} How many bits make one byte?
  \begin{enumerate}[label=\Alph*.]
    \item 8
    \item 4
    \item 16
    \item 1024
  \end{enumerate}
\item \textit{[JKPSI-pattern]} One kilobyte (KB) equals how many bytes?
  \begin{enumerate}[label=\Alph*.]
    \item 1024
    \item 100
    \item 8
    \item 1000 bits
  \end{enumerate}
\item \textit{[SSC CGL-pattern]} Which of the following storage units is the largest?
  \begin{enumerate}[label=\Alph*.]
    \item Terabyte
    \item Gigabyte
    \item Megabyte
    \item Kilobyte
  \end{enumerate}
\item \textit{[JKSSB-pattern]} \textbf{(MSQ: select ALL correct options.)} Which of the following conversions are correct?
  \begin{enumerate}[label=\Alph*.]
    \item 4 bits = 1 nibble
    \item 8 bits = 1 byte
    \item 1024 bytes = 1 KB
    \item 1000 KB = 1 MB
  \end{enumerate}
\item \textit{[SSC CGL-pattern]} \textbf{(MSQ: select ALL correct options.)} Which statements about the bit and the byte are correct?
  \begin{enumerate}[label=\Alph*.]
    \item The bit is the smallest unit of data.
    \item One byte consists of 8 bits.
    \item One nibble consists of 4 bits.
    \item A word always means exactly 8 bits on every machine.
  \end{enumerate}
\item \textit{[PYQ-adapted: JKSSB Website Operator 2026 Q72, provisional key; SRC-04/SRC-07]} What is the smallest unit of data?
  \begin{enumerate}[label=\Alph*.]
\item Bit
   \item Byte
   \item Nibble
   \item Word
  \end{enumerate}

\item \textit{[Original]} How many bits are in one byte?
  \begin{enumerate}[label=\Alph*.]
\item 4
   \item 8
   \item 16
   \item 32
  \end{enumerate}

\item \textit{[Original]} How many nibbles make one byte?
  \begin{enumerate}[label=\Alph*.]
\item One
   \item Two
   \item Four
   \item Eight
  \end{enumerate}

\item \textit{[Original]} Using the binary convention, how many bytes are in 1 KiB?
  \begin{enumerate}[label=\Alph*.]
\item 1000
   \item 1024
   \item 100
   \item 2048
  \end{enumerate}

\item \textit{[Original, MSQ drill]} \textbf{(MSQ: select ALL correct options.)} Which statements are correct?
  \begin{enumerate}[label=\Alph*.]
\item A nibble contains four bits.
   \item A byte contains eight bits.
   \item A bit can represent 0 or 1.
   \item A byte contains four nibbles.
  \end{enumerate}

\end{enumerate}
\lessonheading{49}{Computer Generations \& History}
\begin{enumerate}[label=\textbf{Q\arabic*.}, leftmargin=*]
\item \textit{[PYQ-adapted: JKSSB Website Operator 2026 Q71; official
provisional key; paraphrased]} Which generation is conventionally linked to
the introduction of integrated circuits (ICs)?
  \begin{enumerate}[label=\Alph*.]
    \item First generation
    \item Second generation
    \item Third generation
    \item Fourth generation
  \end{enumerate}
\item \textit{[Original]} Vacuum tubes are conventionally associated with
  which generation?
  \begin{enumerate}[label=\Alph*.]
    \item First
    \item Second
    \item Third
    \item Fourth
  \end{enumerate}
\item \textit{[Original]} Transistors are the principal technology cue for
  the:
  \begin{enumerate}[label=\Alph*.]
    \item First generation
    \item Second generation
    \item Third generation
    \item Fourth generation
  \end{enumerate}
\item \textit{[Original]} Microprocessors are conventionally associated
  with the:
  \begin{enumerate}[label=\Alph*.]
    \item First generation
    \item Second generation
    \item Third generation
    \item Fourth generation
  \end{enumerate}
\item \textit{[Original]} Which sequence correctly lists the main
  technology associations for the first four generations?
  \begin{enumerate}[label=\Alph*.]
    \item Vacuum tubes, transistors, ICs, microprocessors
    \item Transistors, vacuum tubes, microprocessors, ICs
    \item ICs, vacuum tubes, transistors, microprocessors
    \item Microprocessors, ICs, transistors, vacuum tubes
  \end{enumerate}
\item \textit{[Original]} Which statement is \textbf{NOT} correct?
  \begin{enumerate}[label=\Alph*.]
    \item ICs are associated with the third generation.
    \item Microprocessors are associated with the fourth generation.
    \item Transistors are associated with the second generation.
    \item ICs are the conventional first-generation association.
  \end{enumerate}
\item \textit{[Original]} A question explicitly asks for the conventional
  fifth-generation association. Which answer is the best match?
  \begin{enumerate}[label=\Alph*.]
    \item Vacuum tubes and machine language only
    \item Transistors only
    \item Artificial Intelligence and parallel processing
    \item Integrated circuits only
  \end{enumerate}
\item \textit{[Original]} Expand C-DAC.
  \begin{enumerate}[label=\Alph*.]
    \item Centre for Development of Advanced Computing
    \item Council for Digital Applications and Computing
    \item Central Department of Automated Communication
    \item Centre for Data Analysis and Control
  \end{enumerate}
\item \textit{[Original]} Which system is identified by C-DAC as India's
  first supercomputer?
  \begin{enumerate}[label=\Alph*.]
    \item ENIAC
    \item PARAM 8000
    \item UNIVAC I
    \item IBM PC
  \end{enumerate}
\item \textit{[Original]} PARAM 8000 was unveiled in:
  \begin{enumerate}[label=\Alph*.]
    \item 1971
    \item 1981
    \item 1991
    \item 2001
  \end{enumerate}
\item \textit{[Original]} Which organization developed PARAM 8000?
  \begin{enumerate}[label=\Alph*.]
    \item C-DAC
    \item NASA
    \item CERN
    \item IBM
  \end{enumerate}
\item \textit{[Original]} Which statement best describes the status of
  computer-generation dates?
  \begin{enumerate}[label=\Alph*.]
    \item Every textbook uses exactly the same universal date boundaries.
    \item Generation labels are conventional historical groupings; boundaries
      may vary.
    \item Generations are classified only by the monitor type.
    \item A computer cannot contain technologies from more than one period.
  \end{enumerate}
\item \textit{[Original]} Which pairing is mismatched?
  \begin{enumerate}[label=\Alph*.]
    \item First -- vacuum tubes
    \item Second -- transistors
    \item Third -- microprocessors
    \item Fourth -- microprocessors
  \end{enumerate}
\item \textit{[Original]} Evaluate the statements:
  \begin{enumerate}[label=\Roman*.]
    \item ICs are the conventional third-generation association.
    \item Microprocessors are the conventional fourth-generation association.
    \item An AI application alone establishes a universally fixed
      fifth-generation start date.
  \end{enumerate}
  \begin{enumerate}[label=\Alph*.]
    \item I only
    \item I and II only
    \item II and III only
    \item I, II, and III
  \end{enumerate}
\item \textit{[Original]} In a matching exercise, which entry should be
  paired with ``second generation''?
  \begin{enumerate}[label=\Alph*.]
    \item Vacuum tubes
    \item Transistors
    \item Integrated circuits
    \item Microprocessors
  \end{enumerate}
\item \textit{[Original]} Which is a named Indian high-performance-computing
  milestone rather than a generation label?
  \begin{enumerate}[label=\Alph*.]
    \item Fourth generation
    \item PARAM 8000
    \item Transistor generation
    \item Fifth generation
  \end{enumerate}
\item \textit{[Original]} Which statement about the PARAM 8000 fact is
  matched correctly?
  \begin{enumerate}[label=\Alph*.]
    \item PARAM 8000 is the organization; C-DAC is the system.
    \item C-DAC developed PARAM 8000, which was unveiled in 1991.
    \item PARAM 8000 is the third generation of computers.
    \item C-DAC is a type of microprocessor.
  \end{enumerate}
\item \textit{[Original]} A stem says ``introduced integrated circuits''
  but an option says fourth generation. What is the correct response under
  the conventional scheme?
  \begin{enumerate}[label=\Alph*.]
    \item Select third generation; fourth is the microprocessor association.
    \item Select second generation; it introduced transistors.
    \item Select first generation; it introduced vacuum tubes.
    \item Select fifth generation; it introduced AI.
  \end{enumerate}
\item \textit{[Original]} Why is ``AI exists today'' not enough to assign
  every modern computer to a sharply bounded fifth generation?
  \begin{enumerate}[label=\Alph*.]
    \item The fifth-generation label is a broad conventional grouping, not a
      universally fixed boundary.
    \item AI is a vacuum tube.
    \item All modern computers are analog.
    \item Generation labels describe only cabinet size.
  \end{enumerate}
\item \textit{[Original]} Which pair keeps the system, organization, and
  year in their correct roles?
  \begin{enumerate}[label=\Alph*.]
    \item C-DAC -- organization; PARAM 8000 -- system; 1991 -- unveiling year
    \item C-DAC -- system; PARAM 8000 -- year; 1991 -- organization
    \item C-DAC -- generation; PARAM 8000 -- IC; 1991 -- processor
    \item C-DAC -- year; PARAM 8000 -- organization; 1991 -- system
  \end{enumerate}

\item \textit{[JKSSB FAA-pattern]} First-generation computers used which technology?
  \begin{enumerate}[label=\Alph*.]
    \item Vacuum tubes
    \item Microprocessors
    \item Transistors only
    \item Touchscreens
  \end{enumerate}
\item \textit{[JKPSI-pattern]} ENIAC belongs to which computer generation?
  \begin{enumerate}[label=\Alph*.]
    \item First generation
    \item Third generation
    \item Fifth generation
    \item Zero generation printer
  \end{enumerate}
\item \textit{[SSC CGL-pattern]} Which generation introduced microprocessors?
  \begin{enumerate}[label=\Alph*.]
    \item Fourth generation
    \item First generation
    \item Second generation
    \item Zeroth generation
  \end{enumerate}
\item \textit{[JKSSB-pattern]} \textbf{(MSQ: select ALL correct options.)} Which technologies mark the early computer generations?
  \begin{enumerate}[label=\Alph*.]
    \item Vacuum tubes (first generation)
    \item Transistors (second generation)
    \item Integrated circuits (third generation)
    \item Steam engines (first generation)
  \end{enumerate}
\item \textit{[SSC CGL-pattern]} \textbf{(MSQ: select ALL correct options.)} Which statements about computer generations are correct?
  \begin{enumerate}[label=\Alph*.]
    \item First-generation machines used vacuum tubes.
    \item Second-generation machines used transistors.
    \item Fourth-generation machines used microprocessors.
    \item Fifth-generation goals include artificial intelligence.
  \end{enumerate}
\item \textit{[Original]} First-generation electronic computers primarily used:
  \begin{enumerate}[label=\Alph*.]
\item Vacuum tubes
   \item Integrated circuits
   \item Microprocessors
   \item Transistors only
  \end{enumerate}

\item \textit{[Original]} Integrated circuits are most closely associated with which generation?
  \begin{enumerate}[label=\Alph*.]
\item First
   \item Second
   \item Third
   \item Fifth
  \end{enumerate}

\item \textit{[Original]} The microprocessor is most closely associated with which computer generation?
  \begin{enumerate}[label=\Alph*.]
\item First
   \item Second
   \item Third
   \item Fourth
  \end{enumerate}

\item \textit{[Original]} Which technology replaced vacuum tubes in second-generation computers?
  \begin{enumerate}[label=\Alph*.]
\item Transistors
   \item Integrated circuits
   \item Microprocessors
   \item Optical fibers
  \end{enumerate}

\item \textit{[Original, MSQ drill]} \textbf{(MSQ: select ALL correct options.)} Match technology to generation.
  \begin{enumerate}[label=\Alph*.]
\item Vacuum tubes --- first
   \item Transistors --- second
   \item Integrated circuits --- third
   \item Microprocessors --- fourth
  \end{enumerate}

\end{enumerate}
\lessonheading{50}{Data Representation, Number Systems \& Logic}
\begin{enumerate}[label=\textbf{Q\arabic*.}, leftmargin=*]
\item \textit{[PYQ-adapted: JKSSB Computer Instructor/Operator 2025 Q1;
official paper; paraphrased]} What is the decimal ASCII code for the dollar
sign (\$)?
  \begin{enumerate}[label=\Alph*.]
    \item 32
    \item 30
    \item 34
    \item 36
  \end{enumerate}
\item \textit{[PYQ-adapted: JKSSB Computer Instructor/Operator 2025 Q2;
official paper; paraphrased]} Convert the binary number
$11010101.01_2$ to decimal.
  \begin{enumerate}[label=\Alph*.]
    \item $200.25_{10}$
    \item $198.25_{10}$
    \item $213.25_{10}$
    \item $209.25_{10}$
  \end{enumerate}
\item \textit{[PYQ-adapted: JKSSB Computer Instructor/Operator 2025 Q3;
official paper; paraphrased]} Which statement about UTF-8 is correct?
  \begin{enumerate}[label=\Alph*.]
    \item It is fixed-length and uses exactly 8 bits for every character.
    \item It is variable-length and compatible with ASCII in the ASCII range.
    \item It assigns Unicode code points only to Latin letters.
    \item It stores every character using 8 bytes.
  \end{enumerate}
\item \textit{[PYQ-adapted: JKSSB Computer Instructor/Operator 2025 Q4;
official paper; paraphrased]} What is the binary representation of
$278_{10}$?
  \begin{enumerate}[label=\Alph*.]
    \item $100010110_2$
    \item $101010110_2$
    \item $100110110_2$
    \item $110010110_2$
  \end{enumerate}
\item \textit{[PYQ-adapted: JKSSB Website Operator 2026 Q73; official
provisional key; paraphrased]} Which binary numeral has the value
$10_{10}$?
  \begin{enumerate}[label=\Alph*.]
    \item $1010_2$
    \item $1001_2$
    \item $1100_2$
    \item $1110_2$
  \end{enumerate}
\item \textit{[Original]} What is the base of the octal number system?
  \begin{enumerate}[label=\Alph*.]
    \item 2
    \item 8
    \item 10
    \item 16
  \end{enumerate}
\item \textit{[Original]} Which digit is not permitted in an octal numeral?
  \begin{enumerate}[label=\Alph*.]
    \item 0
    \item 5
    \item 7
    \item 8
  \end{enumerate}
\item \textit{[Original]} What is the decimal value of $2F_{16}$?
  \begin{enumerate}[label=\Alph*.]
    \item 31
    \item 32
    \item 47
    \item 63
  \end{enumerate}
\item \textit{[Original]} Convert $1111111_2$ to octal.
  \begin{enumerate}[label=\Alph*.]
    \item $177_8$
    \item $177_2$
    \item $77_8$
    \item $177_{16}$
  \end{enumerate}
\item \textit{[Original]} In BCD, how is the decimal number 25 represented?
  \begin{enumerate}[label=\Alph*.]
    \item $11001_2$
    \item $0010\ 0101$
    \item $0101\ 0010$
    \item $11101_2$
  \end{enumerate}
\item \textit{[Original]} Which statement is \textbf{NOT} correct?
  \begin{enumerate}[label=\Alph*.]
    \item Original ASCII uses 7-bit codes.
    \item Unicode assigns code points across many writing systems.
    \item UTF-8 is variable-length.
    \item BCD and Unicode are interchangeable names for character encoding.
  \end{enumerate}
\item \textit{[Original]} Which option describes the relationship between
  Unicode and UTF-8 correctly?
  \begin{enumerate}[label=\Alph*.]
    \item Unicode assigns code points; UTF-8 encodes them using a
      variable-length byte sequence.
    \item UTF-8 is another name for BCD.
    \item Unicode is a 4-bit code for each decimal digit.
    \item UTF-8 is always fixed at 8 bytes per character.
  \end{enumerate}
\item \textit{[Original]} In common 24-bit RGB, the three 8-bit channels
  represent:
  \begin{enumerate}[label=\Alph*.]
    \item Red, green, and blue
    \item Red, grey, and black
    \item Row, grid, and bit
    \item Read, generate, and store
  \end{enumerate}
\item \textit{[Original]} For inputs $A=1$ and $B=0$, what is the output of
  XOR?
  \begin{enumerate}[label=\Alph*.]
    \item 0
    \item 1
    \item 2
    \item It depends on a clock
  \end{enumerate}
\item \textit{[Original]} Which gate outputs 1 only when both inputs are 1?
  \begin{enumerate}[label=\Alph*.]
    \item OR
    \item XOR
    \item AND
    \item NOT
  \end{enumerate}
\item \textit{[Original]} For $A=1$ and $B=1$, the NAND output is:
  \begin{enumerate}[label=\Alph*.]
    \item 0
    \item 1
    \item 2
    \item Undefined
  \end{enumerate}
\item \textit{[Original]} A NOT gate has:
  \begin{enumerate}[label=\Alph*.]
    \item One input and the opposite output
    \item Two inputs and output 1 only when both are 1
    \item Two inputs and output 1 when they differ
    \item Three inputs representing RGB
  \end{enumerate}
\item \textit{[Original]} Evaluate the statements about a common 24-bit RGB
  representation:
  \begin{enumerate}[label=\Roman*.]
    \item It uses red, green, and blue channels.
    \item Three 8-bit channels total 24 bits per pixel.
    \item Every image format must use exactly 24-bit RGB.
  \end{enumerate}
  \begin{enumerate}[label=\Alph*.]
    \item I only
    \item I and II only
    \item II and III only
    \item I, II, and III
  \end{enumerate}
\item \textit{[Original]} What is $1011_2$ in decimal?
  \begin{enumerate}[label=\Alph*.]
    \item 9
    \item 10
    \item 11
    \item 13
  \end{enumerate}
\item \textit{[Original]} When converting a binary numeral to hexadecimal by
  grouping bits, how many bits are placed in each complete group?
  \begin{enumerate}[label=\Alph*.]
    \item 2
    \item 3
    \item 4
    \item 8
  \end{enumerate}
\item \textit{[Original]} For inputs $A=0$ and $B=0$, what are the outputs
  of AND, OR, XOR, and NAND, in that order?
  \begin{enumerate}[label=\Alph*.]
    \item $0,0,0,1$
    \item $1,0,0,0$
    \item $0,1,1,1$
    \item $1,1,0,1$
  \end{enumerate}

\item \textit{[JKSSB FAA-pattern]} The binary number 1010 equals which decimal value?
  \begin{enumerate}[label=\Alph*.]
    \item 10
    \item 8
    \item 12
    \item 101
  \end{enumerate}
\item \textit{[JKPSI-pattern]} How many distinct values can 3 bits represent?
  \begin{enumerate}[label=\Alph*.]
    \item 8
    \item 3
    \item 6
    \item 9
  \end{enumerate}
\item \textit{[SSC CGL-pattern]} Which logic gate gives output 1 only when all inputs are 1?
  \begin{enumerate}[label=\Alph*.]
    \item AND
    \item OR
    \item NOT for two inputs
    \item Printer gate
  \end{enumerate}
\item \textit{[JKSSB-pattern]} \textbf{(MSQ: select ALL correct options.)} Which binary-decimal conversions are correct?
  \begin{enumerate}[label=\Alph*.]
    \item 101 = 5
    \item 1000 = 8
    \item 111 = 7
    \item 101 = 6
  \end{enumerate}
\item \textit{[SSC CGL-pattern]} \textbf{(MSQ: select ALL correct options.)} Which statements about logic gates are correct?
  \begin{enumerate}[label=\Alph*.]
    \item AND outputs 1 only if all inputs are 1.
    \item OR outputs 1 if any input is 1.
    \item NOT inverts its single input.
    \item AND outputs 1 when all inputs are 0.
  \end{enumerate}
\item \textit{[Original]} What is the decimal value of binary 1010?
  \begin{enumerate}[label=\Alph*.]
\item 8
   \item 10
   \item 12
   \item 14
  \end{enumerate}

\item \textit{[Original]} In hexadecimal notation, the digit A represents decimal:
  \begin{enumerate}[label=\Alph*.]
\item 8
   \item 9
   \item 10
   \item 11
  \end{enumerate}

\item \textit{[Original]} BCD represents each decimal digit using:
  \begin{enumerate}[label=\Alph*.]
\item Four bits
   \item One byte
   \item Three bits
   \item Sixteen bits
  \end{enumerate}

\item \textit{[Original]} Which logic gate outputs 1 only when its two inputs differ?
  \begin{enumerate}[label=\Alph*.]
\item XOR
   \item AND
   \item OR
   \item NOT
  \end{enumerate}

\item \textit{[Original, MSQ drill]} \textbf{(MSQ: select ALL correct options.)} Which statements are correct?
  \begin{enumerate}[label=\Alph*.]
\item One hexadecimal digit represents four binary bits.
   \item A byte contains eight bits.
   \item BCD encodes each decimal digit in four bits.
   \item Decimal 10 is binary 1110.
  \end{enumerate}

\end{enumerate}
\answersection{1}
\answersubheading{01}{What Is a Computer?}
\begin{enumerate}[label=\textbf{Q\arabic*.}, leftmargin=*]
\item \textbf{A.} This includes input, instruction-based processing, output,
  and possible storage. B incorrectly gives the computer independent
  truth-judgement; C reduces it to display and denies storage; D denies
  processing.
\item \textbf{B.} The unprocessed marks are data. A describes organized
  results; C is a stage/result role rather than the raw marks themselves; D
  is the retention function.
\item \textbf{C.} Applying the formula transforms the marks, so it is
  processing. A is input; B is output; D is storage.
\item \textbf{B.} P means Processing. Printing, programming, and peripheral
  are not the IPO stage represented by P.
\item \textbf{B.} IPO is Input, Processing, Output. A reverses the stages; C
  puts output before input; D incorrectly inserts storage at the start.
\item \textbf{A.} Adding Storage to IPO gives IPOS. The other strings do not
  preserve the standard initials in order.
\item \textbf{D.} Saving retains data for later use, which is storage. A
  supplies input; B transforms data; C presents a result.
\item \textbf{C.} Speed cannot establish whether input is true. A, B, and D
  correctly state that results depend on suitable input and instructions.
\item \textbf{A.} GIGO means poor input can lead to poor output. B reverses
  the warning; C falsely promises automatic repair; D incorrectly claims
  numerical data cannot be processed.
\item \textbf{B.} Speed is the ability to perform many operations quickly. A
  is repetition without tiring; C is performing varied tasks; D is retention.
\item \textbf{A.} Diligence describes repeated work without tiring in the
  human sense. B concerns conditional consistency; C is range of tasks; D is
  a data-entry stage.
\item \textbf{B.} Versatility is the ability to perform different tasks with
  different programs. A is tireless repetition; C is retention; D is a
  result/stage.
\item \textbf{A.} Automation means continuing a programmed operation once
  started and supplied with required inputs. B is raw facts; C is a result;
  D describes input-output quality.
\item \textbf{D.} All three statements are correct: data are raw facts or
  symbols, information is purposefully processed data, and a value's role
  depends on the task. A, B, and C each omit at least one true statement.
\item \textbf{A.} Records and appointment scheduling are common healthcare
  applications. B overclaims that software guarantees clinical judgement; C
  confuses display and storage; D incorrectly says processing makes input
  truthful.
\item \textbf{D.} Fast calculation is a capability, not a limitation. A, B,
  and C describe genuine limitations: no independent truth test, need for
  contextual interpretation, and dependence on instructions.
\item \textbf{A.} In a new task the previously displayed average enters the
  system as input. B is transformation; C is retention; D is a characteristic.
\item \textbf{B.} Entered amounts are input; the formula processes them. A
  reverses both roles; C mislabels storage/output; D assigns both the wrong
  stages.
\item \textbf{C.} Results depend on the data and instructions used. A
  confuses speed with fairness; B claims validity despite missing input; D
  wrongly limits output to a monitor.
\item \textbf{A.} The formula can calculate correctly from an incorrect
  entry; the error is at input. B confuses monitor output and input; C
  falsely says storage changes values; D treats speed as validation.

\item \textbf{A.} Data are unprocessed facts while information is processed data with meaning. B reverses both roles; C denies the distinction; D falsely claims information cannot be stored.
\item \textbf{A.} Diligence means tireless repetition. B wrongly gives the machine emotions; C contradicts repetition; D wrongly links accuracy to slowing down.
\item \textbf{D.} Computers have no intuition of their own. Speed (A), accuracy (B), and diligence (C) are standard characteristics.
\item \textbf{A, B, C.} Input, processing, and output are basic functions. Feeling emotions (D) is not a computer function.
\item \textbf{A, B, D.} Input precedes processing, output depends on input and instructions, and processing yields information. C is false because IPO has no S; storage appears only in IPOS.
\item \textbf{C.} Output presents the processed result. Input supplies data, processing transforms it, and storage retains it.

\item \textbf{A.} Diligence means sustained work without fatigue. Versatility is range of uses, accuracy concerns correct results, and automation is operation without continuous manual control.

\item \textbf{B.} Raw facts are data. Information is data organized or processed into a meaningful form; knowledge and output are not synonyms for raw entries.

\item \textbf{C.} A computer follows instructions and does not set its own goals. The other options describe speed, storage, and diligence.

\item \textbf{A, B, C.} Input, processing, and output have distinct roles. Storage retains data or results; it is not the same as input.

\end{enumerate}
\answersubheading{02}{Types of Computers}
\begin{enumerate}[label=\textbf{Q\arabic*.}, leftmargin=*]
\item \textbf{A.} Mainframes emphasize reliable high-volume transactions
  and concurrent users; supercomputers emphasize very large computations.
  B confuses workload class with signal/embedding; C confuses class with
  personal form/use; D wrongly says physical dimensions alone define the
  classes.
\item \textbf{B.} Analog systems represent continuously varying quantities.
  A is discrete representation; C combines analog and digital; D names a
  transaction workload class.
\item \textbf{B.} A laptop commonly represents data digitally. A is
  continuous representation; C is not implied by ordinary digital storage;
  D is a historical size class.
\item \textbf{C.} The arrangement combines continuous measurement with
  digital processing. A ignores the digital stages; B and D describe
  unrelated workload classes.
\item \textbf{A.} Microcomputer is the traditional personal-use class. B is
  a transaction class; C is a historical mid-range multi-user class; D is
  a high-performance computational class.
\item \textbf{B.} Minicomputer is the historical term for a mid-range,
  often multi-user class. A denotes very high-performance computation; C
  denotes personal-use systems; D is signal-based.
\item \textbf{A.} Mainframes are associated with reliable high-volume
  input/output and many transactions/users. B emphasizes large calculations;
  C is built into a product; D handles continuous signals.
\item \textbf{C.} Supercomputers target very large high-performance
  calculations. A is personal-use; B emphasizes transactions; D is a
  historical mid-range category.
\item \textbf{C.} Workload, not cabinet size alone, distinguishes mainframes
  and supercomputers. A, B, and D are correct statements about overlapping
  classification axes.
\item \textbf{A.} A smartphone is digital by representation, a microcomputer
  by personal-use class, and mobile by form. B falsely says analog; C
  confuses portability with supercomputing; D assigns signal type to
  minicomputer.
\item \textbf{A.} The appliance controller is built into a larger device and
  handles a limited task. B is a multi-user transaction class; C is
  high-performance computation; D mixes signal and size labels.
\item \textbf{B.} A general-purpose computer runs many kinds of programs. A
  is a narrower task; C means integrated into a larger product; D is
  signal representation.
\item \textbf{C.} A hybrid system combines analog and digital characteristics.
  A and B misstate the signal classes; D is a workload class, not a signal
  type.
\item \textbf{B.} I and II are correct; III is false because physical size
  alone is not the classification test. A omits II; C incorrectly includes
  III; D also includes III.
\item \textbf{B.} The laptop processes stored data as discrete digital
  values. A confuses the photographed scene with the computer's
  representation; C is an unrelated transaction class; D infers hybrid
  solely from a colour display.
\item \textbf{B.} Many concurrent users completing transactions is the
  mainframe clue. A indicates supercomputing; C indicates embedded
  special-purpose use; D identifies form factor only.
\item \textbf{B.} A very large weather simulation is a high-performance
  computational workload. A better fits transaction volume; C is historical
  mid-range; D is built into a larger product.
\item \textbf{A.} Continuous measurement plus digital processing combines
  analog and digital stages. B confuses a screen with a transaction system;
  C assigns an unrelated workload; D misdefines hybrid.
\item \textbf{A.} Embedded means built into a larger device. B is
  high-performance computation; C is signal-based; D is a historical
  mid-range class.
\item \textbf{B.} Analog/digital is determined by signal representation.
  A concerns workload class; C concerns form factor; D does not define the
  computer's signal handling.

\item \textbf{A.} Supercomputers handle massive scientific workloads. Palmtops (B), embedded controllers (C), and digital diaries (D) lack that capacity.
\item \textbf{A.} A laptop is a portable personal (micro) computer. It is not a mainframe (B), a supercomputer (C), or analog (D).
\item \textbf{A.} PARAM is the Indian supercomputer series. ENIAC (B) was an early American machine; C and D are not supercomputers.
\item \textbf{A, B, C.} Desktops, laptops, and tablets are microcomputers. A mainframe (D) is a larger multi-user class.
\item \textbf{A, B, C.} The three classifications are matched correctly. D is false: supercomputers serve heavy computation, not routine typing.
\item \textbf{B.} Analog computers represent continuous quantities. Digital computers use discrete values; hybrid systems combine analog and digital processing; microcomputer classifies size/use.

\item \textbf{A.} A hybrid computer combines analog and digital characteristics. The other options identify size/use categories or a digital-only class.

\item \textbf{A.} Mainframes serve many concurrent users and transaction workloads. Tablets are personal devices, embedded controllers serve dedicated functions, and analog describes data representation.

\item \textbf{A.} Analog/digital/hybrid describes representation. The remaining sets classify by size/form factor or role.

\item \textbf{A, B, C.} The first three match data-representation categories. Mainframe describes a computer class by scale and workload, not representation.

\end{enumerate}
\answersubheading{03}{Bits, Bytes \& Units}
\begin{enumerate}[label=\textbf{Q\arabic*.}, leftmargin=*]
\item \textbf{B.} A bit is the smallest unit in the conventional hierarchy.
  A is not a basic unit in this list; C is eight bits; D is an
  architecture-dependent processor unit.
\item \textbf{B.} One nibble is four bits. A is two bits; C is one byte; D is
  larger than a nibble.
\item \textbf{C.} One byte contains eight bits. A is one nibble; B and D do
  not match the standard byte definition.
\item \textbf{A.} Two four-bit nibbles make one eight-bit byte. B has too
  few bits; C reverses the relationship; D would make 32 bits.
\item \textbf{B.} Lowercase \texttt{b} denotes bit. A byte uses uppercase
  \texttt{B}; C and D are not meanings of the standard rate symbol.
\item \textbf{C.} Mbps means megabits per second, a rate. MB and GB are
  storage amounts; bytes alone name an amount without a time denominator.
\item \textbf{C.} This lesson uses the traditional 1024 convention. A, B,
  and D are not equal to 1 KB under that convention.
\item \textbf{A.} $3\times1024=3072$ bytes. B uses 1000; C confuses bits
  with the conversion; D adds 3 to 1024 instead of multiplying.
\item \textbf{C.} $2\times1024=2048$ KB. A is half the correct value; B is
  the value for 1 MB; D confuses KB with bytes.
\item \textbf{A.} A byte contains eight bits, so $8/8=1$ MB/s. B omits the
  bit-to-byte conversion; C multiplies instead of dividing; D applies the
  inverse factor.
\item \textbf{B.} $64/8=8$ bytes. A is 32 bits; C is 128 bits; D confuses
  bits with bytes.
\item \textbf{C.} Word size depends on the processor architecture. A and B
  impose unsupported universal sizes; D wrongly equates a word with a
  character.
\item \textbf{D.} Word size is architecture-dependent, not universal. A, B,
  and C state the standard bit/byte conventions correctly.
\item \textbf{A.} MB denotes a byte-based storage amount. B would require
  \texttt{/s}; C changes the unit to megabits per second; D is unrelated to
  file capacity.
\item \textbf{D.} All three statements are true: 4 bits per nibble, 8 bits
  per byte, and two nibbles per byte. A, B, and C each omit a true statement.
\item \textbf{A.} $24/8=3$ MB/s before overhead. B skips conversion; C
  multiplies by eight; D divides by 80 rather than 8.
\item \textbf{A.} Capacity is how much data can be held; transfer rate is
  how much moves per time. B reverses the definitions; C confuses both with
  processor word size; D falsely assumes capacity determines speed.
\item \textbf{A.} Under the 1024 ladder, $1$ GB=$1024$ MB. B is the decimal
  prefix relation; C is not a GB-to-MB conversion; D is a KB-to-byte
  relation.
\item \textbf{A.} Decimal prefixes use powers of 1000, while the traditional
  binary ladder uses powers of 1024. B, C, and D make unrelated claims about
  encoding, unit size, or network speed.
\item \textbf{B.} UTF-8 uses a variable number of bytes, so one byte per
  character is not guaranteed. A, C, and D are all consistent with
  variable-length encoding.

\item \textbf{A.} A byte is 8 bits. 4 bits (B) is a nibble; C and D are larger groupings, not a byte.
\item \textbf{A.} 1 KB is 1024 bytes in the binary convention used by these exams. B, C, and D state other quantities.
\item \textbf{A.} A terabyte exceeds a gigabyte (B), a megabyte (C), and a kilobyte (D).
\item \textbf{A, B, C.} The nibble, byte, and KB conversions are standard. D is false: 1024 KB make 1 MB.
\item \textbf{A, B, C.} The bit is smallest, a byte is 8 bits, and a nibble is 4 bits. D is false because word size varies by machine.
\item \textbf{A.} A bit is a single binary digit. A nibble is four bits, a byte is eight bits, and a word is processor-dependent. The cited answer key is provisional.

\item \textbf{B.} One byte contains eight bits. Four is a nibble; 16 and 32 bits are larger groups.

\item \textbf{B.} A nibble is four bits and a byte is eight bits, so one byte contains two nibbles.

\item \textbf{B.} 1 KiB is 1024 bytes. 1000 is the decimal kilobyte convention; 100 and 2048 are not the stated unit sizes.

\item \textbf{A, B, C.} A nibble is four bits and a byte is eight bits, so a byte has two nibbles, not four.

\end{enumerate}
\answersubheading{49}{Computer Generations \& History}
\begin{enumerate}[label=\textbf{Q\arabic*.}, leftmargin=*]
\item \textbf{C.} ICs are conventionally linked to the third generation.
  First-generation vacuum tubes, second-generation transistors, and
  fourth-generation microprocessors make A, B, and D incorrect.
\item \textbf{A.} Vacuum tubes are the first-generation cue. Transistors,
  ICs, and microprocessors are conventionally associated with the second,
  third, and fourth generations.
\item \textbf{B.} Transistors are the second-generation association. A names
  vacuum tubes; C names ICs; D names microprocessors.
\item \textbf{D.} Microprocessors are the conventional fourth-generation
  cue. A, B, and C name the technologies associated with the first three
  generations.
\item \textbf{A.} The conventional sequence is vacuum tubes, transistors,
  ICs, then microprocessors. B swaps the first two and last two; C and D
  scramble the sequence.
\item \textbf{D.} ICs are not the conventional first-generation technology;
  vacuum tubes are. A, B, and C correctly pair the third, fourth, and second
  generations respectively.
\item \textbf{C.} AI and parallel processing are common fifth-generation
  associations. A, B, and D describe the conventional first-, second-, and
  third-generation cues.
\item \textbf{A.} C-DAC stands for Centre for Development of Advanced
  Computing. B, C, and D are not expansions of the organization's name.
\item \textbf{B.} C-DAC identifies PARAM 8000 as India's first
  supercomputer. ENIAC and UNIVAC I are other historical computers; IBM PC
  is not the named Indian milestone.
\item \textbf{C.} PARAM 8000 was unveiled in 1991. 1971, 1981, and 2001 are
  not its unveiling year.
\item \textbf{A.} C-DAC developed PARAM 8000. NASA, CERN, and IBM are not the
  organization identified for this system.
\item \textbf{B.} Generation dates are conventional groupings, and sources
  may draw boundaries differently. A asserts universal dates; C and D
  reduce classification to monitor type or deny technological overlap.
\item \textbf{C.} Microprocessors are the conventional fourth-generation
  cue, not the third. A, B, and D are the correct first-, second-, and
  fourth-generation pairings.
\item \textbf{B.} Statements I and II are standard conventional
  associations; III is false because AI use does not establish one
  universally fixed start date. A omits II; C includes false III and omits
  I; D also includes false III.
\item \textbf{B.} Transistors are paired with the second generation. Vacuum
  tubes, ICs, and microprocessors correspond to the first, third, and fourth.
\item \textbf{B.} PARAM 8000 is the named system, C-DAC developed it, and
  1991 is the unveiling year. A reverses the system and organization; C
  mislabels all three; D assigns each fact the wrong role.
\item \textbf{B.} C-DAC developed PARAM 8000, unveiled in 1991. A reverses
  system and organization; C misclassifies the system as a generation; D
  assigns all three facts the wrong roles.
\item \textbf{A.} Under the conventional scheme, ICs map to the third
  generation; the fourth is associated with microprocessors. B, C, and D
  give the second-, first-, and fifth-generation association instead.
\item \textbf{A.} The fifth-generation term is a broad conventional grouping,
  not a sharply dated universal category. B mistakes AI for hardware; C
  falsely says modern computers are all analog; D misstates what
  generations classify.
\item \textbf{A.} C-DAC is the organization, PARAM 8000 the system, and
  1991 its unveiling year. B, C, and D assign those three facts to the wrong
  categories.

\item \textbf{A.} Vacuum tubes powered the first generation. B came in the fourth; C in the second; D far later.
\item \textbf{A.} ENIAC is a first-generation machine. B, C, and D place it wrongly.
\item \textbf{A.} Microprocessors define the fourth generation. B, C, and D predate them.
\item \textbf{A, B, C.} The tube, transistor, and IC markers are correct. Steam engines (D) never powered computers.
\item \textbf{A, B, C, D.} All four statements are correct: tubes, transistors, microprocessors, and AI goals in order.
\item \textbf{A.} First-generation computers used vacuum tubes. Transistors, integrated circuits, and microprocessors characterize later generations.

\item \textbf{C.} Third-generation computers used integrated circuits. First used vacuum tubes, second transistors, and fourth microprocessors.

\item \textbf{D.} Fourth-generation computers are associated with microprocessors. The other options correspond to earlier technologies.

\item \textbf{A.} Transistors replaced vacuum tubes in second-generation systems. Integrated circuits came later.

\item \textbf{A, B, C, D.} All four are the standard technology-generation associations.

\end{enumerate}
\answersubheading{50}{Data Representation, Number Systems \& Logic}
\begin{enumerate}[label=\textbf{Q\arabic*.}, leftmargin=*]
\item \textbf{D.} The ASCII decimal code for \texttt{\$} is 36. 32 is the
  space code; 30 and 34 are other control/punctuation values, not the dollar
  sign.
\item \textbf{C.} The integer bits give $128+64+16+4+1=213$; the fractional
  bits give $0/2+1/4=0.25$, so the value is $213.25$. A, B, and D do not
  match the positional sum.
\item \textbf{B.} UTF-8 is variable-length and preserves ASCII values in the
  ASCII range. A falsely says fixed-length; C wrongly limits Unicode to
  Latin; D incorrectly says every character uses eight bytes.
\item \textbf{A.} $278=256+16+4+2$, which sets bits 8, 4, 2, and 1 to give
  $100010110_2$. B, C, and D encode different decimal values.
\item \textbf{A.} $1010_2=8+2=10$. B is 9, C is 12, and D is 14.
\item \textbf{B.} Octal is base 8. A is binary, C decimal, and D hexadecimal.
\item \textbf{D.} Octal digits range from 0 through 7, so 8 is not allowed.
  0, 5, and 7 are valid octal digits.
\item \textbf{C.} $2F_{16}=2(16)+15=47$. A is 31, B is 32, and D is 63;
  none equals the place-value sum.
\item \textbf{A.} Grouping from the right gives $001\ 111\ 111_2$, or
  $1\ 7\ 7_8$. B has the wrong base label; C drops a leading group; D gives
  the wrong base label.
\item \textbf{B.} BCD encodes each decimal digit separately: 2 is 0010 and
  5 is 0101. A is the ordinary binary value 25; C reverses the digit
  nibbles; D is not the two-digit BCD form.
\item \textbf{D.} BCD is a decimal-digit code, not another name for Unicode.
  A correctly describes original ASCII; B correctly describes Unicode; C
  correctly describes UTF-8.
\item \textbf{A.} Unicode assigns code points; UTF-8 is one encoding that
  represents those code points with variable-length byte sequences. B
  conflates UTF-8 and BCD; C confuses Unicode with a four-bit decimal code;
  D falsely claims a fixed eight-byte length.
\item \textbf{A.} The channels are red, green, and blue, with 8 bits each in
  this common 24-bit representation. B, C, and D are not the RGB channel
  names.
\item \textbf{B.} XOR outputs 1 when inputs differ, and 1 and 0 differ. A is
  the AND result; C is not a Boolean output; D incorrectly makes a
  combinational gate depend on a clock.
\item \textbf{C.} AND is 1 only when both inputs are 1. OR is 1 if either
  input is 1; XOR is 1 when they differ; NOT takes one input.
\item \textbf{A.} NAND is the negation of AND; with inputs 1 and 1, AND is 1
  and NAND is 0. B is the AND result; C is not a Boolean value; D is defined
  for these inputs.
\item \textbf{A.} NOT has one input and returns its complement. B describes
  AND; C describes XOR; D confuses a logic gate with colour channels.
\item \textbf{B.} I and II are true: three 8-bit RGB channels total 24 bits.
  III is false because formats can use other channels and bit depths. A
  omits II; C includes false III and omits I; D includes false III.
\item \textbf{C.} $1011_2=8+2+1=11$. A is 9, B is 10, and D is 13.
\item \textbf{C.} One hexadecimal digit represents four binary bits.
  Two bits correspond to one octal digit; three bits form an octal group;
  eight bits make a byte, not a hex digit.
\item \textbf{A.} For inputs 0,0, AND=0, OR=0, XOR=0, and NAND=1. B assigns
  an incorrect AND and NAND; C gives incorrect OR and XOR; D gives
  incorrect AND, OR, and XOR.

\item \textbf{A.} 1010 is 8+2 = 10. B, C, and D miscalculate the positional values.
\item \textbf{A.} Three bits give $2^3 = 8$ values. B, C, and D confuse bits with value counts.
\item \textbf{A.} AND needs all inputs at 1. B needs any input; C inverts one input; D is not a gate.
\item \textbf{A, B, C.} The three conversions are exact. D misreads 101 as 6.
\item \textbf{A, B, C.} The AND, OR, and NOT behaviours are correct. D contradicts the AND rule.
\item \textbf{B.} Binary 1010 equals 8+2=10 in decimal. The other options do not match its place-value sum.

\item \textbf{C.} Hexadecimal digits A through F represent decimal 10 through 15.

\item \textbf{A.} Binary-coded decimal encodes each decimal digit separately in four bits.

\item \textbf{A.} XOR is 1 when inputs differ. AND requires both inputs 1; OR requires at least one; NOT has one input.

\item \textbf{A, B, C.} The first three are correct. Binary 1110 equals decimal 14; decimal 10 is binary 1010.

\end{enumerate}
\unitopener{2}{CPU, Hardware \& Architecture}
\lessonheading{04}{Organization of a Computer}
\begin{enumerate}[label=\textbf{Q\arabic*.}, leftmargin=*]
\item \textit{[Original]} The abbreviation ALU stands for:
  \begin{enumerate}[label=\Alph*.]
    \item Arithmetic Logic Unit
    \item Application Logic Unit
    \item Arithmetic Language Unit
    \item Advanced Logic Unit
  \end{enumerate}
\item \textit{[Original]} The ALU of a computer is responsible for:
  \begin{enumerate}[label=\Alph*.]
    \item Storing data permanently
    \item Controlling input and output devices
    \item Performing arithmetic and logical operations
    \item Supplying power to the system
  \end{enumerate}
\item \textit{[Original]} Which unit directs and coordinates the operation of
  all other units of a computer?
  \begin{enumerate}[label=\Alph*.]
    \item Arithmetic Logic Unit
    \item Memory Unit
    \item Input Unit
    \item Control Unit
  \end{enumerate}
\item \textit{[Original]} The CPU is made up of:
  \begin{enumerate}[label=\Alph*.]
    \item Input unit, output unit and memory
    \item Main memory and secondary storage
    \item ALU, control unit and registers
    \item Monitor, keyboard and mouse
  \end{enumerate}
\item \textit{[Original]} Expand CPU.
  \begin{enumerate}[label=\Alph*.]
    \item Central Processing Unit
    \item Computer Personal Unit
    \item Central Program Utility
    \item Control Processing Unit
  \end{enumerate}
\item \textit{[Original]} Which part of the CPU provides small, fast storage
  used to hold data and addresses during processing?
  \begin{enumerate}[label=\Alph*.]
    \item Registers
    \item Secondary storage
    \item Hard disk
    \item Output unit
  \end{enumerate}
\item \textit{[Original]} The address bus is used to carry:
  \begin{enumerate}[label=\Alph*.]
    \item Data between the CPU and memory
    \item The address of a memory location from the CPU to memory
    \item Control signals from the CPU to devices
    \item Electrical power to the motherboard
  \end{enumerate}
\item \textit{[Original]} Which bus is bidirectional, carrying data and
  instructions in both directions?
  \begin{enumerate}[label=\Alph*.]
    \item Address bus
    \item Data bus
    \item Control bus
    \item Power bus
  \end{enumerate}
\item \textit{[Original]} The control bus carries:
  \begin{enumerate}[label=\Alph*.]
    \item Memory addresses from the CPU
    \item Program data between the CPU and memory
    \item Control and timing signals between the CPU and other units
    \item Only input data from the keyboard
  \end{enumerate}
\item \textit{[Original]} Which statement is \textbf{NOT} correct?
  \begin{enumerate}[label=\Alph*.]
    \item The ALU performs arithmetic and logical operations.
    \item The control unit directs the sequence of operations.
    \item The CPU directly accesses secondary storage to process every
      instruction.
    \item Registers are part of the CPU.
  \end{enumerate}
\item \textit{[Original]} The CPU accesses which storage directly during
  processing?
  \begin{enumerate}[label=\Alph*.]
    \item Secondary storage
    \item Main memory (primary storage)
    \item A pen drive
    \item A hard disk
  \end{enumerate}
\item \textit{[Original]} Which of the following is a secondary storage
  device?
  \begin{enumerate}[label=\Alph*.]
    \item RAM
    \item ROM
    \item Hard disk
    \item Register
  \end{enumerate}
\item \textit{[Original]} Which statement about primary and secondary storage
  is correct?
  \begin{enumerate}[label=\Alph*.]
    \item Main memory holds the data and instructions currently being
      processed.
    \item The CPU processes data directly from secondary storage.
    \item RAM is non-volatile secondary storage.
    \item Secondary storage is faster than the CPU's registers.
  \end{enumerate}
\item \textit{[Original]} The correct order of the instruction cycle is:
  \begin{enumerate}[label=\Alph*.]
    \item Decode, fetch, execute
    \item Execute, decode, fetch
    \item Fetch, execute, decode
    \item Fetch, decode, execute
  \end{enumerate}
\item \textit{[Original]} During the fetch step, the CPU:
  \begin{enumerate}[label=\Alph*.]
    \item Performs the arithmetic operation
    \item Writes the final result to secondary storage
    \item Obtains the next instruction from main memory
    \item Displays the output on the monitor
  \end{enumerate}
\item \textit{[Original]} The stored-program (von Neumann) concept means
  that:
  \begin{enumerate}[label=\Alph*.]
    \item Only data are stored in memory
    \item Programs are stored only in secondary storage
    \item Both program instructions and data are stored in the same main
      memory
    \item Instructions are kept in the output device
  \end{enumerate}
\item \textit{[Original]} The idea that a program's instructions are stored
  in memory together with the data is associated with:
  \begin{enumerate}[label=\Alph*.]
    \item Von Neumann architecture
    \item The address bus
    \item Secondary storage
    \item The output unit
  \end{enumerate}
\item \textit{[Original]} Evaluate the statements:
  \begin{enumerate}[label=\Roman*.]
    \item The address bus carries the address from the CPU to memory.
    \item The data bus is bidirectional.
    \item The control bus carries control signals from the CPU to other
      units.
  \end{enumerate}
  \begin{enumerate}[label=\Alph*.]
    \item I only
    \item I and II only
    \item II and III only
    \item I, II and III
  \end{enumerate}
\item \textit{[Original]} Expansion slots on the motherboard are used to:
  \begin{enumerate}[label=\Alph*.]
    \item Permanently store the operating system
    \item Replace the CPU's control unit
    \item Supply the computer with electric power only
    \item Add expansion cards such as graphics or network cards
  \end{enumerate}
\item \textit{[Original]} Which pair is correctly matched?
  \begin{enumerate}[label=\Alph*.]
    \item Address bus --- bidirectional
    \item Data bus --- unidirectional
    \item Control bus --- carries control signals
    \item Data bus --- carries memory addresses
  \end{enumerate}

\item \textit{[JKSSB FAA-pattern]} Which unit is called the brain of the computer?
  \begin{enumerate}[label=\Alph*.]
    \item CPU
    \item Monitor
    \item Printer
    \item Hard disk
  \end{enumerate}
\item \textit{[JKPSI-pattern]} The ALU performs which kind of operations?
  \begin{enumerate}[label=\Alph*.]
    \item Arithmetic and logical operations
    \item Only printing operations
    \item Only storage operations
    \item Only network routing
  \end{enumerate}
\item \textit{[SSC CGL-pattern]} In a computer, CU stands for:
  \begin{enumerate}[label=\Alph*.]
    \item Control Unit
    \item Central Utility
    \item Computer Upload
    \item Cached User
  \end{enumerate}
\item \textit{[JKSSB-pattern]} \textbf{(MSQ: select ALL correct options.)} Which of the following are parts of the CPU?
  \begin{enumerate}[label=\Alph*.]
    \item Arithmetic and Logic Unit
    \item Control Unit
    \item Registers
    \item Hard disk drive
  \end{enumerate}
\item \textit{[SSC CGL-pattern]} \textbf{(MSQ: select ALL correct options.)} Which statements about computer organization are correct?
  \begin{enumerate}[label=\Alph*.]
    \item The CPU consists of the ALU, the CU, and registers.
    \item Data flows from input to processing to output.
    \item The Control Unit stores files permanently for the user.
    \item Memory holds data and instructions for processing.
  \end{enumerate}
\item \textit{[PYQ-adapted: JKSSB Website Operator 2026 Q65, provisional key; SRC-04/SRC-07]} Which unit directs the sequence of operations and controls data flow inside the CPU?
  \begin{enumerate}[label=\Alph*.]
\item Arithmetic Logic Unit
   \item Control Unit
   \item Main memory
   \item Input unit
  \end{enumerate}

\item \textit{[PYQ-adapted: JKSSB Junior Assistant 2026 Q61, final key; SRC-03/SRC-06]} Which CPU component performs arithmetic and logical operations on binary data?
  \begin{enumerate}[label=\Alph*.]
\item Control Unit
   \item Arithmetic Logic Unit
   \item Program Counter
   \item Input unit
  \end{enumerate}

\item \textit{[Original]} Which bus carries the address of a memory location selected by the CPU?
  \begin{enumerate}[label=\Alph*.]
\item Address bus
   \item Data bus
   \item Control bus
   \item Expansion bus
  \end{enumerate}

\item \textit{[Original]} Which sequence best describes a basic computer system's information flow?
  \begin{enumerate}[label=\Alph*.]
\item Output, processing, input
   \item Input, processing, output
   \item Storage, output, input
   \item Processing, output, input
  \end{enumerate}

\item \textit{[Original, MSQ drill]} \textbf{(MSQ: select ALL correct options.)} Which are functional units of a general-purpose computer?
  \begin{enumerate}[label=\Alph*.]
\item Input unit
   \item Memory unit
   \item Central Processing Unit
   \item Output unit
  \end{enumerate}

\end{enumerate}
\lessonheading{05}{CPU, ALU, Control Unit \& Registers}
\begin{enumerate}[label=\textbf{Q\arabic*.}, leftmargin=*]
\item \textit{[Original]} Expand ALU.
  \begin{enumerate}[label=\Alph*.]
    \item Arithmetic Logic Unit
    \item Advanced Logic Unit
    \item Arithmetic Language Unit
    \item Application Logic Unit
  \end{enumerate}
\item \textit{[Original]} Which CPU component performs arithmetic and
  logical operations on data?
  \begin{enumerate}[label=\Alph*.]
    \item Control Unit
    \item Register file
    \item Arithmetic Logic Unit (ALU)
    \item System bus
  \end{enumerate}
\item \textit{[Original]} The main function of the Control Unit is to:
  \begin{enumerate}[label=\Alph*.]
    \item perform addition and subtraction of numbers
    \item direct the sequence of operations and issue control signals
    \item store intermediate arithmetic results
    \item provide permanent secondary storage
  \end{enumerate}
\item \textit{[Original]} Which of the following is \textbf{NOT} a function
  of the Control Unit?
  \begin{enumerate}[label=\Alph*.]
    \item Decoding instructions
    \item Issuing control signals
    \item Directing the flow of data between units
    \item Performing the addition of two numbers
  \end{enumerate}
\item \textit{[Original]} Which register holds the address of the
  \textbf{next} instruction to be fetched?
  \begin{enumerate}[label=\Alph*.]
    \item Instruction Register (IR)
    \item Program Counter (PC)
    \item Memory Data Register (MDR)
    \item Accumulator (ACC)
  \end{enumerate}
\item \textit{[Original]} Which register holds the instruction that is
  currently being decoded and executed?
  \begin{enumerate}[label=\Alph*.]
    \item Program Counter (PC)
    \item Memory Data Register (MDR)
    \item Instruction Register (IR)
    \item Memory Address Register (MAR)
  \end{enumerate}
\item \textit{[Original]} Which register holds the address of the memory
  location to be accessed?
  \begin{enumerate}[label=\Alph*.]
    \item Memory Data Register (MDR)
    \item Memory Address Register (MAR)
    \item Instruction Register (IR)
    \item Accumulator (ACC)
  \end{enumerate}
\item \textit{[Original]} Which register holds the data being transferred to
  or from memory?
  \begin{enumerate}[label=\Alph*.]
    \item Memory Address Register (MAR)
    \item Program Counter (PC)
    \item Memory Data Register (MDR)
    \item Instruction Register (IR)
  \end{enumerate}
\item \textit{[Original]} The accumulator (ACC) is primarily used to hold:
  \begin{enumerate}[label=\Alph*.]
    \item the address of the next instruction
    \item intermediate arithmetic and logic results
    \item the last instruction fetched from memory
    \item a permanent copy of the operating system
  \end{enumerate}
\item \textit{[Original]} Which is the fastest storage in a computer system?
  \begin{enumerate}[label=\Alph*.]
    \item Registers
    \item L1 cache
    \item RAM
    \item Solid-state drive
  \end{enumerate}
\item \textit{[Original]} Where are registers located?
  \begin{enumerate}[label=\Alph*.]
    \item In cache memory
    \item In main memory (RAM)
    \item In secondary storage
    \item Inside the CPU
  \end{enumerate}
\item \textit{[Original]} Clock speed is measured in hertz (Hz). A processor
  rated at 2.4 GHz completes approximately:
  \begin{enumerate}[label=\Alph*.]
    \item 2.4 million cycles per second
    \item 2.4 billion cycles per second
    \item 2.4 thousand cycles per second
    \item 2.4 trillion cycles per second
  \end{enumerate}
\item \textit{[Original]} Which statement about clock cycles is correct?
  \begin{enumerate}[label=\Alph*.]
    \item Every machine instruction completes in exactly one clock cycle.
    \item All instructions require the same number of clock cycles.
    \item Different instructions may take different numbers of clock cycles.
    \item A clock cycle is measured in bytes.
  \end{enumerate}
\item \textit{[Original]} Evaluate the statements:
  \begin{enumerate}[label=\Roman*.]
    \item The ALU works on binary data using logic gates.
    \item The Control Unit directs the sequence of operations and data flow.
    \item The Control Unit performs the arithmetic calculations.
  \end{enumerate}
  \begin{enumerate}[label=\Alph*.]
    \item I only
    \item I and II only
    \item II and III only
    \item I, II, and III
  \end{enumerate}
\item \textit{[Original]} Which abbreviation is correctly matched with its
  full form?
  \begin{enumerate}[label=\Alph*.]
    \item PC --- Program Counter
    \item IR --- Instruction Recorder
    \item MAR --- Memory Arithmetic Register
    \item MDR --- Main Data Register
  \end{enumerate}
\item \textit{[Original]} The Program Counter holds the address of the
  \underline{\hspace{2.2cm}} instruction to be fetched.
  \begin{enumerate}[label=\Alph*.]
    \item next
    \item current
    \item previously executed
    \item data
  \end{enumerate}
\item \textit{[Original]} Registers used for general temporary storage of
  operands and intermediate results, without a fixed dedicated role, are
  called:
  \begin{enumerate}[label=\Alph*.]
    \item status registers
    \item general-purpose registers
    \item memory address registers
    \item instruction registers
  \end{enumerate}
\item \textit{[Original]} The status (flags) register holds:
  \begin{enumerate}[label=\Alph*.]
    \item the address of the memory location to access
    \item the instruction to be decoded
    \item condition flags such as zero, carry, and sign
    \item the data to be written to memory
  \end{enumerate}
\item \textit{[Original]} A quad-core processor means that the processor:
  \begin{enumerate}[label=\Alph*.]
    \item has four bytes of cache memory
    \item supports four monitors at once
    \item runs four operating systems
    \item contains four independent processing cores on a single chip
  \end{enumerate}
\item \textit{[Original]} In the office desktop, which unit fetches each
  instruction, decodes it, and issues control signals to the ALU, registers,
  and memory?
  \begin{enumerate}[label=\Alph*.]
    \item Arithmetic Logic Unit
    \item Control Unit
    \item Hard disk
    \item Monitor
  \end{enumerate}

\item \textit{[JKSSB FAA-pattern]} Which register holds the address of the next instruction to be executed?
  \begin{enumerate}[label=\Alph*.]
    \item Program Counter
    \item Stack Pointer
    \item Data Register
    \item Printer Buffer
  \end{enumerate}
\item \textit{[JKPSI-pattern]} MAR stands for:
  \begin{enumerate}[label=\Alph*.]
    \item Memory Address Register
    \item Main Access Router
    \item Machine Arithmetic Rule
    \item Monitor Alignment Record
  \end{enumerate}
\item \textit{[SSC CGL-pattern]} Which part of the CPU performs comparison operations such as greater-than checks?
  \begin{enumerate}[label=\Alph*.]
    \item ALU
    \item Monitor
    \item Hard disk
    \item Speaker
  \end{enumerate}
\item \textit{[JKSSB-pattern]} \textbf{(MSQ: select ALL correct options.)} Which of the following are CPU registers?
  \begin{enumerate}[label=\Alph*.]
    \item Program Counter (PC)
    \item Instruction Register (IR)
    \item Memory Address Register (MAR)
    \item USB port
  \end{enumerate}
\item \textit{[SSC CGL-pattern]} \textbf{(MSQ: select ALL correct options.)} Which statements about CPU internals are correct?
  \begin{enumerate}[label=\Alph*.]
    \item The Program Counter holds the next instruction's address.
    \item The Instruction Register holds the current instruction.
    \item The ALU carries out arithmetic and logical operations.
    \item The Control Unit stores user files permanently.
  \end{enumerate}
\item \textit{[PYQ-adapted: SSC CGL 2025 Tier-II Paper-I, 19-Jan-2026 Q6; SRC-18]} Which register holds the instruction currently being decoded?
  \begin{enumerate}[label=\Alph*.]
\item Memory Address Register
   \item Program Counter
   \item Instruction Register
   \item Accumulator
  \end{enumerate}

\item \textit{[PYQ-adapted: JKSSB Junior Assistant 2026 Q64, final key; SRC-03/SRC-06]} Which register normally holds the address of the next instruction to be fetched?
  \begin{enumerate}[label=\Alph*.]
\item Instruction Register
   \item Program Counter
   \item Accumulator
   \item Memory Data Register
  \end{enumerate}

\item \textit{[Original]} Which register temporarily holds the address of a memory location being accessed?
  \begin{enumerate}[label=\Alph*.]
\item MAR
   \item IR
   \item Accumulator
   \item Status register
  \end{enumerate}

\item \textit{[Original]} The clock rate of a processor is commonly expressed in:
  \begin{enumerate}[label=\Alph*.]
\item Bytes per second
   \item Hertz
   \item Pixels per inch
   \item Revolutions per minute
  \end{enumerate}

\item \textit{[Original, MSQ drill]} \textbf{(MSQ: select ALL correct options.)} Match each register with its common role.
  \begin{enumerate}[label=\Alph*.]
\item PC --- address of the next instruction
   \item IR --- current instruction being decoded
   \item MAR --- memory address for a transfer
   \item Accumulator --- intermediate arithmetic/logic result
  \end{enumerate}

\end{enumerate}
\lessonheading{06}{CPU vs GPU \& Microprocessor}
\begin{enumerate}[label=\textbf{Q\arabic*.}, leftmargin=*]
\item \textit{[Original]} Which statement best describes the cores of a CPU?
  \begin{enumerate}[label=\Alph*.]
    \item Many simple cores optimised for throughput
    \item A single storage unit
    \item Few powerful general-purpose cores optimised for latency
    \item No cores at all
  \end{enumerate}
\item \textit{[Original]} A GPU is designed primarily for:
  \begin{enumerate}[label=\Alph*.]
    \item sequential general-purpose work only
    \item parallel throughput tasks such as graphics and pixels
    \item power management
    \item secondary storage
  \end{enumerate}
\item \textit{[Original]} Serial processing means that:
  \begin{enumerate}[label=\Alph*.]
    \item multiple tasks run simultaneously
    \item one task is done at a time, in sequence
    \item work is split across many cores
    \item no instructions are executed
  \end{enumerate}
\item \textit{[Original]} Parallel processing means that:
  \begin{enumerate}[label=\Alph*.]
    \item one task is done at a time in strict order
    \item only graphics work is performed
    \item multiple tasks or operations proceed simultaneously
    \item a GPU is always required
  \end{enumerate}
\item \textit{[Original]} Which statement about GPU cores is correct?
  \begin{enumerate}[label=\Alph*.]
    \item GPU cores are identical to general-purpose CPU cores.
    \item GPU cores replace the CPU entirely.
    \item GPU cores are used only for storing files.
    \item GPU cores are simpler and more numerous, optimised for parallel
      throughput.
  \end{enumerate}
\item \textit{[Original]} GPU stands for:
  \begin{enumerate}[label=\Alph*.]
    \item General Processing Unit
    \item Graphics Processing Unit
    \item Graphical Parallel Unit
    \item Global Processing Utility
  \end{enumerate}
\item \textit{[Original]} CPU stands for:
  \begin{enumerate}[label=\Alph*.]
    \item Central Processing Unit
    \item Computer Peripheral Unit
    \item Central Program Utility
    \item Core Processing Unit
  \end{enumerate}
\item \textit{[Original]} A microprocessor is best described as:
  \begin{enumerate}[label=\Alph*.]
    \item a type of secondary storage
    \item a CPU implemented on a single chip
    \item a complete computer with memory and I/O on one chip
    \item an output device
  \end{enumerate}
\item \textit{[Original]} A microcontroller is best described as:
  \begin{enumerate}[label=\Alph*.]
    \item a CPU on a single chip with no memory or I/O
    \item a large mainframe processor
    \item a GPU with many cores
    \item a microprocessor plus integrated memory and I/O peripherals on one
      chip
  \end{enumerate}
\item \textit{[Original]} Microcontrollers are typically used in:
  \begin{enumerate}[label=\Alph*.]
    \item embedded and special-purpose systems
    \item supercomputers only
    \item plotters and printers only
    \item cloud data centres only
  \end{enumerate}
\item \textit{[Original]} Hyper-threading allows:
  \begin{enumerate}[label=\Alph*.]
    \item the physical core count to double
    \item one physical core to present two logical threads
    \item the clock speed to double
    \item multitasking to be disabled
  \end{enumerate}
\item \textit{[Original]} Evaluate the statements:
  \begin{enumerate}[label=\Roman*.]
    \item A CPU typically has fewer, more powerful cores than a GPU.
    \item A GPU has many simple cores designed for parallel throughput.
    \item GPU cores are fully general-purpose replacements for CPU cores.
  \end{enumerate}
  \begin{enumerate}[label=\Alph*.]
    \item I only
    \item I and II only
    \item II and III only
    \item I, II, and III
  \end{enumerate}
\item \textit{[Original]} L1 cache is:
  \begin{enumerate}[label=\Alph*.]
    \item the largest and slowest cache level
    \item the smallest and fastest level, closest to the core
    \item stored on the hard disk
    \item the same as main memory
  \end{enumerate}
\item \textit{[Original]} Which order lists cache from smallest and fastest to
  largest and slowest?
  \begin{enumerate}[label=\Alph*.]
    \item L3, L2, L1
    \item L2, L1, L3
    \item L1, L2, L3
    \item L1, L3, L2
  \end{enumerate}
\item \textit{[Original]} Clock speed is typically measured in GHz and
  expresses:
  \begin{enumerate}[label=\Alph*.]
    \item storage capacity
    \item memory size
    \item cycles per second
    \item number of cores
  \end{enumerate}
\item \textit{[Original]} A multi-core processor is one that has:
  \begin{enumerate}[label=\Alph*.]
    \item multiple independent processing units (cores) on one chip
    \item multiple monitors
    \item additional RAM
    \item a single slower core
  \end{enumerate}
\item \textit{[Original]} An office desktop has a quad-core 2.4\,GHz CPU,
  16\,GB RAM and an SSD, yet one particular program runs slowly. What is the
  most appropriate diagnosis?
  \begin{enumerate}[label=\Alph*.]
    \item The CPU must be faulty, because more cores always means faster.
    \item The GPU must be missing.
    \item The clock speed must be zero.
    \item A single-threaded program or a memory/disk-I/O bottleneck --- not
      automatically a ``slow CPU.''
  \end{enumerate}
\item \textit{[Original]} Which statement is \textbf{NOT} correct?
  \begin{enumerate}[label=\Alph*.]
    \item A CPU has few, powerful general-purpose cores.
    \item A GPU has many simpler cores for parallel work.
    \item A microprocessor integrates a CPU on a single chip.
    \item A microcontroller is just a CPU with no memory or I/O.
  \end{enumerate}
\item \textit{[Original]} The key difference between a microprocessor and a
  microcontroller is:
  \begin{enumerate}[label=\Alph*.]
    \item clock speed only
    \item colour and size only
    \item the level of integration --- the microcontroller includes memory
      and I/O on the same chip
    \item the microprocessor is analog
  \end{enumerate}
\item \textit{[Original]} Which workload is best suited to a GPU?
  \begin{enumerate}[label=\Alph*.]
    \item a single sequential spreadsheet calculation
    \item storing files
    \item printing a document
    \item processing millions of pixels or parallel floating-point operations
  \end{enumerate}

\item \textit{[JKSSB FAA-pattern]} The Intel 4004 is historically important as the first:
  \begin{enumerate}[label=\Alph*.]
    \item commercial microprocessor.
    \item hard disk drive.
    \item programming language.
    \item optical mouse.
  \end{enumerate}
\item \textit{[JKPSI-pattern]} A GPU is best suited for which kind of workload?
  \begin{enumerate}[label=\Alph*.]
    \item Massively parallel graphics and numeric tasks
    \item Power supply regulation
    \item Paper printing control
    \item Keyboard input scanning
  \end{enumerate}
\item \textit{[SSC CGL-pattern]} Which company developed the Pentium series of processors?
  \begin{enumerate}[label=\Alph*.]
    \item Intel
    \item A paper mill
    \item A textile firm
    \item A shipping company
  \end{enumerate}
\item \textit{[JKSSB-pattern]} \textbf{(MSQ: select ALL correct options.)} Which of the following are associated with microprocessors?
  \begin{enumerate}[label=\Alph*.]
    \item Intel
    \item AMD
    \item ARM-architecture designers
    \item Printer ink cartridges
  \end{enumerate}
\item \textit{[SSC CGL-pattern]} \textbf{(MSQ: select ALL correct options.)} Which statements about processors are correct?
  \begin{enumerate}[label=\Alph*.]
    \item A microprocessor places the CPU on a single chip.
    \item A GPU excels at parallel data-heavy tasks.
    \item Processor clock speed is measured in hertz.
    \item RAM chips sit inside the microprocessor package by definition.
  \end{enumerate}
\item \textit{[Original]} Which processor is generally optimized for a wide range of sequential general-purpose tasks?
  \begin{enumerate}[label=\Alph*.]
\item CPU
   \item GPU
   \item Network interface card
   \item Sound card
  \end{enumerate}

\item \textit{[Original]} A graphics workload applies the same operation to many pixels in parallel. Which component is commonly suited to this work?
  \begin{enumerate}[label=\Alph*.]
\item GPU
   \item Hard disk
   \item Keyboard controller
   \item ROM
  \end{enumerate}

\item \textit{[Original]} A microcontroller is most likely to combine a CPU with:
  \begin{enumerate}[label=\Alph*.]
\item Only an external display
   \item Memory and input/output peripherals on one chip
   \item A separate supercomputer
   \item A printer interface only
  \end{enumerate}

\item \textit{[Original]} Which is a typical example of an embedded-system use for a microcontroller?
  \begin{enumerate}[label=\Alph*.]
\item Dedicated washing-machine control
   \item A large-scale weather simulation
   \item A multi-user banking database server
   \item A general-purpose desktop operating system
  \end{enumerate}

\item \textit{[Original, MSQ drill]} \textbf{(MSQ: select ALL correct options.)} Which statements are generally correct?
  \begin{enumerate}[label=\Alph*.]
\item CPUs are designed for broad general-purpose processing.
   \item GPUs can execute many similar operations in parallel.
   \item Microcontrollers commonly serve embedded control tasks.
   \item A GPU is another name for a hard disk.
  \end{enumerate}

\end{enumerate}
\lessonheading{59}{Current CPU Architecture, Data Representation \& RISC/CISC}
\begin{enumerate}[label=\textbf{Q\arabic*.}, leftmargin=*]
\item \textit{[PYQ-adapted: JKSSB Junior Assistant 2026 Q62; final key;
paraphrased]} Which statement describes von Neumann (stored-program)
architecture?
  \begin{enumerate}[label=\Alph*.]
    \item Instructions and data are stored in two separate memories.
    \item Only data is stored in memory; instructions are not.
    \item Instructions and data are stored together in the same memory.
    \item Only instructions are stored in memory; data is not.
  \end{enumerate}
\item \textit{[Original]} In immediate addressing, the operand is:
  \begin{enumerate}[label=\Alph*.]
    \item held in a general-purpose register.
    \item contained within the instruction itself.
    \item stored in main memory at an address named by the instruction.
    \item the address of the next instruction to fetch.
  \end{enumerate}
\item \textit{[Original]} The cache-hit ratio is defined as:
  \begin{enumerate}[label=\Alph*.]
    \item hits divided by total accesses.
    \item misses divided by total accesses.
    \item hits divided by misses.
    \item total accesses divided by hits.
  \end{enumerate}
\item \textit{[Original]} Why is a cache hit faster than a cache miss?
  \begin{enumerate}[label=\Alph*.]
    \item The cache is located on secondary storage.
    \item The cache stores only instructions, never data.
    \item The cache is smaller and faster and sits closer to the CPU, so its
      access time is lower than main memory.
    \item The cache runs at the same speed as a hard disk.
  \end{enumerate}
\item \textit{[PYQ-adapted: JKSSB Junior Assistant 2026 Q62; final key;
paraphrased]} Which statement about data hazards is correct?
  \begin{enumerate}[label=\Alph*.]
    \item A hazard can arise only when an operand is in memory.
    \item Register addressing removes every possible hazard.
    \item A hazard can arise only with immediate addressing.
    \item A hazard can still arise from a dependency between instructions,
      for example in a pipeline.
  \end{enumerate}
\item \textit{[Original]} Which option correctly contrasts cache and virtual
memory?
  \begin{enumerate}[label=\Alph*.]
    \item Cache is on secondary storage, while virtual memory is in the
      cache.
    \item Cache and virtual memory are the same piece of hardware.
    \item Virtual memory is a single cache line.
    \item Cache is small fast memory between the CPU and main memory;
      virtual memory makes memory appear larger by using secondary storage as
      backing store.
  \end{enumerate}
\item \textit{[PYQ-adapted: JKSSB Junior Assistant 2026 Q64; final key;
paraphrased]} The Program Counter holds:
  \begin{enumerate}[label=\Alph*.]
    \item the address of the instruction currently being executed.
    \item the address of the next instruction to be fetched.
    \item the operand value of the current instruction.
    \item the last instruction that finished executing.
  \end{enumerate}
\item \textit{[Original]} The Instruction Register holds:
  \begin{enumerate}[label=\Alph*.]
    \item the address of the next instruction to fetch.
    \item the instruction currently being fetched, decoded, and executed.
    \item the value of the operand.
    \item the width of the data bus.
  \end{enumerate}
\item \textit{[Original]} Interrupts are normally processed:
  \begin{enumerate}[label=\Alph*.]
    \item only after the entire program has completed.
    \item only during system boot.
    \item only for hardware faults.
    \item at instruction boundaries, during program execution.
  \end{enumerate}
\item \textit{[Original]} Which statement about instruction timing is correct?
  \begin{enumerate}[label=\Alph*.]
    \item Every machine instruction takes exactly one clock cycle.
    \item The clock-cycle count is fixed for the whole program.
    \item Different machine instructions can take different numbers of clock
      cycles.
    \item Machine instructions do not consume clock cycles.
  \end{enumerate}
\item \textit{[PYQ-adapted: JKSSB Junior Assistant 2026 Q65; final key;
paraphrased]} In two's-complement arithmetic, overflow occurs when:
  \begin{enumerate}[label=\Alph*.]
    \item there is a carry out of the most significant bit.
    \item the result of an addition is exactly zero.
    \item the operation is performed on unsigned numbers.
    \item two operands of the same sign produce a result of the opposite
      sign.
  \end{enumerate}
\item \textit{[Original]} When widening a value to more bits, which rule is
correct?
  \begin{enumerate}[label=\Alph*.]
    \item Zero-extension is for signed values and sign-extension is for
      unsigned values.
    \item Zero-extension is for unsigned values and sign-extension is for
      signed values.
    \item Both unsigned and signed values use zero-extension.
    \item Both unsigned and signed values use sign-extension.
  \end{enumerate}
\item \textit{[PYQ-adapted: JKSSB Junior Assistant 2026 Q65; final key;
paraphrased]} In BCD (binary-coded decimal), the decimal number 25 is
represented as:
  \begin{enumerate}[label=\Alph*.]
    \item the plain binary value $11001_2$.
    \item eight bits per decimal digit.
    \item the hexadecimal value $25_{16}$.
    \item $0010\ 0101$, one 4-bit code for each decimal digit.
  \end{enumerate}
\item \textit{[Original]} Which statement about byte-addressability and the
effective address is correct?
  \begin{enumerate}[label=\Alph*.]
    \item The instruction length is always equal to the effective address.
    \item Memory is word-addressable only, so individual bytes have no
      address.
    \item Each byte has its own address, and the effective address is the
      operand's actual location, which need not equal the instruction
      length.
    \item The effective address is always equal to the Program Counter value.
  \end{enumerate}
\item \textit{[Original]} What are the full forms of RISC and CISC?
  \begin{enumerate}[label=\Alph*.]
    \item RISC = Reduced Instruction Set Computer; CISC = Complex
      Instruction Set Computer
    \item RISC = Rapid Instruction Set Computer; CISC = Complex Instruction
      Set Computer
    \item RISC = Reduced Integer Storage Cache; CISC = Central Instruction
      Set Computer
    \item RISC = Random Instruction Set Computing; CISC = Common Instruction
      Set Computer
  \end{enumerate}
\item \textit{[Original]} Which set of features describes a RISC processor?
  \begin{enumerate}[label=\Alph*.]
    \item Many complex, variable-length instructions.
    \item Instructions that always operate directly on memory operands.
    \item Microprogrammed control of highly complex instructions.
    \item Simple, mostly fixed-length instructions, many registers, and a
      load/store architecture.
  \end{enumerate}
\item \textit{[Original]} Which set of features describes a CISC processor?
  \begin{enumerate}[label=\Alph*.]
    \item A few simple fixed-length instructions only.
    \item A large set of complex, variable-length instructions, often
      allowing memory operands.
    \item No general-purpose registers at all.
    \item A strict load/store-only memory access rule.
  \end{enumerate}
\item \textit{[Original]} Which statement is \textbf{NOT} correct?
  \begin{enumerate}[label=\Alph*.]
    \item The Program Counter holds the address of the next instruction to
      fetch.
    \item Interrupts can be handled at instruction boundaries.
    \item All machine instructions take the same number of clock cycles.
    \item Two's-complement overflow can occur when two same-sign operands are
      added.
  \end{enumerate}
\item \textit{[Original]} Which code correctly represents the decimal digit 5
in BCD?
  \begin{enumerate}[label=\Alph*.]
    \item $0101$
    \item $1001$
    \item $101$
    \item $1110$
  \end{enumerate}
\item \textit{[Original]} Evaluate the statements:
  \begin{enumerate}[label=\Roman*.]
    \item In immediate addressing, the operand is part of the instruction.
    \item The cache-hit ratio equals hits divided by total accesses.
    \item Two's-complement overflow is the same as a carry out of the most
      significant bit.
  \end{enumerate}
  \begin{enumerate}[label=\Alph*.]
    \item I and II only
    \item II and III only
    \item I and III only
    \item I, II, and III
  \end{enumerate}
\item \textit{[Original]} The von Neumann bottleneck refers to:
  \begin{enumerate}[label=\Alph*.]
    \item the delay between keyboard input and screen output.
    \item the limited rate at which instructions and data can be transferred
      between the CPU and memory over a shared path.
    \item the time taken for a hard disk to spin up.
    \item the number of general-purpose registers in the CPU.
  \end{enumerate}
\item \textit{[Original]} The main purpose of virtual memory is to:
  \begin{enumerate}[label=\Alph*.]
    \item increase the number of CPU registers.
    \item replace the cache entirely.
    \item reduce the CPU clock speed.
    \item make the available memory appear larger than physical RAM by using
      secondary storage as backing store.
  \end{enumerate}

\item \textit{[JKSSB FAA-pattern]} RISC designs are characterised by:
  \begin{enumerate}[label=\Alph*.]
    \item a small set of simple instructions executed quickly.
    \item hundreds of complex printer commands.
    \item manual paper calculation only.
    \item slower clocks by definition.
  \end{enumerate}
\item \textit{[JKPSI-pattern]} CISC designs are characterised by:
  \begin{enumerate}[label=\Alph*.]
    \item many complex multi-step instructions.
    \item no use of memory at all.
    \item operations only on paper.
    \item absence of any control unit.
  \end{enumerate}
\item \textit{[SSC CGL-pattern]} Which processor family follows the x86 CISC tradition?
  \begin{enumerate}[label=\Alph*.]
    \item Intel Core desktop processors
    \item Paper calculators
    \item Mechanical typewriters
    \item Wall clocks
  \end{enumerate}
\item \textit{[JKSSB-pattern]} \textbf{(MSQ: select ALL correct options.)} Which statements about RISC are correct?
  \begin{enumerate}[label=\Alph*.]
    \item Instructions are simple and uniform.
    \item Most execution uses registers (load-store design).
    \item Pipelining is easy to implement well.
    \item RISC forbids the use of any memory.
  \end{enumerate}
\item \textit{[SSC CGL-pattern]} \textbf{(MSQ: select ALL correct options.)} Which statements comparing RISC and CISC are correct?
  \begin{enumerate}[label=\Alph*.]
    \item RISC uses fewer, simpler instructions than CISC.
    \item CISC instructions can complete complex tasks in fewer lines.
    \item ARM designs follow RISC ideas.
    \item CISC processors contain no control unit.
  \end{enumerate}
\item \textit{[PYQ-adapted: SSC CGL 2025 Tier-II Paper-I, 19-Jan-2026 Q6; SRC-18]} Which register holds the instruction currently being decoded and executed?
  \begin{enumerate}[label=\Alph*.]
\item MAR
   \item PC
   \item IR
   \item Accumulator
  \end{enumerate}

\item \textit{[PYQ-adapted: JKSSB Junior Assistant 2026 Q65, final key; SRC-03/SRC-06]} In BCD, each decimal digit is represented using:
  \begin{enumerate}[label=\Alph*.]
\item Four bits
   \item Eight bits
   \item One bit
   \item Sixteen bits
  \end{enumerate}

\item \textit{[PYQ-adapted: JKSSB Junior Assistant 2026 Q65, final key; SRC-03/SRC-06]} For unsigned widening, which method fills the newly added high-order bits with zeros?
  \begin{enumerate}[label=\Alph*.]
\item Zero extension
   \item Sign extension
   \item Sign inversion
   \item Carry propagation
  \end{enumerate}

\item \textit{[Original]} In a pipelined processor, a data hazard can occur when:
  \begin{enumerate}[label=\Alph*.]
\item An instruction depends on a result not yet produced by an earlier instruction
   \item A printer runs out of paper
   \item A file extension is hidden
   \item A keyboard key is repeated
  \end{enumerate}

\item \textit{[Original, MSQ drill]} \textbf{(MSQ: select ALL correct options.)} Which are standard CPU architecture concepts?
  \begin{enumerate}[label=\Alph*.]
\item Program Counter stores the next instruction address.
   \item Instruction Register holds the current instruction.
   \item A carry-out and signed overflow are always identical.
   \item BCD uses four bits per decimal digit.
  \end{enumerate}

\end{enumerate}
\answersection{2}
\answersubheading{04}{Organization of a Computer}
\begin{enumerate}[label=\textbf{Q\arabic*.}, leftmargin=*]
\item \textbf{A.} ALU stands for Arithmetic Logic Unit. ``Application Logic
  Unit'', ``Arithmetic Language Unit'' and ``Advanced Logic Unit'' are
  invented expansions that do not name the unit correctly.
\item \textbf{C.} The ALU performs arithmetic and logical (comparison)
  operations. Permanent storage (A) is a storage job, control of I/O devices
  (B) is the control unit's work, and supplying power (D) is the power
  supply's role.
\item \textbf{D.} The control unit directs and coordinates all other units,
  including the ALU. Options A, B and C name units that are themselves
  coordinated, not the coordinating unit.
\item \textbf{C.} The CPU consists of the ALU, the control unit and registers.
  Input/output units and memory (A) are separate units; main and secondary
  storage (B) are storage, not the CPU; monitor, keyboard and mouse (D) are
  peripheral devices.
\item \textbf{A.} CPU expands to Central Processing Unit. The other three
  options are incorrect wordings that are not the standard expansion.
\item \textbf{A.} Registers are small, fast storage inside the CPU used to
  hold data and addresses during processing. Secondary storage (B) and a hard
  disk (C) are slower and outside the CPU; the output unit (D) presents
  results.
\item \textbf{B.} The address bus carries the address of a memory location
  from the CPU to memory. Data travel on the data bus (A), control signals on
  the control bus (C), and power is supplied separately (D), not over the
  address bus.
\item \textbf{B.} The data bus is bidirectional, moving data and instructions
  both ways. The address bus (A) is one-way from the CPU, the control bus (C)
  carries control signals, and there is no separate ``power bus'' in this
  three-bus model (D).
\item \textbf{C.} The control bus carries control and timing signals such as
  read and write. Addresses travel on the address bus (A), data on the data
  bus (B), and keyboard input is not what the control bus carries (D).
\item \textbf{C.} This is the incorrect statement: the CPU accesses main
  memory directly, not secondary storage. Statements A, B and D are all true
  of the ALU, the control unit and the CPU respectively.
\item \textbf{B.} The CPU directly accesses main memory (primary storage).
  Secondary storage (A), a pen drive (C) and a hard disk (D) are all
  secondary storage, reached only through main memory and the OS.
\item \textbf{C.} The hard disk is secondary storage. RAM (A), ROM (B) and
  registers (D) are primary storage or CPU storage, not secondary storage.
\item \textbf{A.} Main memory holds the data and instructions currently being
  processed. Option B is the classic error (the CPU does not process directly
  from secondary storage); RAM is volatile, not non-volatile secondary storage
  (C); and registers are faster than secondary storage (D).
\item \textbf{D.} The instruction cycle runs fetch, decode, execute. Options
  A, B and C scramble this fixed order.
\item \textbf{C.} During fetch the CPU obtains the next instruction from main
  memory. Arithmetic is the execute step (A); writing the result back is the
  store step (B); displaying output (D) is not the fetch step.
\item \textbf{C.} The stored-program concept stores both instructions and data
  in the same main memory. Option A wrongly excludes instructions; option B
  wrongly limits programs to secondary storage; option D places instructions
  in an output device.
\item \textbf{A.} Storing instructions and data together in memory is the von
  Neumann architecture idea. The address bus (B), secondary storage (C) and
  the output unit (D) are unrelated to the concept.
\item \textbf{D.} All three statements are correct: the address bus carries
  the address from the CPU to memory, the data bus is bidirectional, and the
  control bus carries control signals. Options A, B and C each omit at least
  one true statement.
\item \textbf{D.} Expansion slots let you add cards such as graphics or
  network cards. They do not permanently store the OS (A), replace the
  control unit (B), or supply power (C).
\item \textbf{C.} The control bus carries control signals. The address bus is
  unidirectional, not bidirectional (A); the data bus is bidirectional, not
  unidirectional (B); and the data bus carries data, not addresses (D).

\item \textbf{A.} The CPU processes instructions and is called the brain. B, C, and D are peripherals or storage, not the processing brain.
\item \textbf{A.} The ALU handles arithmetic and logical work. B, C, and D describe peripheral or unrelated roles.
\item \textbf{A.} CU is the Control Unit. B, C, and D are not the recognised expansion.
\item \textbf{A, B, C.} The ALU, CU, and registers form the CPU. The hard disk (D) is secondary storage outside the CPU.
\item \textbf{A, B, D.} The CPU composition, the data-flow order, and memory's role are correct. C is false: the CU directs operations; it does not permanently store files.
\item \textbf{B.} The Control Unit sequences operations and directs data flow. The ALU performs arithmetic/logic; memory stores data; input supplies data. The cited key is provisional.

\item \textbf{B.} The ALU performs arithmetic and logical operations. The Control Unit directs execution, the Program Counter holds an instruction address, and the input unit supplies data.

\item \textbf{A.} The address bus carries location addresses. The data bus carries data; the control bus carries timing/control signals; expansion buses connect peripherals.

\item \textbf{B.} Input supplies data, processing transforms it, and output presents results. The other sequences reverse that basic order.

\item \textbf{A, B, C, D.} A general-purpose computer has input, memory, processing (CPU), and output functions.

\end{enumerate}
\answersubheading{05}{CPU, ALU, Control Unit \& Registers}
\begin{enumerate}[label=\textbf{Q\arabic*.}, leftmargin=*]
\item \textbf{A.} ALU stands for Arithmetic Logic Unit. B, C, and D are not
  expansions of the abbreviation.
\item \textbf{C.} The Arithmetic Logic Unit performs arithmetic and logical
  operations. The Control Unit directs operations; the register file stores
  values; the system bus carries signals.
\item \textbf{B.} The Control Unit directs the sequence of operations and
  issues control signals. Performing arithmetic is the ALU's job (A); the
  CU does not store results (C) or provide secondary storage (D).
\item \textbf{D.} Adding two numbers is an arithmetic operation performed by
  the ALU, not the Control Unit. Decoding instructions, issuing control
  signals, and directing data flow are all Control Unit functions.
\item \textbf{B.} The Program Counter holds the address of the next
  instruction to be fetched. The IR holds the current instruction; the MDR
  holds data; the accumulator holds intermediate results.
\item \textbf{C.} The Instruction Register holds the instruction currently
  being decoded and executed. The PC holds the next address; the MDR holds
  data; the MAR holds an address.
\item \textbf{B.} The Memory Address Register holds the address of the memory
  location to be accessed. The MDR holds data; the IR holds an instruction;
  the accumulator holds a result.
\item \textbf{C.} The Memory Data Register holds the data being transferred
  to or from memory. The MAR holds the address; the PC holds the next
  instruction address; the IR holds the instruction.
\item \textbf{B.} The accumulator holds intermediate arithmetic and logic
  results. It does not hold the next-instruction address (A), the fetched
  instruction (C), or a permanent copy of the OS (D).
\item \textbf{A.} Registers are the fastest storage, followed by L1 cache,
  RAM, and secondary storage such as an SSD. B, C, and D are progressively
  slower.
\item \textbf{D.} Registers are located inside the CPU. They are not part of
  cache (A), main memory (B), or secondary storage (C).
\item \textbf{B.} One gigahertz is one billion ($10^9$) cycles per second, so
  2.4 GHz is about 2.4 billion cycles per second. A, C, and D misstate the
  giga- scale.
\item \textbf{C.} Different instructions may take different numbers of clock
  cycles, so A and B are false. A clock cycle is a unit of time, not bytes
  (D).
\item \textbf{B.} Statements I and II are correct: the ALU works on binary
  data with logic gates, and the Control Unit directs operations and data
  flow. III is false because the Control Unit does not perform arithmetic.
  A omits II; C includes false III and omits I; D includes false III.
\item \textbf{A.} PC is correctly matched with Program Counter. IR is
  Instruction Register, MAR is Memory Address Register, and MDR is Memory
  Data Register, so B, C, and D are wrong expansions.
\item \textbf{A.} The Program Counter holds the address of the next
  instruction to be fetched. It does not hold the current instruction (B),
  the previously executed one (C), or data (D).
\item \textbf{B.} Registers with no fixed dedicated role, used for general
  operands and intermediate results, are general-purpose registers. Status
  registers hold flags; MAR holds addresses; IR holds instructions.
\item \textbf{C.} The status or flags register holds condition flags such as
  zero, carry, and sign. A describes the MAR; B describes the IR; D
  describes the MDR.
\item \textbf{D.} A quad-core processor contains four independent processing
  cores on a single chip. It does not mean four bytes of cache (A), four
  monitors (B), or four operating systems (C).
\item \textbf{B.} The Control Unit fetches each instruction, decodes it, and
  issues control signals to the other units. The ALU computes; the hard disk
  is secondary storage; the monitor is an output device.

\item \textbf{A.} The Program Counter tracks the next instruction. B, C, and D hold other values, not the next-instruction address.
\item \textbf{A.} MAR is the Memory Address Register. B, C, and D are not the recognised expansion.
\item \textbf{A.} Comparisons are logical operations done by the ALU. B, C, and D are peripherals or storage.
\item \textbf{A, B, C.} The PC, IR, and MAR are registers. A USB port (D) is an interface, not a register.
\item \textbf{A, B, C.} The PC, IR, and ALU roles are stated correctly. D is false: the CU directs work; it does not store files.
\item \textbf{C.} The Instruction Register holds the fetched instruction during decode/execute. MAR holds an address, PC holds the next instruction address, and the accumulator stores intermediate results.

\item \textbf{B.} The Program Counter holds the next instruction address. The IR holds the instruction itself, the accumulator holds results, and MDR holds data transferred to/from memory.

\item \textbf{A.} The Memory Address Register holds a memory address. IR holds an instruction, the accumulator stores results, and a status register stores condition flags.

\item \textbf{B.} Clock rate is cycles per second, measured in hertz. The other units measure data rate, image density, and rotational speed.

\item \textbf{A, B, C, D.} All four register-role pairings are correct.

\end{enumerate}
\answersubheading{06}{CPU vs GPU \& Microprocessor}
\begin{enumerate}[label=\textbf{Q\arabic*.}, leftmargin=*]
\item \textbf{C.} A CPU has few powerful general-purpose cores optimised for
  latency. A describes a GPU; B and D misdescribe the CPU.
\item \textbf{B.} A GPU targets parallel throughput work such as graphics and
  pixels. A names the CPU's general-purpose sequential role; C and D are
  unrelated to GPU design.
\item \textbf{B.} Serial processing performs one task at a time, in sequence.
  A and C describe parallel behaviour; D is false.
\item \textbf{C.} Parallel processing runs several tasks or operations at the
  same time. A is the serial definition; B is too narrow; D is false, since
  parallel work does not always require a GPU.
\item \textbf{D.} GPU cores are simpler and more numerous, optimised for
  parallel throughput. A wrongly equates them with CPU cores; B and C
  misstate their role.
\item \textbf{B.} GPU stands for Graphics Processing Unit. A, C, and D are
  invented expansions.
\item \textbf{A.} CPU stands for Central Processing Unit. B, C, and D are
  invented expansions.
\item \textbf{B.} A microprocessor is a CPU implemented on a single chip, and
  it is general-purpose. A and D name other components; C actually describes
  a microcontroller.
\item \textbf{D.} A microcontroller is a microprocessor plus integrated memory
  and I/O peripherals on one chip. A describes a plain CPU; B and C are
  unrelated devices.
\item \textbf{A.} Microcontrollers are used in embedded and special-purpose
  systems such as appliances and controllers. B, C, and D name other classes
  of machine.
\item \textbf{B.} Hyper-threading lets one physical core present two logical
  threads. It does not double the physical core count (A) or the clock speed
  (C), and it does not disable multitasking (D).
\item \textbf{B.} Statements I and II are correct standard descriptions; III
  is false because GPU cores are simpler and parallel-oriented, not fully
  general-purpose replacements for CPU cores. A omits II; C includes false
  III and omits I; D also includes false III.
\item \textbf{B.} L1 is the smallest and fastest level, closest to the core. A
  and D confuse the levels or their role; C wrongly places cache on the hard
  disk.
\item \textbf{C.} The order from smallest and fastest to largest and slowest is
  L1, L2, L3. A, B, and D scramble this order.
\item \textbf{C.} Clock speed in GHz expresses cycles per second. A, B, and D
  describe storage, memory, or core count, not clock speed.
\item \textbf{A.} A multi-core processor has multiple independent processing
  units (cores) on one chip. B, C, and D describe unrelated properties.
\item \textbf{D.} A single-threaded program uses only one core, and the
  remaining slowness usually comes from a memory or disk-I/O bottleneck, so
  the CPU is not automatically at fault. A overstates the benefit of cores; B
  and C are unrelated to the symptom.
\item \textbf{D.} A microcontroller is not just a CPU: it integrates memory and
  I/O on the same chip. A, B, and C are correct statements, so D is the one
  that is NOT correct.
\item \textbf{C.} The key difference is the level of integration: the
  microcontroller includes memory and I/O on the same chip. A and B name
  secondary attributes; D is false.
\item \textbf{D.} A GPU suits parallel work such as processing millions of
  pixels or parallel floating-point operations. A is a sequential
  general-purpose task (CPU); B and C are not compute-parallel workloads.

\item \textbf{A.} The 4004 was the first commercial microprocessor. B, C, and D name unrelated hardware or software.
\item \textbf{A.} GPUs handle many parallel threads for graphics and numeric work. B, C, and D are not processor workloads.
\item \textbf{A.} Pentium is an Intel family. B, C, and D are non-computer firms.
\item \textbf{A, B, C.} Intel, AMD, and ARM-architecture designers all make processor technology. Ink cartridges (D) are printer consumables.
\item \textbf{A, B, C.} The single-chip CPU, the GPU's parallel strength, and hertz clock units are correct. D is false: RAM is separate main memory.
\item \textbf{A.} A CPU is a general-purpose processor. GPUs specialize in parallel graphics/data workloads; a NIC and sound card are peripheral controllers.

\item \textbf{A.} A GPU has many parallel processing elements suited to graphics operations. Storage, input control, and firmware memory do not perform this computation.

\item \textbf{B.} A microcontroller integrates processing, memory, and peripherals for embedded control. The other choices do not describe that integrated design.

\item \textbf{A.} A washing-machine controller performs a dedicated embedded task. The other workloads typically use general-purpose or high-performance computer systems.

\item \textbf{A, B, C.} The first three describe processor roles. A GPU processes data; a hard disk stores it.

\end{enumerate}
\answersubheading{59}{Current CPU Architecture, Data Representation \& RISC/CISC}
\begin{enumerate}[label=\textbf{Q\arabic*.}, leftmargin=*]
\item \textbf{C.} von Neumann architecture stores instructions and data
  together in the same memory. A describes a separate-memory (Harvard-style)
  design; B and D wrongly keep only one of the two in memory.
\item \textbf{B.} Immediate addressing carries the operand inside the
  instruction encoding. A describes register addressing; C describes direct
  addressing; D confuses the operand with the next address.
\item \textbf{A.} The hit ratio is hits divided by total accesses. B gives the
  miss ratio; C and D invert the intended quantities.
\item \textbf{C.} A cache hit avoids a slow main-memory access because the
  cache is small, fast, and close to the CPU. A places the cache on secondary
  storage; B wrongly limits it to instructions; D gives it disk speed.
\item \textbf{D.} Hazards stem from dependencies between instructions, such as
  a read-after-write dependency exposed by pipelining. A wrongly restricts
  hazards to memory operands; B claims register addressing removes all
  hazards; C wrongly ties hazards to immediate addressing.
\item \textbf{D.} Cache is small fast memory between the CPU and main memory;
  virtual memory makes memory appear larger using secondary storage. A
  misplaces both; B treats them as one device; C reduces virtual memory to a
  single cache line.
\item \textbf{B.} The Program Counter holds the address of the next
  instruction to fetch. A is the standard trap (current instruction); C
  confuses it with an operand; D describes a past instruction.
\item \textbf{B.} The Instruction Register holds the instruction currently
  being fetched, decoded, and executed. A is the PC's role; C names an operand
  value; D names the data-bus width.
\item \textbf{D.} Interrupts are recognised at instruction boundaries during
  execution. A is the trap (only after the whole program); B and C wrongly
  limit interrupts to boot or to hardware faults.
\item \textbf{C.} Different instructions can take different numbers of clock
  cycles. A and B assert a constant count; D denies that instructions use
  cycles at all.
\item \textbf{D.} Overflow is detected when adding two same-sign operands
  yields an opposite-sign result. A is the trap (carry out of the MSB); B and
  C misstate the condition.
\item \textbf{B.} Zero-extension is for unsigned values and sign-extension is
  for signed values. A swaps them (the trap); C and D apply a single rule to
  both cases.
\item \textbf{D.} BCD encodes each decimal digit in 4 bits, so 25 is
  $0010\ 0101$. A is the plain binary value; B uses eight bits per digit; C
  treats the number as hexadecimal.
\item \textbf{C.} Memory is byte-addressable, so each byte has its own
  address, and the effective address is the operand's actual location, which
  need not equal the instruction length. A and D wrongly equate unrelated
  quantities; B denies byte-addressability.
\item \textbf{A.} RISC is Reduced Instruction Set Computer and CISC is
  Complex Instruction Set Computer. B, C, and D give incorrect expansions of
  the abbreviations.
\item \textbf{D.} RISC uses simple, mostly fixed-length instructions, many
  registers, and a load/store architecture. A and C describe CISC traits; B
  contradicts the load/store discipline.
\item \textbf{B.} CISC offers a large set of complex, variable-length
  instructions and often permits memory operands. A and D describe RISC; C is
  false of CISC designs.
\item \textbf{C.} ``All instructions take the same number of clock cycles'' is
  false; the count varies. A, B, and D are all correct statements about the PC,
  interrupts, and overflow.
\item \textbf{A.} The BCD code for the digit 5 is $0101$. B is the code for 9;
  C is the plain binary of 5, not its BCD code; D is not a valid BCD digit.
\item \textbf{A.} Statements I and II are true. III is false because overflow
  is a sign condition, not a carry out of the MSB. B includes false III; C
  includes III and omits II; D includes III.
\item \textbf{B.} The von Neumann bottleneck is the limited rate at which
  instructions and data move between the CPU and memory over a shared path. A,
  C, and D name unrelated delays or quantities.
\item \textbf{D.} Virtual memory makes memory appear larger than physical RAM
  by using secondary storage as backing store. A, B, and C describe unrelated
  or incorrect effects.

\item \textbf{A.} RISC means reduced, fast instructions. B, C, and D misstate the design idea.
\item \textbf{A.} CISC packs complex instructions. B, C, and D deny normal processor parts.
\item \textbf{A.} Intel Core chips continue the x86 CISC line. B, C, and D are non-processor objects.
\item \textbf{A, B, C.} Simple instructions, load-store style, and pipelining ease are correct. D is false: RISC still uses memory via loads/stores.
\item \textbf{A, B, C.} The instruction-count, code-density, and ARM facts are correct. D is false: every CPU has a control unit.
\item \textbf{C.} The Instruction Register holds the current instruction. MAR stores an address, PC the next instruction address, and the accumulator an intermediate result.

\item \textbf{A.} BCD encodes each decimal digit in four bits. It is not the same as converting the whole number to pure binary.

\item \textbf{A.} Zero extension pads an unsigned value with zeros. Sign extension copies the sign bit for signed values.

\item \textbf{A.} A data hazard is a dependency conflict between overlapping instructions. The other choices are peripheral or file-management issues.

\item \textbf{A, B, D.} PC, IR, and BCD statements are correct. Carry-out is not generally equivalent to signed overflow.

\end{enumerate}
\unitopener{3}{Memory, Storage \& Backup}
\lessonheading{07}{Computer Memory Hierarchy}
\begin{enumerate}[label=\textbf{Q\arabic*.}, leftmargin=*]
\item \textit{[Original]} Expand RAM.
  \begin{enumerate}[label=\Alph*.]
    \item Read Access Memory
    \item Random Access Memory
    \item Rapid Access Memory
    \item Random Allocation Memory
  \end{enumerate}
\item \textit{[Original]} Expand ROM.
  \begin{enumerate}[label=\Alph*.]
    \item Random-Only Memory
    \item Read-Once Memory
    \item Read-Only Memory
    \item Runtime Object Memory
  \end{enumerate}
\item \textit{[Original]} Which option correctly lists the memory hierarchy
  in order of decreasing speed (fastest first)?
  \begin{enumerate}[label=\Alph*.]
    \item Registers, cache, main memory (RAM), secondary storage
    \item Cache, registers, main memory (RAM), secondary storage
    \item Main memory (RAM), registers, cache, secondary storage
    \item Secondary storage, main memory (RAM), cache, registers
  \end{enumerate}
\item \textit{[Original]} Which of the following is the fastest storage in a
  computer system?
  \begin{enumerate}[label=\Alph*.]
    \item Main memory (RAM)
    \item L1 cache
    \item Solid-state drive (SSD)
    \item Registers
  \end{enumerate}
\item \textit{[Original]} Cache memory is \underline{\hspace{2.2cm}} than
  main memory (RAM).
  \begin{enumerate}[label=\Alph*.]
    \item slower
    \item faster
    \item larger
    \item cheaper per bit
  \end{enumerate}
\item \textit{[Original]} Which statement correctly distinguishes cache
  memory from main memory (RAM)?
  \begin{enumerate}[label=\Alph*.]
    \item Cache is larger and faster than RAM.
    \item Cache and RAM have the same speed and size.
    \item Cache is smaller and faster than RAM; RAM is larger and slower.
    \item Cache is a permanent backup of the contents of RAM.
  \end{enumerate}
\item \textit{[Original]} Which level of cache is the smallest and fastest?
  \begin{enumerate}[label=\Alph*.]
    \item L1
    \item L2
    \item L3
    \item L4
  \end{enumerate}
\item \textit{[Original]} Among the cache levels L1, L2, and L3, which is the
  largest and slowest (though still faster than main memory)?
  \begin{enumerate}[label=\Alph*.]
    \item L1
    \item L2
    \item Registers
    \item L3
  \end{enumerate}
\item \textit{[Original]} Which of the following is \textbf{NOT} true of
  cache memory?
  \begin{enumerate}[label=\Alph*.]
    \item It stores recently and frequently used data.
    \item It helps exploit locality of reference.
    \item It acts as a permanent backup for the contents of RAM.
    \item It is faster than main memory.
  \end{enumerate}
\item \textit{[Original]} RAM is a \underline{\hspace{2.2cm}} memory, whereas
  ROM is \underline{\hspace{2.2cm}}.
  \begin{enumerate}[label=\Alph*.]
    \item non-volatile; volatile
    \item volatile; non-volatile
    \item read-only; writable
    \item secondary; primary
  \end{enumerate}
\item \textit{[Original]} Which memory loses its contents when the power is
  switched off?
  \begin{enumerate}[label=\Alph*.]
    \item RAM
    \item ROM
    \item Hard disk drive
    \item Solid-state drive
  \end{enumerate}
\item \textit{[Original]} Which of the following is \textbf{primary memory}?
  \begin{enumerate}[label=\Alph*.]
    \item Solid-state drive
    \item Hard disk drive
    \item Magnetic tape
    \item RAM
  \end{enumerate}
\item \textit{[Original]} Which of the following is a \textbf{secondary
  storage} device?
  \begin{enumerate}[label=\Alph*.]
    \item Register
    \item Cache
    \item Solid-state drive (SSD)
    \item RAM
  \end{enumerate}
\item \textit{[Original]} Why does a memory hierarchy exist?
  \begin{enumerate}[label=\Alph*.]
    \item Faster memory is smaller and more expensive per bit, so the
      hierarchy balances speed, cost, and capacity.
    \item All memory types have the same speed and cost.
    \item Secondary storage is the fastest type of memory.
    \item Registers can store the entire contents of main memory.
  \end{enumerate}
\item \textit{[Original]} Evaluate the statements:
  \begin{enumerate}[label=\Roman*.]
    \item Registers are the fastest storage in a computer.
    \item Cache memory is faster than main memory.
    \item RAM is non-volatile and retains its data when the power is off.
  \end{enumerate}
  \begin{enumerate}[label=\Alph*.]
    \item I only
    \item I and II only
    \item II and III only
    \item I, II, and III
  \end{enumerate}
\item \textit{[Original]} The tendency of a processor to reuse recently
  accessed data and nearby data is called:
  \begin{enumerate}[label=\Alph*.]
    \item locality of reference
    \item overflow
    \item fragmentation
    \item deadlock
  \end{enumerate}
\item \textit{[Original]} Temporal locality refers to:
  \begin{enumerate}[label=\Alph*.]
    \item accessing data stored at nearby addresses.
    \item storing data permanently on secondary storage.
    \item reusing the same data or instruction again soon after it was
      accessed.
    \item losing data when the power is switched off.
  \end{enumerate}
\item \textit{[Original]} Which memory is typically the \textbf{smallest} in
  capacity?
  \begin{enumerate}[label=\Alph*.]
    \item Solid-state drive (SSD)
    \item Main memory (RAM)
    \item L1 cache
    \item Registers
  \end{enumerate}
\item \textit{[Original]} Which component retains its contents without
  power, so it can store files between sessions?
  \begin{enumerate}[label=\Alph*.]
    \item Registers
    \item Secondary storage such as an SSD or HDD
    \item Cache memory
    \item RAM
  \end{enumerate}
\item \textit{[Original]} Which pairing is correctly matched?
  \begin{enumerate}[label=\Alph*.]
    \item Registers --- slowest storage in the computer
    \item Cache --- non-volatile permanent backup of RAM
    \item RAM --- non-volatile read-only memory
    \item SSD --- secondary storage that retains data without power
  \end{enumerate}

\item \textit{[JKSSB FAA-pattern]} Which memory in the hierarchy is the fastest?
  \begin{enumerate}[label=\Alph*.]
    \item Registers
    \item Hard disk
    \item Magnetic tape
    \item DVD
  \end{enumerate}
\item \textit{[JKPSI-pattern]} Which option lists memories in decreasing order of speed (fastest first)?
  \begin{enumerate}[label=\Alph*.]
    \item Registers, cache, RAM, hard disk
    \item Hard disk, RAM, cache, registers
    \item RAM, hard disk, registers, cache
    \item Cache, hard disk, registers, RAM
  \end{enumerate}
\item \textit{[SSC CGL-pattern]} Which memory sits between the CPU and main memory to speed up access?
  \begin{enumerate}[label=\Alph*.]
    \item Cache memory
    \item Magnetic tape
    \item Floppy disk
    \item Printer buffer card
  \end{enumerate}
\item \textit{[JKSSB-pattern]} \textbf{(MSQ: select ALL correct options.)} Which of the following are primary (main) memory?
  \begin{enumerate}[label=\Alph*.]
    \item RAM
    \item ROM
    \item Cache memory
    \item Hard disk drive
  \end{enumerate}
\item \textit{[SSC CGL-pattern]} \textbf{(MSQ: select ALL correct options.)} Which statements about the memory hierarchy are correct?
  \begin{enumerate}[label=\Alph*.]
    \item Registers are the fastest storage in the hierarchy.
    \item Cache speeds up access to frequently used data.
    \item Secondary storage is non-volatile.
    \item ROM is volatile working memory.
  \end{enumerate}
\item \textit{[PYQ-adapted: JKSSB Junior Assistant 2026 Q63, final key; SRC-03/SRC-06]} Which memory is typically faster than main memory and holds copies of recently used data?
  \begin{enumerate}[label=\Alph*.]
\item Cache
   \item HDD
   \item Optical disc
   \item Magnetic tape
  \end{enumerate}

\item \textit{[Original]} In the usual speed hierarchy, which is fastest?
  \begin{enumerate}[label=\Alph*.]
\item Registers
   \item RAM
   \item SSD
   \item Magnetic tape
  \end{enumerate}

\item \textit{[Original]} As capacity increases and cost per bit generally falls, memory in a computer is usually:
  \begin{enumerate}[label=\Alph*.]
\item Faster at every level
   \item Slower at each lower level of the hierarchy
   \item Volatile at every level
   \item Removed from the CPU
  \end{enumerate}

\item \textit{[Original]} Which pair is correctly ordered from faster to slower?
  \begin{enumerate}[label=\Alph*.]
\item RAM, registers, cache
   \item Registers, cache, RAM
   \item SSD, RAM, registers
   \item Tape, cache, registers
  \end{enumerate}

\item \textit{[Original, MSQ drill]} \textbf{(MSQ: select ALL correct options.)} Which statements describe the memory hierarchy?
  \begin{enumerate}[label=\Alph*.]
\item Registers are very small and fast.
   \item Cache can reduce average access time.
   \item Secondary storage is generally non-volatile.
   \item Magnetic tape is faster than CPU registers.
  \end{enumerate}

\end{enumerate}
\lessonheading{08}{RAM, ROM \& Cache}
\begin{enumerate}[label=\textbf{Q\arabic*.}, leftmargin=*]
\item \textit{[Original]} Which type of memory loses its contents when the
  computer is switched off?
  \begin{enumerate}[label=\Alph*.]
    \item ROM
    \item Hard disk
    \item RAM
    \item Flash drive
  \end{enumerate}
\item \textit{[Original]} RAM stands for:
  \begin{enumerate}[label=\Alph*.]
    \item Random Access Memory
    \item Read Access Memory
    \item Rapid Access Memory
    \item Read-Only Access Memory
  \end{enumerate}
\item \textit{[Original]} ROM stands for:
  \begin{enumerate}[label=\Alph*.]
    \item Read-Only Memory
    \item Random-Only Memory
    \item Read-Output Memory
    \item Rapid-Only Memory
  \end{enumerate}
\item \textit{[Original]} Which statement is \textbf{NOT} correct?
  \begin{enumerate}[label=\Alph*.]
    \item RAM is volatile.
    \item ROM is non-volatile.
    \item RAM retains its contents after power-off.
    \item ROM is read-only in normal operation.
  \end{enumerate}
\item \textit{[Original]} In the office desktop (EX-01), the main memory that
  holds the running programs and is cleared when the machine is switched off
  is:
  \begin{enumerate}[label=\Alph*.]
    \item ROM
    \item RAM
    \item L1 cache
    \item the SSD
  \end{enumerate}
\item \textit{[Original]} Which of the following is volatile memory?
  \begin{enumerate}[label=\Alph*.]
    \item DRAM
    \item ROM
    \item PROM
    \item EPROM
  \end{enumerate}
\item \textit{[Original]} Which memory is faster and does \textbf{not}
  require periodic refreshing?
  \begin{enumerate}[label=\Alph*.]
    \item DRAM
    \item ROM
    \item Hard disk
    \item SRAM
  \end{enumerate}
\item \textit{[Original]} DRAM differs from SRAM mainly because DRAM:
  \begin{enumerate}[label=\Alph*.]
    \item is faster and needs no refresh
    \item is used for cache memory
    \item is non-volatile
    \item is denser and cheaper but needs periodic refresh
  \end{enumerate}
\item \textit{[Original]} SRAM is commonly used for:
  \begin{enumerate}[label=\Alph*.]
    \item main memory
    \item cache memory
    \item secondary storage
    \item firmware
  \end{enumerate}
\item \textit{[Original]} DRAM is commonly used for:
  \begin{enumerate}[label=\Alph*.]
    \item main memory
    \item cache memory
    \item permanent storage of the start-up code
    \item CPU registers
  \end{enumerate}
\item \textit{[Original]} Which abbreviation--full form pairing is correct?
  \begin{enumerate}[label=\Alph*.]
    \item DRAM --- Direct Random Access Memory
    \item EEPROM --- Electrically Erasable Programmable Read-Only Memory
    \item EPROM --- Electrically Programmable Read-Only Memory
    \item SRAM --- Serial Random Access Memory
  \end{enumerate}
\item \textit{[Original]} Which ROM type can be programmed once by the user
  and cannot be erased?
  \begin{enumerate}[label=\Alph*.]
    \item EPROM
    \item EEPROM
    \item DRAM
    \item PROM
  \end{enumerate}
\item \textit{[Original]} Which ROM type is erased by exposing it to
  ultraviolet light?
  \begin{enumerate}[label=\Alph*.]
    \item PROM
    \item EPROM
    \item EEPROM
    \item Mask ROM
  \end{enumerate}
\item \textit{[Original]} Which memory can be erased and rewritten
  electrically, without removing it from the circuit?
  \begin{enumerate}[label=\Alph*.]
    \item PROM
    \item EPROM
    \item EEPROM
    \item SRAM
  \end{enumerate}
\item \textit{[Original]} Flash memory is best described as:
  \begin{enumerate}[label=\Alph*.]
    \item a volatile form of DRAM
    \item an EEPROM-derived non-volatile memory erased in blocks
    \item the mask-programmed type of ROM
    \item a type of CPU register
  \end{enumerate}
\item \textit{[Original]} Evaluate the statements:
  \begin{enumerate}[label=\Roman*.]
    \item RAM is volatile, while ROM is non-volatile.
    \item ROM is read-only in normal operation, though special programming
      is possible.
    \item Cache memory is slower than main memory.
  \end{enumerate}
  \begin{enumerate}[label=\Alph*.]
    \item I and II only
    \item II and III only
    \item I and III only
    \item I, II, and III
  \end{enumerate}
\item \textit{[Original]} Which cache level is the smallest, fastest, and
  closest to the CPU core?
  \begin{enumerate}[label=\Alph*.]
    \item L1
    \item L2
    \item L3
    \item Main memory
  \end{enumerate}
\item \textit{[Original]} Which cache level is the largest and is typically
  shared across the processor's cores?
  \begin{enumerate}[label=\Alph*.]
    \item L1
    \item L2
    \item ROM
    \item L3
  \end{enumerate}
\item \textit{[Original]} In cache operation, a ``miss'' occurs when:
  \begin{enumerate}[label=\Alph*.]
    \item the required data is found in the cache
    \item data is written to ROM
    \item the required data is not found in the cache
    \item the cache is powered off
  \end{enumerate}
\item \textit{[Original]} A cache handles 100 memory accesses, of which 90
  are hits. The hit ratio is:
  \begin{enumerate}[label=\Alph*.]
    \item 0.10
    \item 1.00
    \item 9.0
    \item 0.90
  \end{enumerate}

\item \textit{[JKSSB FAA-pattern]} RAM is best described as which kind of memory?
  \begin{enumerate}[label=\Alph*.]
    \item Volatile read-write memory
    \item Non-volatile printer memory
    \item Permanent archival tape
    \item Read-only speaker memory
  \end{enumerate}
\item \textit{[JKPSI-pattern]} The startup firmware (BIOS/UEFI code) is stored in:
  \begin{enumerate}[label=\Alph*.]
    \item ROM/Flash memory
    \item RAM working space
    \item Cache tags only
    \item The printer queue
  \end{enumerate}
\item \textit{[SSC CGL-pattern]} Cache memory is normally built from which technology?
  \begin{enumerate}[label=\Alph*.]
    \item SRAM
    \item Magnetic tape
    \item Floppy coating
    \item Paper cards
  \end{enumerate}
\item \textit{[JKSSB-pattern]} \textbf{(MSQ: select ALL correct options.)} Which of the following memories are volatile?
  \begin{enumerate}[label=\Alph*.]
    \item RAM
    \item Cache memory
    \item CPU registers
    \item Masked ROM
  \end{enumerate}
\item \textit{[SSC CGL-pattern]} \textbf{(MSQ: select ALL correct options.)} Which statements about RAM, ROM, and cache are correct?
  \begin{enumerate}[label=\Alph*.]
    \item DRAM needs periodic refreshing.
    \item SRAM is faster than DRAM.
    \item Cache is commonly built with SRAM.
    \item ROM loses its contents when power goes off.
  \end{enumerate}
\item \textit{[PYQ-adapted: JKSSB Sub-Inspector 2022, QCRIAS item 391; SRC-17/SRC-22]} What normally happens to data held only in RAM when power is switched off?
  \begin{enumerate}[label=\Alph*.]
\item It is retained permanently.
   \item It is lost.
   \item It is copied automatically to ROM.
   \item It is stored in the CPU cache.
  \end{enumerate}

\item \textit{[PYQ-adapted: JKSSB FAA 2022, QCRIAS item 551; SRC-17/SRC-21]} Which memory type is traditionally associated with storing startup firmware such as BIOS?
  \begin{enumerate}[label=\Alph*.]
\item RAM
   \item ROM
   \item CPU cache
   \item Register file
  \end{enumerate}

\item \textit{[Original]} Which memory type is volatile?
  \begin{enumerate}[label=\Alph*.]
\item RAM
   \item ROM
   \item SSD
   \item DVD
  \end{enumerate}

\item \textit{[Original]} Which type of RAM is commonly built from flip-flop circuits and does not need periodic refreshing?
  \begin{enumerate}[label=\Alph*.]
\item DRAM
   \item SRAM
   \item ROM
   \item Flash memory
  \end{enumerate}

\item \textit{[Original, MSQ drill]} \textbf{(MSQ: select ALL correct options.)} Which statements are correct?
  \begin{enumerate}[label=\Alph*.]
\item RAM is normally volatile.
   \item ROM is non-volatile.
   \item DRAM requires periodic refresh.
   \item SRAM is slower than DRAM in the memory hierarchy.
  \end{enumerate}

\end{enumerate}
\lessonheading{09}{Virtual Memory}
\begin{enumerate}[label=\textbf{Q\arabic*.}, leftmargin=*]
\item \textit{[Original]} Virtual memory is best described as:
  \begin{enumerate}[label=\Alph*.]
    \item extra physical RAM added to the computer
    \item a technique that uses secondary storage to extend the addressable
      memory beyond physical RAM
    \item a type of cache built into the CPU
    \item a read-only memory used to store the BIOS
  \end{enumerate}
\item \textit{[Original]} Compared with RAM, virtual memory is:
  \begin{enumerate}[label=\Alph*.]
    \item faster
    \item exactly the same speed
    \item slower
    \item always disabled
  \end{enumerate}
\item \textit{[Original]} Which statement about virtual memory is correct?
  \begin{enumerate}[label=\Alph*.]
    \item It is additional physical RAM.
    \item It is disk-backed and slower than RAM.
    \item It is faster than cache.
    \item It removes the need for RAM.
  \end{enumerate}
\item \textit{[Original]} The main reason virtual memory exists is to:
  \begin{enumerate}[label=\Alph*.]
    \item increase the clock speed of the processor
    \item allow programs and data larger than physical RAM to run, while
      giving each process a large contiguous view of memory
    \item replace the CPU with a faster unit
    \item store the firmware permanently
  \end{enumerate}
\item \textit{[Original]} Expand MMU.
  \begin{enumerate}[label=\Alph*.]
    \item Memory Management Unit
    \item Main Memory Utility
    \item Multiprocessor Memory Unit
    \item Machine Mapping Unit
  \end{enumerate}
\item \textit{[Original]} The MMU is responsible for:
  \begin{enumerate}[label=\Alph*.]
    \item translating virtual addresses into physical addresses
    \item storing files permanently on disk
    \item cooling the processor
    \item displaying output on the monitor
  \end{enumerate}
\item \textit{[Original]} In paging, the fixed-size blocks into which a
  program's virtual address space is divided are called:
  \begin{enumerate}[label=\Alph*.]
    \item segments
    \item pages
    \item frames
    \item sectors
  \end{enumerate}
\item \textit{[Original]} In paging, the fixed-size blocks of physical memory
  into which pages are loaded are called:
  \begin{enumerate}[label=\Alph*.]
    \item pages
    \item frames
    \item segments
    \item registers
  \end{enumerate}
\item \textit{[Original]} Which statement correctly distinguishes paging from
  segmentation?
  \begin{enumerate}[label=\Alph*.]
    \item Paging uses variable-size segments; segmentation uses fixed-size
      pages.
    \item Paging divides memory into fixed-size pages and frames; segmentation
      divides a program into variable-size logical segments.
    \item Both paging and segmentation divide memory into equal fixed-size
      blocks.
    \item Segmentation divides memory into equal frames; paging uses
      variable-size logical units.
  \end{enumerate}
\item \textit{[Original]} Segmentation divides a program into:
  \begin{enumerate}[label=\Alph*.]
    \item equal fixed-size pages
    \item variable-size logical segments
    \item physical frames of equal size
    \item cache lines
  \end{enumerate}
\item \textit{[Original]} A page fault occurs when:
  \begin{enumerate}[label=\Alph*.]
    \item a hardware component develops a physical fault
    \item the referenced page is not currently in RAM and must be loaded from
      disk
    \item the contents of RAM become corrupted
    \item the processor overheats
  \end{enumerate}
\item \textit{[Original]} Demand paging loads a page into memory:
  \begin{enumerate}[label=\Alph*.]
    \item only when it is referenced and not already present
    \item only if the whole program is loaded at startup
    \item never, because pages stay on disk
    \item only during the boot process
  \end{enumerate}
\item \textit{[Original]} Which statement about virtual memory is \textbf{NOT}
  correct?
  \begin{enumerate}[label=\Alph*.]
    \item Virtual memory uses secondary storage.
    \item Virtual memory is slower than RAM.
    \item Virtual memory is additional physical RAM.
    \item Virtual memory lets a process use a large contiguous address space.
  \end{enumerate}
\item \textit{[Original]} Evaluate the statements:
  \begin{enumerate}[label=\Roman*.]
    \item Virtual memory uses secondary storage to extend memory.
    \item Virtual memory is slower than RAM.
    \item Virtual memory is extra physical RAM.
  \end{enumerate}
  \begin{enumerate}[label=\Alph*.]
    \item I and II only
    \item I only
    \item II and III only
    \item I, II, and III
  \end{enumerate}
\item \textit{[Original]} Swapping refers to:
  \begin{enumerate}[label=\Alph*.]
    \item moving data or pages between RAM and disk (secondary storage)
    \item exchanging the CPU for a GPU
    \item changing the monitor resolution
    \item moving data from cache to a CPU register
  \end{enumerate}
\item \textit{[Original]} Thrashing is best described as:
  \begin{enumerate}[label=\Alph*.]
    \item excessive paging that leaves little time for useful work
    \item a fast cache hit
    \item a type of segmentation
    \item the automatic addition of extra RAM
  \end{enumerate}
\item \textit{[Original]} Which statement correctly distinguishes cache from
  virtual memory?
  \begin{enumerate}[label=\Alph*.]
    \item Cache is disk-backed; virtual memory sits on the CPU chip.
    \item Cache is small, fast memory close to the CPU; virtual memory uses
      slower, disk-backed storage.
    \item Both cache and virtual memory are forms of secondary storage.
    \item Cache extends RAM using disk; virtual memory speeds up the CPU.
  \end{enumerate}
\item \textit{[Original]} The office desktop (EX-01) has 16\,GB of RAM and
  runs an application whose working set is larger than RAM. What does the
  operating system do?
  \begin{enumerate}[label=\Alph*.]
    \item It refuses to run the application.
    \item It uses virtual memory, keeping some pages on disk and loading them
      on demand.
    \item It converts part of the disk into extra RAM chips.
    \item It raises the clock speed to compensate.
  \end{enumerate}
\item \textit{[Original]} From the perspective of the CPU, which of these is
  the slowest to access?
  \begin{enumerate}[label=\Alph*.]
    \item L1 cache
    \item RAM
    \item disk-backed virtual memory
    \item a CPU register
  \end{enumerate}
\item \textit{[Original]} Which pair is correctly matched?
  \begin{enumerate}[label=\Alph*.]
    \item Paging --- fixed-size pages and frames
    \item Segmentation --- fixed-size pages and frames
    \item Paging --- variable-size logical segments
    \item MMU --- permanent file storage on disk
  \end{enumerate}

\item \textit{[JKSSB FAA-pattern]} Virtual memory extends RAM by using space on the:
  \begin{enumerate}[label=\Alph*.]
    \item hard disk (page/swap file).
    \item printer.
    \item speaker.
    \item keyboard cable.
  \end{enumerate}
\item \textit{[JKPSI-pattern]} The page file (swap file) is stored on the:
  \begin{enumerate}[label=\Alph*.]
    \item hard disk/SSD.
    \item CPU register file.
    \item monitor screen.
    \item mouse sensor.
  \end{enumerate}
\item \textit{[SSC CGL-pattern]} Thrashing occurs when the system:
  \begin{enumerate}[label=\Alph*.]
    \item spends most time swapping pages instead of executing.
    \item prints too many pages at once.
    \item runs only one small program.
    \item keeps the monitor switched off.
  \end{enumerate}
\item \textit{[JKSSB-pattern]} \textbf{(MSQ: select ALL correct options.)} Which of the following relate to virtual memory?
  \begin{enumerate}[label=\Alph*.]
    \item Paging
    \item Swap/page file
    \item Hard-disk space used as memory extension
    \item Printer spooler
  \end{enumerate}
\item \textit{[SSC CGL-pattern]} \textbf{(MSQ: select ALL correct options.)} Which statements about virtual memory are correct?
  \begin{enumerate}[label=\Alph*.]
    \item Paging moves fixed-size blocks between RAM and disk.
    \item A page fault triggers loading the needed page from disk.
    \item Thrashing means excessive paging with little progress.
    \item Virtual memory is a special RAM chip on the motherboard.
  \end{enumerate}
\item \textit{[PYQ-adapted: JKSSB Junior Assistant 2026 Q62, final key; SRC-03/SRC-06]} Virtual memory uses part of which resource to extend the apparent capacity of main memory?
  \begin{enumerate}[label=\Alph*.]
\item Secondary storage
   \item CPU register file
   \item Printer buffer
   \item ROM firmware area
  \end{enumerate}

\item \textit{[Original]} A page currently used by a process is absent from RAM but present in its backing store. The operating system typically responds with a:
  \begin{enumerate}[label=\Alph*.]
\item Page fault
   \item Cache hit
   \item Power failure
   \item Print spool
  \end{enumerate}

\item \textit{[Original]} Virtual memory allows a process to address a space that can be:
  \begin{enumerate}[label=\Alph*.]
\item Larger than the available physical RAM
   \item Stored only in CPU registers
   \item Volatile only when the disk is disconnected
   \item Equal to the size of the cache
  \end{enumerate}

\item \textit{[Original]} Which statement best distinguishes virtual memory from cache?
  \begin{enumerate}[label=\Alph*.]
\item Virtual memory extends the apparent main-memory space; cache speeds up frequent access.
   \item Both are permanent backups.
   \item Cache replaces the operating system.
   \item Virtual memory is a CPU register.
  \end{enumerate}

\item \textit{[Original, MSQ drill]} \textbf{(MSQ: select ALL correct options.)} Which statements about virtual memory are correct?
  \begin{enumerate}[label=\Alph*.]
\item It is managed with operating-system support.
   \item A page fault can occur when a needed page is not in RAM.
   \item Backing storage may be used for non-resident pages.
   \item It makes secondary storage as fast as registers.
  \end{enumerate}

\end{enumerate}
\lessonheading{10}{Storage Devices}
\begin{enumerate}[label=\textbf{Q\arabic*.}, leftmargin=*]
\item \textit{[Original]} Which type of storage is directly accessible by the
  CPU and holds the data currently in use?
  \begin{enumerate}[label=\Alph*.]
    \item Hard disk
    \item DVD
    \item RAM (primary storage)
    \item Magnetic tape
  \end{enumerate}
\item \textit{[Original]} Which of the following is a non-volatile storage
  device, retaining its contents without power?
  \begin{enumerate}[label=\Alph*.]
    \item Hard disk
    \item RAM
    \item Cache
    \item CPU register
  \end{enumerate}
\item \textit{[Original]} Which statement about a hard disk drive (HDD) is
  correct?
  \begin{enumerate}[label=\Alph*.]
    \item It is magnetic and mechanical, with rotating platters and a moving
      read/write head.
    \item It is flash-based and has no moving parts.
    \item It uses a laser to read pits and lands.
    \item It is a sequential-access medium.
  \end{enumerate}
\item \textit{[Original]} Which of the following is a sequential-access
  storage medium?
  \begin{enumerate}[label=\Alph*.]
    \item Magnetic tape
    \item RAM
    \item Hard disk
    \item SSD
  \end{enumerate}
\item \textit{[Original]} Among the following devices, which has the highest
  transfer rate?
  \begin{enumerate}[label=\Alph*.]
    \item Hard disk drive
    \item Solid-state drive (SSD)
    \item CD drive
    \item Magnetic tape
  \end{enumerate}
\item \textit{[Original]} Which option orders the optical media by standard
  single-layer capacity from smallest to largest?
  \begin{enumerate}[label=\Alph*.]
    \item CD $<$ DVD $<$ Blu-ray
    \item DVD $<$ CD $<$ Blu-ray
    \item Blu-ray $<$ DVD $<$ CD
    \item CD $<$ Blu-ray $<$ DVD
  \end{enumerate}
\item \textit{[Original]} The approximate capacity of a standard CD is:
  \begin{enumerate}[label=\Alph*.]
    \item 4.7 GB
    \item 25 GB
    \item 700 MB
    \item 700 GB
  \end{enumerate}
\item \textit{[Original]} A single-layer Blu-ray disc has an approximate
  capacity of:
  \begin{enumerate}[label=\Alph*.]
    \item 25 GB
    \item 700 MB
    \item 4.7 GB
    \item 256 MB
  \end{enumerate}
\item \textit{[Original]} Which statement is \textbf{NOT} correct?
  \begin{enumerate}[label=\Alph*.]
    \item RAM is volatile.
    \item Secondary storage is non-volatile.
    \item A hard disk uses flash memory to store data.
    \item An SSD has no moving parts.
  \end{enumerate}
\item \textit{[Original]} Which device stores data in flash memory with no
  moving parts?
  \begin{enumerate}[label=\Alph*.]
    \item Magnetic tape
    \item CD-ROM
    \item Floppy disk
    \item SSD
  \end{enumerate}
\item \textit{[Original]} What is the full form of SSD?
  \begin{enumerate}[label=\Alph*.]
    \item Secondary Storage Device
    \item Solid-State Drive
    \item Sequential Storage Disk
    \item System Software Drive
  \end{enumerate}
\item \textit{[Original]} Which storage level is best described as an
  archival level, often using robotic tape libraries loaded on demand?
  \begin{enumerate}[label=\Alph*.]
    \item Cache
    \item CPU register
    \item RAM
    \item Tertiary storage
  \end{enumerate}
\item \textit{[Original]} Evaluate the statements:
  \begin{enumerate}[label=\Roman*.]
    \item RAM is volatile.
    \item Secondary storage is non-volatile.
    \item A hard disk drive has no moving parts.
  \end{enumerate}
  \begin{enumerate}[label=\Alph*.]
    \item I only
    \item I and II only
    \item II and III only
    \item I, II, and III
  \end{enumerate}
\item \textit{[Original]} Cloud storage means that data is:
  \begin{enumerate}[label=\Alph*.]
    \item Stored only in local RAM
    \item Stored only on a CD
    \item Stored on remote servers and accessed over the Internet
    \item Stored in CPU registers
  \end{enumerate}
\item \textit{[Original]} Which statement correctly distinguishes an HDD
  from an SSD?
  \begin{enumerate}[label=\Alph*.]
    \item The HDD is flash-based, while the SSD is magnetic.
    \item Both devices use rotating platters.
    \item The SSD is slower and more shock-sensitive than the HDD.
    \item The HDD is magnetic and mechanical, while the SSD is flash-based
      with no moving parts.
  \end{enumerate}
\item \textit{[Original]} What is the full form of HDD?
  \begin{enumerate}[label=\Alph*.]
    \item High Density Disk
    \item High Definition Drive
    \item Hard Disk Drive
    \item Hard Data Device
  \end{enumerate}
\item \textit{[Original]} Which access method do RAM and hard disks commonly
  share?
  \begin{enumerate}[label=\Alph*.]
    \item Sequential access
    \item Random access
    \item Only direct memory access
    \item Purely optical access
  \end{enumerate}
\item \textit{[Original]} Magnetic tape stores and retrieves data:
  \begin{enumerate}[label=\Alph*.]
    \item Randomly
    \item Sequentially
    \item Only in RAM
    \item With a laser
  \end{enumerate}
\item \textit{[Original]} Which pairing of storage device and technology is
  correct?
  \begin{enumerate}[label=\Alph*.]
    \item CD -- magnetic
    \item Blu-ray -- flash
    \item Magnetic tape -- optical
    \item SSD -- flash memory
  \end{enumerate}
\item \textit{[Original]} An office desktop needs a fast boot drive with no
  moving parts and a high transfer rate. Which device is best suited as its
  system drive?
  \begin{enumerate}[label=\Alph*.]
    \item CD-ROM
    \item Magnetic tape
    \item SSD
    \item Floppy disk
  \end{enumerate}

\item \textit{[JKSSB FAA-pattern]} Which of the following is an optical storage device?
  \begin{enumerate}[label=\Alph*.]
    \item DVD
    \item Hard disk
    \item RAM stick
    \item Register file
  \end{enumerate}
\item \textit{[JKPSI-pattern]} A 1.44 MB floppy is an example of which storage class?
  \begin{enumerate}[label=\Alph*.]
    \item Magnetic storage
    \item Optical storage
    \item Semiconductor register storage
    \item Network cloud storage
  \end{enumerate}
\item \textit{[SSC CGL-pattern]} Which device has no moving mechanical parts?
  \begin{enumerate}[label=\Alph*.]
    \item SSD
    \item Hard disk drive
    \item DVD drive
    \item Magnetic tape drive
  \end{enumerate}
\item \textit{[JKSSB-pattern]} \textbf{(MSQ: select ALL correct options.)} Which of the following are secondary storage devices?
  \begin{enumerate}[label=\Alph*.]
    \item Hard disk drive
    \item Solid-state drive
    \item DVD
    \item RAM
  \end{enumerate}
\item \textit{[SSC CGL-pattern]} \textbf{(MSQ: select ALL correct options.)} Which statements about storage devices are correct?
  \begin{enumerate}[label=\Alph*.]
    \item A hard disk stores data magnetically.
    \item An SSD stores data in flash memory.
    \item Optical discs are read with a laser.
    \item CPU registers are secondary storage.
  \end{enumerate}
\item \textit{[PYQ-adapted: JKSSB Website Operator 2026 Q67, provisional key; SRC-04/SRC-07]} Which listed secondary storage device generally has the highest transfer rate?
  \begin{enumerate}[label=\Alph*.]
\item Magnetic tape
   \item HDD
   \item SSD
   \item CD-ROM
  \end{enumerate}

\item \textit{[PYQ-adapted: JKSSB FAA 2022, QCRIAS item 550; SRC-21]} Hard-disk rotational speed is commonly specified in:
  \begin{enumerate}[label=\Alph*.]
\item RPM
   \item Pixels per inch
   \item Hertz of the CPU clock
   \item Bits per byte
  \end{enumerate}

\item \textit{[Original]} Which medium is optical storage?
  \begin{enumerate}[label=\Alph*.]
\item DVD
   \item SSD
   \item HDD
   \item Magnetic tape
  \end{enumerate}

\item \textit{[Original]} Which storage medium is most commonly associated with sequential access and archival backup?
  \begin{enumerate}[label=\Alph*.]
\item Magnetic tape
   \item CPU cache
   \item RAM
   \item Register
  \end{enumerate}

\item \textit{[Original, MSQ drill]} \textbf{(MSQ: select ALL correct options.)} Which are non-volatile secondary storage media?
  \begin{enumerate}[label=\Alph*.]
\item SSD
   \item HDD
   \item DVD
   \item CPU register
  \end{enumerate}

\end{enumerate}
\lessonheading{11}{Backup Types}
\begin{enumerate}[label=\textbf{Q\arabic*.}, leftmargin=*]
\item \textit{[Original]} Which backup type copies all selected files and
  folders every time it runs?
  \begin{enumerate}[label=\Alph*.]
    \item Incremental
    \item Differential
    \item Full
    \item Image
  \end{enumerate}
\item \textit{[Original]} An incremental backup captures changes since the:
  \begin{enumerate}[label=\Alph*.]
    \item last full backup
    \item last backup of any type
    \item last deletion only
    \item start of the operating system
  \end{enumerate}
\item \textit{[Original]} A differential backup captures changes since the:
  \begin{enumerate}[label=\Alph*.]
    \item last backup of any type
    \item last hour
    \item last full backup
    \item last new file only
  \end{enumerate}
\item \textit{[Original]} Which statement correctly distinguishes an
  incremental backup from a differential backup?
  \begin{enumerate}[label=\Alph*.]
    \item Incremental captures changes since the last full backup;
      differential captures changes since the last backup of any type.
    \item Incremental captures changes since the last backup of any type;
      differential captures changes since the last full backup.
    \item Both capture changes since the last full backup.
    \item Both capture only the operating system.
  \end{enumerate}
\item \textit{[Original]} Restoring from an incremental chain typically
  requires:
  \begin{enumerate}[label=\Alph*.]
    \item the last full backup plus every incremental backup in sequence
    \item only the most recent incremental backup
    \item only the most recent differential backup
    \item the image backup alone
  \end{enumerate}
\item \textit{[Original]} Restoring from a differential chain typically
  requires:
  \begin{enumerate}[label=\Alph*.]
    \item the last full backup plus every differential backup
    \item the last full backup plus the most recent differential backup
    \item only the most recent differential backup
    \item only the incremental backups
  \end{enumerate}
\item \textit{[Original]} Which backup type usually takes the most time and
  storage on each run?
  \begin{enumerate}[label=\Alph*.]
    \item Incremental
    \item Differential
    \item Full
    \item None of these
  \end{enumerate}
\item \textit{[Original]} Which backup type is heaviest to create but
  fastest to restore because only one complete set is needed?
  \begin{enumerate}[label=\Alph*.]
    \item Incremental
    \item Differential
    \item Full
    \item File-level
  \end{enumerate}
\item \textit{[Original]} An image (system image) backup is best described
  as:
  \begin{enumerate}[label=\Alph*.]
    \item a copy of only the documents folder
    \item a complete clone of a system, including the OS, settings and data
    \item a copy of changes since the last full backup
    \item a copy of only the most recently deleted files
  \end{enumerate}
\item \textit{[Original]} A file-level backup:
  \begin{enumerate}[label=\Alph*.]
    \item clones the entire disk including the boot sector
    \item backs up selected files and folders rather than the whole system
    \item backs up only the operating system
    \item is another name for RAID
  \end{enumerate}
\item \textit{[Original]} Which of the following is \textbf{NOT} a backup
  type discussed in this lesson?
  \begin{enumerate}[label=\Alph*.]
    \item Full
    \item Incremental
    \item Differential
    \item Defragmentation
  \end{enumerate}
\item \textit{[Original]} The 3-2-1 backup rule means:
  \begin{enumerate}[label=\Alph*.]
    \item 3 copies, 2 different media, 1 copy off-site
    \item 3 media, 2 off-site copies, 1 cloud account
    \item 3 backups per week, 2 per month, 1 per year
    \item 3 full backups, 2 differential, 1 incremental
  \end{enumerate}
\item \textit{[Original]} In the 3-2-1 rule, the number 2 refers to:
  \begin{enumerate}[label=\Alph*.]
    \item two off-site copies
    \item two different types of storage media
    \item two backups per day
    \item two encryption keys
  \end{enumerate}
\item \textit{[Original]} Expand RAID and state its correct relation to
  backup.
  \begin{enumerate}[label=\Alph*.]
    \item Redundant Array of Independent Disks; it provides redundancy but is
      not a backup.
    \item Random Access Internal Disk; it is a backup type.
    \item Redundant Archive of Integrated Disks; it encrypts backups.
    \item Read-Ahead Interface Device; it schedules backups.
  \end{enumerate}
\item \textit{[Original]} Evaluate the statements:
  \begin{enumerate}[label=\Roman*.]
    \item A full backup copies everything selected each time it runs.
    \item A differential backup copies changes since the last full backup.
    \item An incremental backup copies changes since the last full backup.
  \end{enumerate}
  \begin{enumerate}[label=\Alph*.]
    \item I and II only
    \item I and III only
    \item II and III only
    \item I, II and III
  \end{enumerate}
\item \textit{[Original]} Which statement best distinguishes a backup from
  ordinary storage and an archive?
  \begin{enumerate}[label=\Alph*.]
    \item A backup is a recoverable copy; storage is the working copy; an
      archive is old data kept for long-term retention.
    \item Backup, storage and archive are three names for the same thing.
    \item A backup is the file currently being edited.
    \item An archive is a second disk used only for RAID.
  \end{enumerate}
\item \textit{[Original]} Cloud backup is best described as:
  \begin{enumerate}[label=\Alph*.]
    \item storing backup copies on remote servers, giving an off-site copy
    \item the same as RAID 1
    \item storing the backup only on the same local disk
    \item a synonym for an incremental backup
  \end{enumerate}
\item \textit{[Original]} An offline backup is one that is:
  \begin{enumerate}[label=\Alph*.]
    \item disconnected from the network, which helps against ransomware
    \item always connected and updated instantly
    \item not created at all
    \item stored only in main memory
  \end{enumerate}
\item \textit{[Original]} Backup rotation (for example, the
  grandfather--father--son scheme) is used mainly to:
  \begin{enumerate}[label=\Alph*.]
    \item rotate the CPU clock speed
    \item maintain several dated versions so an older clean copy remains
      available
    \item encrypt files automatically
    \item replace full backups with RAID
  \end{enumerate}
\item \textit{[Original]} Which pairing is correct?
  \begin{enumerate}[label=\Alph*.]
    \item Differential --- changes since the last backup of any type
    \item Incremental --- changes since the last full backup
    \item Incremental --- changes since the last backup of any type
    \item Image --- only selected files and folders
  \end{enumerate}

\item \textit{[JKSSB FAA-pattern]} Which backup copies only the data changed since the last backup of any kind?
  \begin{enumerate}[label=\Alph*.]
    \item Incremental backup
    \item Full backup
    \item Printer test page
    \item Disk defragmentation
  \end{enumerate}
\item \textit{[JKPSI-pattern]} Restoring from a full-plus-incremental chain needs the last full backup plus:
  \begin{enumerate}[label=\Alph*.]
    \item every incremental since that full backup.
    \item only the printer driver.
    \item only the first file of the disk.
    \item none of the incrementals.
  \end{enumerate}
\item \textit{[SSC CGL-pattern]} RAID is primarily used to provide:
  \begin{enumerate}[label=\Alph*.]
    \item redundancy and/or faster disk access.
    \item louder speaker output.
    \item brighter monitor colours.
    \item faster keyboard typing.
  \end{enumerate}
\item \textit{[JKSSB-pattern]} \textbf{(MSQ: select ALL correct options.)} Which of the following are recognised backup types?
  \begin{enumerate}[label=\Alph*.]
    \item Full backup
    \item Incremental backup
    \item Differential backup
    \item Disk formatting
  \end{enumerate}
\item \textit{[SSC CGL-pattern]} \textbf{(MSQ: select ALL correct options.)} Which statements about backups are correct?
  \begin{enumerate}[label=\Alph*.]
    \item A full backup copies all selected data.
    \item A differential backup copies changes since the last full backup.
    \item An incremental backup copies changes since the last backup of any kind.
    \item Backup and file synchronisation mean exactly the same thing.
  \end{enumerate}
\item \textit{[Original]} Which backup type copies all selected data each time it runs?
  \begin{enumerate}[label=\Alph*.]
\item Full backup
   \item Incremental backup
   \item Differential backup
   \item Transaction log only
  \end{enumerate}

\item \textit{[Original]} After a full backup, an incremental backup normally copies files changed since:
  \begin{enumerate}[label=\Alph*.]
\item The previous backup of any kind
   \item The previous full backup only
   \item The computer was purchased
   \item The last restore test
  \end{enumerate}

\item \textit{[Original]} Which backup usually grows cumulatively until the next full backup?
  \begin{enumerate}[label=\Alph*.]
\item Differential
   \item Incremental
   \item Mirror deletion
   \item Temporary cache
  \end{enumerate}

\item \textit{[Original]} The 3-2-1 backup rule recommends keeping:
  \begin{enumerate}[label=\Alph*.]
\item Three copies, on two media types, with one copy off-site
   \item Three files, on one device, with two deleted
   \item Two copies only, both online
   \item One copy in RAM and two in cache
  \end{enumerate}

\item \textit{[Original, MSQ drill]} \textbf{(MSQ: select ALL correct options.)} Which practices improve backup resilience?
  \begin{enumerate}[label=\Alph*.]
\item Keep an off-site copy.
   \item Test restoration from backups.
   \item Keep at least one copy isolated from routine access.
   \item Treat RAID as a complete substitute for backup.
  \end{enumerate}

\end{enumerate}
\answersection{3}
\answersubheading{07}{Computer Memory Hierarchy}
\begin{enumerate}[label=\textbf{Q\arabic*.}, leftmargin=*]
\item \textbf{B.} RAM stands for Random Access Memory. A, C, and D are not
  expansions of the abbreviation.
\item \textbf{C.} ROM stands for Read-Only Memory. A, B, and D are not the
  correct full forms.
\item \textbf{A.} Speed decreases from registers to cache to main memory
  (RAM) to secondary storage. B swaps the top two levels; C and D scramble
  the order.
\item \textbf{D.} Registers are the fastest storage in the computer. L1
  cache (B), RAM (A), and an SSD (C) are each progressively slower.
\item \textbf{B.} Cache is faster than main memory. It is slower than RAM is
  false (A); cache is smaller, not larger (C); and cache is more expensive
  per bit, not cheaper (D).
\item \textbf{C.} Cache is smaller and faster than RAM, while RAM is larger
  and slower. A reverses the size relation; B claims they are identical; D
  wrongly calls cache a backup.
\item \textbf{A.} L1 is the smallest and fastest cache level. L2 and L3 are
  successively larger and slower; L4 is not a standard cache level.
\item \textbf{D.} L3 is the largest and slowest of the three cache levels,
  though still faster than main memory. L1 (A) is smallest and fastest; L2
  (B) is the middle level; registers (C) are not a cache level.
\item \textbf{C.} Cache is not a permanent backup; it holds copies of data
  that live in RAM. Storing recent data (A), exploiting locality (B), and
  being faster than main memory (D) are all true.
\item \textbf{B.} RAM is volatile and ROM is non-volatile. A reverses the
  pair; C describes ROM as read-only and RAM as writable instead of the
  volatility distinction; D misclassifies both as storage types.
\item \textbf{A.} RAM is volatile and loses its contents when the power is
  switched off. ROM (B), the hard disk drive (C), and the solid-state drive
  (D) are all non-volatile.
\item \textbf{D.} RAM is primary memory, directly accessible by the CPU. An
  SSD (A), an HDD (B), and magnetic tape (C) are secondary or tertiary
  storage, not primary memory.
\item \textbf{C.} A solid-state drive is secondary storage. A register (A),
  cache (B), and RAM (D) are primary (CPU-side or main memory), not
  secondary storage.
\item \textbf{A.} The hierarchy exists because faster memory is smaller and
  more expensive per bit, so a layered arrangement balances speed, cost, and
  capacity. B claims all memory is equal; C names the slowest level as
  fastest; D overstates the size of registers.
\item \textbf{B.} Statements I and II are correct: registers are the fastest
  storage, and cache is faster than main memory. III is false because RAM is
  volatile. A omits II; C includes false III and omits I; D includes false
  III.
\item \textbf{A.} Reusing recently accessed data and nearby data is locality
  of reference. Overflow (B), fragmentation (C), and deadlock (D) are
  unrelated concepts.
\item \textbf{C.} Temporal locality is reusing the same data or instruction
  again soon after it was accessed. A describes spatial locality; B describes
  persistence on storage; D describes volatility.
\item \textbf{D.} Registers are the smallest storage. L1 cache (C) is the
  next smallest; RAM (B) and an SSD (A) are larger.
\item \textbf{B.} Secondary storage such as an SSD or HDD is non-volatile
  and keeps files between sessions. Registers (A), cache (C), and RAM (D) do
  not retain their contents without power.
\item \textbf{D.} An SSD is secondary storage that retains data without
  power, so this pairing is correct. A mislabels registers as slowest (they
  are fastest); B wrongly calls cache a non-volatile backup; C wrongly calls
  RAM non-volatile and read-only.

\item \textbf{A.} Registers sit at the top of the hierarchy and are fastest. B, C, and D are slow secondary media.
\item \textbf{A.} Speed falls from registers to cache to RAM to hard disk. B, C, and D scramble that order.
\item \textbf{A.} Cache bridges the CPU--RAM speed gap. Tape (B), floppy (C), and printer hardware (D) do not.
\item \textbf{A, B, C.} RAM, ROM, and cache are primary memory. The hard disk (D) is secondary storage.
\item \textbf{A, B, C.} Register speed, cache behaviour, and secondary-storage persistence are correct. D is false: ROM is non-volatile.
\item \textbf{A.} Cache sits between CPU and RAM to reduce access time. HDDs, optical discs, and tape are secondary storage and are slower.

\item \textbf{A.} Registers are closest to the CPU and fastest. RAM, SSD, and tape are progressively slower storage levels.

\item \textbf{B.} Larger, cheaper storage is generally slower, which motivates the hierarchy. Not all levels are volatile, and memory is not removed from the CPU.

\item \textbf{B.} The usual order is registers, cache, then RAM. The other options reverse one or more levels.

\item \textbf{A, B, C.} Registers are small/fast, cache reduces access time, and secondary storage retains data without power. Tape is much slower than registers.

\end{enumerate}
\answersubheading{08}{RAM, ROM \& Cache}
\begin{enumerate}[label=\textbf{Q\arabic*.}, leftmargin=*]
\item \textbf{C.} RAM is volatile, so it loses its contents when power is
  removed. ROM, a hard disk, and a flash drive are non-volatile and retain
  their contents.
\item \textbf{A.} RAM stands for Random Access Memory. B, C, and D are
  invented expansions.
\item \textbf{A.} ROM stands for Read-Only Memory. B, C, and D are invented
  expansions.
\item \textbf{C.} RAM is volatile and does \emph{not} retain its contents
  after power-off, so C is the statement that is NOT correct. A, B, and D
  are all true statements.
\item \textbf{B.} RAM is the volatile main memory holding running programs
  and their data. ROM is non-volatile; L1 cache is smaller working memory;
  the SSD is secondary storage.
\item \textbf{A.} DRAM is a volatile form of RAM. ROM, PROM, and EPROM
  belong to the non-volatile ROM family.
\item \textbf{D.} SRAM is faster than DRAM and needs no refresh, which is
  why it is used for cache. DRAM needs periodic refresh; ROM and a hard disk
  are non-volatile and slower than SRAM.
\item \textbf{D.} DRAM is denser and cheaper than SRAM but needs periodic
  refresh. A describes SRAM; B is the use of SRAM; C is false, since DRAM is
  volatile.
\item \textbf{B.} SRAM is used for cache memory because it is fast and needs
  no refresh. A is DRAM's role; C and D describe other memory classes.
\item \textbf{A.} DRAM serves as main memory: dense and cheap, at the cost of
  periodic refresh. B is SRAM's role; C describes ROM; D is a register, not
  DRAM.
\item \textbf{B.} EEPROM is Electrically Erasable Programmable Read-Only
  Memory. A should be Dynamic, C should be Erasable (UV), and D should be
  Static, so each of the others is a wrong expansion.
\item \textbf{D.} PROM is programmable once and cannot be erased. EPROM and
  EEPROM can be erased; DRAM is a volatile RAM, not a ROM type.
\item \textbf{B.} EPROM is erased by ultraviolet light. PROM cannot be
  erased, EEPROM is erased electrically, and mask ROM is fixed at
  manufacture.
\item \textbf{C.} EEPROM is erased and rewritten electrically, in-circuit.
  PROM cannot be erased, EPROM needs UV light plus removal, and SRAM is
  volatile RAM.
\item \textbf{B.} Flash is an EEPROM-derived non-volatile memory erased in
  blocks. A wrongly makes it volatile DRAM; C describes a fixed ROM type; D
  is a register, not flash.
\item \textbf{A.} Statements I and II are correct: RAM is volatile and ROM
  is non-volatile, and ROM is normally read-only though specially
  programmable. III is false because cache is faster than main memory, not
  slower. B and C each include false III; D includes all three.
\item \textbf{A.} L1 is the smallest, fastest cache and sits closest to the
  core. L2 and L3 are larger and slower; main memory is slower still.
\item \textbf{D.} L3 is the largest cache level and is typically shared
  across cores. L1 and L2 are smaller and faster; ROM is not cache.
\item \textbf{C.} A miss means the data is not found in the cache, so it is
  fetched from main memory. A describes a hit; B and D are unrelated.
\item \textbf{D.} Hit ratio $=$ hits $\div$ total accesses $= 90/100 = 0.90$.
  A is the miss fraction, B would mean every access hit, and C is
  arithmetically wrong.

\item \textbf{A.} RAM is volatile and readable/writable. B, C, and D wrongly make it permanent or tie it to peripherals.
\item \textbf{A.} Firmware lives in non-volatile ROM/Flash. RAM (B) loses contents; C and D are not firmware stores.
\item \textbf{A.} SRAM's speed suits cache. Tape (B), floppy coating (C), and paper (D) are not cache technologies.
\item \textbf{A, B, C.} RAM, cache, and registers lose contents without power. Masked ROM (D) is non-volatile.
\item \textbf{A, B, C.} DRAM refresh, SRAM speed, and SRAM-based cache are correct. D is false: ROM keeps contents without power.
\item \textbf{B.} RAM is volatile, so its contents are lost when power is removed. It is not automatically copied to ROM or cache.

\item \textbf{B.} Firmware must be available at startup and stored in non-volatile memory, traditionally ROM. The other listed memories are not the standard firmware store.

\item \textbf{A.} RAM loses its contents without power. ROM, SSDs, and optical discs are non-volatile.

\item \textbf{B.} SRAM uses flip-flop cells and does not require periodic refresh. DRAM uses capacitors and must be refreshed; ROM and flash are non-volatile.

\item \textbf{A, B, C.} RAM is normally volatile, ROM retains data, and DRAM needs refresh. SRAM is generally faster than DRAM, not slower.

\end{enumerate}
\answersubheading{09}{Virtual Memory}
\begin{enumerate}[label=\textbf{Q\arabic*.}, leftmargin=*]
\item \textbf{B.} Virtual memory uses secondary storage to extend the
  addressable memory beyond physical RAM. A wrongly calls it extra RAM; C
  describes cache; D describes ROM.
\item \textbf{C.} Virtual memory is slower than RAM because it is disk-backed.
  A and B overstate its speed; D is false, since virtual memory is normally in
  use.
\item \textbf{B.} Virtual memory is disk-backed and slower than RAM, extending
  capacity rather than adding RAM. A contradicts this; C is false because cache
  is faster; D is false because RAM is still required.
\item \textbf{B.} Virtual memory lets programs and data larger than physical
  RAM run and gives each process a large contiguous view. A, C, and D describe
  unrelated hardware changes.
\item \textbf{A.} MMU stands for Memory Management Unit. B, C, and D are
  invented expansions.
\item \textbf{A.} The MMU translates virtual addresses into physical addresses
  at run time. B is the job of secondary storage; C and D describe other
  components.
\item \textbf{B.} In paging, the fixed-size virtual blocks are pages. Frames
  (C) are the physical blocks; segments (A) belong to segmentation; sectors (D)
  are disk units.
\item \textbf{B.} The fixed-size physical blocks into which pages are loaded
  are frames. Pages (A) are the virtual blocks; segments (C) belong to
  segmentation; registers (D) are CPU storage.
\item \textbf{B.} Paging uses fixed-size pages and frames, while segmentation
  uses variable-size logical segments. A and D reverse the two; C wrongly makes
  both fixed-size.
\item \textbf{B.} Segmentation divides a program into variable-size logical
  segments. A, C, and D describe paging or cache rather than segmentation.
\item \textbf{B.} A page fault means the referenced page is not currently in
  RAM and must be loaded from disk. A calls it a hardware fault, which is
  wrong; C and D describe unrelated failures.
\item \textbf{A.} Demand paging loads a page only when it is referenced and not
  already present. B describes loading everything at once; C and D are false.
\item \textbf{C.} Virtual memory is \emph{not} additional physical RAM; it is
  disk-backed and slower than RAM. A, B, and D are correct statements, so C is
  the one that is NOT correct.
\item \textbf{A.} Statements I and II are correct; III is false because virtual
  memory is not extra physical RAM but disk-backed storage. B omits II; C
  includes false III and omits I; D includes false III.
\item \textbf{A.} Swapping moves data or pages between RAM and secondary
  storage. B, C, and D describe unrelated operations.
\item \textbf{A.} Thrashing is excessive paging that leaves little time for
  useful work. B, C, and D misdescribe it.
\item \textbf{B.} Cache is small, fast memory close to the CPU, whereas virtual
  memory is slower and disk-backed. A reverses the two; C wrongly makes both
  secondary storage; D swaps their purposes.
\item \textbf{B.} When the working set exceeds RAM, the OS uses virtual memory,
  keeping some pages on disk and loading them on demand. A is false; C is
  impossible (disk is not converted into RAM chips); D does not address memory
  pressure.
\item \textbf{C.} Disk-backed virtual memory is the slowest of these. L1 cache
  (A) and a CPU register (D) are fastest, and RAM (B) is faster than the disk.
\item \textbf{A.} Paging uses fixed-size pages and frames. B and C mismatch the
  block rule; D wrongly assigns file storage to the MMU.

\item \textbf{A.} The page/swap file on disk backs virtual memory. B, C, and D are peripherals, not memory backing.
\item \textbf{A.} The page file resides on disk storage. B, C, and D are not file stores.
\item \textbf{A.} Thrashing is pathological paging overhead. B, C, and D describe unrelated situations.
\item \textbf{A, B, C.} Paging, the swap file, and disk backing all belong to virtual memory. The spooler (D) serves printing.
\item \textbf{A, B, C.} Fixed-size paging, page-fault loading, and the thrashing definition are correct. D is false: virtual memory is a technique, not a chip.
\item \textbf{A.} Virtual memory uses storage space, usually on an SSD/HDD, as backing for memory pages. Registers, printer buffers, and firmware ROM do not provide this backing store.

\item \textbf{A.} A page fault signals that the referenced page is not resident in RAM and must be fetched. A cache hit, power failure, and print spool are different events.

\item \textbf{A.} Virtual addressing can provide a larger apparent address space than physical RAM. Registers and cache do not define the process's full virtual address space.

\item \textbf{A.} Virtual memory uses storage as backing for memory pages, while cache keeps copies of data for speed. Neither is a permanent backup, OS replacement, or register.

\item \textbf{A, B, C.} The OS manages virtual memory, faults locate missing pages, and backing storage can hold non-resident pages. Storage remains much slower than registers.

\end{enumerate}
\answersubheading{10}{Storage Devices}
\begin{enumerate}[label=\textbf{Q\arabic*.}, leftmargin=*]
\item \textbf{C.} Primary storage such as RAM is directly accessible by the
  CPU and holds what is currently in use. Hard disks, DVDs, and magnetic tape
  are secondary storage reached through main memory.
\item \textbf{A.} A hard disk retains its contents without power, so it is
  non-volatile. RAM, cache, and CPU registers are all volatile and lose their
  contents when power is removed.
\item \textbf{A.} An HDD stores data magnetically on rotating platters read
  and written by a moving head; it is mechanical and has moving parts. B
  describes an SSD, C describes an optical disc, and D describes tape.
\item \textbf{A.} Magnetic tape is read and written in order, so it is
  sequential. RAM, a hard disk, and an SSD offer approximately random access
  to any location.
\item \textbf{B.} The SSD is flash-based with no moving parts, so it reaches
  a high transfer rate faster than a mechanical hard disk, an optical CD
  drive, or magnetic tape (KB PYQ-27).
\item \textbf{A.} The standard single-layer capacities are CD $\approx$
  700 MB, DVD $\approx$ 4.7 GB, and Blu-ray $\approx$ 25 GB, so the order is
  CD $<$ DVD $<$ Blu-ray. B, C, and D scramble this sequence.
\item \textbf{C.} A standard CD holds approximately 700 MB. 4.7 GB is a DVD,
  25 GB is a single-layer Blu-ray, and 700 GB is far too large for any CD.
\item \textbf{A.} A single-layer Blu-ray disc holds approximately 25 GB.
  700 MB is a CD, 4.7 GB is a DVD, and 256 MB is smaller than a CD.
\item \textbf{C.} A hard disk stores data magnetically, not in flash memory;
  the statement is false. A, B, and D are all correct statements.
\item \textbf{D.} An SSD uses flash memory and has no moving parts. Magnetic
  tape is magnetic, and a CD-ROM is optical; a floppy disk is magnetic and
  has a rotating platter (and is largely obsolete).
\item \textbf{B.} SSD stands for Solid-State Drive. A, C, and D are invented
  expansions that do not match the standard full form.
\item \textbf{D.} Tertiary storage is the archival level, often implemented
  as robotic tape libraries. Cache, CPU registers, and RAM are primary
  storage levels, not archival.
\item \textbf{B.} Statements I and II are correct. Statement III is false
  because a hard disk has moving parts (spinning platters and a moving head),
  so A omits II, C includes the false III, and D includes the false III.
\item \textbf{C.} Cloud storage holds data on remote servers reached over the
  Internet. A, B, and D wrongly limit the data to local memory, a CD, or CPU
  registers.
\item \textbf{D.} The HDD is magnetic and mechanical with moving parts, while
  the SSD is flash-based with no moving parts. A reverses the technologies,
  B is false for both, and C reverses the speed and shock-resistance facts.
\item \textbf{C.} HDD stands for Hard Disk Drive. A, B, and D are incorrect
  expansions.
\item \textbf{B.} RAM and hard disks both allow any location to be reached in
  approximately equal time, which is random (direct) access. Sequential
  access belongs to tape, and the other options are not access methods.
\item \textbf{B.} Magnetic tape is read and written in order, which is
  sequential access. It is neither random nor optical, and data is not stored
  ``only in RAM.''
\item \textbf{D.} The SSD uses flash memory. A is wrong because a CD is
  optical, B is wrong because Blu-ray is optical, and C is wrong because
  magnetic tape is magnetic, not optical.
\item \textbf{C.} An SSD is flash-based with no moving parts and a high
  transfer rate, ideal for a fast system drive. A CD-ROM and magnetic tape
  are slow, and a floppy disk is obsolete and low-capacity.

\item \textbf{A.} DVDs are read optically. The hard disk (B) is magnetic; C and D are primary/register memory.
\item \textbf{A.} Floppies store data magnetically. B, C, and D name other technologies.
\item \textbf{A.} SSDs use flash with no moving parts. B, C, and D all move media or heads.
\item \textbf{A, B, C.} HDDs, SSDs, and DVDs retain data as secondary storage. RAM (D) is primary volatile memory.
\item \textbf{A, B, C.} Magnetic HDDs, flash SSDs, and laser-read optical discs are correct. D is false: registers are primary CPU storage.
\item \textbf{C.} An SSD has no moving read/write head and generally transfers data faster than HDDs, tape, or optical discs. The cited key is provisional.

\item \textbf{A.} RPM means revolutions per minute. The alternatives measure display density, CPU frequency, and data grouping.

\item \textbf{A.} A DVD stores data optically. SSDs use flash, HDDs use magnetic platters, and tape is magnetic sequential storage.

\item \textbf{A.} Magnetic tape is a removable sequential-access medium commonly used for backup/archive. The others are volatile CPU/main memory.

\item \textbf{A, B, C.} SSDs, HDDs, and DVDs retain data without power. CPU registers are volatile internal storage.

\end{enumerate}
\answersubheading{11}{Backup Types}
\begin{enumerate}[label=\textbf{Q\arabic*.}, leftmargin=*]
\item \textbf{C.} A full backup copies all selected files and folders on every
  run. Incremental (A) and differential (B) copy only changes; an image (D)
  is a whole-system clone, a coverage type rather than the ``copy everything
  each run'' definition.
\item \textbf{B.} An incremental backup captures changes since the last backup
  of any type. A describes differential; C is not a backup rule; D is not a
  basis for change tracking.
\item \textbf{C.} A differential backup captures changes since the last full
  backup. A describes incremental; B is an arbitrary interval; D excludes
  modified files, which is wrong.
\item \textbf{B.} Incremental looks back to the immediately preceding backup
  of any type; differential always looks back to the last full backup. A
  swaps the two definitions; C wrongly makes them identical; D ignores that
  both back up user data, not just the OS.
\item \textbf{A.} An incremental restore needs the last full backup plus every
  incremental applied in sequence, because each holds only its own slice of
  changes. B and C would lose intermediate changes; D ignores the chain
  entirely.
\item \textbf{B.} A differential restore needs the last full backup plus only
  the most recent differential, which already contains all changes since that
  full backup. A overstates the requirement; C omits the full baseline; D
  names the wrong type.
\item \textbf{C.} A full backup repeats all selected data each run, so it is
  the largest and slowest to create. Incremental (A) is the smallest and
  fastest; differential (B) is in between.
\item \textbf{C.} The full backup is heaviest to create but fastest to
  restore, since one complete set suffices. Incremental (A) is fast to create
  but slow to restore; differential (B) is intermediate; file-level (D) does
  not imply a single complete restore set.
\item \textbf{B.} An image is a complete clone of a system, including OS,
  settings and data. A and D describe narrow copies; C describes a
  differential, not an image.
\item \textbf{B.} File-level backup copies selected files and folders rather
  than a whole system. A describes an image; C narrows it incorrectly; D
  confuses file-level backup with RAID.
\item \textbf{D.} Defragmentation reorganises file placement on a disk; it is
  not a backup type. Full, incremental and differential (A, B, C) are the
  three time-based backup types.
\item \textbf{A.} The 3-2-1 rule is 3 copies, 2 different media, 1 copy
  off-site. B swaps the roles of the numbers; C invents a frequency scheme; D
  confuses backup types with the rule.
\item \textbf{B.} In 3-2-1 the 2 stands for two different types of storage
  media. A misreads it as off-site copies (that is the ``1''); C and D attach
  meanings the rule does not carry.
\item \textbf{A.} RAID expands to Redundant Array of Independent Disks, and it
  provides redundancy (uptime), not backup (recovery). B, C and D give wrong
  expansions and wrongly treat RAID as a backup mechanism.
\item \textbf{A.} Statements I and II are correct: full copies everything, and
  differential tracks changes since the last full backup. III is false --- an
  incremental tracks changes since the last backup of any type, not
  specifically the last full. B and D wrongly include III; C omits I.
\item \textbf{A.} A backup is a recoverable copy, storage is the working
  primary copy, and an archive is old data kept for long-term retention. B
  collapses the three into one; C mistakes the live working file for a
  backup; D misidentifies an archive as a RAID disk.
\item \textbf{A.} Cloud backup stores copies on remote servers, supplying an
  off-site copy. B wrongly equates it with RAID; C describes a local copy,
  which is not off-site; D is a category error, since cloud is a location,
  not a backup type.
\item \textbf{A.} An offline backup is disconnected from the network, which
  limits ransomware spread across connected drives. B describes an online
  backup; C denies a backup exists; D describes volatile memory, not backup
  media.
\item \textbf{B.} Rotation keeps several dated backup sets so an older clean
  copy survives if corruption is found late. A is unrelated to CPU clocks; C
  confuses rotation with encryption; D wrongly substitutes RAID for a backup
  strategy.
\item \textbf{C.} An incremental backup is defined as changes since the last
  backup of any type. A assigns the differential definition to the wrong
  type; B assigns the differential's ``last full'' rule to incremental; D
  describes file-level, not image.

\item \textbf{A.} Incremental backups capture changes since the previous backup. B copies everything; C and D are not backups.
\item \textbf{A.} Each incremental in the chain is needed after the full backup. B, C, and D omit required backup sets.
\item \textbf{A.} RAID gives redundancy/performance across disks. B, C, and D describe unrelated peripherals.
\item \textbf{A, B, C.} Full, incremental, and differential are backup types. Formatting (D) erases rather than backs up.
\item \textbf{A, B, C.} The three backup definitions are standard. D is false: sync mirrors current state while backup keeps restorable copies.
\item \textbf{A.} A full backup copies the entire selected set. Incremental copies changes since the latest backup; differential copies changes since the latest full backup.

\item \textbf{A.} Incremental backups record changes since the most recent backup. Differential backups use the last full backup as their reference point.

\item \textbf{A.} Each differential backup includes changes since the last full backup, so the set can grow. Incrementals reset their comparison after each backup.

\item \textbf{A.} The rule calls for three copies, two different media, and one off-site copy. The other choices do not provide this redundancy.

\item \textbf{A, B, C.} Off-site, tested, and isolated copies improve recoverability. RAID improves availability but does not protect against deletion, corruption, or every disaster.

\end{enumerate}
\unitopener{4}{I/O, Peripherals \& Multimedia}
\lessonheading{12}{Input Devices}
\begin{enumerate}[label=\textbf{Q\arabic*.}, leftmargin=*]
\item \textit{[Original]} Expand QR.
  \begin{enumerate}[label=\Alph*.]
    \item Quick Response
    \item Quick Reference
    \item Quality Read
    \item Query Result
  \end{enumerate}
\item \textit{[Original]} Expand RFID.
  \begin{enumerate}[label=\Alph*.]
    \item Rapid File Identification
    \item Random Frequency Interface Device
    \item Radio Frequency Identification
    \item Remote Field Input Device
  \end{enumerate}
\item \textit{[Original]} Which of the following is a text input device?
  \begin{enumerate}[label=\Alph*.]
    \item Keyboard
    \item Microphone
    \item Scanner
    \item Trackball
  \end{enumerate}
\item \textit{[Original]} A pointing device in which the ball is rotated in
  place while the device itself remains stationary is called a:
  \begin{enumerate}[label=\Alph*.]
    \item mouse.
    \item trackball.
    \item joystick.
    \item light pen.
  \end{enumerate}
\item \textit{[Original]} Which statement best distinguishes a touchscreen
  from a touchpad?
  \begin{enumerate}[label=\Alph*.]
    \item A touchscreen is touched directly, whereas a touchpad is a separate
      pad that moves the pointer on the screen.
    \item A touchscreen is an output device, whereas a touchpad is an input
      device.
    \item A touchscreen is used only in laptops, whereas a touchpad is used
      only in ATMs.
    \item A touchscreen and a touchpad are the same device under different
      names.
  \end{enumerate}
\item \textit{[Original]} A pen-shaped input device with a light sensor at its
  tip, pointed at a CRT screen to select or draw, is called a:
  \begin{enumerate}[label=\Alph*.]
    \item stylus.
    \item light pen.
    \item joystick.
    \item digitizer.
  \end{enumerate}
\item \textit{[PYQ JA 2026 Q68, key:final]} Consider the following devices:
  (i) Light pen (ii) Punch-card reader (iii) Plotter (iv) Scanner. Which of
  these are input devices?
  \begin{enumerate}[label=\Alph*.]
    \item (i), (ii) and (iv) only
    \item (i), (iii) and (iv) only
    \item (ii), (iii) and (iv) only
    \item (i), (ii), (iii) and (iv)
  \end{enumerate}
\item \textit{[PYQ JA 2026 Q69, key:final]} Match the device with its correct
  description:
  \begin{enumerate}[label=\Roman*.]
    \item Haptic gloves
    \item Tactile display
    \item Voice Response Unit
    \item Data glove
  \end{enumerate}
  \begin{enumerate}[label=\arabic*.]
    \item Speaks responses to a caller
    \item Worn on the hand to give touch/force feedback
    \item Worn on the hand to capture hand movement as data
    \item Presents touch sensations to the user
  \end{enumerate}
  \begin{enumerate}[label=\Alph*.]
    \item I-2, II-4, III-1, IV-3
    \item I-3, II-4, III-1, IV-2
    \item I-2, II-1, III-4, IV-3
    \item I-4, II-2, III-3, IV-1
  \end{enumerate}
\item \textit{[Original]} Which of the following is \textbf{NOT} an input
  device?
  \begin{enumerate}[label=\Alph*.]
    \item Trackball
    \item Scanner
    \item Drum printer
    \item Keyboard
  \end{enumerate}
\item \textit{[Original]} A scanner converts:
  \begin{enumerate}[label=\Alph*.]
    \item a paper document or picture into a digital image.
    \item a digital image into a printed page.
    \item sound into an electrical signal.
    \item hand movement into digital coordinates.
  \end{enumerate}
\item \textit{[Original]} Which device is a small digital camera connected to
  a computer and used mainly for live video?
  \begin{enumerate}[label=\Alph*.]
    \item Flatbed scanner
    \item Webcam
    \item Touchscreen
    \item Digital projector
  \end{enumerate}
\item \textit{[Original]} A microphone is classified as a(n):
  \begin{enumerate}[label=\Alph*.]
    \item audio input device.
    \item audio output device.
    \item image input device.
    \item pointing input device.
  \end{enumerate}
\item \textit{[Original]} Which device captures a physical trait such as a
  fingerprint or iris for identification?
  \begin{enumerate}[label=\Alph*.]
    \item Biometric sensor
    \item Barcode reader
    \item RFID tag
    \item Joystick
  \end{enumerate}
\item \textit{[Original]} Which of the following uses radio waves between a
  tag and a reader, so that no line of sight is required?
  \begin{enumerate}[label=\Alph*.]
    \item Barcode
    \item QR code
    \item RFID
    \item OCR
  \end{enumerate}
\item \textit{[Original]} A lever that can be moved on two or three axes and
  is used mainly for games is called a:
  \begin{enumerate}[label=\Alph*.]
    \item trackball.
    \item joystick.
    \item touchpad.
    \item stylus.
  \end{enumerate}
\item \textit{[Original]} Which of the following is a \textbf{direct} input
  device?
  \begin{enumerate}[label=\Alph*.]
    \item Keyboard
    \item Mouse
    \item Touchscreen
    \item Trackball
  \end{enumerate}
\item \textit{[Original]} Evaluate the statements:
  \begin{enumerate}[label=\Roman*.]
    \item A light pen is an input device.
    \item A plotter is an output device.
    \item A sound card supplies data to the computer as an input.
  \end{enumerate}
  \begin{enumerate}[label=\Alph*.]
    \item I and II only
    \item II and III only
    \item I and III only
    \item I, II and III
  \end{enumerate}
\item \textit{[Original]} A digitizer (graphics tablet) is used with a
  pen-like tool to:
  \begin{enumerate}[label=\Alph*.]
    \item print large engineering drawings.
    \item convert pen strokes into digital coordinates.
    \item display video on a screen.
    \item produce sound from the computer.
  \end{enumerate}
\item \textit{[Original]} The main function of a mouse is to:
  \begin{enumerate}[label=\Alph*.]
    \item type characters into a document.
    \item control an on-screen pointer to point, click, and drag.
    \item scan printed pages into the computer.
    \item capture live audio from the surroundings.
  \end{enumerate}
\item \textit{[Original]} Which pairing of a device with its type is correctly
  matched?
  \begin{enumerate}[label=\Alph*.]
    \item Microphone --- audio output
    \item Scanner --- image input
    \item Plotter --- input device
    \item Speaker --- input device
  \end{enumerate}

\item \textit{[JKSSB FAA-pattern]} Which device converts hand-drawn strokes directly into computer input?
  \begin{enumerate}[label=\Alph*.]
    \item Digitizer tablet
    \item Plotter
    \item Speaker
    \item Projector
  \end{enumerate}
\item \textit{[JKPSI-pattern]} Which of the following is an input device?
  \begin{enumerate}[label=\Alph*.]
    \item Scanner
    \item Printer
    \item Monitor
    \item Plotter
  \end{enumerate}
\item \textit{[SSC CGL-pattern]} A light pen is used with which device?
  \begin{enumerate}[label=\Alph*.]
    \item Monitor screen
    \item Hard disk
    \item Speaker
    \item Printer cartridge
  \end{enumerate}
\item \textit{[JKSSB-pattern]} \textbf{(MSQ: select ALL correct options.)} Which of the following are input devices?
  \begin{enumerate}[label=\Alph*.]
    \item Keyboard
    \item Mouse
    \item Microphone
    \item Printer
  \end{enumerate}
\item \textit{[SSC CGL-pattern]} \textbf{(MSQ: select ALL correct options.)} Which pairs are correctly matched?
  \begin{enumerate}[label=\Alph*.]
    \item Scanner -- input device
    \item Speaker -- output device
    \item Plotter -- input device
    \item Keyboard -- input device
  \end{enumerate}
\item \textit{[PYQ-adapted: JKSSB Junior Assistant 2026 Q68, final key; SRC-03/SRC-06]} Which device is used to enter hand-drawn strokes directly into a computer?
  \begin{enumerate}[label=\Alph*.]
\item Digitizer tablet
   \item Plotter
   \item Projector
   \item Speaker
  \end{enumerate}

\item \textit{[Original]} Which device captures still images or video as computer input?
  \begin{enumerate}[label=\Alph*.]
\item Webcam
   \item Plotter
   \item Projector
   \item Printer
  \end{enumerate}

\item \textit{[Original]} A user moves a pointer by rolling a ball mounted on a stationary pointing device. Which device is this?
  \begin{enumerate}[label=\Alph*.]
\item Trackball
   \item Joystick
   \item Scanner
   \item Light pen
  \end{enumerate}

\item \textit{[Original]} Which input device is most suitable for recording spoken instructions?
  \begin{enumerate}[label=\Alph*.]
\item Microphone
   \item Monitor
   \item Plotter
   \item Speaker
  \end{enumerate}

\item \textit{[Original, MSQ drill]} \textbf{(MSQ: select ALL correct options.)} Which are input devices?
  \begin{enumerate}[label=\Alph*.]
\item Scanner
   \item Keyboard
   \item Microphone
   \item Plotter
  \end{enumerate}

\end{enumerate}
\lessonheading{13}{OCR, OMR, MICR}
\begin{enumerate}[label=\textbf{Q\arabic*.}, leftmargin=*]
\item \textit{[Original]} Expand OCR.
  \begin{enumerate}[label=\Alph*.]
    \item Optical Character Recognition
    \item Optical Code Recognition
    \item Online Character Recognition
    \item Optical Character Reading
  \end{enumerate}
\item \textit{[Original]} Expand OMR.
  \begin{enumerate}[label=\Alph*.]
    \item Optical Mark Recognition
    \item Optical Memory Reader
    \item Online Mark Recognition
    \item Optical Mark Recording
  \end{enumerate}
\item \textit{[Original]} Expand MICR.
  \begin{enumerate}[label=\Alph*.]
    \item Magnetic Ink Character Recognition
    \item Magnetic Image Character Recognition
    \item Magnetic Ink Code Reader
    \item Machine Ink Character Recognition
  \end{enumerate}
\item \textit{[PYQ WO 2026 Q70, key:provisional]} Which reading technology
  uses magnetic ink characters?
  \begin{enumerate}[label=\Alph*.]
    \item OCR
    \item OMR
    \item MICR
    \item RFID
  \end{enumerate}
\item \textit{[Original]} Which technology is used to read objective (MCQ)
  answer sheets?
  \begin{enumerate}[label=\Alph*.]
    \item OCR
    \item OMR
    \item MICR
    \item Barcode
  \end{enumerate}
\item \textit{[Original]} Which technology reads the characters printed at
  the bottom of a bank cheque?
  \begin{enumerate}[label=\Alph*.]
    \item OMR
    \item OCR
    \item MICR
    \item QR code
  \end{enumerate}
\item \textit{[Original]} Which technology reads printed or scanned text and
  converts it into editable text?
  \begin{enumerate}[label=\Alph*.]
    \item OMR
    \item OCR
    \item MICR
    \item RFID
  \end{enumerate}
\item \textit{[Original]} Which statement correctly distinguishes OCR, OMR,
  and MICR?
  \begin{enumerate}[label=\Alph*.]
    \item OCR reads marks, OMR reads magnetic ink, and MICR reads printed
      text.
    \item OCR reads printed text, OMR reads marks, and MICR reads
      magnetic-ink characters.
    \item OCR reads magnetic ink, OMR reads printed text, and MICR reads
      marks.
    \item All three read printed text and differ only in reading speed.
  \end{enumerate}
\item \textit{[Original]} Which of the following is \textbf{NOT} correctly
  matched?
  \begin{enumerate}[label=\Alph*.]
    \item OCR --- printed or scanned text
    \item OMR --- marks on a form
    \item MICR --- magnetic-ink characters
    \item OMR --- characters printed in magnetic ink
  \end{enumerate}
\item \textit{[Original]} Which is a one-dimensional pattern of parallel
  lines used on product labels?
  \begin{enumerate}[label=\Alph*.]
    \item QR code
    \item Barcode
    \item RFID tag
    \item MICR line
  \end{enumerate}
\item \textit{[Original]} QR stands for:
  \begin{enumerate}[label=\Alph*.]
    \item Quick Recognition
    \item Quick Response
    \item Quantitative Response
    \item Quick Reading
  \end{enumerate}
\item \textit{[Original]} Compared with a barcode, a QR code:
  \begin{enumerate}[label=\Alph*.]
    \item stores less data and is one-dimensional.
    \item stores more data and is two-dimensional.
    \item stores the same data and is read only along the lines.
    \item uses magnetic ink.
  \end{enumerate}
\item \textit{[Original]} RFID stands for:
  \begin{enumerate}[label=\Alph*.]
    \item Radio Frequency Identification
    \item Rapid Frequency Identification
    \item Radio Finger Identification
    \item Remote Frequency Identifier
  \end{enumerate}
\item \textit{[Original]} Which of the following does \textbf{NOT} require
  line of sight to read the tag?
  \begin{enumerate}[label=\Alph*.]
    \item Barcode
    \item QR code
    \item RFID
    \item OCR of a printed page
  \end{enumerate}
\item \textit{[Original]} Evaluate the statements:
  \begin{enumerate}[label=\Roman*.]
    \item OCR converts a scanned image of text into editable text.
    \item OMR is used to read marks on MCQ answer sheets.
    \item MICR uses magnetic ink and is used for bank cheques.
  \end{enumerate}
  \begin{enumerate}[label=\Alph*.]
    \item I and II only
    \item II and III only
    \item I and III only
    \item I, II, and III
  \end{enumerate}
\item \textit{[Original]} Which of the following is \textbf{NOT} an example
  of biometric input?
  \begin{enumerate}[label=\Alph*.]
    \item Fingerprint scanner
    \item Iris scanner
    \item Barcode scanner
    \item Face-recognition camera
  \end{enumerate}
\item \textit{[Original]} Which pairing is correctly matched?
  \begin{enumerate}[label=\Alph*.]
    \item MICR --- MCQ answer sheets
    \item OMR --- bank cheque clearing
    \item OCR --- converting scanned documents into editable text
    \item QR code --- magnetic-ink characters
  \end{enumerate}
\item \textit{[Original]} Biometric sensors are most commonly used for which
  purpose?
  \begin{enumerate}[label=\Alph*.]
    \item Printing photographs
    \item Attendance and authentication
    \item Reading barcodes
    \item Converting scanned text
  \end{enumerate}
\item \textit{[Original]} Which of the following is \textbf{NOT} true of
  OMR?
  \begin{enumerate}[label=\Alph*.]
    \item It reads marks at fixed positions on a form.
    \item It is used for objective answer sheets.
    \item It recognises the characters printed in magnetic ink.
    \item It detects whether a position has been marked.
  \end{enumerate}
\item \textit{[Original]} Match the technology with its use:
  \begin{enumerate}[label=\Roman*.]
    \item OCR
    \item OMR
    \item MICR
  \end{enumerate}
  \begin{enumerate}[label=(\alph*)]
    \item MCQ answer sheets
    \item Bank cheques
    \item Digitising printed books
  \end{enumerate}
  \begin{enumerate}[label=\Alph*.]
    \item I-c, II-a, III-b
    \item I-a, II-b, III-c
    \item I-b, II-c, III-a
    \item I-c, II-b, III-a
  \end{enumerate}

\item \textit{[JKSSB FAA-pattern]} OMR is commonly used to check:
  \begin{enumerate}[label=\Alph*.]
    \item objective answer sheets with shaded bubbles.
    \item handwritten essays for grammar.
    \item printed photographs for colour balance.
    \item audio recordings for noise.
  \end{enumerate}
\item \textit{[JKPSI-pattern]} MICR is mainly used by which industry?
  \begin{enumerate}[label=\Alph*.]
    \item Banking (cheque processing)
    \item Film production
    \item Bakery packaging
    \item Textile weaving
  \end{enumerate}
\item \textit{[SSC CGL-pattern]} OCR stands for:
  \begin{enumerate}[label=\Alph*.]
    \item Optical Character Recognition
    \item Online Computer Repair
    \item Output Copy Register
    \item Office Cable Router
  \end{enumerate}
\item \textit{[JKSSB-pattern]} \textbf{(MSQ: select ALL correct options.)} Which technologies capture data from paper without manual typing?
  \begin{enumerate}[label=\Alph*.]
    \item OCR
    \item OMR
    \item MICR
    \item Loudspeaker playback
  \end{enumerate}
\item \textit{[SSC CGL-pattern]} \textbf{(MSQ: select ALL correct options.)} Which statements about OCR, OMR, and MICR are correct?
  \begin{enumerate}[label=\Alph*.]
    \item OCR converts printed text images into editable text.
    \item OMR detects marked positions such as shaded bubbles.
    \item MICR reads magnetic-ink characters on cheques.
    \item All three devices produce printed output on paper.
  \end{enumerate}
\item \textit{[Original]} Which technology converts printed characters into editable text?
  \begin{enumerate}[label=\Alph*.]
\item OCR
   \item OMR
   \item MICR
   \item MIDI
  \end{enumerate}

\item \textit{[PYQ-adapted: JKSSB Website Operator 2026 Q70, provisional key; SRC-04/SRC-07]} Which technology is used to read characters printed in magnetic ink on bank cheques?
  \begin{enumerate}[label=\Alph*.]
\item OMR
   \item OCR
   \item MICR
   \item QR recognition
  \end{enumerate}

\item \textit{[PYQ-adapted: JKSSB FAA 2022, QCRIAS item 553; SRC-21]} A scanned page must be converted into editable text. Which technology is the best fit?
  \begin{enumerate}[label=\Alph*.]
\item OCR
   \item OMR
   \item MICR
   \item Barcode printing
  \end{enumerate}

\item \textit{[Original]} OMR is particularly useful for:
  \begin{enumerate}[label=\Alph*.]
\item Detecting shaded response bubbles on forms
   \item Reading cheque account numbers in magnetic ink
   \item Converting speech to audio output
   \item Printing a document
  \end{enumerate}

\item \textit{[Original, MSQ drill]} \textbf{(MSQ: select ALL correct options.)} Match each recognition technology to its task.
  \begin{enumerate}[label=\Alph*.]
\item OCR --- printed text recognition
   \item OMR --- mark detection
   \item MICR --- magnetic-ink character reading
   \item OMR --- magnetic-ink character reading
  \end{enumerate}

\end{enumerate}
\lessonheading{14}{Output Devices}
\begin{enumerate}[label=\textbf{Q\arabic*.}, leftmargin=*]
\item \textit{[Original]} Expand MFP.
  \begin{enumerate}[label=\Alph*.]
    \item Multifunction Printer
    \item Multiple File Processor
    \item Main Frame Printer
    \item Multi-Format Plotter
  \end{enumerate}
\item \textit{[Original]} Expand CRT.
  \begin{enumerate}[label=\Alph*.]
    \item Central Read Tube
    \item Cathode Ray Tube
    \item Computer Read Terminal
    \item Cathode Receiver Tube
  \end{enumerate}
\item \textit{[Original]} Which of the following is an output device?
  \begin{enumerate}[label=\Alph*.]
    \item Scanner
    \item Keyboard
    \item Monitor
    \item Light pen
  \end{enumerate}
\item \textit{[Original]} Which of the following is \textbf{NOT} an output
  device?
  \begin{enumerate}[label=\Alph*.]
    \item Monitor
    \item Printer
    \item Scanner
    \item Speaker
  \end{enumerate}
\item \textit{[Original]} Output that is shown on a screen and is not printed
  is called:
  \begin{enumerate}[label=\Alph*.]
    \item hard copy.
    \item soft copy.
    \item backup copy.
    \item master copy.
  \end{enumerate}
\item \textit{[Original]} Output that is printed on paper is called:
  \begin{enumerate}[label=\Alph*.]
    \item soft copy.
    \item screen copy.
    \item hard copy.
    \item video copy.
  \end{enumerate}
\item \textit{[Original]} The older, bulky display technology that uses a
  vacuum tube and an electron beam is:
  \begin{enumerate}[label=\Alph*.]
    \item LCD.
    \item LED.
    \item OLED.
    \item CRT.
  \end{enumerate}
\item \textit{[Original]} Which display technology uses an organic compound so
  that each pixel emits its own light and no backlight is needed?
  \begin{enumerate}[label=\Alph*.]
    \item LCD
    \item LED
    \item OLED
    \item CRT
  \end{enumerate}
\item \textit{[Original]} Which display uses liquid crystal cells and requires
  a backlight?
  \begin{enumerate}[label=\Alph*.]
    \item LCD
    \item CRT
    \item Plotter
    \item Speaker
  \end{enumerate}
\item \textit{[Original]} Evaluate the statements:
  \begin{enumerate}[label=\Roman*.]
    \item A monitor produces soft copy.
    \item A printer produces hard copy.
    \item A projector is an input device.
  \end{enumerate}
  \begin{enumerate}[label=\Alph*.]
    \item I and II only
    \item II and III only
    \item I and III only
    \item I, II and III
  \end{enumerate}
\item \textit{[Original]} A device that throws a computer image onto a large
  screen or wall is a:
  \begin{enumerate}[label=\Alph*.]
    \item monitor.
    \item projector.
    \item plotter.
    \item scanner.
  \end{enumerate}
\item \textit{[PYQ WO 2026 Q69, key:provisional]} Which output device is used
  to produce large engineering drawings?
  \begin{enumerate}[label=\Alph*.]
    \item Laser printer
    \item Plotter
    \item Speaker
    \item Flatbed scanner
  \end{enumerate}
\item \textit{[Original]} Which statement best distinguishes a printer from a
  plotter?
  \begin{enumerate}[label=\Alph*.]
    \item A printer produces standard pages of text and graphics, whereas a
      plotter draws large, precise line drawings.
    \item A printer is an input device, whereas a plotter is an output device.
    \item A printer produces soft copy, whereas a plotter produces hard copy.
    \item A printer and a plotter are the same device under two names.
  \end{enumerate}
\item \textit{[Original]} Which of the following is an audio output device?
  \begin{enumerate}[label=\Alph*.]
    \item Microphone
    \item Speaker
    \item Scanner
    \item Webcam
  \end{enumerate}
\item \textit{[Original]} Headphones are classified as a(n):
  \begin{enumerate}[label=\Alph*.]
    \item audio output device.
    \item audio input device.
    \item image output device.
    \item pointing input device.
  \end{enumerate}
\item \textit{[PYQ JA 2026 Q68, key:final]} Consider the following devices:
  (i) Scanner (ii) Plotter (iii) Drum printer (iv) Light pen. Which of these
  are output devices?
  \begin{enumerate}[label=\Alph*.]
    \item (i) and (iv) only
    \item (ii) and (iii) only
    \item (i), (ii) and (iii) only
    \item (i), (iii) and (iv) only
  \end{enumerate}
\item \textit{[Original]} Which device combines printing, scanning, copying,
  and often faxing in a single unit?
  \begin{enumerate}[label=\Alph*.]
    \item Plotter
    \item Multifunction Printer (MFP)
    \item Dot-matrix printer
    \item Projector
  \end{enumerate}
\item \textit{[Original]} Which of the following is an impact printer?
  \begin{enumerate}[label=\Alph*.]
    \item Laser printer
    \item Inkjet printer
    \item Dot-matrix printer
    \item Thermal printer
  \end{enumerate}
\item \textit{[Original]} A printer that forms the image with toner and fuses
  it onto the paper by heat is a:
  \begin{enumerate}[label=\Alph*.]
    \item dot-matrix printer.
    \item laser printer.
    \item plotter.
    \item projector.
  \end{enumerate}
\item \textit{[Original]} Which pairing of an output device with its function
  is correctly matched?
  \begin{enumerate}[label=\Alph*.]
    \item Plotter --- produces large engineering drawings
    \item Monitor --- produces hard copy
    \item Speaker --- audio input device
    \item Scanner --- output device
  \end{enumerate}

\item \textit{[JKSSB FAA-pattern]} Which of the following is an output device?
  \begin{enumerate}[label=\Alph*.]
    \item Monitor
    \item Keyboard
    \item Mouse
    \item Scanner
  \end{enumerate}
\item \textit{[JKPSI-pattern]} Which device produces hard-copy output?
  \begin{enumerate}[label=\Alph*.]
    \item Printer
    \item Keyboard
    \item Microphone
    \item Joystick
  \end{enumerate}
\item \textit{[SSC CGL-pattern]} Soft-copy output is obtained from a:
  \begin{enumerate}[label=\Alph*.]
    \item monitor screen.
    \item printed register.
    \item plotted chart on paper.
    \item stamped envelope.
  \end{enumerate}
\item \textit{[JKSSB-pattern]} \textbf{(MSQ: select ALL correct options.)} Which of the following are output devices?
  \begin{enumerate}[label=\Alph*.]
    \item Monitor
    \item Printer
    \item Speaker
    \item Keyboard
  \end{enumerate}
\item \textit{[SSC CGL-pattern]} \textbf{(MSQ: select ALL correct options.)} Which statements about output devices are correct?
  \begin{enumerate}[label=\Alph*.]
    \item A monitor shows soft-copy output.
    \item A printer produces hard-copy output.
    \item A speaker gives audio output.
    \item A scanner displays result output.
  \end{enumerate}
\item \textit{[PYQ-adapted: JKSSB Website Operator 2026 Q69, provisional key; SRC-04/SRC-07]} Which output device is designed for producing large engineering drawings?
  \begin{enumerate}[label=\Alph*.]
\item Plotter
   \item Joystick
   \item Scanner
   \item Microphone
  \end{enumerate}

\item \textit{[Original]} A monitor displays a report on screen without making a paper copy. This is:
  \begin{enumerate}[label=\Alph*.]
\item Soft copy output
   \item Hard copy output
   \item Input
   \item Secondary storage
  \end{enumerate}

\item \textit{[Original]} Which device provides visual output on a large screen for an audience?
  \begin{enumerate}[label=\Alph*.]
\item Projector
   \item Scanner
   \item Barcode reader
   \item Trackball
  \end{enumerate}

\item \textit{[Original]} A tactile display primarily presents information through:
  \begin{enumerate}[label=\Alph*.]
\item Raised or refreshable tactile elements
   \item Magnetic ink
   \item Printed barcodes
   \item Audio sampling only
  \end{enumerate}

\item \textit{[Original, MSQ drill]} \textbf{(MSQ: select ALL correct options.)} Which are output devices?
  \begin{enumerate}[label=\Alph*.]
\item Monitor
   \item Speaker
   \item Plotter
   \item Scanner
  \end{enumerate}

\end{enumerate}
\lessonheading{15}{Printers and Displays}
\begin{enumerate}[label=\textbf{Q\arabic*.}, leftmargin=*]
\item \textit{[Original]} Which of the following is an \textbf{impact
  printer}?
  \begin{enumerate}[label=\Alph*.]
    \item Inkjet printer
    \item Laser printer
    \item Dot matrix printer
    \item Thermal printer
  \end{enumerate}
\item \textit{[Original]} A printer that forms characters by striking an
  inked ribbon against the paper is called a:
  \begin{enumerate}[label=\Alph*.]
    \item non-impact printer
    \item impact printer
    \item page printer
    \item plotter
  \end{enumerate}
\item \textit{[Original]} Which of the following is \textbf{NOT} an impact
  printer?
  \begin{enumerate}[label=\Alph*.]
    \item Dot matrix printer
    \item Daisy wheel printer
    \item Drum printer
    \item Inkjet printer
  \end{enumerate}
\item \textit{[Original]} Which printer forms characters and images by
  spraying tiny droplets of liquid ink onto the paper?
  \begin{enumerate}[label=\Alph*.]
    \item Laser printer
    \item Inkjet printer
    \item Thermal printer
    \item Dot matrix printer
  \end{enumerate}
\item \textit{[Original]} Which printer uses a laser beam and toner powder to
  form an image through a xerographic process?
  \begin{enumerate}[label=\Alph*.]
    \item Inkjet printer
    \item Laser printer
    \item Dot matrix printer
    \item Thermal printer
  \end{enumerate}
\item \textit{[Original]} Which of the following is a \textbf{page printer}
  (it prints a whole page at a time)?
  \begin{enumerate}[label=\Alph*.]
    \item Dot matrix printer
    \item Daisy wheel printer
    \item Laser printer
    \item Plotter
  \end{enumerate}
\item \textit{[Original]} Which printer produces an image by applying heat to
  specially treated paper, as used in POS receipts and barcode labels?
  \begin{enumerate}[label=\Alph*.]
    \item Laser printer
    \item Inkjet printer
    \item Thermal printer
    \item Dot matrix printer
  \end{enumerate}
\item \textit{[Original]} A 3D printer produces its output by:
  \begin{enumerate}[label=\Alph*.]
    \item striking an inked ribbon against paper
    \item adding material layer by layer to build an object
    \item spraying liquid ink onto paper
    \item melting toner onto a page
  \end{enumerate}
\item \textit{[Original]} Which output device is used to produce large
  engineering drawings, maps, and CAD line drawings?
  \begin{enumerate}[label=\Alph*.]
    \item Inkjet printer
    \item Plotter
    \item Laser printer
    \item Dot matrix printer
  \end{enumerate}
\item \textit{[Original]} Which statement correctly distinguishes an inkjet
  printer from a laser printer?
  \begin{enumerate}[label=\Alph*.]
    \item An inkjet printer uses toner, while a laser printer uses liquid
      ink.
    \item An inkjet printer sprays liquid ink, while a laser printer uses a
      laser beam and toner.
    \item Both are impact printers.
    \item A laser printer is a character printer, while an inkjet printer is
      a page printer.
  \end{enumerate}
\item \textit{[Original]} A printer that prints an entire line of text at a
  time, such as a drum printer, is called a:
  \begin{enumerate}[label=\Alph*.]
    \item character printer
    \item line printer
    \item page printer
    \item plotter
  \end{enumerate}
\item \textit{[Original]} Expand \textbf{DPI} as used in printing and display.
  \begin{enumerate}[label=\Alph*.]
    \item Dots Per Inch
    \item Data Per Inch
    \item Dots Per Image
    \item Display Pixels Index
  \end{enumerate}
\item \textit{[Original]} The number of pixels a display can show
  horizontally and vertically, written as columns $\times$ rows, is its:
  \begin{enumerate}[label=\Alph*.]
    \item refresh rate
    \item resolution
    \item response time
    \item contrast ratio
  \end{enumerate}
\item \textit{[Original]} The smallest addressable element of a display is
  called a:
  \begin{enumerate}[label=\Alph*.]
    \item pixel
    \item byte
    \item bitmap
    \item raster
  \end{enumerate}
\item \textit{[Original]} The refresh rate of a display is measured in:
  \begin{enumerate}[label=\Alph*.]
    \item hertz (Hz)
    \item milliseconds (ms)
    \item pixels
    \item nits
  \end{enumerate}
\item \textit{[Original]} The response time of a monitor is measured in:
  \begin{enumerate}[label=\Alph*.]
    \item hertz (Hz)
    \item milliseconds (ms)
    \item dots per inch (DPI)
    \item nits
  \end{enumerate}
\item \textit{[Original]} The difference between the brightest white and the
  darkest black a display can produce is described by its:
  \begin{enumerate}[label=\Alph*.]
    \item refresh rate
    \item contrast ratio
    \item resolution
    \item response time
  \end{enumerate}
\item \textit{[Original]} The brightness of a display, that is the amount of
  light it emits, is commonly measured in:
  \begin{enumerate}[label=\Alph*.]
    \item hertz (Hz)
    \item nits
    \item milliseconds (ms)
    \item dots per inch (DPI)
  \end{enumerate}
\item \textit{[Original]} Which statement is \textbf{NOT} correct about
  non-impact printers?
  \begin{enumerate}[label=\Alph*.]
    \item Inkjet printers are non-impact.
    \item Laser printers are non-impact.
    \item Thermal printers are non-impact.
    \item Dot matrix printers are non-impact.
  \end{enumerate}
\item \textit{[Original]} Evaluate the statements:
  \begin{enumerate}[label=\Roman*.]
    \item A dot matrix printer is an impact printer.
    \item A laser printer is a non-impact printer.
    \item A thermal printer is an impact printer.
  \end{enumerate}
  \begin{enumerate}[label=\Alph*.]
    \item I only
    \item I and II only
    \item II and III only
    \item I, II, and III
  \end{enumerate}

\item \textit{[JKSSB FAA-pattern]} Which printer strikes the paper through an inked ribbon?
  \begin{enumerate}[label=\Alph*.]
    \item Dot-matrix printer
    \item Laser printer
    \item Inkjet printer
    \item Thermal printer
  \end{enumerate}
\item \textit{[JKPSI-pattern]} Which printer type generally gives the fastest high-volume office output?
  \begin{enumerate}[label=\Alph*.]
    \item Laser printer
    \item Dot-matrix printer
    \item Daisy-wheel printer
    \item Manual typewriter
  \end{enumerate}
\item \textit{[SSC CGL-pattern]} Screen resolution is measured in:
  \begin{enumerate}[label=\Alph*.]
    \item pixels.
    \item litres.
    \item kilograms.
    \item volts.
  \end{enumerate}
\item \textit{[JKSSB-pattern]} \textbf{(MSQ: select ALL correct options.)} Which of the following are impact printers?
  \begin{enumerate}[label=\Alph*.]
    \item Dot-matrix printer
    \item Daisy-wheel printer
    \item Laser printer
    \item Inkjet printer
  \end{enumerate}
\item \textit{[SSC CGL-pattern]} \textbf{(MSQ: select ALL correct options.)} Which statements about printers and displays are correct?
  \begin{enumerate}[label=\Alph*.]
    \item Laser printers are non-impact printers.
    \item Inkjet printers spray droplets of ink.
    \item Higher pixel counts mean sharper displayed images.
    \item Dot-matrix printers use a laser beam to form characters.
  \end{enumerate}
\item \textit{[Original]} Which printer forms characters by striking an inked ribbon against paper?
  \begin{enumerate}[label=\Alph*.]
\item Dot-matrix printer
   \item Inkjet printer
   \item Laser printer
   \item Thermal printer
  \end{enumerate}

\item \textit{[Original]} Which printer is generally classified as non-impact?
  \begin{enumerate}[label=\Alph*.]
\item Laser printer
   \item Daisy-wheel printer
   \item Dot-matrix printer
   \item Line printer
  \end{enumerate}

\item \textit{[Original]} Which device produces a hard copy?
  \begin{enumerate}[label=\Alph*.]
\item Printer
   \item Monitor
   \item Speaker
   \item Projector
  \end{enumerate}

\item \textit{[Original]} A laser printer typically uses which consumable?
  \begin{enumerate}[label=\Alph*.]
\item Toner
   \item Magnetic ink ribbon
   \item Liquid ink cartridge only
   \item Thermal paper roll only
  \end{enumerate}

\item \textit{[Original, MSQ drill]} \textbf{(MSQ: select ALL correct options.)} Which statements are correct?
  \begin{enumerate}[label=\Alph*.]
\item Dot-matrix printing is impact printing.
   \item Laser printing is non-impact printing.
   \item Inkjet printers place ink droplets on media.
   \item A monitor is a hard-copy device.
  \end{enumerate}

\end{enumerate}
\lessonheading{16}{Computer Ports}
\begin{enumerate}[label=\textbf{Q\arabic*.}, leftmargin=*]
\item \textit{[Original]} Expand USB.
  \begin{enumerate}[label=\Alph*.]
    \item Universal Serial Bus
    \item Universal System Bus
    \item Unified Serial Board
    \item Universal Storage Bus
  \end{enumerate}
\item \textit{[Original]} Expand HDMI.
  \begin{enumerate}[label=\Alph*.]
    \item High-Definition Multimedia Interface
    \item High-Definition Media Input
    \item High-Digital Multimedia Interface
    \item High-Definition Memory Interface
  \end{enumerate}
\item \textit{[Original]} Expand VGA.
  \begin{enumerate}[label=\Alph*.]
    \item Video Graphics Array
    \item Visual Graphics Adapter
    \item Video Gateway Array
    \item Virtual Graphics Array
  \end{enumerate}
\item \textit{[Original]} Expand RJ-45, the connector used with a network
  cable.
  \begin{enumerate}[label=\Alph*.]
    \item Registered Jack 45
    \item Routed Jack 45
    \item Registered Junction 45
    \item Random Jack 45
  \end{enumerate}
\item \textit{[Original]} Expand SATA.
  \begin{enumerate}[label=\Alph*.]
    \item Serial Advanced Technology Attachment
    \item Serial Access Terminal Adapter
    \item System Access Technology Attachment
    \item Serial Advanced Transfer Array
  \end{enumerate}
\item \textit{[Original]} Which of the following is \textbf{NOT} a video
  output port?
  \begin{enumerate}[label=\Alph*.]
    \item HDMI
    \item DisplayPort
    \item VGA
    \item RJ-45
  \end{enumerate}
\item \textit{[Original]} Which USB connector type is \textbf{reversible},
  so it can be inserted either way up?
  \begin{enumerate}[label=\Alph*.]
    \item USB Type-A
    \item USB Type-B
    \item USB Type-C
    \item Micro-USB
  \end{enumerate}
\item \textit{[Original]} Which statement about HDMI and VGA is correct?
  \begin{enumerate}[label=\Alph*.]
    \item Both are analog video ports.
    \item HDMI is digital and carries audio; VGA is analog and carries no
      audio.
    \item Both carry digital audio and video.
    \item HDMI is analog; VGA is digital.
  \end{enumerate}
\item \textit{[Original]} Which connector is the standard for a \textbf{wired
  Ethernet LAN} connection?
  \begin{enumerate}[label=\Alph*.]
    \item RJ-11
    \item RJ-45
    \item USB-C
    \item HDMI
  \end{enumerate}
\item \textit{[Original]} Evaluate the statements:
  \begin{enumerate}[label=\Roman*.]
    \item USB can supply power to a connected device.
    \item USB supports hot swapping.
    \item VGA carries both video and audio.
  \end{enumerate}
  \begin{enumerate}[label=\Alph*.]
    \item I and II only
    \item II and III only
    \item I and III only
    \item I, II, and III
  \end{enumerate}
\item \textit{[Original]} In computer networking, a ``port number'' such as
  80 or 443 refers to:
  \begin{enumerate}[label=\Alph*.]
    \item a physical socket on the network card.
    \item a logical (software) port that identifies a service.
    \item the number of pins in an RJ-45 connector.
    \item the speed of the network cable.
  \end{enumerate}
\item \textit{[Original]} Which port is designed mainly for connecting a
  computer to a monitor and supports high refresh rates and multiple
  monitors?
  \begin{enumerate}[label=\Alph*.]
    \item VGA
    \item DisplayPort
    \item RJ-45
    \item SATA
  \end{enumerate}
\item \textit{[Original]} Which video port is \textbf{analog} and is now
  largely replaced by HDMI and DisplayPort?
  \begin{enumerate}[label=\Alph*.]
    \item VGA
    \item HDMI
    \item DisplayPort
    \item USB-C
  \end{enumerate}
\item \textit{[Original]} Which feature of USB means a device can be
  connected or removed without switching off the computer?
  \begin{enumerate}[label=\Alph*.]
    \item cold boot
    \item hot swapping
    \item fragmentation
    \item multiplexing
  \end{enumerate}
\item \textit{[Original]} Which interface connects an \textbf{internal} hard
  disk or SSD to the motherboard?
  \begin{enumerate}[label=\Alph*.]
    \item SATA
    \item HDMI
    \item RJ-45
    \item 3.5 mm audio jack
  \end{enumerate}
\item \textit{[Original]} On a typical PC, which audio connector adds a
  microphone channel to stereo audio?
  \begin{enumerate}[label=\Alph*.]
    \item TRS
    \item TRRS
    \item RJ-45
    \item VGA
  \end{enumerate}
\item \textit{[Original]} Which pairing of port and typical device is
  correctly matched?
  \begin{enumerate}[label=\Alph*.]
    \item HDMI --- headphones
    \item Ethernet (RJ-45) --- wired network router
    \item VGA --- printer
    \item SATA --- keyboard
  \end{enumerate}
\item \textit{[Original]} Which of the following is a \textbf{storage
  interface} rather than a general-purpose peripheral port?
  \begin{enumerate}[label=\Alph*.]
    \item USB
    \item SATA
    \item HDMI
    \item audio jack
  \end{enumerate}
\item \textit{[Original]} What is \textbf{eSATA} used for?
  \begin{enumerate}[label=\Alph*.]
    \item connecting an external storage drive
    \item connecting a monitor
    \item connecting a microphone
    \item connecting to a wired network
  \end{enumerate}
\item \textit{[Original]} Which statement uses the port terminology
  correctly?
  \begin{enumerate}[label=\Alph*.]
    \item USB is the standard, USB-C is a connector type, and the USB port is
      the socket on the device.
    \item A connector is the socket on the device.
    \item A port and a connector always mean exactly the same thing.
    \item HDMI is an analog video connector.
  \end{enumerate}

\item \textit{[JKSSB FAA-pattern]} Which port is commonly used to connect a modern printer?
  \begin{enumerate}[label=\Alph*.]
    \item USB
    \item VGA
    \item RJ-11 only
    \item Audio jack
  \end{enumerate}
\item \textit{[JKPSI-pattern]} HDMI ports carry which signals?
  \begin{enumerate}[label=\Alph*.]
    \item Audio and video
    \item Only printer ink
    \item Only keyboard strokes
    \item Only power supply current
  \end{enumerate}
\item \textit{[SSC CGL-pattern]} LAN cables are usually connected through which connector?
  \begin{enumerate}[label=\Alph*.]
    \item RJ-45
    \item USB-C charger pin
    \item Audio jack
    \item VGA screw pin
  \end{enumerate}
\item \textit{[JKSSB-pattern]} \textbf{(MSQ: select ALL correct options.)} Which of the following are computer ports/interfaces?
  \begin{enumerate}[label=\Alph*.]
    \item USB
    \item HDMI
    \item Ethernet (RJ-45)
    \item Stapler slot
  \end{enumerate}
\item \textit{[SSC CGL-pattern]} \textbf{(MSQ: select ALL correct options.)} Which statements about ports are correct?
  \begin{enumerate}[label=\Alph*.]
    \item USB connects many peripherals such as printers and drives.
    \item HDMI carries audio and video to displays.
    \item RJ-45 connects Ethernet network cables.
    \item VGA carries digital audio to speakers.
  \end{enumerate}
\item \textit{[Original]} Which connector is commonly used to carry digital video and audio from a computer to a display?
  \begin{enumerate}[label=\Alph*.]
\item HDMI
   \item RJ-11
   \item PS/2
   \item VGA only
  \end{enumerate}

\item \textit{[Original]} Which interface is commonly used for a wired Ethernet connection?
  \begin{enumerate}[label=\Alph*.]
\item RJ-45
   \item HDMI
   \item 3.5 mm audio jack
   \item VGA
  \end{enumerate}

\item \textit{[Original]} Which statement about USB Type-C is correct?
  \begin{enumerate}[label=\Alph*.]
\item Its plug is reversible.
   \item It is an analog-only video connector.
   \item It is used only for telephone lines.
   \item It is a logical network port number.
  \end{enumerate}

\item \textit{[Original]} Which port number is a logical transport-layer endpoint rather than a physical socket?
  \begin{enumerate}[label=\Alph*.]
\item TCP port 80
   \item HDMI socket
   \item USB-A receptacle
   \item VGA connector
  \end{enumerate}

\item \textit{[Original, MSQ drill]} \textbf{(MSQ: select ALL correct options.)} Match the interface to a common use.
  \begin{enumerate}[label=\Alph*.]
\item HDMI --- digital audio/video
   \item RJ-45 --- Ethernet
   \item USB --- peripheral/data connection
   \item VGA --- magnetic-ink character reading
  \end{enumerate}

\end{enumerate}
\lessonheading{65}{Current Multimedia, GUI and Media Formats}
\begin{enumerate}[label=\textbf{Q\arabic*.}, leftmargin=*]
\item \textit{[Original]} Expand GUI.
  \begin{enumerate}[label=\Alph*.]
    \item General User Interface
    \item Graphical User Interface
    \item Graphical Utility Interconnect
    \item Graphic User Interpreter
  \end{enumerate}
\item \textit{[Original]} Expand WYSIWYG.
  \begin{enumerate}[label=\Alph*.]
    \item What You See Is What You Get
    \item What You Save Is What You Get
    \item Where You See Is Where You Go
    \item What You Send Is What You Gain
  \end{enumerate}
\item \textit{[Original]} Expand MP3.
  \begin{enumerate}[label=\Alph*.]
    \item MPEG-3 Audio Layer I
    \item MPEG Audio Video Layer III
    \item Moving Picture Audio Layer III
    \item MPEG-1 Audio Layer III
  \end{enumerate}
\item \textit{[Original]} Expand PNG.
  \begin{enumerate}[label=\Alph*.]
    \item Portable Network Graphics
    \item Personal Network Graphics
    \item Portable Native Graphics
    \item Portable Network Group
  \end{enumerate}
\item \textit{[Original]} Expand JPEG.
  \begin{enumerate}[label=\Alph*.]
    \item Joint Picture Encoding Group
    \item Java Photographic Experts Group
    \item Joint Photographic Experts Group
    \item Joint Photographic Encoding Group
  \end{enumerate}
\item \textit{[Original]} Expand MP4.
  \begin{enumerate}[label=\Alph*.]
    \item MPEG-4 Portable 14
    \item MPEG-4 Picture 4
    \item MPEG-4 Player 4
    \item MPEG-4 Part 14
  \end{enumerate}
\item \textit{[Original]} In digitising an analogue signal, the process of
  measuring the signal at regular intervals is called:
  \begin{enumerate}[label=\Alph*.]
    \item quantization
    \item sampling
    \item synchronization
    \item compression
  \end{enumerate}
\item \textit{[Original]} The process of mapping each measured sample to one of
  a fixed set of discrete levels is called:
  \begin{enumerate}[label=\Alph*.]
    \item sampling
    \item synchronization
    \item quantization
    \item buffering
  \end{enumerate}
\item \textit{[Original]} Which of the following uses \textbf{lossless}
  compression?
  \begin{enumerate}[label=\Alph*.]
    \item JPEG
    \item MP3
    \item PNG
    \item MP4
  \end{enumerate}
\item \textit{[Original]} Which of the following correctly pairs an image
  format with its nature?
  \begin{enumerate}[label=\Alph*.]
    \item PNG --- lossy, does not support transparency
    \item JPEG --- lossless, best for photographs
    \item GIF --- lossless, limited to 256 colours
    \item BMP --- a small web format for photographs
  \end{enumerate}
\item \textit{[Original]} Which image format supports simple animation?
  \begin{enumerate}[label=\Alph*.]
    \item PNG
    \item GIF
    \item JPEG
    \item BMP
  \end{enumerate}
\item \textit{[Original]} Which format is most suitable for photographs?
  \begin{enumerate}[label=\Alph*.]
    \item PNG
    \item GIF
    \item JPEG
    \item TXT
  \end{enumerate}
\item \textit{[Original]} In multimedia, the alignment in time of separate
  streams such as audio and video is called:
  \begin{enumerate}[label=\Alph*.]
    \item compression
    \item sampling
    \item synchronization
    \item quantization
  \end{enumerate}
\item \textit{[Original]} The four classic elements of a GUI are remembered by
  the short form WIMP. Which set is correct?
  \begin{enumerate}[label=\Alph*.]
    \item Windows, Icons, Menus, Pointers
    \item Windows, Images, Mouse, Program
    \item Wires, Icons, Menus, Panels
    \item Windows, Icons, Modes, Printers
  \end{enumerate}
\item \textit{[Original]} A small window that asks the user for input or
  confirms an action, such as Save As or Print, is called a:
  \begin{enumerate}[label=\Alph*.]
    \item desktop
    \item dialogue box
    \item taskbar
    \item shortcut
  \end{enumerate}
\item \textit{[Original]} WYSIWYG means that:
  \begin{enumerate}[label=\Alph*.]
    \item the document is stored on the cloud.
    \item the screen shows the document exactly as it will be printed.
    \item the software corrects spelling automatically.
    \item the file is compressed before saving.
  \end{enumerate}
\item \textit{[Original]} A plotter is best suited for:
  \begin{enumerate}[label=\Alph*.]
    \item printing ordinary text letters.
    \item large engineering drawings and maps.
    \item scanning photographs into a computer.
    \item recording sound during a presentation.
  \end{enumerate}
\item \textit{[Original]} Which of the following is a \textbf{non-impact}
  printer?
  \begin{enumerate}[label=\Alph*.]
    \item Dot-matrix printer
    \item Daisy-wheel printer
    \item Line printer
    \item Laser printer
  \end{enumerate}
\item \textit{[Original]} Which of the following is \textbf{NOT} a lossless
  format?
  \begin{enumerate}[label=\Alph*.]
    \item PNG
    \item GIF
    \item JPEG
    \item ZIP
  \end{enumerate}
\item \textit{[Original]} Evaluate the statements:
  \begin{enumerate}[label=\Roman*.]
    \item MP3 is a lossy audio format.
    \item MP4 is a container format that can hold video and audio together.
    \item PNG is a lossless image format that supports transparency.
  \end{enumerate}
  \begin{enumerate}[label=\Alph*.]
    \item I and II only
    \item I and III only
    \item II and III only
    \item I, II, and III
  \end{enumerate}

\item \textit{[JKSSB FAA-pattern]} GUI stands for:
  \begin{enumerate}[label=\Alph*.]
    \item Graphical User Interface
    \item General Utility Ink
    \item Giant Screen Umbrella
    \item Green Upload Icon
  \end{enumerate}
\item \textit{[JKPSI-pattern]} Which format is commonly used for web images with animation?
  \begin{enumerate}[label=\Alph*.]
    \item GIF
    \item DOCX
    \item XLSX
    \item TXT
  \end{enumerate}
\item \textit{[SSC CGL-pattern]} MP3 is a format for which media?
  \begin{enumerate}[label=\Alph*.]
    \item Audio
    \item Printed text
    \item Spreadsheet data
    \item Database keys
  \end{enumerate}
\item \textit{[JKSSB-pattern]} \textbf{(MSQ: select ALL correct options.)} Which are GUI elements?
  \begin{enumerate}[label=\Alph*.]
    \item Windows
    \item Icons
    \item Menus
    \item Printer gears
  \end{enumerate}
\item \textit{[SSC CGL-pattern]} \textbf{(MSQ: select ALL correct options.)} Which media-format pairs are correctly matched?
  \begin{enumerate}[label=\Alph*.]
    \item .jpg -- image
    \item .mp3 -- audio
    \item .mp4 -- video
    \item .jpg -- spreadsheet
  \end{enumerate}
\item \textit{[PYQ-adapted: JKSSB Website Operator 2026 Q69, provisional key; SRC-04/SRC-07]} Which device is used to produce large engineering drawings?
  \begin{enumerate}[label=\Alph*.]
\item Plotter
   \item Microphone
   \item Scanner
   \item Trackball
  \end{enumerate}

\item \textit{[PYQ-adapted: JKSSB Junior Assistant 2026 Q68, final key; SRC-03/SRC-06]} Which device is an output device?
  \begin{enumerate}[label=\Alph*.]
\item Plotter
   \item Scanner
   \item Light pen
   \item Digitizer tablet
  \end{enumerate}

\item \textit{[Original]} Sampling in digital audio converts a continuous signal into:
  \begin{enumerate}[label=\Alph*.]
\item Measurements taken at discrete time intervals
   \item A printed paper page
   \item A file-system directory
   \item A network address
  \end{enumerate}

\item \textit{[Original]} Which file format commonly stores compressed still images?
  \begin{enumerate}[label=\Alph*.]
\item JPEG
   \item DOCX
   \item CSV
   \item EXE
  \end{enumerate}

\item \textit{[Original, MSQ drill]} \textbf{(MSQ: select ALL correct options.)} Which are multimedia output devices?
  \begin{enumerate}[label=\Alph*.]
\item Monitor
   \item Speaker
   \item Projector
   \item Scanner
  \end{enumerate}

\end{enumerate}
\answersection{4}
\answersubheading{12}{Input Devices}
\begin{enumerate}[label=\textbf{Q\arabic*.}, leftmargin=*]
\item \textbf{A.} QR stands for Quick Response. B, C, and D are not the
  correct full form.
\item \textbf{C.} RFID stands for Radio Frequency Identification. A, B, and D
  are not the correct expansions.
\item \textbf{A.} A keyboard is a text input device. A microphone (B) is
  audio input, a scanner (C) is image input, and a trackball (D) is a
  pointing input device.
\item \textbf{B.} A trackball rotates its ball in place while the device
  stays still. A mouse (A) slides over a surface; a joystick (C) is a lever;
  a light pen (D) uses a light sensor.
\item \textbf{A.} On a touchscreen the display itself is touched directly;
  on a touchpad a separate pad moves the pointer. B is wrong because both are
  input devices; C reverses the typical uses; D denies a real difference.
\item \textbf{B.} The light pen has a light sensor and is pointed at a CRT
  screen. A stylus (A) is used on a tablet, a joystick (C) is a lever, and a
  digitizer (D) is the tablet surface that reads the stylus.
\item \textbf{A.} Light pen, punch-card reader, and scanner are input devices;
  the plotter is an output device. B and C wrongly include the plotter, and D
  includes every device, also wrongly counting the plotter.
\item \textbf{A.} Haptic gloves give touch/force feedback (2), a tactile
  display presents touch sensations (4), a Voice Response Unit speaks to a
  caller (1), and a data glove captures hand movement as data (3). B, C, and
  D scramble these matches.
\item \textbf{C.} A drum printer produces output. A trackball (A), a scanner
  (B), and a keyboard (D) are input devices.
\item \textbf{A.} A scanner turns a paper document or picture into a digital
  image. Turning an image into a printed page (B) is printing, sound into a
  signal (C) is a microphone, and hand movement into coordinates (D) is a
  digitizer.
\item \textbf{B.} A webcam is a small digital camera used for live video. A
  flatbed scanner (A) scans pages, a touchscreen (C) is a pointing input
  device, and a digital projector (D) is an output device.
\item \textbf{A.} A microphone converts sound into a signal, so it is an
  audio input device. A speaker (B) is audio output; a scanner or camera (C)
  is image input; a mouse (D) is a pointing device.
\item \textbf{A.} A biometric sensor reads a physical trait such as a
  fingerprint or iris. A barcode reader (B) reads optical codes, an RFID tag
  (C) works by radio waves, and a joystick (D) is for games and control.
\item \textbf{C.} RFID exchanges radio waves between a tag and a reader with
  no line of sight. A barcode (A) and a QR code (B) are read optically, and
  OCR (D) reads text from a scanned image.
\item \textbf{B.} A joystick is a lever moved on two or three axes. A
  trackball (A) has a rotating ball, a touchpad (C) is a flat finger pad, and
  a stylus (D) is a pen-like tool.
\item \textbf{C.} A touchscreen captures the selection directly on the
  display. A keyboard (A), a mouse (B), and a trackball (D) are indirect
  input devices.
\item \textbf{A.} Statements I and II are true: the light pen is input and
  the plotter is output. III is false because a sound card is an output
  device, not an input. B and C wrongly include III; D includes all three.
\item \textbf{B.} A digitizer converts pen strokes into digital coordinates.
  Printing large drawings (A) is a plotter's job, displaying video (C) is a
  monitor, and producing sound (D) is a speaker.
\item \textbf{B.} A mouse controls the on-screen pointer to point, click, and
  drag. Typing characters (A) is a keyboard, scanning pages (C) is a scanner,
  and capturing audio (D) is a microphone.
\item \textbf{B.} A scanner is an image input device, so this pairing is
  correct. A microphone (A) is input, not output; a plotter (C) is output,
  not input; and a speaker (D) is output, not input.

\item \textbf{A.} A digitizer tablet captures strokes as input. The plotter (B), speaker (C), and projector (D) are output devices.
\item \textbf{A.} A scanner feeds images into the computer. B, C, and D all present output.
\item \textbf{A.} A light pen draws on the monitor screen. B, C, and D are not pen-compatible input surfaces.
\item \textbf{A, B, C.} Keyboard, mouse, and microphone all supply input. The printer (D) produces output.
\item \textbf{A, B, D.} Scanner-input, speaker-output, and keyboard-input are correct. C is false: a plotter is an output device.
\item \textbf{A.} A digitizer tablet is an input device. A plotter, projector, and speaker present output.

\item \textbf{A.} A webcam captures visual input. Plotters, projectors, and printers are output devices.

\item \textbf{A.} A trackball moves the pointer by rotating a ball while the device remains stationary. The other devices use a stick, capture images, or detect a screen position.

\item \textbf{A.} A microphone converts sound to input. A monitor, plotter, and speaker deliver output.

\item \textbf{A, B, C.} Scanner, keyboard, and microphone provide input. A plotter produces output.

\end{enumerate}
\answersubheading{13}{OCR, OMR, MICR}
\begin{enumerate}[label=\textbf{Q\arabic*.}, leftmargin=*]
\item \textbf{A.} OCR stands for Optical Character Recognition. B, C, and D
  are not expansions of the abbreviation.
\item \textbf{A.} OMR stands for Optical Mark Recognition. B, C, and D are
  not the correct full forms.
\item \textbf{A.} MICR stands for Magnetic Ink Character Recognition. B, C,
  and D are not the correct full forms.
\item \textbf{C.} Magnetic ink characters are read by MICR. OCR (A) reads
  printed text, OMR (B) reads marks, and RFID (D) reads radio tags.
\item \textbf{B.} OMR reads marks at fixed positions, which is exactly how
  answer-sheet bubbles are filled. OCR (A) reads text, MICR (C) reads
  magnetic ink, and a barcode (D) reads product labels.
\item \textbf{C.} Cheque characters are printed in magnetic ink, so MICR
  reads them. OMR (A) reads marks, OCR (B) reads ordinary printed text, and
  a QR code (D) is a two-dimensional optical code.
\item \textbf{B.} OCR converts an image of printed or scanned text into
  editable text. OMR (A) reads marks, MICR (C) reads magnetic ink, and RFID
  (D) reads radio tags.
\item \textbf{B.} OCR reads printed text, OMR reads marks, and MICR reads
  magnetic-ink characters. A, C, and D scramble these read-targets
  (TRAP-25).
\item \textbf{D.} OMR reads marks, not magnetic-ink characters; magnetic ink
  belongs to MICR. A, B, and C are all correct matches.
\item \textbf{B.} A barcode is a one-dimensional pattern of parallel lines.
  A QR code (A) is two-dimensional, an RFID tag (C) uses radio waves, and a
  MICR line (D) is magnetic-ink text on a cheque.
\item \textbf{B.} QR stands for Quick Response. A, C, and D are not the
  correct full forms.
\item \textbf{B.} A QR code is two-dimensional and stores more data than a
  barcode. A reverses both facts; C wrongly makes them equal and
  line-limited; D wrongly assigns magnetic ink to a QR code.
\item \textbf{A.} RFID stands for Radio Frequency Identification. B, C, and
  D are not the correct full forms.
\item \textbf{C.} RFID uses radio waves and needs no line of sight. A
  barcode (A), a QR code (B), and OCR of a printed page (D) all require the
  reader to see the pattern.
\item \textbf{D.} All three statements are correct: OCR converts scanned
  text into editable text, OMR reads marks on MCQ sheets, and MICR reads
  magnetic ink on bank cheques. A, B, and C each omit at least one true
  statement.
\item \textbf{C.} A barcode scanner reads a printed code, not a human
  trait, so it is not biometric. A fingerprint scanner (A), an iris scanner
  (B), and a face-recognition camera (D) all identify a person from a
  physical trait.
\item \textbf{C.} OCR is correctly paired with converting scanned documents
  into editable text. A wrongly pairs MICR with answer sheets; B wrongly
  pairs OMR with cheque clearing; D wrongly pairs a QR code with magnetic
  ink.
\item \textbf{B.} Biometric sensors are used for attendance and
  authentication. Printing (A) is output, reading barcodes (C) is an optical
  code task, and converting scanned text (D) is OCR.
\item \textbf{C.} OMR reads marks; it does not recognise magnetic-ink
  characters, which is MICR's role. Statements A, B, and D are all true of
  OMR.
\item \textbf{A.} OCR digitises printed books (I-c), OMR reads MCQ answer
  sheets (II-a), and MICR reads bank cheques (III-b). B, C, and D assign at
  least one technology to the wrong use.

\item \textbf{A.} OMR senses shaded marks on objective sheets. B, C, and D are not mark-reading tasks.
\item \textbf{A.} Banks use MICR for cheques. B, C, and D do not process magnetic-ink characters.
\item \textbf{A.} OCR is Optical Character Recognition. B, C, and D are not the recognised expansion.
\item \textbf{A, B, C.} OCR, OMR, and MICR all digitise paper data. Speaker playback (D) outputs sound instead.
\item \textbf{A, B, C.} The three roles are stated correctly. D is false: all three are input technologies, not printers.
\item \textbf{A.} OCR recognizes printed characters as text. OMR detects marked responses, MICR reads magnetic-ink characters, and MIDI describes music data.

\item \textbf{C.} MICR reads magnetic-ink characters. OMR detects marks, OCR recognizes ordinary printed text, and QR recognition reads two-dimensional codes. The cited key is provisional.

\item \textbf{A.} OCR recognizes character shapes in an image and produces text. OMR detects marks, MICR reads magnetic ink, and barcode printing creates codes rather than extracting text.

\item \textbf{A.} OMR detects marked regions, such as answer bubbles. MICR reads cheque characters; the other choices concern audio output and printing.

\item \textbf{A, B, C.} OCR, OMR, and MICR have distinct roles. Option D assigns MICR's role to OMR.

\end{enumerate}
\answersubheading{14}{Output Devices}
\begin{enumerate}[label=\textbf{Q\arabic*.}, leftmargin=*]
\item \textbf{A.} MFP stands for Multifunction Printer. B, C, and D are not
  the correct expansion.
\item \textbf{B.} CRT stands for Cathode Ray Tube. A, C, and D are not the
  correct full forms.
\item \textbf{C.} A monitor presents processed output as soft copy. A scanner
  (A), a keyboard (B), and a light pen (D) are input devices.
\item \textbf{C.} A scanner takes data in, so it is an input device. A monitor
  (A), a printer (B), and a speaker (D) all present output.
\item \textbf{B.} Output shown on a screen and not printed is soft copy. Hard
  copy (A) is printed; backup copy (C) and master copy (D) are not output
  forms.
\item \textbf{C.} Output printed on paper is hard copy. Soft copy (A) is shown
  on a screen; screen copy (B) and video copy (D) are not the standard terms.
\item \textbf{D.} CRT (Cathode Ray Tube) is the older, bulky display that uses
  a vacuum tube and an electron beam. LCD (A), LED (B), and OLED (C) are flat
  panels.
\item \textbf{C.} OLED (Organic Light Emitting Diode) uses an organic compound
  and each pixel emits its own light, so no backlight is needed. LCD (A) and
  LED (B) need a backlight; CRT (D) is a vacuum-tube technology.
\item \textbf{A.} LCD (Liquid Crystal Display) uses liquid crystal cells and
  requires a backlight. CRT (B) is a vacuum tube; a plotter (C) and a speaker
  (D) are not displays.
\item \textbf{A.} Statements I and II are true: a monitor gives soft copy and
  a printer gives hard copy. III is false because a projector is an output
  device. B and C wrongly include III; D includes all three.
\item \textbf{B.} A projector throws an image onto a screen or wall as soft
  copy. A monitor (A) shows the image on its own screen, a plotter (C) draws
  on paper, and a scanner (D) is an input device.
\item \textbf{B.} A plotter is used for large engineering drawings. A laser
  printer (A) is for ordinary pages, a speaker (C) gives audio, and a flatbed
  scanner (D) is an input device.
\item \textbf{A.} A printer produces standard pages of text and graphics,
  while a plotter draws large, precise line drawings. B wrongly calls a
  printer an input device; C wrongly calls a printer's output soft copy; D
  denies the real difference.
\item \textbf{B.} A speaker converts an electrical signal into sound, so it is
  an audio output device. A microphone (A) is audio input, a scanner (C) is
  image input, and a webcam (D) is video input.
\item \textbf{A.} Headphones produce sound for a single listener, so they are
  an audio output device. Audio input (B) is a microphone; image output (C) is
  a monitor; a pointing input device (D) is a mouse.
\item \textbf{B.} Statement (ii) plotter and (iii) drum printer are output
  devices. (i) scanner and (iv) light pen are input devices, so A, C, and D
  are wrong because each includes at least one input device.
\item \textbf{B.} A Multifunction Printer (MFP) combines printing, scanning,
  copying, and often faxing. A plotter (A), a dot-matrix printer (C), and a
  projector (D) each perform only one output job.
\item \textbf{C.} A dot-matrix printer strikes a ribbon against the paper, so
  it is an impact printer. Laser (A), inkjet (B), and thermal (D) printers
  mark the paper without striking it.
\item \textbf{B.} A laser printer forms the image with toner and fuses it onto
  the paper by heat. A dot-matrix printer (A) is impact, a plotter (C) draws
  line drawings, and a projector (D) shows soft copy.
\item \textbf{A.} A plotter produces large engineering drawings, so this
  pairing is correct. A monitor (B) gives soft copy, not hard copy; a speaker
  (C) is audio output, not input; and a scanner (D) is an input device, not an
  output device.

\item \textbf{A.} The monitor presents output. B, C, and D all feed input into the computer.
\item \textbf{A.} Printers put output on paper. B, C, and D are input devices.
\item \textbf{A.} Screen display is soft copy. B, C, and D are paper-based hard copies.
\item \textbf{A, B, C.} Monitor, printer, and speaker all deliver output. The keyboard (D) supplies input.
\item \textbf{A, B, C.} The soft-copy, hard-copy, and audio roles are correct. D is false: a scanner is an input device.
\item \textbf{A.} A plotter produces large-format drawings. The other devices provide input. The cited answer key is provisional.

\item \textbf{A.} A screen display is soft copy output. Hard copy is printed output; input and storage have different roles.

\item \textbf{A.} A projector displays visual output. The other devices capture input or move a pointer.

\item \textbf{A.} A tactile display conveys information through touch, often with refreshable braille cells. The other options describe unrelated media/encoding.

\item \textbf{A, B, C.} Monitor, speaker, and plotter deliver output. A scanner captures input.

\end{enumerate}
\answersubheading{15}{Printers and Displays}
\begin{enumerate}[label=\textbf{Q\arabic*.}, leftmargin=*]
\item \textbf{C.} A dot matrix printer strikes an inked ribbon against the
  paper, so it is an impact printer. Inkjet (A), laser (B), and thermal (D)
  form the image without striking, so they are non-impact.
\item \textbf{B.} Striking an inked ribbon against the paper is exactly the
  definition of an impact printer. A non-impact printer (A) does not strike
  the paper; a page printer (C) is classified by the unit printed, not by
  contact; a plotter (D) draws lines and is not defined this way.
\item \textbf{D.} The inkjet printer is non-impact because it sprays ink
  without touching the paper. The dot matrix (A), daisy wheel (B), and drum
  printer (C) all strike a ribbon or paper and are impact printers.
\item \textbf{B.} An inkjet printer sprays droplets of liquid ink through
  nozzles. A laser printer (A) uses toner; a thermal printer (C) uses heat;
  a dot matrix printer (D) strikes a ribbon.
\item \textbf{B.} A laser printer uses a laser beam and toner in a
  xerographic process. An inkjet printer (A) sprays liquid ink; a dot matrix
  printer (C) strikes a ribbon; a thermal printer (D) uses heat on thermal
  paper.
\item \textbf{C.} A laser printer prints a whole page at a time, so it is a
  page printer. A dot matrix (A) and a daisy wheel printer (B) are character
  printers; a plotter (D) draws lines and is not a page printer.
\item \textbf{C.} A thermal printer applies heat to specially treated paper,
  which is why it is used for receipts and labels. A laser printer (A) uses
  a laser beam and toner; an inkjet printer (B) sprays liquid ink; a dot
  matrix printer (D) strikes a ribbon.
\item \textbf{B.} A 3D printer works by additive manufacturing, adding
  material layer by layer to build an object. Striking a ribbon (A) is dot
  matrix, spraying ink (C) is inkjet, and melting toner (D) is laser --- all
  two-dimensional paper printing.
\item \textbf{B.} A plotter draws vector line graphics with pens, so it suits
  large engineering drawings, maps, and CAD plans. Inkjet (A) and laser (C)
  printers print dots/characters rather than drawn lines; a dot matrix printer
  (D) is a low-quality character printer.
\item \textbf{B.} An inkjet printer sprays liquid ink, while a laser printer
  uses a laser beam and toner --- the correct contrast. A reverses the
  consumables; C is wrong because both are non-impact; D swaps their print
  units (a laser is a page printer, not a character printer).
\item \textbf{B.} A printer that prints one whole line at a time, such as a
  drum or chain printer, is a line printer. A character printer (A) prints one
  character at a time; a page printer (C) prints a whole page; a plotter (D)
  draws lines and is not a printer of this kind.
\item \textbf{A.} DPI stands for Dots Per Inch, a measure of print/display
  sharpness. B, C, and D are not the correct expansions.
\item \textbf{B.} Resolution is the number of pixels a display shows,
  written as columns $\times$ rows. Refresh rate (A) is the redraws per
  second; response time (C) is the pixel switching time; contrast ratio (D)
  compares white and black levels.
\item \textbf{A.} A pixel (picture element) is the smallest addressable
  element of a display. A byte (B) is a unit of data; a bitmap (C) is an
  image representation; a raster (D) is a pixel grid, not a single element.
\item \textbf{A.} Refresh rate is measured in hertz (Hz), the number of
  redraws per second. Milliseconds (B) measure response time; pixels (C)
  measure resolution; nits (D) measure brightness.
\item \textbf{B.} Response time is measured in milliseconds (ms), the time a
  pixel takes to change state. Hertz (A) measures refresh rate; DPI (C)
  measures sharpness; nits (D) measure brightness.
\item \textbf{B.} The ratio between the brightest white and darkest black is
  the contrast ratio. Refresh rate (A) is redraws per second; resolution (C)
  is the pixel count; response time (D) is the pixel switching time.
\item \textbf{B.} Brightness, the light a display emits, is commonly measured
  in nits (cd/m$^2$). Hertz (A) measures refresh rate; milliseconds (C)
  measure response time; DPI (D) measures sharpness.
\item \textbf{D.} The dot matrix printer is an \textbf{impact} printer,
  because it strikes the ribbon; calling it non-impact is the incorrect
  statement. Inkjet (A), laser (B), and thermal (C) printers are genuinely
  non-impact.
\item \textbf{B.} Statements I and II are correct: a dot matrix printer is
  impact, and a laser printer is non-impact. III is false because a thermal
  printer is non-impact, not impact. A omits II; C includes false III and
  omits I; D includes false III.

\item \textbf{A.} Dot-matrix pins strike a ribbon. B, C, and D are non-impact technologies.
\item \textbf{A.} Laser printers lead in fast bulk printing. B and C are slow impact types; D is not a computer printer.
\item \textbf{A.} Resolution counts pixels. B, C, and D are unrelated physical units.
\item \textbf{A, B.} Dot-matrix and daisy-wheel printers strike the paper. Laser (C) and inkjet (D) are non-impact.
\item \textbf{A, B, C.} The laser, inkjet, and resolution facts are correct. D is false: dot-matrix uses pins, not lasers.
\item \textbf{A.} A dot-matrix printer is an impact printer that strikes a ribbon. Inkjet sprays droplets, laser uses toner and a drum, and thermal printing uses heat-sensitive media.

\item \textbf{A.} A laser printer forms images without striking paper. Daisy-wheel, dot-matrix, and line printers use impact mechanisms.

\item \textbf{A.} A printer produces a physical paper copy. The others provide soft-copy visual or audio output.

\item \textbf{A.} Laser printers use toner. Ink ribbons are common in impact printers, liquid ink is used in inkjets, and thermal paper is used by thermal printers.

\item \textbf{A, B, C.} The first three describe printer technologies. A monitor produces soft-copy output, not hard copy.

\end{enumerate}
\answersubheading{16}{Computer Ports}
\begin{enumerate}[label=\textbf{Q\arabic*.}, leftmargin=*]
\item \textbf{A.} USB stands for Universal Serial Bus. B, C, and D are not
  expansions of the abbreviation.
\item \textbf{A.} HDMI stands for High-Definition Multimedia Interface. B
  drops ``Multimedia''; C changes ``Definition'' to ``Digital''; D uses
  ``Memory'' instead of ``Multimedia''.
\item \textbf{A.} VGA stands for Video Graphics Array. B, C, and D are not the
  correct full forms.
\item \textbf{A.} RJ-45 stands for Registered Jack 45. B, C, and D are not the
  correct expansion of ``RJ''.
\item \textbf{A.} SATA stands for Serial Advanced Technology Attachment (Serial
  ATA). B, C, and D are not the correct expansion.
\item \textbf{D.} RJ-45 is the Ethernet network connector, not a video output
  port. HDMI (A), DisplayPort (B), and VGA (C) are all display ports.
\item \textbf{C.} USB Type-C is reversible and can be inserted either way up.
  Type-A (A), Type-B (B), and Micro-USB (D) are not orientation-independent.
\item \textbf{B.} HDMI is digital and carries audio, while VGA is analog and
  carries video only. A reverses the classification; C wrongly gives VGA
  audio; D swaps the two.
\item \textbf{B.} RJ-45 is the eight-pin connector for a wired Ethernet LAN.
  RJ-11 (A) is the smaller telephone jack; USB-C (C) and HDMI (D) are not
  network connectors.
\item \textbf{A.} Statements I and II are correct: USB supplies power and
  supports hot swapping. III is false because VGA carries video only. B
  includes false III; C includes false III and omits II; D includes false III.
\item \textbf{B.} A port number such as 80 or 443 is a logical (software) port
  that names a service. A wrongly calls it a physical socket; C confuses it
  with pin count; D confuses it with cable speed.
\item \textbf{B.} DisplayPort is designed mainly for computers and monitors
  and supports high refresh rates and multi-monitor setups. VGA (A) is an
  older analog port; RJ-45 (C) is networking; SATA (D) is storage.
\item \textbf{A.} VGA is the analog, video-only port now largely replaced by
  HDMI and DisplayPort. HDMI (B) and DisplayPort (C) are digital; USB-C (D)
  is a connector type, not the analog video port.
\item \textbf{B.} Hot swapping is the ability to attach or remove a device
  while the computer is running. Cold boot (A), fragmentation (C), and
  multiplexing (D) are unrelated.
\item \textbf{A.} SATA connects internal drives (HDD, SSD, optical drive) to
  the motherboard. HDMI (B) is video; RJ-45 (C) is networking; the 3.5 mm jack
  (D) is audio.
\item \textbf{B.} TRRS adds a microphone channel to stereo audio, while TRS
  (A) carries stereo audio only. RJ-45 (C) is networking; VGA (D) is video.
\item \textbf{B.} An Ethernet (RJ-45) port typically connects a wired network
  router, so this pairing is correct. A pairs HDMI with headphones (HDMI is
  for displays); C pairs VGA with a printer; D pairs SATA with a keyboard
  (SATA is for storage drives).
\item \textbf{B.} SATA is a storage interface for drives. USB (A), HDMI (C),
  and the audio jack (D) are general-purpose peripheral or display ports, not
  storage interfaces.
\item \textbf{A.} eSATA (external SATA) connects an external storage drive.
  B, C, and D describe monitor, microphone, and network connections, which use
  other ports.
\item \textbf{A.} This statement is precise: USB is the standard, USB-C is a
  connector type, and a USB port is the socket. B wrongly calls the connector
  the socket; C wrongly says port and connector are identical; D wrongly calls
  HDMI analog (it is digital).

\item \textbf{A.} USB is the standard printer interface. VGA (B) drives displays; C and D do not connect printers.
\item \textbf{A.} HDMI combines audio and video. B, C, and D name signals HDMI does not exclusively carry.
\item \textbf{A.} RJ-45 terminates Ethernet LAN cables. B, C, and D are other connectors.
\item \textbf{A, B, C.} USB, HDMI, and Ethernet are interfaces. A stapler slot (D) is not a computer port.
\item \textbf{A, B, C.} The USB, HDMI, and RJ-45 roles are correct. D is false: VGA carries analog video, not digital audio.
\item \textbf{A.} HDMI commonly carries digital audio and video. RJ-11 is associated with telephone wiring, PS/2 with legacy keyboard/mouse connections, and VGA carries analog video.

\item \textbf{A.} RJ-45 is commonly used for Ethernet. HDMI and VGA serve displays; a 3.5 mm jack serves analog audio.

\item \textbf{A.} USB Type-C is a reversible physical connector. It can carry different signals depending on implementation; the other options confuse it with unrelated interfaces.

\item \textbf{A.} TCP port 80 is a logical endpoint number. The others are physical connectors.

\item \textbf{A, B, C.} The first three pairings are standard uses. MICR, not VGA, reads magnetic-ink characters.

\end{enumerate}
\answersubheading{65}{Current Multimedia, GUI and Media Formats}
\begin{enumerate}[label=\textbf{Q\arabic*.}, leftmargin=*]
\item \textbf{B.} GUI is Graphical User Interface. A, C, and D are not the
  recognised expansion.
\item \textbf{A.} WYSIWYG is What You See Is What You Get. B, C, and D are
  invented phrases.
\item \textbf{D.} MP3 is MPEG-1 Audio Layer III. A names the wrong MPEG layer;
  B and C incorrectly insert ``Video'' and ``Moving Picture''.
\item \textbf{A.} PNG is Portable Network Graphics. B, C, and D are plausible
  but incorrect expansions.
\item \textbf{C.} JPEG is Joint Photographic Experts Group. A, B, and D alter
  the words ``Photographic'', ``Experts'', or ``Encoding''.
\item \textbf{D.} MP4 is MPEG-4 Part 14. A, B, and C substitute ``Portable'',
  ``Picture'', or ``Player'' for ``Part''.
\item \textbf{B.} Measuring a continuous signal at regular intervals is
  sampling. Quantization (A) maps each sample to discrete levels;
  synchronization (C) aligns streams in time; compression (D) reduces file
  size.
\item \textbf{C.} Mapping each sample to one of a fixed set of levels is
  quantization. Sampling (A) is the measurement step; synchronization (B) and
  buffering (D) are unrelated.
\item \textbf{C.} PNG is lossless. JPEG (A), MP3 (B), and MP4 (D) are all lossy
  formats.
\item \textbf{C.} GIF is lossless and limited to 256 colours. A wrongly calls
  PNG lossy and non-transparent; B wrongly calls JPEG lossless; D wrongly
  describes BMP as a small web format.
\item \textbf{B.} GIF supports simple animation. PNG (A) and JPEG (C) are
  static image formats, and BMP (D) does not support animation.
\item \textbf{C.} JPEG is the lossy format designed for photographs. PNG (A)
  and GIF (B) are better for graphics with exact colours, and TXT (D) is a
  plain-text format, not an image format.
\item \textbf{C.} Synchronization keeps audio and video aligned in time.
  Compression (A), sampling (B), and quantization (D) are digitisation or
  size-reduction steps, not timing alignment.
\item \textbf{A.} WIMP = Windows, Icons, Menus, Pointers. B, C, and D replace
  correct elements with wrong ones.
\item \textbf{B.} A dialogue box is a small window that collects input or
  confirms an action. A desktop (A), taskbar (C), and shortcut (D) are other
  GUI features, not input windows.
\item \textbf{B.} WYSIWYG means the screen shows the document exactly as it
  will be printed. A, C, and D describe cloud storage, spell-check, and
  compression instead.
\item \textbf{B.} A plotter draws vector graphics and is used for large
  engineering drawings and maps. It is an output device, so it does not print
  ordinary text (A), scan photographs (C), or record sound (D).
\item \textbf{D.} A laser printer is non-impact; it forms the image without
  striking the paper. A dot-matrix (A), daisy-wheel (B), and line printer (C)
  are all impact printers.
\item \textbf{C.} JPEG is lossy. PNG (A), GIF (B), and ZIP (D) are all
  lossless.
\item \textbf{D.} All three statements are correct: MP3 is lossy audio, MP4 is
  a container holding video and audio, and PNG is lossless and supports
  transparency. A, B, and C each omit at least one true statement.

\item \textbf{A.} GUI is Graphical User Interface. B, C, and D are not the recognised expansion.
\item \textbf{A.} GIF supports simple web animation. B, C, and D are document or text formats.
\item \textbf{A.} MP3 stores audio. B, C, and D are other data kinds.
\item \textbf{A, B, C.} Windows, icons, and menus are GUI parts. Gears (D) are hardware.
\item \textbf{A, B, C.} The image, audio, and video matches are correct. D is false: spreadsheets use .xlsx.
\item \textbf{A.} A plotter produces large-format output. The other devices provide input.

\item \textbf{A.} A plotter outputs drawings. A scanner, light pen, and digitizer tablet provide input.

\item \textbf{A.} Sampling records signal values at discrete times. The other choices are unrelated.

\item \textbf{A.} JPEG is an image format commonly using lossy compression. DOCX is a document package, CSV is delimited text, and EXE is executable.

\item \textbf{A, B, C.} Monitor, speaker, and projector output visual/audio content. A scanner captures input.

\end{enumerate}
\unitopener{5}{Software, Operating Systems, Windows \& Utilities}
\lessonheading{17}{Hardware, Software and Firmware}
\begin{enumerate}[label=\textbf{Q\arabic*.}, leftmargin=*]
\item \textit{[Original]} Expand OS.
  \begin{enumerate}[label=\Alph*.]
    \item Operating System
    \item Open Source
    \item Output Setting
    \item Object Store
  \end{enumerate}
\item \textit{[Original]} Expand BIOS.
  \begin{enumerate}[label=\Alph*.]
    \item Binary Input Output Setup
    \item Basic Input Output System
    \item Basic Internal Operating System
    \item Boot Input Output Source
  \end{enumerate}
\item \textit{[Original]} Which of the following is \textbf{hardware}?
  \begin{enumerate}[label=\Alph*.]
    \item MS Word
    \item A printer driver
    \item The CPU
    \item Windows operating system
  \end{enumerate}
\item \textit{[Original]} Which of the following is \textbf{software}?
  \begin{enumerate}[label=\Alph*.]
    \item Keyboard
    \item Monitor
    \item MS Word
    \item CPU
  \end{enumerate}
\item \textit{[Original]} A device driver is best classified as:
  \begin{enumerate}[label=\Alph*.]
    \item hardware
    \item software
    \item firmware
    \item humanware
  \end{enumerate}
\item \textit{[Original]} The operating system is an example of:
  \begin{enumerate}[label=\Alph*.]
    \item application software
    \item system software
    \item hardware
    \item firmware
  \end{enumerate}
\item \textit{[Original]} Firmware is best described as:
  \begin{enumerate}[label=\Alph*.]
    \item a physical circuit inside the processor
    \item software stored in non-volatile memory inside a hardware device
    \item an application program for editing documents
    \item a person who operates the computer
  \end{enumerate}
\item \textit{[Original]} Which of the following is normally stored in
  non-volatile memory such as ROM or flash and is used to start a device?
  \begin{enumerate}[label=\Alph*.]
    \item Firmware
    \item RAM
    \item A document file
    \item A user account
  \end{enumerate}
\item \textit{[Original]} \textbf{Humanware} refers to:
  \begin{enumerate}[label=\Alph*.]
    \item the physical parts of the computer
    \item the programs stored on disk
    \item the people who use and operate the computer system
    \item the startup code in ROM
  \end{enumerate}
\item \textit{[Original]} Which statement is \textbf{NOT} correct?
  \begin{enumerate}[label=\Alph*.]
    \item A device driver is a hardware component.
    \item The operating system is system software.
    \item Firmware is stored in non-volatile memory.
    \item Humanware is the people component of a system.
  \end{enumerate}
\item \textit{[Original]} Which pairing is correctly matched?
  \begin{enumerate}[label=\Alph*.]
    \item Hardware --- MS Word
    \item Software --- keyboard
    \item Firmware --- BIOS
    \item Humanware --- CPU
  \end{enumerate}
\item \textit{[Original]} Arrange the computing stack from the bottom
  (lowest) to the top (highest).
  \begin{enumerate}[label=\Alph*.]
    \item Application, operating system, device driver, firmware, hardware
    \item Hardware, firmware, device driver, operating system, application
    \item Hardware, application, operating system, firmware, device driver
    \item Firmware, hardware, operating system, device driver, application
  \end{enumerate}
\item \textit{[Original]} Evaluate the statements:
  \begin{enumerate}[label=\Roman*.]
    \item A device driver is a software component.
    \item The operating system is system software.
    \item Firmware is stored on the hardware in non-volatile memory.
  \end{enumerate}
  \begin{enumerate}[label=\Alph*.]
    \item I only
    \item I and II only
    \item II and III only
    \item I, II, and III
  \end{enumerate}
\item \textit{[Original]} Which layer lets the operating system communicate
  with a specific hardware device?
  \begin{enumerate}[label=\Alph*.]
    \item The application software
    \item The device driver
    \item The humanware
    \item The BIOS
  \end{enumerate}
\item \textit{[Original]} Which of the following is an example of
  \textbf{application software}?
  \begin{enumerate}[label=\Alph*.]
    \item The operating system
    \item MS Word
    \item The BIOS
    \item A printer driver
  \end{enumerate}
\item \textit{[Original]} Which of the following is an example of
  \textbf{system software}?
  \begin{enumerate}[label=\Alph*.]
    \item MS Excel
    \item The operating system
    \item A computer game
    \item A photo editor
  \end{enumerate}
\item \textit{[Original]} Firmware is usually stored in:
  \begin{enumerate}[label=\Alph*.]
    \item RAM
    \item ROM or flash memory
    \item the CPU registers
    \item a user folder on the hard disk
  \end{enumerate}
\item \textit{[Original]} Which of the following is \textbf{NOT} a type of
  hardware?
  \begin{enumerate}[label=\Alph*.]
    \item Motherboard
    \item Printer
    \item Device driver
    \item Keyboard
  \end{enumerate}
\item \textit{[Original]} The human component of a computer system --- the
  users, programmers, and administrators --- is called:
  \begin{enumerate}[label=\Alph*.]
    \item hardware
    \item software
    \item firmware
    \item humanware
  \end{enumerate}
\item \textit{[Original]} A printer driver is needed because it:
  \begin{enumerate}[label=\Alph*.]
    \item is the physical printer component
    \item lets the operating system communicate with the printer
    \item is the printer's firmware
    \item permanently stores the printed document
  \end{enumerate}

\item \textit{[JKSSB FAA-pattern]} Which of the following is hardware?
  \begin{enumerate}[label=\Alph*.]
    \item Keyboard
    \item MS Word
    \item Antivirus program
    \item Spreadsheet formula
  \end{enumerate}
\item \textit{[JKPSI-pattern]} Firmware differs from ordinary software because it is:
  \begin{enumerate}[label=\Alph*.]
    \item stored in non-volatile chips and tied to hardware.
    \item always typed fresh by the user.
    \item only heard through speakers.
    \item printed on paper before use.
  \end{enumerate}
\item \textit{[SSC CGL-pattern]} A device driver is an example of:
  \begin{enumerate}[label=\Alph*.]
    \item system software.
    \item an input device.
    \item an output device.
    \item a storage disk.
  \end{enumerate}
\item \textit{[JKSSB-pattern]} \textbf{(MSQ: select ALL correct options.)} Which of the following are hardware?
  \begin{enumerate}[label=\Alph*.]
    \item Monitor
    \item Hard disk
    \item Operating system
    \item Printer
  \end{enumerate}
\item \textit{[SSC CGL-pattern]} \textbf{(MSQ: select ALL correct options.)} Which statements about hardware, software, and firmware are correct?
  \begin{enumerate}[label=\Alph*.]
    \item Hardware refers to physical components.
    \item Software refers to programs and instructions.
    \item Firmware is software embedded for hardware control.
    \item A device driver is a hardware component.
  \end{enumerate}
\item \textit{[Original]} Which term describes the physical components of a computer?
  \begin{enumerate}[label=\Alph*.]
\item Hardware
   \item Software
   \item Firmware
   \item Data
  \end{enumerate}

\item \textit{[PYQ-adapted: JKSSB FAA 2024, QCRIAS item 668; SRC-16/SRC-19]} Firmware is typically stored in which kind of memory?
  \begin{enumerate}[label=\Alph*.]
\item Volatile RAM only
   \item Non-volatile memory such as ROM/flash
   \item CPU cache only
   \item A printer buffer only
  \end{enumerate}

\item \textit{[Original]} A device driver is best classified as:
  \begin{enumerate}[label=\Alph*.]
\item Software that helps the operating system communicate with hardware
   \item A physical input device
   \item A storage medium
   \item A type of CPU register
  \end{enumerate}

\item \textit{[Original]} Which is an example of application software?
  \begin{enumerate}[label=\Alph*.]
\item Word processor
   \item Device driver
   \item Boot firmware
   \item Memory controller hardware
  \end{enumerate}

\item \textit{[Original, MSQ drill]} \textbf{(MSQ: select ALL correct options.)} Which statements correctly classify components?
  \begin{enumerate}[label=\Alph*.]
\item A keyboard is hardware.
   \item A device driver is software.
   \item Firmware is low-level software associated with hardware control.
   \item A spreadsheet file is a CPU register.
  \end{enumerate}

\end{enumerate}
\lessonheading{18}{Software Categories}
\begin{enumerate}[label=\textbf{Q\arabic*.}, leftmargin=*]
\item \textit{[Original]} Software is best defined as:
  \begin{enumerate}[label=\Alph*.]
    \item the physical parts of a computer that can be touched.
    \item a set of programs or instructions that tell the computer what to do.
    \item only the operating system supplied with the machine.
    \item the data kept in secondary storage.
  \end{enumerate}
\item \textit{[Original]} Which of the following is \textbf{system software}?
  \begin{enumerate}[label=\Alph*.]
    \item MS Word
    \item Operating System
    \item A payroll package
    \item A media player
  \end{enumerate}
\item \textit{[Original]} Expand OS.
  \begin{enumerate}[label=\Alph*.]
    \item Open Source
    \item Operating System
    \item Output Storage
    \item Office Suite
  \end{enumerate}
\item \textit{[Original]} Which of the following is \textbf{application
  software}?
  \begin{enumerate}[label=\Alph*.]
    \item A device driver
    \item The operating system
    \item MS Excel
    \item The BIOS
  \end{enumerate}
\item \textit{[Original]} Which one of the following is \textbf{utility
  software}?
  \begin{enumerate}[label=\Alph*.]
    \item An antivirus program
    \item A web browser
    \item MS PowerPoint
    \item A compiler
  \end{enumerate}
\item \textit{[Original]} In the context of programming, expand IDE.
  \begin{enumerate}[label=\Alph*.]
    \item Integrated Development Environment
    \item Internal Data Engine
    \item Integrated Disk Editor
    \item Internet Data Exchange
  \end{enumerate}
\item \textit{[Original]} Compilers, interpreters, and assemblers together
  belong to:
  \begin{enumerate}[label=\Alph*.]
    \item application software
    \item programming software (language translators)
    \item utility software
    \item embedded software
  \end{enumerate}
\item \textit{[Original]} Software that is permanently stored inside a washing
  machine or a microwave oven in order to control it is called:
  \begin{enumerate}[label=\Alph*.]
    \item application software
    \item embedded software
    \item freeware
    \item middleware
  \end{enumerate}
\item \textit{[Original]} Middleware is best described as:
  \begin{enumerate}[label=\Alph*.]
    \item software that connects different applications and systems, sitting
      between the operating system and application software.
    \item the physical wiring that connects two computers.
    \item a free trial version of a software product.
    \item the hardware located inside the CPU.
  \end{enumerate}
\item \textit{[Original]} Which of the following is \textbf{NOT} system
  software?
  \begin{enumerate}[label=\Alph*.]
    \item Operating system
    \item Device driver
    \item MS Word
    \item BIOS
  \end{enumerate}
\item \textit{[Original]} Which statement correctly distinguishes open-source
  software, freeware, and shareware?
  \begin{enumerate}[label=\Alph*.]
    \item All three make the source code freely available.
    \item Open-source makes the source code available; freeware is free of
      cost but its source may be closed; shareware is a trial version that
      later requires payment.
    \item Freeware always provides the source code, and shareware is free
      forever.
    \item Open-source software and shareware mean the same thing.
  \end{enumerate}
\item \textit{[Original]} Open-source software means software:
  \begin{enumerate}[label=\Alph*.]
    \item that is free of cost only.
    \item whose source code is available to view, modify, and redistribute
      under its licence.
    \item that can be used only for a short trial period.
    \item that carries no copyright of any kind.
  \end{enumerate}
\item \textit{[Original]} Freeware is software that:
  \begin{enumerate}[label=\Alph*.]
    \item is distributed free of cost, although its source code may not be
      available.
    \item always includes its full source code.
    \item must be paid for after a fixed trial period.
    \item has no owner and no licence.
  \end{enumerate}
\item \textit{[Original]} Shareware is software that:
  \begin{enumerate}[label=\Alph*.]
    \item is fully free and provides its source code.
    \item is offered as a limited trial and requires payment for continued or
      full use.
    \item forms part of the operating system.
    \item is exactly the same as public-domain software.
  \end{enumerate}
\item \textit{[Original]} Public-domain software:
  \begin{enumerate}[label=\Alph*.]
    \item is protected by strict copyright.
    \item has no copyright and may be used freely.
    \item is a limited trial version.
    \item always requires a licence key before use.
  \end{enumerate}
\item \textit{[Original]} Evaluate the statements:
  \begin{enumerate}[label=\Roman*.]
    \item An operating system is system software.
    \item Freeware and open-source software mean exactly the same thing.
    \item Shareware is usually distributed as a trial version.
  \end{enumerate}
  \begin{enumerate}[label=\Alph*.]
    \item I and II only
    \item I and III only
    \item II and III only
    \item I, II, and III
  \end{enumerate}
\item \textit{[Original]} Expand GPL.
  \begin{enumerate}[label=\Alph*.]
    \item General Public License
    \item GNU General Public License
    \item General Purpose Language
    \item Global Program Loader
  \end{enumerate}
\item \textit{[Original]} The main purpose of utility software is to:
  \begin{enumerate}[label=\Alph*.]
    \item help users create documents and presentations.
    \item perform maintenance tasks such as antivirus scanning, disk
      cleanup, and backup.
    \item translate source code into machine code.
    \item connect two different applications to each other.
  \end{enumerate}
\item \textit{[Original]} Application software is commonly divided into:
  \begin{enumerate}[label=\Alph*.]
    \item general-purpose and specific-purpose (custom) software.
    \item system software and utility software.
    \item freeware and shareware.
    \item primary and secondary software.
  \end{enumerate}
\item \textit{[Original]} Which of the following pairs is correctly matched?
  \begin{enumerate}[label=\Alph*.]
    \item Proprietary software --- source code freely available to everyone
    \item Open-source software --- source code available under a licence
    \item Freeware --- software that must always be paid for
    \item Shareware --- software that never requires payment
  \end{enumerate}

\item \textit{[JKSSB FAA-pattern]} Which of the following is system software?
  \begin{enumerate}[label=\Alph*.]
    \item Operating system
    \item Railway timetable book
    \item Cricket bat
    \item Paper file
  \end{enumerate}
\item \textit{[JKPSI-pattern]} MS Excel is best classified as:
  \begin{enumerate}[label=\Alph*.]
    \item application software.
    \item system boot firmware.
    \item a hardware device.
    \item a network cable.
  \end{enumerate}
\item \textit{[SSC CGL-pattern]} Which of the following is an example of application software?
  \begin{enumerate}[label=\Alph*.]
    \item Word processor
    \item BIOS chip
    \item Motherboard
    \item Power supply
  \end{enumerate}
\item \textit{[JKSSB-pattern]} \textbf{(MSQ: select ALL correct options.)} Which of the following are application software?
  \begin{enumerate}[label=\Alph*.]
    \item Word processor
    \item Spreadsheet package
    \item Presentation package
    \item Operating system kernel
  \end{enumerate}
\item \textit{[SSC CGL-pattern]} \textbf{(MSQ: select ALL correct options.)} Which statements about software categories are correct?
  \begin{enumerate}[label=\Alph*.]
    \item System software manages hardware resources.
    \item Application software serves end-user tasks.
    \item Utility programs assist maintenance and management.
    \item An operating system is application software.
  \end{enumerate}
\item \textit{[PYQ-adapted: JKSSB Sub-Inspector 2022, QCRIAS item 329; SRC-17/SRC-22]} Which software is designed to control computer hardware and provide services for applications?
  \begin{enumerate}[label=\Alph*.]
\item System software
   \item Application software
   \item User data
   \item Presentation theme
  \end{enumerate}

\item \textit{[Original]} Which is a utility program?
  \begin{enumerate}[label=\Alph*.]
\item Disk cleanup tool
   \item Word processor
   \item Presentation file
   \item Device cable
  \end{enumerate}

\item \textit{[Original]} Freeware is primarily classified by its:
  \begin{enumerate}[label=\Alph*.]
\item Price to the user
   \item Source-code license alone
   \item Processor architecture
   \item Storage medium
  \end{enumerate}

\item \textit{[Original]} Open-source software is distinguished primarily by access to its:
  \begin{enumerate}[label=\Alph*.]
\item Source code under a license that permits specified use/modification
   \item Hardware schematics in every case
   \item Free support hotline
   \item User's private data
  \end{enumerate}

\item \textit{[Original, MSQ drill]} \textbf{(MSQ: select ALL correct options.)} Which are software categories or licensing/distribution labels?
  \begin{enumerate}[label=\Alph*.]
\item System software
   \item Application software
   \item Utility software
   \item Open source
  \end{enumerate}

\end{enumerate}
\lessonheading{19}{Operating System}
\begin{enumerate}[label=\textbf{Q\arabic*.}, leftmargin=*]
\item \textit{[PYQ WO 2026 Q61, key:provisional]} The operating system is best
  described as:
  \begin{enumerate}[label=\Alph*.]
    \item a hardware component located inside the CPU.
    \item a layer that provides a user-friendly interface between the user and
      the computer hardware.
    \item an application program such as a word processor.
    \item a type of secondary storage device.
  \end{enumerate}
\item \textit{[PYQ WO 2026 Q62, key:provisional]} While a computer is running,
  the operating system is resident in:
  \begin{enumerate}[label=\Alph*.]
    \item secondary storage
    \item an offline backup tape
    \item primary storage (main memory)
    \item the display adapter
  \end{enumerate}
\item \textit{[PYQ WO 2026 Q63, key:provisional]} Multitasking means:
  \begin{enumerate}[label=\Alph*.]
    \item connecting more than one printer to a computer.
    \item storing more than one file on a disk.
    \item using more than one monitor at a time.
    \item running more than one program at a time.
  \end{enumerate}
\item \textit{[Original]} Expand OS.
  \begin{enumerate}[label=\Alph*.]
    \item Open Source
    \item Online Service
    \item Operating System
    \item Output Storage
  \end{enumerate}
\item \textit{[Original]} Which component is the core part of the operating
  system that interacts directly with the hardware?
  \begin{enumerate}[label=\Alph*.]
    \item Kernel
    \item GUI
    \item Bootstrap
    \item Compiler
  \end{enumerate}
\item \textit{[Original]} Which of the following is \textbf{NOT} a function of
  the operating system?
  \begin{enumerate}[label=\Alph*.]
    \item Process management
    \item Memory management
    \item File management
    \item Designing web pages
  \end{enumerate}
\item \textit{[Original]} GUI stands for:
  \begin{enumerate}[label=\Alph*.]
    \item General User Interface
    \item Graphical Utility Integration
    \item Guided Utility Interface
    \item Graphical User Interface
  \end{enumerate}
\item \textit{[Original]} In which type of interface does the user type text
  commands to give instructions to the computer?
  \begin{enumerate}[label=\Alph*.]
    \item GUI
    \item CLI
    \item Ribbon
    \item Desktop
  \end{enumerate}
\item \textit{[Original]} The process of starting a computer and loading the
  operating system into main memory is called:
  \begin{enumerate}[label=\Alph*.]
    \item Booting
    \item Formatting
    \item Compiling
    \item Defragmenting
  \end{enumerate}
\item \textit{[Original]} Which firmware initialises the hardware and then
  loads the operating system during startup?
  \begin{enumerate}[label=\Alph*.]
    \item Device driver
    \item Kernel
    \item BIOS/UEFI
    \item Compiler
  \end{enumerate}
\item \textit{[Original]} A program that is currently executing is called a:
  \begin{enumerate}[label=\Alph*.]
    \item process
    \item folder
    \item shortcut
    \item file
  \end{enumerate}
\item \textit{[Original]} Which management function decides how much main
  memory each running program receives?
  \begin{enumerate}[label=\Alph*.]
    \item Process management
    \item Memory management
    \item File management
    \item Device management
  \end{enumerate}
\item \textit{[Original]} Which of the following is handled by file management?
  \begin{enumerate}[label=\Alph*.]
    \item Scheduling the CPU among processes
    \item Allocating RAM to programs
    \item Creating, reading, writing, and deleting files and folders
    \item Sending print jobs to a printer
  \end{enumerate}
\item \textit{[Original]} The software that lets the operating system
  communicate with a specific hardware device is called a:
  \begin{enumerate}[label=\Alph*.]
    \item compiler
    \item browser
    \item device driver
    \item spreadsheet
  \end{enumerate}
\item \textit{[Original]} Which pairing is correctly matched?
  \begin{enumerate}[label=\Alph*.]
    \item Kernel --- the core of the operating system
    \item Device driver --- an application for editing text
    \item Bootstrap --- a type of secondary storage
    \item GUI --- a command-line-only interface
  \end{enumerate}
\item \textit{[Original]} Time-sharing allows:
  \begin{enumerate}[label=\Alph*.]
    \item a computer to work without an operating system.
    \item many users to use a single computer at the same time.
    \item files to be stored without a file system.
    \item hardware to run without electricity.
  \end{enumerate}
\item \textit{[Original]} Which of the following is the most common example of
  an open-source operating system?
  \begin{enumerate}[label=\Alph*.]
    \item Windows 11
    \item macOS
    \item Android
    \item MS Word
  \end{enumerate}
\item \textit{[Original]} Evaluate the statements:
  \begin{enumerate}[label=\Roman*.]
    \item The operating system provides a user-friendly interface between the
      user and the hardware.
    \item The kernel is the core component of the operating system.
    \item An operating system is an application program.
  \end{enumerate}
  \begin{enumerate}[label=\Alph*.]
    \item I and II only
    \item I and III only
    \item II and III only
    \item I, II, and III
  \end{enumerate}
\item \textit{[Original]} Logging in with a username and password is part of:
  \begin{enumerate}[label=\Alph*.]
    \item file compression.
    \item disk defragmentation.
    \item security management (authentication).
    \item print spooling.
  \end{enumerate}
\item \textit{[Original]} Why should the operating system and system software
  be updated regularly?
  \begin{enumerate}[label=\Alph*.]
    \item To increase the physical size of the hard disk
    \item To fix security holes and add improvements
    \item To convert the GUI into hardware
    \item To remove the need for a kernel
  \end{enumerate}

\item \textit{[JKSSB FAA-pattern]} Which of the following is an operating system?
  \begin{enumerate}[label=\Alph*.]
    \item Windows
    \item MS Paint brush
    \item Printer paper
    \item Speaker wire
  \end{enumerate}
\item \textit{[JKPSI-pattern]} The core of an operating system is called the:
  \begin{enumerate}[label=\Alph*.]
    \item kernel.
    \item pencil.
    \item envelope.
    \item stamp.
  \end{enumerate}
\item \textit{[SSC CGL-pattern]} Which of the following is an open-source operating system?
  \begin{enumerate}[label=\Alph*.]
    \item Linux
    \item A printed book
    \item A wall clock
    \item A desk lamp
  \end{enumerate}
\item \textit{[JKSSB-pattern]} \textbf{(MSQ: select ALL correct options.)} Which of the following are operating systems?
  \begin{enumerate}[label=\Alph*.]
    \item Windows
    \item Linux
    \item MS Word
    \item Android
  \end{enumerate}
\item \textit{[SSC CGL-pattern]} \textbf{(MSQ: select ALL correct options.)} Which statements about operating systems are correct?
  \begin{enumerate}[label=\Alph*.]
    \item The OS manages processes and memory.
    \item The OS handles file and device management.
    \item Booting loads the OS into memory.
    \item A word processor is an operating system.
  \end{enumerate}
\item \textit{[PYQ-adapted: JKSSB Website Operator 2026 Q61, provisional key; SRC-04/SRC-07]} Which software layer provides common services and an interface between users/applications and computer hardware?
  \begin{enumerate}[label=\Alph*.]
\item Operating system
   \item Spreadsheet
   \item Web page
   \item Device cable
  \end{enumerate}

\item \textit{[PYQ-adapted: JKSSB Junior Assistant 2026 Q70, final key; SRC-03/SRC-06]} Which action is a recognized way an operating system can recover from deadlock?
  \begin{enumerate}[label=\Alph*.]
\item Resource preemption
   \item Increasing screen brightness
   \item Formatting every disk automatically
   \item Changing a document font
  \end{enumerate}

\item \textit{[Original]} Which OS function organizes files into named folders and manages their storage?
  \begin{enumerate}[label=\Alph*.]
\item File-system management
   \item Slide animation
   \item OCR
   \item Mail merge
  \end{enumerate}

\item \textit{[Original]} In a multitasking operating system, the CPU can:
  \begin{enumerate}[label=\Alph*.]
\item Switch among several tasks so they make progress
   \item Execute every program on a separate physical CPU
   \item Store all programs permanently in registers
   \item Eliminate the need for memory
  \end{enumerate}

\item \textit{[Original, MSQ drill]} \textbf{(MSQ: select ALL correct options.)} Which are common operating-system responsibilities?
  \begin{enumerate}[label=\Alph*.]
\item Process scheduling
   \item Memory management
   \item File management
   \item Editing a user's video as its only purpose
  \end{enumerate}

\end{enumerate}
\lessonheading{20}{Windows Explorer}
\begin{enumerate}[label=\textbf{Q\arabic*.}, leftmargin=*]
\item \textit{[Original]} What is Windows Explorer called in Windows 10 and
  Windows 11?
  \begin{enumerate}[label=\Alph*.]
    \item Control Panel
    \item File Explorer
    \item Task Manager
    \item Device Manager
  \end{enumerate}
\item \textit{[Original]} The executable program of Windows Explorer is:
  \begin{enumerate}[label=\Alph*.]
    \item winfile.exe
    \item explore.exe
    \item explorer.exe
    \item manager.exe
  \end{enumerate}
\item \textit{[Original]} Which part of the File Explorer window shows the
  full path of the folder that is currently open?
  \begin{enumerate}[label=\Alph*.]
    \item Address bar
    \item Navigation pane
    \item Search box
    \item Status bar
  \end{enumerate}
\item \textit{[Original]} In a file path, the parts are separated by:
  \begin{enumerate}[label=\Alph*.]
    \item a forward slash (/)
    \item a backslash (\textbackslash)
    \item a colon (:)
    \item a full stop (.)
  \end{enumerate}
\item \textit{[Original]} A folder that is placed inside another folder is
  called a:
  \begin{enumerate}[label=\Alph*.]
    \item root folder
    \item subfolder
    \item drive
    \item shortcut
  \end{enumerate}
\item \textit{[Original]} Which of the following is a valid file extension
  for a Microsoft Word document?
  \begin{enumerate}[label=\Alph*.]
    \item .xlsx
    \item .pptx
    \item .docx
    \item .txt
  \end{enumerate}
\item \textit{[Original]} The extension of a file is used to indicate:
  \begin{enumerate}[label=\Alph*.]
    \item the size of the file in bytes.
    \item the type of the file and the program that opens it.
    \item the date the file was created.
    \item the folder in which the file is stored.
  \end{enumerate}
\item \textit{[Original]} Which of the following statements is \textbf{NOT}
  correct regarding the Recycle Bin?
  \begin{enumerate}[label=\Alph*.]
    \item Files deleted with the Delete key normally go to the Recycle Bin.
    \item A file in the Recycle Bin can be restored to its original location.
    \item Emptying the Recycle Bin deletes its contents permanently.
    \item A file deleted with Shift\,+\,Delete can be restored from the
      Recycle Bin.
  \end{enumerate}
\item \textit{[Original]} Which keyboard shortcut permanently deletes a file
  without sending it to the Recycle Bin?
  \begin{enumerate}[label=\Alph*.]
    \item Delete
    \item Alt\,+\,Delete
    \item Ctrl\,+\,Delete
    \item Shift\,+\,Delete
  \end{enumerate}
\item \textit{[Original]} Which keyboard shortcut is used to rename a selected
  file or folder?
  \begin{enumerate}[label=\Alph*.]
    \item F1
    \item F2
    \item F3
    \item F5
  \end{enumerate}
\item \textit{[Original]} Match the shortcut to its action:
  \begin{enumerate}[label=\Roman*.]
    \item Ctrl\,+\,C
    \item Ctrl\,+\,X
    \item Ctrl\,+\,V
  \end{enumerate}
  \begin{enumerate}[label=\Alph*.]
    \item I --- Paste, II --- Copy, III --- Cut
    \item I --- Copy, II --- Paste, III --- Cut
    \item I --- Cut, II --- Copy, III --- Paste
    \item I --- Copy, II --- Cut, III --- Paste
  \end{enumerate}
\item \textit{[Original]} Which of the following operations leaves the
  original file in its place and creates a duplicate?
  \begin{enumerate}[label=\Alph*.]
    \item Move
    \item Copy
    \item Cut
    \item Rename
  \end{enumerate}
\item \textit{[Original]} Cut followed by Paste performs which operation?
  \begin{enumerate}[label=\Alph*.]
    \item Copy the file
    \item Move the file
    \item Rename the file
    \item Delete the file
  \end{enumerate}
\item \textit{[Original]} When a file is dragged to a folder on a
  \textbf{different} drive, the file is:
  \begin{enumerate}[label=\Alph*.]
    \item moved, because drag always moves.
    \item copied, because Windows assumes a copy between separate devices.
    \item deleted, because the drives are different.
    \item renamed, because the path changes.
  \end{enumerate}
\item \textit{[Original]} When a file is dragged to a folder on the
  \textbf{same} drive, by default the file is:
  \begin{enumerate}[label=\Alph*.]
    \item copied.
    \item moved.
    \item deleted.
    \item sent to the Recycle Bin.
  \end{enumerate}
\item \textit{[Original]} Which of the following is \textbf{NOT} a file
  attribute in Windows?
  \begin{enumerate}[label=\Alph*.]
    \item Hidden
    \item Read-only
    \item System
    \item Recycle
  \end{enumerate}
\item \textit{[Original]} To display hidden files and folders in Windows 10 or
  11, the correct route is:
  \begin{enumerate}[label=\Alph*.]
    \item View tab $\rightarrow$ Show $\rightarrow$ Hidden items
    \item Home tab $\rightarrow$ Copy $\rightarrow$ Hidden items
    \item Share tab $\rightarrow$ Rename $\rightarrow$ Hidden items
    \item Edit tab $\rightarrow$ Properties $\rightarrow$ Hidden items
  \end{enumerate}
\item \textit{[Original]} Which shortcut opens the Properties of a selected
  file or folder?
  \begin{enumerate}[label=\Alph*.]
    \item Ctrl\,+\,P
    \item Alt\,+\,Enter
    \item Ctrl\,+\,Enter
    \item Alt\,+\,P
  \end{enumerate}
\item \textit{[Original]} Evaluate the statements:
  \begin{enumerate}[label=\Roman*.]
    \item Copy creates a duplicate and keeps the original file.
    \item The Delete key sends a file to the Recycle Bin.
    \item A file deleted with Shift\,+\,Delete can be restored from the
      Recycle Bin.
  \end{enumerate}
  \begin{enumerate}[label=\Alph*.]
    \item I only
    \item I and II only
    \item II and III only
    \item I, II, and III
  \end{enumerate}
\item \textit{[Original]} Which pairing is correctly matched?
  \begin{enumerate}[label=\Alph*.]
    \item Extension --- the format of the file data
    \item Address bar --- shows the full path of the current folder
    \item Recycle Bin --- permanently deletes files at once
    \item Ctrl\,+\,A --- renames a selected file
  \end{enumerate}

\item \textit{[JKSSB FAA-pattern]} In Windows, deleted files are first moved to the:
  \begin{enumerate}[label=\Alph*.]
    \item Recycle Bin.
    \item printer tray.
    \item speaker box.
    \item monitor stand.
  \end{enumerate}
\item \textit{[JKPSI-pattern]} Which view in File Explorer shows the folder hierarchy on the left?
  \begin{enumerate}[label=\Alph*.]
    \item Navigation pane
    \item Printer queue
    \item Volume mixer
    \item Sticky notes board
  \end{enumerate}
\item \textit{[SSC CGL-pattern]} A file extension such as .txt mainly indicates the:
  \begin{enumerate}[label=\Alph*.]
    \item file type/format.
    \item file's physical weight.
    \item printer's ink level.
    \item monitor's brightness.
  \end{enumerate}
\item \textit{[JKSSB-pattern]} \textbf{(MSQ: select ALL correct options.)} Which actions can be done in File Explorer?
  \begin{enumerate}[label=\Alph*.]
    \item Copying files
    \item Renaming folders
    \item Creating new folders
    \item Boiling water
  \end{enumerate}
\item \textit{[SSC CGL-pattern]} \textbf{(MSQ: select ALL correct options.)} Which statements about Windows file handling are correct?
  \begin{enumerate}[label=\Alph*.]
    \item Folders can contain files and subfolders.
    \item Deleted files go to the Recycle Bin first.
    \item Files can be restored from the Recycle Bin.
    \item Emptying the Recycle Bin restores all files.
  \end{enumerate}
\item \textit{[Original]} Which path format identifies a file by its drive and folders in Windows?
  \begin{enumerate}[label=\Alph*.]
\item C:\textbackslash Users\textbackslash Asha\textbackslash report.docx
   \item https://example.com
   \item user@example.com
   \item 192.0.2.10
  \end{enumerate}

\item \textit{[Original]} A file extension usually appears:
  \begin{enumerate}[label=\Alph*.]
\item After the final period in a filename
   \item Before the drive letter
   \item In place of the file contents
   \item Only inside the Recycle Bin
  \end{enumerate}

\item \textit{[Original]} Which Windows action normally moves a deleted local file to the Recycle Bin rather than immediately erasing it?
  \begin{enumerate}[label=\Alph*.]
\item Delete
   \item Rename
   \item Refresh
   \item Search
  \end{enumerate}

\item \textit{[Original]} Which operation creates another copy while leaving the original file in place?
  \begin{enumerate}[label=\Alph*.]
\item Copy
   \item Move
   \item Cut
   \item Delete
  \end{enumerate}

\item \textit{[Original, MSQ drill]} \textbf{(MSQ: select ALL correct options.)} Which statements about files/folders are correct?
  \begin{enumerate}[label=\Alph*.]
\item A folder can contain files and other folders.
   \item A path identifies a location in a file system.
   \item Renaming normally changes a file's name, not its stored contents.
   \item Copying a file necessarily deletes the original.
  \end{enumerate}

\end{enumerate}
\lessonheading{21}{Keyboard Shortcuts}
\begin{enumerate}[label=\textbf{Q\arabic*.}, leftmargin=*]
\item \textit{[Original]} In keyboard shortcuts, what does \textbf{Ctrl}
  stand for?
  \begin{enumerate}[label=\Alph*.]
    \item Central
    \item Control
    \item Counter
    \item Compute
  \end{enumerate}
\item \textit{[Original]} A user wants to copy selected text while keeping the
  original in place. Which shortcut should be used?
  \begin{enumerate}[label=\Alph*.]
    \item Ctrl+X
    \item Ctrl+V
    \item Ctrl+C
    \item Ctrl+A
  \end{enumerate}
\item \textit{[Original]} Which shortcut cuts the selected item and removes it
  from its original position?
  \begin{enumerate}[label=\Alph*.]
    \item Ctrl+C
    \item Ctrl+X
    \item Ctrl+V
    \item Ctrl+S
  \end{enumerate}
\item \textit{[Original]} Which shortcut pastes the item held on the
  clipboard?
  \begin{enumerate}[label=\Alph*.]
    \item Ctrl+P
    \item Ctrl+C
    \item Ctrl+V
    \item Ctrl+X
  \end{enumerate}
\item \textit{[Original]} Which shortcut is used to undo the last action?
  \begin{enumerate}[label=\Alph*.]
    \item Ctrl+Y
    \item Ctrl+Z
    \item Ctrl+U
    \item Ctrl+A
  \end{enumerate}
\item \textit{[Original]} Which shortcut redoes an action that was just
  undone?
  \begin{enumerate}[label=\Alph*.]
    \item Ctrl+Z
    \item Ctrl+R
    \item Ctrl+Y
    \item Ctrl+D
  \end{enumerate}
\item \textit{[Original]} Which shortcut selects all the content in the
  current document or window?
  \begin{enumerate}[label=\Alph*.]
    \item Ctrl+A
    \item Ctrl+S
    \item Ctrl+E
    \item Ctrl+L
  \end{enumerate}
\item \textit{[Original]} A user has finished editing a document and wants to
  store the latest changes. Which shortcut saves the file?
  \begin{enumerate}[label=\Alph*.]
    \item Ctrl+P
    \item Ctrl+S
    \item Ctrl+D
    \item Ctrl+O
  \end{enumerate}
\item \textit{[Original]} Which shortcut opens the print dialog?
  \begin{enumerate}[label=\Alph*.]
    \item Ctrl+P
    \item Ctrl+Q
    \item Ctrl+R
    \item Ctrl+F
  \end{enumerate}
\item \textit{[Original]} Which shortcut opens the \textbf{Find} search box?
  \begin{enumerate}[label=\Alph*.]
    \item Ctrl+H
    \item Ctrl+N
    \item Ctrl+F
    \item Ctrl+G
  \end{enumerate}
\item \textit{[Original]} Which shortcut opens the \textbf{Replace} function?
  \begin{enumerate}[label=\Alph*.]
    \item Ctrl+F
    \item Ctrl+H
    \item Ctrl+R
    \item Ctrl+E
  \end{enumerate}
\item \textit{[Original]} Which shortcut switches between the open windows or
  running applications?
  \begin{enumerate}[label=\Alph*.]
    \item Alt+F4
    \item Alt+Tab
    \item Ctrl+Tab
    \item Win+Tab
  \end{enumerate}
\item \textit{[Original]} Which shortcut closes the active window or
  application?
  \begin{enumerate}[label=\Alph*.]
    \item F4
    \item F1
    \item Alt+F4
    \item Ctrl+F4
  \end{enumerate}
\item \textit{[Original]} Which function key opens \textbf{Help}?
  \begin{enumerate}[label=\Alph*.]
    \item F2
    \item F1
    \item F5
    \item F7
  \end{enumerate}
\item \textit{[Original]} In File Explorer, which function key renames the
  selected file or folder?
  \begin{enumerate}[label=\Alph*.]
    \item F2
    \item F1
    \item F5
    \item F12
  \end{enumerate}
\item \textit{[PYQ WO 2026 Q80, key:provisional]} In MS PowerPoint, which key
  starts the slide show from the beginning?
  \begin{enumerate}[label=\Alph*.]
    \item F1
    \item F2
    \item F5
    \item F7
  \end{enumerate}
\item \textit{[Original]} Which function key starts the spelling and grammar
  check in MS Word?
  \begin{enumerate}[label=\Alph*.]
    \item F7
    \item F5
    \item F2
    \item F1
  \end{enumerate}
\item \textit{[Original]} Which of the following correctly matches the
  Windows-key shortcuts to their actions?
  \begin{enumerate}[label=\Alph*.]
    \item Win+D opens File Explorer; Win+E shows the desktop; Win+L refreshes
      the window
    \item Win+D locks the computer; Win+E shows the desktop; Win+L opens File
      Explorer
    \item Win+D shows the desktop; Win+E opens File Explorer; Win+L locks the
      computer
    \item Win+D prints the screen; Win+E pastes text; Win+L logs off
  \end{enumerate}
\item \textit{[Original]} Which of the following key--action pairings is
  \textbf{NOT} correct?
  \begin{enumerate}[label=\Alph*.]
    \item F1 --- Help
    \item F2 --- Rename
    \item F5 --- Save
    \item F7 --- Spelling and grammar check
  \end{enumerate}
\item \textit{[Original]} Evaluate the statements:
  \begin{enumerate}[label=\Roman*.]
    \item Alt+Tab switches between open windows.
    \item Ctrl+Y is used to redo an action.
    \item F1 is used to rename a selected file.
  \end{enumerate}
  \begin{enumerate}[label=\Alph*.]
    \item I only
    \item I and II only
    \item II and III only
    \item I, II, and III
  \end{enumerate}

\item \textit{[JKSSB FAA-pattern]} Which shortcut pastes copied content?
  \begin{enumerate}[label=\Alph*.]
    \item Ctrl+V
    \item Ctrl+P
    \item Ctrl+N
    \item Ctrl+O
  \end{enumerate}
\item \textit{[JKPSI-pattern]} Which shortcut undoes the last action in most Windows programs?
  \begin{enumerate}[label=\Alph*.]
    \item Ctrl+Z
    \item Ctrl+Y repeated
    \item Ctrl+B
    \item Ctrl+U
  \end{enumerate}
\item \textit{[SSC CGL-pattern]} Which function key usually starts a PowerPoint slide show from the beginning?
  \begin{enumerate}[label=\Alph*.]
    \item F5
    \item F2
    \item F7
    \item F9
  \end{enumerate}
\item \textit{[JKSSB-pattern]} \textbf{(MSQ: select ALL correct options.)} Which shortcuts are used for clipboard operations?
  \begin{enumerate}[label=\Alph*.]
    \item Ctrl+C (copy)
    \item Ctrl+X (cut)
    \item Ctrl+V (paste)
    \item Ctrl+S (save)
  \end{enumerate}
\item \textit{[SSC CGL-pattern]} \textbf{(MSQ: select ALL correct options.)} Which shortcut pairs are correctly matched?
  \begin{enumerate}[label=\Alph*.]
    \item Ctrl+S -- save
    \item Ctrl+P -- print
    \item Ctrl+Z -- undo
    \item Ctrl+V -- increase font size
  \end{enumerate}
\item \textit{[Original]} Which shortcut commonly saves the active document in Windows applications?
  \begin{enumerate}[label=\Alph*.]
\item Ctrl+S
   \item Ctrl+P
   \item Ctrl+Z
   \item Ctrl+F
  \end{enumerate}

\item \textit{[Original]} Which shortcut commonly undoes the most recent action?
  \begin{enumerate}[label=\Alph*.]
\item Ctrl+Z
   \item Ctrl+Y
   \item Ctrl+N
   \item Ctrl+O
  \end{enumerate}

\item \textit{[Original]} In many Windows applications, Alt+Tab is used to:
  \begin{enumerate}[label=\Alph*.]
\item Switch between open applications
   \item Save the current document
   \item Insert a hyperlink
   \item Close the computer
  \end{enumerate}

\item \textit{[Original]} Which shortcut usually selects all text or items in the active document/list?
  \begin{enumerate}[label=\Alph*.]
\item Ctrl+A
   \item Ctrl+X
   \item Ctrl+V
   \item Ctrl+K
  \end{enumerate}

\item \textit{[Original, MSQ drill]} \textbf{(MSQ: select ALL correct options.)} Match the common shortcut to its action.
  \begin{enumerate}[label=\Alph*.]
\item Ctrl+C --- copy
   \item Ctrl+V --- paste
   \item Ctrl+X --- cut
   \item Ctrl+P --- save
  \end{enumerate}

\end{enumerate}
\lessonheading{45}{Open Source}
\begin{enumerate}[label=\textbf{Q\arabic*.}, leftmargin=*]
\item \textit{[Original]} Expand GPL.
  \begin{enumerate}[label=\Alph*.]
    \item General Public Licence
    \item GNU General Public License
    \item Global Program Licence
    \item General Purpose License
  \end{enumerate}
\item \textit{[Original]} Expand FSF.
  \begin{enumerate}[label=\Alph*.]
    \item Free Source Forum
    \item Full Source Foundation
    \item Free Software Foundation
    \item Freeware Standards Federation
  \end{enumerate}
\item \textit{[PYQ WO 2026 Q64, key:provisional]} Which of the following is an
  open-source mobile operating system?
  \begin{enumerate}[label=\Alph*.]
    \item iOS
    \item Android
    \item Windows Phone
    \item BlackBerry OS
  \end{enumerate}
\item \textit{[PYQ WO 2026 Q66, key:provisional]} An open-source licence
  generally allows a user to:
  \begin{enumerate}[label=\Alph*.]
    \item use the software free of cost only.
    \item view, change, and share the source code.
    \item run the software for a limited trial period.
    \item use the software without any licence conditions at all.
  \end{enumerate}
\item \textit{[Original]} Source code is:
  \begin{enumerate}[label=\Alph*.]
    \item the machine-readable form executed by the processor.
    \item the human-readable instructions written by a programmer.
    \item the compiled object code distributed to users.
    \item the licence text that accompanies software.
  \end{enumerate}
\item \textit{[Original]} Proprietary software is best described as software
  whose:
  \begin{enumerate}[label=\Alph*.]
    \item source code is kept secret and only the executable is distributed.
    \item source code is public but the price is zero.
    \item source code has no copyright owner.
    \item source code can be freely forked and modified by anyone.
  \end{enumerate}
\item \textit{[Original]} Freeware is software that is:
  \begin{enumerate}[label=\Alph*.]
    \item free to use, but whose source code is usually not provided.
    \item free to study, change, and redistribute with its source code.
    \item available only as a paid trial version.
    \item free of copyright and any restriction.
  \end{enumerate}
\item \textit{[Original]} The term ``free software'' emphasises:
  \begin{enumerate}[label=\Alph*.]
    \item that the software costs nothing to buy.
    \item that the user has the freedom to run, study, change, and share it.
    \item that the software is available only for a trial period.
    \item that the software has no licence at all.
  \end{enumerate}
\item \textit{[Original]} Which statement correctly distinguishes free software
  from freeware?
  \begin{enumerate}[label=\Alph*.]
    \item Free software is about freedom; freeware is free of cost.
    \item Free software is free of cost; freeware is about freedom.
    \item Both terms mean exactly the same thing.
    \item Freeware always includes the source code; free software never does.
  \end{enumerate}
\item \textit{[Original]} Public domain software is software that:
  \begin{enumerate}[label=\Alph*.]
    \item is copyrighted but licensed for public use.
    \item has no copyright owner and may be used without restriction.
    \item may be used free of cost but only during a trial period.
    \item keeps its source code secret from all users.
  \end{enumerate}
\item \textit{[Original]} Which of the following is \textbf{NOT} an example of
  open-source software?
  \begin{enumerate}[label=\Alph*.]
    \item Linux
    \item Mozilla Firefox
    \item LibreOffice
    \item Microsoft Office
  \end{enumerate}
\item \textit{[Original]} In open-source development, a repository is:
  \begin{enumerate}[label=\Alph*.]
    \item a storage location holding the source code and its change history.
    \item the compiler that converts source code into object code.
    \item the licence that the project is released under.
    \item the paid trial version of the software.
  \end{enumerate}
\item \textit{[Original]} Which of the following is a popular platform that
  hosts open-source code repositories?
  \begin{enumerate}[label=\Alph*.]
    \item GitHub
    \item Adobe Photoshop
    \item Microsoft Office
    \item Oracle Database
  \end{enumerate}
\item \textit{[Original]} In community development, a ``fork'' of a project is:
  \begin{enumerate}[label=\Alph*.]
    \item a bug report filed by a user.
    \item the licence text of the project.
    \item a personal copy of the project for experimenting or improving.
    \item the compiled executable distributed to the public.
  \end{enumerate}
\item \textit{[Original]} Evaluate the statements:
  \begin{enumerate}[label=\Roman*.]
    \item Open-source software makes its source code publicly available.
    \item Freeware always includes the source code.
    \item Public domain software has no copyright owner.
  \end{enumerate}
  \begin{enumerate}[label=\Alph*.]
    \item I and II only
    \item I and III only
    \item II and III only
    \item I, II, and III
  \end{enumerate}
\item \textit{[Original]} Which licence is a copyleft licence?
  \begin{enumerate}[label=\Alph*.]
    \item MIT License
    \item Apache License 2.0
    \item GNU General Public License
    \item BSD License
  \end{enumerate}
\item \textit{[Original]} Which licence is a permissive licence?
  \begin{enumerate}[label=\Alph*.]
    \item GNU General Public License
    \item MIT License
    \item GNU AGPL
    \item None of these
  \end{enumerate}
\item \textit{[Original]} Shareware is software that is:
  \begin{enumerate}[label=\Alph*.]
    \item free of cost forever, with its source code provided.
    \item distributed as a free trial that requires payment to continue.
    \item free to modify and share by anyone.
    \item free of copyright and all restrictions.
  \end{enumerate}
\item \textit{[Original]} Commercial software is software that is:
  \begin{enumerate}[label=\Alph*.]
    \item developed and sold for profit.
    \item available only in the public domain.
    \item always open source.
    \item always freeware.
  \end{enumerate}
\item \textit{[Original]} Which pairing is correctly matched?
  \begin{enumerate}[label=\Alph*.]
    \item Open source --- source code kept secret from users
    \item Freeware --- free of cost but source code usually not provided
    \item Public domain --- still protected by full copyright
    \item Shareware --- free to modify and redistribute the source code
  \end{enumerate}

\item \textit{[JKSSB FAA-pattern]} Open-source software is defined by the freedom to:
  \begin{enumerate}[label=\Alph*.]
    \item view, use, modify, and share its source code.
    \item print unlimited paper copies only.
    \item hear its audio louder.
    \item hide its code from everyone.
  \end{enumerate}
\item \textit{[JKPSI-pattern]} Which of the following is open-source software?
  \begin{enumerate}[label=\Alph*.]
    \item LibreOffice
    \item A sealed printer cartridge
    \item A wall clock
    \item A paper diary
  \end{enumerate}
\item \textit{[SSC CGL-pattern]} Linux is an example of which software class?
  \begin{enumerate}[label=\Alph*.]
    \item Open-source operating system
    \item Printer hardware
    \item Speaker accessory
    \item Paper product
  \end{enumerate}
\item \textit{[JKSSB-pattern]} \textbf{(MSQ: select ALL correct options.)} Which are open-source software?
  \begin{enumerate}[label=\Alph*.]
    \item Linux
    \item LibreOffice
    \item VLC media player
    \item A locked proprietary BIOS with no source
  \end{enumerate}
\item \textit{[SSC CGL-pattern]} \textbf{(MSQ: select ALL correct options.)} Which statements about open source are correct?
  \begin{enumerate}[label=\Alph*.]
    \item Source code is publicly available.
    \item Users may modify and redistribute it under its licence.
    \item It is always free of every licence condition.
    \item Open source forbids all commercial use.
  \end{enumerate}
\item \textit{[Original]} Open-source software generally makes which material available under its license?
  \begin{enumerate}[label=\Alph*.]
\item Source code
   \item The user's password database
   \item Private files of all users
   \item A computer's BIOS password
  \end{enumerate}

\item \textit{[Original]} Freeware and open-source are:
  \begin{enumerate}[label=\Alph*.]
\item Different concepts; free price does not necessarily grant source-code rights
   \item Always identical terms
   \item Both hardware categories
   \item Names for trial versions only
  \end{enumerate}

\item \textit{[Original]} Which license-related statement is most accurate?
  \begin{enumerate}[label=\Alph*.]
\item Users must follow the software's license terms.
   \item Open-source software has no license.
   \item Public-domain software is always proprietary.
   \item Freeware may never be copyrighted.
  \end{enumerate}

\item \textit{[Original]} A proprietary application is usually distributed with:
  \begin{enumerate}[label=\Alph*.]
\item Restrictions defined by its owner/license
   \item Automatic public rights to modify its source
   \item No copyright protection
   \item Mandatory source publication
  \end{enumerate}

\item \textit{[Original, MSQ drill]} \textbf{(MSQ: select ALL correct options.)} Which statements about software distribution are correct?
  \begin{enumerate}[label=\Alph*.]
\item Freeware may be free of charge while remaining closed-source.
   \item Open-source software is governed by a license.
   \item Shareware may be distributed as a trial.
   \item All free-of-charge software is public domain.
  \end{enumerate}

\end{enumerate}
\lessonheading{46}{Firmware, BIOS, UEFI and Drivers}
\begin{enumerate}[label=\textbf{Q\arabic*.}, leftmargin=*]
\item \textit{[Original]} Expand BIOS.
  \begin{enumerate}[label=\Alph*.]
    \item Basic Input/Output System
    \item Binary Input/Output System
    \item Basic Internal Operating System
    \item Basic Input/Output Software
  \end{enumerate}
\item \textit{[Original]} Expand UEFI.
  \begin{enumerate}[label=\Alph*.]
    \item Universal Embedded Firmware Interface
    \item Unified Extensible Firmware Interface
    \item Unified External File Interface
    \item Universal Extensible Firmware Integration
  \end{enumerate}
\item \textit{[Original]} Expand POST.
  \begin{enumerate}[label=\Alph*.]
    \item Power-Off System Test
    \item Primary Operating System Test
    \item Power-On Self-Test
    \item Post-Operation Startup Task
  \end{enumerate}
\item \textit{[PYQ WO 2026 Q68, key:provisional]} Which firmware initialises
  the hardware and loads the operating system when a computer is switched on?
  \begin{enumerate}[label=\Alph*.]
    \item Device driver
    \item BIOS
    \item Application program
    \item Compiler
  \end{enumerate}
\item \textit{[Original]} Firmware is best described as:
  \begin{enumerate}[label=\Alph*.]
    \item temporary data stored in RAM while a program runs.
    \item software stored in non-volatile memory that gives a device its
      basic, low-level instructions.
    \item a program that the user opens from the desktop.
    \item the physical circuits inside a processor.
  \end{enumerate}
\item \textit{[Original]} Firmware is normally stored in:
  \begin{enumerate}[label=\Alph*.]
    \item RAM
    \item a CPU register
    \item ROM or flash memory
    \item the cache
  \end{enumerate}
\item \textit{[Original]} What does the POST do at startup?
  \begin{enumerate}[label=\Alph*.]
    \item Loads the user's application programs
    \item Tests whether the basic hardware is present and working
    \item Formats the hard disk
    \item Installs device drivers
  \end{enumerate}
\item \textit{[Original]} Which item comes FIRST in the power-on sequence?
  \begin{enumerate}[label=\Alph*.]
    \item Application programs
    \item Device drivers
    \item Firmware (BIOS/UEFI)
    \item Operating system
  \end{enumerate}
\item \textit{[Original]} A device driver is:
  \begin{enumerate}[label=\Alph*.]
    \item hardware that plugs into the motherboard.
    \item software that lets the operating system communicate with a
      specific device.
    \item the same as the BIOS.
    \item a type of cache memory.
  \end{enumerate}
\item \textit{[Original]} In which order does the computer normally start
  these layers?
  \begin{enumerate}[label=\Alph*.]
    \item application $\rightarrow$ OS $\rightarrow$ driver $\rightarrow$
      firmware
    \item firmware $\rightarrow$ OS $\rightarrow$ driver $\rightarrow$
      application
    \item OS $\rightarrow$ firmware $\rightarrow$ application $\rightarrow$
      driver
    \item driver $\rightarrow$ firmware $\rightarrow$ OS $\rightarrow$
      application
  \end{enumerate}
\item \textit{[Original]} Which of the following is \textbf{NOT} a function
  of the BIOS?
  \begin{enumerate}[label=\Alph*.]
    \item Running the power-on self-test
    \item Initialising the basic hardware
    \item Finding and loading the operating system
    \item Compiling the user's source code
  \end{enumerate}
\item \textit{[Original]} Which statement correctly contrasts BIOS and UEFI?
  \begin{enumerate}[label=\Alph*.]
    \item BIOS is the modern firmware and UEFI is the legacy firmware.
    \item UEFI is the modern firmware; it boots faster, supports large GPT
      disks, and can provide Secure Boot.
    \item BIOS and UEFI are identical in every feature.
    \item UEFI is a device driver, while BIOS is firmware.
  \end{enumerate}
\item \textit{[Original]} Which option is an advantage of UEFI over the
  legacy BIOS?
  \begin{enumerate}[label=\Alph*.]
    \item It removes the need for any firmware.
    \item It supports very large disks using GPT and offers a graphical
      setup.
    \item It works only on 32-bit computers.
    \item It makes device drivers unnecessary.
  \end{enumerate}
\item \textit{[Original]} A student claims that the program which initialises
  the hardware and starts the operating system is a device driver. Which is
  the correct term?
  \begin{enumerate}[label=\Alph*.]
    \item Application software
    \item BIOS/UEFI firmware
    \item Compiler
    \item Device Manager
  \end{enumerate}
\item \textit{[Original]} A firmware update is commonly called:
  \begin{enumerate}[label=\Alph*.]
    \item flashing
    \item formatting
    \item partitioning
    \item compiling
  \end{enumerate}
\item \textit{[Original]} What is the main risk during a firmware update?
  \begin{enumerate}[label=\Alph*.]
    \item The screen resolution may change.
    \item Interrupting power during flashing can leave the device unusable.
    \item The mouse will stop working.
    \item The font size will change.
  \end{enumerate}
\item \textit{[Original]} Evaluate the statements:
  \begin{enumerate}[label=\Roman*.]
    \item Firmware is stored in non-volatile memory.
    \item Device drivers are loaded by the operating system.
    \item The POST runs after the operating system has started.
  \end{enumerate}
  \begin{enumerate}[label=\Alph*.]
    \item I and II only
    \item I and III only
    \item II and III only
    \item I, II, and III
  \end{enumerate}
\item \textit{[Original]} BIOS settings such as boot order and system
  date/time are stored in:
  \begin{enumerate}[label=\Alph*.]
    \item RAM
    \item the CMOS (battery-backed) chip
    \item the cache
    \item a device driver
  \end{enumerate}
\item \textit{[Original]} Which of the following is software that runs before
  the operating system?
  \begin{enumerate}[label=\Alph*.]
    \item A word processor
    \item Firmware (BIOS/UEFI)
    \item A web browser
    \item A device driver
  \end{enumerate}
\item \textit{[Original]} Which pairing is correctly matched?
  \begin{enumerate}[label=\Alph*.]
    \item BIOS --- a device driver for printers
    \item UEFI --- the legacy firmware with a text-only setup
    \item Firmware --- permanent software stored in ROM/flash
    \item POST --- the application that loads the OS
  \end{enumerate}

\item \textit{[JKSSB FAA-pattern]} Which firmware initialises hardware and loads the OS at startup?
  \begin{enumerate}[label=\Alph*.]
    \item BIOS/UEFI
    \item Compiler
    \item Assembler
    \item Media player
  \end{enumerate}
\item \textit{[JKPSI-pattern]} UEFI is best described as:
  \begin{enumerate}[label=\Alph*.]
    \item a modern replacement for legacy BIOS.
    \item a type of printer paper.
    \item a speaker cable.
    \item a keyboard cover.
  \end{enumerate}
\item \textit{[SSC CGL-pattern]} A device driver is:
  \begin{enumerate}[label=\Alph*.]
    \item software that lets the OS control a device.
    \item a hardware screwdriver.
    \item a printer paper tray.
    \item a monitor stand.
  \end{enumerate}
\item \textit{[JKSSB-pattern]} \textbf{(MSQ: select ALL correct options.)} Which tasks happen during system boot?
  \begin{enumerate}[label=\Alph*.]
    \item Power-on self-test (POST)
    \item Loading the operating system
    \item Initialising hardware devices
    \item Printing a test novel
  \end{enumerate}
\item \textit{[SSC CGL-pattern]} \textbf{(MSQ: select ALL correct options.)} Which statements about firmware and drivers are correct?
  \begin{enumerate}[label=\Alph*.]
    \item BIOS/UEFI starts before the OS loads.
    \item Drivers translate OS commands for devices.
    \item UEFI supports larger disks than legacy BIOS.
    \item A compiler loads the operating system at startup.
  \end{enumerate}
\item \textit{[PYQ-adapted: JKSSB FAA 2024, QCRIAS item 668; SRC-16/SRC-19]} Firmware is commonly stored in:
  \begin{enumerate}[label=\Alph*.]
\item Cache
   \item RAM
   \item External printer memory only
   \item ROM or flash memory
  \end{enumerate}

\item \textit{[PYQ-adapted: JKSSB Website Operator 2026 Q68, provisional key; SRC-04/SRC-07]} Which firmware initializes hardware and begins the boot process on a traditional PC?
  \begin{enumerate}[label=\Alph*.]
\item BIOS
   \item Device driver
   \item Spreadsheet
   \item Web browser
  \end{enumerate}

\item \textit{[Original]} UEFI is best described as:
  \begin{enumerate}[label=\Alph*.]
\item A modern firmware interface used during startup
   \item A type of spreadsheet formula
   \item A printer cable
   \item A file-compression format
  \end{enumerate}

\item \textit{[Original]} A device driver primarily enables:
  \begin{enumerate}[label=\Alph*.]
\item The operating system to communicate with a hardware device
   \item A monitor to store files permanently
   \item A user to change the CPU clock by typing a document
   \item A network to replace the operating system
  \end{enumerate}

\item \textit{[Original, MSQ drill]} \textbf{(MSQ: select ALL correct options.)} Which statements are correct?
  \begin{enumerate}[label=\Alph*.]
\item Firmware is low-level software associated with a device.
   \item ROM/flash can retain firmware without power.
   \item BIOS/UEFI participates in startup.
   \item A device driver is a physical circuit board.
  \end{enumerate}

\end{enumerate}
\lessonheading{52}{OS PYQ Traps and System Utilities}
\begin{enumerate}[label=\textbf{Q\arabic*.}, leftmargin=*]
\item \textit{[Original]} A device driver is best classified as:
  \begin{enumerate}[label=\Alph*.]
    \item hardware
    \item system software
    \item application software
    \item firmware
  \end{enumerate}
\item \textit{[Original]} Which of the following is \textbf{NOT} system
  software?
  \begin{enumerate}[label=\Alph*.]
    \item the Windows operating system
    \item a device driver
    \item MS Word
    \item a utility program
  \end{enumerate}
\item \textit{[Original]} Expand DMA.
  \begin{enumerate}[label=\Alph*.]
    \item Direct Memory Allocation
    \item Direct Memory Access
    \item Dynamic Memory Access
    \item Dual Mode Access
  \end{enumerate}
\item \textit{[Original]} During a DMA transfer, data is transferred directly
  between:
  \begin{enumerate}[label=\Alph*.]
    \item the CPU and the I/O device
    \item the I/O device and main memory
    \item cache and registers
    \item secondary storage and the CPU
  \end{enumerate}
\item \textit{[Original]} The main advantage of DMA is that:
  \begin{enumerate}[label=\Alph*.]
    \item the CPU is not involved in each individual byte's transfer
    \item it increases the clock speed of the CPU
    \item it increases the capacity of RAM
    \item it encrypts data during the transfer
  \end{enumerate}
\item \textit{[Original]} Expand RTOS.
  \begin{enumerate}[label=\Alph*.]
    \item Real-Time Operating System
    \item Real-Time Output System
    \item Rapid Transfer Operating System
    \item Runtime Object System
  \end{enumerate}
\item \textit{[Original]} An RTOS is specifically chosen when:
  \begin{enumerate}[label=\Alph*.]
    \item tasks must complete within strict time deadlines
    \item a document has to be typed and formatted
    \item a website has to be designed
    \item a photograph has to be edited
  \end{enumerate}
\item \textit{[Original]} Utility software such as Disk Cleanup and Disk
  Defragmenter is used mainly to:
  \begin{enumerate}[label=\Alph*.]
    \item create office documents
    \item maintain and optimise the system
    \item browse the internet
    \item compile source code into object code
  \end{enumerate}
\item \textit{[Original]} Which of the following is application software?
  \begin{enumerate}[label=\Alph*.]
    \item Windows
    \item MS Word
    \item Disk Defragmenter
    \item a device driver
  \end{enumerate}
\item \textit{[Original]} Which of the following is an operating system?
  \begin{enumerate}[label=\Alph*.]
    \item MS Word
    \item Windows
    \item Google Chrome
    \item MS Excel
  \end{enumerate}
\item \textit{[Original]} Which of the following is a Linux distribution?
  \begin{enumerate}[label=\Alph*.]
    \item Ubuntu
    \item Windows 11
    \item MS-DOS
    \item macOS
  \end{enumerate}
\item \textit{[Original]} Which of the following is \textbf{NOT} an operating
  system?
  \begin{enumerate}[label=\Alph*.]
    \item Linux
    \item UNIX
    \item MS-DOS
    \item MS Access
  \end{enumerate}
\item \textit{[Original]} Evaluate the statements:
  \begin{enumerate}[label=\Roman*.]
    \item A device driver is a software program.
    \item The operating system is system software.
    \item MS Word is system software.
  \end{enumerate}
  \begin{enumerate}[label=\Alph*.]
    \item I only
    \item I and II only
    \item II and III only
    \item I, II, and III
  \end{enumerate}
\item \textit{[Original]} Which of the following is an \textbf{internal}
  hardware component?
  \begin{enumerate}[label=\Alph*.]
    \item monitor
    \item keyboard
    \item motherboard
    \item printer
  \end{enumerate}
\item \textit{[Original]} Which of the following is an \textbf{external}
  (peripheral) device?
  \begin{enumerate}[label=\Alph*.]
    \item a RAM module
    \item the motherboard
    \item a mouse
    \item the CPU
  \end{enumerate}
\item \textit{[Original]} In the Windows Run dialog box, typing
  \texttt{winword} opens:
  \begin{enumerate}[label=\Alph*.]
    \item Notepad
    \item MS Word
    \item WordPad
    \item File Explorer
  \end{enumerate}
\item \textit{[Original]} In the Windows Run box, which command opens the
  Command Prompt?
  \begin{enumerate}[label=\Alph*.]
    \item \texttt{dos}
    \item \texttt{cmd}
    \item \texttt{com}
    \item \texttt{command}
  \end{enumerate}
\item \textit{[Original]} While the computer is running, the operating system
  is held in:
  \begin{enumerate}[label=\Alph*.]
    \item secondary storage
    \item primary memory (RAM)
    \item registers only
    \item the printer buffer
  \end{enumerate}
\item \textit{[Original]} Which statement about the operating system is
  correct?
  \begin{enumerate}[label=\Alph*.]
    \item It is system software that provides a user interface and manages
      the computer's resources.
    \item It is a hardware chip inside the CPU.
    \item It is application software for typing documents.
    \item It is the same thing as a device driver.
  \end{enumerate}
\item \textit{[Original]} Expand DOS.
  \begin{enumerate}[label=\Alph*.]
    \item Disk Operating System
    \item Data Operating System
    \item Digital Operating System
    \item Direct Operating System
  \end{enumerate}

\item \textit{[JKSSB FAA-pattern]} Which program schedules tasks and manages hardware in a computer?
  \begin{enumerate}[label=\Alph*.]
    \item Operating system
    \item Media player
    \item Photo editor
    \item Web comic
  \end{enumerate}
\item \textit{[JKPSI-pattern]} A system call is a request made to the:
  \begin{enumerate}[label=\Alph*.]
    \item operating-system kernel.
    \item printer paper.
    \item speaker cone.
    \item mouse cable.
  \end{enumerate}
\item \textit{[SSC CGL-pattern]} Which utility frees disk space by removing unneeded files?
  \begin{enumerate}[label=\Alph*.]
    \item Disk Cleanup
    \item Screen Magnifier
    \item Volume Mixer
    \item Paint Brush
  \end{enumerate}
\item \textit{[JKSSB-pattern]} \textbf{(MSQ: select ALL correct options.)} Which are operating-system functions?
  \begin{enumerate}[label=\Alph*.]
    \item Process management
    \item Memory management
    \item File management
    \item Paper manufacturing
  \end{enumerate}
\item \textit{[SSC CGL-pattern]} \textbf{(MSQ: select ALL correct options.)} Which statements about OS utilities are correct?
  \begin{enumerate}[label=\Alph*.]
    \item Disk defragmentation reorganises scattered file parts.
    \item Antivirus scans for malicious code.
    \item Backup utilities keep restorable copies.
    \item Formatting a disk installs a new monitor.
  \end{enumerate}
\item \textit{[PYQ-adapted: JKSSB Sub-Inspector 2022, QCRIAS item 329; SRC-17/SRC-22]} Which software category controls computer operations and coordinates hardware?
  \begin{enumerate}[label=\Alph*.]
\item System software
   \item Application software
   \item User data
   \item Presentation theme
  \end{enumerate}

\item \textit{[Original]} In Linux, which command commonly lists files in a directory?
  \begin{enumerate}[label=\Alph*.]
\item ls
   \item cd
   \item pwd
   \item rm
  \end{enumerate}

\item \textit{[Original]} Which Windows utility is intended to remove selected temporary files and free disk space?
  \begin{enumerate}[label=\Alph*.]
\item Disk Cleanup
   \item Paint
   \item Notepad
   \item Character Map
  \end{enumerate}

\item \textit{[Original]} Which type of operating system is designed to respond within defined timing constraints?
  \begin{enumerate}[label=\Alph*.]
\item Real-time operating system
   \item Batch-only database
   \item Spreadsheet system
   \item Web browser
  \end{enumerate}

\item \textit{[Original, MSQ drill]} \textbf{(MSQ: select ALL correct options.)} Match the command to its common Linux role.
  \begin{enumerate}[label=\Alph*.]
\item ls --- list directory entries
   \item pwd --- show current working directory
   \item cd --- change directory
   \item rm --- rename a directory only
  \end{enumerate}

\end{enumerate}
\lessonheading{60}{Current OS, Processes, Linux and Utilities}
\begin{enumerate}[label=\textbf{Q\arabic*.}, leftmargin=*]
\item \textit{[PYQ-adapted: JKSSB Junior Assistant 2026 Q70; final key;
  paraphrased]} A system is in a deadlock. Which method recovers from the
  deadlock by taking a resource away from one process and giving it to
  another?
  \begin{enumerate}[label=\Alph*.]
    \item Resource preemption
    \item Mutual exclusion
    \item Hold and wait
    \item Circular wait
  \end{enumerate}
\item \textit{[Original]} Which of the following is \textbf{NOT} one of the
  four necessary conditions for deadlock?
  \begin{enumerate}[label=\Alph*.]
    \item Mutual exclusion
    \item Hold and wait
    \item Preemption
    \item Circular wait
  \end{enumerate}
\item \textit{[PYQ-adapted: JKSSB Junior Assistant 2026 Q71; final key;
  paraphrased]} A low-priority task holds a resource that a high-priority
  task needs, and meanwhile a medium-priority task preempts the low-priority
  task and runs. This situation is called:
  \begin{enumerate}[label=\Alph*.]
    \item priority inversion
    \item priority inheritance
    \item deadlock
    \item fragmentation
  \end{enumerate}
\item \textit{[Original]} Which technique resolves priority inversion by
  temporarily raising the priority of the task that holds the needed
  resource?
  \begin{enumerate}[label=\Alph*.]
    \item Time slicing
    \item Priority inheritance
    \item Spooling
    \item Paging
  \end{enumerate}
\item \textit{[Original]} Which of the following does \textbf{NOT} by itself
  prevent priority inversion?
  \begin{enumerate}[label=\Alph*.]
    \item Priority inheritance
    \item Priority ceiling protocol
    \item Round-robin scheduling
    \item Priority boosting
  \end{enumerate}
\item \textit{[Original]} A ready process is repeatedly passed over and waits
  for an indefinite time because higher-priority processes keep arriving.
  This condition is called:
  \begin{enumerate}[label=\Alph*.]
    \item starvation
    \item deadlock
    \item thrashing
    \item fragmentation
  \end{enumerate}
\item \textit{[Original]} Which scheduling technique gradually increases the
  priority of a long-waiting process so that it eventually runs?
  \begin{enumerate}[label=\Alph*.]
    \item Aging
    \item Swapping
    \item Spooling
    \item Caching
  \end{enumerate}
\item \textit{[Original]} In preemptive priority scheduling, the CPU is taken
  away from the running process when:
  \begin{enumerate}[label=\Alph*.]
    \item the running process finishes only.
    \item a process with a higher priority becomes ready.
    \item the running process performs an input/output operation only.
    \item the running process requests more memory only.
  \end{enumerate}
\item \textit{[Original]} Expand RTOS.
  \begin{enumerate}[label=\Alph*.]
    \item Real-Time Operating System
    \item Rapid Task Operating System
    \item Remote Terminal Operating System
    \item Read-Only Operating System
  \end{enumerate}
\item \textit{[Original]} Evaluate the statements:
  \begin{enumerate}[label=\Roman*.]
    \item In priority inheritance, a low-priority task holding a needed
      resource temporarily inherits the priority of the high-priority task
      waiting for it.
    \item Round-robin scheduling by itself prevents priority inversion.
    \item Starvation means a process may wait indefinitely while other
      processes are repeatedly served.
  \end{enumerate}
  \begin{enumerate}[label=\Alph*.]
    \item I and II only
    \item II and III only
    \item I and III only
    \item I, II, and III
  \end{enumerate}
\item \textit{[Original]} The core part of the Linux operating system, which
  manages processes, memory, and devices, is called the:
  \begin{enumerate}[label=\Alph*.]
    \item Linux kernel
    \item Linux shell
    \item GNOME desktop
    \item GRUB bootloader
  \end{enumerate}
\item \textit{[Original]} In Linux, which pairing of a directory with its
  purpose is correct?
  \begin{enumerate}[label=\Alph*.]
    \item \texttt{/root} --- the home directory of the superuser (root user)
    \item \texttt{/home} --- the home directory of the superuser
    \item \texttt{/root} --- the home directories of ordinary users
    \item \texttt{/home} --- the root of the whole file system
  \end{enumerate}
\item \textit{[Original]} Which file system is native to Linux and is commonly
  its default?
  \begin{enumerate}[label=\Alph*.]
    \item NTFS
    \item ext4
    \item FAT32
    \item exFAT
  \end{enumerate}
\item \textit{[Original]} Which file system is the default for modern Windows
  installations?
  \begin{enumerate}[label=\Alph*.]
    \item ext4
    \item NTFS
    \item ext3
    \item HFS+
  \end{enumerate}
\item \textit{[Original]} What is the Windows Store (Microsoft Store)?
  \begin{enumerate}[label=\Alph*.]
    \item A file compression utility
    \item A disk defragmentation tool
    \item An online marketplace used to download and install apps for Windows
    \item A Linux package manager
  \end{enumerate}
\item \textit{[Original]} Which statement correctly contrasts Notepad and
  WordPad?
  \begin{enumerate}[label=\Alph*.]
    \item Notepad supports rich text formatting, while WordPad is plain text
      only.
    \item Notepad is a plain-text editor, while WordPad supports basic
      formatting such as bold and font styles.
    \item Both are spreadsheet programs.
    \item WordPad cannot save a file in any format.
  \end{enumerate}
\item \textit{[Original]} Which built-in Windows utility deletes unnecessary
  files such as temporary files and the Recycle Bin contents to free disk
  space?
  \begin{enumerate}[label=\Alph*.]
    \item Disk Cleanup
    \item Disk Defragmenter
    \item Device Manager
    \item Task Scheduler
  \end{enumerate}
\item \textit{[Original]} Why should a solid-state drive (SSD) generally
  \textbf{NOT} be defragmented?
  \begin{enumerate}[label=\Alph*.]
    \item It has no moving parts, so defragmenting gives no speed benefit and
      only adds unnecessary write wear; TRIM is used instead.
    \item It is faster than RAM and does not need any maintenance.
    \item It is volatile and loses data when defragmented.
    \item It stores data only in RAM.
  \end{enumerate}
\item \textit{[Original]} Which Windows utility is used to scan documents and
  send or receive faxes (the imaging and fax tool)?
  \begin{enumerate}[label=\Alph*.]
    \item Windows Fax and Scan
    \item WordPad
    \item Disk Cleanup
    \item Windows Store
  \end{enumerate}
\item \textit{[Original]} What does the command \texttt{xcopy /s} do?
  \begin{enumerate}[label=\Alph*.]
    \item Copies files and subdirectories, but excludes empty ones.
    \item Deletes files and subdirectories permanently.
    \item Renames a single file only.
    \item Formats a disk.
  \end{enumerate}

\item \textit{[JKSSB FAA-pattern]} In Linux, which command lists directory contents?
  \begin{enumerate}[label=\Alph*.]
    \item ls
    \item print
    \item speak
    \item paint
  \end{enumerate}
\item \textit{[JKPSI-pattern]} A process waiting for input/output is in which state?
  \begin{enumerate}[label=\Alph*.]
    \item Blocked/waiting state
    \item Printed state
    \item Speaker state
    \item Paper state
  \end{enumerate}
\item \textit{[SSC CGL-pattern]} Linux is best described as:
  \begin{enumerate}[label=\Alph*.]
    \item an open-source multi-user operating system.
    \item a printer model.
    \item a paper size.
    \item a speaker brand.
  \end{enumerate}
\item \textit{[JKSSB-pattern]} \textbf{(MSQ: select ALL correct options.)} Which are Linux commands/utilities?
  \begin{enumerate}[label=\Alph*.]
    \item ls (list files)
    \item cd (change directory)
    \item mkdir (make directory)
    \item boil-water
  \end{enumerate}
\item \textit{[SSC CGL-pattern]} \textbf{(MSQ: select ALL correct options.)} Which statements about OS processes are correct?
  \begin{enumerate}[label=\Alph*.]
    \item A process is a program in execution.
    \item The scheduler picks which ready process runs next.
    \item A blocked process waits for an event such as I/O.
    \item A process is a printed paper file.
  \end{enumerate}
\item \textit{[PYQ-adapted: JKSSB Junior Assistant 2026 Q70, final key; SRC-03/SRC-06]} Which deadlock recovery method takes a resource away from a process to resolve a resource cycle?
  \begin{enumerate}[label=\Alph*.]
\item Resource preemption
   \item Increasing display resolution
   \item Formatting the keyboard
   \item Renaming a user account
  \end{enumerate}

\item \textit{[Original]} In Linux, which command prints the current working directory?
  \begin{enumerate}[label=\Alph*.]
\item pwd
   \item ls
   \item cd
   \item rm
  \end{enumerate}

\item \textit{[Original]} Which command is commonly used to list directory contents in Linux?
  \begin{enumerate}[label=\Alph*.]
\item ls
   \item pwd
   \item mkdir
   \item rmdir
  \end{enumerate}

\item \textit{[Original]} Which scheduling policy gives each ready process a fixed time slice in turn?
  \begin{enumerate}[label=\Alph*.]
\item Round-robin
   \item First-come, first-served only
   \item Shortest job first only
   \item Non-preemptive batch sorting
  \end{enumerate}

\item \textit{[Original, MSQ drill]} \textbf{(MSQ: select ALL correct options.)} Which are common operating-system concepts?
  \begin{enumerate}[label=\Alph*.]
\item A process is a program in execution.
   \item Deadlock involves processes waiting on resources.
   \item A file system organizes stored files.
   \item A device driver is a physical keyboard.
  \end{enumerate}

\end{enumerate}
\answersection{5}
\answersubheading{17}{Hardware, Software and Firmware}
\begin{enumerate}[label=\textbf{Q\arabic*.}, leftmargin=*]
\item \textbf{A.} OS stands for Operating System. Open Source (B), Output
  Setting (C), and Object Store (D) are not expansions of the abbreviation.
\item \textbf{B.} BIOS stands for Basic Input Output System. A, C, and D are
  invented expansions.
\item \textbf{C.} The CPU is a physical, tangible part, so it is hardware.
  MS Word (A) and the Windows operating system (D) are software, and a printer
  driver (B) is software, not hardware.
\item \textbf{C.} MS Word is a program, so it is software. The keyboard (A),
  monitor (B), and CPU (D) are all physical hardware.
\item \textbf{B.} A device driver is a program, so it is software even though
  it is named after a device (TRAP-10). It is not hardware (A); it is not
  lower-level embedded code (C); and it is not a person (D).
\item \textbf{B.} The operating system is system software. It is not
  application software (A); it is not physical hardware (C); and it is not
  firmware embedded in a device (D).
\item \textbf{B.} Firmware is software stored in non-volatile memory inside a
  hardware device, for example the BIOS. It is not merely a circuit (A); it is
  not a document application (C); and it is not a person (D).
\item \textbf{A.} Firmware is the software kept in non-volatile memory such as
  ROM or flash and used to start and control the device. RAM (B) is volatile;
  a document file (C) and a user account (D) are not startup code.
\item \textbf{C.} Humanware, or peopleware, is the human component of a
  computer system. The physical parts are hardware (A); stored programs are
  software (B); and startup code in ROM is firmware (D).
\item \textbf{A.} A device driver is software, not a hardware component, so
  statement A is the incorrect one (TRAP-10). B, C, and D are all correct.
\item \textbf{C.} Firmware is correctly paired with the BIOS, which is
  firmware. A wrongly pairs hardware with MS Word (software); B wrongly pairs
  software with the keyboard (hardware); D wrongly pairs humanware with the
  CPU (hardware).
\item \textbf{B.} Bottom to top, the stack is hardware, firmware, device
  driver, operating system, application. A is fully reversed; C and D place
  the layers out of order.
\item \textbf{D.} All three statements are correct: a device driver is a
  software component; the operating system is system software; and firmware
  is stored on the hardware in non-volatile memory. Options A, B, and C each
  omit one or more true statements.
\item \textbf{B.} The device driver is the layer that lets the operating
  system communicate with a specific device. Application software (A) sits
  above the OS, humanware (C) is people, and the BIOS (D) is firmware that
  starts the machine, not the translator for one device.
\item \textbf{B.} MS Word is application software, written to let the user do
  a task. The operating system (A) is system software; the BIOS (C) is
  firmware; and a printer driver (D) is system software (a driver).
\item \textbf{B.} The operating system is system software. MS Excel (A), a
  computer game (C), and a photo editor (D) are all application software.
\item \textbf{B.} Firmware is usually held in ROM or flash memory, which is
  non-volatile. RAM (A) is volatile; CPU registers (C) are volatile working
  storage; and a user folder on the hard disk (D) is ordinary file storage,
  not the home of firmware.
\item \textbf{C.} A device driver is software, not hardware (TRAP-10). The
  motherboard (A), printer (B), and keyboard (D) are all physical hardware.
\item \textbf{D.} The human component is called humanware (or peopleware).
  Hardware (A) is physical, software (B) is programs, and firmware (C) is
  embedded code --- none of these means the people.
\item \textbf{B.} A printer driver lets the operating system communicate with
  the printer by translating commands for that model. It is not a physical
  part (A); it is not the printer's firmware (C); and it does not permanently
  store the document (D) --- files live on secondary storage.

\item \textbf{A.} The keyboard is a physical device. B, C, and D are programs or program content.
\item \textbf{A.} Firmware sits in non-volatile memory controlling hardware. B, C, and D describe non-firmware ideas.
\item \textbf{A.} Drivers are system software linking OS and devices. B, C, and D wrongly call software a device.
\item \textbf{A, B, D.} Monitor, hard disk, and printer are physical. The OS (C) is software.
\item \textbf{A, B, C.} The three definitions are standard. D is false: a driver is software, not hardware.
\item \textbf{A.} Hardware is the physical equipment. Software is instructions, firmware is low-level software stored in non-volatile memory, and data is processed content.

\item \textbf{B.} Firmware must be available at startup and is stored in non-volatile memory, commonly flash/EEPROM. The other listed memories do not serve as its normal persistent store.

\item \textbf{A.} A driver is software that controls or communicates with a device. It is not hardware, storage, or a register.

\item \textbf{A.} A word processor helps users perform tasks and is application software. A driver and firmware are system-level software; a memory controller is hardware.

\item \textbf{A, B, C.} Keyboard, driver, and firmware classifications are correct. A spreadsheet file is data/document content, not a register.

\end{enumerate}
\answersubheading{18}{Software Categories}
\begin{enumerate}[label=\textbf{Q\arabic*.}, leftmargin=*]
\item \textbf{B.} Software is a set of programs or instructions that tell the
  computer what to do. A describes hardware; C wrongly narrows software to the
  operating system; D describes data, not programs.
\item \textbf{B.} The operating system is system software because it manages
  the hardware. MS Word (A) and a payroll package (C) are application
  software; a media player (D) is ordinarily treated as application software.
\item \textbf{B.} OS stands for Operating System. A (Open Source) is a
  licensing term; C and D are not the correct expansions.
\item \textbf{C.} MS Excel is application software because it helps the user
  perform a task. A device driver (A), the operating system (B), and the BIOS
  (D) are all system software.
\item \textbf{A.} An antivirus program performs the maintenance task of
  scanning for and removing viruses, so it is utility software. A web browser
  (B) and MS PowerPoint (C) are application software; a compiler (D) is
  programming software (language translator).
\item \textbf{A.} IDE stands for Integrated Development Environment, which
  combines an editor, a compiler or interpreter, and a debugger. B, C, and D
  are not the correct expansions.
\item \textbf{B.} Compilers, interpreters, and assemblers are language
  translators, so they are programming software. They are not application
  software (A), not utility software (C), and not embedded software (D).
\item \textbf{B.} Software stored inside a device to control that device is
  embedded software, usually held as firmware. Application software (A) is
  for user tasks; freeware (C) is a licensing model, not a device role;
  middleware (D) connects systems.
\item \textbf{A.} Middleware connects different applications and systems and
  sits between the operating system and application software, offering
  services such as messaging and authentication. B describes cabling; C
  describes shareware; D describes hardware.
\item \textbf{C.} MS Word is application software, so it is not system
  software. The operating system (A), a device driver (B), and the BIOS (D)
  are all system software.
\item \textbf{B.} Open-source gives access to the source code, freeware is
  free of cost but may keep its source closed, and shareware is a trial that
  later requires payment. A merges all three; C falsely equates freeware with
  open-source and calls shareware free forever; D equates open-source with
  shareware.
\item \textbf{B.} Open-source software publishes its source code and permits
  viewing, modification, and redistribution under its licence. A describes
  only price, not source availability; C describes shareware; D describes
  public-domain software.
\item \textbf{A.} Freeware is distributed free of cost, but its source code
  may remain closed. B wrongly claims the source is always included; C
  describes shareware; D is incorrect because freeware still has an owner and
  a licence.
\item \textbf{B.} Shareware is offered as a limited trial and requires
  payment for continued or full use. A describes open-source/free software; C
  describes system software; D confuses shareware with public-domain
  software.
\item \textbf{B.} Public-domain software has no copyright and may be used
  freely. A is the opposite; C describes shareware; D describes proprietary
  software that needs activation.
\item \textbf{B.} Statements I and III are correct: an operating system is
  system software, and shareware is usually a trial version. Statement II is
  false because freeware (free of cost) is not the same as open-source
  (source code available). A includes false II; C omits I and includes false
  II; D includes false II.
\item \textbf{B.} GPL stands for GNU General Public License. A is incomplete
  (it drops ``GNU''); C and D are not the correct expansions.
\item \textbf{B.} Utility software performs maintenance and housekeeping
  tasks such as antivirus scanning, disk cleanup, and backup. A describes
  application software; C describes a compiler (programming software); D
  describes middleware.
\item \textbf{A.} Application software is divided into general-purpose
  (ready-made packages) and specific-purpose or custom software. B and C are
  other software groups or licensing terms, not the two kinds of application
  software; D is not a standard classification.
\item \textbf{B.} Open-source software makes its source code available under
  a licence, so this pair is correct. A wrongly says proprietary source is
  openly available; C wrongly says freeware must be paid for; D wrongly says
  shareware never requires payment.

\item \textbf{A.} The OS manages hardware directly. B, C, and D are non-software objects.
\item \textbf{A.} Excel serves the end user, so it is application software. B, C, and D are firmware, hardware, or cabling.
\item \textbf{A.} A word processor is application software. B, C, and D are hardware/firmware.
\item \textbf{A, B, C.} The three packages serve users directly. The kernel (D) is system software.
\item \textbf{A, B, C.} The three category roles are correct. D is false: the OS is system software.
\item \textbf{A.} System software manages hardware and provides services. Application software performs user tasks; data and themes are not software categories for hardware control.

\item \textbf{A.} Disk cleanup is a system utility. A word processor is application software; a presentation file is data, and a cable is hardware.

\item \textbf{A.} Freeware is distributed at no monetary cost, but this alone does not establish open-source freedoms. Processor architecture and storage medium are unrelated.

\item \textbf{A.} Open-source licenses grant rights concerning source code. They do not guarantee free support, hardware schematics, or access to user data.

\item \textbf{A, B, C, D.} System/application/utility classify roles; open source describes source availability and license terms.

\end{enumerate}
\answersubheading{19}{Operating System}
\begin{enumerate}[label=\textbf{Q\arabic*.}, leftmargin=*]
\item \textbf{B.} The operating system provides a user-friendly interface
  layer between the user and the computer hardware. A calls it hardware when it
  is software; C calls it an application when it is system software; D calls it
  storage, which it is not.
\item \textbf{C.} While running, the OS is resident in primary storage (main
  memory). Secondary storage (A) and offline tape (B) are non-volatile stores;
  the display adapter (D) is an output device, not where the OS runs.
\item \textbf{D.} Multitasking means running more than one program at a time.
  Connecting several printers (A), storing several files (B), and using several
  monitors (C) are unrelated to multitasking.
\item \textbf{C.} OS stands for Operating System. A, B, and D are not the
  correct expansions of the abbreviation.
\item \textbf{A.} The kernel is the core part of the OS that interacts directly
  with the hardware. The GUI (B) is just the interface; the bootstrap (C) is a
  startup loader; a compiler (D) is a language tool, not the OS core.
\item \textbf{D.} Designing web pages is an application task, not an OS
  function. Process (A), memory (B), and file (C) management are all core OS
  functions.
\item \textbf{D.} GUI stands for Graphical User Interface. A, B, and C are not
  the correct full forms.
\item \textbf{B.} A CLI (Command Line Interface) takes typed text commands. A
  GUI (A) uses icons and clicks; the Ribbon (C) is a GUI toolbar; the Desktop
  (D) is a GUI workspace.
\item \textbf{A.} Starting the computer and loading the OS into main memory is
  called booting. Formatting (B) prepares storage; compiling (C) translates
  source code; defragmenting (D) reorganises disk data.
\item \textbf{C.} BIOS/UEFI initialises the hardware and then loads the OS
  during startup. A device driver (A) controls one device; the kernel (B) is
  the OS core loaded by the boot process; a compiler (D) translates programs.
\item \textbf{A.} A program in execution is a process. A folder (B), a shortcut
  (C), and a file (D) are stored items, not running programs.
\item \textbf{B.} Memory management decides how much main memory each program
  gets. Process management (A) schedules the CPU; file management (C) handles
  files; device management (D) controls input/output devices.
\item \textbf{C.} File management creates, reads, writes, and deletes files and
  folders. Scheduling (A), memory allocation (B), and print handling (D) belong
  to process, memory, and device management respectively.
\item \textbf{C.} A device driver lets the OS communicate with a specific
  hardware device. A compiler (A) translates code; a browser (B) is an
  internet application; a spreadsheet (D) is an application program.
\item \textbf{A.} The kernel is the core of the operating system, so this
  pairing is correct. A device driver (B) is not a text editor; the bootstrap
  (C) is not storage; a GUI (D) is not a command-line-only interface.
\item \textbf{B.} Time-sharing lets many users use a single computer at the
  same time. A computer cannot work without an OS (A), files need a file
  system (C), and hardware needs power (D).
\item \textbf{C.} Android is the most common example of an open-source
  operating system and is built on the Linux kernel. Windows 11 (A) and macOS
  (B) are proprietary; MS Word (D) is an application, not an OS.
\item \textbf{A.} Statements I and II are correct: the OS provides a
  user-friendly interface, and the kernel is its core. III is false because the
  OS is system software, not an application. B, C, and D all include the false
  statement III.
\item \textbf{C.} Logging in with a username and password is authentication, a
  part of security management. Compression (A), defragmentation (B), and
  spooling (D) are unrelated to login.
\item \textbf{B.} Updates fix security holes and add improvements. They cannot
  enlarge the physical disk (A), turn a GUI into hardware (C), or remove the
  need for a kernel (D).

\item \textbf{A.} Windows manages hardware and programs. B, C, and D are tools, consumables, or parts.
\item \textbf{A.} The kernel is the OS core. B, C, and D are unrelated objects.
\item \textbf{A.} Linux is open source. B, C, and D are not software at all.
\item \textbf{A, B, D.} Windows, Linux, and Android are operating systems. MS Word (C) is application software.
\item \textbf{A, B, C.} Resource management, file/device handling, and boot loading are OS roles. D is false: a word processor is an application.
\item \textbf{A.} The operating system manages resources and provides services to applications/users. The other options are an application, content, or hardware. The cited key is provisional.

\item \textbf{A.} Resource preemption can reclaim resources to break a deadlock. The other actions do not resolve process-resource deadlock.

\item \textbf{A.} File-system management organizes stored files/directories. The other options belong to presentation, text recognition, and document automation.

\item \textbf{A.} Multitasking schedules CPU time among tasks. It does not require one physical CPU per program or eliminate storage/memory.

\item \textbf{A, B, C.} Scheduling processes, managing memory, and managing files are OS responsibilities. Video editing is an application task.

\end{enumerate}
\answersubheading{20}{Windows Explorer}
\begin{enumerate}[label=\textbf{Q\arabic*.}, leftmargin=*]
\item \textbf{B.} Windows Explorer is presented as File Explorer in Windows 10
  and Windows 11. Control Panel (A), Task Manager (C), and Device Manager (D)
  are different Windows tools.
\item \textbf{C.} The program file is explorer.exe. The other spellings
  (winfile.exe, explore.exe, manager.exe) are not the Explorer executable.
\item \textbf{A.} The address bar shows the full path of the open folder. The
  navigation pane (B) shows the tree of drives and folders; the search box (C)
  finds files; the status bar (D) shows item counts.
\item \textbf{B.} Path components are separated by a backslash. A forward
  slash (A) is used on the web, a colon (C) follows the drive letter, and a
  full stop (D) separates a file name from its extension.
\item \textbf{B.} A folder inside another folder is a subfolder. The root
  folder (A) is the top-most folder of a drive; a drive (C) is a storage
  location; a shortcut (D) is a link to an item.
\item \textbf{C.} A Word document has the extension .docx. .xlsx (A) is Excel,
  .pptx (B) is PowerPoint, and .txt (D) is plain text.
\item \textbf{B.} The extension indicates the file's type and the program that
  opens it. Size (A), creation date (C), and folder location (D) are shown in
  Properties but are not what the extension indicates.
\item \textbf{D.} Shift\,+\,Delete deletes permanently and bypasses the
  Recycle Bin, so the file cannot be restored from it. A, B, and C are all
  correct statements about the Recycle Bin.
\item \textbf{D.} Shift\,+\,Delete deletes permanently. The plain Delete key
  (A) sends a file to the Recycle Bin; Alt\,+\,Delete (B) and Ctrl\,+\,Delete
  (C) are not the permanent-delete shortcut.
\item \textbf{B.} F2 renames the selected item. F1 (A) is help, F3 (C) opens
  search, and F5 (D) refreshes the window.
\item \textbf{D.} Ctrl\,+\,C is Copy, Ctrl\,+\,X is Cut, and Ctrl\,+\,V is
  Paste. A, B, and C each mismatch at least one shortcut to its action.
\item \textbf{B.} Copy duplicates an item and keeps the original. Move (A) and
  Cut (C) remove the original from its place; Rename (D) only changes the name.
\item \textbf{B.} Cut followed by Paste is the standard move operation. Copy
  (A) leaves the original; Rename (C) changes only the name; Delete (D) sends
  the file to the Recycle Bin.
\item \textbf{B.} Dragging to a different drive copies, because Windows
  assumes a copy between separate devices. This is not always a move (A); the
  file is not deleted (C); and it is not renamed (D).
\item \textbf{B.} Dragging within the same drive moves the file by default. It
  is not copied (A); it is not deleted (C); and it is not sent to the Recycle
  Bin (D).
\item \textbf{D.} ``Recycle'' is not a file attribute; the Recycle Bin is a
  folder. Hidden (A), Read-only (B), and System (C) are genuine file
  attributes.
\item \textbf{A.} Hidden items are shown via View $\rightarrow$ Show
  $\rightarrow$ Hidden items. The Home, Share, and Edit routes in B, C, and D
  are not the correct path.
\item \textbf{B.} Alt\,+\,Enter opens Properties for the selected item.
  Ctrl\,+\,P (A) prints, Ctrl\,+\,Enter (C) is not the Properties shortcut,
  and Alt\,+\,P (D) is unrelated.
\item \textbf{B.} Statements I and II are correct: Copy keeps the original,
  and Delete sends a file to the Recycle Bin. III is false because
  Shift\,+\,Delete bypasses the Recycle Bin and cannot be restored from it.
  A omits II; C includes the false III and omits I; D includes the false III.
\item \textbf{B.} The address bar shows the full path of the current folder,
  so this pairing is correct. A wrongly equates an extension with the format
  (the extension is the label, the format is the encoding); C wrongly says the
  Recycle Bin deletes at once; D wrongly assigns Rename to Ctrl\,+\,A (which
  selects all).

\item \textbf{A.} Deletions rest in the Recycle Bin until emptied. B, C, and D are not file stores.
\item \textbf{A.} The navigation pane shows the folder tree. B, C, and D are unrelated panels.
\item \textbf{A.} Extensions identify the file's type. B, C, and D describe physical or device states.
\item \textbf{A, B, C.} Copy, rename, and new-folder actions are Explorer basics. D is not a computer task.
\item \textbf{A, B, C.} Folder nesting, Recycle Bin staging, and restore are correct. D is false: emptying removes files instead of restoring them.
\item \textbf{A.} A drive-letter path with folders identifies a Windows file location. The other options are a URL, email address, and IP address.

\item \textbf{A.} An extension is the filename suffix after the final period, such as `.docx`. It does not replace contents or require the Recycle Bin.

\item \textbf{A.} Delete commonly moves a file to the Recycle Bin. Rename changes its name, Refresh updates a view, and Search locates items.

\item \textbf{A.} Copy creates a duplicate. Move relocates the original; Cut prepares it for moving; Delete removes it from its current location.

\item \textbf{A, B, C.} Folders can contain items, paths locate them, and renaming changes the name. Copying leaves the original in place.

\end{enumerate}
\answersubheading{21}{Keyboard Shortcuts}
\begin{enumerate}[label=\textbf{Q\arabic*.}, leftmargin=*]
\item \textbf{B.} \textbf{Ctrl} is short for \textbf{Control}. A, C, and D
  (\textit{Central}, \textit{Counter}, \textit{Compute}) are not the expansion
  of the abbreviation.
\item \textbf{C.} \textit{Ctrl+C} copies the selection and leaves the original
  in place. \textit{Ctrl+X} (A) cuts and removes the original; \textit{Ctrl+V}
  (B) pastes; \textit{Ctrl+A} (D) selects all.
\item \textbf{B.} \textit{Ctrl+X} cuts the selection and removes it from its
  original position. \textit{Ctrl+C} (A) copies and keeps the original;
  \textit{Ctrl+V} (C) pastes; \textit{Ctrl+S} (D) saves.
\item \textbf{C.} \textit{Ctrl+V} pastes the clipboard contents. \textit{Ctrl+P}
  (A) prints; \textit{Ctrl+C} (B) copies; \textit{Ctrl+X} (D) cuts.
\item \textbf{B.} \textit{Ctrl+Z} undoes the last action. \textit{Ctrl+Y} (A)
  redoes instead; \textit{Ctrl+U} (C) is underline in many editors; \textit{Ctrl+A}
  (D) selects all.
\item \textbf{C.} \textit{Ctrl+Y} redoes the action that was just undone.
  \textit{Ctrl+Z} (A) undoes; \textit{Ctrl+R} (B) refreshes or right-aligns
  depending on the application; \textit{Ctrl+D} (D) is a different command.
\item \textbf{A.} \textit{Ctrl+A} selects all the content in the current
  document or window. \textit{Ctrl+S} (B) saves; \textit{Ctrl+E} (C) centres
  text in many editors; \textit{Ctrl+L} (D) is a different command.
\item \textbf{B.} \textit{Ctrl+S} saves the file. \textit{Ctrl+P} (A) prints;
  \textit{Ctrl+D} (C) is a different command; \textit{Ctrl+O} (D) opens a file.
\item \textbf{A.} \textit{Ctrl+P} opens the print dialog. \textit{Ctrl+Q} (B)
  quits in some applications; \textit{Ctrl+R} (C) refreshes or right-aligns;
  \textit{Ctrl+F} (D) opens Find.
\item \textbf{C.} \textit{Ctrl+F} opens the Find search box. \textit{Ctrl+H} (A)
  opens Replace; \textit{Ctrl+N} (B) creates a new file; \textit{Ctrl+G} (D)
  is a different command.
\item \textbf{B.} \textit{Ctrl+H} opens the Replace function. \textit{Ctrl+F} (A)
  only finds; \textit{Ctrl+R} (C) refreshes or right-aligns; \textit{Ctrl+E} (D)
  is a different command.
\item \textbf{B.} \textit{Alt+Tab} switches between open windows or running
  applications. \textit{Alt+F4} (A) closes the active window; \textit{Ctrl+Tab}
  (C) switches tabs in some applications; \textit{Win+Tab} (D) opens the task
  view, not the classic window switcher.
\item \textbf{C.} \textit{Alt+F4} closes the active window or application. F4
  alone (A) and F1 (B) do other things; \textit{Ctrl+F4} (D) closes a document
  tab or child window in some applications, not the window itself.
\item \textbf{B.} \textit{F1} opens Help. \textit{F2} (A) renames; \textit{F5}
  (C) refreshes or starts a slide show; \textit{F7} (D) runs the spell check.
\item \textbf{A.} \textit{F2} renames the selected file or folder. \textit{F1}
  (B) opens Help; \textit{F5} (C) refreshes; \textit{F12} (D) is
  \textit{Save As} in MS Word, not rename.
\item \textbf{C.} In PowerPoint, \textit{F5} starts the slide show from the
  beginning. \textit{F1} (A) is Help, \textit{F2} (B) is Rename, and \textit{F7}
  (D) is the spell check.
\item \textbf{A.} \textit{F7} starts the spelling and grammar check in MS Word.
  \textit{F5} (B) refreshes or starts a slide show; \textit{F2} (C) renames;
  \textit{F1} (D) opens Help.
\item \textbf{C.} \textit{Win+D} shows the desktop, \textit{Win+E} opens File
  Explorer, and \textit{Win+L} locks the computer. A, B, and D all mismatch the
  keys with the wrong actions.
\item \textbf{C.} \textit{F5} does not save; it refreshes the window or, in
  PowerPoint, starts a slide show. Saving is \textit{Ctrl+S}. The other three
  pairings --- F1 Help (A), F2 Rename (B), and F7 Spelling and grammar check
  (D) --- are all correct.
\item \textbf{B.} Statements I and II are correct: \textit{Alt+Tab} switches
  between open windows, and \textit{Ctrl+Y} redoes. III is false, because
  \textit{F1} opens Help --- it is \textit{F2} that renames. A omits II; C
  includes the false III and omits I; D includes the false III.

\item \textbf{A.} Ctrl+V pastes. B prints, C opens new, and D opens files.
\item \textbf{A.} Ctrl+Z undoes. B redoes, C bolds, and D underlines.
\item \textbf{A.} F5 launches the show from the start. B, C, and D serve rename, spell-check, or other roles.
\item \textbf{A, B, C.} Copy, cut, and paste are clipboard operations. Ctrl+S (D) saves files instead.
\item \textbf{A, B, C.} Save, print, and undo pairs are correct. D is false: Ctrl+V pastes.
\item \textbf{A.} Ctrl+S saves. Ctrl+P prints, Ctrl+Z undoes, and Ctrl+F searches.

\item \textbf{A.} Ctrl+Z is Undo. Ctrl+Y commonly redoes, Ctrl+N creates a new item, and Ctrl+O opens a file.

\item \textbf{A.} Alt+Tab switches among open windows. The other actions use different commands or controls.

\item \textbf{A.} Ctrl+A selects all. Ctrl+X cuts, Ctrl+V pastes, and Ctrl+K commonly inserts a hyperlink in Office applications.

\item \textbf{A, B, C.} Ctrl+C copies, Ctrl+V pastes, and Ctrl+X cuts. Ctrl+P opens Print; Ctrl+S saves.

\end{enumerate}
\answersubheading{45}{Open Source}
\begin{enumerate}[label=\textbf{Q\arabic*.}, leftmargin=*]
\item \textbf{B.} GPL stands for GNU General Public License, the well-known
  copyleft licence. A, C, and D are not the correct expansion.
\item \textbf{C.} FSF stands for Free Software Foundation, founded by Richard
  Stallman. A, B, and D are not the correct full forms.
\item \textbf{B.} Android is the open-source mobile operating system, built on
  the Linux kernel. iOS (A) is proprietary; Windows Phone (C) and BlackBerry OS
  (D) are not the open-source answer.
\item \textbf{B.} An open-source licence lets a user view, change, and share
  the source code. A describes freeware; C describes shareware; D wrongly
  claims there are no licence conditions.
\item \textbf{B.} Source code is the human-readable set of instructions written
  by a programmer. A and C describe the machine-readable executable or object
  code; D describes the licence text, not the code.
\item \textbf{A.} Proprietary software keeps its source code secret and
  distributes only the executable. B and C describe open source and public
  domain; D describes an open-source permission.
\item \textbf{A.} Freeware is free of cost, but its source code is usually not
  provided. B describes free software; C describes shareware; D describes
  public domain.
\item \textbf{B.} ``Free software'' means freedom --- to run, study, change,
  and share the software, not simply a zero price. A is the freeware sense; C
  describes shareware; D is false.
\item \textbf{A.} Free software is about freedom; freeware is only free of
  cost. B reverses the two; C wrongly equates them; D is false in both parts.
\item \textbf{B.} Public domain software has no copyright owner and may be used
  without restriction. A describes open source; C describes shareware; D
  describes proprietary software.
\item \textbf{D.} Microsoft Office is proprietary. Linux (A), Mozilla Firefox
  (B), and LibreOffice (C) are all open-source projects.
\item \textbf{A.} A repository stores the source code and its change history,
  typically under version control. B describes a compiler; C a licence; D
  shareware.
\item \textbf{A.} GitHub is a popular repository-hosting platform (along with
  GitLab, Bitbucket, and SourceForge). B, C, and D are proprietary software
  products, not code-hosting platforms.
\item \textbf{C.} A fork is a contributor's personal copy of a project used
  for experimenting or improving. A is an issue; B is the licence; D is the
  executable.
\item \textbf{B.} Statements I and III are correct. II is false because
  freeware usually does \emph{not} include the source code. A includes false
  II; C omits true I; D includes false II.
\item \textbf{C.} The GNU General Public License is a copyleft licence,
  requiring derivatives to stay under the same licence. MIT (A), Apache 2.0
  (B), and BSD (D) are permissive.
\item \textbf{B.} The MIT License is permissive, allowing reuse with few
  conditions. The GNU GPL (A) and GNU AGPL (C) are copyleft; D is therefore
  wrong.
\item \textbf{B.} Shareware is a free trial that requires payment to continue.
  A describes free software; C describes open source; D describes public
  domain.
\item \textbf{A.} Commercial software is developed and sold for profit. It is
  often proprietary but not always open source (B and C false); it is not
  necessarily freeware (D false).
\item \textbf{B.} Freeware is free of cost but its source code is usually not
  provided, so this pairing is correct. A is wrong (open source shares the
  source); C is wrong (public domain has no copyright); D is wrong (shareware
  hides its source).

\item \textbf{A.} Open source grants code freedoms. B and C limit it to physical acts; D is the opposite (closed source).
\item \textbf{A.} LibreOffice is open source. B, C, and D are hardware or paper objects.
\item \textbf{A.} Linux is the flagship open-source OS. B, C, and D are hardware or paper.
\item \textbf{A, B, C.} Linux, LibreOffice, and VLC are open source. A locked BIOS (D) is proprietary.
\item \textbf{A, B.} Public code and licence-based sharing are correct. C is false (licences impose conditions); D is false (commercial use is often allowed).
\item \textbf{A.} Open-source licenses provide rights to inspect/use/modify source under their terms. They do not expose private user data or credentials.

\item \textbf{A.} Freeware concerns price; open source concerns source access and license freedoms. They may overlap but are not synonymous.

\item \textbf{A.} Open-source and other software are distributed under licenses/terms. The other claims are false generalizations.

\item \textbf{A.} Proprietary licensing typically restricts use/modification according to the owner's terms. Source publication and modification rights are not automatic.

\item \textbf{A, B, C.} Freeware, open-source, and shareware describe different distribution/licensing properties. No-cost distribution does not automatically place software in the public domain.

\end{enumerate}
\answersubheading{46}{Firmware, BIOS, UEFI and Drivers}
\begin{enumerate}[label=\textbf{Q\arabic*.}, leftmargin=*]
\item \textbf{A.} BIOS stands for Basic Input/Output System. B, C, and D are
  not the correct expansion of the abbreviation.
\item \textbf{B.} UEFI stands for Unified Extensible Firmware Interface. A, C,
  and D are not the correct full forms.
\item \textbf{C.} POST stands for Power-On Self-Test. A, B, and D are not the
  correct expansions.
\item \textbf{B.} BIOS is the firmware that initialises the hardware and
  loads the operating system. A device driver (A) is loaded by the OS; an
  application (C) and a compiler (D) are higher-level software, not startup
  firmware.
\item \textbf{B.} Firmware is software stored in non-volatile memory that
  gives a device its basic low-level instructions. A describes RAM data; C
  describes ordinary application software; D describes hardware circuitry.
\item \textbf{C.} Firmware is normally stored in ROM or flash memory, which
  are non-volatile. RAM (A), a register (B), and cache (D) are all volatile
  and empty at power-on.
\item \textbf{B.} The POST tests whether the basic hardware is present and
  working. Loading applications (A), formatting a disk (C), and installing
  drivers (D) are not POST functions.
\item \textbf{C.} Firmware (BIOS/UEFI) runs first at power-on. The operating
  system (D), device drivers (B), and applications (A) each start later in
  the chain.
\item \textbf{B.} A device driver is software that lets the OS communicate
  with a specific device. A describes hardware; C is wrong because the BIOS
  is firmware, not a driver; D is wrong because a driver is not cache.
\item \textbf{B.} The correct order is firmware, then OS, then driver, then
  application. A, C, and D all place a later layer before an earlier one.
\item \textbf{D.} Compiling source code is done by a compiler after the OS
  and application are running, so it is not a BIOS function. Running POST
  (A), initialising hardware (B), and loading the OS (C) are all BIOS jobs.
\item \textbf{B.} UEFI is the modern firmware: it boots faster, supports
  large GPT disks, and can provide Secure Boot. A reverses the ages; C
  wrongly calls them identical; D wrongly calls UEFI a device driver.
\item \textbf{B.} UEFI supports very large disks using GPT and offers a
  graphical setup. It does not remove the need for firmware (A), is not
  limited to 32-bit systems (C), and does not remove device drivers (D).
\item \textbf{B.} The firmware that initialises hardware and starts the OS
  is the BIOS/UEFI. Application software (A), a compiler (C), and Device
  Manager (D) are not the startup firmware.
\item \textbf{A.} A firmware update is commonly called flashing. Formatting
  (B) prepares storage; partitioning (C) divides a disk; compiling (D)
  translates source code.
\item \textbf{B.} Interrupting power during flashing can leave the device
  unusable. Screen resolution (A), the mouse (C), and font size (D) are not
  the risk of a firmware update.
\item \textbf{A.} Statements I and II are correct: firmware is non-volatile,
  and drivers are loaded by the OS. III is false because POST runs before
  the OS, not after. B and C include false III; D includes false III.
\item \textbf{B.} BIOS settings are stored in the CMOS battery-backed chip.
  RAM (A) is volatile; the cache (C) holds recent data; a device driver (D)
  controls a device, not BIOS settings.
\item \textbf{B.} Firmware (BIOS/UEFI) runs before the operating system. A
  word processor (A), a web browser (C), and a device driver (D) all run
  after the OS has started.
\item \textbf{C.} Firmware is permanent software stored in ROM/flash, so this
  pairing is correct. A wrongly calls BIOS a printer driver; B calls UEFI
  the legacy text-only firmware; D calls POST the application that loads the
  OS.

\item \textbf{A.} BIOS/UEFI boots the machine. Compilers (B) and assemblers (C) translate code; D plays media.
\item \textbf{A.} UEFI succeeds legacy BIOS. B, C, and D are unrelated physical items.
\item \textbf{A.} Drivers are control software. B, C, and D are hardware objects.
\item \textbf{A, B, C.} POST, OS loading, and device init are boot steps. Printing (D) is not.
\item \textbf{A, B, C.} The boot order, driver role, and UEFI advantage are correct. D is false: firmware, not compilers, boots the OS.
\item \textbf{D.} Firmware must persist without power and is commonly stored in ROM/flash. Cache and RAM are volatile; printer memory is not the general answer.

\item \textbf{A.} BIOS firmware initializes hardware and starts booting. Drivers are loaded for device operation; spreadsheet/browser are applications. The cited key is provisional.

\item \textbf{A.} UEFI is firmware interface technology used to initialize the system and start the boot process. The other choices are unrelated.

\item \textbf{A.} Drivers let the OS control/communicate with hardware. The other options confuse storage, applications, and system software.

\item \textbf{A, B, C.} Firmware is software stored persistently; BIOS/UEFI participates in boot. A driver is software, not a circuit board.

\end{enumerate}
\answersubheading{52}{OS PYQ Traps and System Utilities}
\begin{enumerate}[label=\textbf{Q\arabic*.}, leftmargin=*]
\item \textbf{B.} A device driver is software --- specifically system
  software, because it extends the OS's control of a device. A is the
  classic trap (TRAP-10); a driver controls hardware but is not itself
  hardware. C is wrong because the driver is not a user application, and D
  confuses a driver with firmware.
\item \textbf{C.} MS Word is application software, not system software.
  Windows (A) is an operating system, a device driver (B) is system
  software, and a utility program (D) is a maintenance category of system
  software --- all are system software.
\item \textbf{B.} DMA stands for Direct Memory Access. A, C, and D are not
  the correct expansions.
\item \textbf{B.} DMA moves data directly between the I/O device and main
  memory. A is the near-neighbour trap: the CPU starts and ends the transfer
  but does not sit in the data path. C describes internal CPU storage, and D
  is not the DMA path.
\item \textbf{A.} DMA's benefit is that the CPU does not handle each byte of
  the transfer, so it is free for other work. B is wrong --- DMA does not
  change CPU clock speed; C is wrong --- DMA does not change RAM size; D is
  wrong --- DMA does not encrypt data.
\item \textbf{A.} RTOS stands for Real-Time Operating System. B, C, and D are
  invented expansions.
\item \textbf{A.} An RTOS is chosen when tasks must meet strict deadlines.
  Typing a document (B), designing a website (C), and editing a photograph
  (D) are ordinary application tasks without hard timing deadlines.
\item \textbf{B.} Utility software maintains and optimises the system. A
  describes application software, C describes a browser (application), and D
  describes a translator such as a compiler.
\item \textbf{B.} MS Word is application software. Windows (A) is an
  operating system, Disk Defragmenter (C) is a utility program, and a device
  driver (D) is system software.
\item \textbf{B.} Windows is an operating system. MS Word (A), Chrome (C),
  and MS Excel (D) are application programs, not operating systems.
\item \textbf{A.} Ubuntu is a Linux distribution. Windows 11 (B), MS-DOS
  (C), and macOS (D) are operating systems but not Linux distributions.
\item \textbf{D.} MS Access is a database application, not an operating
  system. Linux (A), UNIX (B), and MS-DOS (C) are all operating systems. The
  trap is that all four options are familiar computing names.
\item \textbf{B.} Statements I and II are correct: a device driver is a
  software program, and the OS is system software. III is false because MS
  Word is application software, so A omits II, C includes the false III and
  omits I, and D includes the false III.
\item \textbf{C.} The motherboard is inside the system unit. A monitor (A),
  a keyboard (B), and a printer (D) are external peripherals.
\item \textbf{C.} A mouse is an external (peripheral) device. A RAM module
  (A), the motherboard (B), and the CPU (D) are all internal components.
\item \textbf{B.} Typing \texttt{winword} in the Run box opens MS Word.
  \texttt{notepad} (A), \texttt{wordpad} (C), and \texttt{explorer} (D) are
  opened by different commands; WordPad is the nearest neighbour and a common
  wrong choice.
\item \textbf{B.} \texttt{cmd} opens the Command Prompt. \texttt{dos} (A) and
  \texttt{com} (C) are not the command, and \texttt{command} (D) is the
  legacy MS-DOS shell, not the modern Command Prompt.
\item \textbf{B.} While running, the OS is held in primary memory (RAM)
  (PYQ-22). Secondary storage (A) holds it permanently but does not execute
  it; registers (C) are too small; (D) is unrelated.
\item \textbf{A.} The OS is system software that gives a user interface and
  manages resources (PYQ-21). B is false --- it is not hardware; C is false
  --- it is not an application; D is false --- a device driver is a small
  part of the OS's device control, not the whole OS.
\item \textbf{A.} DOS stands for Disk Operating System. B, C, and D are not
  the correct full forms.

\item \textbf{A.} The OS manages tasks and hardware. B, C, and D are applications or content.
\item \textbf{A.} System calls reach the kernel. B, C, and D are hardware, not call targets.
\item \textbf{A.} Disk Cleanup deletes junk files. B, C, and D serve display, audio, or drawing roles.
\item \textbf{A, B, C.} Process, memory, and file management are OS jobs. Manufacturing (D) is not.
\item \textbf{A, B, C.} The defragment, antivirus, and backup roles are correct. D is false: formatting prepares disks; it installs no monitor.
\item \textbf{A.} System software, including the OS, manages resources and hardware. Application software performs user tasks; data and themes are not that software category.

\item \textbf{A.} `ls` lists directory contents. `cd` changes directory, `pwd` prints the current path, and `rm` removes files.

\item \textbf{A.} Disk Cleanup removes selected temporary/unneeded files. The other utilities edit/draw text or inspect characters.

\item \textbf{A.} An RTOS is designed for bounded timing behavior. The other options are not operating-system classes.

\item \textbf{A, B, C.} The first three are common command roles. `rm` removes files/directories; renaming is commonly done with `mv`.

\end{enumerate}
\answersubheading{60}{Current OS, Processes, Linux and Utilities}
\begin{enumerate}[label=\textbf{Q\arabic*.}, leftmargin=*]
\item \textbf{A.} Resource preemption recovers from deadlock by forcibly
  taking a resource from one process and giving it to another, breaking the
  cycle. B, C, and D are necessary \emph{conditions} for deadlock, not
  recovery methods.
\item \textbf{C.} The four necessary conditions are mutual exclusion, hold and
  wait, \emph{no preemption}, and circular wait. ``Preemption'' is used to
  \emph{recover} from deadlock, so it is not a condition. A, B, and D are
  genuine conditions.
\item \textbf{A.} A high-priority task waiting behind a low-priority holder
  that is itself preempted by a medium-priority task is priority inversion. B
  is the cure, not the problem; C requires a circular resource dependency; D
  concerns memory, not priorities.
\item \textbf{B.} Priority inheritance temporarily raises the holder to the
  waiting task's priority so it releases the resource quickly. Time slicing
  (A), spooling (C), and paging (D) do not address priority inversion.
\item \textbf{C.} Round-robin gives each ready process a turn but ignores
  resource holding, so by itself it does not prevent priority inversion.
  Priority inheritance (A), the priority ceiling protocol (B), and priority
  boosting (D) all address inversion.
\item \textbf{A.} Being repeatedly passed over while others are served is
  starvation. Deadlock (B) is a circular wait with no progress possible;
  thrashing (C) is excessive paging; fragmentation (D) concerns memory or
  disk layout.
\item \textbf{A.} Aging gradually raises the priority of a long-waiting
  process so it eventually runs. Swapping (B), spooling (C), and caching (D)
  are unrelated storage techniques.
\item \textbf{B.} Preemptive priority scheduling interrupts the running
  process as soon as a higher-priority process becomes ready. A and C describe
  non-preemptive behaviour; D is not a preemption trigger.
\item \textbf{A.} RTOS expands to Real-Time Operating System, designed to meet
  time deadlines. B, C, and D are incorrect expansions.
\item \textbf{C.} Statements I and III are true. II is false: round-robin by
  itself does not prevent priority inversion. A includes false II; B includes
  II and omits I; D includes II.
\item \textbf{A.} The core that manages processes, memory, devices, and file
  systems is the Linux kernel. The shell (B) is the command interpreter; GNOME
  (C) is a desktop environment; GRUB (D) is a bootloader.
\item \textbf{A.} \texttt{/root} is the home directory of the superuser;
  \texttt{/home} holds ordinary users' home directories. B reverses the two;
  C mislabels \texttt{/root}; D mislabels \texttt{/home} as the top of the
  file system (which is \texttt{/}).
\item \textbf{B.} ext4 is the native, common default file system for Linux.
  NTFS (A) is the Windows default; FAT32 (C) is used for broad compatibility;
  exFAT (D) is a removable-media file system, not the Linux native default.
\item \textbf{B.} NTFS is the default for modern Windows installations. ext4
  (A) and ext3 (C) are Linux file systems; HFS+ (D) is an older Apple file
  system.
\item \textbf{C.} The Windows Store (Microsoft Store) is the built-in online
  marketplace for downloading and installing apps. It is not a compression
  tool (A), a defragmenter (B), or a Linux package manager (D).
\item \textbf{B.} Notepad is a plain-text editor, while WordPad supports basic
  rich-text formatting such as bold and fonts. A swaps them; C mislabels both
  as spreadsheets; D is false, since WordPad can save files (for example in
  \texttt{.rtf}).
\item \textbf{A.} Disk Cleanup frees disk space by removing temporary files
  and the Recycle Bin contents. The Defragmenter (B) rearranges file
  fragments; Device Manager (C) manages hardware; Task Scheduler (D) runs
  tasks on a schedule.
\item \textbf{A.} An SSD has no moving parts, so defragmenting gives no speed
  benefit and only adds write wear; Windows uses TRIM instead. B is
  irrelevant to defragmentation; C is false (an SSD is non-volatile); D is
  false.
\item \textbf{A.} Windows Fax and Scan is the imaging and fax utility used to
  scan documents and send or receive faxes. WordPad (B) is a text editor;
  Disk Cleanup (C) frees space; the Windows Store (D) installs apps.
\item \textbf{A.} \texttt{xcopy /s} copies the named directory including its
  subdirectories, excluding empty ones. It does not delete files (B), rename a
  file (C), or format a disk (D).

\item \textbf{A.} The ls command lists files. B, C, and D are not Linux commands.
\item \textbf{A.} Waiting processes are blocked. B, C, and D are not process states.
\item \textbf{A.} Linux is open source and multi-user. B, C, and D name unrelated objects.
\item \textbf{A, B, C.} ls, cd, and mkdir are standard utilities. D is not a command.
\item \textbf{A, B, C.} The process definition, scheduler role, and blocked state are correct. D is false: processes are runtime entities.
\item \textbf{A.} Resource preemption reclaims a resource to break deadlock. The other actions do not change process-resource allocation.

\item \textbf{A.} `pwd` prints the current directory. `ls` lists entries, `cd` changes directory, and `rm` removes files.

\item \textbf{A.} `ls` lists directory entries. `pwd` prints the path; `mkdir` and `rmdir` create/remove directories.

\item \textbf{A.} Round-robin cycles through ready processes using a time quantum. The other options do not specify that rotating time-slice policy.

\item \textbf{A, B, C.} Processes, deadlocks, and file systems are OS concepts. A driver is software that helps control a device.

\end{enumerate}
\unitopener{6}{Programming \& Data Structures}
\lessonheading{51}{Programming Foundations, Compiler and Data Structures}
\begin{enumerate}[label=\textbf{Q\arabic*.}, leftmargin=*]
\item \textit{[Original]} Expand OOP.
  \begin{enumerate}[label=\Alph*.]
    \item Object-Oriented Programming
    \item Object-Oriented Processing
    \item Open-Output Programming
    \item Operational Object Programming
  \end{enumerate}
\item \textit{[Original]} Expand LIFO.
  \begin{enumerate}[label=\Alph*.]
    \item Last In First Out
    \item Last In First Open
    \item Large Input Fast Output
    \item Last Input First Order
  \end{enumerate}
\item \textit{[Original]} Which language is the only one a computer executes
  directly?
  \begin{enumerate}[label=\Alph*.]
    \item Assembly language
    \item Machine language
    \item High-level language
    \item Pseudocode
  \end{enumerate}
\item \textit{[Original]} Assembly language uses short codes such as ADD and
  MOV, which are called:
  \begin{enumerate}[label=\Alph*.]
    \item mnemonics
    \item keywords
    \item variables
    \item functions
  \end{enumerate}
\item \textit{[Original]} Which translator converts an assembly-language
  program into machine code?
  \begin{enumerate}[label=\Alph*.]
    \item Compiler
    \item Interpreter
    \item Assembler
    \item Linker
  \end{enumerate}
\item \textit{[Original]} Which statement correctly distinguishes a compiler
  from an interpreter?
  \begin{enumerate}[label=\Alph*.]
    \item A compiler translates line by line, while an interpreter translates
      the whole program at once.
    \item A compiler translates the whole program at once, while an
      interpreter translates and executes it line by line.
    \item Both translate the whole program at once.
    \item A compiler executes source code directly without any translation.
  \end{enumerate}
\item \textit{[Original]} Source code is best described as:
  \begin{enumerate}[label=\Alph*.]
    \item the binary instructions executed directly by the CPU.
    \item the program written by the programmer in a high-level or assembly
      language.
    \item the machine-readable output produced by a compiler.
    \item firmware stored in read-only memory.
  \end{enumerate}
\item \textit{[Original]} The machine-readable code produced by a compiler
  from source code is called:
  \begin{enumerate}[label=\Alph*.]
    \item source code
    \item object code
    \item pseudocode
    \item natural-language code
  \end{enumerate}
\item \textit{[Original]} An algorithm is best defined as:
  \begin{enumerate}[label=\Alph*.]
    \item a pictorial representation using standard symbols.
    \item a finite, step-by-step procedure to solve a problem.
    \item the binary code executed by the CPU.
    \item a software tool for finding errors.
  \end{enumerate}
\item \textit{[Original]} In a flowchart, which symbol is used for a decision?
  \begin{enumerate}[label=\Alph*.]
    \item Rectangle
    \item Parallelogram
    \item Diamond
    \item Oval
  \end{enumerate}
\item \textit{[Original]} The process of finding and correcting errors in a
  program is called:
  \begin{enumerate}[label=\Alph*.]
    \item compiling
    \item debugging
    \item linking
    \item encoding
  \end{enumerate}
\item \textit{[Original]} A program compiles and runs but produces a wrong
  result. This is an example of a:
  \begin{enumerate}[label=\Alph*.]
    \item syntax error
    \item logical error
    \item runtime error
    \item translation error
  \end{enumerate}
\item \textit{[Original]} In object-oriented programming, a class is:
  \begin{enumerate}[label=\Alph*.]
    \item an instance created from an object.
    \item a blueprint or template from which objects are created.
    \item a permanent secondary-storage location.
    \item a machine-language instruction.
  \end{enumerate}
\item \textit{[Original]} Which OOP feature means ``one name, many forms''?
  \begin{enumerate}[label=\Alph*.]
    \item Encapsulation
    \item Inheritance
    \item Polymorphism
    \item Abstraction
  \end{enumerate}
\item \textit{[Original]} A stack follows which order?
  \begin{enumerate}[label=\Alph*.]
    \item FIFO
    \item LIFO
    \item Random
    \item Sorted
  \end{enumerate}
\item \textit{[Original]} In a stack, the insertion and deletion operations
  are called:
  \begin{enumerate}[label=\Alph*.]
    \item enqueue and dequeue
    \item push and pop
    \item insert and remove
    \item add and delete
  \end{enumerate}
\item \textit{[Original]} A queue follows which order?
  \begin{enumerate}[label=\Alph*.]
    \item LIFO
    \item FIFO
    \item Random
    \item Reverse
  \end{enumerate}
\item \textit{[Original]} Which data structure is hierarchical, consisting of
  a root node with child nodes?
  \begin{enumerate}[label=\Alph*.]
    \item Array
    \item Stack
    \item Tree
    \item Queue
  \end{enumerate}
\item \textit{[Original]} Which of the following is \textbf{NOT} true of a
  compiler?
  \begin{enumerate}[label=\Alph*.]
    \item It translates the whole program at once.
    \item It produces object code.
    \item It translates and executes the program line by line.
    \item It reports errors after translating the whole program.
  \end{enumerate}
\item \textit{[Original]} Evaluate the statements:
  \begin{enumerate}[label=\Roman*.]
    \item An array stores elements of the same type in contiguous memory
      locations.
    \item A linked list stores its elements in contiguous memory locations.
    \item A stack follows the Last In, First Out order.
  \end{enumerate}
  \begin{enumerate}[label=\Alph*.]
    \item I and II only
    \item II and III only
    \item I and III only
    \item I, II, and III
  \end{enumerate}

\item \textit{[JKSSB FAA-pattern]} A compiler translates a program:
  \begin{enumerate}[label=\Alph*.]
    \item as a whole before execution.
    \item one printed page at a time.
    \item only during paper printing.
    \item without reading the code.
  \end{enumerate}
\item \textit{[JKPSI-pattern]} Which data structure works on a last-in-first-out basis?
  \begin{enumerate}[label=\Alph*.]
    \item Stack
    \item Queue
    \item Printer tray
    \item Hard disk
  \end{enumerate}
\item \textit{[SSC CGL-pattern]} Which of the following is a high-level programming language?
  \begin{enumerate}[label=\Alph*.]
    \item Python
    \item Machine code
    \item Printer ink
    \item Speaker wire
  \end{enumerate}
\item \textit{[JKSSB-pattern]} \textbf{(MSQ: select ALL correct options.)} Which are programming languages?
  \begin{enumerate}[label=\Alph*.]
    \item C
    \item Java
    \item Python
    \item Hard disk
  \end{enumerate}
\item \textit{[SSC CGL-pattern]} \textbf{(MSQ: select ALL correct options.)} Which statements about translators and structures are correct?
  \begin{enumerate}[label=\Alph*.]
    \item Compilers translate whole programs at once.
    \item Interpreters translate line by line.
    \item A queue follows first-in-first-out order.
    \item A stack follows first-in-first-out order.
  \end{enumerate}
\item \textit{[Original]} Which translator converts an entire high-level program before execution?
  \begin{enumerate}[label=\Alph*.]
\item Compiler
   \item Interpreter
   \item Text editor
   \item Linker only
  \end{enumerate}

\item \textit{[Original]} Which data structure follows the LIFO principle?
  \begin{enumerate}[label=\Alph*.]
\item Stack
   \item Queue
   \item Array only
   \item Graph
  \end{enumerate}

\item \textit{[Original]} A written step-by-step procedure for solving a problem is an:
  \begin{enumerate}[label=\Alph*.]
\item Algorithm
   \item Operating system
   \item Database record
   \item Device driver
  \end{enumerate}

\item \textit{[Original]} Which language is closest to the processor's native instruction representation?
  \begin{enumerate}[label=\Alph*.]
\item Machine language
   \item High-level language
   \item Natural language
   \item Spreadsheet formula language only
  \end{enumerate}

\item \textit{[Original, MSQ drill]} \textbf{(MSQ: select ALL correct options.)} Which pairings are correct?
  \begin{enumerate}[label=\Alph*.]
\item Stack --- LIFO
   \item Queue --- FIFO
   \item Compiler --- translates a program before execution
   \item Debugging --- intentionally adds errors
  \end{enumerate}

\end{enumerate}
\answersection{6}
\answersubheading{51}{Programming Foundations, Compiler and Data Structures}
\begin{enumerate}[label=\textbf{Q\arabic*.}, leftmargin=*]
\item \textbf{A.} OOP stands for Object-Oriented Programming. B, C, and D are
  not the accepted expansion of the abbreviation.
\item \textbf{A.} LIFO stands for Last In, First Out. B, C, and D are not the
  correct full forms.
\item \textbf{B.} Machine language is the only language the CPU executes
  directly. Assembly (A) must be assembled, a high-level language (C) must be
  compiled or interpreted, and pseudocode (D) is not executable at all.
\item \textbf{A.} Short codes such as ADD and MOV are called mnemonics.
  Keywords (B) are words reserved by a high-level language; variables (C) and
  functions (D) are program elements, not instruction codes.
\item \textbf{C.} An assembler translates assembly language into machine
  code. A compiler (A) and an interpreter (B) translate high-level
  languages; a linker (D) combines object files but does not translate
  assembly source.
\item \textbf{B.} A compiler translates the whole program at once, while an
  interpreter translates and executes it line by line. A reverses the two;
  C wrongly claims both work on the whole program; D is false because a
  compiler does translate the source.
\item \textbf{B.} Source code is the program written by the programmer in a
  high-level or assembly language and is human-readable. A describes machine
  code, C describes object code, and D describes firmware.
\item \textbf{B.} The compiler's machine-readable output is called object
  code. Source code (A) is the human-written input, pseudocode (C) is a
  design notation, and natural-language code (D) is not a technical term.
\item \textbf{B.} An algorithm is a finite, step-by-step procedure to solve a
  problem. A describes a flowchart, C describes machine code, and D describes
  a debugger.
\item \textbf{C.} A decision is shown by a diamond (rhombus). A rectangle (A)
  is a process, a parallelogram (B) is input/output, and an oval (D) is a
  terminator for start or end.
\item \textbf{B.} Finding and correcting errors is debugging. Compiling (A)
  translates source code, linking (C) joins object files, and encoding (D)
  is unrelated to error correction.
\item \textbf{B.} A program that compiles and runs but gives a wrong result
  has a logical error. A syntax error (A) would stop compilation, a runtime
  error (C) occurs during execution, and a translation error (D) is not a
  standard category.
\item \textbf{B.} A class is a blueprint or template from which objects are
  created. A reverses the relation (an object is an instance of a class), C
  confuses a class with storage, and D confuses it with a machine
  instruction.
\item \textbf{C.} ``One name, many forms'' describes polymorphism.
  Encapsulation (A) binds data and methods, inheritance (B) reuses a parent
  class, and abstraction (D) hides implementation detail.
\item \textbf{B.} A stack follows LIFO (Last In, First Out). FIFO (A) is the
  queue order; random (C) and sorted (D) are not the stack rule.
\item \textbf{B.} Stack insertion and deletion are called push and pop.
  Enqueue and dequeue (A) belong to a queue; insert/remove (C) and
  add/delete (D) are generic terms, not the stack operations.
\item \textbf{B.} A queue follows FIFO (First In, First Out). LIFO (A) is the
  stack order; random (C) and reverse (D) are not queue rules.
\item \textbf{C.} A root node with child nodes describes a tree. An array
  (A) is a flat indexed collection, a stack (B) is linear LIFO, and a queue
  (D) is linear FIFO.
\item \textbf{C.} A compiler does not translate and execute line by line;
  that is the interpreter's method, so C is NOT true. Translating the whole
  program at once (A), producing object code (B), and reporting errors after
  the full translation (D) are all true of a compiler.
\item \textbf{C.} Statements I and III are correct: an array stores same-type
  elements in contiguous memory, and a stack is LIFO. II is false because a
  linked list does not store its elements contiguously; it links separate
  nodes by pointers. A includes false II, B includes false II, and D includes
  false II.

\item \textbf{A.} Compilers convert the whole program first. B, C, and D describe non-compiler behaviour.
\item \textbf{A.} Stacks are LIFO. Queues (B) are FIFO; C and D are hardware.
\item \textbf{A.} Python is high level. Machine code (B) is low level; C and D are materials.
\item \textbf{A, B, C.} C, Java, and Python are languages. A hard disk (D) is storage hardware.
\item \textbf{A, B, C.} Compiler, interpreter, and queue facts are correct. D is false: stacks are LIFO.
\item \textbf{A.} A compiler translates a program into another form before execution. An interpreter executes through translation as it runs; an editor edits text; a linker combines object modules.

\item \textbf{A.} A stack is last-in, first-out. A queue is FIFO; arrays and graphs are not defined by LIFO behavior.

\item \textbf{A.} An algorithm is a finite procedure for solving a problem. The other terms describe software/system/data objects.

\item \textbf{A.} Machine language is represented in processor instructions. High-level languages are more abstract; natural language is not processor code.

\item \textbf{A, B, C.} The first three are correct. Debugging identifies and fixes errors rather than adding them.

\end{enumerate}
\unitopener{7}{Office Applications \& File Formats}
\lessonheading{22}{MS Word Fundamentals}
\begin{enumerate}[label=\textbf{Q\arabic*.}, leftmargin=*]
\item \textit{[Original]} MS Word is an example of a:
  \begin{enumerate}[label=\Alph*.]
    \item spreadsheet program.
    \item word processor.
    \item database management system.
    \item presentation program.
  \end{enumerate}
\item \textit{[Original]} A document created in the modern version of MS Word
  is saved by default with which file extension?
  \begin{enumerate}[label=\Alph*.]
    \item .xlsx
    \item .pptx
    \item .docx
    \item .txt
  \end{enumerate}
\item \textit{[Original]} Expand QAT as used in the MS Word window.
  \begin{enumerate}[label=\Alph*.]
    \item Quick Access Toolbar
    \item Quick Action Tab
    \item Qualified Access Table
    \item Quick Alignment Tool
  \end{enumerate}
\item \textit{[Original]} The Ribbon in MS Word replaced which older
  interface element?
  \begin{enumerate}[label=\Alph*.]
    \item The status bar
    \item Menus and toolbars
    \item The document ruler
    \item The keyboard shortcuts
  \end{enumerate}
\item \textit{[Original]} The Font and Paragraph groups are found on which
  Ribbon tab?
  \begin{enumerate}[label=\Alph*.]
    \item Insert
    \item View
    \item Home
    \item Layout
  \end{enumerate}
\item \textit{[Original]} By default, which three commands appear on the
  Quick Access Toolbar?
  \begin{enumerate}[label=\Alph*.]
    \item Save, Undo, and Redo
    \item Bold, Italic, and Underline
    \item Copy, Cut, and Paste
    \item New, Open, and Print
  \end{enumerate}
\item \textit{[Original]} Which keyboard shortcut saves the current MS Word
  document?
  \begin{enumerate}[label=\Alph*.]
    \item Ctrl+N
    \item Ctrl+S
    \item Ctrl+P
    \item Ctrl+O
  \end{enumerate}
\item \textit{[Original]} Which keyboard shortcut opens the Print screen in
  MS Word?
  \begin{enumerate}[label=\Alph*.]
    \item Ctrl+S
    \item Ctrl+E
    \item Ctrl+P
    \item Ctrl+Enter
  \end{enumerate}
\item \textit{[Original]} Which of the following is \textbf{NOT} true of the
  Quick Access Toolbar?
  \begin{enumerate}[label=\Alph*.]
    \item It appears in the top-left corner of the window.
    \item It can be customised by adding or removing commands.
    \item It is fixed and cannot be customised in any way.
    \item It can be moved below the Ribbon.
  \end{enumerate}
\item \textit{[Original]} Which paragraph alignment makes the text flush with
  \textbf{both} the left and right margins?
  \begin{enumerate}[label=\Alph*.]
    \item Left
    \item Centre
    \item Right
    \item Justify
  \end{enumerate}
\item \textit{[Original]} Which keyboard shortcut is used to \textbf{centre}
  the selected paragraph?
  \begin{enumerate}[label=\Alph*.]
    \item Ctrl+L
    \item Ctrl+E
    \item Ctrl+R
    \item Ctrl+J
  \end{enumerate}
\item \textit{[Original]} A hanging indent:
  \begin{enumerate}[label=\Alph*.]
    \item pushes only the first line of a paragraph inward.
    \item pushes every line except the first inward.
    \item centres every line of the paragraph.
    \item adds equal space before and after the paragraph.
  \end{enumerate}
\item \textit{[Original]} Which feature should be used for a list in which
  the \textbf{order} of the items matters?
  \begin{enumerate}[label=\Alph*.]
    \item Bullets
    \item Numbering
    \item Text wrapping
    \item Justification
  \end{enumerate}
\item \textit{[Original]} In a table, the intersection of a row and a column
  is called a:
  \begin{enumerate}[label=\Alph*.]
    \item cell
    \item field
    \item record
    \item margin
  \end{enumerate}
\item \textit{[Original]} Evaluate the statements:
  \begin{enumerate}[label=\Roman*.]
    \item The Home tab contains the Font and Paragraph groups.
    \item By default the Quick Access Toolbar shows Save, Undo, and Redo.
    \item A Page Break allows different headers and footers in different
      parts of a document.
  \end{enumerate}
  \begin{enumerate}[label=\Alph*.]
    \item I only
    \item I and II only
    \item II and III only
    \item I, II, and III
  \end{enumerate}
\item \textit{[Original]} Which area repeats at the \textbf{bottom} of every
  page in a document?
  \begin{enumerate}[label=\Alph*.]
    \item Header
    \item Footer
    \item Body
    \item Margin
  \end{enumerate}
\item \textit{[Original]} Which set correctly distinguishes the two page
  orientations?
  \begin{enumerate}[label=\Alph*.]
    \item Portrait is wider than tall; Landscape is taller than wide.
    \item Portrait is taller than wide; Landscape is wider than tall.
    \item Both orientations are taller than wide.
    \item Orientation has no effect on page dimensions.
  \end{enumerate}
\item \textit{[Original]} Which break must be used to give one part of a
  document a \textbf{different header} from the rest?
  \begin{enumerate}[label=\Alph*.]
    \item Page Break
    \item Section Break
    \item Line Break
    \item Column Break
  \end{enumerate}
\item \textit{[Original]} The setting that controls how text flows around a
  picture is called:
  \begin{enumerate}[label=\Alph*.]
    \item text wrapping
    \item text alignment
    \item line spacing
    \item indentation
  \end{enumerate}
\item \textit{[Original]} Which option correctly matches the MS Word command
  with its purpose?
  \begin{enumerate}[label=\Alph*.]
    \item Justify --- lines up text on the right margin only
    \item Footer --- repeats at the top of every page
    \item Print preview --- shows how pages will appear before printing
    \item Bullets --- numbers the items of an ordered list
  \end{enumerate}

\item \textit{[JKSSB FAA-pattern]} Which feature combines a standard letter with an address list to print personalised copies?
  \begin{enumerate}[label=\Alph*.]
    \item Mail Merge
    \item Word Count
    \item Zoom slider
    \item Ruler guide
  \end{enumerate}
\item \textit{[JKPSI-pattern]} Which shortcut saves a Word document?
  \begin{enumerate}[label=\Alph*.]
    \item Ctrl+S
    \item Ctrl+P
    \item Ctrl+B
    \item Ctrl+E
  \end{enumerate}
\item \textit{[SSC CGL-pattern]} Which alignment centres each paragraph on the page?
  \begin{enumerate}[label=\Alph*.]
    \item Center alignment
    \item Left alignment
    \item Right alignment
    \item Justify with gaps
  \end{enumerate}
\item \textit{[JKSSB-pattern]} \textbf{(MSQ: select ALL correct options.)} Which are valid text-formatting options in MS Word?
  \begin{enumerate}[label=\Alph*.]
    \item Bold
    \item Italic
    \item Underline
    \item Paper weight
  \end{enumerate}
\item \textit{[SSC CGL-pattern]} \textbf{(MSQ: select ALL correct options.)} Which statements about MS Word are correct?
  \begin{enumerate}[label=\Alph*.]
    \item A .docx file is the default Word format.
    \item Track Changes records edits for review.
    \item Headers repeat content at the top of pages.
    \item Saving always prints the document.
  \end{enumerate}
\item \textit{[PYQ-adapted: SSC CGL 2025 Tier-II Paper-I, 19-Jan-2026 Q14; SRC-18]} A report starts a new section on the next page. Which setting can be changed
  for that section without changing the previous section?
  \begin{enumerate}[label=\Alph*.]
\item AutoCorrect replacement list
   \item Page orientation and margins
   \item Normal style's font face
   \item Document Properties metadata
  \end{enumerate}

\item \textit{[PYQ-adapted: JKSSB Junior Assistant 2026 Q73, key:final; SRC-03/SRC-06]} Which break is needed when a later part of a Word document must use a different
  header or footer?
  \begin{enumerate}[label=\Alph*.]
\item Section Break
   \item Page Break
   \item Line Break
   \item Column Break
  \end{enumerate}

\item \textit{[PYQ-adapted: JKSSB Website Operator 2026 Q77, key:provisional; SRC-04/SRC-07]} Which Word feature combines a main document with a recipient list to produce
  personalized letters?
  \begin{enumerate}[label=\Alph*.]
\item Track Changes
   \item Mail Merge
   \item SmartArt
   \item Word Count
  \end{enumerate}

\item \textit{[Original]} Which shortcut opens the Find and Replace dialog in Word?
  \begin{enumerate}[label=\Alph*.]
\item Ctrl+H
   \item Ctrl+P
   \item Ctrl+B
   \item Ctrl+K
  \end{enumerate}

\item \textit{[Original, MSQ drill]} \textbf{(MSQ: select ALL correct options.)} Which settings can be specific to
  a Word section?
  \begin{enumerate}[label=\Alph*.]
\item Page orientation
   \item Page margins
   \item Headers and footers
   \item The font definition of the Normal style
  \end{enumerate}

\end{enumerate}
\lessonheading{23}{MS Word High-Yield}
\begin{enumerate}[label=\textbf{Q\arabic*.}, leftmargin=*]
\item \textit{[Original]} Which keyboard shortcut opens the Find and Replace
  dialog box in MS Word?
  \begin{enumerate}[label=\Alph*.]
    \item Ctrl+F
    \item Ctrl+H
    \item Ctrl+G
    \item Ctrl+K
  \end{enumerate}
\item \textit{[Original]} Which function key starts the spelling and grammar
  check in MS Word?
  \begin{enumerate}[label=\Alph*.]
    \item F5
    \item F7
    \item F9
    \item F11
  \end{enumerate}
\item \textit{[Original]} Which shortcut is used to insert a hyperlink in
  MS Word?
  \begin{enumerate}[label=\Alph*.]
    \item Ctrl+K
    \item Ctrl+L
    \item Ctrl+H
    \item Ctrl+J
  \end{enumerate}
\item \textit{[Original]} In MS Word, \textbf{Ctrl+F} is used to:
  \begin{enumerate}[label=\Alph*.]
    \item find and replace text.
    \item find (locate) text without replacing it.
    \item insert a hyperlink.
    \item open the print options.
  \end{enumerate}
\item \textit{[Original]} In MS Word, a \textbf{bookmark} is best described as:
  \begin{enumerate}[label=\Alph*.]
    \item a clickable link to a web page.
    \item a named location or selection in the document.
    \item a note attached to a point in the document.
    \item a recorded edit that stores the author's name.
  \end{enumerate}
\item \textit{[Original]} Which feature in Word allows
you to combine letters with a database of addresses?
  \begin{enumerate}[label=\Alph*.]
    \item Track Changes
    \item Header \& Footer
    \item AutoCorrect
    \item Mail Merge
  \end{enumerate}
\item \textit{[Original]} What is the process in Microsoft
Word that combines names and addresses with a standard document called?
  \begin{enumerate}[label=\Alph*.]
    \item Mail Merge
    \item Document Linking
    \item Data Consolidation
    \item Template Binding
  \end{enumerate}
\item \textit{[Original]} In a mail merge, the \textbf{main document} is:
  \begin{enumerate}[label=\Alph*.]
    \item the list of names and addresses.
    \item the standard letter or template that contains the merge fields.
    \item the printer settings used for the merged output.
    \item the folder in which the merged files are stored.
  \end{enumerate}
\item \textit{[Original]} Which of the following breaks must be used so that a
  single Word document can have \textbf{different headers and footers} in
  different parts of the document?
  \begin{enumerate}[label=\Alph*.]
    \item Page Break
    \item Line Break
    \item Column Break
    \item Section Break
  \end{enumerate}
\item \textit{[Original]} Which of the following statements
is true about MS Word?
  \begin{enumerate}[label=\Alph*.]
    \item Using Section Breaks allows different headers and footers in the same
      document.
    \item Using Page Break allows you to have different headers and footers in
      the same document.
    \item Headers and footers cannot be changed once set in a document.
    \item Headers and footers are automatically deleted when you close the
      document.
  \end{enumerate}
\item \textit{[Original]} Which of the following statements
is \textbf{NOT} correct about MS Word?
  \begin{enumerate}[label=\Alph*.]
    \item A content control can restrict editing in a document; however it
      cannot enforce data validation.
    \item You can apply different first-line indents and hanging indents to
      different paragraphs within the same style.
    \item The Navigation Pane allows searching by headings, pages, or results,
      but cannot directly display comments.
    \item In Track Changes, accepting a change permanently removes the author's
      metadata associated with that change.
  \end{enumerate}
\item \textit{[Original]} A \textbf{comment} in a Word document is used to:
  \begin{enumerate}[label=\Alph*.]
    \item create a clickable link to another file.
    \item attach a reviewer's note without changing the document text.
    \item record every edit made to the document.
    \item mark a location so a hyperlink can jump to it.
  \end{enumerate}
\item \textit{[Original]} When \textbf{Track Changes} is enabled in MS Word,
  each recorded edit stores:
  \begin{enumerate}[label=\Alph*.]
    \item only the page number on which the edit occurs.
    \item the author's name and the time of the change.
    \item the printer used to print the document.
    \item the file name of the document only.
  \end{enumerate}
\item \textit{[Original]} Evaluate the following statements:
  \begin{enumerate}[label=\Roman*.]
    \item Ctrl+H opens the Find and Replace dialog box.
    \item Track Changes records the author of each edit.
    \item A Page Break allows different headers and footers in the same
      document.
  \end{enumerate}
  \begin{enumerate}[label=\Alph*.]
    \item I and II only
    \item I and III only
    \item II and III only
    \item I, II and III
  \end{enumerate}
\item \textit{[Original]} Applying a \textbf{Heading} style to a paragraph in
  MS Word enables the paragraph to:
  \begin{enumerate}[label=\Alph*.]
    \item appear in the Navigation Pane and a Table of Contents.
    \item be printed on a separate page automatically.
    \item be tracked as a revision with the author's name.
    \item be exported to PDF only.
  \end{enumerate}
\item \textit{[Original]} Expand \textbf{PDF}, the fixed-layout format to which
  a Word document can be exported.
  \begin{enumerate}[label=\Alph*.]
    \item Printed Document File
    \item Portable Document Format
    \item Public Data Format
    \item Personal Document Folder
  \end{enumerate}
\item \textit{[Original]} Which pair of MS Word shortcuts is correctly matched?
  \begin{enumerate}[label=\Alph*.]
    \item Ctrl+S --- Print; Ctrl+P --- Save
    \item Ctrl+S --- Save; Ctrl+P --- Print
    \item Ctrl+S --- Spelling; Ctrl+P --- Paste
    \item Ctrl+S --- Select all; Ctrl+P --- Print preview
  \end{enumerate}
\item \textit{[Original]} All of the following are MS Word review or editing
  features \textbf{EXCEPT}:
  \begin{enumerate}[label=\Alph*.]
    \item Track Changes
    \item Comments
    \item Bookmarks
    \item PivotTable
  \end{enumerate}
\item \textit{[Original]} Which keyboard shortcut in Microsoft
Word is used to remove the document window split?
  \begin{enumerate}[label=\Alph*.]
    \item Ctrl + Shift + S
    \item Alt + Shift + C
    \item Ctrl + Alt + C
    \item Shift + F6
  \end{enumerate}
\item \textit{[Original]} Which of the following pairs is correctly matched?
  \begin{enumerate}[label=\Alph*.]
    \item Ctrl+H --- run the spelling and grammar check
    \item F7 --- open Find and Replace
    \item Ctrl+K --- insert a hyperlink
    \item Ctrl+F --- open the print options
  \end{enumerate}

\item \textit{[JKSSB FAA-pattern]} Which key runs spelling and grammar checking in MS Word?
  \begin{enumerate}[label=\Alph*.]
    \item F7
    \item F5
    \item F2
    \item F9
  \end{enumerate}
\item \textit{[JKPSI-pattern]} Which feature lets different document sections use different headers?
  \begin{enumerate}[label=\Alph*.]
    \item Section Break
    \item Page Colour
    \item Zoom Out
    \item Word Wrap
  \end{enumerate}
\item \textit{[SSC CGL-pattern]} Which option shows a document exactly as it will print?
  \begin{enumerate}[label=\Alph*.]
    \item Print Layout view
    \item Draft invoice pad
    \item Paper tray
    \item Printer cable
  \end{enumerate}
\item \textit{[JKSSB-pattern]} \textbf{(MSQ: select ALL correct options.)} Which items can appear in a Word header or footer?
  \begin{enumerate}[label=\Alph*.]
    \item Page numbers
    \item Document title text
    \item Date fields
    \item The computer's power supply
  \end{enumerate}
\item \textit{[SSC CGL-pattern]} \textbf{(MSQ: select ALL correct options.)} Which statements about MS Word are correct?
  \begin{enumerate}[label=\Alph*.]
    \item Ctrl+B applies bold formatting.
    \item A table organises text in rows and columns.
    \item Portrait and landscape are page orientations.
    \item F7 starts a slide show.
  \end{enumerate}
\item \textit{[PYQ-adapted: SSC CGL 2025 Tier-II Paper-I, 19-Jan-2026 Q16; SRC-18]} A user changes one heading's appearance but wants other headings using the same style to remain unchanged. What is the safest approach?
  \begin{enumerate}[label=\Alph*.]
\item Modify only that instance's direct formatting
   \item Redefine the shared style
   \item Change the document template for every file
   \item Use AutoCorrect
  \end{enumerate}

\item \textit{[Original]} Which shortcut commonly opens the Find and Replace dialog in Word?
  \begin{enumerate}[label=\Alph*.]
\item Ctrl+H
   \item Ctrl+P
   \item Ctrl+S
   \item Ctrl+L
  \end{enumerate}

\item \textit{[Original]} Which key commonly starts spelling and grammar checking in Word?
  \begin{enumerate}[label=\Alph*.]
\item F1
   \item F5
   \item F7
   \item F12
  \end{enumerate}

\item \textit{[Original]} Which feature records insertions and deletions so a reviewer can accept or reject them?
  \begin{enumerate}[label=\Alph*.]
\item Track Changes
   \item Word Count
   \item Mail Merge
   \item Page Color
  \end{enumerate}

\item \textit{[Original, MSQ drill]} \textbf{(MSQ: select ALL correct options.)} Which features support document review/collaboration?
  \begin{enumerate}[label=\Alph*.]
\item Comments
   \item Track Changes
   \item Accept/Reject Changes
   \item Mail Merge
  \end{enumerate}

\end{enumerate}
\lessonheading{24}{Excel Structure}
\begin{enumerate}[label=\textbf{Q\arabic*.}, leftmargin=*]
\item \textit{[Original]} What is an Excel workbook?
  \begin{enumerate}[label=\Alph*.]
    \item A single page of rows and columns
    \item The Excel file that can contain one or more worksheets
    \item A formula that calculates a result
    \item The bar that shows the contents of a cell
  \end{enumerate}
\item \textit{[Original]} What is an Excel worksheet?
  \begin{enumerate}[label=\Alph*.]
    \item The file in which Excel data is saved
    \item A single sheet of rows and columns inside a workbook
    \item A block of selected cells
    \item The name given to a cell
  \end{enumerate}
\item \textit{[Original]} In a worksheet, the intersection of a row and a
  column is called a:
  \begin{enumerate}[label=\Alph*.]
    \item range
    \item cell
    \item worksheet
    \item workbook
  \end{enumerate}
\item \textit{[Original]} Which of the following is a valid cell address in
  Excel?
  \begin{enumerate}[label=\Alph*.]
    \item 5B
    \item B5
    \item B:5
    \item B\_5
  \end{enumerate}
\item \textit{[Original]} The currently selected cell in a worksheet is
  called the:
  \begin{enumerate}[label=\Alph*.]
    \item active cell
    \item name cell
    \item anchor cell
    \item formula cell
  \end{enumerate}
\item \textit{[Original]} Which element displays the contents of the active
  cell, including a formula?
  \begin{enumerate}[label=\Alph*.]
    \item Name Box
    \item Formula Bar
    \item Sheet tab
    \item Status bar
  \end{enumerate}
\item \textit{[Original]} The Name Box, located to the left of the Formula
  Bar, normally displays:
  \begin{enumerate}[label=\Alph*.]
    \item the contents of the active cell
    \item the address of the active cell or the name of the selected range
    \item the sum of the selected numbers
    \item the name of the workbook
  \end{enumerate}
\item \textit{[Original]} A group of two or more selected cells, such as
  A1:C10, is called a:
  \begin{enumerate}[label=\Alph*.]
    \item workbook
    \item range
    \item tab
    \item field
  \end{enumerate}
\item \textit{[Original]} What is the total number of rows in a worksheet in
  Excel 2007 and later?
  \begin{enumerate}[label=\Alph*.]
    \item 65,536
    \item 1,024,000
    \item 1,048,576
    \item 16,384
  \end{enumerate}
\item \textit{[Original]} What is the total number of columns in a worksheet
  in Excel 2007 and later?
  \begin{enumerate}[label=\Alph*.]
    \item 256
    \item 1,024
    \item 16,384
    \item 1,048,576
  \end{enumerate}
\item \textit{[Original]} The last column of an Excel 2007 and later
  worksheet is labelled:
  \begin{enumerate}[label=\Alph*.]
    \item Z
    \item ZZ
    \item IV
    \item XFD
  \end{enumerate}
\item \textit{[Original]} Which statement about a workbook and a worksheet is
  correct?
  \begin{enumerate}[label=\Alph*.]
    \item A workbook is a page inside a worksheet.
    \item A worksheet is the Excel file, and a workbook is a part of it.
    \item A workbook is the Excel file and can contain one or more
      worksheets.
    \item A workbook can contain only one worksheet, and no more.
  \end{enumerate}
\item \textit{[Original]} Which of the following is \textbf{NOT} true of a
  worksheet in Excel?
  \begin{enumerate}[label=\Alph*.]
    \item It is made up of rows and columns.
    \item It has a tab at the bottom of the Excel window.
    \item It is the same as the workbook file.
    \item It can hold text, numbers, and formulas.
  \end{enumerate}
\item \textit{[Original]} Which of the following represents a valid
  \textbf{range} in Excel?
  \begin{enumerate}[label=\Alph*.]
    \item B5
    \item A1:C10
    \item \$A\$1
    \item =SUM(A1:A2)
  \end{enumerate}
\item \textit{[Original]} In Excel, by default, text is aligned to the
  \underline{\hspace{1.6cm}} of a cell, while numbers are aligned to the
  \underline{\hspace{1.6cm}}.
  \begin{enumerate}[label=\Alph*.]
    \item right; left
    \item left; right
    \item centre; left
    \item left; centre
  \end{enumerate}
\item \textit{[Original]} Which of the following is \textbf{NOT} a basic data
  type that can be entered in an Excel cell?
  \begin{enumerate}[label=\Alph*.]
    \item Text
    \item Number
    \item Formula
    \item Row
  \end{enumerate}
\item \textit{[Original]} In Excel 2007 and later, the default file extension
  of a workbook is:
  \begin{enumerate}[label=\Alph*.]
    \item .xls
    \item .xlsx
    \item .docx
    \item .csv
  \end{enumerate}
\item \textit{[Original]} Evaluate the following statements about Excel:
  \begin{enumerate}[label=\Roman*.]
    \item A workbook can contain more than one worksheet.
    \item The Name Box shows the address of the active cell.
    \item The Formula Bar shows the contents of the active cell.
  \end{enumerate}
  \begin{enumerate}[label=\Alph*.]
    \item I and II only
    \item II and III only
    \item I and III only
    \item I, II, and III
  \end{enumerate}
\item \textit{[Original]} How many cells are there in the range A1:B10?
  \begin{enumerate}[label=\Alph*.]
    \item 10
    \item 12
    \item 20
    \item 100
  \end{enumerate}
\item \textit{[Original]} Which pairing is correctly matched?
  \begin{enumerate}[label=\Alph*.]
    \item Workbook --- a single page of rows and columns
    \item Worksheet --- the Excel file saved on disk
    \item Active cell --- the currently selected cell
    \item Range --- a single cell address such as B5
  \end{enumerate}

\item \textit{[JKSSB FAA-pattern]} A worksheet cell is identified by its:
  \begin{enumerate}[label=\Alph*.]
    \item cell address (column letter + row number).
    \item printer name.
    \item monitor size.
    \item speaker volume.
  \end{enumerate}
\item \textit{[JKPSI-pattern]} An Excel workbook is a collection of:
  \begin{enumerate}[label=\Alph*.]
    \item worksheets.
    \item printer trays.
    \item wall posters.
    \item envelope stamps.
  \end{enumerate}
\item \textit{[SSC CGL-pattern]} The intersection of a row and a column in Excel is called a:
  \begin{enumerate}[label=\Alph*.]
    \item cell.
    \item bell.
    \item shell.
    \item well.
  \end{enumerate}
\item \textit{[JKSSB-pattern]} \textbf{(MSQ: select ALL correct options.)} Which are parts of the Excel window?
  \begin{enumerate}[label=\Alph*.]
    \item Formula bar
    \item Name box
    \item Worksheet grid
    \item Engine radiator
  \end{enumerate}
\item \textit{[SSC CGL-pattern]} \textbf{(MSQ: select ALL correct options.)} Which statements about Excel structure are correct?
  \begin{enumerate}[label=\Alph*.]
    \item Columns are labelled with letters.
    \item Rows are numbered.
    \item The active cell receives typed input.
    \item A workbook holds only one cell.
  \end{enumerate}
\item \textit{[Original]} The intersection of a worksheet row and column is a:
  \begin{enumerate}[label=\Alph*.]
\item Cell
   \item Workbook
   \item Formula bar
   \item Ribbon
  \end{enumerate}

\item \textit{[Original]} A collection of worksheets saved in one Excel file is a:
  \begin{enumerate}[label=\Alph*.]
\item Workbook
   \item Cell range
   \item Chart area
   \item Field
  \end{enumerate}

\item \textit{[Original]} Which box displays the address of the active cell and can be used to jump to a cell reference?
  \begin{enumerate}[label=\Alph*.]
\item Name Box
   \item Status bar
   \item Formula bar
   \item Title bar
  \end{enumerate}

\item \textit{[Original]} In the reference B7, which part identifies the column?
  \begin{enumerate}[label=\Alph*.]
\item B
   \item 7
   \item B7
   \item Neither part
  \end{enumerate}

\item \textit{[Original, MSQ drill]} \textbf{(MSQ: select ALL correct options.)} Which Excel terms are matched correctly?
  \begin{enumerate}[label=\Alph*.]
\item Workbook --- file containing worksheets
   \item Worksheet --- grid of rows and columns
   \item Cell --- intersection of a row and column
   \item Formula bar --- identifies only the active worksheet tab
  \end{enumerate}

\end{enumerate}
\lessonheading{25}{Excel Formulas}
\begin{enumerate}[label=\textbf{Q\arabic*.}, leftmargin=*]
\item \textit{[Original]} Which statement best describes a formula in
  MS Excel?
  \begin{enumerate}[label=\Alph*.]
    \item A fixed value typed into a cell that is never recalculated.
    \item An instruction, beginning with \texttt{=}, that calculates a cell's
      value from operands such as references, constants, and functions.
    \item A permanent record of the cell's last result.
    \item A formatting rule applied to a range of cells.
  \end{enumerate}
\item \textit{[Original]} Which symbol is used for multiplication in an Excel
  formula?
  \begin{enumerate}[label=\Alph*.]
    \item \texttt{x}
    \item \texttt{*}
    \item \texttt{.}
    \item \texttt{/}
  \end{enumerate}
\item \textit{[Original]} In an Excel cell reference, what does the dollar
  sign \texttt{\$} do?
  \begin{enumerate}[label=\Alph*.]
    \item It locks the part of the reference that follows it, so that part
      does not change when the formula is copied.
    \item It converts the cell value to currency format.
    \item It marks the cell as read-only.
    \item It divides the cell value by 100.
  \end{enumerate}
\item \textit{[PYQ WO 2026 Q78, key:provisional]} In MS Excel, a formula must
  begin with which character?
  \begin{enumerate}[label=\Alph*.]
    \item \texttt{+}
    \item \texttt{=}
    \item \texttt{@}
    \item \texttt{\#}
  \end{enumerate}
\item \textit{[Original]} Which of the following is an absolute reference?
  \begin{enumerate}[label=\Alph*.]
    \item \texttt{A1}
    \item \texttt{\$A\$1}
    \item \texttt{A\$1}
    \item \texttt{\$A1}
  \end{enumerate}
\item \textit{[Original]} Which operator is used in an Excel formula to join
  (concatenate) two pieces of text?
  \begin{enumerate}[label=\Alph*.]
    \item \texttt{\&}
    \item \texttt{+}
    \item \texttt{:}
    \item \texttt{@}
  \end{enumerate}
\item \textit{[Original]} In the formula \texttt{=A1+\$A\$1}, which reference
  will \textbf{NOT} change when the formula is copied to another cell?
  \begin{enumerate}[label=\Alph*.]
    \item \texttt{A1}
    \item \texttt{\$A\$1}
    \item Both references change.
    \item Neither reference changes.
  \end{enumerate}
\item \textit{[Original]} Cell B1 contains the formula \texttt{=A1}. The
  formula is copied to cell B2. What does B2 contain?
  \begin{enumerate}[label=\Alph*.]
    \item \texttt{=A1}
    \item \texttt{=A2}
    \item \texttt{=B2}
    \item \texttt{=\$A\$1}
  \end{enumerate}
\item \textit{[Original]} Cell B1 contains the formula \texttt{=\$A\$1}. The
  formula is copied to cell B5. What does B5 contain?
  \begin{enumerate}[label=\Alph*.]
    \item \texttt{=\$A\$1}
    \item \texttt{=\$A\$5}
    \item \texttt{=\$A5}
    \item \texttt{=A\$1}
  \end{enumerate}
\item \textit{[Original]} Cell B1 contains the formula \texttt{=\$A1}. The
  formula is copied to cell C2. What does C2 contain?
  \begin{enumerate}[label=\Alph*.]
    \item \texttt{=\$A1}
    \item \texttt{=\$A2}
    \item \texttt{=\$B2}
    \item \texttt{=\$B1}
  \end{enumerate}
\item \textit{[Original]} Cell B1 contains the formula \texttt{=A\$1}. The
  formula is copied to cell C2. What does C2 contain?
  \begin{enumerate}[label=\Alph*.]
    \item \texttt{=B\$1}
    \item \texttt{=B\$2}
    \item \texttt{=A\$1}
    \item \texttt{=C\$1}
  \end{enumerate}
\item \textit{[Original]} Cell C2 contains the formula \texttt{=A2*\$B\$1}.
  The formula is copied to cell C3. What does C3 contain?
  \begin{enumerate}[label=\Alph*.]
    \item \texttt{=A2*\$B\$1}
    \item \texttt{=A3*\$B\$1}
    \item \texttt{=A2*\$B\$2}
    \item \texttt{=A3*\$B\$2}
  \end{enumerate}
\item \textit{[Original]} What does the Excel error \texttt{\#DIV/0!} mean?
  \begin{enumerate}[label=\Alph*.]
    \item A number is divided by zero or by an empty cell.
    \item A referenced cell has been deleted.
    \item The wrong type of operand was used.
    \item A function name is misspelled.
  \end{enumerate}
\item \textit{[Original]} What does the Excel error \texttt{\#REF!} mean?
  \begin{enumerate}[label=\Alph*.]
    \item A reference is invalid because the referenced cell was deleted.
    \item A formula divides by zero.
    \item A text value was used where a number was expected.
    \item The formula is missing its leading \texttt{=}.
  \end{enumerate}
\item \textit{[Original]} What does the Excel error \texttt{\#VALUE!} mean?
  \begin{enumerate}[label=\Alph*.]
    \item The wrong type of operand was used, such as text where a number is
      needed.
    \item A referenced file was not found.
    \item A referenced cell was deleted.
    \item The requested value is not available.
  \end{enumerate}
\item \textit{[Original]} Which of the following is \textbf{NOT} true of a
  relative reference such as \texttt{A1}?
  \begin{enumerate}[label=\Alph*.]
    \item It is stored as an offset from the formula's own cell.
    \item It shifts when the formula is copied to another cell.
    \item It always points to the same cell no matter where it is copied.
    \item It is the default type of reference in Excel.
  \end{enumerate}
\item \textit{[Original]} Which of the following correctly states the effect
  of the dollar sign in a cell reference?
  \begin{enumerate}[label=\Alph*.]
    \item \texttt{\$A1} locks both the column and the row.
    \item \texttt{A\$1} locks the row, and \texttt{\$A1} locks the column.
    \item \texttt{A\$1} locks the column, and \texttt{\$A1} locks the row.
    \item \texttt{\$A\$1} locks only the row.
  \end{enumerate}
\item \textit{[Original]} Evaluate the following statements:
  \begin{enumerate}[label=\Roman*.]
    \item A cell entry that begins with \texttt{+} is still treated as a
      formula by Excel.
    \item An absolute reference such as \texttt{\$A\$1} does not change when
      the formula is copied.
    \item A relative reference such as \texttt{A1} changes when the formula is
      copied.
  \end{enumerate}
  \begin{enumerate}[label=\Alph*.]
    \item I only
    \item II and III only
    \item I and II only
    \item I, II, and III
  \end{enumerate}
\item \textit{[Original]} In an Excel formula, what does the operator
  \texttt{<>} mean?
  \begin{enumerate}[label=\Alph*.]
    \item Not equal to
    \item Equal to
    \item Greater than
    \item Less than or equal to
  \end{enumerate}
\item \textit{[Original]} What is the result of the Excel formula
  \texttt{=2+3*4}?
  \begin{enumerate}[label=\Alph*.]
    \item 14
    \item 20
    \item 24
    \item 9
  \end{enumerate}

\item \textit{[JKSSB FAA-pattern]} Every Excel formula must begin with:
  \begin{enumerate}[label=\Alph*.]
    \item =
    \item +
    \item \%
    \item \#
  \end{enumerate}
\item \textit{[JKPSI-pattern]} In =A1+\$B\$2, the reference \$B\$2 is which type?
  \begin{enumerate}[label=\Alph*.]
    \item Absolute reference
    \item Printer reference
    \item Audio reference
    \item Paper reference
  \end{enumerate}
\item \textit{[SSC CGL-pattern]} Which operator joins two text strings in Excel?
  \begin{enumerate}[label=\Alph*.]
    \item \&
    \item +
    \item \%
    \item =
  \end{enumerate}
\item \textit{[JKSSB-pattern]} \textbf{(MSQ: select ALL correct options.)} Which are valid Excel formula elements?
  \begin{enumerate}[label=\Alph*.]
    \item Cell references
    \item Arithmetic operators
    \item Built-in functions
    \item Printer ink colours
  \end{enumerate}
\item \textit{[SSC CGL-pattern]} \textbf{(MSQ: select ALL correct options.)} Which statements about Excel formulas are correct?
  \begin{enumerate}[label=\Alph*.]
    \item Formulas recalculate when inputs change.
    \item \$ signs lock rows/columns in absolute references.
    \item Copying a relative formula adjusts its addresses.
    \item Formulas must end with a full stop.
  \end{enumerate}
\item \textit{[Original]} Which symbol makes both the row and column absolute in the reference \$A\$1?
  \begin{enumerate}[label=\Alph*.]
\item Dollar sign (\$)
   \item Percent sign (\%)
   \item Hash sign (\#)
   \item Ampersand (\&)
  \end{enumerate}

\item \textit{[Original]} When the formula =A1 is copied one column to the right, a relative reference normally becomes:
  \begin{enumerate}[label=\Alph*.]
\item =B1
   \item =\$A\$1
   \item =A\$1
   \item =A2
  \end{enumerate}

\item \textit{[Original]} Which reference locks the row but allows the column to change?
  \begin{enumerate}[label=\Alph*.]
\item A\$1
   \item \$A1
   \item \$A\$1
   \item A1
  \end{enumerate}

\item \textit{[Original]} Which operator combines text values in many Excel formulas?
  \begin{enumerate}[label=\Alph*.]
\item \&
   \item \^{}
   \item /
   \item >
  \end{enumerate}

\item \textit{[Original, MSQ drill]} \textbf{(MSQ: select ALL correct options.)} Which are valid Excel cell references?
  \begin{enumerate}[label=\Alph*.]
\item A1
   \item \$B\$7
   \item C\$4
   \item \$D8
  \end{enumerate}

\end{enumerate}
\lessonheading{26}{Excel Functions}
\begin{enumerate}[label=\textbf{Q\arabic*.}, leftmargin=*]
\item \textit{[PYQ JA 2026 Q74, key:final]} An Excel cell contains a nested IF
  formula:
  \begin{center}
    \texttt{=IF(A1>50,"High",IF(A1>25,"Medium","Low"))}
  \end{center}
  Cell A1 is blank. What will the formula display? 
  \begin{enumerate}[label=\Alph*.]
    \item High
    \item Medium
    \item Low
    \item \texttt{\#VALUE!}
  \end{enumerate}
\item \textit{[Original]} Which function returns the average (arithmetic mean)
  of the numbers in a range?
  \begin{enumerate}[label=\Alph*.]
    \item \texttt{AVERAGE}
    \item \texttt{MEAN}
    \item \texttt{TOTAL}
    \item \texttt{MEDIAN}
  \end{enumerate}
\item \textit{[Original]} Which function counts only the cells in a range that
  contain numbers?
  \begin{enumerate}[label=\Alph*.]
    \item \texttt{COUNT}
    \item \texttt{COUNTA}
    \item \texttt{COUNTIF}
    \item \texttt{SUM}
  \end{enumerate}
\item \textit{[Original]} A range A1:A10 contains 6 numbers, 2 text labels,
  and 2 blank cells. What does \texttt{=COUNTA(A1:A10)} return?
  \begin{enumerate}[label=\Alph*.]
    \item 6
    \item 8
    \item 10
    \item 2
  \end{enumerate}
\item \textit{[Original]} Using the same range (6 numbers, 2 text labels, 2
  blanks), what does \texttt{=COUNT(A1:A10)} return?
  \begin{enumerate}[label=\Alph*.]
    \item 6
    \item 8
    \item 10
    \item 4
  \end{enumerate}
\item \textit{[Original]} Which function counts every non-empty cell in a
  range, whether it holds a number or text?
  \begin{enumerate}[label=\Alph*.]
    \item \texttt{COUNTA}
    \item \texttt{COUNT}
    \item \texttt{COUNTBLANK}
    \item \texttt{SUM}
  \end{enumerate}
\item \textit{[Original]} Which function returns the largest number in a
  range?
  \begin{enumerate}[label=\Alph*.]
    \item \texttt{MAX}
    \item \texttt{MIN}
    \item \texttt{HIGH}
    \item \texttt{TOP}
  \end{enumerate}
\item \textit{[Original]} Which function returns the smallest number in a
  range?
  \begin{enumerate}[label=\Alph*.]
    \item \texttt{MIN}
    \item \texttt{MAX}
    \item \texttt{MINIMUM}
    \item \texttt{LEAST}
  \end{enumerate}
\item \textit{[Original]} Which function adds together all the valid numbers
  in a range?
  \begin{enumerate}[label=\Alph*.]
    \item \texttt{SUM}
    \item \texttt{ADD}
    \item \texttt{TOTAL}
    \item \texttt{PLUS}
  \end{enumerate}
\item \textit{[Original]} Which keyboard shortcut inserts the AutoSum function
  in MS Excel?
  \begin{enumerate}[label=\Alph*.]
    \item \texttt{Alt + =}
    \item \texttt{Ctrl + =}
    \item \texttt{Shift + F3}
    \item \texttt{Ctrl + S}
  \end{enumerate}
\item \textit{[Original]} By default, which formula does the AutoSum button
  insert?
  \begin{enumerate}[label=\Alph*.]
    \item \texttt{=SUM(...)}
    \item \texttt{=AVERAGE(...)}
    \item \texttt{=COUNT(...)}
    \item \texttt{=MAX(...)}
  \end{enumerate}
\item \textit{[Original]} Which statement correctly distinguishes a function
  from a formula in MS Excel?
  \begin{enumerate}[label=\Alph*.]
    \item A formula must start with \texttt{=}, while a function is a named
      built-in operation used inside a formula.
    \item A function must start with \texttt{=}, while a formula is a named
      built-in operation.
    \item A formula and a function are exactly the same thing.
    \item A formula can never contain a function.
  \end{enumerate}
\item \textit{[Original]} Which error value appears when a formula divides a
  number by zero or by a blank cell?
  \begin{enumerate}[label=\Alph*.]
    \item \texttt{\#DIV/0!}
    \item \texttt{\#VALUE!}
    \item \texttt{\#NAME?}
    \item \texttt{\#REF!}
  \end{enumerate}
\item \textit{[Original]} Which error value appears when a formula refers to a
  name or function that Excel does not recognise, such as a misspelt function
  name?
  \begin{enumerate}[label=\Alph*.]
    \item \texttt{\#NAME?}
    \item \texttt{\#REF!}
    \item \texttt{\#N/A}
    \item \texttt{\#NUM!}
  \end{enumerate}
\item \textit{[Original]} Which error value appears when a formula uses the
  wrong type of data --- for example, when text is added to a number?
  \begin{enumerate}[label=\Alph*.]
    \item \texttt{\#VALUE!}
    \item \texttt{\#NAME?}
    \item \texttt{\#DIV/0!}
    \item \texttt{\#N/A}
  \end{enumerate}
\item \textit{[Original]} Which error value appears when a formula refers to a
  cell that has been deleted, so the reference is no longer valid?
  \begin{enumerate}[label=\Alph*.]
    \item \texttt{\#REF!}
    \item \texttt{\#NAME?}
    \item \texttt{\#VALUE!}
    \item \texttt{\#NULL!}
  \end{enumerate}
\item \textit{[Original]} Which error value means that a value needed by the
  formula is not available, such as when a lookup finds no match?
  \begin{enumerate}[label=\Alph*.]
    \item \texttt{\#N/A}
    \item \texttt{\#NUM!}
    \item \texttt{\#REF!}
    \item \texttt{\#DIV/0!}
  \end{enumerate}
\item \textit{[Original]} Which of the following is \textbf{NOT} a standard
  Excel error value?
  \begin{enumerate}[label=\Alph*.]
    \item \texttt{\#DIV/0!}
    \item \texttt{\#VALUE!}
    \item \texttt{\#NAME?}
    \item \texttt{\#TOTAL!}
  \end{enumerate}
\item \textit{[Original]} Evaluate the statements about MS Excel:
  \begin{enumerate}[label=\Roman*.]
    \item \texttt{COUNT} counts only numeric cells, whereas \texttt{COUNTA}
      counts all non-empty cells.
    \item In a numeric comparison, a blank cell is treated as zero.
    \item A formula must always begin with the equal sign (\texttt{=}).
  \end{enumerate}
  \begin{enumerate}[label=\Alph*.]
    \item I and II only
    \item II and III only
    \item I and III only
    \item I, II, and III
  \end{enumerate}
\item \textit{[Original]} Match the function with what it returns:
  \begin{enumerate}[label=\Roman*.]
    \item \texttt{SUM}
    \item \texttt{MAX}
    \item \texttt{COUNT}
  \end{enumerate}
  \begin{enumerate}[label=(\alph*)]
    \item the largest number in a range
    \item the total of the numbers in a range
    \item the number of numeric cells in a range
  \end{enumerate}
  \begin{enumerate}[label=\Alph*.]
    \item I-b, II-a, III-c
    \item I-a, II-b, III-c
    \item I-c, II-a, III-b
    \item I-b, II-c, III-a
  \end{enumerate}

\item \textit{[JKSSB FAA-pattern]} Which function totals a range of numbers?
  \begin{enumerate}[label=\Alph*.]
    \item SUM
    \item PRINT
    \item SPEAK
    \item DRAW
  \end{enumerate}
\item \textit{[JKPSI-pattern]} Which function counts only the cells holding numbers?
  \begin{enumerate}[label=\Alph*.]
    \item COUNT
    \item COUNTA for numbers only
    \item PRINTER
    \item SPEAKER
  \end{enumerate}
\item \textit{[SSC CGL-pattern]} Which function returns the largest value in a range?
  \begin{enumerate}[label=\Alph*.]
    \item MAX
    \item MIN for largest
    \item AVG for largest
    \item SUM for largest
  \end{enumerate}
\item \textit{[JKSSB-pattern]} \textbf{(MSQ: select ALL correct options.)} Which Excel functions work with conditions?
  \begin{enumerate}[label=\Alph*.]
    \item IF
    \item COUNTIF
    \item SUMIF
    \item PRINTIF
  \end{enumerate}
\item \textit{[SSC CGL-pattern]} \textbf{(MSQ: select ALL correct options.)} Which statements about Excel functions are correct?
  \begin{enumerate}[label=\Alph*.]
    \item AVERAGE returns the arithmetic mean.
    \item MIN returns the smallest value.
    \item IF chooses between two results by a test.
    \item SUM prints the worksheet.
  \end{enumerate}
\item \textit{[PYQ-adapted: SSC CGL 2025 Tier-II Paper-I, 19-Jan-2026 Q4; SRC-18]} Which Excel function counts cells that satisfy one specified condition?
  \begin{enumerate}[label=\Alph*.]
\item COUNT
   \item COUNTBLANK
   \item COUNTIF
   \item SUMIF
  \end{enumerate}

\item \textit{[PYQ-adapted: SSC CGL 2025 Tier-II Paper-I, 19-Jan-2026 Q9; SRC-18]} A teacher needs to show only rows where marks exceed 75 while keeping the
  remaining records in the worksheet. Which tool should be used?
  \begin{enumerate}[label=\Alph*.]
\item Sort
   \item Filter
   \item Remove Duplicates
   \item Subtotal
  \end{enumerate}

\item \textit{[Original]} Which function counts non-empty cells, including cells containing text?
  \begin{enumerate}[label=\Alph*.]
\item COUNTA
   \item COUNT
   \item COUNTBLANK
   \item SUM
  \end{enumerate}

\item \textit{[Original]} In a formula, which function tests a condition and returns one value when it
  is true and another when it is false?
  \begin{enumerate}[label=\Alph*.]
\item MAX
   \item IF
   \item AVERAGE
   \item ROUND
  \end{enumerate}

\item \textit{[Original, MSQ drill]} \textbf{(MSQ: select ALL correct options.)} Which statements about Excel
  counting functions are correct?
  \begin{enumerate}[label=\Alph*.]
\item COUNT counts cells containing numbers.
   \item COUNTA counts non-empty cells.
   \item COUNTBLANK counts empty cells.
   \item COUNTIF sums the values that meet a criterion.
  \end{enumerate}

\end{enumerate}
\lessonheading{27}{Excel Data Tools}
\begin{enumerate}[label=\textbf{Q\arabic*.}, leftmargin=*]
\item \textit{[Original]} The Excel feature that rearranges rows into a chosen
  order, such as A to Z or Z to A, is called:
  \begin{enumerate}[label=\Alph*.]
    \item Filter
    \item Sort
    \item Goal Seek
    \item Protect Sheet
  \end{enumerate}
\item \textit{[Original]} Which keyboard shortcut turns AutoFilter on or off for
  the selected range?
  \begin{enumerate}[label=\Alph*.]
    \item \textbf{Ctrl}$+$\textbf{L}
    \item \textbf{Ctrl}$+$\textbf{Shift}$+$\textbf{L}
    \item \textbf{Alt}$+$\textbf{F}
    \item \textbf{Ctrl}$+$\textbf{F}
  \end{enumerate}
\item \textit{[Original]} Which keyboard shortcut converts the selected range
  into an Excel table?
  \begin{enumerate}[label=\Alph*.]
    \item \textbf{Ctrl}$+$\textbf{T}
    \item \textbf{Ctrl}$+$\textbf{F}
    \item \textbf{Ctrl}$+$\textbf{Shift}$+$\textbf{T}
    \item \textbf{Alt}$+$\textbf{T}
  \end{enumerate}
\item \textit{[Original]} When a filter is applied to a column, the rows that do
  not match the condition are:
  \begin{enumerate}[label=\Alph*.]
    \item permanently deleted from the worksheet
    \item temporarily hidden, but not deleted
    \item moved to a new worksheet
    \item converted into a PivotTable
  \end{enumerate}
\item \textit{[Original]} Which of the following is \textbf{NOT} true about
  sorting in Excel?
  \begin{enumerate}[label=\Alph*.]
    \item Data can be sorted by more than one level using Add Level.
    \item Data can be sorted by cell colour or font colour.
    \item A custom list can sort by days of the week or months.
    \item Ascending order arranges numbers from largest to smallest.
  \end{enumerate}
\item \textit{[Original]} Which feature is provided by an Excel table but
  \textbf{not} by a plain range of cells?
  \begin{enumerate}[label=\Alph*.]
    \item Structured references such as \texttt{Table1[Sales]}
    \item The ability to enter text
    \item The ability to enter formulas
    \item The ability to print the data
  \end{enumerate}
\item \textit{[Original]} Which statement about an Excel table and a PivotTable
  is correct?
  \begin{enumerate}[label=\Alph*.]
    \item A PivotTable is a type of Excel table.
    \item An Excel table can serve as the source data for a PivotTable.
    \item A PivotTable permanently changes the source data.
    \item An Excel table cannot be filtered.
  \end{enumerate}
\item \textit{[Original]} Which chart type is best used to show the parts of a
  single whole, such as the share of each subject in a student's total marks?
  \begin{enumerate}[label=\Alph*.]
    \item Line chart
    \item Scatter (XY) chart
    \item Pie chart
    \item Column chart
  \end{enumerate}
\item \textit{[Original]} When the source data of a chart is changed, the chart:
  \begin{enumerate}[label=\Alph*.]
    \item remains unchanged until it is re-created
    \item updates automatically
    \item is deleted
    \item is converted into a PivotTable
  \end{enumerate}
\item \textit{[Original]} In a chart, the component that names each data series
  is the:
  \begin{enumerate}[label=\Alph*.]
    \item legend
    \item value axis
    \item plot area
    \item data label
  \end{enumerate}
\item \textit{[Original]} Which of the following is \textbf{NOT} one of the four
  areas into which fields are dragged in a PivotTable?
  \begin{enumerate}[label=\Alph*.]
    \item Rows
    \item Columns
    \item Values
    \item Insert
  \end{enumerate}
\item \textit{[Original]} After the source data of a PivotTable is edited, what
  must be done to display the new summary?
  \begin{enumerate}[label=\Alph*.]
    \item Refresh the PivotTable
    \item Restart the computer
    \item Nothing --- it updates automatically
    \item Delete and rename the worksheet
  \end{enumerate}
\item \textit{[Original]} Evaluate the statements:
  \begin{enumerate}[label=\Roman*.]
    \item A PivotTable summarises data but is not the source table itself.
    \item A PivotChart is linked to a PivotTable.
    \item A filter permanently deletes the rows it hides.
  \end{enumerate}
  \begin{enumerate}[label=\Alph*.]
    \item I only
    \item I and II only
    \item II and III only
    \item I, II, and III
  \end{enumerate}
\item \textit{[Original]} To count the rows that satisfy two conditions at the
  same time, which function is designed specifically for the purpose?
  \begin{enumerate}[label=\Alph*.]
    \item COUNT
    \item COUNTIFS
    \item SUM
    \item AVERAGE
  \end{enumerate}
\item \textit{[Original]} Which is the correct general form of the COUNTIFS
  function?
  \begin{enumerate}[label=\Alph*.]
    \item \texttt{COUNTIFS(range1, range2, crit1, crit2)}
    \item \texttt{COUNTIFS(range1, crit1, range2, crit2)}
    \item \texttt{COUNTIFS(crit1, crit2)}
    \item \texttt{COUNTIFS(range, crit1+crit2)}
  \end{enumerate}
\item \textit{[Original]} Which statement about COUNTIFS and SUMPRODUCT is
  correct?
  \begin{enumerate}[label=\Alph*.]
    \item COUNTIFS cannot count with two criteria at once.
    \item SUMPRODUCT can only add numbers and cannot count rows.
    \item COUNTIFS applies AND logic across paired range-and-criteria
      arguments, and SUMPRODUCT can also count multi-criterion rows using
      boolean arrays.
    \item Both functions can apply only OR logic.
  \end{enumerate}
\item \textit{[Original]} Which printing setting repeats the header row at the
  top of every printed page?
  \begin{enumerate}[label=\Alph*.]
    \item Print Area
    \item Print Titles
    \item Freeze Panes
    \item Page Break Preview
  \end{enumerate}
\item \textit{[Original]} In Excel, are gridlines printed by default?
  \begin{enumerate}[label=\Alph*.]
    \item Yes, always
    \item No, not unless Print gridlines is enabled
    \item Only in landscape orientation
    \item Only when a Print Area is set
  \end{enumerate}
\item \textit{[Original]} Which statement about protection in Excel is correct?
  \begin{enumerate}[label=\Alph*.]
    \item A locked cell cannot be edited even when the sheet is unprotected.
    \item Turning on Protect Workbook stops individual cells from being edited.
    \item A cell locking setting takes effect only after the sheet is
      protected.
    \item Protection encrypts the workbook file with a password.
  \end{enumerate}
\item \textit{[Original]} Which what-if analysis tool finds the input value
  required to reach a specified target result?
  \begin{enumerate}[label=\Alph*.]
    \item Goal Seek
    \item Scenario Manager
    \item Data Table
    \item PivotChart
  \end{enumerate}

\item \textit{[JKSSB FAA-pattern]} Which Excel tool arranges rows alphabetically or numerically?
  \begin{enumerate}[label=\Alph*.]
    \item Sort
    \item Print Preview
    \item Spell Pad
    \item Paint Brush
  \end{enumerate}
\item \textit{[JKPSI-pattern]} Which tool shows only rows meeting chosen criteria?
  \begin{enumerate}[label=\Alph*.]
    \item Filter
    \item Printer
    \item Speaker
    \item Scanner
  \end{enumerate}
\item \textit{[SSC CGL-pattern]} A PivotTable is used to:
  \begin{enumerate}[label=\Alph*.]
    \item summarise and analyse large datasets.
    \item boil drinking water.
    \item print envelopes.
    \item play music files.
  \end{enumerate}
\item \textit{[JKSSB-pattern]} \textbf{(MSQ: select ALL correct options.)} Which Excel tools help analyse data?
  \begin{enumerate}[label=\Alph*.]
    \item Sort
    \item Filter
    \item PivotTable
    \item Paperweight
  \end{enumerate}
\item \textit{[SSC CGL-pattern]} \textbf{(MSQ: select ALL correct options.)} Which statements about Excel data tools are correct?
  \begin{enumerate}[label=\Alph*.]
    \item Sorting reorders rows by chosen columns.
    \item Filtering hides rows that fail the criteria.
    \item Charts visualise worksheet data.
    \item Sorting deletes duplicate files from disk.
  \end{enumerate}
\item \textit{[PYQ-adapted: SSC CGL 2025 Tier-II Paper-I, 19-Jan-2026 Q9; SRC-18]} A teacher needs to show only rows where marks exceed 75 while keeping other records. Which tool should be used?
  \begin{enumerate}[label=\Alph*.]
\item Filter
   \item Remove Duplicates
   \item Delete Sheet
   \item Sort A to Z
  \end{enumerate}

\item \textit{[Original]} Which Excel feature summarizes and reorganizes a large data table by categories?
  \begin{enumerate}[label=\Alph*.]
\item PivotTable
   \item AutoCorrect
   \item Mail Merge
   \item Slide Master
  \end{enumerate}

\item \textit{[Original]} Which chart is generally suitable for showing a trend over time?
  \begin{enumerate}[label=\Alph*.]
\item Line chart
   \item Pie chart
   \item Doughnut chart
   \item Radar chart
  \end{enumerate}

\item \textit{[Original]} Which tool adds drop-down controls to choose which records are displayed?
  \begin{enumerate}[label=\Alph*.]
\item AutoFilter
   \item Goal Seek
   \item Spell Check
   \item Page Break Preview
  \end{enumerate}

\item \textit{[Original, MSQ drill]} \textbf{(MSQ: select ALL correct options.)} Which statements about Excel data tools are correct?
  \begin{enumerate}[label=\Alph*.]
\item Sorting changes record order.
   \item Filtering hides non-matching rows from the current view.
   \item A PivotTable can aggregate data by fields.
   \item Removing duplicates only changes font formatting.
  \end{enumerate}

\end{enumerate}
\lessonheading{28}{Access Database Basics}
\begin{enumerate}[label=\textbf{Q\arabic*.}, leftmargin=*]
\item \textit{[Original]} Expand DBMS.
  \begin{enumerate}[label=\Alph*.]
    \item Data Backup Management System
    \item Database Management System
    \item Digital Base Management System
    \item Data Block Mapping System
  \end{enumerate}
\item \textit{[Original]} Microsoft Access is best described as a:
  \begin{enumerate}[label=\Alph*.]
    \item spreadsheet program
    \item word processor
    \item database management system (DBMS)
    \item presentation program
  \end{enumerate}
\item \textit{[Original]} Which statement correctly distinguishes a database
  from a DBMS?
  \begin{enumerate}[label=\Alph*.]
    \item A database is the software, and a DBMS is the data it stores.
    \item A database is the organised collection of data, and a DBMS is the
      software that manages it.
    \item A database and a DBMS are two names for the same thing.
    \item A DBMS is a type of spreadsheet used to store data.
  \end{enumerate}
\item \textit{[Original]} Microsoft Access is part of which application suite?
  \begin{enumerate}[label=\Alph*.]
    \item Microsoft Office
    \item Adobe Creative Cloud
    \item Google Workspace
    \item LibreOffice
  \end{enumerate}
\item \textit{[Original]} In a database table, a row is called a:
  \begin{enumerate}[label=\Alph*.]
    \item field
    \item table
    \item record
    \item query
  \end{enumerate}
\item \textit{[Original]} In a database table, a column is called a:
  \begin{enumerate}[label=\Alph*.]
    \item record
    \item field
    \item report
    \item form
  \end{enumerate}
\item \textit{[Original]} A collection of related records stored together in
  rows and columns is called a:
  \begin{enumerate}[label=\Alph*.]
    \item field
    \item record
    \item table
    \item data type
  \end{enumerate}
\item \textit{[Original]} Which Access object stores the actual data of a
  database?
  \begin{enumerate}[label=\Alph*.]
    \item Query
    \item Form
    \item Report
    \item Table
  \end{enumerate}
\item \textit{[Original]} Which Access object is used to present data in a
  printable layout?
  \begin{enumerate}[label=\Alph*.]
    \item Table
    \item Query
    \item Form
    \item Report
  \end{enumerate}
\item \textit{[Original]} Which Access object provides a screen through which
  users can enter and view data conveniently?
  \begin{enumerate}[label=\Alph*.]
    \item Form
    \item Table
    \item Index
    \item Data type
  \end{enumerate}
\item \textit{[Original]} Which Access data type automatically inserts a
  unique number for each new record?
  \begin{enumerate}[label=\Alph*.]
    \item Short Text
    \item Number
    \item AutoNumber
    \item Yes/No
  \end{enumerate}
\item \textit{[Original]} Which Access data type is best suited to store only
  a True or False value?
  \begin{enumerate}[label=\Alph*.]
    \item Number
    \item Yes/No
    \item Date/Time
    \item Short Text
  \end{enumerate}
\item \textit{[Original]} Which Access data type is best suited to store a
  student's name?
  \begin{enumerate}[label=\Alph*.]
    \item Number
    \item Date/Time
    \item Short Text
    \item AutoNumber
  \end{enumerate}
\item \textit{[Original]} A Validation Rule of \texttt{Between 0 And 100} on a
  Marks field ensures that:
  \begin{enumerate}[label=\Alph*.]
    \item only values from 0 to 100 are accepted.
    \item the field cannot be left blank.
    \item the field is indexed for faster search.
    \item each record gets a unique number.
  \end{enumerate}
\item \textit{[Original]} Which field property holds the message that is
  displayed when a Validation Rule is violated?
  \begin{enumerate}[label=\Alph*.]
    \item Default Value
    \item Input Mask
    \item Validation Text
    \item Field Size
  \end{enumerate}
\item \textit{[Original]} Which field property, when set to Yes, stops a field
  from being left blank?
  \begin{enumerate}[label=\Alph*.]
    \item Required
    \item Format
    \item Caption
    \item Indexed
  \end{enumerate}
\item \textit{[Original]} The main purpose of an index on a field is to:
  \begin{enumerate}[label=\Alph*.]
    \item check that entered values are valid.
    \item speed up searching and sorting.
    \item store the data permanently.
    \item uniquely identify each record by itself.
  \end{enumerate}
\item \textit{[Original]} Which of the following is \textbf{NOT} true about a
  primary key?
  \begin{enumerate}[label=\Alph*.]
    \item It uniquely identifies each record.
    \item It cannot be left blank.
    \item It can contain duplicate values.
    \item It is often an AutoNumber field.
  \end{enumerate}
\item \textit{[Original]} Evaluate the statements:
  \begin{enumerate}[label=\Roman*.]
    \item A field is a column in a table.
    \item A record is a row in a table.
    \item A table is a collection of related records.
  \end{enumerate}
  \begin{enumerate}[label=\Alph*.]
    \item I only
    \item I and II only
    \item II and III only
    \item I, II, and III
  \end{enumerate}
\item \textit{[Original]} Which pairing is correctly matched?
  \begin{enumerate}[label=\Alph*.]
    \item Field --- one complete entry in a table
    \item Record --- one kind of information in a table
    \item AutoNumber --- a data type that adds a unique number to each new
      record
    \item DBMS --- the organised collection of data
  \end{enumerate}

\item \textit{[JKSSB FAA-pattern]} In a database table, each row is called a:
  \begin{enumerate}[label=\Alph*.]
    \item record.
    \item printer.
    \item speaker.
    \item cable.
  \end{enumerate}
\item \textit{[JKPSI-pattern]} In a database table, each column is called a:
  \begin{enumerate}[label=\Alph*.]
    \item field.
    \item forest.
    \item folder glue.
    \item paper clip.
  \end{enumerate}
\item \textit{[SSC CGL-pattern]} DBMS stands for:
  \begin{enumerate}[label=\Alph*.]
    \item Database Management System
    \item Data Binding Machine Shop
    \item Daily Book Market Stall
    \item Direct Board Mail Service
  \end{enumerate}
\item \textit{[JKSSB-pattern]} \textbf{(MSQ: select ALL correct options.)} Which are components of a database table?
  \begin{enumerate}[label=\Alph*.]
    \item Records (rows)
    \item Fields (columns)
    \item Primary key values
    \item Printer ink levels
  \end{enumerate}
\item \textit{[SSC CGL-pattern]} \textbf{(MSQ: select ALL correct options.)} Which statements about databases are correct?
  \begin{enumerate}[label=\Alph*.]
    \item A table stores data in rows and columns.
    \item A record describes one entity such as one student.
    \item A field holds one attribute such as a name.
    \item A database table is a printed paper form.
  \end{enumerate}
\item \textit{[Original]} In a relational database table, a row is commonly called a:
  \begin{enumerate}[label=\Alph*.]
\item Record
   \item Field
   \item Query
   \item Form
  \end{enumerate}

\item \textit{[Original]} In MS Access, a field corresponds most closely to a table:
  \begin{enumerate}[label=\Alph*.]
\item Column
   \item Row
   \item Report
   \item Database file
  \end{enumerate}

\item \textit{[Original]} Which data type is most appropriate for a field storing a calendar date?
  \begin{enumerate}[label=\Alph*.]
\item Date/Time
   \item Currency
   \item Attachment
   \item Yes/No
  \end{enumerate}

\item \textit{[Original]} A DBMS primarily helps users:
  \begin{enumerate}[label=\Alph*.]
\item Define, store, retrieve, and manage data
   \item Design computer circuits
   \item Print only slide presentations
   \item Replace all operating systems
  \end{enumerate}

\item \textit{[Original, MSQ drill]} \textbf{(MSQ: select ALL correct options.)} Which are common database objects in Access?
  \begin{enumerate}[label=\Alph*.]
\item Tables
   \item Queries
   \item Forms
   \item Reports
  \end{enumerate}

\end{enumerate}
\lessonheading{29}{Access Keys and Relationships}
\begin{enumerate}[label=\textbf{Q\arabic*.}, leftmargin=*]
\item \textit{[Original]} In a database, the abbreviation \textbf{PK} stands
  for:
  \begin{enumerate}[label=\Alph*.]
    \item Public Key
    \item Primary Key
    \item Parent Key
    \item Private Key
  \end{enumerate}
\item \textit{[Original]} In a database, the abbreviation \textbf{FK} stands
  for:
  \begin{enumerate}[label=\Alph*.]
    \item Field Key
    \item Full Key
    \item Foreign Key
    \item Foreign Column
  \end{enumerate}
\item \textit{[Original]} Which key uniquely identifies each record in a
  table?
  \begin{enumerate}[label=\Alph*.]
    \item Foreign key
    \item Alternate key
    \item Primary key
    \item Secondary key
  \end{enumerate}
\item \textit{[Original]} Which of the following is \textbf{NOT} a
  characteristic of a primary key?
  \begin{enumerate}[label=\Alph*.]
    \item It uniquely identifies each record.
    \item It cannot be left blank.
    \item It may contain duplicate values.
    \item A table has at most one of it.
  \end{enumerate}
\item \textit{[Original]} A field in one table that refers to the primary key
  of another table is called a:
  \begin{enumerate}[label=\Alph*.]
    \item primary key
    \item composite key
    \item foreign key
    \item candidate key
  \end{enumerate}
\item \textit{[Original]} Which key is made up of \textbf{two or more fields
  used together} to identify a record?
  \begin{enumerate}[label=\Alph*.]
    \item Primary key
    \item Composite key
    \item Foreign key
    \item Alternate key
  \end{enumerate}
\item \textit{[Original]} Which key may contain duplicate values and may be
  left blank?
  \begin{enumerate}[label=\Alph*.]
    \item Primary key
    \item Composite key
    \item Foreign key
    \item Candidate key
  \end{enumerate}
\item \textit{[Original]} How many primary keys can a single table have?
  \begin{enumerate}[label=\Alph*.]
    \item Exactly one
    \item At most one
    \item At least two
    \item Any number
  \end{enumerate}
\item \textit{[Original]} Which type of relationship is the \textbf{most
  common} in MS Access?
  \begin{enumerate}[label=\Alph*.]
    \item One-to-one
    \item One-to-many
    \item Many-to-many
    \item Zero-to-one
  \end{enumerate}
\item \textit{[Original]} In a one-to-many relationship, the ``many'' side of
  the relationship holds the:
  \begin{enumerate}[label=\Alph*.]
    \item primary key
    \item foreign key
    \item composite key
    \item alternate key
  \end{enumerate}
\item \textit{[Original]} Which relationship type cannot be stored directly
  and must be implemented with a junction table?
  \begin{enumerate}[label=\Alph*.]
    \item One-to-one
    \item One-to-many
    \item Many-to-many
    \item One-to-zero
  \end{enumerate}
\item \textit{[Original]} A junction table typically uses as its primary key:
  \begin{enumerate}[label=\Alph*.]
    \item a single AutoNumber field only
    \item a composite key formed from the foreign keys of both tables
    \item the primary key of the first table
    \item the primary key of the second table
  \end{enumerate}
\item \textit{[Original]} Referential integrity ensures that:
  \begin{enumerate}[label=\Alph*.]
    \item every table has a composite key.
    \item a foreign-key value matches an existing primary-key value in the
      related table.
    \item all fields in a table have the same data type.
    \item every table has more than one primary key.
  \end{enumerate}
\item \textit{[Original]} Which option automatically deletes the related child
  records when a parent record is deleted?
  \begin{enumerate}[label=\Alph*.]
    \item Cascade Update
    \item Cascade Delete
    \item Referential integrity
    \item Composite key
  \end{enumerate}
\item \textit{[Original]} Evaluate the statements:
  \begin{enumerate}[label=\Roman*.]
    \item A primary key can contain duplicate values.
    \item A foreign key can contain duplicate values.
    \item A table can have more than one primary key.
  \end{enumerate}
  \begin{enumerate}[label=\Alph*.]
    \item I only
    \item II only
    \item I and II only
    \item II and III only
  \end{enumerate}
\item \textit{[Original]} The file extension \texttt{.mdb} is the MS Access
  database format used by:
  \begin{enumerate}[label=\Alph*.]
    \item Access 2003 and earlier versions
    \item Access 2007 and later versions
    \item Access 2016 only
    \item all current versions
  \end{enumerate}
\item \textit{[Original]} The file extension \texttt{.accdb} is:
  \begin{enumerate}[label=\Alph*.]
    \item the older Access format from before 2003
    \item the newer, default Access format used from Access 2007 onward
    \item a spreadsheet format
    \item a text format
  \end{enumerate}
\item \textit{[Original]} Which of the following is \textbf{NOT} correct about
  the Access file formats?
  \begin{enumerate}[label=\Alph*.]
    \item \texttt{.mdb} is the older format.
    \item \texttt{.accdb} is the newer format.
    \item \texttt{.mdb} is the newer format introduced with Access 2007.
    \item \texttt{.accdb} supports attachment and multi-valued fields.
  \end{enumerate}
\item \textit{[Original]} Which pairing is correctly matched?
  \begin{enumerate}[label=\Alph*.]
    \item Primary key --- may contain duplicate values
    \item Foreign key --- uniquely identifies each record in its own table
    \item Composite key --- a primary key made of two or more fields
    \item Junction table --- stores a one-to-one relationship
  \end{enumerate}
\item \textit{[Original]} A candidate key that was \textbf{not} chosen as the
  primary key is called a(n):
  \begin{enumerate}[label=\Alph*.]
    \item foreign key
    \item alternate key
    \item composite key
    \item junction key
  \end{enumerate}

\item \textit{[JKSSB FAA-pattern]} A primary key in a table must be:
  \begin{enumerate}[label=\Alph*.]
    \item unique and never null.
    \item repeated in every row.
    \item left blank for safety.
    \item shared with the printer.
  \end{enumerate}
\item \textit{[JKPSI-pattern]} A foreign key is used to:
  \begin{enumerate}[label=\Alph*.]
    \item link a table to another table's primary key.
    \item lock the computer cabinet.
    \item open the office door.
    \item start the printer motor.
  \end{enumerate}
\item \textit{[SSC CGL-pattern]} Which relationship connects one record to many records?
  \begin{enumerate}[label=\Alph*.]
    \item One-to-many
    \item One-to-one printer
    \item Many-to-none paper
    \item All-to-printer
  \end{enumerate}
\item \textit{[JKSSB-pattern]} \textbf{(MSQ: select ALL correct options.)} Which are valid key types in databases?
  \begin{enumerate}[label=\Alph*.]
    \item Primary key
    \item Foreign key
    \item Candidate key
    \item Door key
  \end{enumerate}
\item \textit{[SSC CGL-pattern]} \textbf{(MSQ: select ALL correct options.)} Which statements about keys are correct?
  \begin{enumerate}[label=\Alph*.]
    \item A primary key uniquely identifies each record.
    \item A foreign key references another table's key.
    \item Duplicate primary-key values are not allowed.
    \item Every table must contain a printer serial number.
  \end{enumerate}
\item \textit{[Original]} A primary key should identify each record in a table:
  \begin{enumerate}[label=\Alph*.]
\item Uniquely
   \item With a duplicate value in every row
   \item Only by display color
   \item Using a blank value
  \end{enumerate}

\item \textit{[Original]} A foreign key is used primarily to:
  \begin{enumerate}[label=\Alph*.]
\item Reference a related record in another table
   \item Encrypt the database file
   \item Sort a report alphabetically
   \item Store an image only
  \end{enumerate}

\item \textit{[Original]} A key made from two or more fields is called a:
  \begin{enumerate}[label=\Alph*.]
\item Composite key
   \item Foreign form
   \item Query key
   \item Null key
  \end{enumerate}

\item \textit{[Original]} Referential integrity helps ensure that:
  \begin{enumerate}[label=\Alph*.]
\item Relationships between tables remain valid
   \item Every field is printed in bold
   \item Passwords are stored as plain text
   \item A query has no criteria
  \end{enumerate}

\item \textit{[Original, MSQ drill]} \textbf{(MSQ: select ALL correct options.)} Which statements about keys are correct?
  \begin{enumerate}[label=\Alph*.]
\item A primary key uniquely identifies a row.
   \item A foreign key can reference a primary key in another table.
   \item A composite key contains multiple fields.
   \item A primary key should contain duplicate values.
  \end{enumerate}

\end{enumerate}
\lessonheading{30}{Access Queries, Forms and Reports}
\begin{enumerate}[label=\textbf{Q\arabic*.}, leftmargin=*]
\item \textit{[Original]} The abbreviation \textbf{SQL} stands for:
  \begin{enumerate}[label=\Alph*.]
    \item Standard Query Language
    \item Structured Query Language
    \item Simple Query Language
    \item System Query Language
  \end{enumerate}
\item \textit{[Original]} Which MS Access object stores the actual data in
  rows and columns?
  \begin{enumerate}[label=\Alph*.]
    \item Query
    \item Form
    \item Table
    \item Report
  \end{enumerate}
\item \textit{[Original]} Which MS Access object is used to ask a question
  and retrieve selected data from the tables?
  \begin{enumerate}[label=\Alph*.]
    \item Form
    \item Query
    \item Report
    \item Table
  \end{enumerate}
\item \textit{[PYQ WO 2026 Q79, key:provisional]} A query that retrieves
  data on the basis of a condition is called a:
  \begin{enumerate}[label=\Alph*.]
    \item Delete query
    \item Update query
    \item Append query
    \item Select query
  \end{enumerate}
\item \textit{[Original]} Which type of query changes the existing values in
  records?
  \begin{enumerate}[label=\Alph*.]
    \item Select query
    \item Update query
    \item Append query
    \item Crosstab query
  \end{enumerate}
\item \textit{[Original]} Which type of query adds records from one table to
  another?
  \begin{enumerate}[label=\Alph*.]
    \item Update query
    \item Delete query
    \item Append query
    \item Select query
  \end{enumerate}
\item \textit{[Original]} Which type of query removes records that match a
  condition?
  \begin{enumerate}[label=\Alph*.]
    \item Delete query
    \item Select query
    \item Update query
    \item Append query
  \end{enumerate}
\item \textit{[Original]} Which of the following does \textbf{NOT} modify the
  stored data?
  \begin{enumerate}[label=\Alph*.]
    \item Update query
    \item Append query
    \item Delete query
    \item Select query
  \end{enumerate}
\item \textit{[Original]} Which MS Access object provides an on-screen
  interface for entering and viewing data one record at a time?
  \begin{enumerate}[label=\Alph*.]
    \item Table
    \item Query
    \item Form
    \item Report
  \end{enumerate}
\item \textit{[Original]} Which MS Access object presents data in a
  formatted, printable layout?
  \begin{enumerate}[label=\Alph*.]
    \item Form
    \item Report
    \item Query
    \item Table
  \end{enumerate}
\item \textit{[Original]} Which statement about an Access report is correct?
  \begin{enumerate}[label=\Alph*.]
    \item A report is used to enter new records.
    \item A report is read-only and is meant for printing.
    \item A report stores the actual data of the database.
    \item A report modifies records that match a condition.
  \end{enumerate}
\item \textit{[Original]} Which option gives the correct order of Access
  objects in which data is stored, selected, entered and printed?
  \begin{enumerate}[label=\Alph*.]
    \item Query $\rightarrow$ Table $\rightarrow$ Report $\rightarrow$ Form
    \item Table $\rightarrow$ Form $\rightarrow$ Query $\rightarrow$ Report
    \item Table $\rightarrow$ Query $\rightarrow$ Form $\rightarrow$ Report
    \item Form $\rightarrow$ Table $\rightarrow$ Query $\rightarrow$ Report
  \end{enumerate}
\item \textit{[Original]} Forms and reports in Access are usually built on:
  \begin{enumerate}[label=\Alph*.]
    \item another form.
    \item a table or a query.
    \item a report only.
    \item the Navigation Pane.
  \end{enumerate}
\item \textit{[Original]} In Access, a lookup field is used to:
  \begin{enumerate}[label=\Alph*.]
    \item delete records that match a condition.
    \item print a formatted summary of the data.
    \item choose a value from a list supplied by a table or query.
    \item create a new table automatically.
  \end{enumerate}
\item \textit{[Original]} Which of the following is \textbf{NOT} an object in
  an MS Access database?
  \begin{enumerate}[label=\Alph*.]
    \item Table
    \item Query
    \item Report
    \item Spreadsheet
  \end{enumerate}
\item \textit{[Original]} Evaluate the statements:
  \begin{enumerate}[label=\Roman*.]
    \item A Select query modifies the stored data.
    \item An Update query changes existing records.
    \item A Form is used to print a formatted report.
  \end{enumerate}
  \begin{enumerate}[label=\Alph*.]
    \item I and II only
    \item II only
    \item I and III only
    \item I, II and III
  \end{enumerate}
\item \textit{[Original]} The grid used in Access Query Design view to build
  a query by choosing fields and criteria is called:
  \begin{enumerate}[label=\Alph*.]
    \item QBE (Query By Example)
    \item SQL
    \item QWERTY
    \item PDF
  \end{enumerate}
\item \textit{[Original]} Which SQL keyword is used to retrieve data from a
  table?
  \begin{enumerate}[label=\Alph*.]
    \item UPDATE
    \item SELECT
    \item DELETE
    \item INSERT
  \end{enumerate}
\item \textit{[Original]} In the SQL statement
  ``\texttt{SELECT Name FROM Students WHERE Marks > 60}'', the
  \texttt{WHERE} clause:
  \begin{enumerate}[label=\Alph*.]
    \item names the table to be searched.
    \item lists the columns to be displayed.
    \item gives the condition that the rows must satisfy.
    \item sorts the result in ascending order.
  \end{enumerate}
\item \textit{[Original]} Which pairing is correctly matched?
  \begin{enumerate}[label=\Alph*.]
    \item Table --- asks a question and retrieves data
    \item Query --- stores the actual data
    \item Form --- provides a screen for data entry
    \item Report --- modifies records that match a condition
  \end{enumerate}

\item \textit{[JKSSB FAA-pattern]} Which Access object retrieves selected data from tables?
  \begin{enumerate}[label=\Alph*.]
    \item Query
    \item Wallpaper
    \item Ringtone
    \item Screensaver
  \end{enumerate}
\item \textit{[JKPSI-pattern]} Which Access object provides a screen for easy data entry?
  \begin{enumerate}[label=\Alph*.]
    \item Form
    \item Printer tray
    \item Speaker cone
    \item Mouse pad
  \end{enumerate}
\item \textit{[SSC CGL-pattern]} SQL stands for:
  \begin{enumerate}[label=\Alph*.]
    \item Structured Query Language
    \item Simple Paper List
    \item Speaker Quality Level
    \item Screen Quiet Lamp
  \end{enumerate}
\item \textit{[JKSSB-pattern]} \textbf{(MSQ: select ALL correct options.)} Which are Access database objects?
  \begin{enumerate}[label=\Alph*.]
    \item Tables
    \item Queries
    \item Forms
    \item Cricket bats
  \end{enumerate}
\item \textit{[SSC CGL-pattern]} \textbf{(MSQ: select ALL correct options.)} Which statements about Access objects are correct?
  \begin{enumerate}[label=\Alph*.]
    \item Queries fetch filtered data.
    \item Forms simplify data entry.
    \item Reports present data for printing.
    \item Queries physically print paper.
  \end{enumerate}
\item \textit{[Original]} Which Access object retrieves selected records based on criteria?
  \begin{enumerate}[label=\Alph*.]
\item Query
   \item Form
   \item Report
   \item Macro
  \end{enumerate}

\item \textit{[Original]} Which object is primarily designed for a formatted, printable summary of data?
  \begin{enumerate}[label=\Alph*.]
\item Report
   \item Table
   \item Relationship
   \item Module
  \end{enumerate}

\item \textit{[Original]} Which Access object is commonly used as a user-friendly screen for entering records?
  \begin{enumerate}[label=\Alph*.]
\item Form
   \item Report
   \item Index
   \item Relationship
  \end{enumerate}

\item \textit{[Original]} A SELECT query is generally used to:
  \begin{enumerate}[label=\Alph*.]
\item Retrieve records meeting conditions
   \item Delete the database application
   \item Change the computer's BIOS
   \item Print a PowerPoint slide
  \end{enumerate}

\item \textit{[Original, MSQ drill]} \textbf{(MSQ: select ALL correct options.)} Match each Access object to its typical role.
  \begin{enumerate}[label=\Alph*.]
\item Table --- store records
   \item Query --- retrieve selected data
   \item Form --- data entry/view interface
   \item Report --- formatted presentation of data
  \end{enumerate}

\end{enumerate}
\lessonheading{31}{PowerPoint Basics}
\begin{enumerate}[label=\textbf{Q\arabic*.}, leftmargin=*]
\item \textit{[Original]} Expand PPTX.
  \begin{enumerate}[label=\Alph*.]
    \item PowerPoint Text Presentation
    \item PowerPoint XML Presentation
    \item Portable Point Transfer XML
    \item Present Text XML
  \end{enumerate}
\item \textit{[Original]} The default file extension of a PowerPoint
  presentation is:
  \begin{enumerate}[label=\Alph*.]
    \item \texttt{.ppsx}
    \item \texttt{.potx}
    \item \texttt{.pptx}
    \item \texttt{.pptm}
  \end{enumerate}
\item \textit{[Original]} A single page of a PowerPoint presentation is
  called a:
  \begin{enumerate}[label=\Alph*.]
    \item worksheet
    \item slide
    \item document
    \item record
  \end{enumerate}
\item \textit{[Original]} The dotted-line box on a slide that holds text, a
  picture, a chart, or a table is called a:
  \begin{enumerate}[label=\Alph*.]
    \item placeholder
    \item cell
    \item field
    \item frame border
  \end{enumerate}
\item \textit{[Original]} The arrangement of placeholders on a slide is
  called the slide:
  \begin{enumerate}[label=\Alph*.]
    \item theme
    \item layout
    \item master
    \item template
  \end{enumerate}
\item \textit{[Original]} A coordinated set of colours, fonts, and effects
  applied from the Design tab is called a:
  \begin{enumerate}[label=\Alph*.]
    \item theme
    \item placeholder
    \item macro
    \item handout
  \end{enumerate}
\item \textit{[Original]} Which of the following is \textbf{NOT} true of the
  Slide Master?
  \begin{enumerate}[label=\Alph*.]
    \item It controls the design of all linked layouts and slides.
    \item A change to its font or logo applies to every linked slide.
    \item It is opened from the View tab through Slide Master.
    \item It prints audience handouts with six slides per page.
  \end{enumerate}
\item \textit{[Original]} Which master controls only the printed handout
  layouts?
  \begin{enumerate}[label=\Alph*.]
    \item Slide Master
    \item Handout Master
    \item Title Master
    \item Animation Master
  \end{enumerate}
\item \textit{[Original]} Which view is the default editing view of
  PowerPoint?
  \begin{enumerate}[label=\Alph*.]
    \item Slide Sorter
    \item Reading View
    \item Normal view
    \item Slide Show view
  \end{enumerate}
\item \textit{[Original]} In which view are thumbnails of all slides shown
  together so that slides can be reordered by drag and drop?
  \begin{enumerate}[label=\Alph*.]
    \item Normal view
    \item Notes Page view
    \item Slide Sorter view
    \item Outline View
  \end{enumerate}
\item \textit{[Original]} Speaker notes are best described as:
  \begin{enumerate}[label=\Alph*.]
    \item printed sheets of six slides given to the audience.
    \item the speaker's reminders below a slide, unseen by the audience.
    \item the movement effect between two slides.
    \item the top-level slide controlling all layouts.
  \end{enumerate}
\item \textit{[Original]} Printed sheets carrying several slides per page for
  the audience are called:
  \begin{enumerate}[label=\Alph*.]
    \item notes
    \item handouts
    \item placeholders
    \item transitions
  \end{enumerate}
\item \textit{[Original]} Which key starts the slide show from the beginning?
  \begin{enumerate}[label=\Alph*.]
    \item F1
    \item F5
    \item F7
    \item F9
  \end{enumerate}
\item \textit{[Original]} Which key combination starts the slide show from
  the current slide?
  \begin{enumerate}[label=\Alph*.]
    \item Shift+F5
    \item Ctrl+F5
    \item Alt+F5
    \item Ctrl+Shift+F5
  \end{enumerate}
\item \textit{[Original]} The movement effect applied \textit{between} two
  slides is called a \underline{\hspace{2.2cm}}, whereas the movement of an
  object \textit{on} a slide is called \underline{\hspace{2.2cm}}.
  \begin{enumerate}[label=\Alph*.]
    \item animation; transition
    \item transition; animation
    \item layout; theme
    \item master; handout
  \end{enumerate}
\item \textit{[Original]} Which pairing is correctly matched?
  \begin{enumerate}[label=\Alph*.]
    \item \texttt{.ppsx} --- editable working presentation
    \item \texttt{.potx} --- PowerPoint template
    \item \texttt{.pptx} --- file that opens directly as a slide show
    \item \texttt{.pptm} --- plain-text outline export
  \end{enumerate}
\item \textit{[Original]} Which layout carries no placeholders, so every
  object on it must be inserted by hand?
  \begin{enumerate}[label=\Alph*.]
    \item Title and Content
    \item Two Content
    \item Comparison
    \item Blank
  \end{enumerate}
\item \textit{[Original]} Which statement about handouts is
  \textbf{incorrect}?
  \begin{enumerate}[label=\Alph*.]
    \item Handouts can carry 1, 2, 3, 4, 6, or 9 slides per page.
    \item The 3-per-page layout adds ruled lines for the audience to write on.
    \item Six slides per page is a common exam-tested choice.
    \item Handouts are the speaker's hidden reminders shown in Presenter View.
  \end{enumerate}
\item \textit{[Original]} Evaluate the statements:
  \begin{enumerate}[label=\Roman*.]
    \item Normal view is the default editing view of PowerPoint.
    \item Slide Sorter view shows thumbnails of all slides for reordering.
    \item Reading View is the top-level slide controlling all layouts.
  \end{enumerate}
  \begin{enumerate}[label=\Alph*.]
    \item I only
    \item I and II only
    \item II and III only
    \item I, II, and III
  \end{enumerate}
\item \textit{[Original]} Match the shortcut or key with its slide-show
  action: 1.~F5, 2.~Shift+F5, 3.~Esc, 4.~B.
  \begin{enumerate}[label=\Alph*.]
    \item 1--from beginning, 2--from current slide, 3--end show, 4--black screen
    \item 1--from current slide, 2--from beginning, 3--black screen, 4--end show
    \item 1--end show, 2--black screen, 3--from beginning, 4--from current slide
    \item 1--black screen, 2--end show, 3--from current slide, 4--from beginning
  \end{enumerate}

\item \textit{[JKSSB FAA-pattern]} Each page of a PowerPoint presentation is called a:
  \begin{enumerate}[label=\Alph*.]
    \item slide.
    \item printer.
    \item speaker.
    \item cable.
  \end{enumerate}
\item \textit{[JKPSI-pattern]} Which view shows all slides as thumbnails for reordering?
  \begin{enumerate}[label=\Alph*.]
    \item Slide Sorter view
    \item Printer queue view
    \item Volume mixer view
    \item Calculator view
  \end{enumerate}
\item \textit{[SSC CGL-pattern]} The default PowerPoint file extension is:
  \begin{enumerate}[label=\Alph*.]
    \item .pptx
    \item .docx
    \item .xlsx
    \item .mp3x
  \end{enumerate}
\item \textit{[JKSSB-pattern]} \textbf{(MSQ: select ALL correct options.)} Which items can a PowerPoint slide contain?
  \begin{enumerate}[label=\Alph*.]
    \item Text boxes
    \item Images
    \item Charts
    \item Engine oil
  \end{enumerate}
\item \textit{[SSC CGL-pattern]} \textbf{(MSQ: select ALL correct options.)} Which statements about PowerPoint are correct?
  \begin{enumerate}[label=\Alph*.]
    \item Slides can hold text, images, and charts.
    \item Slide Sorter view helps reorder slides.
    \item F5 starts the show from the beginning.
    \item A slide is a printed paper register.
  \end{enumerate}
\item \textit{[Original]} A PowerPoint slide layout controls the:
  \begin{enumerate}[label=\Alph*.]
\item Arrangement of placeholders on a slide
   \item Order of files in a folder
   \item Worksheet cell address
   \item Email server protocol
  \end{enumerate}

\item \textit{[Original]} Which PowerPoint feature stores common formatting and repeated slide elements?
  \begin{enumerate}[label=\Alph*.]
\item Slide Master
   \item Formula bar
   \item Navigation Pane
   \item Mail Merge
  \end{enumerate}

\item \textit{[Original]} A placeholder on a slide is intended to:
  \begin{enumerate}[label=\Alph*.]
\item Hold content such as text, pictures, or charts
   \item Store an operating system
   \item Save an email password
   \item Execute a database query
  \end{enumerate}

\item \textit{[Original]} Which view is designed to rearrange slide thumbnails quickly?
  \begin{enumerate}[label=\Alph*.]
\item Slide Sorter
   \item Notes Page
   \item Reading View
   \item Outline view
  \end{enumerate}

\item \textit{[Original, MSQ drill]} \textbf{(MSQ: select ALL correct options.)} Which are common PowerPoint views?
  \begin{enumerate}[label=\Alph*.]
\item Normal view
   \item Slide Sorter view
   \item Notes Page view
   \item Slide Master view
  \end{enumerate}

\end{enumerate}
\lessonheading{32}{PowerPoint Animation and Transition}
\begin{enumerate}[label=\textbf{Q\arabic*.}, leftmargin=*]
\item \textit{[Original]} In PowerPoint, an animation is applied to:
  \begin{enumerate}[label=\Alph*.]
    \item the whole presentation at once
    \item an object or text on a slide
    \item the change from one slide to the next
    \item the file name of the presentation
  \end{enumerate}
\item \textit{[Original]} In PowerPoint, a transition is applied to:
  \begin{enumerate}[label=\Alph*.]
    \item a single object on a slide
    \item the movement from one slide to the next
    \item the notes pane
    \item the spelling and grammar checker
  \end{enumerate}
\item \textit{[PYQ WO 2026 Q80, key:provisional]} Which key starts the slide
  show \textbf{from the beginning}?
  \begin{enumerate}[label=\Alph*.]
    \item F5
    \item F7
    \item F9
    \item Shift+F5
  \end{enumerate}
\item \textit{[Original]} Which shortcut starts the slide show \textbf{from the
  current slide}?
  \begin{enumerate}[label=\Alph*.]
    \item F5
    \item Ctrl+F5
    \item Shift+F5
    \item Alt+F5
  \end{enumerate}
\item \textit{[Original]} In PowerPoint, \textbf{Alt+F5} is commonly used to:
  \begin{enumerate}[label=\Alph*.]
    \item end the slide show
    \item start the presentation in Presenter View
    \item spell-check the current slide
    \item open the Notes pane
  \end{enumerate}
\item \textit{[Original]} Which key \textbf{ends} a running slide show?
  \begin{enumerate}[label=\Alph*.]
    \item F5
    \item Esc
    \item Enter
    \item Tab
  \end{enumerate}
\item \textit{[Original]} Which of the following is \textbf{NOT} an animation
  effect category in PowerPoint?
  \begin{enumerate}[label=\Alph*.]
    \item Entrance
    \item Exit
    \item Transition
    \item Emphasis
  \end{enumerate}
\item \textit{[Original]} Which animation effect moves an object \textbf{along a
  defined path} on the slide?
  \begin{enumerate}[label=\Alph*.]
    \item Entrance
    \item Exit
    \item Emphasis
    \item Motion Path
  \end{enumerate}
\item \textit{[Original]} Which animation effect is used to make an object
  \textbf{appear} on the slide?
  \begin{enumerate}[label=\Alph*.]
    \item Entrance
    \item Exit
    \item Emphasis
    \item Motion Path
  \end{enumerate}
\item \textit{[Original]} An effect that draws attention to an object already
  visible on the slide, \textbf{without removing it}, is called:
  \begin{enumerate}[label=\Alph*.]
    \item Entrance
    \item Exit
    \item Emphasis
    \item Transition
  \end{enumerate}
\item \textit{[Original]} Which transition smoothly moves and animates matching
  objects as one slide changes into the next?
  \begin{enumerate}[label=\Alph*.]
    \item Fade
    \item Morph
    \item Push
    \item Wipe
  \end{enumerate}
\item \textit{[Original]} The Notes pane in PowerPoint is mainly used to:
  \begin{enumerate}[label=\Alph*.]
    \item print several slides on one page
    \item store speaker notes for the presenter
    \item animate objects on the slide
    \item change the theme colours
  \end{enumerate}
\item \textit{[Original]} A handout in PowerPoint is intended to:
  \begin{enumerate}[label=\Alph*.]
    \item control the slide transitions
    \item give the audience printed slides, often several per page
    \item start the slide show from the first slide
    \item edit the Slide Master
  \end{enumerate}
\item \textit{[Original]} Which view displays small thumbnails of all slides so
  they can be reordered?
  \begin{enumerate}[label=\Alph*.]
    \item Slide Sorter view
    \item Reading View
    \item Notes Page view
    \item Slide Show view
  \end{enumerate}
\item \textit{[Original]} Which of the following is \textbf{NOT} a standard view
  in PowerPoint?
  \begin{enumerate}[label=\Alph*.]
    \item Normal view
    \item Slide Sorter view
    \item Reading View
    \item Formula View
  \end{enumerate}
\item \textit{[Original]} Evaluate the statements:
  \begin{enumerate}[label=\Roman*.]
    \item An animation is applied to objects on a slide.
    \item A transition is applied between slides.
    \item Shift+F5 starts the slide show from the current slide.
  \end{enumerate}
  \begin{enumerate}[label=\Alph*.]
    \item I and II only
    \item II and III only
    \item I and III only
    \item I, II, and III
  \end{enumerate}
\item \textit{[PYQ-adapted]} When a SmartArt graphic is converted to text, what
  is lost and must be re-applied or re-entered?
  \begin{enumerate}[label=\Alph*.]
    \item nothing; the design is preserved exactly
    \item the graphical layout --- it becomes a plain bulleted list
    \item the file name of the presentation
    \item the slide transitions
  \end{enumerate}
\item \textit{[Original]} During a running slide show, pressing the
\textbf{B} key will:
  \begin{enumerate}[label=\Alph*.]
    \item end the slide show
    \item turn the screen black
    \item go back to the previous slide
    \item open Presenter View
  \end{enumerate}
\item \textit{[Original]} Which key is correctly matched with its PowerPoint
  function?
  \begin{enumerate}[label=\Alph*.]
    \item F5 --- spelling and grammar check
    \item F7 --- start the slide show
    \item Esc --- end the slide show
    \item Alt+F5 --- end the slide show
  \end{enumerate}
\item \textit{[Original]} Which pairing of a PowerPoint feature with its function
  is correct?
  \begin{enumerate}[label=\Alph*.]
    \item Morph --- a transition between similar slides
    \item Notes pane --- an animation effect
    \item Motion Path --- a transition between slides
    \item Handout --- starts the show from the current slide
  \end{enumerate}

\item \textit{[JKSSB FAA-pattern]} Slide transitions control movement between:
  \begin{enumerate}[label=\Alph*.]
    \item one slide and the next.
    \item the printer and the paper.
    \item the speaker and the ear.
    \item the mouse and the pad.
  \end{enumerate}
\item \textit{[JKPSI-pattern]} Custom Animation is applied to:
  \begin{enumerate}[label=\Alph*.]
    \item objects within a slide.
    \item the power cable.
    \item the printer ribbon.
    \item the monitor stand.
  \end{enumerate}
\item \textit{[SSC CGL-pattern]} Which key starts the slide show from the current slide?
  \begin{enumerate}[label=\Alph*.]
    \item Shift+F5
    \item F7
    \item F2
    \item Ctrl+S
  \end{enumerate}
\item \textit{[JKSSB-pattern]} \textbf{(MSQ: select ALL correct options.)} Which are PowerPoint animation/transition effects?
  \begin{enumerate}[label=\Alph*.]
    \item Fade transition
    \item Fly In animation
    \item Wipe transition
    \item Paper jam
  \end{enumerate}
\item \textit{[SSC CGL-pattern]} \textbf{(MSQ: select ALL correct options.)} Which statements about animation are correct?
  \begin{enumerate}[label=\Alph*.]
    \item Transitions act between slides.
    \item Animations act on objects inside slides.
    \item Effect timing can be adjusted.
    \item Animations print extra paper copies.
  \end{enumerate}
\item \textit{[Original]} A transition effect occurs when:
  \begin{enumerate}[label=\Alph*.]
\item Moving from one slide to another
   \item A word is selected in a text box
   \item A spreadsheet is recalculated
   \item A file is renamed
  \end{enumerate}

\item \textit{[PYQ-adapted: JKSSB Website Operator 2026 Q80, provisional key; SRC-04/SRC-07]} Which key starts a PowerPoint slide show from the beginning?
  \begin{enumerate}[label=\Alph*.]
\item F5
   \item Shift+F5
   \item F7
   \item Ctrl+S
  \end{enumerate}

\item \textit{[Original]} Which shortcut commonly starts a slide show from the current slide?
  \begin{enumerate}[label=\Alph*.]
\item Shift+F5
   \item F5
   \item Ctrl+P
   \item Alt+F4
  \end{enumerate}

\item \textit{[Original]} An animation effect is applied to:
  \begin{enumerate}[label=\Alph*.]
\item An object or content within a slide
   \item The change between two files in a folder
   \item A worksheet's row height only
   \item The computer boot process
  \end{enumerate}

\item \textit{[Original, MSQ drill]} \textbf{(MSQ: select ALL correct options.)} Which statements are correct?
  \begin{enumerate}[label=\Alph*.]
\item F5 starts the slide show from the beginning.
   \item Shift+F5 starts from the current slide.
   \item A transition is between slides.
   \item An animation is always the same as a transition.
  \end{enumerate}

\end{enumerate}
\lessonheading{33}{PDF and File Formats}
\begin{enumerate}[label=\textbf{Q\arabic*.}, leftmargin=*]
\item \textit{[Original]} Expand PDF.
  \begin{enumerate}[label=\Alph*.]
    \item Portable Document Format
    \item Personal Document File
    \item Portable Data File
    \item Printed Document Format
  \end{enumerate}
\item \textit{[Original]} Which file format stores a document in a fixed
  layout so that it looks the same on every device and printer?
  \begin{enumerate}[label=\Alph*.]
    \item CSV
    \item PDF
    \item TXT
    \item DOCX
  \end{enumerate}
\item \textit{[Original]} In the file name \texttt{map.png}, the ``png'' part
  is the \underline{\hspace{1.8cm}} and PNG is the \underline{\hspace{1.8cm}}.
  \begin{enumerate}[label=\Alph*.]
    \item format; extension
    \item extension; format
    \item folder; drive
    \item format; version
  \end{enumerate}
\item \textit{[Original]} In a scanned (image-only) PDF, the text on each
  page is:
  \begin{enumerate}[label=\Alph*.]
    \item stored as an image, so it cannot be directly selected or searched
    \item stored as editable text, so it can be edited freely
    \item stored as a spreadsheet, so it can be sorted
    \item stored as a vector font, so it can be resized
  \end{enumerate}
\item \textit{[Original]} Expand OCR.
  \begin{enumerate}[label=\Alph*.]
    \item Optical Code Reader
    \item Optical Character Recognition
    \item Online Content Retrieval
    \item Output Character Rendering
  \end{enumerate}
\item \textit{[Original]} OCR is used to:
  \begin{enumerate}[label=\Alph*.]
    \item compress an image without losing quality
    \item convert an image of text into machine-readable, editable text
    \item protect a document with a password
    \item convert a spreadsheet into a slide show
  \end{enumerate}
\item \textit{[Original]} Which file format is the default word-processing
  format of Microsoft Word?
  \begin{enumerate}[label=\Alph*.]
    \item XLSX
    \item PPTX
    \item DOCX
    \item CSV
  \end{enumerate}
\item \textit{[Original]} Which extension belongs to a Microsoft Excel
  workbook?
  \begin{enumerate}[label=\Alph*.]
    \item \texttt{.xlsx}
    \item \texttt{.pptx}
    \item \texttt{.docx}
    \item \texttt{.png}
  \end{enumerate}
\item \textit{[Original]} Expand CSV.
  \begin{enumerate}[label=\Alph*.]
    \item Computer System Values
    \item Comma-Separated Values
    \item Common Spreadsheet Version
    \item Column-Sorted Variables
  \end{enumerate}
\item \textit{[Original]} Which of the following is a \textbf{plain-text}
  file format?
  \begin{enumerate}[label=\Alph*.]
    \item \texttt{.docx}
    \item \texttt{.pptx}
    \item \texttt{.txt}
    \item \texttt{.pdf}
  \end{enumerate}
\item \textit{[Original]} Which of the following is \textbf{NOT} an image
  file extension?
  \begin{enumerate}[label=\Alph*.]
    \item \texttt{.png}
    \item \texttt{.jpg}
    \item \texttt{.gif}
    \item \texttt{.docx}
  \end{enumerate}
\item \textit{[Original]} Protecting a PDF with an encryption password mainly
  provides:
  \begin{enumerate}[label=\Alph*.]
    \item confidentiality, so only an authorised reader can see the content
    \item authenticity, so the signer is verified
    \item integrity, so changes are detected
    \item compression, so the file becomes smaller
  \end{enumerate}
\item \textit{[Original]} A digital signature on a PDF mainly provides:
  \begin{enumerate}[label=\Alph*.]
    \item confidentiality of the content
    \item authenticity and integrity
    \item lossless image compression
    \item automatic OCR of the pages
  \end{enumerate}
\item \textit{[Original]} Which statement about a digital signature is
  correct?
  \begin{enumerate}[label=\Alph*.]
    \item It encrypts the file so only the signer can open it.
    \item It verifies the signer and detects any change made after signing.
    \item It converts a scanned PDF into a text PDF.
    \item It permanently removes sensitive content from the file.
  \end{enumerate}
\item \textit{[Original]} Evaluate the statements:
  \begin{enumerate}[label=\Roman*.]
    \item Encryption provides confidentiality.
    \item A digital signature provides authenticity and integrity.
    \item A digital signature encrypts the file so that no one else can read
      it.
  \end{enumerate}
  \begin{enumerate}[label=\Alph*.]
    \item I and II only
    \item I and III only
    \item II and III only
    \item I, II, and III
  \end{enumerate}
\item \textit{[Original]} Redaction of a PDF means:
  \begin{enumerate}[label=\Alph*.]
    \item compressing the file to reduce its size
    \item permanently removing sensitive content before sharing the file
    \item adding a digital signature to the file
    \item converting the file to PDF/A for archiving
  \end{enumerate}
\item \textit{[Original]} PDF/A is designed mainly for:
  \begin{enumerate}[label=\Alph*.]
    \item editing the text of a document
    \item the long-term archiving of documents
    \item scanning printed pages into images
    \item compressing photographs
  \end{enumerate}
\item \textit{[Original]} Which image format commonly uses lossless
  compression and supports transparency?
  \begin{enumerate}[label=\Alph*.]
    \item PNG
    \item JPEG
    \item BMP
    \item TXT
  \end{enumerate}
\item \textit{[Original]} Which pairing is correctly matched?
  \begin{enumerate}[label=\Alph*.]
    \item XLSX --- a Microsoft PowerPoint presentation
    \item DOCX --- a Microsoft Excel workbook
    \item PPTX --- a Microsoft PowerPoint presentation
    \item CSV --- a fixed-layout document format
  \end{enumerate}
\item \textit{[Original]} Which statement best distinguishes a file extension
  from a file format?
  \begin{enumerate}[label=\Alph*.]
    \item The extension is the suffix in the file name, while the format is
      the internal structure of the data.
    \item The extension is the internal structure of the data, while the
      format is the suffix in the file name.
    \item The extension and the format always mean exactly the same thing.
    \item The extension is always decided by the folder in which the file is
      stored.
  \end{enumerate}

\item \textit{[JKSSB FAA-pattern]} PDF stands for:
  \begin{enumerate}[label=\Alph*.]
    \item Portable Document Format
    \item Printer Disk File
    \item Personal Data Folder
    \item Public Display Frame
  \end{enumerate}
\item \textit{[JKPSI-pattern]} Which program is designed to read PDF files?
  \begin{enumerate}[label=\Alph*.]
    \item Adobe Reader/Acrobat
    \item Calculator
    \item Paint sprayer
    \item Door alarm
  \end{enumerate}
\item \textit{[SSC CGL-pattern]} Which format keeps document layout fixed across devices?
  \begin{enumerate}[label=\Alph*.]
    \item PDF
    \item Unformatted plain chatter
    \item Printer beep code
    \item Speaker buzz track
  \end{enumerate}
\item \textit{[JKSSB-pattern]} \textbf{(MSQ: select ALL correct options.)} Which are common document/file formats?
  \begin{enumerate}[label=\Alph*.]
    \item .pdf
    \item .docx
    \item .txt
    \item .fryingpan
  \end{enumerate}
\item \textit{[SSC CGL-pattern]} \textbf{(MSQ: select ALL correct options.)} Which statements about PDF and formats are correct?
  \begin{enumerate}[label=\Alph*.]
    \item PDFs preserve layout across devices.
    \item PDFs can hold text, images, and links.
    \item .docx is an editable Word format.
    \item PDFs can only be opened on printers.
  \end{enumerate}
\item \textit{[Original]} Which file extension is commonly associated with a modern Word document?
  \begin{enumerate}[label=\Alph*.]
\item .docx
   \item .xlsx
   \item .pptx
   \item .accdb
  \end{enumerate}

\item \textit{[Original]} A PDF is primarily designed to preserve:
  \begin{enumerate}[label=\Alph*.]
\item Document layout across devices
   \item Spreadsheet formulas as the only content
   \item Executable program instructions
   \item Audio channels only
  \end{enumerate}

\item \textit{[Original]} Which file format is commonly used for plain comma-separated tabular data?
  \begin{enumerate}[label=\Alph*.]
\item CSV
   \item PPTX
   \item PNG
   \item EXE
  \end{enumerate}

\item \textit{[Original]} Which format is an image format?
  \begin{enumerate}[label=\Alph*.]
\item PNG
   \item DOCX
   \item XLSX
   \item PPTX
  \end{enumerate}

\item \textit{[Original, MSQ drill]} \textbf{(MSQ: select ALL correct options.)} Which extension-to-file-type pairings are correct?
  \begin{enumerate}[label=\Alph*.]
\item .docx --- word-processing document
   \item .xlsx --- spreadsheet workbook
   \item .pptx --- presentation
   \item .png --- database file
  \end{enumerate}

\end{enumerate}
\lessonheading{54}{Office Shortcut and Product Trap Lab}
\begin{enumerate}[label=\textbf{Q\arabic*.}, leftmargin=*]
\item \textit{[Original]} In MS Word, the key \textbf{F7} is used for:
  \begin{enumerate}[label=\Alph*.]
    \item Help
    \item Spelling \& Grammar (spell check)
    \item the Thesaurus
    \item Find and Replace
  \end{enumerate}
\item \textit{[Original]} In MS Word, \textbf{Ctrl+H} opens:
  \begin{enumerate}[label=\Alph*.]
    \item Find only
    \item Find and Replace
    \item Insert Hyperlink
    \item Spelling \& Grammar
  \end{enumerate}
\item \textit{[Original]} Which shortcut inserts a hyperlink in MS Word?
  \begin{enumerate}[label=\Alph*.]
    \item Ctrl+H
    \item Ctrl+K
    \item Ctrl+L
    \item Ctrl+J
  \end{enumerate}
\item \textit{[Original]} In MS Word, \textbf{Ctrl+X} followed by
  \textbf{Ctrl+V} performs:
  \begin{enumerate}[label=\Alph*.]
    \item copy and paste
    \item cut and paste
    \item undo and redo
    \item find and replace
  \end{enumerate}
\item \textit{[Original]} Which MS Word view shows the document exactly as it
  will appear on the printed page?
  \begin{enumerate}[label=\Alph*.]
    \item Web Layout
    \item Outline
    \item Print Layout
    \item Draft
  \end{enumerate}
\item \textit{[Original]} Track Changes in MS Word is used to:
  \begin{enumerate}[label=\Alph*.]
    \item change the page orientation of the document
    \item record insertions, deletions, and formatting edits for later
      accept or reject
    \item check spelling and grammar automatically
    \item compress the pictures in the document
  \end{enumerate}
\item \textit{[Original]} Which statement about Track Changes in MS Word is
  \textbf{NOT} true?
  \begin{enumerate}[label=\Alph*.]
    \item It records insertions and deletions.
    \item Reviewers can accept or reject each change.
    \item It stores the author's name and the time of the edit as metadata.
    \item It permanently deletes the original text so that it can never be
      restored.
  \end{enumerate}
\item \textit{[Original]} Which type of break is required so that the pages
  after it can have a \textbf{different header or footer} in MS Word?
  \begin{enumerate}[label=\Alph*.]
    \item Page Break
    \item Section Break
    \item Column Break
    \item Line Break
  \end{enumerate}
\item \textit{[Original]} A Word feature that combines a main document with
  an address database to produce many personalised letters is called:
  \begin{enumerate}[label=\Alph*.]
    \item Track Changes
    \item AutoCorrect
    \item Mail Merge
    \item Find and Replace
  \end{enumerate}
\item \textit{[Original]} In MS Word, which shortcut removes a document window
  split?
  \begin{enumerate}[label=\Alph*.]
    \item Alt+Shift+C
    \item Ctrl+Shift+C
    \item Ctrl+Alt+C
    \item Alt+Ctrl+V
  \end{enumerate}
\item \textit{[Original]} Which key starts a PowerPoint slide show from the
  \textbf{beginning}?
  \begin{enumerate}[label=\Alph*.]
    \item F5
    \item F7
    \item F1
    \item F2
  \end{enumerate}
\item \textit{[Original]} In PowerPoint, \textbf{Shift+F5} is used to:
  \begin{enumerate}[label=\Alph*.]
    \item end the slide show
    \item start the slide show from the current slide
    \item start the presentation in Presenter View
    \item insert a new slide
  \end{enumerate}
\item \textit{[Original]} In PowerPoint, \textbf{Alt+F5} starts the
  presentation in:
  \begin{enumerate}[label=\Alph*.]
    \item Slide Sorter view
    \item Presenter View
    \item Reading view
    \item Slide Master view
  \end{enumerate}
\item \textit{[Original]} Which key ends a running PowerPoint slide show?
  \begin{enumerate}[label=\Alph*.]
    \item Esc
    \item Enter
    \item F5
    \item Ctrl+H
  \end{enumerate}
\item \textit{[Original]} In PowerPoint, which of the following existing slide
  contents \textbf{cannot} be directly converted into SmartArt?
  \begin{enumerate}[label=\Alph*.]
    \item a bulleted list
    \item a block of text
    \item a table of data
    \item indented text
  \end{enumerate}
\item \textit{[Original]} In PowerPoint, the \textbf{Handouts} print layout is
  used to:
  \begin{enumerate}[label=\Alph*.]
    \item edit the speaker notes of each slide
    \item print several slides on one page for the audience
    \item change the design of all slides at once
    \item switch the presentation between Portrait and Landscape
  \end{enumerate}
\item \textit{[Original]} In PowerPoint, the \textbf{Slide Master} controls:
  \begin{enumerate}[label=\Alph*.]
    \item only the title slide of the presentation
    \item the design, fonts, colours, and layout of all slides that use it
    \item the spell-check dictionary
    \item the printing settings only
  \end{enumerate}
\item \textit{[Original]} In PowerPoint, changing the \textbf{orientation}
  switches the pages between:
  \begin{enumerate}[label=\Alph*.]
    \item relative and absolute
    \item Portrait and Landscape
    \item Draft and Print Layout
    \item Normal and Slide Sorter
  \end{enumerate}
\item \textit{[Original]} In MS Excel, the reference \textbf{\$A\$1} is a:
  \begin{enumerate}[label=\Alph*.]
    \item relative reference
    \item absolute reference
    \item mixed reference with only the row fixed
    \item named range
  \end{enumerate}
\item \textit{[Original]} Evaluate the statements:
  \begin{enumerate}[label=\Roman*.]
    \item MS-DOS is an operating system.
    \item MS Access is a spreadsheet.
    \item MS Word is a word processor.
  \end{enumerate}
  \begin{enumerate}[label=\Alph*.]
    \item I only
    \item I and III only
    \item II and III only
    \item I, II, and III
  \end{enumerate}

\item \textit{[JKSSB FAA-pattern]} Which shortcut selects all content in a document?
  \begin{enumerate}[label=\Alph*.]
    \item Ctrl+A
    \item Ctrl+S
    \item Ctrl+P
    \item Ctrl+N
  \end{enumerate}
\item \textit{[JKPSI-pattern]} Which shortcut opens the Find dialog in most programs?
  \begin{enumerate}[label=\Alph*.]
    \item Ctrl+F
    \item Ctrl+B
    \item Ctrl+U
    \item Ctrl+E
  \end{enumerate}
\item \textit{[SSC CGL-pattern]} Which shortcut creates a new document?
  \begin{enumerate}[label=\Alph*.]
    \item Ctrl+N
    \item Ctrl+P
    \item Ctrl+S
    \item Ctrl+V
  \end{enumerate}
\item \textit{[JKSSB-pattern]} \textbf{(MSQ: select ALL correct options.)} Which shortcut pairs are correctly matched?
  \begin{enumerate}[label=\Alph*.]
    \item Ctrl+N -- new document
    \item Ctrl+O -- open file
    \item Ctrl+S -- save file
    \item Ctrl+X -- paste content
  \end{enumerate}
\item \textit{[SSC CGL-pattern]} \textbf{(MSQ: select ALL correct options.)} Which product matches are correct?
  \begin{enumerate}[label=\Alph*.]
    \item MS Word -- word processing
    \item MS Excel -- spreadsheets
    \item MS PowerPoint -- presentations
    \item MS Excel -- slide-show presentations
  \end{enumerate}
\item \textit{[Original]} Which shortcut commonly opens the Print dialog in Office applications?
  \begin{enumerate}[label=\Alph*.]
\item Ctrl+P
   \item Ctrl+H
   \item Ctrl+K
   \item Ctrl+Z
  \end{enumerate}

\item \textit{[Original]} Which Word shortcut commonly inserts a hyperlink?
  \begin{enumerate}[label=\Alph*.]
\item Ctrl+K
   \item Ctrl+L
   \item Ctrl+M
   \item Ctrl+W
  \end{enumerate}

\item \textit{[Original]} Which key starts a PowerPoint slide show from the current slide?
  \begin{enumerate}[label=\Alph*.]
\item Shift+F5
   \item F5
   \item F7
   \item Ctrl+P
  \end{enumerate}

\item \textit{[Original]} Which key commonly starts spelling and grammar review in Word?
  \begin{enumerate}[label=\Alph*.]
\item F7
   \item F5
   \item F1
   \item F12
  \end{enumerate}

\item \textit{[Original, MSQ drill]} \textbf{(MSQ: select ALL correct options.)} Match common shortcuts to actions.
  \begin{enumerate}[label=\Alph*.]
\item Ctrl+S --- save
   \item Ctrl+P --- print
   \item Ctrl+Z --- undo
   \item Ctrl+C --- paste
  \end{enumerate}

\end{enumerate}
\lessonheading{55}{File Formats and Named Technology Associations}
\begin{enumerate}[label=\textbf{Q\arabic*.}, leftmargin=*]
\item \textit{[Original]} Expand PNG.
  \begin{enumerate}[label=\Alph*.]
    \item Portable Network Graphics
    \item Personal Network Graphics
    \item Portable New Graphic
    \item Public Network Graphics
  \end{enumerate}
\item \textit{[Original]} PNG is best classified as:
  \begin{enumerate}[label=\Alph*.]
    \item an image format
    \item a document format
    \item a database format
    \item a spreadsheet format
  \end{enumerate}
\item \textit{[Original]} Expand PDF.
  \begin{enumerate}[label=\Alph*.]
    \item Personal Document File
    \item Portable Document Format
    \item Portable Data Format
    \item Printed Document Format
  \end{enumerate}
\item \textit{[Original]} PDF is chiefly known for:
  \begin{enumerate}[label=\Alph*.]
    \item keeping a fixed layout that looks the same on every device
    \item storing editable spreadsheet formulas
    \item holding a relational database with keys
    \item compressing audio with loss
  \end{enumerate}
\item \textit{[Original]} Which of the following is primarily a \textbf{PDF
  reader}?
  \begin{enumerate}[label=\Alph*.]
    \item Mozilla Firefox
    \item Adobe Acrobat Reader
    \item TeamViewer
    \item Microsoft Excel
  \end{enumerate}
\item \textit{[Original]} Firefox is developed by:
  \begin{enumerate}[label=\Alph*.]
    \item Microsoft
    \item Adobe
    \item Mozilla
    \item Google
  \end{enumerate}
\item \textit{[Original]} Firefox is a:
  \begin{enumerate}[label=\Alph*.]
    \item web browser
    \item PDF reader
    \item spreadsheet program
    \item database program
  \end{enumerate}
\item \textit{[Original]} PARAM is associated with:
  \begin{enumerate}[label=\Alph*.]
    \item ISRO
    \item C-DAC
    \item Microsoft
    \item Mozilla
  \end{enumerate}
\item \textit{[Original]} Expand C-DAC.
  \begin{enumerate}[label=\Alph*.]
    \item Centre for Data and Computing
    \item Centre for Development of Advanced Computing
    \item Central Department of Advanced Computers
    \item Committee for Digital and Computer Applications
  \end{enumerate}
\item \textit{[Original]} Which extension is used by a Microsoft Excel
  workbook?
  \begin{enumerate}[label=\Alph*.]
    \item \texttt{.docx}
    \item \texttt{.xlsx}
    \item \texttt{.pptx}
    \item \texttt{.accdb}
  \end{enumerate}
\item \textit{[Original]} Which extension denotes a Microsoft Word document?
  \begin{enumerate}[label=\Alph*.]
    \item \texttt{.docx}
    \item \texttt{.xlsx}
    \item \texttt{.pptx}
    \item \texttt{.txt}
  \end{enumerate}
\item \textit{[Original]} The extension \texttt{.accdb} is used by a database
  created in:
  \begin{enumerate}[label=\Alph*.]
    \item Microsoft Excel
    \item Microsoft Word
    \item Microsoft Access
    \item Microsoft PowerPoint
  \end{enumerate}
\item \textit{[Original]} CSV stands for:
  \begin{enumerate}[label=\Alph*.]
    \item Common System Values
    \item Comma-Separated Values
    \item Comma-Stored Variables
    \item Central Spreadsheet Values
  \end{enumerate}
\item \textit{[Original]} A CSV file stores data as:
  \begin{enumerate}[label=\Alph*.]
    \item a table in plain text, with values separated by commas
    \item a compressed image of a page
    \item a database with primary and foreign keys
    \item a formatted document with fonts and colours
  \end{enumerate}
\item \textit{[Original]} Which of the following is \textbf{NOT} an image
  file format?
  \begin{enumerate}[label=\Alph*.]
    \item PNG
    \item JPEG
    \item GIF
    \item CSV
  \end{enumerate}
\item \textit{[Original]} Evaluate the statements:
  \begin{enumerate}[label=\Roman*.]
    \item PNG is an image format.
    \item PDF preserves the layout of a document across devices.
    \item \texttt{.accdb} is the extension of a Microsoft Word file.
  \end{enumerate}
  \begin{enumerate}[label=\Alph*.]
    \item I only
    \item I and II only
    \item II and III only
    \item I, II, and III
  \end{enumerate}
\item \textit{[Original]} The extension \texttt{.mdb} denotes:
  \begin{enumerate}[label=\Alph*.]
    \item an older Microsoft Access database file
    \item a Microsoft Word document
    \item a PDF document
    \item a portable image file
  \end{enumerate}
\item \textit{[Original]} Match List-I with List-II:
  \begin{enumerate}[label=(\roman*)]
    \item PNG
    \item PARAM
    \item Firefox
    \item \texttt{.pptx}
  \end{enumerate}
  \begin{enumerate}[label=(\alph*)]
    \item Mozilla web browser
    \item Image format
    \item PowerPoint presentation
    \item C-DAC supercomputer
  \end{enumerate}
  Which option gives all correct matches?
  \begin{enumerate}[label=\Alph*.]
    \item (i)-b, (ii)-d, (iii)-a, (iv)-c
    \item (i)-a, (ii)-b, (iii)-c, (iv)-d
    \item (i)-d, (ii)-c, (iii)-b, (iv)-a
    \item (i)-b, (ii)-c, (iii)-a, (iv)-d
  \end{enumerate}
\item \textit{[Original]} Which of the following is a
  \textbf{remote-support} tool?
  \begin{enumerate}[label=\Alph*.]
    \item Adobe Acrobat Reader
    \item Mozilla Firefox
    \item TeamViewer
    \item Microsoft PowerPoint
  \end{enumerate}
\item \textit{[Original]} Which statement about \texttt{.png} and PNG is
  correct?
  \begin{enumerate}[label=\Alph*.]
    \item \texttt{.png} is the extension and PNG is the format.
    \item PNG is the extension and \texttt{.png} is the format.
    \item The two are the same thing, so either name can be used for both.
    \item \texttt{.png} is the format and PNG is the file name.
  \end{enumerate}

\item \textit{[JKSSB FAA-pattern]} Which extension belongs to a Word document?
  \begin{enumerate}[label=\Alph*.]
    \item .docx
    \item .xlsx
    \item .pptx
    \item .mp3
  \end{enumerate}
\item \textit{[JKPSI-pattern]} Which extension belongs to an Excel workbook?
  \begin{enumerate}[label=\Alph*.]
    \item .xlsx
    \item .docx
    \item .pptx
    \item .txt-paper
  \end{enumerate}
\item \textit{[SSC CGL-pattern]} Which format is best for web photographs with many colours?
  \begin{enumerate}[label=\Alph*.]
    \item JPEG
    \item Plain text
    \item Printer beep
    \item Speaker hum
  \end{enumerate}
\item \textit{[JKSSB-pattern]} \textbf{(MSQ: select ALL correct options.)} Which extensions denote image files?
  \begin{enumerate}[label=\Alph*.]
    \item .jpg
    \item .png
    \item .gif
    \item .docx
  \end{enumerate}
\item \textit{[SSC CGL-pattern]} \textbf{(MSQ: select ALL correct options.)} Which format pairs are correctly matched?
  \begin{enumerate}[label=\Alph*.]
    \item .mp3 -- audio
    \item .mp4 -- video
    \item .pdf -- fixed-layout document
    \item .xlsx -- printed poster paper
  \end{enumerate}
\item \textit{[PYQ-adapted: JKSSB FAA 2022, QCRIAS item 554; SRC-21]} Which listed device is a flash-memory device?
  \begin{enumerate}[label=\Alph*.]
\item Hard disk drive
   \item RAM
   \item USB drive
   \item Blu-ray disc
  \end{enumerate}

\item \textit{[Original]} Which extension is commonly associated with a PowerPoint presentation?
  \begin{enumerate}[label=\Alph*.]
\item .pptx
   \item .xlsx
   \item .docx
   \item .accdb
  \end{enumerate}

\item \textit{[Original]} Which statement correctly distinguishes a file extension from a file format?
  \begin{enumerate}[label=\Alph*.]
\item `.png` is an extension; PNG is the format name.
   \item PNG is an extension and `.png` is the format.
   \item They always indicate the file's actual contents with certainty.
   \item An extension is a storage device.
  \end{enumerate}

\item \textit{[Original]} Which file type is typically plain-text tabular data and does not preserve rich workbook formatting?
  \begin{enumerate}[label=\Alph*.]
\item CSV
   \item XLSX
   \item PPTX
   \item DOCX
  \end{enumerate}

\item \textit{[Original, MSQ drill]} \textbf{(MSQ: select ALL correct options.)} Which pairs are correctly matched?
  \begin{enumerate}[label=\Alph*.]
\item .docx --- Word document
   \item .xlsx --- Excel workbook
   \item .pptx --- PowerPoint presentation
   \item .png --- database file
  \end{enumerate}

\end{enumerate}
\lessonheading{61}{Current Office, Email and Statement Evaluation}
\begin{enumerate}[label=\textbf{Q\arabic*.}, leftmargin=*]
\item \textit{[Original]} In MS Word, content controls (drop-down list, date
  picker, check box) are inserted from which tab?
  \begin{enumerate}[label=\Alph*.]
    \item Developer tab
    \item Review tab
    \item Mailings tab
    \item References tab
  \end{enumerate}
\item \textit{[Original]} Which of the following statements about MS Word is
  \textbf{NOT} correct?
  \begin{enumerate}[label=\Alph*.]
    \item A content control can enforce limited data validation, such as a
      drop-down list or a date picker.
    \item A hanging indent shifts every line of the paragraph except the first
      line.
    \item A reviewer's comments on a sentence appear directly in the Navigation
      Pane.
    \item Accepting a tracked change keeps the new text and removes the
      revision metadata attached to that change.
  \end{enumerate}
\item \textit{[Original]} Which pairing of indent type and effect is correctly
  matched?
  \begin{enumerate}[label=\Alph*.]
    \item First-line indent --- indents every line except the first
    \item Hanging indent --- indents only the first line of the paragraph
    \item First-line indent --- indents only the first line of the paragraph
    \item Hanging indent --- centres the first line and left-aligns the rest
  \end{enumerate}
\item \textit{[Original]} The Word Navigation Pane offers which three views?
  \begin{enumerate}[label=\Alph*.]
    \item Headings, Pages, Results
    \item Comments, Footnotes, Endnotes
    \item Draft, Outline, Print Layout
    \item Insert, Design, References
  \end{enumerate}
\item \textit{[Original]} When Track Changes is switched on, each recorded
  change stores:
  \begin{enumerate}[label=\Alph*.]
    \item only the page number of the edit.
    \item the author's name and the date and time of the edit.
    \item only the final text with no author information.
    \item the printer name and the paper size.
  \end{enumerate}
\item \textit{[Original]} A report needs no header on the title page but a
  running header on every later page. Which break makes this possible?
  \begin{enumerate}[label=\Alph*.]
    \item Page Break
    \item Column Break
    \item Section Break
    \item Line Break
  \end{enumerate}
\item \textit{[Original]} Given
  \texttt{=IF(A1>50,"High",IF(A1>25,"Medium","Low"))} with A1 containing 30,
  the formula returns:
  \begin{enumerate}[label=\Alph*.]
    \item High
    \item Medium
    \item Low
    \item 0
  \end{enumerate}
\item \textit{[Original]} With the same formula
  \texttt{=IF(A1>50,"High",IF(A1>25,"Medium","Low"))}, if A1 is blank the
  result is:
  \begin{enumerate}[label=\Alph*.]
    \item High, because a blank cell is read as 100.
    \item Medium, because a blank cell skips the first test.
    \item Low, because a blank cell is treated as 0 in the comparison.
    \item \texttt{\#VALUE!}, because a blank cell always causes an error.
  \end{enumerate}
\item \textit{[Original]} Cell A1 contains the text ``abc''. Which statement
  is correct?
  \begin{enumerate}[label=\Alph*.]
    \item \texttt{=A1+5} returns 5.
    \item \texttt{ISNUMBER(A1)} returns TRUE.
    \item \texttt{=A1+5} returns \texttt{\#VALUE!} because text is not a
      number.
    \item Text is treated as 0 in every arithmetic operation.
  \end{enumerate}
\item \textit{[Original]} Which is the correct general form of COUNTIFS for
  counting rows that satisfy all the given conditions?
  \begin{enumerate}[label=\Alph*.]
    \item \texttt{COUNTIFS(criteria1, range1, criteria2, range2)}
    \item \texttt{COUNTIFS(range1, criteria1, range2, criteria2)}
    \item \texttt{COUNTIFS(range1 + range2, criteria1 AND criteria2)}
    \item \texttt{COUNTIFS(range1; criteria1; TRUE)}
  \end{enumerate}
\item \textit{[Original]} Which formula counts the rows in which A1:A10 exceeds
  50 \textbf{and} B1:B10 contains ``Yes''?
  \begin{enumerate}[label=\Alph*.]
    \item \texttt{=SUM(A1:A10>50, B1:B10="Yes")}
    \item \texttt{=SUMPRODUCT((A1:A10>50)*(B1:B10="Yes"))}
    \item \texttt{=AVERAGE(A1:A10, B1:B10)}
    \item \texttt{=COUNT(A1:A10, B1:B10)}
  \end{enumerate}
\item \textit{[Original]} Evaluate the statements:
  \begin{enumerate}[label=\Roman*.]
    \item COUNTIFS counts only the rows satisfying all the given
      range-and-criteria pairs.
    \item SUMPRODUCT can only add numbers and cannot count rows meeting
      multiple conditions.
    \item In a numeric comparison, a blank cell is treated as zero.
  \end{enumerate}
  \begin{enumerate}[label=\Alph*.]
    \item I and III only
    \item II only
    \item I, II, and III
    \item I and II only
  \end{enumerate}
\item \textit{[Original]} An email rule (filter) that moves all mail from one
  sender into a dedicated folder consists of:
  \begin{enumerate}[label=\Alph*.]
    \item a condition plus an action, applied automatically to mail.
    \item a digital signature plus encryption, applied manually.
    \item a Bcc field plus a Cc field, applied once.
    \item a Section Break plus a Page Break, applied to pages.
  \end{enumerate}
\item \textit{[Original]} Which statement correctly distinguishes a digital
  signature from encryption?
  \begin{enumerate}[label=\Alph*.]
    \item A digital signature encrypts the message body so only the recipient
      can read it.
    \item Encryption proves the sender's identity; a digital signature hides
      the content.
    \item A digital signature provides authenticity and integrity, while
      encryption provides confidentiality.
    \item A digital signature and encryption are different names for the same
      operation.
  \end{enumerate}
\item \textit{[Original]} Expand POP3.
  \begin{enumerate}[label=\Alph*.]
    \item Point-to-Point Protocol version 3
    \item Post Office Protocol version 3
    \item Packet Ordering Protocol version 3
    \item Private Office Protocol version 3
  \end{enumerate}
\item \textit{[Original]} Which statement about mail protocols is correct?
  \begin{enumerate}[label=\Alph*.]
    \item POP3 sends mail; SMTP downloads mail to the local device.
    \item IMAP downloads mail and always deletes it from the server; POP3
      keeps devices synchronised.
    \item SMTP sends mail; POP3 downloads mail to the local device; IMAP keeps
      mail on the server and synchronises devices.
    \item POP3 keeps mail synced on the server across all devices.
  \end{enumerate}
\item \textit{[Original]} A message is sent with addresses in the To, Cc, and
  Bcc fields. Who can see the Bcc addresses?
  \begin{enumerate}[label=\Alph*.]
    \item Every recipient, since Bcc is visible to all.
    \item Only the sender; the Bcc addresses are hidden from the other
      recipients.
    \item Only the Cc recipients; the sender cannot see them.
    \item Nobody, since Bcc addresses are deleted before sending.
  \end{enumerate}
\item \textit{[Original]} In PowerPoint, which content type CANNOT be converted
  to a SmartArt graphic in one step without re-entry?
  \begin{enumerate}[label=\Alph*.]
    \item A bulleted list of text
    \item A picture inserted on the slide
    \item Plain paragraph text with levels
    \item A numbered list of points
  \end{enumerate}
\item \textit{[Original]} Evaluate the statements:
  \begin{enumerate}[label=\Roman*.]
    \item A Section Break, not a Page Break, allows different headers and
      footers in one document.
    \item Addresses in the Bcc field are hidden from the sender.
    \item A digital signature encrypts the message body for confidentiality.
  \end{enumerate}
  \begin{enumerate}[label=\Alph*.]
    \item I only
    \item I and II only
    \item II and III only
    \item I, II, and III
  \end{enumerate}
\item \textit{[Original]} Which pairing is correctly matched?
  \begin{enumerate}[label=\Alph*.]
    \item Page Break --- starts a new section with its own headers and
      footers
    \item COUNTIFS --- handles one condition only; SUMPRODUCT --- only adds
      numbers
    \item First-line indent --- shifts only the first line; Hanging indent
      --- shifts every line except the first
    \item POP3 --- syncs mail on the server; IMAP --- downloads and deletes
      from the server
  \end{enumerate}

\item \textit{[JKSSB FAA-pattern]} Which shortcut italicises selected text in Word?
  \begin{enumerate}[label=\Alph*.]
    \item Ctrl+I
    \item Ctrl+P
    \item Ctrl+S
    \item Ctrl+N
  \end{enumerate}
\item \textit{[JKPSI-pattern]} The Bcc field differs from To/Cc because Bcc addresses are:
  \begin{enumerate}[label=\Alph*.]
    \item hidden from other recipients.
    \item printed in bold automatically.
    \item read aloud by speakers.
    \item stored on paper only.
  \end{enumerate}
\item \textit{[SSC CGL-pattern]} Which statement about email attachments is correct?
  \begin{enumerate}[label=\Alph*.]
    \item Unexpected attachments can carry malware and need caution.
    \item Every attachment is always safe to open.
    \item Attachments increase printer ink levels.
    \item Attachments can only be text on paper.
  \end{enumerate}
\item \textit{[JKSSB-pattern]} \textbf{(MSQ: select ALL correct options.)} Which statements about MS Office files are correct?
  \begin{enumerate}[label=\Alph*.]
    \item .docx holds Word documents.
    \item .xlsx holds Excel workbooks.
    \item .pptx holds PowerPoint presentations.
    \item .docx holds video films.
  \end{enumerate}
\item \textit{[SSC CGL-pattern]} \textbf{(MSQ: select ALL correct options.)} Consider the statements: I. POP3 downloads mail to one machine. II. IMAP syncs mail across devices. III. SMTP prints paper letters. Which are correct?
  \begin{enumerate}[label=\Alph*.]
    \item I and II only
    \item III only
    \item I, II, and III
    \item II and III only
  \end{enumerate}
\item \textit{[PYQ-adapted: SSC CGL 2025 Tier-II Paper-I, 19-Jan-2026 Q14; SRC-18]} Which setting can be changed for one Word section without changing the previous section?
  \begin{enumerate}[label=\Alph*.]
\item Page orientation and margins
   \item AutoCorrect replacement list
   \item Normal style definition
   \item File Properties metadata
  \end{enumerate}

\item \textit{[PYQ-adapted: SSC CGL 2025 Tier-II Paper-I, 19-Jan-2026 Q4; SRC-18]} Which Excel function counts cells that meet one criterion?
  \begin{enumerate}[label=\Alph*.]
\item COUNT
   \item COUNTBLANK
   \item COUNTIF
   \item SUMIF
  \end{enumerate}

\item \textit{[PYQ-adapted: JKSSB Junior Assistant 2026 Q76, final key; SRC-03/SRC-06]} Which email field hides listed recipients from other recipients?
  \begin{enumerate}[label=\Alph*.]
\item To
   \item Cc
   \item Bcc
   \item Subject
  \end{enumerate}

\item \textit{[Original]} Which feature in Word records edits for later review?
  \begin{enumerate}[label=\Alph*.]
\item Track Changes
   \item Mail Merge
   \item AutoSum
   \item Slide Sorter
  \end{enumerate}

\item \textit{[Original, MSQ drill]} \textbf{(MSQ: select ALL correct options.)} Which distinctions are accurate?
  \begin{enumerate}[label=\Alph*.]
\item Section Break can enable different page setup for a new section.
   \item Bcc conceals recipients from other recipients.
   \item COUNTIF counts cells meeting criteria.
   \item Digital signatures are the same as message encryption.
  \end{enumerate}

\end{enumerate}
\answersection{7}
\answersubheading{22}{MS Word Fundamentals}
\begin{enumerate}[label=\textbf{Q\arabic*.}, leftmargin=*]
\item \textbf{B.} MS Word is a word processor, used to create, edit, format,
  and print text documents. A spreadsheet (A) is Excel; a DBMS (C) is Access;
  a presentation program (D) is PowerPoint.
\item \textbf{C.} The modern default Word format uses the .docx extension.
  .xlsx (A) is Excel, .pptx (B) is PowerPoint, and .txt (D) is plain text,
  not the Word document format.
\item \textbf{A.} QAT stands for Quick Access Toolbar. B, C, and D are not
  expansions of the abbreviation.
\item \textbf{B.} The Ribbon replaced the older menus and toolbars. The
  status bar (A) and ruler (C) still exist alongside the Ribbon, and
  shortcuts (D) are keyboard actions, not the interface the Ribbon replaced.
\item \textbf{C.} The Font and Paragraph groups are on the Home tab. Insert
  (A) is for tables and pictures; View (B) is for document views; Layout (D)
  is for margins, orientation, and breaks.
\item \textbf{A.} By default the Quick Access Toolbar shows Save, Undo, and
  Redo. Bold/Italic/Underline (B) and Copy/Cut/Paste (C) are Home-tab
  commands, and New/Open/Print (D) live in Backstage view, not on the QAT by
  default.
\item \textbf{B.} Ctrl+S saves the document. Ctrl+N (A) creates a new
  document; Ctrl+P (C) prints; Ctrl+O (D) opens a document.
\item \textbf{C.} Ctrl+P opens the Print screen. Ctrl+S (A) saves; Ctrl+E
  (B) centres text; Ctrl+Enter (D) inserts a page break.
\item \textbf{C.} The Quick Access Toolbar \emph{can} be customised, so the
  claim that it is fixed and cannot be customised is false. It sits at the
  top-left (A), is customisable (B), and can be moved below the Ribbon (D),
  all of which are true.
\item \textbf{D.} Justify makes text flush with both left and right margins.
  Left (A) is flush on the left only, Right (C) on the right only, and
  Centre (B) places each line in the middle.
\item \textbf{B.} Ctrl+E centres the paragraph. Ctrl+L (A) is left, Ctrl+R
  (C) is right, and Ctrl+J (D) is justify.
\item \textbf{B.} A hanging indent pushes every line except the first inward,
  as used in reference lists. A describes a first-line indent; C describes
  centring; D describes paragraph spacing, not indentation.
\item \textbf{B.} Numbering is used when the order of items matters. Bullets
  (A) are for unordered lists; text wrapping (C) controls flow around
  objects; justification (D) is an alignment, not a list feature.
\item \textbf{A.} The intersection of a row and a column in a table is a
  cell. Field (B) and record (C) are database terms; margin (D) is the blank
  space around the page.
\item \textbf{B.} Statements I and II are correct: the Home tab holds the
  Font and Paragraph groups, and the QAT shows Save, Undo, and Redo by
  default. III is false because a Page Break does not change headers and
  footers --- that needs a Section Break. A omits II; C includes false III
  and omits I; D includes false III.
\item \textbf{B.} The footer repeats at the bottom of every page. The header
  (A) repeats at the top; the body (C) is the main text area; the margin (D)
  is the blank space around the page.
\item \textbf{B.} Portrait is taller than wide and Landscape is wider than
  tall. A reverses the two; C is wrong because Landscape is wider than tall;
  D is false because orientation changes the page dimensions.
\item \textbf{B.} A Section Break creates a new section that can carry its
  own headers and footers. A Page Break (A) only starts a new page; a Line
  Break (C) moves to the next line; a Column Break (D) moves to the next
  column, none of which changes headers and footers.
\item \textbf{A.} Text wrapping controls how text flows around a picture.
  Text alignment (B), line spacing (C), and indentation (D) affect the text
  itself, not the flow around an image.
\item \textbf{C.} Print preview shows how the pages will appear before
  printing. A is wrong because Justify flushes both margins, not the right
  only; B is wrong because a footer repeats at the bottom, not the top; D is
  wrong because bullets are for unordered lists, while numbering serves
  ordered lists.

\item \textbf{A.} Mail Merge joins letters with address data. B counts words; C and D adjust the view.
\item \textbf{A.} Ctrl+S saves. B prints, C bolds, and D centres text.
\item \textbf{A.} Center alignment centres lines. B, C, and D align differently.
\item \textbf{A, B, C.} Bold, italic, and underline format text. Paper weight (D) is a physical property.
\item \textbf{A, B, C.} The default format, Track Changes, and headers are correct. D is false: saving does not print.
\item \textbf{B.} Page setup belongs to each section, so orientation and margins can differ.
  AutoCorrect is application-wide, the Normal style applies through the document,
  and Properties describe the file rather than one section.

\item \textbf{A.} A Section Break creates a new section with independent header/footer
  settings. A Page Break only starts a page; line and column breaks change text
  flow without creating a section.

\item \textbf{B.} Mail Merge combines a main document with a recipient data source. Track
  Changes records edits, SmartArt creates diagrams, and Word Count reports text
  statistics. The cited key is provisional.

\item \textbf{A.} Ctrl+H opens Find and Replace. Ctrl+P opens Print, Ctrl+B toggles bold, and
  Ctrl+K inserts a hyperlink.

\item \textbf{A, B, C.} Orientation, margins, and headers/footers can vary by section. The
  Normal style is a document-level style, not a per-section setting.

\end{enumerate}
\answersubheading{23}{MS Word High-Yield}
\begin{enumerate}[label=\textbf{Q\arabic*.}, leftmargin=*]
\item \textbf{B.} Ctrl+H opens Find and Replace. A (Ctrl+F) opens Find only;
  C (Ctrl+G) opens Go To; D (Ctrl+K) inserts a hyperlink.
\item \textbf{B.} F7 starts the spelling and grammar check. A (F5) starts a
  PowerPoint slide show; C (F9) and D (F11) serve other functions, not Word
  proofing.
\item \textbf{A.} Ctrl+K inserts a hyperlink. B (Ctrl+L) left-aligns; C
  (Ctrl+H) opens Find and Replace; D (Ctrl+J) justifies text.
\item \textbf{B.} Ctrl+F locates text without replacing it. A describes
  Ctrl+H; C describes Ctrl+K; D describes Ctrl+P.
\item \textbf{B.} A bookmark is a named location or selection used as a jump
  target. A describes a hyperlink; C describes a comment; D describes a
  tracked change.
\item \textbf{D.} Mail Merge combines a letter with a database of addresses.
  A (Track Changes) records edits; B (Header \& Footer) repeats content on
  pages; C (AutoCorrect) fixes typing errors as you type.
\item \textbf{A.} The process that combines names and addresses with a
  standard document is Mail Merge. B, C and D are not real Word feature names.
\item \textbf{B.} The main document is the standard letter or template that
  holds the merge fields. A is the data source; C and D describe settings and
  storage, not the main document.
\item \textbf{D.} A Section Break creates a new section, which can carry its
  own headers and footers. A Page Break (A), a Line Break (B) and a Column
  Break (C) do not create a new section.
\item \textbf{A.} Section Breaks allow different headers and footers in the
  same document. B is false because a Page Break keeps the same section
  formatting; C and D are false because headers/footers are fully editable and
  are not deleted on closing.
\item \textbf{A.} This statement is NOT correct: a content control \emph{can}
  enforce data validation (for example, restricting entry to a list or a date
  range). B is correct (per-paragraph indents can differ even under one
  style); C is correct (comments appear in the Reviewing/Comments pane, not the
  Navigation Pane); D is correct (accepting a change removes its revision
  metadata).
\item \textbf{B.} A comment attaches a reviewer's note without changing the
  document text. A describes a hyperlink; C describes Track Changes; D
  describes a bookmark.
\item \textbf{B.} Track Changes stores the author's name and the time of each
  edit as revision metadata. A, C and D describe unrelated data that is not
  stored per change.
\item \textbf{A.} Statements I and II are correct: Ctrl+H opens Find and
  Replace, and Track Changes records each edit's author. III is false because a
  Page Break cannot create a new header/footer section. B and D include false
  III; C omits I and includes false III.
\item \textbf{A.} A Heading style gives the paragraph an outline level, so it
  appears in the Navigation Pane and a Table of Contents. B, C and D describe
  unrelated behaviours.
\item \textbf{B.} PDF stands for Portable Document Format. A, C and D are not
  the correct full form.
\item \textbf{B.} Ctrl+S saves the document and Ctrl+P opens the print
  options. A reverses them; C and D misname the commands.
\item \textbf{D.} A PivotTable is an Excel feature, not a Word review or
  editing feature. Track Changes (A), Comments (B) and Bookmarks (C) are all
  Word features.
\item \textbf{B.} Alt+Shift+C removes the document window split. A
  (Ctrl+Shift+S) is not that command; C (Ctrl+Alt+C) inserts the copyright
  symbol; D (Shift+F6) moves between panes. Ctrl+Alt+S toggles the split on and
  off but is not the documented ``remove'' shortcut.
\item \textbf{C.} Ctrl+K inserts a hyperlink. A (Ctrl+H) opens Find and
  Replace; B (F7) runs the spelling and grammar check; D (Ctrl+F) opens Find.

\item \textbf{A.} F7 opens spell/grammar checking. F5 (B) runs slide shows; C and D serve other roles.
\item \textbf{A.} Section Breaks restart headers/footers. B, C, and D do not divide sections.
\item \textbf{A.} Print Layout previews the printed page. B, C, and D are non-view objects.
\item \textbf{A, B, C.} Page numbers, titles, and date fields belong in headers/footers. The power supply (D) is hardware.
\item \textbf{A, B, C.} Bold shortcut, tables, and orientations are correct. D is false: F7 checks spelling; F5 runs shows.
\item \textbf{A.} Direct formatting changes only the selected instance. Redefining a shared style changes all text using it; templates and AutoCorrect do not isolate one heading.

\item \textbf{A.} Ctrl+H opens Find and Replace. Ctrl+P prints, Ctrl+S saves, and Ctrl+L commonly aligns text left.

\item \textbf{C.} F7 starts spelling and grammar review in common Word versions. F1 opens help, F5 navigates/goes to, and F12 opens Save As.

\item \textbf{A.} Track Changes records edits for review. The other features count words, create personalized documents, or change page appearance.

\item \textbf{A, B, C.} Comments and tracked changes support review; Accept/Reject resolves tracked edits. Mail Merge generates personalized documents.

\end{enumerate}
\answersubheading{24}{Excel Structure}
\begin{enumerate}[label=\textbf{Q\arabic*.}, leftmargin=*]
\item \textbf{B.} A workbook is the Excel file and can contain one or more
  worksheets. A describes a worksheet; C describes a formula; D describes the
  Formula Bar.
\item \textbf{B.} A worksheet is a single sheet of rows and columns inside a
  workbook. A describes a workbook; C describes a range; D describes a cell's
  address, not a worksheet.
\item \textbf{B.} A cell is the intersection of a row and a column. A range
  (A) is a block of cells; a worksheet (C) is the whole sheet; a workbook (D)
  is the file.
\item \textbf{B.} A cell address is written as the column letter followed by
  the row number, so B5 is valid. A reverses the order; C and D use invalid
  separators.
\item \textbf{A.} The selected cell is the active cell. ``Name cell'', ``anchor
  cell'', and ``formula cell'' are not Excel terms.
\item \textbf{B.} The Formula Bar shows the contents of the active cell,
  including any formula. The Name Box (A) shows the address; a sheet tab (C)
  switches worksheets; the status bar (D) shows status information.
\item \textbf{B.} The Name Box shows the address of the active cell or the
  name of the selected range. Contents appear in the Formula Bar (A); a sum
  is shown in the status bar (C); the workbook name is in the title bar (D).
\item \textbf{B.} A group of two or more cells is a range. A workbook (A) is
  the file; a tab (C) switches sheets; a field (D) is a database term, not an
  Excel block of cells.
\item \textbf{C.} A modern worksheet has 1,048,576 rows. 65,536 (A) was the
  old \texttt{.xls} limit; 16,384 (D) is the column count, not the row count.
\item \textbf{C.} A modern worksheet has 16,384 columns. 256 (A) was the old
  \texttt{.xls} limit; 1,048,576 (D) is the row count.
\item \textbf{D.} The last column in Excel 2007 and later is XFD. IV (C) was
  the last column of the older \texttt{.xls} format; Z (A) and ZZ (B) are not
  the final column.
\item \textbf{C.} A workbook is the Excel file and can contain one or more
  worksheets. A and B swap the two terms; D is false because a workbook can
  hold many worksheets.
\item \textbf{C.} It is \emph{not} true that a worksheet is the same as the
  workbook file --- the workbook is the file and the worksheet is a page
  inside it. A, B, and D are all true statements.
\item \textbf{B.} A range has two corner addresses joined by a colon, so
  A1:C10 is a range. B5 (A) and \$A\$1 (C) are single-cell references, not
  ranges; D is a function formula.
\item \textbf{B.} By default text is left-aligned and numbers are
  right-aligned. A reverses the two; C and D misalign text.
\item \textbf{D.} A row is a structural line of the grid, not a data type
  entered in a cell. Text (A), number (B), and formula (C) are all valid
  cell data types.
\item \textbf{B.} The default workbook extension in Excel 2007 and later is
  .xlsx. .xls (A) is the older format; .docx (C) is a Word document; .csv (D)
  is plain text data.
\item \textbf{D.} All three statements are correct: a workbook can hold
  multiple worksheets, the Name Box shows the active cell's address, and the
  Formula Bar shows the active cell's contents. A, B, and C each omit a true
  statement.
\item \textbf{C.} A1:B10 covers two columns (A, B) and ten rows (1--10), so it
  has $2 \times 10 = 20$ cells. A counts only the rows; B and D are wrong
  products.
\item \textbf{C.} The active cell is the currently selected cell, so this
  pairing is correct. A swaps workbook and worksheet; B swaps the file with a
  sheet; D wrongly calls a single address a range.

\item \textbf{A.} Addresses like B5 locate cells. B, C, and D are unrelated device properties.
\item \textbf{A.} Workbooks contain worksheets. B, C, and D are non-Excel objects.
\item \textbf{A.} The row-column crossing is a cell. B, C, and D are look-alike distractors.
\item \textbf{A, B, C.} Formula bar, Name box, and grid are Excel parts. A radiator (D) is not.
\item \textbf{A, B, C.} Lettered columns, numbered rows, and the active cell are correct. D is false: workbooks hold whole sheets.
\item \textbf{A.} A cell is the intersection of a row and column. A workbook contains worksheets; the formula bar and Ribbon are interface elements.

\item \textbf{A.} A workbook is the Excel file that can contain worksheets. A range is a set of cells; chart area and field have different meanings.

\item \textbf{A.} The Name Box displays the active cell reference and accepts a destination. The Formula bar displays cell content; the status and title bars report other information.

\item \textbf{A.} The letter identifies the column; the number identifies the row.

\item \textbf{A, B, C.} The first three definitions are correct. The Formula bar displays or edits the active cell's contents; worksheet tabs identify sheets.

\end{enumerate}
\answersubheading{25}{Excel Formulas}
\begin{enumerate}[label=\textbf{Q\arabic*.}, leftmargin=*]
\item \textbf{B.} A formula is an instruction, beginning with \texttt{=}, that
  calculates a cell's value from operands. A describes a constant; C ignores
  that the cell stores the formula, not a frozen result; D describes
  formatting, not a formula.
\item \textbf{B.} Multiplication uses the asterisk \texttt{*}. A (\texttt{x})
  and C (\texttt{.}) are not Excel operators, and D (\texttt{/}) is the
  division operator.
\item \textbf{A.} The \texttt{\$} locks the part of the reference that
  directly follows it, so that part is fixed when the formula is copied. B, C,
  and D describe formatting or access settings, not the meaning of \texttt{\$}
  in a reference.
\item \textbf{B.} A formula must begin with the equal sign \texttt{=} (TRAP-30,
  PYQ-38). A, C, and D are the trap distractors: a formula does not start with
  \texttt{+}, \texttt{@}, or \texttt{\#}.
\item \textbf{B.} An absolute reference locks both the column and the row with
  two dollar signs: \texttt{\$A\$1}. A (\texttt{A1}) is relative; C
  (\texttt{A\$1}) locks only the row; D (\texttt{\$A1}) locks only the column;
  so each is a mixed or relative reference, not absolute.
\item \textbf{A.} The ampersand \texttt{\&} concatenates (joins) text. B
  (\texttt{+}) adds numbers; C (\texttt{:}) is the range operator; D
  (\texttt{@}) is not an Excel operator.
\item \textbf{B.} In \texttt{=A1+\$A\$1}, only \texttt{\$A\$1} is absolute, so
  it stays the same when the formula is copied. A (\texttt{A1}) is relative
  and shifts; C is wrong because A1 does change; D is wrong because A1 changes.
\item \textbf{B.} \texttt{=A1} is relative; copying from B1 to B2 moves down
  one row, so the reference becomes \texttt{=A2}. A would require an absolute
  reference; C has the wrong column; D is the unchanged absolute form, not the
  result of copying a relative reference.
\item \textbf{A.} \texttt{=\$A\$1} is absolute: both the column and the row
  are locked, so copying to B5 leaves it \texttt{=\$A\$1}. B and C wrongly
  change the row; D wrongly changes the column.
\item \textbf{B.} In \texttt{=\$A1} the column is locked and the row is
  relative. Copying from B1 to C2 moves down one row, so the row follows to 2
  and the column stays A: \texttt{=\$A2}. A ignores the row change; C moves
  the locked column; D moves the column without following the row.
\item \textbf{A.} In \texttt{=A\$1} the row is locked and the column is
  relative. Copying from B1 to C2 moves one column right, so the column
  follows from A to B and the row stays 1: \texttt{=B\$1}. B wrongly changes
  the locked row; C ignores the column change; D keeps the locked column A
  instead of following it.
\item \textbf{B.} \texttt{=A2*\$B\$1}: the price reference \texttt{A2} is
  relative and follows the copy down to \texttt{A3}, while \texttt{\$B\$1} is
  absolute and stays fixed. A does not advance the price; C wrongly moves the
  locked rate; D moves both, which only happens if the rate were relative.
\item \textbf{A.} \texttt{\#DIV/0!} appears when a number is divided by zero
  or by an empty cell. B describes \texttt{\#REF!}; C describes
  \texttt{\#VALUE!}; D describes \texttt{\#NAME?}.
\item \textbf{A.} \texttt{\#REF!} means a reference is invalid because the
  referenced cell was deleted. B describes \texttt{\#DIV/0!}; C describes
  \texttt{\#VALUE!}; D describes a missing \texttt{=}, which does not produce
  \texttt{\#REF!}.
\item \textbf{A.} \texttt{\#VALUE!} means the wrong type of operand was used,
  such as text where a number is needed. B is not an Excel error of this kind;
  C describes \texttt{\#REF!}; D describes \texttt{\#N/A}.
\item \textbf{C.} A relative reference does \emph{not} always point to the
  same cell; it shifts when the formula is copied, so C is the false
  statement. A, B, and D are all true of relative references.
\item \textbf{B.} A \texttt{\$} before a letter locks the column, and a
  \texttt{\$} before a number locks the row, so \texttt{A\$1} locks the row
  while \texttt{\$A1} locks the column. A wrongly says \texttt{\$A1} locks
  both; C swaps the two; D wrongly says \texttt{\$A\$1} locks only the row.
\item \textbf{B.} Statement II (absolute references do not change) and
  statement III (relative references do change) are correct. Statement I is
  false: under the exam rule a formula must start with \texttt{=}. A omits III;
  C includes the false I; D includes the false I.
\item \textbf{A.} The operator \texttt{<>} means ``not equal to''. B is the
  equality operator \texttt{=}; C is \texttt{>}; D is \texttt{<=}.
\item \textbf{A.} Multiplication is evaluated before addition, so
  \texttt{3*4} is 12 and \texttt{2+12} is 14. B (20) would require
  \texttt{=(2+3)*4}; C (24) would require \texttt{=2*3*4}; D (9) would require
  \texttt{=(2+3)+4}.

\item \textbf{A.} The equals sign opens a formula. The other symbols do not mark a formula.
\item \textbf{A.} Dollar signs fix the reference absolutely. B, C, and D are not reference types.
\item \textbf{A.} The ampersand concatenates text. B adds numbers; C and D serve other roles.
\item \textbf{A, B, C.} References, operators, and functions build formulas. Ink colours (D) are irrelevant.
\item \textbf{A, B, C.} Auto-recalculation and reference behaviour are correct. D is false: no full stop is needed.
\item \textbf{A.} The dollar signs lock the column and row when a formula is copied. The other symbols do not create absolute references.

\item \textbf{A.} A relative reference shifts with the copy location, so the column changes from A to B. The other options lock a dimension or change the row.

\item \textbf{A.} A\$1 fixes row 1 but leaves the column relative. \$A1 locks the column; \$A\$1 locks both; A1 locks neither.

\item \textbf{A.} The ampersand concatenates text. The other symbols represent exponentiation, division, and comparison.

\item \textbf{A, B, C, D.} All four are valid A1-style references; dollar signs indicate absolute row/column components.

\end{enumerate}
\answersubheading{26}{Excel Functions}
\begin{enumerate}[label=\textbf{Q\arabic*.}, leftmargin=*]
\item \textbf{C.} With A1 blank, Excel reads the blank as 0 in a numeric
  comparison, so A1 is neither greater than 50 nor greater than 25; both tests
  fail and the final value \texttt{Low} is returned. A and B would need A1 to
  be above 50 or above 25; D is wrong because a comparison of a blank does not
  produce \texttt{\#VALUE!} (PYQ-14).
\item \textbf{A.} \texttt{AVERAGE} returns the arithmetic mean of the numbers
  in a range. \texttt{MEAN} (B) is not an Excel function, \texttt{TOTAL} (C)
  is not a function, and \texttt{MEDIAN} (D) returns the middle value, not the
  mean.
\item \textbf{A.} \texttt{COUNT} counts only the cells that contain numbers.
  \texttt{COUNTA} (B) also counts text, \texttt{COUNTIF} (C) counts cells
  meeting a condition, and \texttt{SUM} (D) adds values instead of counting
  them.
\item \textbf{B.} \texttt{COUNTA} counts every non-empty cell: the 6 numbers
  plus the 2 text labels make 8. The blank cells are not counted. 6 (A) is the
  \texttt{COUNT} result, 10 (C) counts the blanks too, and 2 (D) counts only
  the blanks.
\item \textbf{A.} \texttt{COUNT} counts numeric cells only, so it returns 6.
  8 (B) is the \texttt{COUNTA} result, 10 (C) is the total number of cells,
  and 4 (D) is not produced by either function.
\item \textbf{A.} \texttt{COUNTA} counts every non-empty cell, whether it
  holds a number or text. \texttt{COUNT} (B) would leave out the text,
  \texttt{COUNTBLANK} (C) counts only empty cells, and \texttt{SUM} (D) adds
  numeric values.
\item \textbf{A.} \texttt{MAX} returns the largest number in a range.
  \texttt{MIN} (B) returns the smallest, and \texttt{HIGH} (C) and
  \texttt{TOP} (D) are not Excel functions.
\item \textbf{A.} \texttt{MIN} returns the smallest number in a range.
  \texttt{MAX} (B) returns the largest, and \texttt{MINIMUM} (C) and
  \texttt{LEAST} (D) are not Excel functions.
\item \textbf{A.} \texttt{SUM} adds the numbers in a range. \texttt{ADD} (B),
  \texttt{TOTAL} (C), and \texttt{PLUS} (D) are not Excel functions.
\item \textbf{A.} AutoSum is inserted with \texttt{Alt + =}. \texttt{Ctrl + =}
  (B) is not the AutoSum shortcut, \texttt{Shift + F3} (C) opens the Insert
  Function dialog, and \texttt{Ctrl + S} (D) saves the workbook.
\item \textbf{A.} AutoSum inserts a \texttt{SUM} formula by default, written
  as \texttt{=SUM(...)} over the adjacent range. \texttt{=AVERAGE(...)} (B),
  \texttt{=COUNT(...)} (C), and \texttt{=MAX(...)} (D) are available only from
  the AutoSum drop-down, not by default.
\item \textbf{A.} A formula must start with \texttt{=} and may use a function;
  a function is a named built-in operation used inside a formula. B reverses
  the two definitions, C wrongly calls them identical, and D is false because
  formulas commonly contain functions.
\item \textbf{A.} Dividing by zero, or by a blank cell (read as zero), gives
  \texttt{\#DIV/0!}. \texttt{\#VALUE!} (B) is a wrong data type,
  \texttt{\#NAME?} (C) is an unrecognised name, and \texttt{\#REF!} (D) is an
  invalid reference.
\item \textbf{A.} An unrecognised name, such as a misspelt function name,
  gives \texttt{\#NAME?}. \texttt{\#REF!} (B) is a deleted reference,
  \texttt{\#N/A} (C) is a value not available, and \texttt{\#NUM!} (D) is an
  invalid number.
\item \textbf{A.} Using the wrong type of data, such as text where a number is
  needed, gives \texttt{\#VALUE!}. \texttt{\#NAME?} (B) is an unrecognised
  name, \texttt{\#DIV/0!} (C) is division by zero, and \texttt{\#N/A} (D) is a
  value not available.
\item \textbf{A.} A reference to a deleted cell is no longer valid and gives
  \texttt{\#REF!}. \texttt{\#NAME?} (B) is an unrecognised name,
  \texttt{\#VALUE!} (C) is a wrong data type, and \texttt{\#NULL!} (D) is the
  intersection of non-intersecting ranges.
\item \textbf{A.} A lookup that finds no match leaves the needed value
  unavailable, which gives \texttt{\#N/A}. \texttt{\#NUM!} (B) is an invalid
  number, \texttt{\#REF!} (C) is an invalid reference, and \texttt{\#DIV/0!}
  (D) is division by zero.
\item \textbf{D.} \texttt{\#TOTAL!} is not an Excel error value. The other
  three are all real: \texttt{\#DIV/0!} (A), \texttt{\#VALUE!} (B), and
  \texttt{\#NAME?} (C).
\item \textbf{D.} All three statements are correct: \texttt{COUNT} counts
  numbers while \texttt{COUNTA} counts all non-empty cells, a blank is read as
  zero in a numeric comparison, and every formula begins with \texttt{=}. A,
  B, and C each omit at least one true statement.
\item \textbf{A.} \texttt{SUM} returns the total (I-b), \texttt{MAX} returns
  the largest number (II-a), and \texttt{COUNT} returns the number of numeric
  cells (III-c). B, C, and D each assign at least one function to the wrong
  result.

\item \textbf{A.} SUM adds ranges. B, C, and D are not totalling functions.
\item \textbf{A.} COUNT tallies numeric cells. COUNTA (B) counts non-empty cells of any kind; C and D are devices.
\item \textbf{A.} MAX returns the largest value. B finds the smallest; C averages; D totals.
\item \textbf{A, B, C.} IF, COUNTIF, and SUMIF all test conditions. PRINTIF (D) is not an Excel function.
\item \textbf{A, B, C.} The AVERAGE, MIN, and IF roles are correct. D is false: SUM calculates; it does not print.
\item \textbf{C.} COUNTIF counts cells meeting a criterion. COUNT counts numeric cells
  without a criterion, COUNTBLANK counts empty cells, and SUMIF adds matching
  values rather than counting them.

\item \textbf{B.} Filter temporarily displays rows that meet criteria. Sort changes order,
  Remove Duplicates deletes repeated records, and Subtotal inserts grouped
  summaries.

\item \textbf{A.} COUNTA counts non-empty cells. COUNT counts numeric cells, COUNTBLANK
  counts empty cells, and SUM adds numeric values.

\item \textbf{B.} IF evaluates a logical test and selects between two results. MAX returns
  the largest value, AVERAGE calculates a mean, and ROUND changes numeric
  precision.

\item \textbf{A, B, C.} COUNT, COUNTA, and COUNTBLANK perform the stated counts. COUNTIF
  counts matching cells; SUMIF, not COUNTIF, sums their values.

\end{enumerate}
\answersubheading{27}{Excel Data Tools}
\begin{enumerate}[label=\textbf{Q\arabic*.}, leftmargin=*]
\item \textbf{B.} Sorting rearranges rows into a chosen order. A filter (A)
  hides rows rather than ordering them; Goal Seek (C) is a what-if tool;
  Protect Sheet (D) is a protection command.
\item \textbf{B.} \textbf{Ctrl}$+$\textbf{Shift}$+$\textbf{L} toggles
  AutoFilter. Ctrl+L (A), Alt+F (C), and Ctrl+F (D) do not turn AutoFilter on
  or off.
\item \textbf{A.} \textbf{Ctrl}$+$\textbf{T} converts a range into an Excel
  table. Ctrl+F (B) opens Find; Ctrl+Shift+T (C) and Alt+T (D) are not the
  table shortcut.
\item \textbf{B.} A filter temporarily hides the non-matching rows and does not
  remove them. It does not delete (A), move them to another sheet (C), or
  create a PivotTable (D); clearing the filter restores every row.
\item \textbf{D.} This is the \textbf{NOT} item. Ascending order arranges
  numbers from smallest to largest, so ``largest to smallest'' is descending and
  option D is false. Multi-level sort (A), sorting by colour (B), and custom
  lists (C) are all genuine Excel sorting features.
\item \textbf{A.} A table gives structured references such as
  \texttt{Table1[Sales]} and auto-expands; a plain range uses cell addresses.
  Entering text (B), entering formulas (C), and printing (D) are possible in a
  plain range too.
\item \textbf{B.} An Excel table is source data that a PivotTable can
  summarise. A PivotTable is not a kind of table (A); it never changes the
  source data (C); and a table can be filtered (D).
\item \textbf{C.} A pie chart shows the parts of one whole. A line chart (A)
  shows a trend over time; a scatter chart (B) shows a relationship between two
  numeric variables; a column chart (D) compares values across categories.
\item \textbf{B.} A chart is linked to its source range, so it updates
  automatically when the data changes. It is not frozen (A), deleted (C), or
  converted into a PivotTable (D).
\item \textbf{A.} The legend names each data series. The value axis (B) carries
  the numbers, the plot area (C) is where the data is drawn, and a data label
  (D) shows the value at a point.
\item \textbf{D.} This is a \textbf{NOT} item. The four PivotTable areas are
  Rows, Columns, Values, and Filters; ``Insert'' (D) is not one of them. Rows
  (A), Columns (B), and Values (C) are all genuine areas.
\item \textbf{A.} A PivotTable stores a summary view and does not track the
  source automatically, so it must be Refreshed. It does not update by itself
  (C); restarting the computer (B) and deleting the worksheet (D) are
  irrelevant.
\item \textbf{B.} Statements I and II are correct: a PivotTable summarises data
  without being the source table, and a PivotChart is linked to a PivotTable.
  Statement III is false because a filter hides rows temporarily --- it never
  deletes them. Options A, C, and D all include the false statement III.
\item \textbf{B.} COUNTIFS is built to count rows meeting several conditions at
  once. COUNT (A), SUM (C), and AVERAGE (D) do not take paired range-criteria
  arguments for multi-criterion counting.
\item \textbf{B.} COUNTIFS takes \textbf{paired} range-and-criteria arguments:
  \texttt{COUNTIFS(range1, crit1,} \texttt{range2, crit2)}. Option A reverses the
  pairing;
  option C gives only criteria with no ranges; option D tries to add two
  criteria into one argument.
\item \textbf{C.} COUNTIFS applies AND logic across paired range-and-criteria
  arguments, and SUMPRODUCT can also count multi-criterion rows using boolean
  arrays --- the stable distinction behind PYQ-15 and TRAP-14. COUNTIFS can
  count two criteria at once (A is false); SUMPRODUCT can count as well as add
  (B is false); neither applies only OR logic (D is false).
\item \textbf{B.} Print Titles repeats chosen header rows or columns on every
  printed page. Print Area (A) limits what is printed; Freeze Panes (C) is a
  screen-view aid; Page Break Preview (D) only shows page breaks.
\item \textbf{B.} Gridlines are \textbf{not} printed by default; the Print
  gridlines option must be enabled. Options A, C, and D each state an incorrect
  condition.
\item \textbf{C.} A cell is Locked by default, but locking takes effect only
  after the sheet is protected. Option A is backwards --- a lock has no effect
  while the sheet is unprotected; Protect Workbook (B) protects structure, not
  individual cells; protection (D) is not file encryption.
\item \textbf{A.} Goal Seek finds the input required to reach a target result.
  Scenario Manager (B) stores alternative input sets; a Data Table (C) shows
  results for one or two changing inputs; a PivotChart (D) is a chart on a
  PivotTable, not a what-if tool.

\item \textbf{A.} Sort orders rows. B previews printing; C and D are unrelated tools.
\item \textbf{A.} Filter displays matching rows. B, C, and D are hardware devices.
\item \textbf{A.} PivotTables summarise big tables. B, C, and D are non-analysis tasks.
\item \textbf{A, B, C.} Sort, Filter, and PivotTable analyse data. A paperweight (D) is not a tool.
\item \textbf{A, B, C.} Sort, Filter, and Chart roles are correct. D is false: sorting reorders; it does not delete files.
\item \textbf{A.} Filter displays rows meeting criteria without deleting the rest. Remove Duplicates removes records; deleting a sheet removes content; sorting changes order.

\item \textbf{A.} A PivotTable summarizes and rearranges data by fields. The other features belong to text correction, Word, or PowerPoint.

\item \textbf{A.} A line chart makes change over an ordered time axis easy to see. Pie/doughnut charts emphasize parts of a whole; radar charts compare dimensions.

\item \textbf{A.} AutoFilter provides criteria drop-downs. Goal Seek solves for an input; the other tools serve unrelated functions.

\item \textbf{A, B, C.} Sorting changes order, filtering changes visible rows, and PivotTables aggregate. Remove Duplicates changes records, not just formatting.

\end{enumerate}
\answersubheading{28}{Access Database Basics}
\begin{enumerate}[label=\textbf{Q\arabic*.}, leftmargin=*]
\item \textbf{B.} DBMS stands for Database Management System. A, C, and D are
  not the expansions of the abbreviation.
\item \textbf{C.} Access is a database management system (DBMS). A
  spreadsheet (A) and a word processor (B) are different Office tools, and a
  presentation program (D) is PowerPoint, not Access.
\item \textbf{B.} A database is the organised collection of data, while a
  DBMS is the software that manages it. A reverses the two; C wrongly says
  they are the same; D wrongly calls a DBMS a spreadsheet.
\item \textbf{A.} Microsoft Access is part of the Microsoft Office suite.
  Adobe Creative Cloud (B), Google Workspace (C), and LibreOffice (D) are
  other software families, not the suite that contains Access.
\item \textbf{C.} A row is one record --- one complete entry. A field (A) is
  a column, a table (B) is the whole collection, and a query (D) is an
  object that retrieves data.
\item \textbf{B.} A column is one field --- one kind of information. A record
  (A) is a row, a report (C) is a printable output object, and a form (D) is
  a data-entry screen.
\item \textbf{C.} A table is a collection of related records arranged in rows
  and columns. A field (A) is a single column, a record (B) is a single row,
  and a data type (D) describes the kind of value a field holds.
\item \textbf{D.} The table stores the actual data. A query (A) retrieves
  data, a form (B) is for entry and viewing, and a report (C) is for
  printing; none of these stores the raw data.
\item \textbf{D.} A report presents data in a printable layout. A table (A)
  stores data, a query (B) retrieves selected data, and a form (C) is a
  data-entry screen.
\item \textbf{A.} A form provides a screen to enter and view data. A table
  (B) stores the data, an index (C) speeds up search, and a data type (D)
  describes the kind of value a field holds.
\item \textbf{C.} AutoNumber automatically inserts a unique number for each
  new record. Short Text (A) holds text, Number (B) holds values for
  calculation, and Yes/No (D) holds True or False.
\item \textbf{B.} Yes/No stores only a True or False value. Number (A) is for
  calculated values, Date/Time (C) is for dates and times, and Short Text (D)
  is for text.
\item \textbf{C.} Short Text is best for a name, which is text. Number (A)
  is for calculations, Date/Time (B) is for dates and times, and AutoNumber
  (D) auto-fills a unique number.
\item \textbf{A.} The rule \texttt{Between 0 And 100} accepts only values in
  that range and rejects others. Required (B) stops blank entries, an index
  (C) speeds up search, and a unique number per record (D) is AutoNumber's
  behaviour --- none of these is what a range rule does.
\item \textbf{C.} Validation Text is the message shown when a Validation Rule
  is violated. Default Value (A) pre-fills a value, Input Mask (B) fixes the
  entry pattern, and Field Size (D) limits how much a field can hold.
\item \textbf{A.} The Required property, set to Yes, stops a field from being
  left blank. Format (B) and Caption (C) affect display, and Indexed (D)
  controls whether the field is indexed for speed.
\item \textbf{B.} An index speeds up searching and sorting. Validation (A) is
  a separate check, storage (C) is the job of the table, and unique
  identification (D) is the job of the primary key --- an index by itself
  does not make a field unique.
\item \textbf{C.} A primary key \textbf{cannot} contain duplicate values, so
  saying that it can is false. It does uniquely identify each record (A),
  cannot be left blank (B), and is often an AutoNumber field (D).
\item \textbf{D.} All three statements are correct: a field is a column, a
  record is a row, and a table is a collection of related records. A, B, and
  C each omit at least one correct statement.
\item \textbf{C.} AutoNumber is a data type that adds a unique number to each
  new record, so this pairing is correct. A mislabels a field (a field is a
  column, not a complete entry); B mislabels a record (a record is a row, not
  one kind of information); D mislabels a DBMS (a DBMS is the software, not
  the data).

\item \textbf{A.} Rows are records. B, C, and D are hardware, not row names.
\item \textbf{A.} Columns are fields. B, C, and D are unrelated words.
\item \textbf{A.} DBMS is Database Management System. B, C, and D are not the recognised expansion.
\item \textbf{A, B, C.} Records, fields, and key values structure tables. Ink levels (D) are printer data.
\item \textbf{A, B, C.} The table, record, and field facts are correct. D is false: tables are digital structures.
\item \textbf{A.} A row is a record; a column is a field. Queries retrieve/process data and forms provide an interface.

\item \textbf{A.} A field is a column storing one attribute. A row is a record; reports and database files are higher-level objects.

\item \textbf{A.} Date/Time stores date and time values. Currency stores money; Attachment stores files; Yes/No stores Boolean values.

\item \textbf{A.} A DBMS manages databases and access to their data. The other options describe unrelated tasks.

\item \textbf{A, B, C, D.} Tables store data, queries retrieve/process it, forms support entry/viewing, and reports format output.

\end{enumerate}
\answersubheading{29}{Access Keys and Relationships}
\begin{enumerate}[label=\textbf{Q\arabic*.}, leftmargin=*]
\item \textbf{B.} PK stands for Primary Key. A (Public Key) and D (Private
  Key) come from cryptography, and C (Parent Key) is not the standard
  expansion.
\item \textbf{C.} FK stands for Foreign Key. A, B, and D are not the standard
  expansions of the abbreviation.
\item \textbf{C.} The primary key uniquely identifies each record. A foreign
  key (A) only links to another table; an alternate key (B) is an unused
  candidate key; a secondary key (D) is not the identifying key.
\item \textbf{C.} A primary key must be unique, so it \textbf{cannot} contain
  duplicate values. Statements A, B, and D are all true of a primary key.
\item \textbf{C.} A field that refers to another table's primary key is the
  foreign key. A primary key (A) identifies records in its own table; a
  composite key (B) is a multi-field primary key; a candidate key (D) is a key
  eligible to become the primary key.
\item \textbf{B.} A composite (compound) key is a primary key made of two or
  more fields together. A primary key (A) may be a single field; a foreign key
  (C) links tables; an alternate key (D) is an unchosen candidate key.
\item \textbf{C.} A foreign key may repeat and may be blank. Primary and
  composite keys must be unique and not blank, and a candidate key is a
  potential primary key, so none of them may repeat.
\item \textbf{B.} A table can have \textbf{at most one} primary key (it may
  have none while being designed). ``Exactly one'' (A) is too strong; C and D
  are wrong because a table cannot have two or more primary keys.
\item \textbf{B.} One-to-many is the most common relationship in Access. A
  one-to-one (A) is the least common; a many-to-many (C) is built indirectly;
  zero-to-one (D) is not a standard Access relationship type.
\item \textbf{B.} The ``many'' side holds the foreign key. The primary key
  (A) sits on the ``one'' side; a composite key (C) is a multi-field primary
  key; an alternate key (D) is relevant only to candidate keys.
\item \textbf{C.} A many-to-many relationship cannot be stored directly and
  needs a junction table. One-to-one (A) and one-to-many (B) are stored
  directly; one-to-zero (D) is not a standard type.
\item \textbf{B.} A junction table usually uses a composite key formed from
  the two foreign keys. A single AutoNumber field (A) is not the typical
  junction key, and C and D name only one of the two linked tables.
\item \textbf{B.} Referential integrity requires a foreign-key value to match
  an existing primary-key value (or be blank). A (composite keys are optional)
  is unrelated; C describes data typing; D is false because a table has at
  most one primary key.
\item \textbf{B.} Cascade Delete removes the related child records when a
  parent is deleted. Cascade Update (A) changes child values when the parent
  key changes; referential integrity (C) is the overall rule; a composite key
  (D) is unrelated.
\item \textbf{B.} Only statement II is correct: a foreign key may contain
  duplicate values. Statement I is false (a primary key is unique), and
  statement III is false (a table has at most one primary key). A, C, and D
  all include a false statement.
\item \textbf{A.} \texttt{.mdb} is the older format used by Access 2003 and
  earlier. B describes \texttt{.accdb}; C and D are false because
  \texttt{.mdb} is not the modern default.
\item \textbf{B.} \texttt{.accdb} is the newer, default Access format used
  from Access 2007 onward. A is wrong (it is the newer format, not older); C
  and D name unrelated file types.
\item \textbf{C.} Statement C is incorrect: \texttt{.mdb} is the older
  format, not the newer one. A, B, and D are all correct statements about the
  two formats.
\item \textbf{C.} A composite key is a primary key made of two or more
  fields, so this pairing is correct. A is wrong (a primary key cannot
  repeat); B is wrong (a foreign key does not identify records in its own
  table); D is wrong (a junction table implements many-to-many, not
  one-to-one).
\item \textbf{B.} A candidate key not chosen as the primary key is an
  alternate key. A foreign key (A) links tables; a composite key (C) is a
  multi-field primary key; ``junction key'' (D) is not a standard term.

\item \textbf{A.} Primary keys uniquely identify rows. B, C, and D violate uniqueness or relevance.
\item \textbf{A.} Foreign keys reference related tables. B, C, and D describe physical locks, not database keys.
\item \textbf{A.} One-to-many links one row with many rows. B, C, and D are not relationship types.
\item \textbf{A, B, C.} Primary, foreign, and candidate keys are database concepts. A door key (D) is physical.
\item \textbf{A, B, C.} Uniqueness, referencing, and no-duplicate rules are correct. D is false: printers are unrelated.
\item \textbf{A.} A primary key uniquely identifies records and cannot be null. The other choices defeat reliable identification.

\item \textbf{A.} A foreign key links a record to a key in another table and supports referential integrity. The alternatives are unrelated.

\item \textbf{A.} A composite key uses multiple fields together to identify a record. The other terms are not standard key categories.

\item \textbf{A.} Referential integrity constrains relationships, such as preventing unmatched foreign-key references. The other options are unrelated.

\item \textbf{A, B, C.} The first three statements are correct. Duplicate primary-key values make records ambiguous.

\end{enumerate}
\answersubheading{30}{Access Queries, Forms and Reports}
\begin{enumerate}[label=\textbf{Q\arabic*.}, leftmargin=*]
\item \textbf{B.} SQL stands for Structured Query Language. A, C and D are
  not the correct expansions of the abbreviation.
\item \textbf{C.} A table stores the actual data in rows and columns. A
  query (A), a form (B) and a report (D) all work on that data; they do not
  store it.
\item \textbf{B.} A query asks a question and retrieves selected data from
  the tables. A form (A) is a data-entry screen; a report (C) is printed
  output; a table (D) is the store.
\item \textbf{D.} A query that retrieves data on the basis of a condition is
  a Select query. Delete (A), Update (B) and Append (C) are action queries
  that modify data.
\item \textbf{B.} An Update query changes the existing values in records. A
  Select query (A) only retrieves data; an Append query (C) adds records; a
  crosstab query (D) summarises data rather than editing it.
\item \textbf{C.} An Append query adds records from one table to another. An
  Update query (A) changes existing rows; a Delete query (B) removes rows; a
  Select query (D) retrieves rows.
\item \textbf{A.} A Delete query removes the records that match a condition.
  Select (B) retrieves, Update (C) changes, and Append (D) adds records.
\item \textbf{D.} A Select query only retrieves data, so it does not modify
  the stored data. Update (A), Append (B) and Delete (C) are action queries
  that all change the data.
\item \textbf{C.} A form provides an on-screen interface to enter and view
  data. A table (A) stores data, a query (B) retrieves it, and a report (D)
  prints it.
\item \textbf{B.} A report presents data in a formatted, printable layout. A
  form (A) is for on-screen data entry; a query (C) retrieves data; a table
  (D) stores it.
\item \textbf{B.} A report is read-only and is meant for printing and
  sharing. Entering records (A) is a form's job; storing data (C) is a
  table's job; modifying records (D) is an action query's job.
\item \textbf{C.} The sequence is Table (stores) $\rightarrow$ Query
  (selects) $\rightarrow$ Form (enters) $\rightarrow$ Report (prints). A, B
  and D scramble this order.
\item \textbf{B.} Forms and reports are usually based on a table or a query.
  A form is not built on another form (A); a report is not built only on a
  report (C); the Navigation Pane (D) is just a list of objects, not a data
  source.
\item \textbf{C.} A lookup field lets the user choose a value from a list
  supplied by a table or query. Deleting records (A) is a Delete query's job;
  printing a summary (B) is a report's job; a lookup does not create a table
  (D).
\item \textbf{D.} A spreadsheet is not an object in an MS Access database.
  Table (A), query (B) and report (C) are all standard Access objects.
\item \textbf{B.} Only statement II is correct: an Update query changes
  existing records. Statement I is false --- a Select query retrieves data,
  it does not modify it. Statement III is false --- a form enters and views
  data; printing is the job of a report.
\item \textbf{A.} The Query Design grid is called QBE (Query By Example).
  SQL (B) is the language behind the query; QWERTY (C) is a keyboard layout;
  PDF (D) is a file format.
\item \textbf{B.} \texttt{SELECT} retrieves data from a table. \texttt{UPDATE}
  (A) changes rows, \texttt{DELETE} (C) removes rows, and \texttt{INSERT}
  (D) adds rows.
\item \textbf{C.} The \texttt{WHERE} clause gives the condition that the
  rows must satisfy. The \texttt{FROM} clause names the table (A); the field
  list after \texttt{SELECT} names the columns (B); sorting is done with
  \texttt{ORDER BY}, not \texttt{WHERE} (D).
\item \textbf{C.} A form provides a screen for data entry, so this pairing is
  correct. A mislabels a table as asking questions; B mislabels a query as
  storing data; D mislabels a report as modifying records.

\item \textbf{A.} Queries select data. B, C, and D are display or sound items.
\item \textbf{A.} Forms give entry screens. B, C, and D are hardware parts.
\item \textbf{A.} SQL is Structured Query Language. B, C, and D are not the recognised expansion.
\item \textbf{A, B, C.} Tables, queries, and forms are Access objects. A cricket bat (D) is not.
\item \textbf{A, B, C.} The query, form, and report roles are correct. D is false: reports handle printing, not queries.
\item \textbf{A.} A query selects or manipulates data based on criteria. Forms support interaction, reports format output, and macros automate actions.

\item \textbf{A.} A report formats data for viewing or printing. Tables store records; relationships connect tables; modules contain code.

\item \textbf{A.} Forms provide an interface to enter or view data. Reports present formatted output; indexes and relationships support data organization.

\item \textbf{A.} A SELECT query returns records from one or more tables based on criteria. The other choices are unrelated.

\item \textbf{A, B, C, D.} All four are standard roles of Access objects.

\end{enumerate}
\answersubheading{31}{PowerPoint Basics}
\begin{enumerate}[label=\textbf{Q\arabic*.}, leftmargin=*]
\item \textbf{B.} PPTX stands for PowerPoint XML Presentation. A, C, and D
  are not the correct expansion.
\item \textbf{C.} The default working extension is \texttt{.pptx}.
  \texttt{.ppsx} (A) is the show file, \texttt{.potx} (B) is the template,
  and \texttt{.pptm} (D) is the macro-enabled file.
\item \textbf{B.} One page of a presentation is a slide. Worksheet (A)
  belongs to Excel, document (C) to Word, and record (D) to a database.
\item \textbf{A.} The dotted-line content box is a placeholder. Cell (B) is
  a spreadsheet unit, field (C) a database unit, and frame border (D) is not
  a PowerPoint term.
\item \textbf{B.} The arrangement of placeholders is the layout. A theme (A)
  is colours, fonts, and effects; the master (C) controls design; a template
  (D) is a starter file.
\item \textbf{A.} Colours, fonts, and effects applied from the Design tab
  form a theme. A placeholder (B) holds content, a macro (C) automates
  steps, and a handout (D) is a printed sheet.
\item \textbf{D.} The Slide Master controls linked design, is edited once
  for all slides, and opens from the View tab, so A, B, and C are true.
  Printing six-slides-per-page handouts is the job of the Handout Master,
  not the Slide Master.
\item \textbf{B.} The Handout Master controls only printed handout layouts.
  The Slide Master (A) controls on-screen slides; Title Master (C) and
  Animation Master (D) are not PowerPoint masters.
\item \textbf{C.} Normal view is the default editing view. Slide Sorter (A)
  manages order, Reading View (B) reviews, and Slide Show view (D) presents.
\item \textbf{C.} Slide Sorter view shows all thumbnails for drag-and-drop
  reordering. Normal (A) edits one slide, Notes Page (B) pairs a slide with
  notes, and Outline View (D) shows text only.
\item \textbf{B.} Notes are the speaker's hidden reminders below a slide.
  Handouts (A) serve the audience, a transition (C) moves between slides,
  and the top-level controlling slide (D) is the Slide Master.
\item \textbf{B.} Multi-slide printed sheets for the audience are handouts.
  Notes (A) serve the speaker, placeholders (C) hold content, and
  transitions (D) move between slides.
\item \textbf{B.} F5 starts the show from the beginning. F1 (A) opens Help,
  F7 (C) runs spelling check, and F9 (D) recalculates or has no show role.
\item \textbf{A.} Shift+F5 starts from the current slide. The other
  combinations (B, C, D) do not carry this meaning in PowerPoint.
\item \textbf{B.} Between-slide movement is a transition; on-slide object
  movement is an animation. A reverses the pair; C and D name unrelated
  design and print terms.
\item \textbf{B.} \texttt{.potx} is the PowerPoint template, so this pairing
  is correct. A mislabels \texttt{.ppsx}, which is the show file; C mislabels
  \texttt{.pptx}, which is the editable file; D mislabels \texttt{.pptm},
  which is macro-enabled.
\item \textbf{D.} Blank carries no placeholders. Title and Content (A), Two
  Content (B), and Comparison (C) all supply ready placeholders.
\item \textbf{D.} Handout page counts (A), ruled 3-per-page lines (B), and
  the common six-per-page choice (C) are all true. Hidden reminders shown in
  Presenter View are speaker notes, not handouts, so D is the incorrect
  statement.
\item \textbf{B.} Statements I and II are correct: Normal is the default
  editing view and Slide Sorter reorders thumbnails. III is false because
  the top-level controlling slide is the Slide Master, not Reading View. A
  omits II; C and D include false III.
\item \textbf{A.} F5 starts from the beginning, Shift+F5 from the current
  slide, Esc ends the show, and B blanks to a black screen. B, C, and D
  scramble these pairings.

\item \textbf{A.} Presentation pages are slides. B, C, and D are hardware.
\item \textbf{A.} Slide Sorter shows thumbnails. B, C, and D are unrelated views.
\item \textbf{A.} .pptx is PowerPoint's format. B, C, and D belong to Word, Excel, or audio.
\item \textbf{A, B, C.} Text, images, and charts fill slides. Engine oil (D) is not slide content.
\item \textbf{A, B, C.} Slide contents, Sorter use, and F5 are correct. D is false: slides are digital pages.
\item \textbf{A.} A layout arranges slide placeholders. The other options concern file management, Excel, or email.

\item \textbf{A.} Slide Master controls shared design and layout elements. The alternatives are Excel/Word features.

\item \textbf{A.} Placeholders reserve positions for slide content. They do not store an OS, credentials, or run database queries.

\item \textbf{A.} Slide Sorter presents thumbnails for rearranging slides. Notes Page emphasizes speaker notes; Reading View presents the show; Outline view focuses on text.

\item \textbf{A, B, C, D.} All four are PowerPoint views used for editing, arranging, notes, or shared design.

\end{enumerate}
\answersubheading{32}{PowerPoint Animation and Transition}
\begin{enumerate}[label=\textbf{Q\arabic*.}, leftmargin=*]
\item \textbf{B.} An animation is applied to an object or text on a slide. The
  whole-presentation option (A) describes a transition applied to all slides;
  C describes a transition; D is unrelated.
\item \textbf{B.} A transition is applied to the movement from one slide to the
  next. A describes an animation; C is the notes pane; D is the spelling tool.
\item \textbf{A.} F5 starts the slide show from the beginning. F7 (B) is spell
  check; F9 (C) is not the slide-show shortcut; Shift+F5 (D) starts from the
  current slide.
\item \textbf{C.} Shift+F5 starts from the current slide. F5 (A) starts from
  the beginning; Ctrl+F5 (B) and Alt+F5 (D) are not the current-slide shortcut
  (Alt+F5 opens Presenter View).
\item \textbf{B.} Alt+F5 starts the presentation in Presenter View. Esc (A)
  ends the show; spelling is on F7, not Alt+F5 (C); opening notes is not the
  Alt+F5 action (D).
\item \textbf{B.} Esc ends a running slide show. F5 (A) starts it; Enter (C)
  advances to the next slide; Tab (D) does not end the show.
\item \textbf{C.} Transition is not an animation category; it is the effect
  between slides. Entrance (A), Exit (B), and Emphasis (D) are the animation
  effect categories.
\item \textbf{D.} Motion Path moves an object along a defined path. Entrance
  (A) makes an object appear; Exit (B) makes it leave; Emphasis (C) draws
  attention without moving it.
\item \textbf{A.} Entrance makes an object appear. Exit (B) makes it leave;
  Emphasis (C) highlights an object already visible; Motion Path (D) moves it
  along a path.
\item \textbf{C.} Emphasis draws attention to an object that stays on the
  slide. Entrance (A) and Exit (B) add or remove the object; Transition (D) is
  between slides, not on an object.
\item \textbf{B.} Morph smoothly moves and animates matching objects across
  slides. Fade (A), Push (C), and Wipe (D) are ordinary transitions that do not
  match objects across slides.
\item \textbf{B.} The Notes pane stores speaker notes for the presenter.
  Printing several slides per page is a handout (A); animation is set on the
  Animations tab (C); theme colours are changed from the Design tab (D).
\item \textbf{B.} A handout gives the audience printed slides, usually several
  per page. Transitions are controlled on the Transitions tab (A); starting the
  show is F5 (C); the Slide Master is editing background design, not a handout
  (D).
\item \textbf{A.} Slide Sorter view shows thumbnails of all slides for
  reordering. Reading View (B) reviews the show in a window; Notes Page (C)
  pairs a slide with notes; Slide Show view (D) presents full screen.
\item \textbf{D.} Formula View is not a PowerPoint view. Normal (A), Slide
  Sorter (B), and Reading (C) views all exist in PowerPoint.
\item \textbf{D.} All three statements are correct: animation applies to objects
  on a slide, a transition applies between slides, and Shift+F5 starts from the
  current slide. A, B, and C each omit at least one true statement.
\item \textbf{B.} Converting SmartArt to text turns it into a plain bulleted
  list, so the graphical layout is lost and must be re-applied. The design is
  not preserved exactly (A); the file name (C) and transitions (D) are
  unaffected by the conversion.
\item \textbf{B.} Pressing B turns the screen black during a running show
  (W turns it white). Esc (A) ends the show; P or the left arrow (C) goes
  back; Presenter View (D) is started with Alt+F5 before presenting.
\item \textbf{C.} Esc correctly ends the slide show. F5 (A) starts the show,
  not spell check; F7 (B) is spell check, not the show; Alt+F5 (D) opens
  Presenter View, not end the show.
\item \textbf{A.} Morph is a transition that works between similar slides. The
  Notes pane is for speaker notes, not animation (B); Motion Path is an
  animation, not a transition (C); a handout is a printed audience layout, not
  a way to start the show (D).

\item \textbf{A.} Transitions bridge consecutive slides. B, C, and D pair unrelated hardware.
\item \textbf{A.} Animations move slide objects. B, C, and D are physical parts without animation settings.
\item \textbf{A.} Shift+F5 starts from the current slide. B checks spelling; C renames; D saves.
\item \textbf{A, B, C.} Fade, Fly In, and Wipe are effects. A paper jam (D) is a printer fault.
\item \textbf{A, B, C.} The transition/animation split and timing control are correct. D is false: effects do not print.
\item \textbf{A.} Transitions apply between slides. Animations apply to objects within a slide; the other actions are unrelated.

\item \textbf{A.} F5 starts from the beginning. Shift+F5 starts from the current slide, F7 commonly checks spelling, and Ctrl+S saves. The cited key is provisional.

\item \textbf{A.} Shift+F5 starts at the current slide. F5 starts at the beginning; Ctrl+P prints; Alt+F4 closes the active window.

\item \textbf{A.} Animations affect objects on a slide; transitions govern movement between slides. The other choices are unrelated.

\item \textbf{A, B, C.} The first three describe standard functions. Animation and transition refer to different scopes.

\end{enumerate}
\answersubheading{33}{PDF and File Formats}
\begin{enumerate}[label=\textbf{Q\arabic*.}, leftmargin=*]
\item \textbf{A.} PDF stands for Portable Document Format. B, C, and D are
  not expansions of the abbreviation.
\item \textbf{B.} PDF is a fixed-layout format that preserves fonts, images,
  and positions, so it looks the same on every device and printer. CSV (A)
  and TXT (C) are plain-text formats; DOCX (D) is re-flowable and its layout
  can shift with the device or font.
\item \textbf{B.} The suffix \texttt{png} is the file-name extension, while
  PNG is the actual format. A reverses the two; C mixes up folders and
  drives; D calls the second term a version rather than a format.
\item \textbf{A.} A scanned PDF stores each page as an image, so its text is
  not directly selectable or searchable; OCR is needed. B is false for a
  scan; C and D describe other data types, not a scanned page.
\item \textbf{B.} OCR stands for Optical Character Recognition. A, C, and D
  are not the correct full forms.
\item \textbf{B.} OCR converts the picture of the text into machine-readable,
  editable text. Compression (A) is a different process; password protection
  (C) is encryption; and format conversion from spreadsheet to slide show (D)
  is not OCR.
\item \textbf{C.} DOCX is the default word-processing format of MS Word.
  XLSX (A) is an Excel workbook; PPTX (B) is a PowerPoint presentation; CSV
  (D) is plain-text comma-separated data.
\item \textbf{A.} \texttt{.xlsx} is a Microsoft Excel workbook. PPTX (B) is
  PowerPoint; DOCX (C) is Word; PNG (D) is an image format.
\item \textbf{B.} CSV stands for Comma-Separated Values. A, C, and D are not
  the correct full forms.
\item \textbf{C.} \texttt{.txt} is a plain-text format. DOCX (A) and PPTX (B)
  are formatted Office formats; PDF (D) is a fixed-layout format.
\item \textbf{D.} \texttt{.docx} is a Word document, not an image extension.
  PNG (A), JPG (B), and GIF (C) are all image file extensions.
\item \textbf{A.} Encryption provides confidentiality, so only an authorised
  reader with the correct password or key can see the content. Authenticity
  (B) and integrity (C) are the goals of a digital signature; compression (D)
  is unrelated to password protection.
\item \textbf{B.} A digital signature provides authenticity (the signer is
  verified) and integrity (changes are detected). Confidentiality (A) is the
  goal of encryption; lossless compression (C) is a property of PNG, not a
  signature; OCR (D) is a separate text-recognition process.
\item \textbf{B.} A digital signature verifies the signer and detects any
  change made after signing. A wrongly says the signature encrypts the file;
  C confuses it with OCR; D describes redaction.
\item \textbf{A.} Statements I and II are correct: encryption gives
  confidentiality, and a digital signature gives authenticity and integrity.
  III is false because a digital signature does not encrypt the file. B and C
  include the false III; D includes all three.
\item \textbf{B.} Redaction permanently removes sensitive content before the
  file is shared. Compression (A) only reduces size; a digital signature (C)
  proves origin and integrity; PDF/A (D) is an archiving format.
\item \textbf{B.} PDF/A is a standardised PDF version for long-term
  archiving, so the document renders the same in the future. Editing (A),
  scanning (C), and photo compression (D) are not its purpose.
\item \textbf{A.} PNG commonly uses lossless compression and supports
  transparency. JPEG (B) is lossy; BMP (C) is usually raw and larger; TXT (D)
  is not an image format.
\item \textbf{C.} PPTX is a Microsoft PowerPoint presentation, so this pairing
  is correct. XLSX (A) is an Excel workbook, not PowerPoint; DOCX (B) is a
  Word document, not Excel; CSV (D) is a plain-text data format, not a
  fixed-layout document.
\item \textbf{A.} The extension is the suffix in the file name (for example
  \texttt{.png}), while the format is the internal structure of the data
  (for example PNG). B reverses the two; C wrongly says they are identical;
  D wrongly ties the extension to the folder.

\item \textbf{A.} PDF is Portable Document Format. B, C, and D are not the recognised expansion.
\item \textbf{A.} Reader/Acrobat opens PDFs. B, C, and D are unrelated tools.
\item \textbf{A.} PDF fixes layout everywhere. B, C, and D are not document formats.
\item \textbf{A, B, C.} PDF, DOCX, and TXT are real formats. D is not a file format.
\item \textbf{A, B, C.} Layout fidelity, PDF contents, and DOCX editability are correct. D is false: readers open PDFs on computers.
\item \textbf{A.} .docx is Word, .xlsx is Excel, .pptx is PowerPoint, and .accdb is Access.

\item \textbf{A.} PDF is a page-description/document format intended to preserve layout. It is not a spreadsheet-only or executable format.

\item \textbf{A.} CSV stores delimited text rows and columns. PPTX is a presentation package, PNG is an image format, and EXE is an executable extension.

\item \textbf{A.} PNG is a raster image format. DOCX, XLSX, and PPTX are Office document formats.

\item \textbf{A, B, C.} The first three pairings are correct. PNG is an image format, not a database file.

\end{enumerate}
\answersubheading{54}{Office Shortcut and Product Trap Lab}
\begin{enumerate}[label=\textbf{Q\arabic*.}, leftmargin=*]
\item \textbf{B.} In Word, F7 opens Spelling \& Grammar. A is the nearest
  neighbour: F1 opens Help. C is Shift+F7 (Thesaurus), and D is Ctrl+H (Find
  and Replace).
\item \textbf{B.} Ctrl+H opens Find and Replace, so text can be changed. A is
  Ctrl+F (find only), C is Ctrl+K (Insert Hyperlink), and D is F7 (Spelling
  \& Grammar).
\item \textbf{B.} Ctrl+K inserts a hyperlink. A is Ctrl+H (Find and Replace),
  C is Ctrl+L (align left), and D is Ctrl+J (justify).
\item \textbf{B.} Ctrl+X cuts the selection and Ctrl+V pastes it, so together
  they cut and paste. A is copy (Ctrl+C then Ctrl+V); C is undo
  (Ctrl+Z) and redo (Ctrl+Y); D is Ctrl+H (Find and Replace).
\item \textbf{C.} Print Layout shows the document exactly as it will appear on
  the printed page. Web Layout (A), Outline (B), and Draft (D) do not show the
  true printed appearance.
\item \textbf{B.} Track Changes records insertions, deletions, and formatting
  edits so a reviewer can accept or reject them (PYQ-12). A is page setup, C
  is the spell-check tool, and D is picture compression --- none is Track
  Changes.
\item \textbf{D.} This is false. Track Changes preserves the original text and
  marks the edit; the change is applied only when a reviewer accepts it. A, B,
  and C are all true of Track Changes.
\item \textbf{B.} A Section Break creates a new section, which is what allows a
  different header or footer on the following pages (PYQ-13, TRAP-13). A Page
  Break (A) keeps the same section and the same header/footer; column and line
  breaks (C, D) do not create sections.
\item \textbf{C.} Combining a main document with an address database is Mail
  Merge (PYQ-37, TRAP-29). Track Changes (A) records edits, AutoCorrect (B)
  fixes typing, and Find and Replace (D) substitutes text.
\item \textbf{A.} Alt+Shift+C removes a document window split (PYQ-46);
  Ctrl+Alt+S also toggles the split in Word. B, C, and D are not the removal
  shortcut.
\item \textbf{A.} F5 starts a PowerPoint slide show from the beginning
  (PYQ-40, TRAP-32). F7 (B) checks spelling, F1 (C) opens Help, and F2 (D)
  renames or edits a cell.
\item \textbf{B.} Shift+F5 starts the show from the current slide. A is Esc
  (end the show), C is Alt+F5 (Presenter View), and D is Ctrl+M (new slide).
\item \textbf{B.} Alt+F5 starts the presentation in Presenter View, showing
  speaker notes to the presenter. A is Slide Sorter, C is Reading view, and D
  is Slide Master --- different views and tools.
\item \textbf{A.} Esc ends a running slide show. Enter (B) advances to the next
  slide, F5 (C) starts the show, and Ctrl+H (D) hides the pointer during a
  show (in Word it would open Find and Replace).
\item \textbf{C.} A table of data cannot be converted directly to SmartArt; the
  command is disabled for it, and its text must be copied into a text box
  first (PYQ-17). Bulleted lists (A), blocks of text (B), and indented text
  (D) convert directly.
\item \textbf{B.} The Handouts layout prints several slides on one page for the
  audience. A describes editing notes, C describes the Slide Master, and D
  describes the orientation setting.
\item \textbf{B.} The Slide Master controls the design, fonts, colours, and
  layout of all slides that use it. A is false --- it affects all dependent
  slides, not only the title; C and D name unrelated settings.
\item \textbf{B.} Orientation switches between Portrait and Landscape. A names
  Excel reference types, C names Word views, and D names PowerPoint views.
\item \textbf{B.} Both the column and the row are locked in \$A\$1, so it is an
  absolute reference. A is the relative form A1; C describes a mixed reference
  with the row fixed (A\$1); D is a different Excel feature.
\item \textbf{B.} Statements I and III are correct: MS-DOS is an operating
  system, and MS Word is a word processor. II is false because MS Access is a
  database management system, not a spreadsheet. A omits III; C includes the
  false II; D includes the false II.

\item \textbf{A.} Ctrl+A selects all. B saves, C prints, and D creates new documents.
\item \textbf{A.} Ctrl+F finds text. B bolds, C underlines, and D centres text.
\item \textbf{A.} Ctrl+N opens a new document. B, C, and D print, save, or paste.
\item \textbf{A, B, C.} New, open, and save pairs are correct. D is false: Ctrl+X cuts.
\item \textbf{A, B, C.} The application matches are correct. D is false: a spreadsheet is not a presentation tool.
\item \textbf{A.} Ctrl+P opens Print. Ctrl+H opens Find/Replace, Ctrl+K inserts a hyperlink, and Ctrl+Z undoes.

\item \textbf{A.} Ctrl+K opens the hyperlink action in common Office versions. The other shortcuts have different or application-dependent uses.

\item \textbf{A.} Shift+F5 starts from the current slide; F5 starts from the beginning. F7 and Ctrl+P are not that action.

\item \textbf{A.} F7 invokes spelling/grammar review in common Word versions. F1 is Help, F5 is navigation/go-to, and F12 opens Save As.

\item \textbf{A, B, C.} Ctrl+S saves, Ctrl+P prints, Ctrl+Z undoes, and Ctrl+C copies. Ctrl+V pastes.

\end{enumerate}
\answersubheading{55}{File Formats and Named Technology Associations}
\begin{enumerate}[label=\textbf{Q\arabic*.}, leftmargin=*]
\item \textbf{A.} PNG stands for Portable Network Graphics. B, C, and D are
  not the recognised expansion of the abbreviation.
\item \textbf{A.} PNG is an image format. It is not a document format (B), a
  database format (C), or a spreadsheet format (D).
\item \textbf{B.} PDF stands for Portable Document Format. A, C, and D are
  not the correct full forms.
\item \textbf{A.} PDF is a fixed-layout format, so a document keeps the same
  layout on every device. It does not hold spreadsheet formulas (B), a
  relational database (C), or lossy audio (D).
\item \textbf{B.} Adobe Acrobat Reader is a PDF reader. Firefox (A) is a web
  browser, TeamViewer (C) is a remote-support tool, and Excel (D) is a
  spreadsheet program --- none is primarily a PDF reader.
\item \textbf{C.} Firefox is developed by Mozilla. Microsoft (A) makes Edge,
  Adobe (B) made PDF and Acrobat Reader, and Google (D) makes Chrome.
\item \textbf{A.} Firefox is a web browser. It is not a PDF reader (B), a
  spreadsheet program (C), or a database program (D).
\item \textbf{B.} PARAM is the family of supercomputers developed by C-DAC.
  ISRO (A) is the space agency, Microsoft (C) is a software company, and
  Mozilla (D) makes Firefox --- none is associated with PARAM.
\item \textbf{B.} C-DAC = Centre for Development of Advanced Computing. A, C,
  and D are not the correct expansion.
\item \textbf{B.} Microsoft Excel uses \texttt{.xlsx}. \texttt{.docx} (A) is
  Word, \texttt{.pptx} (C) is PowerPoint, and \texttt{.accdb} (D) is Access.
\item \textbf{A.} Microsoft Word uses \texttt{.docx}. \texttt{.xlsx} (B) is
  Excel, \texttt{.pptx} (C) is PowerPoint, and \texttt{.txt} (D) is plain
  text.
\item \textbf{C.} \texttt{.accdb} is the newer Microsoft Access database
  format. It is not an Excel (A), Word (B), or PowerPoint (D) file.
\item \textbf{B.} CSV stands for Comma-Separated Values. A, C, and D are not
  the correct expansion of the abbreviation.
\item \textbf{A.} A CSV file stores a table as plain text with comma
  separators. It is not a compressed page image (B), a keyed database (C), or
  a formatted document (D).
\item \textbf{D.} CSV is a plain-text, comma-separated table format, not an
  image format. PNG (A), JPEG (B), and GIF (C) are all image formats.
\item \textbf{B.} Statements I and II are correct: PNG is an image format,
  and PDF preserves layout across devices. III is false because
  \texttt{.accdb} is a Microsoft Access database file, not a Word file.
\item \textbf{A.} \texttt{.mdb} is the older Microsoft Access database
  format. It is not a Word document (B), a PDF (C), or an image file (D).
\item \textbf{A.} PNG is an image format (i-b); PARAM is a C-DAC supercomputer
  (ii-d); Firefox is a Mozilla web browser (iii-a); \texttt{.pptx} is a
  PowerPoint presentation (iv-c). B scrambles all four; C reverses the
  associations; D wrongly assigns PARAM to the presentation and
  \texttt{.pptx} to the supercomputer.
\item \textbf{C.} TeamViewer is a remote-support tool. Adobe Acrobat Reader
  (A) is a PDF reader, Firefox (B) is a web browser, and PowerPoint (D) is a
  presentation program.
\item \textbf{A.} \texttt{.png} is the extension (the letters after the dot)
  and PNG (Portable Network Graphics) is the format. B and D reverse the two;
  C wrongly treats the extension and the format as identical.

\item \textbf{A.} .docx is Word's format. B, C, and D belong to Excel, PowerPoint, or audio.
\item \textbf{A.} .xlsx is Excel's format. B and C belong to Word and PowerPoint; D is not a format.
\item \textbf{A.} JPEG suits colour photos. B, C, and D are not image formats.
\item \textbf{A, B, C.} JPG, PNG, and GIF are images. DOCX (D) is a Word document.
\item \textbf{A, B, C.} The audio, video, and PDF matches are correct. D is false: .xlsx is a spreadsheet.
\item \textbf{C.} A USB drive uses flash memory. HDDs use magnetic storage, RAM is working memory, and Blu-ray is optical.

\item \textbf{A.} .pptx is PowerPoint; .xlsx is Excel, .docx is Word, and .accdb is Access.

\item \textbf{A.} The dot-suffix in a filename is an extension; PNG names the format. Extensions can be misleading and do not guarantee contents.

\item \textbf{A.} CSV stores delimited text data. XLSX is a workbook format; PPTX and DOCX are presentation and word-processing formats.

\item \textbf{A, B, C.} The first three are correct. PNG is an image format, not a database format.

\end{enumerate}
\answersubheading{61}{Current Office, Email and Statement Evaluation}
\begin{enumerate}[label=\textbf{Q\arabic*.}, leftmargin=*]
\item \textbf{A.} Content controls are inserted from the Developer tab (enabled
  via Word Options, Customize Ribbon). The Review tab (B) holds Track Changes
  and comments; Mailings (C) holds Mail Merge; References (D) holds the Table
  of Contents and citations.
\item \textbf{C.} Comments are read in the markup/comments pane, not directly
  in the Navigation Pane, so C is the NOT-correct statement. A is true (a
  drop-down list or date picker restricts entry); B is true (hanging indent
  shifts all lines except the first); D is true (accepting keeps the text and
  strips that change's metadata).
\item \textbf{C.} A first-line indent shifts only the first line of the
  paragraph. A reverses the two indents; B describes a first-line indent while
  calling it hanging; D describes no standard indent.
\item \textbf{A.} The Navigation Pane shows Headings, Pages, and Results
  (search matches). B lists review/markup items, not pane views; C lists Word
  document views; D lists ribbon tabs.
\item \textbf{B.} Each tracked change records insertions, deletions, and
  formatting changes together with the author's name and the date and time. A
  omits the author and change type; C wrongly claims no metadata is stored; D
  names printer settings, which are unrelated.
\item \textbf{C.} Headers, footers, margins, and orientation belong to a
  section, so only a Section Break creates a section that can carry its own
  header. A Page Break stays in the same section; column (B) and line (D)
  breaks do not start new sections.
\item \textbf{B.} 30 is not greater than 50 (first test FALSE), but 30 is
  greater than 25 (second test TRUE), so the nested IF returns Medium. A needs
  A1 above 50; C needs both tests FALSE; D is not a value this formula
  produces.
\item \textbf{C.} A blank cell is treated as 0 in a numeric comparison, so
  both tests are FALSE and the formula falls through to Low. A invents a value
  of 100; B wrongly claims the first test is skipped; D is wrong because a
  blank does not produce an error here.
\item \textbf{C.} Text is not a number, so arithmetic on it fails with
  \texttt{\#VALUE!}. A is wrong (no numeric result exists); B is wrong
  (ISNUMBER returns FALSE for text); D overstates the rule (text is not
  silently zero in arithmetic).
\item \textbf{B.} COUNTIFS takes paired arguments: range1, criteria1, range2,
  criteria2, and so on, with AND logic across pairs. A swaps the pair order;
  C invents an addition syntax; D invents a semicolon/TRUE syntax.
\item \textbf{B.} Multiplying the two boolean arrays gives 1 only where both
  conditions hold, and SUMPRODUCT sums those products, hence counts the rows.
  A misuses SUM over comparisons; C averages instead of counting; D counts
  non-empty cells without applying the conditions.
\item \textbf{A.} I is true (COUNTIFS uses AND logic over its pairs) and III
  is true (blank reads as 0). II is false: SUMPRODUCT multiplies arrays and
  can count multi-criterion rows as well as add. B keeps only the false
  statement; C and D both include false II.
\item \textbf{A.} A rule pairs a condition (sender, subject, keywords) with an
  action (move, flag, forward, delete), applied automatically. B mixes in
  security features; C mixes in recipient fields; D mixes in Word breaks.
\item \textbf{C.} A digital signature gives authenticity and integrity (who
  sent it, unaltered), while encryption gives confidentiality (only the
  intended reader can read it). A wrongly claims the signature encrypts; B
  swaps the two roles; D wrongly calls them the same operation.
\item \textbf{B.} POP3 stands for Post Office Protocol version 3. A, C, and D
  are not the correct expansion of the abbreviation.
\item \textbf{C.} SMTP sends mail, POP3 downloads it to the local device
  (typically removing it from the server), and IMAP leaves mail on the server
  and synchronises devices. A swaps POP3 and SMTP; B swaps POP3 and IMAP; D
  wrongly claims POP3 syncs on the server.
\item \textbf{B.} The sender always sees the full recipient list; Bcc addresses
  are hidden from the other recipients (To, Cc, and fellow Bcc). A wrongly
  makes Bcc visible; C wrongly hides Bcc from the sender; D wrongly claims the
  addresses are deleted.
\item \textbf{B.} Convert to SmartArt works directly on text (bulleted,
  numbered, or levelled lists in A, C, D). A picture has no text levels to
  convert, so it cannot become SmartArt in one step without re-entry.
\item \textbf{A.} Only I is correct: a Section Break enables different headers
  and footers. II is false (Bcc is hidden from the other recipients, never
  from the sender). III is false (a signature verifies sender and integrity;
  encryption provides confidentiality). B, C, and D each include a false
  statement.
\item \textbf{C.} A first-line indent shifts only the first line and a hanging
  indent shifts every line except the first. A assigns section behaviour to a
  Page Break; B misstates both COUNTIFS (it handles several criteria) and
  SUMPRODUCT (it can count as well as add); D swaps POP3 and IMAP.

\item \textbf{A.} Ctrl+I italicises. B prints, C saves, and D opens new documents.
\item \textbf{A.} Bcc hides recipients from each other. B, C, and D describe unrelated effects.
\item \textbf{A.} Unknown attachments need caution. B is dangerously false; C and D are nonsense effects.
\item \textbf{A, B, C.} The DOCX, XLSX, and PPTX matches are correct. D is false: videos use media formats.
\item \textbf{A.} I and II state the POP3/IMAP roles correctly, while III is false because SMTP sends mail; it does not print. B, C, and D wrongly include false III.
\item \textbf{A.} Page setup is section-specific. AutoCorrect is application-wide; styles and Properties are not unique page-layout settings for one section.

\item \textbf{C.} COUNTIF counts matching cells. COUNT counts numeric cells without a criterion; COUNTBLANK counts empties; SUMIF sums matching values.

\item \textbf{C.} Bcc conceals those recipients from other recipients. To and Cc recipients are visible; Subject identifies the topic.

\item \textbf{A.} Track Changes records revisions. Mail Merge personalizes documents; AutoSum is an Excel feature; Slide Sorter is a PowerPoint view.

\item \textbf{A, B, C.} Section setup, Bcc, and COUNTIF are correctly described. Digital signatures support authenticity/integrity; encryption supports confidentiality.

\end{enumerate}
\unitopener{8}{Internet, Web \& Email}
\lessonheading{34}{Internet, WWW and Browser}
\begin{enumerate}[label=\textbf{Q\arabic*.}, leftmargin=*]
\item \textit{[Original]} Which option correctly describes the Internet?
  \begin{enumerate}[label=\Alph*.]
    \item A single webpage stored on one computer
    \item A global network of connected computers and networks
    \item A software used to open webpages
    \item A website that finds other webpages
  \end{enumerate}
\item \textit{[Original]} The World Wide Web (WWW) is:
  \begin{enumerate}[label=\Alph*.]
    \item the same thing as the Internet
    \item a service running on the Internet, consisting of linked web pages
      and sites
    \item the company that provides the Internet connection
    \item a device that requests pages from a server
  \end{enumerate}
\item \textit{[Original]} Which of the following correctly matches each item
  to its description?
  \begin{enumerate}[label=\Alph*.]
    \item Webpage --- the whole collection of pages; Website --- one single
      document
    \item Website --- a collection of related webpages; Webpage --- a single
      document; Homepage --- the main or starting page
    \item Homepage --- the whole collection; Webpage --- the connection
      provider
    \item Website --- clickable text; Webpage --- the page-requesting device
  \end{enumerate}
\item \textit{[Original]} A browser is:
  \begin{enumerate}[label=\Alph*.]
    \item software used to open and view websites and webpages
    \item a website that finds pages for typed keywords
    \item the company that provides the Internet connection
    \item a computer that supplies pages to clients
  \end{enumerate}
\item \textit{[Original]} Which of the following is a web browser?
  \begin{enumerate}[label=\Alph*.]
    \item Google
    \item Bing
    \item Mozilla Firefox
    \item Yahoo
  \end{enumerate}
\item \textit{[Original]} Which of the following is a search engine?
  \begin{enumerate}[label=\Alph*.]
    \item Google Chrome
    \item Mozilla Firefox
    \item Microsoft Edge
    \item Bing
  \end{enumerate}
\item \textit{[Original]} Which statement correctly distinguishes a browser
  from a search engine?
  \begin{enumerate}[label=\Alph*.]
    \item A browser is a website and a search engine is software.
    \item A browser is software that opens pages; a search engine is a
      website that finds pages for keywords.
    \item A browser provides the Internet connection; a search engine stores
      websites.
    \item A browser and a search engine are the same thing.
  \end{enumerate}
\item \textit{[Original]} In the client-server model of the Web, the client:
  \begin{enumerate}[label=\Alph*.]
    \item supplies webpages to browsers
    \item provides the Internet connection to homes
    \item requests a service or resource, such as a webpage
    \item stores websites for many users
  \end{enumerate}
\item \textit{[Original]} A web server:
  \begin{enumerate}[label=\Alph*.]
    \item requests pages from a browser
    \item stores websites and sends pages to the clients that ask for them
    \item provides the Internet connection to offices
    \item finds pages matching typed keywords
  \end{enumerate}
\item \textit{[PYQ WO 2026 Q74, key:provisional]} Expand ISP.
  \begin{enumerate}[label=\Alph*.]
    \item Internet Service Provider
    \item Intranet Software Protocol
    \item International Server Program
    \item Internet System Page
  \end{enumerate}
\item \textit{[Original]} An ISP is needed because it:
  \begin{enumerate}[label=\Alph*.]
    \item displays webpages on the screen
    \item provides the Internet connection to homes and offices
    \item jumps to another page when clicked
    \item writes webpages in HTML
  \end{enumerate}
\item \textit{[Original]} A hyperlink is:
  \begin{enumerate}[label=\Alph*.]
    \item the address of a webpage
    \item the company providing the Internet connection
    \item clickable text or an image that jumps to another page or section
    \item software that opens websites
  \end{enumerate}
\item \textit{[Original]} Expand WWW.
  \begin{enumerate}[label=\Alph*.]
    \item World Web Window
    \item Wide World Web
    \item World Wide Web
    \item Web World Wide
  \end{enumerate}
\item \textit{[Original]} Expand URL.
  \begin{enumerate}[label=\Alph*.]
    \item Uniform Resource Locator
    \item Universal Register Loader
    \item Unified Request Link
    \item Uniform Revision Log
  \end{enumerate}
\item \textit{[Original]} Expand HTTP.
  \begin{enumerate}[label=\Alph*.]
    \item HighText Transfer Program
    \item HyperText Transfer Protocol
    \item Home Tool Transfer Page
    \item Hyperlink Text Transport
  \end{enumerate}
\item \textit{[Original]} Which of the following is \textbf{NOT} true about
  the Web and its access tools?
  \begin{enumerate}[label=\Alph*.]
    \item A browser opens and displays webpages.
    \item A search engine finds pages matching typed keywords.
    \item The Internet and the World Wide Web are exactly the same thing.
    \item An ISP provides the Internet connection.
  \end{enumerate}
\item \textit{[Original]} The homepage of a website is:
  \begin{enumerate}[label=\Alph*.]
    \item the company that gives the Internet connection
    \item the main or starting page of the website
    \item every page of the website taken together
    \item the clickable text that jumps to another site
  \end{enumerate}
\item \textit{[Original]} When a user types a site address and presses Enter,
  the correct order of events is:
  \begin{enumerate}[label=\Alph*.]
    \item server displays the page, browser requests it, ISP writes it
    \item browser (client) requests the page, server sends it back, browser
      displays it
    \item ISP requests the page, hyperlink supplies it, browser stores it
    \item search engine provides the connection, server clicks the link
  \end{enumerate}
\item \textit{[Original]} Evaluate the statements:
  \begin{enumerate}[label=\Roman*.]
    \item The Internet is a global network of networks; the WWW is one
      service running on it.
    \item Google Chrome is a browser and Google (search) is a search engine.
    \item A server requests pages while a client supplies them.
  \end{enumerate}
  \begin{enumerate}[label=\Alph*.]
    \item I only
    \item I and II only
    \item II and III only
    \item I, II, and III
  \end{enumerate}
\item \textit{[Original]} Which pairing is correctly matched?
  \begin{enumerate}[label=\Alph*.]
    \item Browser --- Google (search)
    \item Search engine --- Mozilla Firefox
    \item ISP --- company providing the Internet connection
    \item Hyperlink --- address bar software that opens pages
  \end{enumerate}

\item \textit{[JKSSB FAA-pattern]} The Internet is best described as:
  \begin{enumerate}[label=\Alph*.]
    \item a global network of interconnected networks.
    \item a single office printer.
    \item one personal diary.
    \item a television channel.
  \end{enumerate}
\item \textit{[JKPSI-pattern]} Which of the following is a web browser?
  \begin{enumerate}[label=\Alph*.]
    \item Mozilla Firefox
    \item MS Paint
    \item Calculator
    \item Notepad paper
  \end{enumerate}
\item \textit{[SSC CGL-pattern]} WWW stands for:
  \begin{enumerate}[label=\Alph*.]
    \item World Wide Web
    \item World Wrestling Watch
    \item Wide Window Wood
    \item Web Water Well
  \end{enumerate}
\item \textit{[JKSSB-pattern]} \textbf{(MSQ: select ALL correct options.)} Which of the following are web browsers?
  \begin{enumerate}[label=\Alph*.]
    \item Google Chrome
    \item Mozilla Firefox
    \item Microsoft Edge
    \item MS Excel
  \end{enumerate}
\item \textit{[SSC CGL-pattern]} \textbf{(MSQ: select ALL correct options.)} Which statements about the Internet and browsers are correct?
  \begin{enumerate}[label=\Alph*.]
    \item The Internet links networks worldwide.
    \item A browser displays web pages.
    \item An ISP provides Internet access.
    \item A browser is a hardware device.
  \end{enumerate}
\item \textit{[Original]} The World Wide Web is best described as:
  \begin{enumerate}[label=\Alph*.]
\item A system of interlinked resources accessed over the Internet
   \item The physical cables of the Internet
   \item A single search engine
   \item A computer's operating system
  \end{enumerate}

\item \textit{[Original]} Which application retrieves and displays web pages?
  \begin{enumerate}[label=\Alph*.]
\item Web browser
   \item Search engine
   \item Spreadsheet
   \item Mail server
  \end{enumerate}

\item \textit{[Original]} A search engine primarily helps users:
  \begin{enumerate}[label=\Alph*.]
\item Find indexed web content
   \item Route packets at the network layer
   \item Edit local spreadsheets
   \item Store BIOS firmware
  \end{enumerate}

\item \textit{[Original]} A single document available on the Web is a:
  \begin{enumerate}[label=\Alph*.]
\item Web page
   \item Website
   \item Browser
   \item Domain registrar
  \end{enumerate}

\item \textit{[Original, MSQ drill]} \textbf{(MSQ: select ALL correct options.)} Which statements are correct?
  \begin{enumerate}[label=\Alph*.]
\item The Internet is a global network of networks.
   \item The Web uses Internet infrastructure.
   \item A browser and a search engine are different tools.
   \item Every Internet service is a web page.
  \end{enumerate}

\end{enumerate}
\lessonheading{35}{URL, DNS, HTTP and HTTPS}
\begin{enumerate}[label=\textbf{Q\arabic*.}, leftmargin=*]
\item \textit{[Original]} Expand URL.
  \begin{enumerate}[label=\Alph*.]
    \item Uniform Resource Locator
    \item Universal Record Link
    \item Unified Reference Locator
    \item Uniform Retrieval Language
  \end{enumerate}
\item \textit{[Original]} Expand DNS.
  \begin{enumerate}[label=\Alph*.]
    \item Domain Name System
    \item Data Network Service
    \item Digital Node Scheme
    \item Distributed Number Sequence
  \end{enumerate}
\item \textit{[Original]} Expand HTTPS.
  \begin{enumerate}[label=\Alph*.]
    \item Hypertext Transfer Protocol Secure
    \item High Transfer Text Protocol System
    \item Hyperlink Text Transmission Protocol Standard
    \item Host Transfer Protocol Service
  \end{enumerate}
\item \textit{[Original]} In the URL
  \texttt{https://www.example.in:443/study/computer?topic=dns}, which part
  is the scheme?
  \begin{enumerate}[label=\Alph*.]
    \item \texttt{https}
    \item \texttt{www.example.in}
    \item \texttt{:443}
    \item \texttt{/study/computer}
  \end{enumerate}
\item \textit{[Original]} In \texttt{mail.example.com}, which label is the
  subdomain?
  \begin{enumerate}[label=\Alph*.]
    \item \texttt{.com}
    \item \texttt{example}
    \item \texttt{mail}
    \item \texttt{mailto}
  \end{enumerate}
\item \textit{[Original]} Which of the following is a Top-Level Domain
  (TLD)?
  \begin{enumerate}[label=\Alph*.]
    \item \texttt{www}
    \item \texttt{.in}
    \item \texttt{/study}
    \item \texttt{?topic=dns}
  \end{enumerate}
\item \textit{[Original]} The default port for HTTP is
  \underline{\hspace{1.2cm}}, and for HTTPS it is
  \underline{\hspace{1.2cm}}.
  \begin{enumerate}[label=\Alph*.]
    \item 443; 80
    \item 80; 443
    \item 25; 110
    \item 80; 80
  \end{enumerate}
\item \textit{[Original]} Which part of a URL carries extra parameters such
  as \texttt{?topic=dns\&lang=en}?
  \begin{enumerate}[label=\Alph*.]
    \item Scheme
    \item Domain name
    \item Path
    \item Query
  \end{enumerate}
\item \textit{[Original]} Which statement correctly distinguishes a URL from
  a domain name?
  \begin{enumerate}[label=\Alph*.]
    \item A URL is the full address of one resource; the domain name is only
      the host part inside it.
    \item A domain name is the full address of one resource; the URL is only
      the host part inside it.
    \item A URL and a domain name are identical in every part.
    \item A URL names the server hardware; a domain name names the page
      content.
  \end{enumerate}
\item \textit{[Original]} IPv4 is a \underline{\hspace{1.2cm}} address,
  whereas IPv6 is a \underline{\hspace{1.2cm}} address.
  \begin{enumerate}[label=\Alph*.]
    \item 128-bit; 32-bit
    \item 32-bit; 128-bit
    \item 64-bit; 256-bit
    \item 32-bit; 256-bit
  \end{enumerate}
\item \textit{[Original]} The main function of DNS is to:
  \begin{enumerate}[label=\Alph*.]
    \item translate domain names into IP addresses.
    \item encrypt web pages during transfer.
    \item assign a port number to every browser.
    \item carry e-mail between mail servers.
  \end{enumerate}
\item \textit{[Original]} Which of the following is \textbf{NOT} true of DNS?
  \begin{enumerate}[label=\Alph*.]
    \item It maps domain names to IP addresses.
    \item It uses a distributed hierarchy of servers.
    \item It caches answers to speed up repeat visits.
    \item It encrypts the contents of web pages.
  \end{enumerate}
\item \textit{[Original]} When a domain name is not cached, the resolver
  obtains its IP address by:
  \begin{enumerate}[label=\Alph*.]
    \item asking the DNS hierarchy (root, then TLD, then authoritative
      server) for the address.
    \item guessing the address from the URL path.
    \item reading the page content over HTTP first.
    \item copying the port number into the address field.
  \end{enumerate}
\item \textit{[Original]} HTTP stands for:
  \begin{enumerate}[label=\Alph*.]
    \item Hypertext Transfer Protocol
    \item Hyperlink Transmission Process
    \item Host Text Transfer Program
    \item High-speed Transport Protocol
  \end{enumerate}
\item \textit{[Original]} TLS stands for:
  \begin{enumerate}[label=\Alph*.]
    \item Transport Layer Security
    \item Transfer Link Service
    \item Terminal Login System
    \item Transmission Line Splitter
  \end{enumerate}
\item \textit{[Original]} Which statement about HTTP and HTTPS is correct?
  \begin{enumerate}[label=\Alph*.]
    \item HTTP encrypts pages by default; HTTPS removes that encryption.
    \item HTTPS is HTTP protected by TLS; plain HTTP provides no encryption.
    \item HTTP uses port 443 and HTTPS uses port 80.
    \item HTTPS is a page-design language unrelated to HTTP.
  \end{enumerate}
\item \textit{[Original]} A browser shows a padlock for
  \texttt{https://www.example.in} because:
  \begin{enumerate}[label=\Alph*.]
    \item the TLS handshake authenticated the server and encrypted the
      channel.
    \item DNS refused to resolve the name.
    \item the page was transferred over port 80.
    \item the query string was removed from the URL.
  \end{enumerate}
\item \textit{[Original]} Which sequence gives the correct flow when a user
  types a secure URL and presses Enter?
  \begin{enumerate}[label=\Alph*.]
    \item URL $\rightarrow$ DNS $\rightarrow$ IP $\rightarrow$ server
      $\rightarrow$ TLS handshake $\rightarrow$ page.
    \item IP $\rightarrow$ URL $\rightarrow$ DNS $\rightarrow$ TLS handshake
      $\rightarrow$ server $\rightarrow$ page.
    \item URL $\rightarrow$ IP $\rightarrow$ page $\rightarrow$ DNS
      $\rightarrow$ server $\rightarrow$ TLS handshake.
    \item DNS $\rightarrow$ TLS handshake $\rightarrow$ URL $\rightarrow$ IP
      $\rightarrow$ page $\rightarrow$ server.
  \end{enumerate}
\item \textit{[Original]} Evaluate the statements:
  \begin{enumerate}[label=\Roman*.]
    \item DNS translates domain names into IP addresses.
    \item IPv6 is a 128-bit address.
    \item Plain HTTP encrypts the data it carries.
  \end{enumerate}
  \begin{enumerate}[label=\Alph*.]
    \item I only
    \item I and II only
    \item II and III only
    \item I, II, and III
  \end{enumerate}
\item \textit{[Original]} Which pairing is correctly matched?
  \begin{enumerate}[label=\Alph*.]
    \item Scheme --- \texttt{www.example.in}
    \item Port 80 --- default for HTTPS
    \item Query --- parameters after \texttt{?} in a URL
    \item IPv6 --- 256-bit address
  \end{enumerate}

\item \textit{[JKSSB FAA-pattern]} Expand URL.
  \begin{enumerate}[label=\Alph*.]
    \item Uniform Resource Locator
    \item United Road Line
    \item Universal Radio Link
    \item User Ring Light
  \end{enumerate}
\item \textit{[JKPSI-pattern]} DNS primarily converts:
  \begin{enumerate}[label=\Alph*.]
    \item domain names into IP addresses.
    \item printers into scanners.
    \item speakers into monitors.
    \item paper into ink.
  \end{enumerate}
\item \textit{[SSC CGL-pattern]} The S in HTTPS indicates the connection is:
  \begin{enumerate}[label=\Alph*.]
    \item secured/encrypted.
    \item slower by design.
    \item shared publicly.
    \item switched off.
  \end{enumerate}
\item \textit{[JKSSB-pattern]} \textbf{(MSQ: select ALL correct options.)} Which parts can appear in a URL?
  \begin{enumerate}[label=\Alph*.]
    \item Protocol (http/https)
    \item Domain name
    \item Path to the page
    \item Printer ink level
  \end{enumerate}
\item \textit{[SSC CGL-pattern]} \textbf{(MSQ: select ALL correct options.)} Which statements about URL, DNS, and HTTP(S) are correct?
  \begin{enumerate}[label=\Alph*.]
    \item A URL locates a web resource.
    \item DNS resolves names to IP addresses.
    \item HTTPS encrypts web traffic.
    \item HTTP stands for High Telephone Power.
  \end{enumerate}
\item \textit{[Original]} In a URL, the scheme identifies the:
  \begin{enumerate}[label=\Alph*.]
\item Access method or protocol
   \item User's computer model
   \item Page's font size
   \item Local file extension only
  \end{enumerate}

\item \textit{[Original]} Which service translates a domain name into an IP address?
  \begin{enumerate}[label=\Alph*.]
\item DNS
   \item DHCP
   \item SMTP
   \item FTP
  \end{enumerate}

\item \textit{[Original]} HTTPS adds which security protocol to web communication?
  \begin{enumerate}[label=\Alph*.]
\item TLS
   \item POP3
   \item OMR
   \item RAID
  \end{enumerate}

\item \textit{[Original]} Which component of the URL `example.org` is the top-level domain?
  \begin{enumerate}[label=\Alph*.]
\item .org
   \item example
   \item https
   \item /
  \end{enumerate}

\item \textit{[Original, MSQ drill]} \textbf{(MSQ: select ALL correct options.)} Which statements about web addressing/security are correct?
  \begin{enumerate}[label=\Alph*.]
\item DNS can resolve domain names to IP addresses.
   \item HTTPS uses TLS to protect a connection.
   \item A URL can include a scheme and host.
   \item DNS assigns an email password to a user.
  \end{enumerate}

\end{enumerate}
\lessonheading{36}{Web Browsing, Search, Upload and Download}
\begin{enumerate}[label=\textbf{Q\arabic*.}, leftmargin=*]
\item \textit{[Original]} Expand URL.
  \begin{enumerate}[label=\Alph*.]
    \item Uniform Resource Locator
    \item Universal Record Link
    \item Unified Resource Library
    \item Uniform Retrieval Language
  \end{enumerate}
\item \textit{[Original]} The software used to open and view web pages, such
  as Google Chrome or Mozilla Firefox, is called a:
  \begin{enumerate}[label=\Alph*.]
    \item search engine
    \item web browser
    \item file converter
    \item download manager
  \end{enumerate}
\item \textit{[Original]} Which is the correct order of the three stages of a
  search engine?
  \begin{enumerate}[label=\Alph*.]
    \item Ranking, indexing, crawling
    \item Crawling, indexing, ranking
    \item Indexing, crawling, ranking
    \item Crawling, ranking, indexing
  \end{enumerate}
\item \textit{[Original]} In a search engine, the automatic programs that
  visit pages and follow their links to discover new pages are called:
  \begin{enumerate}[label=\Alph*.]
    \item compilers
    \item crawlers (spiders)
    \item drivers
    \item firewalls
  \end{enumerate}
\item \textit{[Original]} A result on the search page that carries an ``Ad''
  label appears there because:
  \begin{enumerate}[label=\Alph*.]
    \item it ranked highest on relevance
    \item an advertiser paid for the placement
    \item it was visited most recently
    \item it is stored in the browser cache
  \end{enumerate}
\item \textit{[Original]} Which statement correctly distinguishes a sponsored
  result from an organic result?
  \begin{enumerate}[label=\Alph*.]
    \item Sponsored results are ranked by relevance; organic results are paid placements.
    \item Sponsored results are paid placements with an Ad label; organic results are ranked by relevance.
    \item Sponsored results come from the browser cache; organic results come from cookies.
    \item Sponsored results and organic results are produced by the same ranking step.
  \end{enumerate}
\item \textit{[Original]} A small file a website stores on the computer to
  remember login state and preferences is called a:
  \begin{enumerate}[label=\Alph*.]
    \item cookie
    \item cache copy
    \item history entry
    \item bookmark
  \end{enumerate}
\item \textit{[Original]} The browser cache stores:
  \begin{enumerate}[label=\Alph*.]
    \item the list of visited page addresses by date
    \item saved shortcuts to favourite pages
    \item copies of page parts such as images and styles for faster repeat visits
    \item login names and passwords for every website
  \end{enumerate}
\item \textit{[Original]} Which pairing correctly matches each item with what
  it stores? (near-neighbour set)
  \begin{enumerate}[label=\Alph*.]
    \item Cookie --- copies of page images; Cache --- list of visited addresses; History --- login state
    \item Cookie --- visitor data such as login and cart; Cache --- copies of page files; History --- list of visited addresses
    \item Cookie --- list of visited addresses; Cache --- visitor login data; History --- copies of page images
    \item Cookie --- saved page shortcuts; Cache --- login passwords; History --- advertised results
  \end{enumerate}
\item \textit{[Original]} Which of the following is \textbf{NOT} true of the
  browser cache?
  \begin{enumerate}[label=\Alph*.]
    \item It stores copies of page parts such as images and styles.
    \item It is managed by the browser to make repeat visits faster.
    \item It acts as a permanent backup of all website login records.
    \item Clearing it does not sign the user out of websites.
  \end{enumerate}
\item \textit{[Original]} Downloading means:
  \begin{enumerate}[label=\Alph*.]
    \item sending a file from the computer to the internet
    \item copying a file from the internet to the computer
    \item converting a file from one format to another
    \item deleting a file from the server permanently
  \end{enumerate}
\item \textit{[Original]} Uploading means:
  \begin{enumerate}[label=\Alph*.]
    \item copying a file from the internet to the computer
    \item saving a web page for offline reading
    \item sending a file from the computer to the internet
    \item clearing the browsing history
  \end{enumerate}
\item \textit{[Original]} A clerk saves a form PDF from the office portal to
  the desktop, then submits the filled form back on the portal. The two steps
  are respectively:
  \begin{enumerate}[label=\Alph*.]
    \item upload, then download
    \item download, then upload
    \item convert, then bookmark
    \item browse, then index
  \end{enumerate}
\item \textit{[Original]} File conversion means:
  \begin{enumerate}[label=\Alph*.]
    \item moving a file from the computer to a server
    \item changing a file from one format to another while keeping the content
    \item copying a file from the internet to the computer
    \item saving the address of a page for quick return
  \end{enumerate}
\item \textit{[Original]} The extension `.pdf' in a file name denotes the:
  \begin{enumerate}[label=\Alph*.]
    \item transfer direction of the file
    \item search rank of the file
    \item format the file is stored in
    \item website that created the file
  \end{enumerate}
\item \textit{[Original]} Which keyboard shortcut opens a new tab in the
  browser?
  \begin{enumerate}[label=\Alph*.]
    \item Ctrl+T
    \item Ctrl+H
    \item Ctrl+D
    \item Ctrl+P
  \end{enumerate}
\item \textit{[Original]} Which shortcut pair is correctly matched?
  \begin{enumerate}[label=\Alph*.]
    \item Ctrl+D --- open browsing history; Ctrl+H --- bookmark this page
    \item Ctrl+D --- bookmark this page; Ctrl+H --- open browsing history
    \item Ctrl+T --- clear the cache; Ctrl+D --- delete cookies
    \item Ctrl+H --- open a new tab; Ctrl+T --- open browsing history
  \end{enumerate}
\item \textit{[Original]} Evaluate the statements:
  \begin{enumerate}[label=\Roman*.]
    \item Crawling collects pages before the index is built.
    \item Ranking orders the matching pages after the search words are typed.
    \item Sponsored results are ordered purely by relevance with no payment involved.
  \end{enumerate}
  \begin{enumerate}[label=\Alph*.]
    \item I only
    \item I and II only
    \item II and III only
    \item I, II, and III
  \end{enumerate}
\item \textit{[Original]} Evaluate the statements:
  \begin{enumerate}[label=\Roman*.]
    \item Downloading copies a file from the internet to the computer.
    \item Uploading sends a file from the computer to the internet.
    \item Converting a .docx file to PDF automatically uploads it to the server.
  \end{enumerate}
  \begin{enumerate}[label=\Alph*.]
    \item I only
    \item I and II only
    \item II and III only
    \item I, II, and III
  \end{enumerate}
\item \textit{[Original]} Which option correctly describes what clearing each
  store does?
  \begin{enumerate}[label=\Alph*.]
    \item Clearing cookies signs the user out of sites; clearing the cache only slows the next visit; clearing history removes the visited-address record.
    \item Clearing the cache signs the user out of sites; clearing cookies speeds up page loading; clearing history deletes saved page images.
    \item Clearing history signs the user out of sites; clearing cookies deletes page images; clearing the cache removes bookmarks.
    \item Clearing cookies deletes the browsing history; clearing history deletes cookies; clearing the cache deletes bookmarks.
  \end{enumerate}

\item \textit{[JKSSB FAA-pattern]} Downloading means transferring data:
  \begin{enumerate}[label=\Alph*.]
    \item from the Internet to your computer.
    \item from your computer to the printer paper.
    \item from the speaker to the ear.
    \item from the mouse to the pad.
  \end{enumerate}
\item \textit{[JKPSI-pattern]} Which of the following is a search engine?
  \begin{enumerate}[label=\Alph*.]
    \item Google
    \item MS Paint
    \item Calculator
    \item Doorbell
  \end{enumerate}
\item \textit{[SSC CGL-pattern]} Uploading means transferring data:
  \begin{enumerate}[label=\Alph*.]
    \item from your computer to the Internet/server.
    \item from paper into the printer.
    \item from the screen into the keyboard.
    \item from the speaker into the microphone wire.
  \end{enumerate}
\item \textit{[JKSSB-pattern]} \textbf{(MSQ: select ALL correct options.)} Which actions happen in a web browser?
  \begin{enumerate}[label=\Alph*.]
    \item Entering a URL to visit pages
    \item Searching with keywords
    \item Bookmarking favourite pages
    \item Refuelling a vehicle
  \end{enumerate}
\item \textit{[SSC CGL-pattern]} \textbf{(MSQ: select ALL correct options.)} Which statements about browsing are correct?
  \begin{enumerate}[label=\Alph*.]
    \item Downloading brings files to your device.
    \item Uploading sends files from your device.
    \item Bookmarks save page addresses for reuse.
    \item Uploading and downloading are the same direction.
  \end{enumerate}
\item \textit{[Original]} Downloading a file means transferring it:
  \begin{enumerate}[label=\Alph*.]
\item From a remote system to the local device
   \item From the local device to a remote system
   \item From RAM to the CPU register file only
   \item From a monitor to a printer
  \end{enumerate}

\item \textit{[Original]} A website crawler primarily:
  \begin{enumerate}[label=\Alph*.]
\item Discovers and visits web pages for indexing
   \item Encrypts every email attachment
   \item Replaces the browser's operating system
   \item Assigns IP addresses on a LAN
  \end{enumerate}

\item \textit{[Original]} Browser cache is used mainly to:
  \begin{enumerate}[label=\Alph*.]
\item Temporarily store resources for faster reuse
   \item Permanently back up a user's documents
   \item Replace DNS
   \item Prevent all malware infections
  \end{enumerate}

\item \textit{[Original]} Which action sends a local file to a website or remote service?
  \begin{enumerate}[label=\Alph*.]
\item Upload
   \item Download
   \item Refresh
   \item Bookmark
  \end{enumerate}

\item \textit{[Original, MSQ drill]} \textbf{(MSQ: select ALL correct options.)} Which statements are correct?
  \begin{enumerate}[label=\Alph*.]
\item Browser history can record visited pages.
   \item Cookies can store site-related data in a browser.
   \item Upload sends data to a remote service.
   \item Browser cache is a guaranteed disaster-recovery backup.
  \end{enumerate}

\end{enumerate}
\lessonheading{37}{Email Basics}
\begin{enumerate}[label=\textbf{Q\arabic*.}, leftmargin=*]
\item \textit{[Original]} In an email address, the @ symbol separates:
  \begin{enumerate}[label=\Alph*.]
    \item the domain name and the host name
    \item the user name and the domain name
    \item the sender name and the recipient name
    \item the two parts of the domain name
  \end{enumerate}
\item \textit{[Original]} Expand Cc as used in email.
  \begin{enumerate}[label=\Alph*.]
    \item Closed Copy
    \item Carbon Copy
    \item Common Copy
    \item Certified Copy
  \end{enumerate}
\item \textit{[Original]} Expand Bcc as used in email.
  \begin{enumerate}[label=\Alph*.]
    \item Blind Carbon Copy
    \item Bulk Carbon Copy
    \item Basic Client Copy
    \item Backup Courtesy Copy
  \end{enumerate}
\item \textit{[Original]} Which statement about Bcc is correct?
  \begin{enumerate}[label=\Alph*.]
    \item Bcc recipients are hidden from each other, while To and Cc addresses remain visible to all.
    \item Bcc hides the sender's address from every recipient.
    \item Bcc hides the To and Cc lists from everyone who receives the message.
    \item Bcc recipients do not receive the message; they only receive a notification.
  \end{enumerate}
\item \textit{[Original]} A new message arrives in the mailbox. It is stored in the:
  \begin{enumerate}[label=\Alph*.]
    \item Drafts folder
    \item Sent folder
    \item Inbox folder
    \item Trash folder
  \end{enumerate}
\item \textit{[Original]} A message is typed, saved midway, and closed without sending. It waits in the:
  \begin{enumerate}[label=\Alph*.]
    \item Inbox
    \item Drafts
    \item Sent
    \item Spam
  \end{enumerate}
\item \textit{[Original]} Copies of messages already sent are kept in the:
  \begin{enumerate}[label=\Alph*.]
    \item Drafts folder
    \item Spam folder
    \item Trash folder
    \item Sent folder
  \end{enumerate}
\item \textit{[Original]} Unwanted bulk mail filtered out of the Inbox is held in the:
  \begin{enumerate}[label=\Alph*.]
    \item Drafts folder
    \item Spam (junk) folder
    \item Sent folder
    \item Outbox archive
  \end{enumerate}
\item \textit{[Original]} Messages the user has deleted are held, until the folder is emptied, in the:
  \begin{enumerate}[label=\Alph*.]
    \item Trash (deleted items) folder
    \item Drafts folder
    \item Inbox folder
    \item Sent folder
  \end{enumerate}
\item \textit{[Original]} Which of the following is \textbf{NOT} a standard mail folder?
  \begin{enumerate}[label=\Alph*.]
    \item Inbox
    \item Drafts
    \item Sent
    \item Compiler
  \end{enumerate}
\item \textit{[Original]} The short topic line shown in the message list before a message is opened is the:
  \begin{enumerate}[label=\Alph*.]
    \item subject
    \item attachment
    \item domain
    \item protocol
  \end{enumerate}
\item \textit{[Original]} The main written content of a message is the \underline{\hspace{2.2cm}}, and a file carried along with the message is the \underline{\hspace{2.2cm}}.
  \begin{enumerate}[label=\Alph*.]
    \item subject; domain
    \item body; attachment
    \item attachment; body
    \item protocol; server
  \end{enumerate}
\item \textit{[Original]} Which action answers only the sender of a received message?
  \begin{enumerate}[label=\Alph*.]
    \item Reply
    \item Reply All
    \item Forward
    \item Delete
  \end{enumerate}
\item \textit{[Original]} Which action answers the sender plus all visible (To and Cc) recipients?
  \begin{enumerate}[label=\Alph*.]
    \item Forward
    \item Reply All
    \item Save as Draft
    \item Mark as Spam
  \end{enumerate}
\item \textit{[Original]} Which action sends a received message onward to new recipients?
  \begin{enumerate}[label=\Alph*.]
    \item Reply
    \item Reply All
    \item Forward
    \item Archive
  \end{enumerate}
\item \textit{[Original]} Using email through a web browser (for example Gmail or Yahoo Mail) is called \underline{\hspace{2.2cm}}, while an installed mail application (for example Outlook or Thunderbird) is called \underline{\hspace{2.2cm}}.
  \begin{enumerate}[label=\Alph*.]
    \item webmail; an email client
    \item an email client; webmail
    \item spam; trash
    \item a protocol; a domain
  \end{enumerate}
\item \textit{[Original]} Which pair correctly matches the incoming-mail protocol to its behaviour?
  \begin{enumerate}[label=\Alph*.]
    \item POP3 --- keeps mail synced on the server; IMAP --- downloads and usually removes it
    \item POP3 --- downloads mail to the local machine; IMAP --- keeps mail synced on the server
    \item POP3 --- sends outgoing mail; IMAP --- filters spam
    \item POP3 --- hides Bcc addresses; IMAP --- shows Bcc addresses
  \end{enumerate}
\item \textit{[Original]} Expand IMAP.
  \begin{enumerate}[label=\Alph*.]
    \item Internet Message Access Protocol
    \item Internal Mail Attachment Program
    \item Internet Media Archive Process
    \item Integrated Message Alert Protocol
  \end{enumerate}
\item \textit{[Original]} Evaluate the statements:
  \begin{enumerate}[label=\Roman*.]
    \item The @ in an email address separates the user name (left) from the domain name (right).
    \item Bcc recipients are hidden from each other, while To and Cc addresses stay visible to all.
    \item A message saved in Drafts has already been sent to its recipients.
  \end{enumerate}
  \begin{enumerate}[label=\Alph*.]
    \item I and II only
    \item II and III only
    \item I and III only
    \item I, II, and III
  \end{enumerate}
\item \textit{[Original]} Which pairing is correctly matched?
  \begin{enumerate}[label=\Alph*.]
    \item Subject --- a file carried with the message
    \item Drafts --- copies of messages already sent
    \item Reply All --- answers the sender plus all visible To and Cc recipients
    \item POP3 --- keeps mail synced on the server for all devices
  \end{enumerate}

\item \textit{[JKSSB FAA-pattern]} In an email address, the symbol @ separates the:
  \begin{enumerate}[label=\Alph*.]
    \item user name and the domain name.
    \item subject and the message body.
    \item sender and every stamp.
    \item keyboard and the monitor.
  \end{enumerate}
\item \textit{[JKPSI-pattern]} Which folder holds incoming emails?
  \begin{enumerate}[label=\Alph*.]
    \item Inbox
    \item Printer tray
    \item Speaker box
    \item Mouse pad
  \end{enumerate}
\item \textit{[SSC CGL-pattern]} The Cc field in an email is used to:
  \begin{enumerate}[label=\Alph*.]
    \item send a copy to additional recipients.
    \item print extra paper copies.
    \item increase speaker volume.
    \item delete the message.
  \end{enumerate}
\item \textit{[JKSSB-pattern]} \textbf{(MSQ: select ALL correct options.)} Which fields can an email message contain?
  \begin{enumerate}[label=\Alph*.]
    \item To
    \item Cc
    \item Subject
    \item Petrol level
  \end{enumerate}
\item \textit{[SSC CGL-pattern]} \textbf{(MSQ: select ALL correct options.)} Which statements about email are correct?
  \begin{enumerate}[label=\Alph*.]
    \item Every email address contains @.
    \item Attachments can carry files with a message.
    \item The Inbox receives incoming mail.
    \item Email travels only inside one keyboard.
  \end{enumerate}
\item \textit{[Original]} Which email field hides listed recipients from other recipients?
  \begin{enumerate}[label=\Alph*.]
\item Bcc
   \item To
   \item Cc
   \item Subject
  \end{enumerate}

\item \textit{[Original]} A file sent along with an email is called an:
  \begin{enumerate}[label=\Alph*.]
\item Attachment
   \item Header
   \item Signature
   \item Folder
  \end{enumerate}

\item \textit{[Original]} Which email action sends a received message to another recipient?
  \begin{enumerate}[label=\Alph*.]
\item Forward
   \item Reply
   \item Draft
   \item Archive
  \end{enumerate}

\item \textit{[Original]} In an email address, the `@` symbol separates the:
  \begin{enumerate}[label=\Alph*.]
\item User name and domain name
   \item Scheme and file path
   \item Subject and message body
   \item Password and server port
  \end{enumerate}

\item \textit{[Original, MSQ drill]} \textbf{(MSQ: select ALL correct options.)} Which statements about email fields are correct?
  \begin{enumerate}[label=\Alph*.]
\item To lists primary recipients.
   \item Cc lists visible copied recipients.
   \item Bcc conceals listed recipients from other recipients.
   \item Subject contains the attachment bytes by definition.
  \end{enumerate}

\end{enumerate}
\lessonheading{38}{SMTP, POP and IMAP}
\begin{enumerate}[label=\textbf{Q\arabic*.}, leftmargin=*]
\item \textit{[Original]} Expand SMTP.
  \begin{enumerate}[label=\Alph*.]
    \item Simple Mail Transfer Protocol
    \item Secure Mail Transmission Process
    \item Server Message Transfer Program
    \item Simple Message Transport Procedure
  \end{enumerate}
\item \textit{[Original]} Expand POP3.
  \begin{enumerate}[label=\Alph*.]
    \item Private Office Protocol 3
    \item Post Office Protocol version 3
    \item Personal Online Program 3
    \item Protected Output Process 3
  \end{enumerate}
\item \textit{[Original]} Expand IMAP.
  \begin{enumerate}[label=\Alph*.]
    \item Internet Message Access Protocol
    \item Internal Mail Attachment Program
    \item Integrated Messaging Application Process
    \item Internet Managed Address Protocol
  \end{enumerate}
\item \textit{[Original]} Which protocol is used to \textbf{send} mail from
  the sender's program toward the recipient's mail server?
  \begin{enumerate}[label=\Alph*.]
    \item SMTP
    \item POP3
    \item IMAP
    \item SSL
  \end{enumerate}
\item \textit{[Original]} Which protocol \textbf{downloads} received mail to
  one local machine, usually removing it from the server?
  \begin{enumerate}[label=\Alph*.]
    \item SMTP
    \item POP3
    \item IMAP
    \item TLS
  \end{enumerate}
\item \textit{[Original]} Which protocol keeps mail on the server and
  \textbf{syncs} every device's view of folders, read marks, and deletions?
  \begin{enumerate}[label=\Alph*.]
    \item SMTP
    \item POP3
    \item IMAP
    \item SSL
  \end{enumerate}
\item \textit{[Original]} A clerk reads the same office mailbox from a phone
  and a desktop and wants both to show the same inbox. Which protocol fits?
  \begin{enumerate}[label=\Alph*.]
    \item POP3, because it downloads to one machine
    \item IMAP, because it syncs each device with the server copy
    \item SMTP, because it sends mail out
    \item SSL, because it encrypts the channel
  \end{enumerate}
\item \textit{[Original]} Which statement correctly distinguishes POP3 from
  IMAP?
  \begin{enumerate}[label=\Alph*.]
    \item POP3 syncs every device with the server; IMAP downloads to one machine.
    \item POP3 downloads to one machine and usually clears the server; IMAP syncs every device with the server copy.
    \item POP3 and IMAP both send mail out through the sender's server.
    \item POP3 keeps the server copy; IMAP always deletes it at once.
  \end{enumerate}
\item \textit{[Original]} Trace one message: sender $\rightarrow$ ? $\rightarrow$
  recipient server $\rightarrow$ ? $\rightarrow$ reader. Which pair fills the
  gaps?
  \begin{enumerate}[label=\Alph*.]
    \item POP3 then SMTP
    \item IMAP then SMTP
    \item SMTP then POP3 or IMAP
    \item SSL then TLS
  \end{enumerate}
\item \textit{[Original]} Which of the following is \textbf{NOT} true of SMTP?
  \begin{enumerate}[label=\Alph*.]
    \item It sends mail from the sender's program to the mail server.
    \item It forwards mail between mail servers.
    \item It downloads received mail into the recipient's inbox for reading.
    \item It carries the message toward the recipient's mailbox.
  \end{enumerate}
\item \textit{[Original]} BCC stands for blind carbon copy. Who can see the
  BCC recipients?
  \begin{enumerate}[label=\Alph*.]
    \item Everyone on the To and Cc lists
    \item The other BCC recipients
    \item Only the sender; BCC recipients are hidden from each other
    \item Nobody, including the sender
  \end{enumerate}
\item \textit{[Original]} Which statement about To, Cc, and BCC is correct?
  \begin{enumerate}[label=\Alph*.]
    \item To and Cc are visible to all recipients; BCC is hidden from the other recipients.
    \item BCC recipients are shown to everyone on the To list.
    \item Cc hides its recipients from the To list.
    \item The sender does not know who is on BCC.
  \end{enumerate}
\item \textit{[Original]} A mail rule moves all newsletters to one folder
  based on the sender and subject. What is this feature?
  \begin{enumerate}[label=\Alph*.]
    \item A filter (rule) sorting incoming mail
    \item A digital signature on outgoing mail
    \item Encryption of the message body
    \item A blind carbon copy
  \end{enumerate}
\item \textit{[Original]} Which statement about a mail signature is correct?
  \begin{enumerate}[label=\Alph*.]
    \item It is a fixed sender block added to outgoing messages and does not encrypt them.
    \item It scrambles the body so only the recipient can read it.
    \item It sorts incoming mail into folders automatically.
    \item It provides the secure transmission channel.
  \end{enumerate}
\item \textit{[Original]} Which pairing correctly matches the security
  mechanism to its job?
  \begin{enumerate}[label=\Alph*.]
    \item Digital signature --- secrecy; encryption --- sender identity
    \item Digital signature --- authenticity and integrity; encryption --- confidentiality
    \item Digital signature --- secure channel; encryption --- folder sorting
    \item Digital signature --- bulk filtering; encryption --- sender hiding
  \end{enumerate}
\item \textit{[Original]} Evaluate the statements:
  \begin{enumerate}[label=\Roman*.]
    \item SMTP sends mail; POP3 and IMAP receive it.
    \item POP3 downloads to one machine while IMAP syncs with the server.
    \item A digital signature encrypts the body so only the recipient can read it.
  \end{enumerate}
  \begin{enumerate}[label=\Alph*.]
    \item I only
    \item I and II only
    \item II and III only
    \item I, II, and III
  \end{enumerate}
\item \textit{[Original]} Secure transmission of email uses:
  \begin{enumerate}[label=\Alph*.]
    \item SSL/TLS
    \item POP3/IMAP
    \item BCC/Cc
    \item Filters/rules
  \end{enumerate}
\item \textit{[Original]} Dozens of identical unwanted advertising messages
  arrive overnight in bulk. What is this?
  \begin{enumerate}[label=\Alph*.]
    \item Spam
    \item Phishing
    \item Spoofing
    \item Encryption
  \end{enumerate}
\item \textit{[Original]} A message pretends to be the clerk's bank, warns
  that the account is blocked, and asks for the password through a link.
  What is this?
  \begin{enumerate}[label=\Alph*.]
    \item Spam
    \item Phishing
    \item A mail filter
    \item A digital signature
  \end{enumerate}
\item \textit{[Original]} A message shows the boss's name as sender, but
  hovering over the link reveals an unknown address and the attachment has
  an unexpected file. Which combination best describes the threat?
  \begin{enumerate}[label=\Alph*.]
    \item Spoofed sender with a malicious link and attachment
    \item A genuine signed message with a safe attachment
    \item A spam filter sorting newsletters
    \item An encrypted channel over SSL/TLS
  \end{enumerate}

\item \textit{[JKSSB FAA-pattern]} SMTP is used mainly to:
  \begin{enumerate}[label=\Alph*.]
    \item send outgoing emails.
    \item print paper letters.
    \item play music files.
    \item scan photographs.
  \end{enumerate}
\item \textit{[JKPSI-pattern]} Expand POP3.
  \begin{enumerate}[label=\Alph*.]
    \item Post Office Protocol version 3
    \item Personal Office Program 3
    \item Packet Output Process 3
    \item Primary Online Protocol 3
  \end{enumerate}
\item \textit{[SSC CGL-pattern]} IMAP differs from POP3 because IMAP:
  \begin{enumerate}[label=\Alph*.]
    \item keeps mail synced on the server for multiple devices.
    \item prints every message twice.
    \item deletes the Internet connection.
    \item works only without any password.
  \end{enumerate}
\item \textit{[JKSSB-pattern]} \textbf{(MSQ: select ALL correct options.)} Which are email protocols?
  \begin{enumerate}[label=\Alph*.]
    \item SMTP
    \item POP3
    \item IMAP
    \item HTTP-print
  \end{enumerate}
\item \textit{[SSC CGL-pattern]} \textbf{(MSQ: select ALL correct options.)} Which statements about mail protocols are correct?
  \begin{enumerate}[label=\Alph*.]
    \item SMTP sends outgoing mail.
    \item POP3 downloads mail to the local machine.
    \item IMAP syncs mail across devices via the server.
    \item SMTP is used to print paper letters.
  \end{enumerate}
\item \textit{[PYQ-adapted: JKSSB Sub-Inspector 2022, QCRIAS item 399; SRC-17/SRC-22]} Which protocol is commonly used by an email client to retrieve messages from a remote server?
  \begin{enumerate}[label=\Alph*.]
\item SMTP
   \item POP3
   \item SNMP
   \item UDP
  \end{enumerate}

\item \textit{[Original]} Which protocol is primarily used to send email from a client to a mail server?
  \begin{enumerate}[label=\Alph*.]
\item SMTP
   \item IMAP
   \item POP3
   \item DNS
  \end{enumerate}

\item \textit{[Original]} Which protocol typically keeps messages synchronized between a client and a mail server?
  \begin{enumerate}[label=\Alph*.]
\item IMAP
   \item POP3 only
   \item FTP
   \item ARP
  \end{enumerate}

\item \textit{[Original]} TLS applied to SMTP is intended to protect:
  \begin{enumerate}[label=\Alph*.]
\item The transport connection between mail systems
   \item The physical paper envelope
   \item The computer's RAM contents at shutdown
   \item Printer output resolution
  \end{enumerate}

\item \textit{[Original, MSQ drill]} \textbf{(MSQ: select ALL correct options.)} Match the email protocol to its common role.
  \begin{enumerate}[label=\Alph*.]
\item SMTP --- sending/transfer
   \item POP3 --- retrieval/download
   \item IMAP --- server mailbox access/synchronization
   \item DNS --- composing email text
  \end{enumerate}

\end{enumerate}
\lessonheading{39}{E-Banking}
\begin{enumerate}[label=\textbf{Q\arabic*.}, leftmargin=*]
\item \textit{[Original]} Expand UPI.
  \begin{enumerate}[label=\Alph*.]
    \item Unified Payments Interface
    \item Universal Payment Internet
    \item United Payments Infrastructure
    \item Uniform Personal Identification
  \end{enumerate}
\item \textit{[Original]} Expand OTP.
  \begin{enumerate}[label=\Alph*.]
    \item One-Time Password
    \item Online Transfer Protocol
    \item Official Transaction PIN
    \item Open Terminal Payment
  \end{enumerate}
\item \textit{[Original]} Expand CVV.
  \begin{enumerate}[label=\Alph*.]
    \item Card Verification Value
    \item Customer Validation Voucher
    \item Centralised Virtual Vault
    \item Card Validity Version
  \end{enumerate}
\item \textit{[Original]} Which code is entered in the phone to authorise a
  UPI payment?
  \begin{enumerate}[label=\Alph*.]
    \item UPI PIN
    \item ATM PIN
    \item CVV
    \item IFSC
  \end{enumerate}
\item \textit{[Original]} The ATM PIN is used:
  \begin{enumerate}[label=\Alph*.]
    \item with the physical card at an ATM or shop terminal.
    \item to authorise a UPI payment in the phone app.
    \item as the CVV printed on the back of the card.
    \item as a one-time code that expires after one transaction.
  \end{enumerate}
\item \textit{[Original]} Which of the following is \textbf{NOT} true of an
  OTP?
  \begin{enumerate}[label=\Alph*.]
    \item It is sent to the registered mobile number, usually by SMS.
    \item It is valid for one transaction only and expires quickly.
    \item It is printed on the back of the debit card and reused for years.
    \item It must never be shared with a caller claiming to be the bank.
  \end{enumerate}
\item \textit{[Original]} Where is the CVV found, and how long does it last?
  \begin{enumerate}[label=\Alph*.]
    \item Sent by SMS; valid for one transaction only.
    \item Printed on the card; reusable for the life of the card.
    \item Set per account in the UPI app; changed only by the holder.
    \item Displayed on the ATM screen; valid for one withdrawal.
  \end{enumerate}
\item \textit{[Original]} A login asks for the account password and then an
  SMS OTP. This is an example of:
  \begin{enumerate}[label=\Alph*.]
    \item single-factor authentication (password only).
    \item two-factor authentication (know + have).
    \item card-present verification with CVV.
    \item biometric authentication with a fingerprint.
  \end{enumerate}
\item \textit{[Original]} MFA stands for:
  \begin{enumerate}[label=\Alph*.]
    \item Multi-Factor Authentication
    \item Managed Fund Account
    \item Mobile Financial Application
    \item Minimum Fee Agreement
  \end{enumerate}
\item \textit{[Original]} Before approving a QR payment, the payer must
  verify:
  \begin{enumerate}[label=\Alph*.]
    \item the merchant name and amount shown on their own phone screen.
    \item the caller's voice on an incoming customer-care call.
    \item the OTP dictated by the shopkeeper.
    \item the QR sticker colour and size on the counter.
  \end{enumerate}
\item \textit{[Original]} A seller receives a collect request marked
  ``refund --- approve with UPI PIN to receive''. The safe action is to:
  \begin{enumerate}[label=\Alph*.]
    \item approve it at once so the refund arrives faster.
    \item approve it and share the OTP with the sender for confirmation.
    \item decline or ignore it, because entering the UPI PIN sends money.
    \item forward the request link to a friend for a second opinion.
  \end{enumerate}
\item \textit{[Original]} Which statement about receiving money on UPI is
  correct?
  \begin{enumerate}[label=\Alph*.]
    \item The receiver must enter the UPI PIN to accept the amount.
    \item The receiver must share the OTP so the sender can complete it.
    \item The receiver only shares a VPA or QR; no PIN is entered.
    \item The receiver must type the CVV of the sender's card.
  \end{enumerate}
\item \textit{[Original]} Which of the following will a genuine bank
  \textbf{never} ask for on a call or SMS?
  \begin{enumerate}[label=\Alph*.]
    \item The branch address for a visit.
    \item The OTP, CVV, ATM PIN, or UPI PIN.
    \item The account holder's full name for identification.
    \item The preferred language for customer service.
  \end{enumerate}
\item \textit{[Original]} Evaluate the statements:
  \begin{enumerate}[label=\Roman*.]
    \item The ATM PIN authorises UPI payments made in the phone app.
    \item The CVV is printed on the card and is reusable.
    \item An OTP is valid for one transaction only and expires quickly.
  \end{enumerate}
  \begin{enumerate}[label=\Alph*.]
    \item I only
    \item I and II only
    \item II and III only
    \item I, II, and III
  \end{enumerate}
\item \textit{[Original]} Evaluate the statements:
  \begin{enumerate}[label=\Roman*.]
    \item Entering the UPI PIN on a collect request sends the approver's
      money.
    \item A password plus an SMS OTP is two-factor authentication.
    \item The bank never asks for OTP, CVV, or any PIN by call or link.
  \end{enumerate}
  \begin{enumerate}[label=\Alph*.]
    \item I only
    \item I and II only
    \item II and III only
    \item I, II, and III
  \end{enumerate}
\item \textit{[Original]} Which pair correctly matches the code to its
  channel?
  \begin{enumerate}[label=\Alph*.]
    \item ATM PIN --- authorises UPI payments in the phone app.
    \item UPI PIN --- authorises a UPI payment on the phone.
    \item CVV --- arrives by SMS and expires after one use.
    \item OTP --- printed on the card and reused for years.
  \end{enumerate}
\item \textit{[Original]} A caller asks the customer to install a
  screen-sharing app ``to complete KYC''. The correct response is to:
  \begin{enumerate}[label=\Alph*.]
    \item install it at once so the account is not blocked.
    \item install it but hide the phone while typing the OTP.
    \item refuse, because screen sharing exposes OTPs and PINs to fraud.
    \item share the OTP first and then install the app.
  \end{enumerate}
\item \textit{[Original]} Which practice keeps net banking safest on a
  shared or public connection?
  \begin{enumerate}[label=\Alph*.]
    \item Banking over public Wi-Fi and staying logged in for convenience.
    \item Typing the bank address, checking the padlock (https), and logging
      out afterwards.
    \item Opening the bank site only through SMS links.
    \item Saving the ATM PIN in the browser for faster login.
  \end{enumerate}
\item \textit{[Original]} UPI transfers are best described as:
  \begin{enumerate}[label=\Alph*.]
    \item instant transfers, available 24 $\times$ 7, each authorised by the
      UPI PIN.
    \item batched transfers settled only on working days at fixed hours.
    \item card payments that require the CVV printed on the card.
    \item branch-only transfers requiring a paper slip.
  \end{enumerate}
\item \textit{[Original]} QR stands for:
  \begin{enumerate}[label=\Alph*.]
    \item Quick Response
    \item Quantum Relay
    \item Queued Receipt
    \item Quoted Rate
  \end{enumerate}

\item \textit{[JKSSB FAA-pattern]} OTP in e-banking stands for:
  \begin{enumerate}[label=\Alph*.]
    \item One-Time Password
    \item Office Table Paper
    \item Online Ticket Printer
    \item Open Teller Pocket
  \end{enumerate}
\item \textit{[JKPSI-pattern]} NEFT/RTGS are systems for:
  \begin{enumerate}[label=\Alph*.]
    \item electronic fund transfer between banks.
    \item printing paper passbooks only.
    \item polishing coins.
    \item delivering parcels.
  \end{enumerate}
\item \textit{[SSC CGL-pattern]} Which device dispenses cash using a debit card?
  \begin{enumerate}[label=\Alph*.]
    \item ATM
    \item Printer
    \item Scanner
    \item Speaker
  \end{enumerate}
\item \textit{[JKSSB-pattern]} \textbf{(MSQ: select ALL correct options.)} Which services belong to e-banking?
  \begin{enumerate}[label=\Alph*.]
    \item Internet banking
    \item Mobile banking
    \item ATM services
    \item Bullock-cart transport
  \end{enumerate}
\item \textit{[SSC CGL-pattern]} \textbf{(MSQ: select ALL correct options.)} Which statements about safe e-banking are correct?
  \begin{enumerate}[label=\Alph*.]
    \item Never share OTPs with callers.
    \item Use official bank apps and sites.
    \item Check transaction SMS alerts.
    \item Share your PIN freely with strangers.
  \end{enumerate}
\item \textit{[Original]} An OTP is best described as:
  \begin{enumerate}[label=\Alph*.]
\item A one-time verification code
   \item A permanent account password
   \item A bank account number
   \item A card's magnetic strip
  \end{enumerate}

\item \textit{[Original]} Which detail should never be shared with an unsolicited caller asking to authorize a payment?
  \begin{enumerate}[label=\Alph*.]
\item UPI PIN
   \item Public merchant name
   \item Payment date
   \item Transaction category
  \end{enumerate}

\item \textit{[Original]} A collect request in a payment app should be treated as:
  \begin{enumerate}[label=\Alph*.]
\item A request to approve a payment, which must be checked carefully
   \item Proof that money has already arrived
   \item A harmless receipt
   \item A request to disclose an OTP to the sender
  \end{enumerate}

\item \textit{[Original]} Which security practice adds a second verification factor beyond a password?
  \begin{enumerate}[label=\Alph*.]
\item Multi-factor authentication
   \item Reusing the same password
   \item Disabling screen lock
   \item Saving a PIN in a public note
  \end{enumerate}

\item \textit{[Original, MSQ drill]} \textbf{(MSQ: select ALL correct options.)} Which are safe digital-payment practices?
  \begin{enumerate}[label=\Alph*.]
\item Verify the recipient before confirming.
   \item Never disclose a UPI PIN or OTP to another person.
   \item Reject unexpected collect requests.
   \item Enter a PIN to receive money into an account.
  \end{enumerate}

\end{enumerate}
\lessonheading{53}{Networking Legacy Services and Web Standards}
\begin{enumerate}[label=\textbf{Q\arabic*.}, leftmargin=*]
\item \textit{[Original]} Which switching method reserves a dedicated path
  for the whole session before data moves?
  \begin{enumerate}[label=\Alph*.]
    \item Packet switching
    \item Circuit switching
    \item Message flooding
    \item Address switching
  \end{enumerate}
\item \textit{[Original]} Which switching method breaks data into packets
  that travel independently, with no path reserved in advance?
  \begin{enumerate}[label=\Alph*.]
    \item Circuit switching
    \item Leased-line switching
    \item Packet switching
    \item Port switching
  \end{enumerate}
\item \textit{[Original]} Expand Telnet in its networking sense.
  \begin{enumerate}[label=\Alph*.]
    \item Telephone Network Terminal
    \item TELecommunication NETwork
    \item Television Network Transfer
    \item Terminal Network Telephone
  \end{enumerate}
\item \textit{[Original]} Which service provides remote login to another
  computer so the user can run commands on it?
  \begin{enumerate}[label=\Alph*.]
    \item FTP
    \item SMTP
    \item Telnet
    \item DNS
  \end{enumerate}
\item \textit{[Original]} Expand FTP.
  \begin{enumerate}[label=\Alph*.]
    \item Fast Text Processing
    \item File Transfer Protocol
    \item Full Terminal Procedure
    \item Forward Transmission Path
  \end{enumerate}
\item \textit{[Original]} Which service is used to upload and download files
  between computers on a network?
  \begin{enumerate}[label=\Alph*.]
    \item Telnet
    \item FTP
    \item SMTP
    \item IMAP
  \end{enumerate}
\item \textit{[Original]} Match the mail protocol to its role: which protocol
  \textbf{sends} mail from the sender toward the receiver's server?
  \begin{enumerate}[label=\Alph*.]
    \item SMTP
    \item POP3
    \item IMAP
    \item FTP
  \end{enumerate}
\item \textit{[Original]} Which protocol downloads mail to the local machine,
  typically removing it from the server?
  \begin{enumerate}[label=\Alph*.]
    \item SMTP
    \item IMAP
    \item POP3
    \item Telnet
  \end{enumerate}
\item \textit{[Original]} Which protocol keeps mail on the server and syncs
  it across several devices?
  \begin{enumerate}[label=\Alph*.]
    \item POP3
    \item SMTP
    \item FTP
    \item IMAP
  \end{enumerate}
\item \textit{[Original]} In an e-mail address, the symbol @ separates:
  \begin{enumerate}[label=\Alph*.]
    \item domain name and host name
    \item user name and domain name
    \item sender name and receiver name
    \item protocol name and address
  \end{enumerate}
\item \textit{[Original]} Expand W3C.
  \begin{enumerate}[label=\Alph*.]
    \item World Wide Web Consortium
    \item World Web Communication Centre
    \item Wide World Web Committee
    \item Web World Control Council
  \end{enumerate}
\item \textit{[Original]} Which of the following correctly describes DHTML?
  \begin{enumerate}[label=\Alph*.]
    \item A separate replacement language that replaces HTML entirely
    \item HTML combined with styling and scripting for dynamic pages
    \item A protocol for transferring files between servers
    \item A hardware socket for network cables
  \end{enumerate}
\item \textit{[Original]} Which of the following is \textbf{NOT} a correct
  statement about this lesson's services?
  \begin{enumerate}[label=\Alph*.]
    \item Telnet provides remote login to another computer.
    \item FTP transfers files between computers.
    \item SMTP sends mail toward the receiver's server.
    \item POP3 keeps mail synced on the server across devices.
  \end{enumerate}
\item \textit{[Original]} Which of the following statements is
  \textbf{INCORRECT}?
  \begin{enumerate}[label=\Alph*.]
    \item IPv6 addresses are 128 bits long.
    \item A digital signature proves authenticity and integrity of a message.
    \item ICMP is used to transfer application data between hosts.
    \item BGP routes traffic between autonomous systems.
  \end{enumerate}
\item \textit{[Original]} How many layers does the OSI reference model have?
  \begin{enumerate}[label=\Alph*.]
    \item Four
    \item Five
    \item Six
    \item Seven
  \end{enumerate}
\item \textit{[Original]} Which pairing is correctly matched?
  \begin{enumerate}[label=\Alph*.]
    \item Physical port --- a number identifying a network service
    \item Logical port --- a hardware socket for cables
    \item Physical port --- a hardware socket for cables or devices
    \item URL --- the numeric address of a device
  \end{enumerate}
\item \textit{[Original]} IPv4 addresses are 32 bits long, whereas IPv6
  addresses are \underline{\hspace{2.2cm}} long.
  \begin{enumerate}[label=\Alph*.]
    \item 64-bit
    \item 128-bit
    \item 256-bit
    \item 512-bit
  \end{enumerate}
\item \textit{[Original]} Which sequence correctly orders the chain from a
  typed web address to the reached machine?
  \begin{enumerate}[label=\Alph*.]
    \item IP address, DNS, URL
    \item URL typed, DNS resolves the domain name, browser reaches the IP address
    \item DNS typed, URL resolves the IP address, browser reaches the domain name
    \item URL reaches the IP address, which resolves DNS
  \end{enumerate}
\item \textit{[Original]} Which statement correctly distinguishes a physical
  port from a logical port?
  \begin{enumerate}[label=\Alph*.]
    \item A physical port is a service number; a logical port is a hardware socket.
    \item Both are hardware sockets for cables.
    \item A physical port is a hardware socket; a logical port is a number
      identifying a service or process.
    \item Both are numbers identifying a service or process.
  \end{enumerate}
\item \textit{[Original]} Evaluate the statements:
  \begin{enumerate}[label=\Roman*.]
    \item POP3 downloads mail to the local machine.
    \item A digital signature encrypts the message body so only the recipient can read it.
    \item HTTP/3 runs over QUIC over UDP, not over TCP.
  \end{enumerate}
  \begin{enumerate}[label=\Alph*.]
    \item I only
    \item I and III only
    \item II and III only
    \item I, II, and III
  \end{enumerate}

\item \textit{[JKSSB FAA-pattern]} Email existed on early networks mainly through which service model?
  \begin{enumerate}[label=\Alph*.]
    \item Store-and-forward messaging servers
    \item Live printer sharing only
    \item Speaker broadcasting only
    \item Paper postal vans
  \end{enumerate}
\item \textit{[JKPSI-pattern]} Telnet was used to:
  \begin{enumerate}[label=\Alph*.]
    \item operate a remote computer over the network.
    \item print posters locally.
    \item play audio CDs.
    \item charge mobile phones.
  \end{enumerate}
\item \textit{[SSC CGL-pattern]} HTML is used to:
  \begin{enumerate}[label=\Alph*.]
    \item structure web pages.
    \item print paper bills.
    \item repair hard disks.
    \item polish screens.
  \end{enumerate}
\item \textit{[JKSSB-pattern]} \textbf{(MSQ: select ALL correct options.)} Which are legacy network services?
  \begin{enumerate}[label=\Alph*.]
    \item Telnet
    \item FTP
    \item Gopher-era information menus
    \item Touchscreen gestures
  \end{enumerate}
\item \textit{[SSC CGL-pattern]} \textbf{(MSQ: select ALL correct options.)} Which statements about web standards are correct?
  \begin{enumerate}[label=\Alph*.]
    \item HTML structures page content.
    \item A browser renders HTML into visible pages.
    \item URLs locate web resources.
    \item HTML is a printer ink formula.
  \end{enumerate}
\item \textit{[PYQ-adapted: JKSSB JKP Constable 2024, QCRIAS item 765; SRC-02]} Which organization develops standards for the World Wide Web?
  \begin{enumerate}[label=\Alph*.]
\item W3C
   \item WHO
   \item OPEC
   \item UNESCO
  \end{enumerate}

\item \textit{[PYQ-adapted: JKSSB Sub-Inspector 2022, QCRIAS item 399; SRC-17/SRC-22]} Which protocol is commonly used by an email client to retrieve messages from a remote server?
  \begin{enumerate}[label=\Alph*.]
\item SMTP
   \item POP3
   \item SNMP
   \item UDP
  \end{enumerate}

\item \textit{[Original]} Which Internet service transfers files between a client and server?
  \begin{enumerate}[label=\Alph*.]
\item FTP
   \item DNS
   \item ARP
   \item ICMP
  \end{enumerate}

\item \textit{[Original]} Which legacy protocol provides remote terminal access without encryption by default?
  \begin{enumerate}[label=\Alph*.]
\item Telnet
   \item HTTPS
   \item SFTP
   \item IMAPS
  \end{enumerate}

\item \textit{[Original, MSQ drill]} \textbf{(MSQ: select ALL correct options.)} Match protocol/service to its common role.
  \begin{enumerate}[label=\Alph*.]
\item DNS --- name resolution
   \item SMTP --- email sending/transfer
   \item FTP --- file transfer
   \item ICMP --- web page markup
  \end{enumerate}

\end{enumerate}
\lessonheading{63}{Current Web, HTML, JavaScript, VBScript and ASP.NET}
\begin{enumerate}[label=\textbf{Q\arabic*.}, leftmargin=*]
\item \textit{[Original]} Expand HTML.
  \begin{enumerate}[label=\Alph*.]
    \item HyperText Markup Language
    \item HighText Machine Language
    \item Hyperlink and Text Markup Language
    \item Home Tool Markup Language
  \end{enumerate}
\item \textit{[Original]} Expand DHTML.
  \begin{enumerate}[label=\Alph*.]
    \item Dynamic HyperText Markup Language
    \item Digital HyperText Markup Language
    \item Distributed HyperText Markup Language
    \item Direct HyperText Manipulation Language
  \end{enumerate}
\item \textit{[Original]} Expand XML.
  \begin{enumerate}[label=\Alph*.]
    \item eXtensible Markup Language
    \item eXecutable Machine Language
    \item eXtra Markup Linkage
    \item eXternal Media Language
  \end{enumerate}
\item \textit{[Original]} An HTML document is best described as:
  \begin{enumerate}[label=\Alph*.]
    \item a compiled binary file
    \item an ASCII-text document
    \item a database table
    \item an encrypted archive
  \end{enumerate}
\item \textit{[Original]} Which tag produces the largest heading?
  \begin{enumerate}[label=\Alph*.]
    \item \texttt{<h6>}
    \item \texttt{<h3>}
    \item \texttt{<h1>}
    \item \texttt{<head>}
  \end{enumerate}
\item \textit{[Original]} The \texttt{<marquee>} tag is used to:
  \begin{enumerate}[label=\Alph*.]
    \item display the largest heading
    \item scroll text across the page
    \item declare a variable
    \item store one user's data on the server
  \end{enumerate}
\item \textit{[Original]} In \texttt{<font size="4">Hello</font>}, which part
  is the attribute?
  \begin{enumerate}[label=\Alph*.]
    \item \texttt{font}
    \item \texttt{Hello}
    \item \texttt{size="4"}
    \item \texttt{</font>}
  \end{enumerate}
\item \textit{[Original]} Which of the following is \textbf{NOT} a valid HTML
  tag?
  \begin{enumerate}[label=\Alph*.]
    \item \texttt{h1}
    \item \texttt{p}
    \item \texttt{marquee}
    \item \texttt{dimtext}
  \end{enumerate}
\item \textit{[Original]} DHTML is best described as:
  \begin{enumerate}[label=\Alph*.]
    \item a completely separate language that replaces HTML
    \item HTML combined with style sheets and scripts for interactive pages
    \item a database query language for web servers
    \item a compiled language for desktop applications
  \end{enumerate}
\item \textit{[Original]} Which statement correctly distinguishes XML from HTML?
  \begin{enumerate}[label=\Alph*.]
    \item XML uses a fixed tag set for display; HTML uses user-defined tags for data.
    \item HTML structures and displays pages with a fixed tag set; XML carries data with user-defined tags.
    \item XML and HTML are identical in purpose and tag sets.
    \item HTML carries data between servers; XML renders pages in the browser.
  \end{enumerate}
\item \textit{[Original]} Which statement correctly distinguishes Java from
  JavaScript?
  \begin{enumerate}[label=\Alph*.]
    \item JavaScript is a cut-down version of Java for desktops.
    \item Java runs in the browser as a page script; JavaScript compiles to bytecode.
    \item Java is a full compiled programming language; JavaScript is a browser scripting language.
    \item Java and JavaScript are two names for the same language.
  \end{enumerate}
\item \textit{[Original]} JavaScript was developed at:
  \begin{enumerate}[label=\Alph*.]
    \item Microsoft
    \item Netscape
    \item Sun Microsystems hardware division
    \item W3C standards office
  \end{enumerate}
\item \textit{[Original]} In JavaScript, \texttt{"Good" + "Morning"} evaluates to:
  \begin{enumerate}[label=\Alph*.]
    \item \texttt{"GoodMorning"}
    \item \texttt{"Good Morning"} with a space added automatically
    \item a numeric sum
    \item an error, since \texttt{+} cannot join strings
  \end{enumerate}
\item \textit{[Original]} In VBScript, the statement \texttt{Dim city} is used to:
  \begin{enumerate}[label=\Alph*.]
    \item scroll text across the page
    \item declare a variable named \texttt{city}
    \item join two strings
    \item store data for one user on the server
  \end{enumerate}
\item \textit{[Original]} VBScript stands for:
  \begin{enumerate}[label=\Alph*.]
    \item Visual Basic Scripting Edition
    \item Verified Binary Script
    \item Virtual Browser Script
    \item Variable Basic Session
  \end{enumerate}
\item \textit{[Original]} In ASP.NET, the \texttt{Session} object stores:
  \begin{enumerate}[label=\Alph*.]
    \item data shared by all users of the application
    \item data for one user across multiple pages
    \item the fixed tag set of HTML
    \item the combined results of search engines
  \end{enumerate}
\item \textit{[Original]} Dogpile is an example of a:
  \begin{enumerate}[label=\Alph*.]
    \item search engine with its own index
    \item meta-search engine
    \item web browser
    \item markup language
  \end{enumerate}
\item \textit{[Original]} Which pairing is correctly matched?
  \begin{enumerate}[label=\Alph*.]
    \item \texttt{h1} --- smallest heading
    \item \texttt{marquee} --- declares a variable
    \item \texttt{Dim} --- VBScript variable declaration
    \item \texttt{Session} --- shared by all users
  \end{enumerate}
\item \textit{[Original]} Evaluate the statements:
  \begin{enumerate}[label=\Roman*.]
    \item HTML is an ASCII-text document.
    \item JavaScript uses the \texttt{+} operator to concatenate strings.
    \item Dogpile keeps its own page index like Google.
  \end{enumerate}
  \begin{enumerate}[label=\Alph*.]
    \item I only
    \item I and II only
    \item II and III only
    \item I, II, and III
  \end{enumerate}
\item \textit{[Original]} A page must greet each visitor by name on every page
  after login, keeping each visitor's name separate. Which ASP.NET facility
  fits, and why?
  \begin{enumerate}[label=\Alph*.]
    \item \texttt{Session}, because it stores data for one user across pages
    \item Application state, because it stores data for one user across pages
    \item \texttt{marquee}, because it scrolls the name across each page
    \item \texttt{Dim}, because it shares data across all visitors
  \end{enumerate}

\item \textit{[JKSSB FAA-pattern]} HTML tags are enclosed within which brackets?
  \begin{enumerate}[label=\Alph*.]
    \item Angle brackets $< >$
    \item Curly printer brackets
    \item Speaker cones
    \item Paper clips
  \end{enumerate}
\item \textit{[JKPSI-pattern]} JavaScript in web pages mainly adds:
  \begin{enumerate}[label=\Alph*.]
    \item interactive behaviour in the browser.
    \item printer ink refills.
    \item monitor glass coating.
    \item keyboard batteries.
  \end{enumerate}
\item \textit{[SSC CGL-pattern]} ASP.NET is a framework for building:
  \begin{enumerate}[label=\Alph*.]
    \item dynamic web applications.
    \item paper airplanes.
    \item wooden furniture.
    \item road bridges.
  \end{enumerate}
\item \textit{[JKSSB-pattern]} \textbf{(MSQ: select ALL correct options.)} Which technologies run in the web browser?
  \begin{enumerate}[label=\Alph*.]
    \item HTML
    \item CSS
    \item JavaScript
    \item Diesel engine
  \end{enumerate}
\item \textit{[SSC CGL-pattern]} \textbf{(MSQ: select ALL correct options.)} Which statements about web technologies are correct?
  \begin{enumerate}[label=\Alph*.]
    \item HTML structures page content.
    \item CSS styles page appearance.
    \item JavaScript adds interactivity.
    \item HTML prints pages without any browser.
  \end{enumerate}
\item \textit{[PYQ-adapted: JKSSB Sub-Inspector 2022, QCRIAS item 359; SRC-17/SRC-22]} Which markup language is commonly used to add dynamic effects to a web page in older exam terminology?
  \begin{enumerate}[label=\Alph*.]
\item DHTML
   \item XML only
   \item C language
   \item Assembly language
  \end{enumerate}

\item \textit{[Original]} Which language primarily structures the content of a web page?
  \begin{enumerate}[label=\Alph*.]
\item HTML
   \item SMTP
   \item SQL only
   \item BIOS
  \end{enumerate}

\item \textit{[Original]} Which technology is primarily used to add client-side interactivity to a web page?
  \begin{enumerate}[label=\Alph*.]
\item JavaScript
   \item JPEG
   \item FTP
   \item ROM
  \end{enumerate}

\item \textit{[Original]} XML is primarily designed to:
  \begin{enumerate}[label=\Alph*.]
\item Mark up and structure data
   \item Route IP packets
   \item Replace all programming languages
   \item Print magnetic-ink characters
  \end{enumerate}

\item \textit{[Original, MSQ drill]} \textbf{(MSQ: select ALL correct options.)} Which are web-related technologies/languages?
  \begin{enumerate}[label=\Alph*.]
\item HTML
   \item JavaScript
   \item XML
   \item SMTP
  \end{enumerate}

\end{enumerate}
\answersection{8}
\answersubheading{34}{Internet, WWW and Browser}
\begin{enumerate}[label=\textbf{Q\arabic*.}, leftmargin=*]
\item \textbf{B.} The Internet is a global network of connected computers and
  networks. A is a webpage; C is a browser; D is a search engine.
\item \textbf{B.} The WWW is one service on the Internet: linked pages and
  sites. A collapses Internet and Web; C describes an ISP; D describes a
  client.
\item \textbf{B.} A website is the collection, a webpage is one document, and
  the homepage is the main page. A swaps website and webpage; C and D attach
  unrelated roles (connection provider, clickable text, requesting device).
\item \textbf{A.} A browser is software that opens and views pages. B is a
  search engine; C is an ISP; D is a server.
\item \textbf{C.} Mozilla Firefox is a browser. Google (A), Bing (B), and
  Yahoo (D) are search engines, not browsers.
\item \textbf{D.} Bing is a search engine. Chrome (A), Firefox (B), and Edge
  (C) are browsers, not search engines.
\item \textbf{B.} A browser is software that opens pages; a search engine is
  a website that finds pages for keywords. A reverses the two; C assigns the
  ISP's and server's roles; D claims they are identical.
\item \textbf{C.} The client requests a service or resource. A and D describe
  a server; B describes an ISP.
\item \textbf{B.} A web server stores websites and sends pages to clients. A
  reverses the direction; C describes an ISP; D describes a search engine.
\item \textbf{A.} ISP stands for Internet Service Provider (WO 2026 Q74).
  B, C, and D are not expansions of the abbreviation.
\item \textbf{B.} The ISP provides the Internet connection. A is the
  browser's job; C is a hyperlink; D describes HTML authorship.
\item \textbf{C.} A hyperlink is clickable text or an image jumping to
  another page. A is a URL; B is an ISP; D is a browser.
\item \textbf{C.} WWW stands for World Wide Web. A, B, and D rearrange or
  misstate the words.
\item \textbf{A.} URL stands for Uniform Resource Locator. B, C, and D are
  not expansions of the abbreviation.
\item \textbf{B.} HTTP stands for HyperText Transfer Protocol. A, C, and D
  are not the correct full form.
\item \textbf{C.} The Internet and the WWW are not the same thing; the Web is
  one service on the Internet, so this statement is false. A, B, and D are
  each true.
\item \textbf{B.} The homepage is the main or starting page. A describes an
  ISP; C describes the whole website; D describes a hyperlink.
\item \textbf{B.} The browser (client) requests, the server sends the page
  back, and the browser displays it. A, C, and D scramble the roles and the
  order.
\item \textbf{B.} Statements I and II are correct: the Web is one service on
  the Internet, and Chrome is a browser while Google (search) is a search
  engine. III is false because the client requests and the server supplies,
  not the reverse. A omits II; C includes false III and omits I; D includes
  false III.
\item \textbf{C.} An ISP is the company providing the Internet connection, so
  this pairing is correct. A mislabels Google (search) as a browser; B
  mislabels Firefox as a search engine; D mislabels a hyperlink as browser
  software.

\item \textbf{A.} The Internet joins networks globally. B, C, and D are single objects, not the Internet.
\item \textbf{A.} Firefox browses the web. B, C, and D are applications or objects, not browsers.
\item \textbf{A.} WWW is World Wide Web. B, C, and D are not the recognised expansion.
\item \textbf{A, B, C.} Chrome, Firefox, and Edge browse the web. Excel (D) is a spreadsheet.
\item \textbf{A, B, C.} The network, browser, and ISP roles are correct. D is false: a browser is software.
\item \textbf{A.} The Web is an information system of linked resources that uses the Internet. It is not the network infrastructure, one search engine, or an OS.

\item \textbf{A.} A browser requests and displays web resources. A search engine helps locate resources; the other choices serve different roles.

\item \textbf{A.} Search engines index and retrieve web information. Routing is a network function; spreadsheets and firmware storage are unrelated.

\item \textbf{A.} A web page is an individual resource/document. A website is a collection of related pages; the other options are software/services.

\item \textbf{A, B, C.} The Internet supports many services; the Web is one of them. Browsers and search engines also have distinct roles.

\end{enumerate}
\answersubheading{35}{URL, DNS, HTTP and HTTPS}
\begin{enumerate}[label=\textbf{Q\arabic*.}, leftmargin=*]
\item \textbf{A.} URL stands for Uniform Resource Locator. B, C, and D are
  not expansions of the abbreviation.
\item \textbf{A.} DNS stands for Domain Name System. B, C, and D are not the
  correct full forms.
\item \textbf{A.} HTTPS stands for Hypertext Transfer Protocol Secure. B,
  C, and D rearrange or invent words that are not in the expansion.
\item \textbf{A.} The scheme is the protocol word before the colon
  (\texttt{https}). B is the host, C is the port, and D is the path.
\item \textbf{C.} \texttt{mail} is the subdomain prefix. A is the TLD, B
  is the registered domain, and D is a different scheme name.
\item \textbf{B.} \texttt{.in} is a Top-Level Domain. A is a subdomain, C
  is a path, and D is a query.
\item \textbf{B.} HTTP defaults to port 80 and HTTPS to port 443. A reverses
  the pair; C gives mail ports; D claims both use 80.
\item \textbf{D.} Extra parameters after \texttt{?} form the query. The
  scheme (A) names the protocol, the domain (B) names the host, and the
  path (C) locates the page.
\item \textbf{A.} The URL is the whole address of one resource while the
  domain is only its host part. B reverses the relation; C claims identity;
  D misdescribes both terms.
\item \textbf{B.} IPv4 is 32-bit and IPv6 is 128-bit. A reverses them; C
  and D use 64-bit and the 256-bit distractor, which is never correct for
  IPv6.
\item \textbf{A.} DNS translates domain names into IP addresses. It does not
  encrypt pages (B), assign ports (C), or carry mail (D).
\item \textbf{D.} DNS never encrypts page contents; that is the job of TLS
  under HTTPS. Mapping names (A), hierarchy (B), and caching (C) are all
  true of DNS.
\item \textbf{A.} The resolver walks the DNS hierarchy from root to TLD to
  the authoritative server. B, C, and D describe no real resolution step.
\item \textbf{A.} HTTP stands for Hypertext Transfer Protocol. B, C, and D
  are not the correct full forms.
\item \textbf{A.} TLS stands for Transport Layer Security. B, C, and D are
  not expansions of the abbreviation.
\item \textbf{B.} HTTPS is HTTP protected by TLS, while plain HTTP has no
  encryption. A reverses the security; C swaps the ports; D denies the
  HTTP family relation.
\item \textbf{A.} The padlock means the TLS handshake authenticated the
  server and encrypted the channel. B describes a failure; C names the plain
  HTTP port; D is irrelevant to the padlock.
\item \textbf{A.} The browser parses the URL, DNS resolves the host to an
  IP, the browser connects to the server, the TLS handshake secures the
  channel, and then the page is fetched. B, C, and D scramble that order.
\item \textbf{B.} Statements I and II are correct: DNS maps names to IPs,
  and IPv6 is 128-bit. III is false because plain HTTP provides no
  encryption. A omits II; C includes false III and omits I; D includes false
  III.
\item \textbf{C.} The query is the parameter string after \texttt{?}. A
  mislabels the host as the scheme; B assigns port 80 to HTTPS instead of
  HTTP; D repeats the 256-bit error for IPv6.

\item \textbf{A.} URL is Uniform Resource Locator. B, C, and D are not the recognised expansion.
\item \textbf{A.} DNS maps names to addresses. B, C, and D describe impossible conversions.
\item \textbf{A.} HTTPS is secured HTTP. B, C, and D misread the letter S.
\item \textbf{A, B, C.} Protocol, domain, and path form URLs. Ink levels (D) never appear in URLs.
\item \textbf{A, B, C.} The URL, DNS, and HTTPS facts are correct. D is false: HTTP is HyperText Transfer Protocol.
\item \textbf{A.} The scheme, such as https, indicates the access protocol. The other details are not what a URL scheme encodes.

\item \textbf{A.} DNS resolves names to records such as IP addresses. DHCP assigns network configuration; SMTP sends mail; FTP transfers files.

\item \textbf{A.} HTTPS uses HTTP protected by TLS. POP3 receives email, OMR reads marks, and RAID organizes disks.

\item \textbf{A.} `.org` is the top-level domain. `example` is the domain label; https is a scheme, and slash separates path components.

\item \textbf{A, B, C.} DNS resolves names, HTTPS uses TLS, and URLs include components such as scheme and host. DNS does not manage passwords.

\end{enumerate}
\answersubheading{36}{Web Browsing, Search, Upload and Download}
\begin{enumerate}[label=\textbf{Q\arabic*.}, leftmargin=*]
\item \textbf{A.} URL stands for Uniform Resource Locator. B, C, and D are
  not the correct full form.
\item \textbf{B.} Chrome and Firefox are web browsers: software that opens
  and displays pages. A search engine (A) finds pages; a converter (C) or
  download manager (D) does neither.
\item \textbf{B.} The pipeline order is crawling, then indexing, then
  ranking. A starts at the end; C builds the index before collecting pages;
  D ranks before the index exists.
\item \textbf{B.} Crawlers (spiders) visit pages and follow links. Compilers
  (A) translate code; drivers (C) control devices; firewalls (D) filter
  traffic.
\item \textbf{B.} An Ad-labelled result is a paid placement. Relevance rank
  (A) describes organic results; recency (C) describes history; stored copies
  (D) describe the cache.
\item \textbf{B.} Sponsored = paid placement with an Ad label; organic =
  ranked by relevance. A swaps the two; C brings in cache and cookies; D
  claims one shared step.
\item \textbf{A.} A cookie is the website-stored file remembering login and
  preferences. A cache copy (B) is page parts; a history entry (C) is a
  visited address; a bookmark (D) is a saved shortcut.
\item \textbf{C.} The cache keeps copies of page parts for faster repeat
  visits. A describes history; B describes bookmarks; D describes nothing the
  cache stores.
\item \textbf{B.} Cookie = visitor data, cache = page-file copies, history =
  visited addresses. A, C, and D each assign at least one store the wrong
  content.
\item \textbf{C.} The cache is not a backup of login records; it holds page
  copies. A, B, and D are each true statements about the cache.
\item \textbf{B.} Downloading copies from the internet to the computer.
  A describes uploading; C describes conversion; D describes deletion.
\item \textbf{C.} Uploading sends from the computer to the internet.
  A describes downloading; B describes offline saving; D describes clearing
  history.
\item \textbf{B.} Saving the portal form to the desktop is a download;
  submitting the filled form back is an upload. A reverses the order; C and
  D name unrelated steps.
\item \textbf{B.} Conversion changes the format while keeping the content.
  A describes uploading; C describes downloading; D describes bookmarking.
\item \textbf{C.} The extension names the format the file is stored in.
  It says nothing about transfer direction (A), search rank (B), or the
  creating site (D).
\item \textbf{A.} Ctrl+T opens a new tab. Ctrl+H (B) opens history; Ctrl+D
  (C) bookmarks the page; Ctrl+P (D) prints.
\item \textbf{B.} Ctrl+D bookmarks the page and Ctrl+H opens history. A
  swaps the two; C and D assign cache, cookie, tab, and history actions to
  the wrong keys.
\item \textbf{B.} I and II are correct: crawling precedes the index, and
  ranking follows the typed words. III is false because sponsored results
  involve payment, not pure relevance. A omits II; C and D include false III.
\item \textbf{B.} I and II are correct: download = internet to computer,
  upload = computer to internet. III is false because conversion only
  changes the format; the file must still be uploaded separately. A omits
  II; C and D include false III.
\item \textbf{A.} Cookies hold logins (clearing signs out), the cache holds
  page copies (clearing only slows the next visit), and history holds visited
  addresses. B, C, and D each attach at least one effect to the wrong store.

\item \textbf{A.} Downloads flow to the user's machine. B, C, and D describe unrelated local movements.
\item \textbf{A.} Google searches the web. B, C, and D are tools or objects, not search engines.
\item \textbf{A.} Uploads leave the user's machine for a server. B, C, and D reverse or confuse the direction.
\item \textbf{A, B, C.} URLs, search, and bookmarks are browser actions. Refuelling (D) is not.
\item \textbf{A, B, C.} The direction definitions and bookmark use are correct. D is false: the directions are opposite.
\item \textbf{A.} Download means receiving data from a remote service. Sending data from the local device is upload.

\item \textbf{A.} A crawler discovers pages that may be indexed by a search engine. The other options describe unrelated security, OS, or network services.

\item \textbf{A.} Cache can speed repeat access by storing copies temporarily. It is not a permanent backup, DNS replacement, or complete security tool.

\item \textbf{A.} Upload sends data from local to remote. Download receives remote data; refresh reloads a page; bookmark saves a link.

\item \textbf{A, B, C.} History records visits, cookies store site data, and upload sends data outward. Cache is temporary and is not a guaranteed backup.

\end{enumerate}
\answersubheading{37}{Email Basics}
\begin{enumerate}[label=\textbf{Q\arabic*.}, leftmargin=*]
\item \textbf{B.} The @ separates the user name (left) from the domain name
  (right). A wrongly splits the domain side; C confuses the address with the
  sender/recipient roles; D claims the @ sits inside the domain.
\item \textbf{B.} Cc is Carbon Copy. A, C, and D are not the email full form
  of the abbreviation.
\item \textbf{A.} Bcc is Blind Carbon Copy. B, C, and D invent expansions
  that are not used in email.
\item \textbf{A.} Bcc recipients receive the message but are hidden from each
  other, while To and Cc stay visible to all. B is wrong because the sender
  always knows the list; C is wrong because To/Cc remain visible; D is wrong
  because Bcc recipients receive the full message.
\item \textbf{C.} New arrivals wait in the Inbox. Drafts (A) holds unsent
  work; Sent (B) holds sent copies; Trash (D) holds deleted messages.
\item \textbf{B.} A composed message saved without sending waits in Drafts.
  Inbox (A) holds received mail; Sent (C) holds sent copies; Spam (D) holds
  filtered junk.
\item \textbf{D.} Sent keeps copies of messages already sent. Drafts (A)
  holds unsent work; Spam (B) holds junk; Trash (C) holds deleted messages.
\item \textbf{B.} Filtered unwanted bulk mail is held in Spam (junk). Drafts
  (A) holds unsent work; Sent (C) holds sent copies; D is not a mail folder.
\item \textbf{A.} Deleted messages wait in Trash (deleted items) until it is
  emptied. Drafts (B), Inbox (C), and Sent (D) hold unsent, received, and
  sent mail respectively.
\item \textbf{D.} Inbox, Drafts, and Sent are standard mail folders; a
  compiler (D) is a program-translation tool, not a folder.
\item \textbf{A.} The subject is the short topic line shown in the list. An
  attachment (B) is a separate file; the domain (C) is part of the address;
  a protocol (D) is a mail-transfer rule.
\item \textbf{B.} The body is the main written content and the attachment is
  the file carried with the message. A swaps in subject/domain; C swaps the
  two terms; D names protocol/server, which are not message parts.
\item \textbf{A.} Reply answers only the sender. Reply All (B) answers all
  visible recipients; Forward (C) sends the message onward; Delete (D)
  discards it.
\item \textbf{B.} Reply All answers the sender plus all visible To and Cc
  recipients. Forward (A) sends onward to new recipients; C saves unsent
  work; D moves the message to junk.
\item \textbf{C.} Forward sends a received message onward to new recipients.
  Reply (A) answers only the sender; Reply All (B) answers all visible
  recipients; Archive (D) files the message away.
\item \textbf{A.} Mail in a browser is webmail and an installed mail
  application is an email client. B swaps the two; C names two folders; D
  names protocol/domain, which are not usage modes.
\item \textbf{B.} POP3 downloads mail to the local machine (usually removing
  it from the server) while IMAP keeps mail synced on the server. A swaps
  the two; C mislabels both as sending/filtering; D confuses protocols with
  Bcc visibility.
\item \textbf{A.} IMAP is Internet Message Access Protocol. B, C, and D are
  not the correct full form.
\item \textbf{A.} Statements I and II are correct: @ divides user name from
  domain name, and Bcc hides recipients from each other while To/Cc stay
  visible. III is false because Drafts holds messages not yet sent. B and C
  each include false III; D includes false III.
\item \textbf{C.} Reply All answers the sender plus all visible To and Cc
  recipients, so this pairing is correct. A mislabels the subject as a file
  (that is an attachment); B mislabels Drafts (sent copies belong in Sent);
  D mislabels POP3 (syncing on the server is IMAP; POP3 downloads).

\item \textbf{A.} The @ divides user and domain parts. B, C, and D describe separations @ does not perform.
\item \textbf{A.} The Inbox collects incoming mail. B, C, and D are hardware, not mail folders.
\item \textbf{A.} Cc sends carbon copies to others. B, C, and D describe unrelated actions.
\item \textbf{A, B, C.} To, Cc, and Subject are email fields. A petrol level (D) is not.
\item \textbf{A, B, C.} The @ rule, attachments, and Inbox role are correct. D is false: email travels across networks.
\item \textbf{A.} Bcc hides those recipients from other recipients. To and Cc recipients are visible; Subject is the message topic.

\item \textbf{A.} An attachment is a file included with a message. Headers carry routing/message metadata; a signature identifies the sender; a folder organizes mail.

\item \textbf{A.} Forward sends a received message onward. Reply responds to its sender; Draft stores an unfinished message; Archive files a message.

\item \textbf{A.} The local part identifies the mailbox/user; the domain identifies the mail domain. The other pairs are unrelated.

\item \textbf{A, B, C.} To, Cc, and Bcc serve recipient roles. Subject is a text topic field; attachments are separate message parts.

\end{enumerate}
\answersubheading{38}{SMTP, POP and IMAP}
\begin{enumerate}[label=\textbf{Q\arabic*.}, leftmargin=*]
\item \textbf{A.} SMTP stands for Simple Mail Transfer Protocol. B, C, and D
  are not the correct expansion.
\item \textbf{B.} POP3 stands for Post Office Protocol version 3. A, C, and
  D are not the correct full forms.
\item \textbf{A.} IMAP stands for Internet Message Access Protocol. B, C,
  and D are not the correct expansions.
\item \textbf{A.} SMTP sends mail from the sender's program toward the
  recipient's server. POP3 (B) and IMAP (C) receive mail; SSL (D) is a secure
  channel, not a sending protocol.
\item \textbf{B.} POP3 downloads received mail to one local machine, usually
  removing it from the server. SMTP (A) sends; IMAP (C) syncs with the
  server; TLS (D) is a secure channel.
\item \textbf{C.} IMAP keeps mail on the server and syncs folders, read
  marks, and deletions across devices. SMTP (A) sends; POP3 (B) downloads to
  one machine; SSL (D) protects transmission.
\item \textbf{B.} IMAP suits several devices because every device mirrors
  the server copy. A describes POP3, which strands mail on one machine; C
  sends rather than receives; D protects the channel rather than syncing.
\item \textbf{B.} POP3 downloads to one machine (usually clearing the
  server) while IMAP syncs every device with the server copy. A swaps the
  two; C wrongly makes both sending protocols; D swaps the server-copy
  behaviour.
\item \textbf{C.} The sender's side and server forwarding use SMTP, and the
  recipient collects with POP3 or IMAP. A and B put the receiving protocols
  on the sending side; D names the secure channel, not the mail flow.
\item \textbf{C.} SMTP does not download mail for reading; that is the job
  of POP3 or IMAP. Sending to the server (A), forwarding between servers
  (B), and carrying the message toward the mailbox (D) are all true of SMTP.
\item \textbf{C.} Only the sender knows the BCC list; BCC recipients are
  hidden from each other. A and B wrongly expose the list; D wrongly hides
  it even from the sender.
\item \textbf{A.} To and Cc are visible to all recipients while BCC stays
  hidden from the other recipients. B exposes BCC; C wrongly hides Cc; D
  wrongly claims the sender is unaware of BCC.
\item \textbf{A.} Sorting incoming mail by sender or subject is a filter
  (rule). A signature (B) labels outgoing mail; encryption (C) hides content;
  BCC (D) addresses hidden recipients.
\item \textbf{A.} A signature is a fixed sender block on outgoing messages
  and provides no encryption. B describes encryption; C describes a filter;
  D describes SSL/TLS.
\item \textbf{B.} A digital signature gives authenticity and integrity while
  encryption gives confidentiality. A swaps the two jobs; C and D attach
  unrelated sorting, channel, filtering, and hiding jobs.
\item \textbf{B.} Statements I and II are correct: SMTP sends while POP3 and
  IMAP receive, and POP3 downloads where IMAP syncs. III is false because a
  signature proves the sender but does not encrypt the body. A omits II; C
  includes false III and omits I; D includes false III.
\item \textbf{A.} SSL/TLS provide the secure email transmission channel.
  POP3/IMAP (B) receive mail; BCC/Cc (C) address it; filters/rules (D) sort
  it.
\item \textbf{A.} Bulk unwanted advertising mail is spam. Phishing (B) is a
  targeted fake-trust fraud; spoofing (C) is a forged sender; encryption (D)
  is a protection, not a threat.
\item \textbf{B.} Impersonating the bank with urgency plus a link to steal
  the password is phishing. Spam (A) is bulk junk; a filter (C) sorts mail;
  a signature (D) labels outgoing mail.
\item \textbf{A.} A forged familiar sender plus a mismatched link and a
  suspicious attachment is a spoofed message carrying a malicious link and
  attachment. B wrongly calls it genuine and safe; C describes sorting; D
  describes the secure channel.

\item \textbf{A.} SMTP transfers outgoing mail. B, C, and D are unrelated device tasks.
\item \textbf{A.} POP3 is Post Office Protocol version 3. B, C, and D are not the recognised expansion.
\item \textbf{A.} IMAP's server sync suits many devices. B, C, and D misstate the protocol.
\item \textbf{A, B, C.} SMTP, POP3, and IMAP handle mail. D is not a mail protocol.
\item \textbf{A, B, C.} The send, download, and sync roles are correct. D is false: SMTP sends mail; it does not print.
\item \textbf{B.} POP3 retrieves mail from a server. SMTP sends/transfers email; SNMP manages network devices; UDP is a transport protocol.

\item \textbf{A.} SMTP handles mail submission/transfer. IMAP and POP3 retrieve/access mail; DNS resolves names.

\item \textbf{A.} IMAP is designed for server-based mailbox access and synchronization. POP3 commonly downloads messages; FTP transfers files; ARP resolves local addresses.

\item \textbf{A.} TLS can protect a connection hop. It does not by itself provide end-to-end message encryption across every mail server.

\item \textbf{A, B, C.} SMTP, POP3, and IMAP have the listed mail roles. DNS resolves domain names, not message composition.

\end{enumerate}
\answersubheading{39}{E-Banking}
\begin{enumerate}[label=\textbf{Q\arabic*.}, leftmargin=*]
\item \textbf{A.} UPI stands for Unified Payments Interface. B, C, and D
  are not the correct full form.
\item \textbf{A.} OTP stands for One-Time Password. B, C, and D are not
  expansions of the abbreviation.
\item \textbf{A.} CVV stands for Card Verification Value. B, C, and D are
  not the correct full form.
\item \textbf{A.} The UPI PIN authorises a UPI payment in the phone. The
  ATM PIN (B) works with the physical card, the CVV (C) is the printed card
  code, and IFSC (D) identifies a branch, not a payer.
\item \textbf{A.} The ATM PIN is typed with the physical card at an ATM or
  shop terminal. B describes the UPI PIN; C describes the CVV; D describes
  the OTP.
\item \textbf{C.} No OTP is printed on the card; that describes the CVV.
  Being sent to the registered mobile (A), single-use and quick expiry (B),
  and never shared with callers (D) are all true of an OTP.
\item \textbf{B.} The CVV is printed on the card and reusable for the life
  of the card. A describes the OTP; C describes the UPI PIN; D is not a
  real code.
\item \textbf{B.} Password (something you know) plus SMS OTP (something you
  have) is two-factor authentication. A ignores the OTP factor; C and D name
  factors not used here (CVV, fingerprint).
\item \textbf{A.} MFA stands for Multi-Factor Authentication. B, C, and D
  are not expansions of the abbreviation.
\item \textbf{A.} The payer verifies the merchant name and amount on their
  own screen before entering the UPI PIN. A caller's voice (B) proves
  nothing; the shopkeeper never needs the OTP (C); sticker appearance (D)
  can be faked.
\item \textbf{C.} Approving a collect request with the UPI PIN sends money,
  so an unexpected ``refund'' request must be declined or ignored. A would
  pay the fraudster; B adds OTP-sharing to the loss; D spreads the fraud
  link.
\item \textbf{C.} Receiving on UPI needs no PIN: the receiver shares a VPA
  or QR. A reverses the rule (PIN = paying); B hands the OTP to the
  fraudster; D misuses the CVV, which belongs to card payments.
\item \textbf{B.} A genuine bank never asks for OTP, CVV, ATM PIN, or UPI
  PIN. A branch address (A), the holder's name (C), or a language choice
  (D) are ordinary service questions, not secrets.
\item \textbf{C.} II and III are correct: the CVV is printed and reusable,
  and the OTP is single-use. I is false because the ATM PIN does not
  authorise UPI payments; that is the UPI PIN's job. A and B include false
  I; D includes false I.
\item \textbf{D.} All three are correct: a collect approval pays (I),
  password plus OTP is two factors (II), and the bank never asks for secret
  codes (III). A, B, and C each omit a true statement.
\item \textbf{B.} The UPI PIN authorises a UPI payment on the phone. A
  swaps the ATM PIN onto the UPI channel; C gives the OTP's behaviour to
  the CVV; D gives the CVV's printing to the OTP.
\item \textbf{C.} A screen-sharing app exposes OTPs and PINs typed on the
  screen, so it must be refused. A and B both install the fraudster's tool;
  D shares the OTP as well.
\item \textbf{B.} Typing the address, checking the padlock (https), and
  logging out is the safe habit. A keeps the session exposed on untrusted
  Wi-Fi; C trusts SMS links, a phishing channel; D stores a secret in the
  browser.
\item \textbf{A.} UPI is instant, 24 $\times$ 7, and each payment is
  authorised by the UPI PIN. B describes batched NEFT-style settlement; C
  describes card payments; D contradicts phone-based e-banking.
\item \textbf{A.} QR stands for Quick Response. B, C, and D are not the
  correct expansion.

\item \textbf{A.} OTP is a One-Time Password. B, C, and D are not the recognised expansion.
\item \textbf{A.} NEFT/RTGS move funds electronically. B, C, and D are non-transfer chores.
\item \textbf{A.} ATMs dispense cash. B, C, and D are unrelated peripherals.
\item \textbf{A, B, C.} Internet, mobile, and ATM channels are e-banking. D is not a banking channel.
\item \textbf{A, B, C.} OTP secrecy, official channels, and alert checks are safe practice. D is unsafe and false.
\item \textbf{A.} An OTP is a temporary one-time code for authentication. It is not a permanent password or payment identifier.

\item \textbf{A.} A UPI PIN authorizes transactions and must remain secret. Public transaction context does not substitute for an authentication secret.

\item \textbf{A.} Approving a collect request may authorize a debit. Check the payee and purpose; it is not proof of receipt and never requires sharing an OTP.

\item \textbf{A.} MFA combines distinct factors. The other choices weaken security.

\item \textbf{A, B, C.} Verify recipient details and protect secrets. Receiving money does not require disclosing or entering a UPI PIN to a requester.

\end{enumerate}
\answersubheading{53}{Networking Legacy Services and Web Standards}
\begin{enumerate}[label=\textbf{Q\arabic*.}, leftmargin=*]
\item \textbf{B.} Circuit switching reserves a dedicated path for the whole
  session before data moves. Packet switching (A) sends independent packets
  with no reserved path; C and D are not switching methods in this lesson.
\item \textbf{C.} Packet switching breaks data into independently travelling
  packets with no reserved path. Circuit switching (A) reserves a path;
  B and D are not recognised switching methods.
\item \textbf{B.} Telnet is TELecommunication NETwork. A, C, and D are not
  the correct expansion.
\item \textbf{C.} Telnet provides remote login for running commands on
  another computer. FTP (A) transfers files; SMTP (B) sends mail; DNS (D)
  resolves names.
\item \textbf{B.} FTP is File Transfer Protocol. A, C, and D are not the
  correct full form.
\item \textbf{B.} FTP uploads and downloads files between computers. Telnet
  (A) gives remote login; SMTP (C) sends mail; IMAP (D) syncs mail kept on
  the server.
\item \textbf{A.} SMTP sends mail from the sender toward the receiver's
  server. POP3 (B) downloads mail locally; IMAP (C) syncs server-side mail;
  FTP (D) moves files, not mail.
\item \textbf{C.} POP3 downloads mail to the local machine, typically
  removing it from the server. SMTP (A) sends mail; IMAP (B) syncs mail on
  the server; Telnet (D) is remote login.
\item \textbf{D.} IMAP keeps mail on the server and syncs it across devices.
  POP3 (A) downloads locally instead; SMTP (B) sends mail; FTP (C) transfers
  files.
\item \textbf{B.} The @ sign separates the user name from the domain name.
  A swaps in host wording; C invents sender/receiver halves; D invents a
  protocol half.
\item \textbf{A.} W3C is the World Wide Web Consortium. B, C, and D are not
  the correct expansion.
\item \textbf{B.} DHTML is HTML combined with styling and scripting for
  dynamic pages. A wrongly calls it a replacement language; C describes FTP;
  D describes a physical port.
\item \textbf{D.} POP3 downloads mail locally; keeping mail synced on the
  server is IMAP, so D is the NOT-correct statement. A, B, and C each state
  their service's true role.
\item \textbf{C.} ICMP is a control/diagnostic protocol and carries no
  application data, so C is the incorrect statement. A states the true
  128-bit length; B states the true signature role; D states the true
  inter-domain BGP role.
\item \textbf{D.} The OSI model has seven layers. Four (A) is the TCP/IP
  model count planted as a near-neighbour; B and C are wrong counts.
\item \textbf{C.} A physical port is a hardware socket for cables or devices.
  A and B swap physical with logical; D confuses the URL with the IP
  address.
\item \textbf{B.} IPv6 addresses are 128-bit long. The 256-bit option (C) is
  the planted real-paper distractor; A and D are wrong lengths.
\item \textbf{B.} The chain is: type the URL, DNS resolves the domain name,
  the browser reaches the IP address. A reverses the order; C and D scramble
  which component resolves which.
\item \textbf{C.} A physical port is a hardware socket while a logical port
  is a service/process number. A reverses the pair; B and D collapse the
  distinction to one side.
\item \textbf{B.} Statements I and III are correct: POP3 downloads locally,
  and HTTP/3 runs over QUIC over UDP. II is false because a digital
  signature proves authenticity/integrity while encryption gives
  confidentiality. A omits III; C and D include false II.

\item \textbf{A.} Early mail used store-and-forward servers. B, C, and D are not messaging models.
\item \textbf{A.} Telnet gave remote terminal access. B, C, and D are local physical tasks.
\item \textbf{A.} HTML marks up web pages. B, C, and D are unrelated chores.
\item \textbf{A, B, C.} Telnet, FTP, and Gopher menus are legacy services. Gestures (D) are modern input acts.
\item \textbf{A, B, C.} The HTML, browser, and URL facts are correct. D is false: HTML is a markup language.
\item \textbf{A.} W3C is the World Wide Web Consortium, which develops Web standards. The other organizations have different mandates.

\item \textbf{B.} POP3 retrieves email. SMTP sends/transfers it; SNMP manages network devices; UDP is a transport protocol.

\item \textbf{A.} FTP is a file-transfer protocol. DNS resolves names, ARP resolves local addresses, and ICMP carries control/diagnostic messages.

\item \textbf{A.} Telnet is a legacy plaintext remote-terminal protocol. HTTPS, SFTP, and IMAPS are secure protocol variants/services.

\item \textbf{A, B, C.} DNS, SMTP, and FTP match their roles. ICMP is for network control/diagnostics, not HTML markup.

\end{enumerate}
\answersubheading{63}{Current Web, HTML, JavaScript, VBScript and ASP.NET}
\begin{enumerate}[label=\textbf{Q\arabic*.}, leftmargin=*]
\item \textbf{A.} HTML stands for HyperText Markup Language. B, C, and D are
  not the correct expansion.
\item \textbf{A.} DHTML stands for Dynamic HyperText Markup Language. B, C,
  and D substitute the wrong first word.
\item \textbf{A.} XML stands for eXtensible Markup Language. B, C, and D are
  not the correct full form.
\item \textbf{B.} An HTML document is plain ASCII text, writable in any text
  editor. A claims it is compiled binary; C and D name unrelated formats.
\item \textbf{C.} \texttt{h1} is the largest heading (\texttt{h6} smallest).
  A and B name smaller headings; D names the header section element, not a
  heading level.
\item \textbf{B.} \texttt{marquee} scrolls text across the page. A describes
  \texttt{h1}; C describes \texttt{Dim}; D describes \texttt{Session}.
\item \textbf{C.} The attribute sits inside the opening tag as
  \texttt{name="value"}, so \texttt{size="4"} is the attribute. A is the tag;
  B is the content; D is the closing tag.
\item \textbf{D.} \texttt{dimtext} is invented and not a defined HTML tag.
  \texttt{h1} (A), \texttt{p} (B), and \texttt{marquee} (C) are all valid tags.
\item \textbf{B.} DHTML is HTML plus style sheets plus scripts, making pages
  interactive after loading. A wrongly calls it a replacement language; C and
  D name unrelated language roles.
\item \textbf{B.} HTML displays pages with a fixed tag set while XML carries
  data with user-defined tags. A reverses the two; C claims they are identical;
  D reverses their roles.
\item \textbf{C.} Java is a full compiled language and JavaScript is a browser
  scripting language. A wrongly calls JavaScript a cut-down Java; B swaps where
  each runs; D claims they are the same language.
\item \textbf{B.} JavaScript was developed at Netscape for the Navigator
  browser. A is the home of VBScript/ASP.NET; C and D are distractors from
  nearby names.
\item \textbf{A.} The \texttt{+} operator concatenates strings, giving
  \texttt{"GoodMorning"} with no added space. B invents an automatic space; C
  treats text as a numeric sum; D denies that \texttt{+} joins strings.
\item \textbf{B.} \texttt{Dim} declares a variable in VBScript. A describes
  \texttt{marquee}; C describes JavaScript \texttt{+}; D describes
  \texttt{Session}.
\item \textbf{A.} VBScript stands for Visual Basic Scripting Edition. B, C,
  and D are invented expansions.
\item \textbf{B.} \texttt{Session} stores data for one user across multiple
  pages. A describes Application state; C and D name unrelated concepts.
\item \textbf{B.} Dogpile is a meta-search engine: it combines other engines'
  results and keeps no index of its own. A describes Google/Bing; C and D name
  unrelated categories.
\item \textbf{C.} \texttt{Dim} is indeed the VBScript variable declaration.
  A reverses the heading ladder (\texttt{h1} is largest); B assigns the wrong
  role to \texttt{marquee}; D wrongly says \texttt{Session} is shared.
\item \textbf{B.} Statements I and II are correct: HTML is an ASCII-text
  document, and JavaScript \texttt{+} concatenates strings. III is false
  because Dogpile keeps no index of its own. A omits II; C and D include false
  III.
\item \textbf{A.} \texttt{Session} is the per-user store across pages, so each
  visitor's name stays separate. B names the shared store, which would mix
  visitors; C scrolls text but stores nothing; D is a declaration statement,
  not a per-user store.

\item \textbf{A.} Tags sit inside $< >$. B, C, and D are unrelated objects.
\item \textbf{A.} JavaScript drives browser interactivity. B, C, and D are physical non-web items.
\item \textbf{A.} ASP.NET builds dynamic web apps. B, C, and D are non-software constructions.
\item \textbf{A, B, C.} HTML, CSS, and JavaScript run in browsers. A diesel engine (D) does not.
\item \textbf{A, B, C.} The structure, style, and interactivity roles are correct. D is false: a browser renders HTML.
\item \textbf{A.} DHTML describes combining HTML, scripting, and style changes for dynamic page behavior. The other choices do not name that presentation concept.

\item \textbf{A.} HTML marks up web content. SMTP transfers email, SQL queries databases, and BIOS is firmware.

\item \textbf{A.} JavaScript is commonly used for client-side behavior. JPEG is an image format, FTP transfers files, and ROM is memory.

\item \textbf{A.} XML structures and describes data. Packet routing, replacing all languages, and MICR printing are unrelated.

\item \textbf{A, B, C.} HTML, JavaScript, and XML are web technologies/languages. SMTP is an email-transfer protocol.

\end{enumerate}
\unitopener{9}{Networking}
\lessonheading{40}{Networking Basics}
\begin{enumerate}[label=\textbf{Q\arabic*.}, leftmargin=*]
\item \textit{[Original]} Expand NIC.
  \begin{enumerate}[label=\Alph*.]
    \item Network Internet Cable
    \item National Information Centre
    \item Network Interface Card
    \item Node Interconnection Circuit
  \end{enumerate}
\item \textit{[Original]} Expand LAN.
  \begin{enumerate}[label=\Alph*.]
    \item Large Area Network
    \item Local Area Network
    \item Linked Access Network
    \item Long Antenna Network
  \end{enumerate}
\item \textit{[Original]} In a computer network, which statement correctly
  distinguishes a client from a server?
  \begin{enumerate}[label=\Alph*.]
    \item The client provides the service; the server requests it.
    \item The client requests a service; the server provides it.
    \item Client and server both only forward signals without storing data.
    \item The server connects the computer to the cable; the client blocks
      unauthorised access.
  \end{enumerate}
\item \textit{[Original]} Which network type covers a single building or
  campus, such as an office floor?
  \begin{enumerate}[label=\Alph*.]
    \item PAN
    \item LAN
    \item MAN
    \item WAN
  \end{enumerate}
\item \textit{[Original]} A bank links its branches spread across one city.
  This network is an example of a:
  \begin{enumerate}[label=\Alph*.]
    \item PAN
    \item LAN
    \item MAN
    \item WLAN built only with infrared remotes
  \end{enumerate}
\item \textit{[Original]} Which of the following is \textbf{NOT} true of an
  intranet?
  \begin{enumerate}[label=\Alph*.]
    \item It is a private network inside one organisation.
    \item It is usable only by the organisation's own staff.
    \item It is the worldwide public network open to anyone with access.
    \item It may carry the organisation's internal portal.
  \end{enumerate}
\item \textit{[Original]} Which topology connects every node to one central
  device, so that all data passes through the centre?
  \begin{enumerate}[label=\Alph*.]
    \item Bus
    \item Star
    \item Ring
    \item Mesh
  \end{enumerate}
\item \textit{[Original]} In which topology does a break in the single main
  cable bring down the whole network?
  \begin{enumerate}[label=\Alph*.]
    \item Bus
    \item Star
    \item Mesh
    \item Tree
  \end{enumerate}
\item \textit{[Original]} In a ring topology, each node:
  \begin{enumerate}[label=\Alph*.]
    \item connects to one central hub that stores all the data.
    \item connects to exactly two neighbours, forming a closed circle.
    \item has a direct link to every other node in the network.
    \item joins the network only through a satellite link.
  \end{enumerate}
\item \textit{[Original]} Which topology is the most reliable but also the
  costliest, because nodes have direct links to many other nodes?
  \begin{enumerate}[label=\Alph*.]
    \item Bus
    \item Star
    \item Mesh
    \item Ring
  \end{enumerate}
\item \textit{[Original]} A large campus joins each department's star network
  through one central backbone. This arrangement is a:
  \begin{enumerate}[label=\Alph*.]
    \item bus topology
    \item ring topology
    \item mesh topology
    \item tree topology
  \end{enumerate}
\item \textit{[Original]} Which of the following lists only guided
  (wired) transmission media?
  \begin{enumerate}[label=\Alph*.]
    \item Twisted-pair cable, coaxial cable, fibre-optic cable
    \item Radio waves, microwaves, satellite links
    \item Infrared, radio waves, microwaves
    \item Bluetooth, infrared, satellite links
  \end{enumerate}
\item \textit{[Original]} Which guided medium carries light signals and is the
  fastest of the three cables?
  \begin{enumerate}[label=\Alph*.]
    \item Twisted-pair cable
    \item Coaxial cable
    \item Fibre-optic cable
    \item Telephone cord with an infrared LED
  \end{enumerate}
\item \textit{[Original]} Which device broadcasts every incoming signal to
  all of its ports?
  \begin{enumerate}[label=\Alph*.]
    \item Hub
    \item Switch
    \item Router
    \item Firewall
  \end{enumerate}
\item \textit{[Original]} Which device forwards each message only to the
  addressed port, and is therefore smarter and faster than a hub?
  \begin{enumerate}[label=\Alph*.]
    \item Hub
    \item Repeater
    \item Switch
    \item Modem
  \end{enumerate}
\item \textit{[Original]} Which device regenerates a weak signal so it can
  travel further, without choosing any path?
  \begin{enumerate}[label=\Alph*.]
    \item Router
    \item Repeater
    \item Gateway
    \item Firewall
  \end{enumerate}
\item \textit{[Original]} Which device connects different networks and
  chooses the best path for each packet?
  \begin{enumerate}[label=\Alph*.]
    \item Hub
    \item Switch
    \item Bridge
    \item Router
  \end{enumerate}
\item \textit{[Original]} Evaluate the statements:
  \begin{enumerate}[label=\Roman*.]
    \item A bridge joins two LAN segments and filters traffic between them.
    \item A gateway translates between networks using different protocols.
    \item A firewall blocks unauthorised access while letting legitimate
      traffic pass.
  \end{enumerate}
  \begin{enumerate}[label=\Alph*.]
    \item I only
    \item I and II only
    \item II and III only
    \item I, II, and III
  \end{enumerate}
\item \textit{[Original]} Which pairing is correctly matched?
  \begin{enumerate}[label=\Alph*.]
    \item Modem --- blocks unauthorised access to the network
    \item Access point --- lets Wi-Fi devices join the wired network
    \item Router --- regenerates a weak signal along a long cable
    \item NIC --- translates between networks using different protocols
  \end{enumerate}
\item \textit{[Original]} An office connects to its ISP over a telephone line
  that cannot carry digital signals directly. Which device performs the
  required digital--analogue conversion?
  \begin{enumerate}[label=\Alph*.]
    \item Switch
    \item Modem
    \item Bridge
    \item Access point
  \end{enumerate}

\item \textit{[JKSSB FAA-pattern]} LAN stands for:
  \begin{enumerate}[label=\Alph*.]
    \item Local Area Network
    \item Large Audio Node
    \item Long Antenna Net
    \item Light Access Number
  \end{enumerate}
\item \textit{[JKPSI-pattern]} Which device connects different networks together?
  \begin{enumerate}[label=\Alph*.]
    \item Router
    \item Printer
    \item Speaker
    \item Scanner
  \end{enumerate}
\item \textit{[SSC CGL-pattern]} Which network spans the largest area?
  \begin{enumerate}[label=\Alph*.]
    \item WAN
    \item LAN
    \item PAN
    \item Room cable
  \end{enumerate}
\item \textit{[JKSSB-pattern]} \textbf{(MSQ: select ALL correct options.)} Which devices forward network data?
  \begin{enumerate}[label=\Alph*.]
    \item Hub
    \item Switch
    \item Router
    \item Toaster
  \end{enumerate}
\item \textit{[SSC CGL-pattern]} \textbf{(MSQ: select ALL correct options.)} Which statements about networks are correct?
  \begin{enumerate}[label=\Alph*.]
    \item A LAN covers a small area such as an office.
    \item A WAN spans large distances.
    \item A router joins different networks.
    \item A printer routes packets between networks.
  \end{enumerate}
\item \textit{[PYQ-adapted: JKSSB FAA 2022, QCRIAS item 547; SRC-21]} The Internet is an example of which network type by geographic scope?
  \begin{enumerate}[label=\Alph*.]
\item LAN
   \item MAN
   \item WAN
   \item PAN
  \end{enumerate}

\item \textit{[PYQ-adapted: JKSSB Sub-Inspector 2022, QCRIAS item 408; SRC-17/SRC-22]} Which device directs traffic between networks?
  \begin{enumerate}[label=\Alph*.]
\item Router
   \item Hub
   \item Keyboard
   \item Plotter
  \end{enumerate}

\item \textit{[Original]} Which topology connects each node to a central device?
  \begin{enumerate}[label=\Alph*.]
\item Star
   \item Ring
   \item Bus
   \item Mesh only
  \end{enumerate}

\item \textit{[Original]} A device that connects a computer to a network is commonly a:
  \begin{enumerate}[label=\Alph*.]
\item NIC
   \item Plotter
   \item OCR reader
   \item Projector
  \end{enumerate}

\item \textit{[Original, MSQ drill]} \textbf{(MSQ: select ALL correct options.)} Which statements are correct?
  \begin{enumerate}[label=\Alph*.]
\item A LAN covers a limited local area.
   \item A router forwards traffic between networks.
   \item A switch commonly connects devices within a LAN.
   \item A hub makes routing decisions using IP addresses.
  \end{enumerate}

\end{enumerate}
\lessonheading{41}{Networking Protocols}
\begin{enumerate}[label=\textbf{Q\arabic*.}, leftmargin=*]
\item \textit{[Original]} Expand TCP.
  \begin{enumerate}[label=\Alph*.]
    \item Transfer Control Program
    \item Transmission Control Protocol
    \item Transport Communication Process
    \item Tele Communication Procedure
  \end{enumerate}
\item \textit{[Original]} Expand DNS.
  \begin{enumerate}[label=\Alph*.]
    \item Domain Name System
    \item Data Network Service
    \item Digital Number Sequence
    \item Direct Node Setup
  \end{enumerate}
\item \textit{[Original]} Which pair correctly matches the transport protocol
  with its behaviour?
  \begin{enumerate}[label=\Alph*.]
    \item TCP --- fast, connectionless, no retransmission
    \item UDP --- reliable, ordered, connection-oriented
    \item TCP --- reliable, ordered, connection-oriented
    \item UDP --- ordered delivery with guaranteed timing
  \end{enumerate}
\item \textit{[Original]} TCP guarantees:
  \begin{enumerate}[label=\Alph*.]
    \item fast arrival time of every packet.
    \item complete and ordered delivery of bytes.
    \item encryption of the message body.
    \item translation of names into addresses.
  \end{enumerate}
\item \textit{[Original]} Which of the following is \textbf{NOT} true of UDP?
  \begin{enumerate}[label=\Alph*.]
    \item It is connectionless.
    \item It offers best-effort delivery without retransmission.
    \item It guarantees ordered delivery of all datagrams.
    \item It is preferred for streaming and DNS queries.
  \end{enumerate}
\item \textit{[Original]} IPv4 uses \underline{\hspace{2.2cm}} addresses,
  whereas IPv6 uses \underline{\hspace{2.2cm}} addresses.
  \begin{enumerate}[label=\Alph*.]
    \item 128-bit; 32-bit
    \item 32-bit; 128-bit
    \item 64-bit; 256-bit
    \item 32-bit; 256-bit
  \end{enumerate}
\item \textit{[Original]} Which statement about IPv6 is correct?
  \begin{enumerate}[label=\Alph*.]
    \item IPv6 addresses are 256-bit long.
    \item IPv6 uses four decimal octets.
    \item IPv6 addresses are 128-bit long with a fixed 40-byte header.
    \item IPv6 headers vary between 20 and 60 bytes.
  \end{enumerate}
\item \textit{[Original]} ICMP is best described as:
  \begin{enumerate}[label=\Alph*.]
    \item a file-transfer protocol with its own port number.
    \item a control and diagnostic protocol carrying errors and probes.
    \item a mail-sending protocol for outgoing messages.
    \item an encryption layer for secure web pages.
  \end{enumerate}
\item \textit{[Original]} Which tool pair belongs to ICMP?
  \begin{enumerate}[label=\Alph*.]
    \item ping and traceroute
    \item mail merge and autosum
    \item FTP and Telnet
    \item ping and mail merge
  \end{enumerate}
\item \textit{[Original]} The service that translates
  \texttt{jkssb.nic.in} into an IP address is:
  \begin{enumerate}[label=\Alph*.]
    \item DHCP
    \item ICMP
    \item DNS
    \item SMTP
  \end{enumerate}
\item \textit{[Original]} A newly connected office desktop receives its IP
  address, mask, gateway, and DNS server automatically. This service is:
  \begin{enumerate}[label=\Alph*.]
    \item DNS
    \item DHCP
    \item ICMP
    \item FTP
  \end{enumerate}
\item \textit{[Original]} Which of the following is correctly matched?
  \begin{enumerate}[label=\Alph*.]
    \item HTTP --- port 443, encrypted with SSL/TLS
    \item HTTPS --- port 80, plain text
    \item HTTP --- port 80, plain text
    \item HTTPS --- port 25, used for sending mail
  \end{enumerate}
\item \textit{[Original]} SMTP is used to:
  \begin{enumerate}[label=\Alph*.]
    \item send outgoing mail to the server.
    \item download received mail to the local machine.
    \item keep mail synced on the server.
    \item translate a domain name into an address.
  \end{enumerate}
\item \textit{[Original]} Which pairing correctly matches the mail protocol
  with its role?
  \begin{enumerate}[label=\Alph*.]
    \item SMTP --- receives and syncs mail on the server
    \item POP3 --- downloads received mail to the local machine
    \item IMAP --- sends outgoing mail to the server
    \item SMTP --- translates names into IP addresses
  \end{enumerate}
\item \textit{[Original]} Evaluate the statements:
  \begin{enumerate}[label=\Roman*.]
    \item TCP is connection-oriented and delivers bytes in order.
    \item IPv6 addresses are 128-bit long.
    \item ICMP transfers application data over a dedicated port.
  \end{enumerate}
  \begin{enumerate}[label=\Alph*.]
    \item I only
    \item I and II only
    \item II and III only
    \item I, II, and III
  \end{enumerate}
\item \textit{[Original]} Which layer pair is correctly matched in the OSI
  model?
  \begin{enumerate}[label=\Alph*.]
    \item TCP and UDP --- Network layer
    \item IP and ICMP --- Transport layer
    \item HTTP, FTP, SMTP, DNS --- Application layer
    \item IP and HTTP --- Physical layer
  \end{enumerate}
\item \textit{[Original]} In the OSI model, the correct order from layer 1 to
  layer 7 is:
  \begin{enumerate}[label=\Alph*.]
    \item Application, Presentation, Session, Transport, Network, Data Link,
      Physical
    \item Physical, Data Link, Network, Transport, Session, Presentation,
      Application
    \item Physical, Network, Data Link, Session, Transport, Application,
      Presentation
    \item Data Link, Physical, Transport, Network, Application, Session,
      Presentation
  \end{enumerate}
\item \textit{[Original]} Which port pairing is correctly matched?
  \begin{enumerate}[label=\Alph*.]
    \item FTP --- 23; Telnet --- 21
    \item SMTP --- 25; DNS --- 53
    \item HTTP --- 443; HTTPS --- 80
    \item POP3 --- 143; IMAP --- 110
  \end{enumerate}
\item \textit{[Original]} Which of the following is \textbf{INCORRECT}?
  \begin{enumerate}[label=\Alph*.]
    \item FTP moves files between computers (port 21).
    \item Telnet opens a plain-text remote login session (port 23).
    \item DHCP assigns addresses to devices joining the network.
    \item ICMP uses port 80 to transfer web pages.
  \end{enumerate}
\item \textit{[Original]} An office PC cannot open any website, but a ping to
  the gateway succeeds. Which conclusion follows?
  \begin{enumerate}[label=\Alph*.]
    \item The network path works, so the fault lies above it (DNS, proxy,
      or application settings).
    \item The ping proves the web server has sent the page.
    \item ICMP has delivered the web page to the browser.
    \item The Translate layer of the OSI model has failed.
  \end{enumerate}

\item \textit{[JKSSB FAA-pattern]} TCP/IP is best described as:
  \begin{enumerate}[label=\Alph*.]
    \item the core protocol suite of the Internet.
    \item a printer ink formula.
    \item a type of monitor glass.
    \item a keyboard layout.
  \end{enumerate}
\item \textit{[JKPSI-pattern]} An IP address identifies a:
  \begin{enumerate}[label=\Alph*.]
    \item device on a network.
    \item printed page number.
    \item speaker's volume.
    \item book's author.
  \end{enumerate}
\item \textit{[SSC CGL-pattern]} IPv4 addresses are how many bits long?
  \begin{enumerate}[label=\Alph*.]
    \item 32
    \item 8
    \item 128
    \item 256
  \end{enumerate}
\item \textit{[JKSSB-pattern]} \textbf{(MSQ: select ALL correct options.)} Which are network protocols?
  \begin{enumerate}[label=\Alph*.]
    \item TCP/IP
    \item HTTP
    \item FTP
    \item Stapler pin
  \end{enumerate}
\item \textit{[SSC CGL-pattern]} \textbf{(MSQ: select ALL correct options.)} Which statements about protocols are correct?
  \begin{enumerate}[label=\Alph*.]
    \item TCP gives reliable ordered delivery.
    \item IP handles addressing and routing.
    \item FTP transfers files between systems.
    \item An IP address identifies a printed page.
  \end{enumerate}
\item \textit{[PYQ-adapted: JKSSB JKP Constable 2024, QCRIAS item 747; SRC-02]} In HTTPS, what does the letter S stand for?
  \begin{enumerate}[label=\Alph*.]
\item Selected
   \item Secure
   \item Software
   \item System
  \end{enumerate}

\item \textit{[PYQ-adapted: JKSSB JKP Constable 2024, QCRIAS item 765; SRC-02]} Which organization develops standards for the World Wide Web?
  \begin{enumerate}[label=\Alph*.]
\item W3C
   \item WHO
   \item OPEC
   \item UNESCO
  \end{enumerate}

\item \textit{[Original]} Which protocol provides reliable, ordered delivery between applications?
  \begin{enumerate}[label=\Alph*.]
\item TCP
   \item UDP
   \item ARP
   \item ICMP
  \end{enumerate}

\item \textit{[Original]} Which protocol translates a domain name into an IP address?
  \begin{enumerate}[label=\Alph*.]
\item DNS
   \item SMTP
   \item FTP
   \item POP3
  \end{enumerate}

\item \textit{[Original, MSQ drill]} \textbf{(MSQ: select ALL correct options.)} Which are application-layer protocols/services?
  \begin{enumerate}[label=\Alph*.]
\item HTTP
   \item SMTP
   \item FTP
   \item Ethernet framing
  \end{enumerate}

\end{enumerate}
\lessonheading{62}{Current Internet, Networking and Cloud IaaS}
\begin{enumerate}[label=\textbf{Q\arabic*.}, leftmargin=*]
\item \textit{[Original]} Expand TCP.
  \begin{enumerate}[label=\Alph*.]
    \item Transfer Control Program
    \item Transmission Control Protocol
    \item Transport Communication Process
    \item Tele Communication Protocol
  \end{enumerate}
\item \textit{[Original]} Which statement correctly describes TCP delivery
  guarantees?
  \begin{enumerate}[label=\Alph*.]
    \item TCP guarantees ordered and reliable delivery, but provides no timing
      guarantee.
    \item TCP guarantees delivery within a fixed time but not the order.
    \item TCP delivers datagrams without setup and without retransmission.
    \item TCP transfers only control messages and uses no acknowledgements.
  \end{enumerate}
\item \textit{[Original]} Which pairing of transport protocol and property is
  correctly matched?
  \begin{enumerate}[label=\Alph*.]
    \item UDP --- guarantees ordered delivery of datagrams
    \item TCP --- connectionless with no retransmission
    \item UDP --- connectionless, fast, with no ordering or reliability
      guarantee
    \item TCP --- transfers only diagnostic messages using no ports
  \end{enumerate}
\item \textit{[Original]} Expand ICMP.
  \begin{enumerate}[label=\Alph*.]
    \item Internet Control Message Protocol
    \item Internal Computer Mail Protocol
    \item Internet Communication Management Process
    \item Inter-Connected Media Protocol
  \end{enumerate}
\item \textit{[Original]} Which of the following is \textbf{NOT} true of ICMP?
  \begin{enumerate}[label=\Alph*.]
    \item It carries control and diagnostic messages about the network.
    \item It supports tools such as ping and traceroute.
    \item It transfers application data such as files and web pages.
    \item It reports errors such as destination unreachable.
  \end{enumerate}
\item \textit{[Original]} Expand BGP.
  \begin{enumerate}[label=\Alph*.]
    \item Basic Gateway Process
    \item Border Gateway Protocol
    \item Broadband Group Protocol
    \item Binary Gateway Program
  \end{enumerate}
\item \textit{[Original]} BGP is best described as:
  \begin{enumerate}[label=\Alph*.]
    \item an inter-domain routing protocol that routes between autonomous
      systems.
    \item a protocol for routing inside a single small local network.
    \item a mail-transfer protocol that downloads mail to the local device.
    \item a diagnostic protocol that carries no routing information.
  \end{enumerate}
\item \textit{[Original]} Which statement correctly distinguishes BGP from
  OSPF and RIP?
  \begin{enumerate}[label=\Alph*.]
    \item BGP routes between autonomous systems; OSPF and RIP route inside a
      single domain.
    \item OSPF routes between autonomous systems; BGP routes inside one LAN.
    \item RIP is inter-domain; BGP and OSPF are mail protocols.
    \item All three transfer application data over ports.
  \end{enumerate}
\item \textit{[Original]} IPv6 addresses are:
  \begin{enumerate}[label=\Alph*.]
    \item 32-bit long.
    \item 64-bit long.
    \item 128-bit long.
    \item 256-bit long.
  \end{enumerate}
\item \textit{[Original]} Which statement about IPv6 is correct?
  \begin{enumerate}[label=\Alph*.]
    \item IPv6 uses a fixed 40-byte header and has no broadcast address.
    \item IPv6 uses a fixed 20-byte header and relies on broadcast delivery.
    \item IPv6 uses 256-bit addresses with a broadcast address for all nodes.
    \item IPv6 is a routing protocol that runs between autonomous systems.
  \end{enumerate}
\item \textit{[Original]} The reason IPv6 removes the need for NAT is:
  \begin{enumerate}[label=\Alph*.]
    \item its vastly larger 128-bit address space.
    \item its use of TCP instead of UDP.
    \item its support for broadcast delivery.
    \item its ability to transfer application data.
  \end{enumerate}
\item \textit{[Original]} HTTP/3 differs from HTTP/1.1 and HTTP/2 in that it:
  \begin{enumerate}[label=\Alph*.]
    \item runs over TCP like the earlier versions.
    \item runs over QUIC over UDP instead of TCP.
    \item transfers only control messages with no ports.
    \item routes between autonomous systems on the backbone.
  \end{enumerate}
\item \textit{[Original]} Expand DNS.
  \begin{enumerate}[label=\Alph*.]
    \item Domain Name System
    \item Data Network Service
    \item Digital Number Sequence
    \item Direct Node Switch
  \end{enumerate}
\item \textit{[Original]} Which statement correctly describes DNS?
  \begin{enumerate}[label=\Alph*.]
    \item It translates domain names into IP addresses through a hierarchy of
      servers.
    \item It routes packets between autonomous systems on the backbone.
    \item It downloads mail from the server to the local device.
    \item It encrypts message bodies for confidentiality.
  \end{enumerate}
\item \textit{[Original]} Match each Internet-governance actor to its role:
  \begin{enumerate}[label=\Alph*.]
    \item PIR --- operates the .org registry; ICANN --- coordinates domain
      names and IP addresses globally.
    \item PIR --- coordinates all IP addresses; ICANN --- operates the .org
      registry.
    \item PIR --- routes between autonomous systems; ICANN --- transfers mail.
    \item PIR --- assigns TCP ports; ICANN --- sends diagnostic messages.
  \end{enumerate}
\item \textit{[Original]} ARPANET is best described as:
  \begin{enumerate}[label=\Alph*.]
    \item the 1969 packet-switched ancestor of the Internet.
    \item a circuit-switched telephone network with dedicated calls.
    \item a cloud service model managing networking hardware.
    \item a mail protocol that syncs mail across devices.
  \end{enumerate}
\item \textit{[Original]} Which statement correctly distinguishes packet
  switching from circuit switching?
  \begin{enumerate}[label=\Alph*.]
    \item Packet switching splits data into addressed packets taking possibly
      different routes; circuit switching holds one dedicated path.
    \item Circuit switching splits data into packets; packet switching holds a
      dedicated circuit.
    \item Both hold a dedicated path for the whole session.
    \item Both deliver packets with guaranteed timing.
  \end{enumerate}
\item \textit{[Original]} Consider the assertion and reason about a router.
  Assertion (A): A router connects different networks and forwards packets
  between them. Reason (R): A router reads the destination IP address and
  chooses the next hop from its routing table.
  \begin{enumerate}[label=\Alph*.]
    \item Both A and R are true, and R is the correct explanation of A.
    \item Both A and R are true, but R is not the correct explanation of A.
    \item A is true but R is false.
    \item A is false but R is true.
  \end{enumerate}
\item \textit{[Original]} Evaluate the statements:
  \begin{enumerate}[label=\Roman*.]
    \item TCP guarantees ordered and reliable delivery but no timing
      guarantee.
    \item ICMP transfers application data such as files and web pages.
    \item In IaaS, the provider manages the networking hardware while the
      customer manages the operating system, applications, and data.
  \end{enumerate}
  \begin{enumerate}[label=\Alph*.]
    \item I and III only
    \item II only
    \item I, II, and III
    \item I and II only
  \end{enumerate}
\item \textit{[Original]} Which pairing is correctly matched?
  \begin{enumerate}[label=\Alph*.]
    \item BGP --- routes inside one LAN; ICMP --- transfers files
    \item HTTP/3 --- uses QUIC over UDP; IPv6 --- 128-bit with no broadcast
    \item IaaS provider --- manages the customer OS and data; PIR ---
      coordinates TCP ports
    \item UDP --- guarantees order; TCP --- connectionless and fast
  \end{enumerate}

\item \textit{[JKSSB FAA-pattern]} In cloud computing, IaaS stands for:
  \begin{enumerate}[label=\Alph*.]
    \item Infrastructure as a Service
    \item Internet and Speaker
    \item Ink and Stationery
    \item Input and Sound
  \end{enumerate}
\item \textit{[JKPSI-pattern]} In the IaaS model, the provider manages the:
  \begin{enumerate}[label=\Alph*.]
    \item physical infrastructure (servers, storage, networking).
    \item user's lunch menu.
    \item office wall paint.
    \item printer paper stock.
  \end{enumerate}
\item \textit{[SSC CGL-pattern]} IPv6 addresses are how many bits long?
  \begin{enumerate}[label=\Alph*.]
    \item 128
    \item 32
    \item 8
    \item 16
  \end{enumerate}
\item \textit{[JKSSB-pattern]} \textbf{(MSQ: select ALL correct options.)} Which are cloud service models?
  \begin{enumerate}[label=\Alph*.]
    \item IaaS
    \item PaaS
    \item SaaS
    \item Paper-as-a-Service
  \end{enumerate}
\item \textit{[SSC CGL-pattern]} \textbf{(MSQ: select ALL correct options.)} Which statements about cloud and networking are correct?
  \begin{enumerate}[label=\Alph*.]
    \item IaaS rents virtualised infrastructure.
    \item SaaS delivers ready-to-use software online.
    \item BGP routes traffic between networks.
    \item IPv6 uses 32-bit addresses.
  \end{enumerate}
\item \textit{[PYQ-adapted: JKSSB Junior Assistant 2026 Q78, final key; SRC-03/SRC-06]} How many bits are in an IPv6 address?
  \begin{enumerate}[label=\Alph*.]
\item 32
   \item 64
   \item 128
   \item 256
  \end{enumerate}

\item \textit{[PYQ-adapted: SSC CGL 2025 Tier-II Paper-I, 19-Jan-2026 Q3; SRC-18]} Which OSI layer routes packets across networks using logical addressing?
  \begin{enumerate}[label=\Alph*.]
\item Data Link
   \item Session
   \item Network
   \item Transport
  \end{enumerate}

\item \textit{[PYQ-adapted: JKSSB Junior Assistant 2026 Q79, final key; SRC-03/SRC-06]} HTTP/3 uses which transport protocol through QUIC?
  \begin{enumerate}[label=\Alph*.]
\item UDP
   \item TCP
   \item ICMP
   \item ARP
  \end{enumerate}

\item \textit{[Original]} Which cloud service model gives a customer virtual machines, storage, and networking while the customer manages the operating system?
  \begin{enumerate}[label=\Alph*.]
\item IaaS
   \item SaaS
   \item PaaS
   \item BIOS
  \end{enumerate}

\item \textit{[Original, MSQ drill]} \textbf{(MSQ: select ALL correct options.)} Which statements are correct?
  \begin{enumerate}[label=\Alph*.]
\item IPv6 uses 128-bit addresses.
   \item Routers forward packets between networks.
   \item UDP is connectionless.
   \item HTTP/3 uses TCP through QUIC.
  \end{enumerate}

\end{enumerate}
\answersection{9}
\answersubheading{40}{Networking Basics}
\begin{enumerate}[label=\textbf{Q\arabic*.}, leftmargin=*]
\item \textbf{C.} NIC stands for Network Interface Card. A, B, and D are not
  expansions of the abbreviation.
\item \textbf{B.} LAN stands for Local Area Network. A, C, and D are not the
  correct full forms.
\item \textbf{B.} The client requests a service and the server provides it.
  A swaps the two roles; C describes neither role (forwarding without
  storing is not client--server); D swaps in the NIC and firewall roles.
\item \textbf{B.} One building or campus is a LAN. A PAN (A) covers one
  person, a MAN (C) covers one city, and a WAN (D) spans countries or the
  world.
\item \textbf{C.} Branches linked across one city form a MAN. A PAN (A) is
  one person's devices, a LAN (B) is one building, and D misdescribes both
  the medium and the range of a WLAN.
\item \textbf{C.} The worldwide public network is the Internet, not an
  intranet. Being private to one organisation (A), restricted to its staff
  (B), and carrying the internal portal (D) are all true of an intranet.
\item \textbf{B.} A star joins every node to one central device (hub or
  switch). A bus (A) shares one cable, a ring (C) joins neighbours in a
  circle, and a mesh (D) uses many direct links.
\item \textbf{A.} In a bus every node shares one main cable, so one break
  fails the whole network. A star (B) isolates a spoke failure, a mesh (C)
  reroutes over other paths, and a tree (D) fails branch-wise on the
  backbone, not on one shared cable.
\item \textbf{B.} In a ring each node joins exactly two neighbours in a
  closed circle. A describes a star, C describes a mesh, and D confuses the
  topology with a transmission medium.
\item \textbf{C.} A mesh is the most reliable and the costliest because of
  its many direct links. A bus (A) is cheapest, a star (B) uses one centre,
  and a ring (D) is a single circle.
\item \textbf{D.} Stars joined through a central backbone form a tree. A bus
  (A) is one shared cable, a ring (B) is a neighbour circle, and a mesh (C)
  is many direct links.
\item \textbf{A.} Twisted pair, coaxial, and fibre-optic are the three guided
  (wired) media. B, C, and D each list unguided (wireless) media: radio
  waves, microwaves, satellite links, infrared, or Bluetooth.
\item \textbf{C.} Fibre-optic cable carries light signals and is the fastest
  of the three cables. Twisted pair (A) is copper and slowest, coaxial (B) is
  shielded copper, and D combines two unrelated items.
\item \textbf{A.} A hub broadcasts every incoming signal to all ports. A
  switch (B) forwards only to the addressed port, a router (C) paths packets
  between networks, and a firewall (D) blocks unauthorised access.
\item \textbf{C.} A switch forwards only to the addressed port. A hub (A)
  broadcasts to all, a repeater (B) only regenerates the signal, and a modem
  (D) converts digital--analogue for the line.
\item \textbf{B.} A repeater regenerates a weak signal without choosing a
  path. A router (A) chooses paths, a gateway (C) translates protocols, and
  a firewall (D) guards access.
\item \textbf{D.} A router joins different networks and chooses the best
  path per packet. A hub (A) broadcasts inside one network, a switch (B)
  forwards inside one network, and a bridge (C) joins two LAN segments.
\item \textbf{D.} All three statements are correct: a bridge joins two LAN
  segments with filtering (I), a gateway translates between different
  protocols (II), and a firewall blocks unauthorised access (III). A omits
  II and III; B omits III; C omits I.
\item \textbf{B.} An access point lets Wi-Fi devices join the wired network.
  A swaps in the firewall's role for the modem, C swaps in the repeater's
  role for the router, and D swaps in the gateway's role for the NIC.
\item \textbf{B.} A modem (modulator-demodulator) converts between digital
  computer signals and analogue line signals. A switch (A) forwards inside
  the LAN, a bridge (C) joins LAN segments, and an access point (D) admits
  Wi-Fi devices.

\item \textbf{A.} LAN is Local Area Network. B, C, and D are not the recognised expansion.
\item \textbf{A.} Routers link networks. B, C, and D are peripherals, not inter-network devices.
\item \textbf{A.} WANs cover the widest areas. LAN (B) and PAN (C) are smaller; D is not a network class.
\item \textbf{A, B, C.} Hubs, switches, and routers move network data. A toaster (D) does not.
\item \textbf{A, B, C.} The LAN/WAN split and router role are correct. D is false: printers do not route packets.
\item \textbf{C.} The Internet spans large geographic areas and interconnected networks, so it is a WAN. LAN, MAN, and PAN cover smaller scopes.

\item \textbf{A.} A router forwards packets between networks. A hub repeats local signals; keyboard is input and plotter is output.

\item \textbf{A.} In a star topology each node connects to a central hub/switch. Ring and bus use different arrangements; mesh uses multiple inter-node paths.

\item \textbf{A.} A Network Interface Card (NIC) provides network connectivity. The others are output or recognition devices.

\item \textbf{A, B, C.} LAN, router, and switch roles are correctly described. A hub repeats signals and does not route using IP addresses.

\end{enumerate}
\answersubheading{41}{Networking Protocols}
\begin{enumerate}[label=\textbf{Q\arabic*.}, leftmargin=*]
\item \textbf{B.} TCP stands for Transmission Control Protocol. A, C, and D
  are not the correct expansion.
\item \textbf{A.} DNS stands for Domain Name System. B, C, and D are not the
  correct full forms.
\item \textbf{C.} TCP is reliable, ordered, and connection-oriented. A
  describes UDP, not TCP; B describes TCP, not UDP; D wrongly promises UDP
  timing guarantees.
\item \textbf{B.} TCP guarantees complete, ordered delivery through
  sequencing and retransmission; it does not promise speed (A), encryption
  (C), or name resolution (D).
\item \textbf{C.} UDP gives no ordering guarantee, so this statement is NOT
  true. Connectionless operation (A), best-effort delivery (B), and use for
  streaming and DNS (D) are all true of UDP.
\item \textbf{B.} IPv4 uses 32-bit addresses and IPv6 uses 128-bit addresses
  (TRM-10). A reverses them; C and D use the 256-bit distractor (TRAP-18).
\item \textbf{C.} IPv6 is 128-bit with a fixed 40-byte header (PYQ-18).
  256-bit (A) is the known trap; four decimal octets (B) describe IPv4; a
  20--60 byte header (D) describes IPv4.
\item \textbf{B.} ICMP carries control and diagnostic messages (errors,
  reachability) only. It is not a file-transfer (A), mail-sending (C), or
  encryption (D) protocol.
\item \textbf{A.} Ping and traceroute are the two ICMP diagnostic tools. Mail
  merge and autosum (B, D) are Office features; FTP and Telnet (C) are
  separate application protocols.
\item \textbf{C.} DNS translates a domain name into its IP address. DHCP (A)
  assigns addresses; ICMP (B) reports errors; SMTP (D) sends mail.
\item \textbf{B.} Automatically handing out the address, mask, gateway, and
  DNS server is DHCP. DNS (A) resolves names; ICMP (C) diagnoses; FTP (D)
  moves files.
\item \textbf{C.} HTTP is port 80 and plain text. A gives HTTP the HTTPS
  port and property; B reverses the pair; D gives HTTPS the SMTP port and
  role.
\item \textbf{A.} SMTP (Simple Mail Transfer Protocol) sends outgoing mail.
  Downloading (B) is POP3; syncing on the server (C) is IMAP; name
  translation (D) is DNS.
\item \textbf{B.} POP3 downloads received mail to the local machine. A gives
  SMTP the IMAP role; C gives IMAP the SMTP role; D gives SMTP the DNS role.
\item \textbf{B.} Statements I and II are correct: TCP is connection-oriented
  and ordered, and IPv6 is 128-bit. III is false because ICMP transfers no
  application data and uses no ports (TRAP-19). A omits II; C includes false
  III and omits I; D includes false III.
\item \textbf{C.} HTTP, FTP, SMTP, and DNS are Application-layer (layer 7)
  protocols. TCP and UDP belong to the Transport layer, not the Network
  layer (A); IP and ICMP belong to the Network layer, not Transport (B); the
  Physical layer (D) carries bits, not IP or HTTP.
\item \textbf{B.} The OSI order is Physical, Data Link, Network, Transport,
  Session, Presentation, Application. A is the exact reverse; C and D
  scramble the middle layers.
\item \textbf{B.} SMTP is port 25 and DNS is port 53. A swaps FTP (21) and
  Telnet (23); C swaps HTTP (80) and HTTPS (443); D swaps POP3 (110) and
  IMAP (143).
\item \textbf{D.} ICMP uses no ports and transfers no web pages, so this is
  INCORRECT. FTP on 21 (A), Telnet on 23 (B), and DHCP address assignment
  (C) are all correct.
\item \textbf{A.} A successful ping proves the Network-layer path works
  (ICMP control traffic), so the fault lies above it: DNS, proxy, or
  application settings. B and C wrongly treat ICMP as page delivery; D names
  a layer that does not exist.

\item \textbf{A.} TCP/IP carries Internet traffic. B, C, and D name unrelated materials or layouts.
\item \textbf{A.} IPs label network devices. B, C, and D are non-network identifiers.
\item \textbf{A.} IPv4 uses 32 bits. B is too short; C is IPv6's length; D is not an IP size.
\item \textbf{A, B, C.} TCP/IP, HTTP, and FTP are protocols. A stapler pin (D) is not.
\item \textbf{A, B, C.} The TCP, IP, and FTP roles are correct. D is false: IPs identify devices, not pages.
\item \textbf{B.} The S in HTTPS means Secure. The other options are not expansions of HTTPS.

\item \textbf{A.} W3C is the World Wide Web Consortium. The other organizations have different mandates.

\item \textbf{A.} TCP provides reliable, ordered transport. UDP is connectionless without that delivery guarantee; ARP and ICMP serve other network functions.

\item \textbf{A.} DNS resolves names. SMTP sends mail, FTP transfers files, and POP3 retrieves email.

\item \textbf{A, B, C.} HTTP, SMTP, and FTP are application-layer protocols. Ethernet framing belongs to the data-link layer.

\end{enumerate}
\answersubheading{62}{Current Internet, Networking and Cloud IaaS}
\begin{enumerate}[label=\textbf{Q\arabic*.}, leftmargin=*]
\item \textbf{B.} TCP stands for Transmission Control Protocol. A, C, and D
  are not the correct expansion.
\item \textbf{A.} TCP gives ordered and reliable delivery with no timing
  promise (PYQ-18). B invents a timing guarantee and drops ordering; C
  describes UDP; D describes ICMP.
\item \textbf{C.} UDP is connectionless and fast with neither ordering nor
  reliability. A moves a TCP property onto UDP; B reverses TCP into
  connectionless; D describes ICMP, not TCP.
\item \textbf{A.} ICMP stands for Internet Control Message Protocol. B, C,
  and D are not the correct full forms.
\item \textbf{C.} ICMP never transfers application data such as files or web
  pages; it carries control and diagnostic messages only (TRAP-19). A, B,
  and D are all true of ICMP.
\item \textbf{B.} BGP stands for Border Gateway Protocol. A, C, and D are not
  the correct expansion.
\item \textbf{A.} BGP is the inter-domain protocol routing between autonomous
  systems (PYQ-19, TRAP-20). B describes an interior protocol; C describes a
  mail protocol; D denies that BGP carries routing information.
\item \textbf{A.} BGP is inter-domain while OSPF and RIP are intra-domain.
  B swaps the scopes; C mislabels all three; D claims routing protocols
  transfer application data.
\item \textbf{C.} IPv6 addresses are 128-bit long (PYQ-18, TRAP-18). A is the
  IPv4 size; B is an intermediate distractor; D doubles the correct size.
\item \textbf{A.} IPv6 has a fixed 40-byte header and no broadcast address
  (multicast replaces it). B gives the IPv4 header size and a broadcast
  claim; C uses the 256-bit distractor with a broadcast; D mislabels IPv6 as
  a routing protocol.
\item \textbf{A.} The 128-bit space is so large that address-sharing via NAT
  is unnecessary (PYQ-19). B, C, and D name mechanisms unrelated to address
  supply.
\item \textbf{B.} HTTP/3 runs over QUIC over UDP, unlike HTTP/1.1 and HTTP/2
  over TCP (PYQ-19, TRAP-21). A states the old stack; C describes ICMP; D
  describes BGP.
\item \textbf{A.} DNS stands for Domain Name System. B, C, and D are not the
  correct expansion.
\item \textbf{A.} DNS maps names to IP addresses hierarchically. B describes
  BGP; C describes POP3; D describes encryption.
\item \textbf{A.} PIR operates the .org registry while ICANN coordinates
  names and numbers globally. B swaps the two roles; C and D assign routing,
  mail, port, and diagnostic roles neither body holds.
\item \textbf{A.} ARPANET (1969) is the packet-switched ancestor of the
  Internet. B describes circuit-switched telephony; C describes IaaS; D
  describes a mail protocol.
\item \textbf{A.} Packet switching uses independent addressed packets while
  circuit switching holds one dedicated path. B swaps the two; C and D
  describe properties neither model has as stated.
\item \textbf{A.} Both the assertion and the reason are true, and reading the
  destination IP from the routing table is exactly how the router forwards
  between networks, so R correctly explains A. B denies the explanatory link;
  C and D each falsely mark one part false. A switch (MAC, one network) or a
  hub (repeats all) fits neither part.
\item \textbf{A.} Statements I and III are correct: TCP gives order and
  reliability without timing, and IaaS splits responsibility with networking
  hardware at the provider (PYQ-18, PYQ-20, TRAP-22). II is false because
  ICMP never transfers application data (TRAP-19). B selects only the false
  statement; C includes it; D pairs a true statement with the false one.
\item \textbf{B.} HTTP/3 uses QUIC over UDP and IPv6 is 128-bit with no
  broadcast, so both halves match. A swaps BGP inside a LAN and gives ICMP a
  file-transfer role; C gives the IaaS provider the customer OS and data and
  misstates PIR's role; D swaps the TCP/UDP properties.

\item \textbf{A.} IaaS is Infrastructure as a Service. B, C, and D are not the recognised expansion.
\item \textbf{A.} Providers own the physical layer in IaaS. B, C, and D are non-cloud chores.
\item \textbf{A.} IPv6 uses 128 bits. B is IPv4's length; C and D are not IP sizes.
\item \textbf{A, B, C.} IaaS, PaaS, and SaaS are the three models. D is not a service model.
\item \textbf{A, B, C.} The IaaS, SaaS, and BGP facts are correct. D is false: IPv6 is 128-bit.
\item \textbf{C.} IPv6 addresses are 128 bits. IPv4 addresses are 32 bits.

\item \textbf{C.} The Network layer provides logical addressing and routing. Data Link handles local frames; Session manages dialogues; Transport provides end-to-end services.

\item \textbf{A.} HTTP/3 uses QUIC over UDP. TCP is used by earlier HTTP versions; ICMP and ARP have other network roles.

\item \textbf{A.} IaaS provides infrastructure resources; the customer manages the guest OS and applications. SaaS delivers an application; PaaS provides a managed platform.

\item \textbf{A, B, C.} The first three are correct. QUIC runs over UDP, not TCP.

\end{enumerate}
\unitopener{10}{Cybersecurity \& IT Law}
\lessonheading{42}{Malware}
\begin{enumerate}[label=\textbf{Q\arabic*.}, leftmargin=*]
\item \textit{[Original]} The term ``malware'' is short for:
  \begin{enumerate}[label=\Alph*.]
    \item managed software
    \item malicious software
    \item manual software
    \item mainframe software
  \end{enumerate}
\item \textit{[Original]} Which malware attaches itself to a host program or
  file and needs the user to open that file before it can run?
  \begin{enumerate}[label=\Alph*.]
    \item Virus
    \item Worm
    \item Bot
    \item Firewall
  \end{enumerate}
\item \textit{[Original]} Which malware spreads by itself over a network,
  with no host file and no human action?
  \begin{enumerate}[label=\Alph*.]
    \item Trojan
    \item Virus
    \item Worm
    \item Adware
  \end{enumerate}
\item \textit{[Original]} A program pretends to be a free game. The user
  installs it willingly, and the hidden malicious part then steals data. It
  does not replicate itself. This is a:
  \begin{enumerate}[label=\Alph*.]
    \item worm
    \item virus
    \item Trojan
    \item logic bomb
  \end{enumerate}
\item \textit{[Original]} Which option correctly matches each malware to its
  spreading method?
  \begin{enumerate}[label=\Alph*.]
    \item Virus --- self-spreading with no host; worm --- needs a host file
      and a click; Trojan --- floods the screen with ads
    \item Virus --- needs a host file and a user action; worm --- standalone
      and self-spreading; Trojan --- disguised as legitimate, no
      self-replication
    \item Virus --- disguised as legitimate software; worm --- demands
      payment; Trojan --- records keystrokes
    \item Virus --- hides the attacker deep in the system; worm --- mines
      cryptocurrency; Trojan --- waits for a trigger date
  \end{enumerate}
\item \textit{[Original]} Which of the following is \textbf{NOT} true of a
  Trojan?
  \begin{enumerate}[label=\Alph*.]
    \item It pretends to be useful or legitimate software.
    \item It relies on the user installing or opening it.
    \item It replicates itself automatically across networks.
    \item It can open a backdoor or install further malware.
  \end{enumerate}
\item \textit{[Original]} A clerk opens an invoice attachment; all documents
  turn unreadable and a message demands payment for the unlock key. This is:
  \begin{enumerate}[label=\Alph*.]
    \item ransomware
    \item adware
    \item worm
    \item rootkit
  \end{enumerate}
\item \textit{[Original]} Which malware secretly observes the user (browsing
  habits, logins, personal details) and sends the data to a third party
  without causing visible damage?
  \begin{enumerate}[label=\Alph*.]
    \item Adware
    \item Spyware
    \item Logic bomb
    \item Bot
  \end{enumerate}
\item \textit{[Original]} A computer is flooded with pop-ups, banners, and
  redirected searches. This best indicates:
  \begin{enumerate}[label=\Alph*.]
    \item spyware
    \item adware
    \item rootkit
    \item backdoor
  \end{enumerate}
\item \textit{[Original]} A tool records every key pressed on a machine and
  forwards the log to an attacker, who harvests passwords and card numbers.
  This tool is a:
  \begin{enumerate}[label=\Alph*.]
    \item keylogger
    \item firewall
    \item logic bomb
    \item bot
  \end{enumerate}
\item \textit{[Original]} Which pairing correctly distinguishes spyware from
  adware?
  \begin{enumerate}[label=\Alph*.]
    \item Spyware shows pop-up ads; adware hides and steals passwords.
    \item Spyware hides and steals data quietly; adware openly floods the
      machine with ads.
    \item Spyware demands payment to unlock files; adware records
      keystrokes.
    \item Spyware spreads by itself over networks; adware needs a host file.
  \end{enumerate}
\item \textit{[Original]} Malware that hides itself and other malware deep
  inside the system, often with administrator-level control, is a:
  \begin{enumerate}[label=\Alph*.]
    \item rootkit
    \item worm
    \item keylogger
    \item bot
  \end{enumerate}
\item \textit{[Original]} A secret entry path left into a system, bypassing
  normal login and checks so an attacker can return later, is called a:
  \begin{enumerate}[label=\Alph*.]
    \item backdoor
    \item firewall
    \item backup
    \item patch
  \end{enumerate}
\item \textit{[Original]} The word ``botnet'' joins ``robot'' and
  ``network''. A botnet is a:
  \begin{enumerate}[label=\Alph*.]
    \item network of compromised machines controlled remotely by one master
    \item single anti-virus program shared over a network
    \item backup network used to restore encrypted files
    \item firewall setting that blocks advertisements
  \end{enumerate}
\item \textit{[Original]} Malicious code already planted on a machine waits
  for a fixed date and then deletes files. This is a:
  \begin{enumerate}[label=\Alph*.]
    \item logic bomb
    \item worm
    \item Trojan
    \item keylogger
  \end{enumerate}
\item \textit{[Original]} A computer suddenly runs hot and slow because an
  attacker is secretly mining cryptocurrency on it. This is:
  \begin{enumerate}[label=\Alph*.]
    \item cryptojacking
    \item adware
    \item logic bomb
    \item backdoor
  \end{enumerate}
\item \textit{[Original]} Evaluate the statements:
  \begin{enumerate}[label=\Roman*.]
    \item A virus needs a host file and human action to spread.
    \item A worm spreads by itself over networks without a host file.
    \item A Trojan replicates itself automatically across networks.
  \end{enumerate}
  \begin{enumerate}[label=\Alph*.]
    \item I only
    \item I and II only
    \item II and III only
    \item I, II, and III
  \end{enumerate}
\item \textit{[Original]} Evaluate the statements:
  \begin{enumerate}[label=\Roman*.]
    \item Ransomware encrypts files and demands payment for the unlock key.
    \item A rootkit hides the intruder deep in the system.
    \item A firewall removes malware already present inside a machine.
  \end{enumerate}
  \begin{enumerate}[label=\Alph*.]
    \item I only
    \item I and II only
    \item II and III only
    \item I, II, and III
  \end{enumerate}
\item \textit{[Original]} Which defence best defeats a ransomware demand,
  since data can be restored without paying?
  \begin{enumerate}[label=\Alph*.]
    \item Regular offline backups
    \item Opening every attachment quickly
    \item Disabling the firewall
    \item Downloading from untrusted sites
  \end{enumerate}
\item \textit{[Original]} Which pairing is correctly matched?
  \begin{enumerate}[label=\Alph*.]
    \item Worm --- needs a host file and a user click to spread
    \item Adware --- secretly records every keystroke
    \item Cryptojacking --- secretly mines cryptocurrency on the victim's
      machine
    \item Trojan --- self-replicating network traveller
  \end{enumerate}

\item \textit{[JKSSB FAA-pattern]} A computer virus needs which of the following to spread?
  \begin{enumerate}[label=\Alph*.]
    \item A host program or file to attach to
    \item Fresh printer ink
    \item A louder speaker
    \item A bigger monitor
  \end{enumerate}
\item \textit{[JKPSI-pattern]} Which malware demands payment to unlock files?
  \begin{enumerate}[label=\Alph*.]
    \item Ransomware
    \item Calculator
    \item Notepad
    \item Paint brush
  \end{enumerate}
\item \textit{[SSC CGL-pattern]} Antivirus software protects a computer against:
  \begin{enumerate}[label=\Alph*.]
    \item malicious programs.
    \item power cuts.
    \item keyboard dust only.
    \item printer paper jams.
  \end{enumerate}
\item \textit{[JKSSB-pattern]} \textbf{(MSQ: select ALL correct options.)} Which of the following are malware?
  \begin{enumerate}[label=\Alph*.]
    \item Virus
    \item Worm
    \item Trojan horse
    \item Antivirus scanner
  \end{enumerate}
\item \textit{[SSC CGL-pattern]} \textbf{(MSQ: select ALL correct options.)} Which statements about malware are correct?
  \begin{enumerate}[label=\Alph*.]
    \item Worms spread without a host program.
    \item Trojans disguise themselves as useful software.
    \item Ransomware encrypts files and demands payment.
    \item Antivirus updates increase infection risk.
  \end{enumerate}
\item \textit{[PYQ-adapted: JKSSB FAA 2022, QCRIAS item 556; SRC-21]} Which virus type attaches itself to executable files and activates when an infected file runs?
  \begin{enumerate}[label=\Alph*.]
\item Direct-action file virus
   \item Boot-sector virus
   \item Macro virus
   \item Overwrite virus only
  \end{enumerate}

\item \textit{[PYQ-adapted: JKSSB FAA 2024, QCRIAS item 664; SRC-16/SRC-19]} Which option is an antivirus product?
  \begin{enumerate}[label=\Alph*.]
\item Trojan horse
   \item Worm
   \item Macro virus
   \item Norton
  \end{enumerate}

\item \textit{[Original]} Malware that encrypts a victim's files and demands payment is commonly called:
  \begin{enumerate}[label=\Alph*.]
\item Ransomware
   \item Adware
   \item Cookie
   \item Device driver
  \end{enumerate}

\item \textit{[Original]} Which malware is commonly disguised as legitimate software and relies on a user running it?
  \begin{enumerate}[label=\Alph*.]
\item Trojan horse
   \item Worm
   \item Firewall
   \item Patch
  \end{enumerate}

\item \textit{[Original, MSQ drill]} \textbf{(MSQ: select ALL correct options.)} Which are malware categories?
  \begin{enumerate}[label=\Alph*.]
\item Virus
   \item Worm
   \item Trojan horse
   \item Firewall
  \end{enumerate}

\end{enumerate}
\lessonheading{43}{Hacking, Phishing and Spoofing}
\begin{enumerate}[label=\textbf{Q\arabic*.}, leftmargin=*]
\item \textit{[Original]} Expand OTP in the context of account security.
  \begin{enumerate}[label=\Alph*.]
    \item One-Time Password
    \item Online Transfer Protocol
    \item Official Transaction Permit
    \item Open Terminal Process
  \end{enumerate}
\item \textit{[Original]} Which term means gaining unauthorised access to
  a computer, account, or network?
  \begin{enumerate}[label=\Alph*.]
    \item Hacking
    \item Backup
    \item Compiling
    \item Updating
  \end{enumerate}
\item \textit{[Original]} What separates an ethical security test from
  hacking?
  \begin{enumerate}[label=\Alph*.]
    \item The test uses a faster computer
    \item The test is done with the owner's written permission
    \item The test is done at night
    \item The test steals no money
  \end{enumerate}
\item \textit{[Original]} A fake bank mail that imitates the bank and asks
  the reader to log in through the enclosed link is an example of:
  \begin{enumerate}[label=\Alph*.]
    \item phishing
    \item backup
    \item caching
    \item compiling
  \end{enumerate}
\item \textit{[Original]} Forging the ``From'' address so a mail displays
  a trusted name that never sent it is called:
  \begin{enumerate}[label=\Alph*.]
    \item email spoofing
    \item mail merge
    \item archiving
    \item forwarding
  \end{enumerate}
\item \textit{[Original]} Which set correctly matches each attack term to
  its meaning?
  \begin{enumerate}[label=\Alph*.]
    \item Hacking --- forged identity; Phishing --- unauthorised access; Spoofing --- fraud by imitation
    \item Hacking --- unauthorised access; Phishing --- fraud by imitation asking the victim to act; Spoofing --- forged identity
    \item Hacking --- fraud by imitation; Phishing --- forged identity; Spoofing --- unauthorised access
    \item Hacking --- backup copy; Phishing --- system update; Spoofing --- password reset
  \end{enumerate}
\item \textit{[Original]} A dictionary attack differs from a brute-force
  attack in that the dictionary attack:
  \begin{enumerate}[label=\Alph*.]
    \item tries every possible character combination
    \item watches the keyboard while the victim types
    \item tries each word in a ready list of common passwords
    \item forges the sender's address
  \end{enumerate}
\item \textit{[Original]} Trying every possible character combination until
  a password is found is called a:
  \begin{enumerate}[label=\Alph*.]
    \item brute-force attack
    \item mail merge
    \item backup operation
    \item cache refresh
  \end{enumerate}
\item \textit{[Original]} Watching the keyboard or screen while the victim
  types a password is called:
  \begin{enumerate}[label=\Alph*.]
    \item shoulder surfing
    \item phishing
    \item spoofing
    \item caching
  \end{enumerate}
\item \textit{[Original]} Which of the following is \textbf{NOT} a sign of
  a phishing message?
  \begin{enumerate}[label=\Alph*.]
    \item An urgent tone demanding immediate action
    \item A link whose real address differs from the genuine website
    \item A sender address that does not match the claimed organisation
    \item A message received by typing the bank's known address manually
  \end{enumerate}
\item \textit{[Original]} Tricking a person, rather than breaking software,
  into revealing credentials is called:
  \begin{enumerate}[label=\Alph*.]
    \item social engineering
    \item defragmentation
    \item compilation
    \item formatting
  \end{enumerate}
\item \textit{[Original]} Inventing a false identity and story, such as a
  bank officer verifying records, to extract an OTP is an example of:
  \begin{enumerate}[label=\Alph*.]
    \item pretexting
    \item caching
    \item indexing
    \item backing up
  \end{enumerate}
\item \textit{[Original]} Offering a free prize or free software to trigger
  a click or download is called:
  \begin{enumerate}[label=\Alph*.]
    \item baiting
    \item compiling
    \item sorting
    \item archiving
  \end{enumerate}
\item \textit{[Original]} Fraudulent text messages (SMS) that imitate a
  trusted sender are called \underline{\hspace{1.6cm}}.
  Fraudulent voice calls of the same kind are called
  \underline{\hspace{1.6cm}}.
  \begin{enumerate}[label=\Alph*.]
    \item smishing; vishing
    \item caching; compiling
    \item backing up; indexing
    \item surfing; browsing
  \end{enumerate}
\item \textit{[Original]} A caller demands the OTP to ``unblock'' an
  account. What is the correct response?
  \begin{enumerate}[label=\Alph*.]
    \item Share the OTP at once because the account is urgent
    \item Refuse, disconnect, and verify through the published helpline
    \item Read out all card numbers to prove identity
    \item Call back the number given by the caller immediately
  \end{enumerate}
\item \textit{[Original]} Before clicking a link in an unexpected message,
  the safest check is to:
  \begin{enumerate}[label=\Alph*.]
    \item hover over or long-press the link to read the real address, or type the known address manually
    \item click quickly before the offer expires
    \item forward the message to friends for advice
    \item reply asking whether the link is genuine
  \end{enumerate}
\item \textit{[Original]} An unexpected attachment arrives from a known
  name. What should the user do first?
  \begin{enumerate}[label=\Alph*.]
    \item Open it at once and enable macros if asked
    \item Confirm with the sender through a separate channel before opening
    \item Forward it to all contacts
    \item Delete the antivirus so the file opens faster
  \end{enumerate}
\item \textit{[Original]} In the attack chain, a weakness is called a
  \underline{\hspace{2.2cm}} and the method using it is called an
  \underline{\hspace{2.2cm}}. The harm delivered is called the
  \underline{\hspace{2.2cm}}.
  \begin{enumerate}[label=\Alph*.]
    \item vulnerability; exploit; payload
    \item payload; vulnerability; exploit
    \item exploit; payload; vulnerability
    \item attack; backup; cache
  \end{enumerate}
\item \textit{[Original]} Evaluate the statements:
  \begin{enumerate}[label=\Roman*.]
    \item Phishing imitates a trusted sender to steal credentials or money.
    \item Spoofing forges a sender address, caller number, or website address.
    \item A genuine bank routinely asks customers for their OTP by phone call.
  \end{enumerate}
  \begin{enumerate}[label=\Alph*.]
    \item I only
    \item I and II only
    \item II and III only
    \item I, II, and III
  \end{enumerate}
\item \textit{[Original]} Which pairing is correctly matched?
  \begin{enumerate}[label=\Alph*.]
    \item Vulnerability --- the method that uses a weakness
    \item Exploit --- a weakness that can be misused
    \item Payload --- the harmful element delivered by the attack
    \item Phishing --- a forged identity with no fraudulent message
  \end{enumerate}

\item \textit{[JKSSB FAA-pattern]} Phishing attacks mainly aim to steal:
  \begin{enumerate}[label=\Alph*.]
    \item passwords and personal details.
    \item printer ink cartridges.
    \item monitor stands.
    \item keyboard dust covers.
  \end{enumerate}
\item \textit{[JKPSI-pattern]} Spoofing means:
  \begin{enumerate}[label=\Alph*.]
    \item faking an identity such as an email sender.
    \item printing extra copies.
    \item increasing speaker bass.
    \item polishing the screen.
  \end{enumerate}
\item \textit{[SSC CGL-pattern]} Which sign suggests a phishing email?
  \begin{enumerate}[label=\Alph*.]
    \item Urgent requests for passwords via unknown links
    \item Correct official sender address you verified
    \item A letter from a known neighbour by post
    \item Your own handwritten diary note
  \end{enumerate}
\item \textit{[JKSSB-pattern]} \textbf{(MSQ: select ALL correct options.)} Which are signs of a phishing attempt?
  \begin{enumerate}[label=\Alph*.]
    \item Unknown sender demanding login details
    \item Suspicious links or attachments
    \item Threats of account closure within hours
    \item A stamped paper electricity bill
  \end{enumerate}
\item \textit{[SSC CGL-pattern]} \textbf{(MSQ: select ALL correct options.)} Which statements about hacking and phishing are correct?
  \begin{enumerate}[label=\Alph*.]
    \item Phishing uses fake messages to steal credentials.
    \item Spoofing hides the attacker's true identity.
    \item Strong unique passwords reduce account theft.
    \item Sharing OTPs with callers is safe practice.
  \end{enumerate}
\item \textit{[Original]} A fraudulent message that impersonates a trusted organization to steal credentials is an example of:
  \begin{enumerate}[label=\Alph*.]
\item Phishing
   \item Defragmentation
   \item Encryption
   \item Load balancing
  \end{enumerate}

\item \textit{[Original]} Forging the apparent sender address of an email is called:
  \begin{enumerate}[label=\Alph*.]
\item Spoofing
   \item Patching
   \item Hashing
   \item Sandboxing
  \end{enumerate}

\item \textit{[Original]} An attack that tricks a person into revealing confidential information primarily exploits:
  \begin{enumerate}[label=\Alph*.]
\item Social engineering
   \item Disk formatting
   \item Data compression
   \item Routing
  \end{enumerate}

\item \textit{[Original]} A software weakness that could be exploited by an attacker is a:
  \begin{enumerate}[label=\Alph*.]
\item Vulnerability
   \item Patch
   \item Backup
   \item Checksum
  \end{enumerate}

\item \textit{[Original, MSQ drill]} \textbf{(MSQ: select ALL correct options.)} Which are common signs of phishing?
  \begin{enumerate}[label=\Alph*.]
\item A mismatched sender/domain
   \item An urgent request to disclose a password
   \item A link whose destination differs from its displayed text
   \item A routine operating-system update from the configured updater
  \end{enumerate}

\end{enumerate}
\lessonheading{44}{Cybersecurity Prevention}
\begin{enumerate}[label=\textbf{Q\arabic*.}, leftmargin=*]
\item \textit{[Original]} Expand MFA.
  \begin{enumerate}[label=\Alph*.]
    \item Multi-Factor Authentication
    \item Maximum File Access
    \item Managed Firewall Authority
    \item Manual Forensic Audit
  \end{enumerate}
\item \textit{[Original]} In the CIA triad, confidentiality means:
  \begin{enumerate}[label=\Alph*.]
    \item only authorised persons can read the data.
    \item the data is accurate and unaltered.
    \item the systems are reachable when needed.
    \item every file has a backup copy.
  \end{enumerate}
\item \textit{[Original]} In the CIA triad, integrity means:
  \begin{enumerate}[label=\Alph*.]
    \item the data is accurate and has not been changed by unauthorised means.
    \item no password is ever required.
    \item the network is always switched on.
    \item the antivirus database is up to date.
  \end{enumerate}
\item \textit{[Original]} In the CIA triad, availability means:
  \begin{enumerate}[label=\Alph*.]
    \item the data and systems are reachable when needed.
    \item every user shares one password.
    \item all files are encrypted.
    \item no backup is required.
  \end{enumerate}
\item \textit{[Original]} A virus signature is:
  \begin{enumerate}[label=\Alph*.]
    \item a stored pattern of a known virus used for recognition.
    \item a handwritten sign on a printed file.
    \item a password written on a USB drive.
    \item a firewall rule for blocking websites.
  \end{enumerate}
\item \textit{[Original]} Real-time scanning checks files:
  \begin{enumerate}[label=\Alph*.]
    \item only once a year during audit.
    \item as they are opened, downloaded, or copied.
    \item after the computer is switched off.
    \item only when the printer is used.
  \end{enumerate}
\item \textit{[Original]} Quarantine differs from deletion in that quarantine:
  \begin{enumerate}[label=\Alph*.]
    \item isolates a suspicious file so it cannot run, allowing later restore or delete.
    \item permanently removes the file at once.
    \item sends the file to all contacts.
    \item installs the file automatically.
  \end{enumerate}
\item \textit{[Original]} The main job of a firewall is to:
  \begin{enumerate}[label=\Alph*.]
    \item remove viruses already saved on the disk.
    \item monitor network traffic and block what its rules forbid.
    \item restore files lost to ransomware.
    \item write strong passwords for users.
  \end{enumerate}
\item \textit{[Original]} Which of the following is \textbf{NOT} a good
  preventive habit?
  \begin{enumerate}[label=\Alph*.]
    \item Installing software updates promptly.
    \item Scanning a stranger's USB drive before opening its files.
    \item Clicking every link in unexpected prize-winning messages.
    \item Using a different strong password for each account.
  \end{enumerate}
\item \textit{[Original]} A patch is best described as:
  \begin{enumerate}[label=\Alph*.]
    \item a small vendor fix for a known security weakness.
    \item a new printer cable.
    \item a type of computer virus.
    \item a backup copy of a file.
  \end{enumerate}
\item \textit{[Original]} Which password is strongest?
  \begin{enumerate}[label=\Alph*.]
    \item rahul123
    \item password
    \item 12345678
    \item RaHuL\#2026!Desk
  \end{enumerate}
\item \textit{[Original]} MFA protects an account mainly because:
  \begin{enumerate}[label=\Alph*.]
    \item it demands a second independent proof even if the password leaks.
    \item it makes the password shorter.
    \item it removes the need for any password.
    \item it deletes all viruses on the disk.
  \end{enumerate}
\item \textit{[Original]} A backup protects which leg of the CIA triad most
  directly?
  \begin{enumerate}[label=\Alph*.]
    \item Confidentiality
    \item Integrity
    \item Availability
    \item Authentication
  \end{enumerate}
\item \textit{[Original]} Good backup practice requires that copies be:
  \begin{enumerate}[label=\Alph*.]
    \item regular and kept offline or off-site, with restores tested.
    \item stored only on the same disk as the original.
    \item made once and never checked again.
    \item shared with strangers for safety.
  \end{enumerate}
\item \textit{[Original]} Before entering bank details on a website, the user
  should check for:
  \begin{enumerate}[label=\Alph*.]
    \item HTTPS and the padlock symbol.
    \item blinking advertisements.
    \item a prize-winning pop-up.
    \item an unknown shortened link.
  \end{enumerate}
\item \textit{[Original]} Home and office Wi-Fi should use:
  \begin{enumerate}[label=\Alph*.]
    \item open access with no password.
    \item WPA2 or WPA3 encryption with a strong password.
    \item the router's default admin password forever.
    \item automatic connection to every hotspot.
  \end{enumerate}
\item \textit{[Original]} The safest way to handle an unknown USB drive is to:
  \begin{enumerate}[label=\Alph*.]
    \item let it auto-run and open every file at once.
    \item scan it with antivirus before opening files, with autorun disabled.
    \item copy its programs to the office desktop unseen.
    \item lend it to every colleague first.
  \end{enumerate}
\item \textit{[Original]} Evaluate the statements:
  \begin{enumerate}[label=\Roman*.]
    \item A firewall removes viruses already saved on the disk.
    \item Antivirus software detects and removes malicious files.
    \item An offline backup lets work resume after ransomware locks files.
  \end{enumerate}
  \begin{enumerate}[label=\Alph*.]
    \item I only
    \item II and III only
    \item I and II only
    \item I, II, and III
  \end{enumerate}
\item \textit{[Original]} Which pairing is correctly matched?
  \begin{enumerate}[label=\Alph*.]
    \item Firewall --- removes infected files from the disk
    \item Quarantine --- permanently deletes a file at once
    \item MFA --- a second proof of identity beyond the password
    \item Backup --- blocks unauthorised network traffic
  \end{enumerate}
\item \textit{[Original]} An office PC is infected through a USB drive that
  launched a program by itself, and months of work are locked by ransomware.
  Which two controls, together, would have prevented the worst damage?
  \begin{enumerate}[label=\Alph*.]
    \item Disabling autorun with USB scanning, plus regular offline backups.
    \item A longer screensaver delay plus a bigger monitor.
    \item Deleting the antivirus and switching the firewall off.
    \item Sharing one password among all staff plus open public Wi-Fi.
  \end{enumerate}

\item \textit{[JKSSB FAA-pattern]} A firewall is used to:
  \begin{enumerate}[label=\Alph*.]
    \item monitor and control network traffic.
    \item print fire-safety posters.
    \item cool the monitor glass.
    \item charge the keyboard battery.
  \end{enumerate}
\item \textit{[JKPSI-pattern]} Two-factor authentication improves security by:
  \begin{enumerate}[label=\Alph*.]
    \item requiring two different proofs of identity.
    \item printing two copies of each file.
    \item using two monitors at once.
    \item typing every password twice slowly.
  \end{enumerate}
\item \textit{[SSC CGL-pattern]} Which password is strongest?
  \begin{enumerate}[label=\Alph*.]
    \item A long mix of letters, digits, and symbols
    \item Your own first name
    \item 123456
    \item The word password
  \end{enumerate}
\item \textit{[JKSSB-pattern]} \textbf{(MSQ: select ALL correct options.)} Which practices improve computer security?
  \begin{enumerate}[label=\Alph*.]
    \item Updating software regularly
    \item Using strong unique passwords
    \item Enabling a firewall
    \item Sharing passwords with friends
  \end{enumerate}
\item \textit{[SSC CGL-pattern]} \textbf{(MSQ: select ALL correct options.)} Which statements about prevention are correct?
  \begin{enumerate}[label=\Alph*.]
    \item Backups limit ransomware damage.
    \item Antivirus should be kept updated.
    \item Unknown attachments should not be opened.
    \item Public Wi-Fi needs no caution at all.
  \end{enumerate}
\item \textit{[Original]} Which security property means that information is accessible only to authorized parties?
  \begin{enumerate}[label=\Alph*.]
\item Confidentiality
   \item Availability
   \item Redundancy
   \item Compression
  \end{enumerate}

\item \textit{[Original]} A firewall primarily helps to:
  \begin{enumerate}[label=\Alph*.]
\item Control network traffic according to rules
   \item Create a backup copy of every document
   \item Repair broken monitor pixels
   \item Convert scanned pages to editable text
  \end{enumerate}

\item \textit{[Original]} Why should security updates be installed promptly?
  \begin{enumerate}[label=\Alph*.]
\item They can fix known vulnerabilities
   \item They guarantee that no attack can occur
   \item They replace all backups
   \item They make passwords unnecessary
  \end{enumerate}

\item \textit{[Original]} Which control is most directly used to restore files after accidental deletion or disk failure?
  \begin{enumerate}[label=\Alph*.]
\item Backup
   \item Firewall
   \item Screen saver
   \item Spam filter
  \end{enumerate}

\item \textit{[Original, MSQ drill]} \textbf{(MSQ: select ALL correct options.)} Which practices strengthen account security?
  \begin{enumerate}[label=\Alph*.]
\item Use multi-factor authentication.
   \item Use unique strong passwords.
   \item Install security updates.
   \item Reuse one password on every service.
  \end{enumerate}

\end{enumerate}
\lessonheading{56}{Cybersecurity, IT Act and Security Controls}
\begin{enumerate}[label=\textbf{Q\arabic*.}, leftmargin=*]
\item \textit{[Original]} Expand MFA.
  \begin{enumerate}[label=\Alph*.]
    \item Multi-Factor Authentication
    \item Multiple-File Access
    \item Managed Firewall Authority
    \item Master File Archive
  \end{enumerate}
\item \textit{[Original]} SSL stands for:
  \begin{enumerate}[label=\Alph*.]
    \item Secure Sockets Layer
    \item System Security Lock
    \item Shared Signature Log
    \item Synchronised Server Link
  \end{enumerate}
\item \textit{[Original]} The Information Technology Act of India was
  enacted in the year:
  \begin{enumerate}[label=\Alph*.]
    \item 1998
    \item 2000
    \item 2008
    \item 2013
  \end{enumerate}
\item \textit{[Original]} Which of the following is \textbf{NOT} within the
  scope of the Information Technology Act, 2000 as covered in this lesson?
  \begin{enumerate}[label=\Alph*.]
    \item Legal validity of electronic records and electronic signatures
    \item Electronic governance
    \item Cyber contraventions and offences
    \item Regulation of printer hardware prices and typing-test fees
  \end{enumerate}
\item \textit{[Original]} Which term is the \textbf{umbrella} term for
  malicious software, under which virus, worm, Trojan, and ransomware fall?
  \begin{enumerate}[label=\Alph*.]
    \item Malware
    \item Firewall
    \item Patch
    \item Backup
  \end{enumerate}
\item \textit{[Original]} Which malware attaches itself to a host file and
  needs human action, such as opening an infected file, to spread?
  \begin{enumerate}[label=\Alph*.]
    \item Worm
    \item Trojan horse
    \item Virus
    \item Ransomware
  \end{enumerate}
\item \textit{[Original]} Which malware copies itself and spreads across a
  network without a host file and without waiting for the user to act?
  \begin{enumerate}[label=\Alph*.]
    \item Virus
    \item Worm
    \item Trojan horse
    \item Spyware-free utility
  \end{enumerate}
\item \textit{[Original]} Which statement is \textbf{NOT} true of a Trojan
  horse?
  \begin{enumerate}[label=\Alph*.]
    \item It pretends to be useful software while hiding a harmful purpose.
    \item It opens a backdoor for stealing passwords or controlling the machine.
    \item It replicates itself across networks without any user action.
    \item It relies on the user installing it under a false disguise.
  \end{enumerate}
\item \textit{[Original]} Ransomware is best described as malware that:
  \begin{enumerate}[label=\Alph*.]
    \item filters network traffic against a set of rules.
    \item encrypts or locks the victim's files and demands payment for release.
    \item isolates suspicious files for later inspection.
    \item fixes software flaws by installing vendor updates.
  \end{enumerate}
\item \textit{[Original]} Evaluate the statements:
  \begin{enumerate}[label=\Roman*.]
    \item A virus needs a host file and human action to spread.
    \item A worm is standalone and spreads across a network unaided.
    \item A Trojan horse replicates itself automatically like a worm.
  \end{enumerate}
  \begin{enumerate}[label=\Alph*.]
    \item I only
    \item I and II only
    \item II and III only
    \item I, II, and III
  \end{enumerate}
\item \textit{[Original]} Placing a suspicious file in quarantine means:
  \begin{enumerate}[label=\Alph*.]
    \item deleting it permanently beyond recovery.
    \item isolating it so it cannot run or spread, pending a decision.
    \item backing it up to a separate disk for safekeeping.
    \item encrypting it for secure transmission to another user.
  \end{enumerate}
\item \textit{[Original]} Phishing is:
  \begin{enumerate}[label=\Alph*.]
    \item a fake message impersonating a trusted sender to steal passwords or money.
    \item a barrier that filters incoming and outgoing network traffic.
    \item a vendor update that fixes known security flaws.
    \item a protocol that encrypts data during transmission.
  \end{enumerate}
\item \textit{[Original]} A customer receives a mail reading ``your account
  is blocked; click here and enter your OTP to restore it.'' The correct
  anti-phishing response is to:
  \begin{enumerate}[label=\Alph*.]
    \item click at once and enter the OTP before the deadline.
    \item forward the OTP to the sender for verification.
    \item ignore the link, share no OTP, and verify through the known official channel.
    \item download the attached ``security tool'' and run it.
  \end{enumerate}
\item \textit{[Original]} Encryption primarily provides:
  \begin{enumerate}[label=\Alph*.]
    \item confidentiality, by converting readable data so only authorised persons can read it.
    \item proof of the sender's identity without hiding the message.
    \item filtering of network packets at the boundary.
    \item isolation of suspicious files from the system.
  \end{enumerate}
\item \textit{[Original]} A digital signature primarily provides:
  \begin{enumerate}[label=\Alph*.]
    \item secrecy of the message content from eavesdroppers.
    \item removal of malware already present on the machine.
    \item proof of who sent the message and that it was not altered.
    \item an encrypted channel for data in transit.
  \end{enumerate}
\item \textit{[Original]} SSL/TLS protect:
  \begin{enumerate}[label=\Alph*.]
    \item files already infected on the disk, by cleaning the malware.
    \item data in transit, by providing a secure encrypted channel.
    \item login accounts, by demanding an OTP at every step.
    \item software flaws, by installing vendor-supplied updates.
  \end{enumerate}
\item \textit{[Original]} Which control filters incoming and outgoing
  network traffic against a set of rules and blocks unauthorised access?
  \begin{enumerate}[label=\Alph*.]
    \item Firewall
    \item Norton
    \item McAfee
    \item Quick Heal
  \end{enumerate}
\item \textit{[Original]} Patching means:
  \begin{enumerate}[label=\Alph*.]
    \item demanding payment to release locked files.
    \item installing vendor-supplied updates that fix known security flaws.
    \item adding a second login check such as an OTP.
    \item copying data to a separate store for recovery after loss.
  \end{enumerate}
\item \textit{[Original]} An attacker obtains a user's password, but the
  account also requires an OTP sent to the user's phone, which the attacker
  does not have. The login fails. This scenario shows the value of:
  \begin{enumerate}[label=\Alph*.]
    \item quarantine
    \item ransomware
    \item Multi-Factor Authentication (MFA)
    \item phishing
  \end{enumerate}
\item \textit{[Original]} Which pairing is correctly matched?
  \begin{enumerate}[label=\Alph*.]
    \item Quarantine --- deletes the file permanently beyond recovery
    \item Kaspersky --- a network barrier that filters traffic by rules
    \item Encryption --- proves the sender without hiding the message
    \item Patching --- fixes security flaws by installing updates
  \end{enumerate}

\item \textit{[JKSSB FAA-pattern]} The Information Technology Act of India was enacted in the year:
  \begin{enumerate}[label=\Alph*.]
    \item 2000
    \item 1990
    \item 2010
    \item 2020
  \end{enumerate}
\item \textit{[JKPSI-pattern]} Which of the following is punishable as hacking under the IT Act?
  \begin{enumerate}[label=\Alph*.]
    \item Unauthorised access causing damage to computer systems
    \item Printing your own document twice
    \item Increasing monitor brightness
    \item Charging a mobile phone
  \end{enumerate}
\item \textit{[SSC CGL-pattern]} A digital signature is used to verify the:
  \begin{enumerate}[label=\Alph*.]
    \item authenticity and integrity of a message.
    \item weight of a parcel.
    \item colour of a printer.
    \item volume of a speaker.
  \end{enumerate}
\item \textit{[JKSSB-pattern]} \textbf{(MSQ: select ALL correct options.)} Which are security controls for computer systems?
  \begin{enumerate}[label=\Alph*.]
    \item Strong passwords
    \item Firewalls
    \item Regular software updates
    \item Sharing OTPs with callers
  \end{enumerate}
\item \textit{[SSC CGL-pattern]} \textbf{(MSQ: select ALL correct options.)} Which statements about the IT Act and security are correct?
  \begin{enumerate}[label=\Alph*.]
    \item The IT Act addresses cybercrime in India.
    \item Digital signatures get legal recognition under the Act.
    \item Hacking with damage is punishable under the Act.
    \item The IT Act bans the use of all computers.
  \end{enumerate}
\item \textit{[PYQ-adapted: JKSSB FAA 2022, QCRIAS item 556; SRC-21]} Which virus type attaches itself to executable files such as .EXE or .COM files?
  \begin{enumerate}[label=\Alph*.]
\item Direct-action file virus
   \item Boot-sector virus
   \item Macro virus
   \item Overwrite virus only
  \end{enumerate}

\item \textit{[PYQ-adapted: JKSSB FAA 2024, QCRIAS item 664; SRC-16/SRC-19]} Which option is antivirus software?
  \begin{enumerate}[label=\Alph*.]
\item Trojan horse
   \item Worm
   \item Macro virus
   \item Norton
  \end{enumerate}

\item \textit{[Original]} Which security control filters network traffic according to rules?
  \begin{enumerate}[label=\Alph*.]
\item Firewall
   \item Spreadsheet
   \item Defragmenter
   \item OCR
  \end{enumerate}

\item \textit{[Original]} A file that appears harmless but contains a hidden malicious purpose is a:
  \begin{enumerate}[label=\Alph*.]
\item Trojan horse
   \item Patch
   \item Backup
   \item Firewall
  \end{enumerate}

\item \textit{[Original, MSQ drill]} \textbf{(MSQ: select ALL correct options.)} Which are security controls?
  \begin{enumerate}[label=\Alph*.]
\item Antivirus software
   \item Firewall
   \item Multi-factor authentication
   \item Phishing email
  \end{enumerate}

\end{enumerate}
\lessonheading{64}{Current Cybersecurity, Cryptography and Cyber Law}
\begin{enumerate}[label=\textbf{Q\arabic*.}, leftmargin=*]
\item \textit{[Original]} A parasitic (file-infector) virus is best described as:
  \begin{enumerate}[label=\Alph*.]
    \item a virus that attaches itself to an executable file and runs when the host program runs
    \item a worm that spreads across a network without any host program
    \item a Trojan horse that disguises itself as useful software without self-replicating
    \item a flood of traffic from many machines that overwhelms a server
  \end{enumerate}
\item \textit{[Original]} Expand DES.
  \begin{enumerate}[label=\Alph*.]
    \item Data Encryption Standard
    \item Digital Encryption System
    \item Distributed Encoding Scheme
    \item Data Evaluation Service
  \end{enumerate}
\item \textit{[Original]} DES processes data in blocks of:
  \begin{enumerate}[label=\Alph*.]
    \item 32 bits
    \item 56 bits
    \item 64 bits
    \item 128 bits
  \end{enumerate}
\item \textit{[Original]} Which of the following is \textbf{NOT} true of a digital signature?
  \begin{enumerate}[label=\Alph*.]
    \item It proves the sender's authenticity.
    \item It proves the message was not altered in transit.
    \item It provides integrity of the message.
    \item It encrypts the message body so only the recipient can read it.
  \end{enumerate}
\item \textit{[Original]} Encryption primarily provides \underline{\hspace{2.2cm}}, whereas a digital signature primarily provides \underline{\hspace{2.2cm}}.
  \begin{enumerate}[label=\Alph*.]
    \item confidentiality; authenticity and integrity
    \item authenticity; confidentiality
    \item availability; confidentiality
    \item integrity; availability
  \end{enumerate}
\item \textit{[Original]} Which firewall type examines only per-packet headers such as source/destination address, port, and protocol?
  \begin{enumerate}[label=\Alph*.]
    \item Packet-filtering firewall
    \item Stateful-inspection firewall
    \item Application (proxy) firewall
    \item Personal anti-virus program
  \end{enumerate}
\item \textit{[Original]} A firewall that remembers open sessions and drops packets that do not belong to a known session is a:
  \begin{enumerate}[label=\Alph*.]
    \item packet-filtering firewall
    \item stateful-inspection firewall
    \item boot-sector virus
    \item Trojan horse
  \end{enumerate}
\item \textit{[Original]} Which firewall acts as a go-between, inspecting application-layer content and commands?
  \begin{enumerate}[label=\Alph*.]
    \item Packet-filtering firewall
    \item Stateful-inspection firewall
    \item Application (proxy) firewall
    \item Spam filter of the mail client only
  \end{enumerate}
\item \textit{[Original]} A fraudulent mail styled as a bank KYC notice asks a clerk for her OTP on a fake login page. This is:
  \begin{enumerate}[label=\Alph*.]
    \item phishing
    \item ransomware
    \item DDoS
    \item disk defragmentation
  \end{enumerate}
\item \textit{[Original]} A clerk's files are suddenly unreadable, with a message demanding payment for the decryption key. This is:
  \begin{enumerate}[label=\Alph*.]
    \item phishing
    \item ransomware
    \item DDoS
    \item hacking without any malicious goal
  \end{enumerate}
\item \textit{[Original]} Which of the following correctly distinguishes the four threats?
  \begin{enumerate}[label=\Alph*.]
    \item Phishing floods a server; ransomware tricks the user; DDoS encrypts files; hacking defragments disks.
    \item Phishing tricks the user into revealing credentials; ransomware locks files for ransom; DDoS floods a service; hacking is unauthorised access generally.
    \item Phishing encrypts files; ransomware floods a server; DDoS reveals OTPs; hacking is a firewall type.
    \item Phishing is a firewall; ransomware is an anti-virus; DDoS is encryption; hacking is a backup.
  \end{enumerate}
\item \textit{[Original]} Expand DDoS.
  \begin{enumerate}[label=\Alph*.]
    \item Distributed Denial of Service
    \item Direct Disk Operating System
    \item Digital Data on Server
    \item Dynamic Domain of Service
  \end{enumerate}
\item \textit{[Original]} Evaluate the statements:
  \begin{enumerate}[label=\Roman*.]
    \item DES is a symmetric cipher: the same key encrypts and decrypts.
    \item DES works on 64-bit blocks of data.
    \item A digital signature by itself keeps the message body secret from everyone except the recipient.
  \end{enumerate}
  \begin{enumerate}[label=\Alph*.]
    \item I only
    \item I and II only
    \item II and III only
    \item I, II, and III
  \end{enumerate}
\item \textit{[Original]} Which pairing of firewall type and function is correctly matched?
  \begin{enumerate}[label=\Alph*.]
    \item Packet-filtering --- inspects application content as a go-between
    \item Stateful inspection --- tracks connection state in addition to headers
    \item Application (proxy) --- decides each packet by headers only
    \item Packet-filtering --- disinfects infected executable files
  \end{enumerate}
\item \textit{[Original]} Which of the following is a governance concern of the Internet at awareness level?
  \begin{enumerate}[label=\Alph*.]
    \item Which country's law applies when an act crosses borders
    \item The colour scheme of a single office website
    \item The brand of printer attached to the office desktop
    \item The seating plan of the computer laboratory
  \end{enumerate}
\item \textit{[Original]} Which data-protection principle requires that personal data be used only for the purpose stated at collection?
  \begin{enumerate}[label=\Alph*.]
    \item Purpose limitation (stated-purpose use)
    \item Unlimited retention
    \item Open sharing without consent
    \item Default publication
  \end{enumerate}
\item \textit{[Original]} Copying licensed software without permission is, at awareness level, an example of:
  \begin{enumerate}[label=\Alph*.]
    \item copyright infringement
    \item fair dealing for private study
    \item first ownership
    \item lawful backup under licence
  \end{enumerate}
\item \textit{[Original]} Match the copyright concept to what it covers: (a) meaning of copyright, (b) first ownership, (c) infringement, (d) fair dealing. Which match is correct?
  \begin{enumerate}[label=\Alph*.]
    \item (a)---punishment for wilful violation; (b)---permitted study uses; (c)---who owns first; (d)---owner's exclusive rights
    \item (a)---owner's exclusive rights; (b)---who owns the work first; (c)---what counts as violation; (d)---limited permitted uses such as study
    \item (a)---who owns first; (b)---owner's rights; (c)---permitted uses; (d)---what counts as violation
    \item (a)---permitted uses; (b)---violation; (c)---owner's rights; (d)---first ownership
  \end{enumerate}
\item \textit{[Original]} Expand ICANN.
  \begin{enumerate}[label=\Alph*.]
    \item Internet Corporation for Assigned Names and Numbers
    \item International Computer and Network Node
    \item Internet Code for Access and Navigation
    \item Indian Council for Assigned Network Names
  \end{enumerate}
\item \textit{[Original]} An office server slows to a halt under a flood of traffic from many machines, so genuine users cannot reach it. The correct response pair is:
  \begin{enumerate}[label=\Alph*.]
    \item identify it as DDoS and filter/block the flood at the network barrier (firewall/ISP level)
    \item identify it as phishing and ask users to change their OTPs only
    \item identify it as a file-infector virus and reinstall the printer driver
    \item identify it as fair dealing and take no technical action
  \end{enumerate}

\item \textit{[JKSSB FAA-pattern]} Encryption is used to provide:
  \begin{enumerate}[label=\Alph*.]
    \item confidentiality of data.
    \item louder speaker output.
    \item brighter monitor light.
    \item faster paper printing.
  \end{enumerate}
\item \textit{[JKPSI-pattern]} A digital signature provides which guarantees?
  \begin{enumerate}[label=\Alph*.]
    \item Authenticity and integrity
    \item Free printer ink
    \item Louder ringtones
    \item Heavier paper
  \end{enumerate}
\item \textit{[SSC CGL-pattern]} Which law deals with cybercrime in India?
  \begin{enumerate}[label=\Alph*.]
    \item The Information Technology Act
    \item The Paper Weight Act
    \item The Speaker Volume Act
    \item The Ink Colour Act
  \end{enumerate}
\item \textit{[JKSSB-pattern]} \textbf{(MSQ: select ALL correct options.)} Which are symmetric-key concepts or examples?
  \begin{enumerate}[label=\Alph*.]
    \item Same key encrypts and decrypts
    \item AES cipher
    \item Fast bulk-data encryption
    \item Public-private key pairs
  \end{enumerate}
\item \textit{[SSC CGL-pattern]} \textbf{(MSQ: select ALL correct options.)} Which statements about cryptography are correct?
  \begin{enumerate}[label=\Alph*.]
    \item Encryption hides message meaning.
    \item Decryption restores the original message.
    \item Hashing checks data integrity.
    \item Encryption makes data print faster.
  \end{enumerate}
\item \textit{[PYQ-adapted: SSC CGL 2025 Tier-II Paper-I, 19-Jan-2026 Q5; SRC-18]} Which option provides message-level email encryption to protect content from unauthorized access?
  \begin{enumerate}[label=\Alph*.]
\item S/MIME
   \item IMAP
   \item SSL alone
   \item SMTP
  \end{enumerate}

\item \textit{[PYQ-adapted: SSC CGL 2025 Tier-II Paper-I, 19-Jan-2026 Q11; SRC-18]} What is a common purpose of a Trojan horse?
  \begin{enumerate}[label=\Alph*.]
\item Disguise malicious activity as a benign program
   \item Self-replicate as a worm
   \item Filter network packets as a firewall
   \item Repair operating-system vulnerabilities
  \end{enumerate}

\item \textit{[Original]} Which property does encryption primarily provide when only authorized parties can read the content?
  \begin{enumerate}[label=\Alph*.]
\item Confidentiality
   \item Availability
   \item Non-repudiation only
   \item Compression
  \end{enumerate}

\item \textit{[Original]} A digital signature primarily helps verify:
  \begin{enumerate}[label=\Alph*.]
\item Authenticity and integrity
   \item Secrecy of the message content by itself
   \item Internet connection speed
   \item File size reduction
  \end{enumerate}

\item \textit{[Original, MSQ drill]} \textbf{(MSQ: select ALL correct options.)} Match security concept to purpose.
  \begin{enumerate}[label=\Alph*.]
\item Encryption --- confidentiality
   \item Digital signature --- authenticity/integrity
   \item Firewall --- traffic control
   \item Antivirus --- guaranteed prevention of every attack
  \end{enumerate}

\end{enumerate}
\answersection{10}
\answersubheading{42}{Malware}
\begin{enumerate}[label=\textbf{Q\arabic*.}, leftmargin=*]
\item \textbf{B.} Malware is short for malicious software. A, C, and D are
  not expansions of the term.
\item \textbf{A.} A virus attaches to a host file and needs the user to open
  it. A worm (B) needs neither host nor click; a bot (C) is a remotely
  controlled machine; a firewall (D) is a protective barrier, not malware.
\item \textbf{C.} A worm is standalone and self-spreading over networks. A
  Trojan (A) relies on disguise and user installation; a virus (B) needs a
  host and a click; adware (D) floods the screen with ads.
\item \textbf{C.} Disguise as a useful program plus no self-replication is a
  Trojan. A worm (A) spreads alone; a virus (B) needs a host file; a logic
  bomb (D) waits for a trigger rather than tricking the user.
\item \textbf{B.} Virus = host + user action; worm = standalone and
  self-spreading; Trojan = disguised, no replication. A swaps virus and worm
  and mislabels the Trojan; C and D attach wrong behaviours (payment,
  keystrokes, hiding, mining, trigger) to these three.
\item \textbf{C.} A Trojan does not replicate itself; that is the false
  statement. Pretending to be legitimate (A), relying on user installation
  (B), and opening backdoors or installing more malware (D) are all true of
  Trojans.
\item \textbf{A.} Locked files plus a payment demand is ransomware. Adware
  (B) shows ads; a worm (C) self-spreads over networks; a rootkit (D) hides
  the intruder rather than demanding money.
\item \textbf{B.} Silent observation and theft of user data is spyware.
  Adware (A) shows visible ads; a logic bomb (C) waits for a trigger; a bot
  (D) is a remotely controlled machine.
\item \textbf{B.} Flooded pop-ups and redirected searches indicate adware.
  Spyware (A) hides; a rootkit (C) conceals deeply; a backdoor (D) is a
  secret re-entry path.
\item \textbf{A.} Recording every keystroke to steal credentials is a
  keylogger. A firewall (B) filters traffic; a logic bomb (C) waits for a
  trigger; a bot (D) is a controlled machine.
\item \textbf{B.} Spyware hides and steals quietly while adware openly shows
  ads. A reverses the pair; C describes ransomware and a keylogger; D
  describes a worm and a virus.
\item \textbf{A.} Deep concealment with high-level control is a rootkit. A
  worm (B) self-spreads; a keylogger (C) records keystrokes; a bot (D) is a
  controlled machine.
\item \textbf{A.} A secret entry bypassing normal checks is a backdoor. A
  firewall (B) filters traffic; a backup (C) is a stored copy for recovery;
  a patch (D) is a software fix.
\item \textbf{A.} A botnet is a network of compromised machines obeying one
  master, used for spam and large attacks. B, C, and D describe an
  anti-virus share, a backup network, and a firewall setting, none of which
  is a botnet.
\item \textbf{A.} Planted code that waits for a trigger (here a date) is a
  logic bomb. A worm (B) spreads actively; a Trojan (C) deceives the user;
  a keylogger (D) records keystrokes.
\item \textbf{A.} Secret mining on the victim's machine is cryptojacking.
  Adware (B) shows ads; a logic bomb (C) waits for a trigger; a backdoor
  (D) is a secret entry.
\item \textbf{B.} Statements I and II are correct: a virus needs a host and
  human action, and a worm self-spreads without a host. III is false because
  a Trojan does not replicate itself. A omits II; C includes false III and
  omits I; D includes false III.
\item \textbf{B.} Statements I and II are correct: ransomware encrypts and
  demands payment, and a rootkit hides the intruder. III is false because a
  firewall filters traffic and does not remove malware already inside. A
  omits II; C includes false III and omits I; D includes false III.
\item \textbf{A.} Regular offline backups let the victim restore data
  without paying. Opening unknown attachments (B), disabling the firewall
  (C), and untrusted downloads (D) invite infection rather than defeating
  it.
\item \textbf{C.} Cryptojacking is the secret mining of cryptocurrency on the
  victim's machine. A mislabels the worm (it needs neither host nor click);
  B mislabels adware (it shows ads; keystroke recording is a keylogger); D
  mislabels the Trojan (it does not self-replicate).

\item \textbf{A.} Viruses attach to hosts to spread. B, C, and D are hardware traits, not spread mechanisms.
\item \textbf{A.} Ransomware extorts payment for locked files. B, C, and D are harmless utilities.
\item \textbf{A.} Antivirus fights malicious code. B, C, and D are physical issues beyond its scope.
\item \textbf{A, B, C.} Viruses, worms, and Trojans are malware. The scanner (D) fights them instead.
\item \textbf{A, B, C.} The worm, Trojan, and ransomware facts are correct. D is false: updates reduce risk.
\item \textbf{A.} A direct-action file virus infects executable files such as .EXE/.COM files. Boot-sector viruses target boot records; macro viruses use document macros.

\item \textbf{D.} Norton is antivirus software. A Trojan, worm, and macro virus are malware types.

\item \textbf{A.} Ransomware encrypts or locks data and demands a ransom. Adware displays advertisements; a cookie is browser data; a driver is software for a device.

\item \textbf{A.} A Trojan appears benign but carries a hidden malicious purpose. A worm self-propagates; firewalls and patches are defensive controls.

\item \textbf{A, B, C.} Virus, worm, and Trojan horse are malware. A firewall is a security control.

\end{enumerate}
\answersubheading{43}{Hacking, Phishing and Spoofing}
\begin{enumerate}[label=\textbf{Q\arabic*.}, leftmargin=*]
\item \textbf{A.} OTP stands for One-Time Password. B, C, and D are not
  expansions of the abbreviation.
\item \textbf{A.} Hacking means gaining unauthorised access. Backup (B),
  compiling (C), and updating (D) are legitimate operations.
\item \textbf{B.} Written permission from the owner is what makes a
  security test ethical; without it the same act is hacking. Speed (A),
  timing (C), and whether money is taken (D) do not decide the matter.
\item \textbf{A.} A fake mail imitating a trusted sender with a call to
  action is phishing. Backup, caching, and compiling (B, C, D) are
  unrelated computer operations.
\item \textbf{A.} Forging the ``From'' address is email spoofing. Mail
  merge (B), archiving (C), and forwarding (D) are legitimate mail
  operations.
\item \textbf{B.} Hacking is unauthorised access, phishing is fraud by
  imitation asking the victim to act, and spoofing is the forged
  identity. A swaps all three meanings; C rotates them; D substitutes
  unrelated operations.
\item \textbf{C.} A dictionary attack tries each word in a ready list of
  common passwords. Trying every combination (A) is brute force;
  watching the keyboard (B) is shoulder surfing; forging an address (D)
  is spoofing.
\item \textbf{A.} Trying every possible combination is a brute-force
  attack. Mail merge (B), backup (C), and cache refresh (D) are
  unrelated.
\item \textbf{A.} Observing typing over the victim's shoulder is
  shoulder surfing. Phishing (B) is fraud by imitation, spoofing (C) is
  forged identity, and caching (D) is a memory technique.
\item \textbf{D.} Typing the bank's known address manually is the safe
  habit, not a phishing sign. Urgency (A), a mismatched link (B), and a
  mismatched sender (C) are all genuine warning signs.
\item \textbf{A.} Manipulating the person rather than the software is
  social engineering. Defragmentation (B), compilation (C), and
  formatting (D) are unrelated technical operations.
\item \textbf{A.} A invented authoritative story to extract details is
  pretexting. Caching (B), indexing (C), and backing up (D) are
  unrelated.
\item \textbf{A.} A free-prize lure that triggers a click or download
  is baiting. Compiling (B), sorting (C), and archiving (D) are
  unrelated.
\item \textbf{A.} SMS phishing is smishing and voice-call phishing is
  vishing. The other pairs (B, C, D) name unrelated operations.
\item \textbf{B.} Never share an OTP; refuse, disconnect, and verify
  through the published helpline. Sharing it (A), reading out card
  numbers (C), or calling back the given number (D) all complete the
  fraud.
\item \textbf{A.} Previewing the real address, or typing the known
  address manually, is the safe check. Clicking quickly (B), forwarding
  (C), or replying (D) all risk completing the attack.
\item \textbf{B.} Verify with the sender through a separate channel
  first, since the displayed name may be spoofed. Opening at once (A),
  forwarding (C), or disabling protection (D) all risk the payload.
\item \textbf{A.} The chain is vulnerability (weakness), exploit
  (method), attack (act), payload (harm delivered). B and C scramble
  the order; D substitutes unrelated terms.
\item \textbf{B.} Statements I and II are correct definitions of
  phishing and spoofing. III is false because no genuine bank asks for
  an OTP by phone call. A omits II; C includes false III and omits I;
  D includes false III.
\item \textbf{C.} The payload is the harmful element the attack
  delivers. A swaps vulnerability with exploit; B reverses the same
  pair; D describes spoofing alone as phishing.

\item \textbf{A.} Phishing harvests credentials and data. B, C, and D are hardware theft, not phishing goals.
\item \textbf{A.} Spoofing forges identity. B, C, and D are unrelated device actions.
\item \textbf{A.} Urgent credential demands flag phishing. B, C, and D are legitimate or offline items.
\item \textbf{A, B, C.} Unknown senders, shady links, and threats mark phishing. A paper bill (D) is offline mail.
\item \textbf{A, B, C.} The phishing, spoofing, and password facts are correct. D is false and unsafe.
\item \textbf{A.} Phishing uses deceptive communication to obtain sensitive information. The other terms describe storage, security, or traffic-management processes.

\item \textbf{A.} Spoofing falsifies an identity or source. Patching updates software; hashing computes a digest; sandboxing isolates execution.

\item \textbf{A.} Social engineering manipulates people rather than relying solely on a technical vulnerability. The other options are unrelated operations.

\item \textbf{A.} A vulnerability is a weakness; an exploit takes advantage of it. A patch fixes defects, a backup preserves data, and a checksum detects changes.

\item \textbf{A, B, C.} Sender mismatches, credential pressure, and deceptive links are warning signs. A routine configured update is not by itself evidence of phishing.

\end{enumerate}
\answersubheading{44}{Cybersecurity Prevention}
\begin{enumerate}[label=\textbf{Q\arabic*.}, leftmargin=*]
\item \textbf{A.} MFA stands for Multi-Factor Authentication. B, C, and D are
  not expansions of the abbreviation.
\item \textbf{A.} Confidentiality means only authorised persons can read the
  data. B describes integrity; C describes availability; D describes backup.
\item \textbf{A.} Integrity means the data is accurate and unaltered by
  unauthorised means. B, C, and D describe unrelated states.
\item \textbf{A.} Availability means the data and systems are reachable when
  needed. B, C, and D describe password sharing, encryption, and skipping
  backups, none of which defines availability.
\item \textbf{A.} A signature is a stored pattern of a known virus used for
  recognition. B, C, and D describe handwriting, a written password, and a
  firewall rule.
\item \textbf{B.} Real-time scanning checks files as they are opened,
  downloaded, or copied. A, C, and D name occasions when no scanning occurs.
\item \textbf{A.} Quarantine isolates a suspicious file so it cannot run, and
  it can later be restored or deleted. B describes deletion; C and D describe
  spreading and installing malware.
\item \textbf{B.} A firewall monitors network traffic and blocks what its
  rules forbid. A is the antivirus job; C is the backup job; D is password
  management.
\item \textbf{C.} Clicking every link in unexpected prize messages invites
  phishing and malware. Prompt updates (A), scanning USB drives (B), and
  distinct strong passwords (D) are all good habits.
\item \textbf{A.} A patch is a small vendor fix for a known security
  weakness. B, C, and D describe hardware, malware, and a backup copy.
\item \textbf{D.} \texttt{RaHuL\#2026!Desk} is long and mixes cases, digits,
  and symbols. A, B, and C are short, predictable, or commonly attacked.
\item \textbf{A.} MFA demands a second independent proof, so a leaked
  password alone is not enough. B and C misstate what MFA does; D is the
  antivirus job.
\item \textbf{C.} Backup restores data after loss, protecting availability.
  Confidentiality (A) is protected by passwords and encryption; integrity (B)
  by clean copies and signatures; authentication (D) is not a leg of the
  triad.
\item \textbf{A.} Backups must be regular, offline or off-site, with restores
  tested. B keeps both copies in one failure zone; C never verifies; D gives
  copies to strangers.
\item \textbf{A.} HTTPS and the padlock show an encrypted session before
  sensitive entry. B, C, and D are signs of advertising, scams, or untrusted
  links.
\item \textbf{B.} WPA2 or WPA3 with a strong password protects the wireless
  link. A leaves traffic open; C keeps a well-known password; D joins hostile
  networks automatically.
\item \textbf{B.} Scan first with autorun disabled, then open only what is
  needed. A lets malware launch itself; C spreads it unseen; D multiplies the
  exposure.
\item \textbf{B.} Statements II and III are correct: antivirus removes
  malicious files, and an offline backup restores work after ransomware. I is
  false because a firewall blocks traffic; it does not clean infected files.
  A omits II and III; C includes false I; D includes false I.
\item \textbf{C.} MFA is a second proof of identity beyond the password. A
  swaps in the antivirus job; B describes deletion, not quarantine; D swaps in
  the firewall job.
\item \textbf{A.} Disabling autorun with USB scanning blocks the entry route,
  and offline backups undo the ransomware damage. B, C, and D either do
  nothing for security or actively worsen it.

\item \textbf{A.} Firewalls filter network traffic. B, C, and D describe unrelated physical chores.
\item \textbf{A.} Two factors mean two identity proofs. B, C, and D duplicate objects rather than proofs.
\item \textbf{A.} Length plus mixed characters resists guessing. B, C, and D are classic weak choices.
\item \textbf{A, B, C.} Updates, strong passwords, and firewalls protect systems. Sharing passwords (D) weakens security.
\item \textbf{A, B, C.} Backups, updates, and attachment caution are correct. D is false: public Wi-Fi needs care.
\item \textbf{A.} Confidentiality restricts disclosure to authorized parties. Availability concerns access when needed; redundancy and compression are different concepts.

\item \textbf{A.} A firewall filters or controls network traffic. Backups, monitor repair, and OCR are separate functions.

\item \textbf{A.} Updates often patch known weaknesses but cannot guarantee complete security or replace backups/authentication.

\item \textbf{A.} A backup provides a recoverable copy. Firewalls filter traffic, screen savers lock/display, and spam filters classify email.

\item \textbf{A, B, C.} MFA, unique passwords, and updates strengthen security. Password reuse increases the impact of a compromised account.

\end{enumerate}
\answersubheading{56}{Cybersecurity, IT Act and Security Controls}
\begin{enumerate}[label=\textbf{Q\arabic*.}, leftmargin=*]
\item \textbf{A.} MFA stands for Multi-Factor Authentication. B, C, and D
  combine familiar words but are not the expansion of the abbreviation.
\item \textbf{A.} SSL stands for Secure Sockets Layer. B, C, and D are
  invented expansions using security words.
\item \textbf{B.} The Information Technology Act is of the year 2000. A
  predates it; C is the year of later amendments, not a fresh Act; D is a
  still later year with no such Act.
\item \textbf{D.} Electronic records and signatures (A), e-governance (B),
  and cyber contraventions and offences (C) are all within the Act's scope;
  regulating hardware prices and typing-test fees (D) is not.
\item \textbf{A.} Malware is the umbrella term; virus, worm, Trojan, and
  ransomware are its types. A firewall (B) filters traffic, a patch (C)
  fixes flaws, and a backup (D) recovers data --- none is the umbrella term.
\item \textbf{C.} The virus attaches to a host file and needs human action
  to spread. A worm (A) is standalone and self-spreading; a Trojan (B)
  disguises itself and opens a backdoor; ransomware (D) locks files for
  payment.
\item \textbf{B.} The worm is standalone, self-copying, and needs no host
  file or user action. A virus (A) needs both a host and an action; a
  Trojan (C) needs installation under disguise; D is not a malware type.
\item \textbf{C.} A Trojan pretends to be useful (A), opens a backdoor
  (B), and relies on disguised installation (D) --- but it does not
  replicate itself, so C is the false statement.
\item \textbf{B.} Ransomware encrypts or locks files and demands payment.
  Filtering traffic (A) is the firewall's job; isolating files (C) is
  quarantine; installing updates (D) is patching.
\item \textbf{B.} Statements I and II are correct: the virus needs a host
  plus action, and the worm is standalone and self-spreading. III is false
  because the Trojan does not replicate. A omits II; C and D include false
  III.
\item \textbf{B.} Quarantine isolates a suspicious file so it cannot run,
  pending a decision. Deletion (A) removes it; backup (C) copies it for
  recovery; encryption for sending (D) protects it in transit.
\item \textbf{A.} Phishing is a fake message impersonating a trusted sender
  to steal credentials or money. B describes a firewall; C describes
  patching; D describes SSL/TLS.
\item \textbf{C.} The correct response is to avoid the link, share no OTP,
  and verify through the known official channel. A, B, and D each comply
  with the attacker's instructions and hand over the credential or install
  the payload.
\item \textbf{A.} Encryption gives confidentiality by hiding content from
  unauthorised readers. Proving identity (B) is the digital signature;
  packet filtering (C) is the firewall; isolation (D) is quarantine.
\item \textbf{C.} A digital signature proves the sender and intactness. It
  does not hide content (A --- that is encryption), remove malware (B ---
  that is antivirus), or create the transit channel itself (D --- that is
  SSL/TLS).
\item \textbf{B.} SSL/TLS provide a secure encrypted channel for data in
  transit. They do not clean infected files (A --- antivirus does that),
  add login factors (C --- MFA does that), or install updates (D ---
  patching does that).
\item \textbf{A.} The firewall filters traffic and blocks unauthorised
  access. B, C, and D are commercial antivirus brand names listed only as
  distractors; none of them is the network-traffic barrier control.
\item \textbf{B.} Patching installs vendor updates fixing known flaws.
  Demanding payment (A) is ransomware; adding a login check (C) is MFA;
  copying data for recovery (D) is backup.
\item \textbf{C.} The second factor (OTP on the phone) blocks the attacker
  despite the stolen password --- the value of Multi-Factor
  Authentication. Quarantine (A) isolates files; ransomware (B) locks
  files; phishing (D) is the deception itself.
\item \textbf{D.} Patching fixes security flaws by installing updates, so
  this pairing is correct. A misdescribes quarantine (isolation, not
  deletion); B attaches the firewall's job to a brand-name distractor;
  C gives the digital signature's job to encryption.

\item \textbf{A.} The IT Act was passed in 2000. B, C, and D name other years.
\item \textbf{A.} Hacking with damage is an offence. B, C, and D are harmless local acts.
\item \textbf{A.} Digital signatures prove authenticity and integrity. B, C, and D are physical measurements.
\item \textbf{A, B, C.} Passwords, firewalls, and updates protect systems. Sharing OTPs (D) destroys security.
\item \textbf{A, B, C.} The cybercrime coverage, signature recognition, and hacking penalty are correct. D is false: the Act regulates use; it bans nothing of the sort.
\item \textbf{A.} A direct-action file virus infects executable files. Boot-sector viruses target boot records; macro viruses use document macros.

\item \textbf{D.} Norton is an antivirus product. A Trojan, worm, and macro virus are malware.

\item \textbf{A.} A firewall filters network traffic. The other choices are an application, storage utility, and text-recognition technology.

\item \textbf{A.} A Trojan disguises malicious behavior as benign software. Patches, backups, and firewalls are defensive measures.

\item \textbf{A, B, C.} Antivirus, firewall, and MFA are controls. A phishing email is an attack vector.

\end{enumerate}
\answersubheading{64}{Current Cybersecurity, Cryptography and Cyber Law}
\begin{enumerate}[label=\textbf{Q\arabic*.}, leftmargin=*]
\item \textbf{A.} A parasitic (file-infector) virus attaches to an executable
  file and runs when the host runs. B describes a worm (no host needed); C
  describes a Trojan (disguise, no self-replication); D describes DDoS.
\item \textbf{A.} DES is the Data Encryption Standard. B, C, and D are not
  the correct full form.
\item \textbf{C.} DES works on 64-bit blocks. A (32) and D (128) are other
  ciphers' sizes; B (56) is the effective key size, not the block size.
\item \textbf{D.} A signature proves authenticity (A), non-alteration (B),
  and integrity (C); it does not hide the body --- that is encryption. D is
  the false statement (TRAP-15).
\item \textbf{A.} Encryption gives confidentiality; a signature gives
  authenticity and integrity. B swaps the pair; C and D use availability,
  which belongs to the DDoS threat, not to these two verbs.
\item \textbf{A.} Per-packet header checks are packet filtering. Stateful (B)
  adds connection state; the proxy (C) inspects application content; an
  anti-virus (D) disinfects files rather than filtering packets.
\item \textbf{B.} Tracking connection state and dropping stray packets is
  stateful inspection. A decides per packet without sessions; C and D are
  malware types, not firewalls.
\item \textbf{C.} The go-between inspecting application content is the
  application (proxy) firewall. A checks headers only; B tracks sessions; D
  is only a mail-client feature.
\item \textbf{A.} A fake bank/KYC mail harvesting OTPs is phishing. B locks
  files for ransom; C floods a server; D is unrelated maintenance.
\item \textbf{B.} Locked files plus a payment demand is ransomware. A steals
  credentials via trickery; C floods a service; D describes no concrete
  mechanism.
\item \textbf{B.} This option assigns each threat its correct mechanism. A, C,
  and D scramble the definitions (flood = DDoS, trick = phishing, lock =
  ransomware; a firewall or backup is not a threat).
\item \textbf{A.} DDoS is Distributed Denial of Service; the distributed
  prefix (many machines) is the tested point. B, C, and D are not expansions.
\item \textbf{B.} I and II are correct: DES is symmetric and uses 64-bit
  blocks. III is false because secrecy of the body is encryption, not
  signature (TRAP-15). A omits II; C and D include false III.
\item \textbf{B.} Stateful inspection tracks connection state plus headers.
  A and C swap the packet-filter and proxy roles; D assigns disinfection
  (anti-virus work) to a packet filter.
\item \textbf{A.} Cross-border applicable law (jurisdiction) is a genuine
  governance concern, with security, privacy, and access. B, C, and D are
  single-office trivia, not governance issues.
\item \textbf{A.} Purpose limitation restricts use to the stated purpose.
  Unlimited retention (B), sharing without consent (C), and default
  publication (D) all violate protection principles.
\item \textbf{A.} Unlicensed copying of software is infringement. Fair dealing
  (B) covers limited uses such as study; first ownership (C) is who owns
  first; a lawful backup (D) requires licence permission.
\item \textbf{B.} Meaning = owner's exclusive rights; first ownership = who
  owns first; infringement = what counts as violation; fair dealing = limited
  permitted uses such as study. A, C, and D rotate the matches.
\item \textbf{A.} ICANN is the Internet Corporation for Assigned Names and
  Numbers. B, C, and D are not the correct full form.
\item \textbf{A.} A multi-machine flood blocking access is DDoS, met at the
  network barrier. B answers phishing, not flooding; C answers a file virus
  with a printer driver; D treats an attack as a copyright exception.

\item \textbf{A.} Encryption keeps data confidential. B, C, and D are unrelated device effects.
\item \textbf{A.} Signatures prove who sent data and that it is unaltered. B, C, and D are non-security freebies.
\item \textbf{A.} The IT Act covers cybercrime. B, C, and D are invented distractors.
\item \textbf{A, B, C.} Shared-key use, AES, and bulk speed mark symmetric ciphers. Key pairs (D) belong to asymmetric crypto.
\item \textbf{A, B, C.} The encryption, decryption, and hashing roles are correct. D is false: encryption does not speed printing.
\item \textbf{A.} S/MIME can encrypt the message itself. IMAP retrieves mail, SMTP transfers it, and TLS/SSL on one connection hop is transport protection rather than end-to-end message encryption.

\item \textbf{A.} A Trojan hides a malicious purpose inside something that appears benign. Self-replication describes a worm; the other options describe defensive controls.

\item \textbf{A.} Encryption protects confidentiality. Availability concerns access; signatures can support non-repudiation/integrity; compression reduces size.

\item \textbf{A.} A digital signature supports sender authenticity and message integrity. Encryption, not signing alone, provides confidentiality.

\item \textbf{A, B, C.} The first three are standard roles. Antivirus reduces risk but cannot guarantee prevention of every attack.

\end{enumerate}
\unitopener{11}{Governance \& Integration}
\lessonheading{47}{Humanware and IT Governance}
\begin{enumerate}[label=\textbf{Q\arabic*.}, leftmargin=*]
\item \textit{[Original]} What does the term \textbf{humanware} refer to in a
  computer system?
  \begin{enumerate}[label=\Alph*.]
    \item The physical parts of the computer
    \item The programs and instructions
    \item The people involved with the computer system
    \item The firmware stored in the ROM
  \end{enumerate}
\item \textit{[Original]} Which of the following is \textbf{NOT} a category of
  humanware?
  \begin{enumerate}[label=\Alph*.]
    \item Users
    \item Operators
    \item System administrators
    \item Firmware
  \end{enumerate}
\item \textit{[Original]} Expand ICT.
  \begin{enumerate}[label=\Alph*.]
    \item Information and Computer Technology
    \item Information and Communication Technology
    \item Internet and Communication Technology
    \item Integrated Computer Technology
  \end{enumerate}
\item \textit{[Original]} Which component of a computer system consists of the
  people who design, operate and use it?
  \begin{enumerate}[label=\Alph*.]
    \item Hardware
    \item Software
    \item Humanware
    \item Firmware
  \end{enumerate}
\item \textit{[Original]} Which statement correctly distinguishes
  \textbf{IT governance} from \textbf{e-governance}?
  \begin{enumerate}[label=\Alph*.]
    \item IT governance is governance \emph{of} IT; e-governance is IT
      \emph{in} governance.
    \item IT governance is IT \emph{in} governance; e-governance is
      governance \emph{of} IT.
    \item They mean exactly the same thing.
    \item IT governance is a type of computer hardware.
  \end{enumerate}
\item \textit{[Original]} E-governance refers to:
  \begin{enumerate}[label=\Alph*.]
    \item the use of ICT by government to deliver services and interact with
      citizens.
    \item the management of a company's internal IT resources.
    \item the repair of government computers.
    \item the sale of software to the government.
  \end{enumerate}
\item \textit{[Original]} In e-governance, G2C stands for:
  \begin{enumerate}[label=\Alph*.]
    \item Government to Company
    \item Government to Citizen
    \item Government to Community
    \item Government to Court
  \end{enumerate}
\item \textit{[Original]} Which pairing of an e-governance interaction type
  with its meaning is correctly matched?
  \begin{enumerate}[label=\Alph*.]
    \item G2B --- Government to Employee
    \item G2G --- Government to Citizen
    \item G2E --- Government to Employee
    \item G2C --- Government to Business
  \end{enumerate}
\item \textit{[Original]} Which of the following is a benefit of using IT in
  governance?
  \begin{enumerate}[label=\Alph*.]
    \item Increased transparency and accountability
    \item Increased use of paperwork
    \item Slower processing of applications
    \item Greater need for citizens to visit offices in person
  \end{enumerate}
\item \textit{[Original]} Evaluate the statements:
  \begin{enumerate}[label=\Roman*.]
    \item E-governance uses ICT to deliver government services to citizens.
    \item IT governance means using IT to govern the state.
    \item G2G means Government to Government.
  \end{enumerate}
  \begin{enumerate}[label=\Alph*.]
    \item I only
    \item I and II only
    \item I and III only
    \item I, II and III
  \end{enumerate}
\item \textit{[Original]} An \textbf{electronic record} is:
  \begin{enumerate}[label=\Alph*.]
    \item a printed paper file kept in a government office.
    \item information created, stored and retrieved in electronic form.
    \item a hardware device for storing paper.
    \item a type of computer virus.
  \end{enumerate}
\item \textit{[Original]} Which of the following is an example of an
  electronic record?
  \begin{enumerate}[label=\Alph*.]
    \item A scanned land document stored on a computer
    \item A typed letter kept only on paper
    \item A printed certificate
    \item A handwritten office register
  \end{enumerate}
\item \textit{[Original]} Authentication is the process of:
  \begin{enumerate}[label=\Alph*.]
    \item granting permission to a verified user.
    \item verifying that a user or system is who or what it claims to be.
    \item encrypting a message before sending it.
    \item backing up a database.
  \end{enumerate}
\item \textit{[Original]} Which statement correctly distinguishes
  \textbf{authentication} from \textbf{authorization}?
  \begin{enumerate}[label=\Alph*.]
    \item Authentication grants permission; authorization proves identity.
    \item Authentication proves identity; authorization grants permission.
    \item Both authentication and authorization prove identity.
    \item Both authentication and authorization grant permission.
  \end{enumerate}
\item \textit{[Original]} Expand OTP.
  \begin{enumerate}[label=\Alph*.]
    \item One-Time Password
    \item Online Transfer Password
    \item Open Transaction Protocol
    \item One-Time Program
  \end{enumerate}
\item \textit{[Original]} Two-factor authentication means using:
  \begin{enumerate}[label=\Alph*.]
    \item two passwords of different lengths.
    \item two different authentication factors.
    \item two separate computers to log in.
    \item two firewalls on the same network.
  \end{enumerate}
\item \textit{[Original]} A fingerprint used for login is an example of which
  authentication factor?
  \begin{enumerate}[label=\Alph*.]
    \item Something you know
    \item Something you have
    \item Something you are
    \item Something you forget
  \end{enumerate}
\item \textit{[Original]} Which of the following is an example of a
  \textbf{digital public service}?
  \begin{enumerate}[label=\Alph*.]
    \item Applying for a certificate through an online government portal
    \item Buying goods from a private online shop
    \item Sending a personal email to a friend
    \item Using a word processor to type a letter
  \end{enumerate}
\item \textit{[Original]} Who is responsible for managing user accounts and
  access rights on a computer system?
  \begin{enumerate}[label=\Alph*.]
    \item The end user
    \item The system administrator
    \item The data-entry operator
    \item The programmer
  \end{enumerate}
\item \textit{[Original]} Which pairing is correctly matched?
  \begin{enumerate}[label=\Alph*.]
    \item User --- writes and maintains the software
    \item Operator --- runs the computer system and routine jobs
    \item Programmer --- manages the network
    \item Technician --- uses the finished system to do office work
  \end{enumerate}

\item \textit{[JKSSB FAA-pattern]} Humanware refers to:
  \begin{enumerate}[label=\Alph*.]
    \item the people who use and manage computer systems.
    \item a printer hardware model.
    \item a type of monitor glass.
    \item a network cable brand.
  \end{enumerate}
\item \textit{[JKPSI-pattern]} G2C in e-governance stands for:
  \begin{enumerate}[label=\Alph*.]
    \item Government to Citizen
    \item Government to Company printer
    \item Gateway to Cable
    \item Government to Cloud speaker
  \end{enumerate}
\item \textit{[SSC CGL-pattern]} Which of the following is an e-governance service?
  \begin{enumerate}[label=\Alph*.]
    \item Online filing of certificates and forms
    \item Polishing office furniture
    \item Watering office plants
    \item Painting office walls
  \end{enumerate}
\item \textit{[JKSSB-pattern]} \textbf{(MSQ: select ALL correct options.)} Which roles fall under humanware?
  \begin{enumerate}[label=\Alph*.]
    \item System users
    \item Programmers
    \item System administrators
    \item Printer ink cartridges
  \end{enumerate}
\item \textit{[SSC CGL-pattern]} \textbf{(MSQ: select ALL correct options.)} Which statements about IT governance and humanware are correct?
  \begin{enumerate}[label=\Alph*.]
    \item Humanware means the people around the system.
    \item E-governance delivers government services digitally.
    \item Trained staff reduce security mistakes.
    \item Hardware alone guarantees good governance.
  \end{enumerate}
\item \textit{[Original]} In an information system, the people who use, operate, or manage the system are sometimes collectively called:
  \begin{enumerate}[label=\Alph*.]
\item Humanware
   \item Firmware
   \item Shareware
   \item Middleware
  \end{enumerate}

\item \textit{[Original]} IT governance is primarily concerned with:
  \begin{enumerate}[label=\Alph*.]
\item Ensuring IT supports organizational goals and is responsibly managed
   \item Increasing screen resolution only
   \item Replacing all staff with computers
   \item Removing accountability for systems
  \end{enumerate}

\item \textit{[Original]} Which role is most associated with managing user accounts and access permissions?
  \begin{enumerate}[label=\Alph*.]
\item System administrator
   \item End user only
   \item Printer operator only
   \item Search engine crawler
  \end{enumerate}

\item \textit{[Original]} An end user is someone who primarily:
  \begin{enumerate}[label=\Alph*.]
\item Uses an application or system to perform a task
   \item Writes every operating-system kernel
   \item Manufactures all processors
   \item Routes every Internet packet
  \end{enumerate}

\item \textit{[Original, MSQ drill]} \textbf{(MSQ: select ALL correct options.)} Which are stakeholders in IT governance?
  \begin{enumerate}[label=\Alph*.]
\item Senior management
   \item IT staff
   \item System users
   \item External auditors or regulators where applicable
  \end{enumerate}

\end{enumerate}
\lessonheading{57}{Digital Platforms and E-Governance}
\begin{enumerate}[label=\textbf{Q\arabic*.}, leftmargin=*]
\item \textit{[Original]} Expand e-NAM.
  \begin{enumerate}[label=\Alph*.]
    \item National Agricultural Marketing
    \item National Agriculture Market
    \item New Agriculture Market
    \item National Agro Market
  \end{enumerate}
\item \textit{[Original]} Expand NeGP.
  \begin{enumerate}[label=\Alph*.]
    \item National e-Governance Programme
    \item New e-Government Plan
    \item National e-Governance Plan
    \item National Electronic Government Project
  \end{enumerate}
\item \textit{[Original]} Expand NeGD.
  \begin{enumerate}[label=\Alph*.]
    \item National e-Governance Department
    \item National e-Governance Division
    \item National e-Government Directorate
    \item National e-Governance Directorate
  \end{enumerate}
\item \textit{[Original]} In e-Governance, G2C stands for:
  \begin{enumerate}[label=\Alph*.]
    \item Government to Company
    \item Government to Community
    \item Government to Citizen
    \item Government to Country
  \end{enumerate}
\item \textit{[Original]} e-NAM is primarily associated with which field?
  \begin{enumerate}[label=\Alph*.]
    \item Online banking
    \item Telemedicine
    \item Railway reservation
    \item Agricultural marketing
  \end{enumerate}
\item \textit{[Original]} e-NAM networks which existing markets to create a
  unified national market?
  \begin{enumerate}[label=\Alph*.]
    \item Stock exchanges
    \item APMC mandis
    \item Shopping malls
    \item Supermarkets
  \end{enumerate}
\item \textit{[Original]} e-NAM was launched in which year?
  \begin{enumerate}[label=\Alph*.]
    \item 2014
    \item 2015
    \item 2016
    \item 2017
  \end{enumerate}
\item \textit{[Original]} Which of the following is an example of a G2C
  (Government to Citizen) service?
  \begin{enumerate}[label=\Alph*.]
    \item GeM
    \item e-Office
    \item PFMS
    \item UMANG
  \end{enumerate}
\item \textit{[Original]} The Government e-Marketplace (GeM) is best
  classified as:
  \begin{enumerate}[label=\Alph*.]
    \item G2C
    \item G2B
    \item G2G
    \item G2E
  \end{enumerate}
\item \textit{[Original]} e-Office, which enables paperless working between
  government departments, is an example of:
  \begin{enumerate}[label=\Alph*.]
    \item G2C
    \item G2B
    \item G2G
    \item G2E
  \end{enumerate}
\item \textit{[Original]} Which of the following is \textbf{NOT} one of the
  four standard e-Governance interaction categories?
  \begin{enumerate}[label=\Alph*.]
    \item G2C
    \item G2B
    \item G2G
    \item G2S
  \end{enumerate}
\item \textit{[Original]} Digital India was launched on:
  \begin{enumerate}[label=\Alph*.]
    \item 1 July 2015
    \item 15 August 2014
    \item 1 January 2016
    \item 26 January 2015
  \end{enumerate}
\item \textit{[Original]} Which of the following is \textbf{NOT} one of the
  nine pillars of the Digital India programme?
  \begin{enumerate}[label=\Alph*.]
    \item Broadband Highways
    \item e-Kranti: Electronic Delivery of Services
    \item Information for All
    \item Digital Signature Infrastructure
  \end{enumerate}
\item \textit{[Original]} e-Kranti is best described as:
  \begin{enumerate}[label=\Alph*.]
    \item a separate country-wide broadband scheme
    \item the organisation that manages e-Governance projects
    \item Pillar 5 of Digital India and the upgraded National e-Governance
      Plan (NeGP 2.0)
    \item the portal for online government procurement
  \end{enumerate}
\item \textit{[Original]} Which organisation was created by MeitY to support
  the programme management and implementation of e-Governance projects?
  \begin{enumerate}[label=\Alph*.]
    \item NeGP
    \item e-Kranti
    \item NeGD
    \item NASSCOM
  \end{enumerate}
\item \textit{[Original]} Consider the statements:
  \begin{enumerate}[label=\Roman*.]
    \item The National e-Governance Plan (NeGP) was approved in May 2006.
    \item e-Kranti is a pillar of the Digital India programme.
    \item The National e-Governance Division (NeGD) is a plan approved by
      the Cabinet.
  \end{enumerate}
  \begin{enumerate}[label=\Alph*.]
    \item I and II only
    \item II and III only
    \item I and III only
    \item I, II and III
  \end{enumerate}
\item \textit{[Original]} The National Portal of India, a single-window
  gateway to government information and services, is available at:
  \begin{enumerate}[label=\Alph*.]
    \item india.gov.in
    \item negd.gov.in
    \item digitalindia.gov.in
    \item gem.gov.in
  \end{enumerate}
\item \textit{[Original]} Which pairing of a named platform with its purpose
  is correctly matched?
  \begin{enumerate}[label=\Alph*.]
    \item e-NAM --- online government procurement of goods
    \item DigiLocker --- secure digital wallet for authentic documents
    \item UMANG --- agricultural commodity trading
    \item GeM --- digital document storage
  \end{enumerate}
\item \textit{[Original]} Which of the following provides authenticity and
  integrity of an electronic record?
  \begin{enumerate}[label=\Alph*.]
    \item Encryption
    \item Digital signature
    \item Compression
    \item Backup
  \end{enumerate}
\item \textit{[Original]} Match each platform with its purpose:
  \begin{enumerate}[label=(\roman*)]
    \item e-NAM
    \item MyGov
    \item e-Office
  \end{enumerate}
  \begin{enumerate}[label=(\alph*)]
    \item Citizen engagement and participative governance
    \item Online agricultural commodity trading
    \item Paperless official file management
  \end{enumerate}
  Which option gives all correct matches?
  \begin{enumerate}[label=\Alph*.]
    \item (i)-a, (ii)-b, (iii)-c
    \item (i)-b, (ii)-a, (iii)-c
    \item (i)-c, (ii)-a, (iii)-b
    \item (i)-b, (ii)-c, (iii)-a
  \end{enumerate}

\item \textit{[JKSSB FAA-pattern]} UMANG is a platform for accessing:
  \begin{enumerate}[label=\Alph*.]
    \item government services on mobile.
    \item cinema tickets only.
    \item printer drivers only.
    \item music albums only.
  \end{enumerate}
\item \textit{[JKPSI-pattern]} DigiLocker provides citizens with:
  \begin{enumerate}[label=\Alph*.]
    \item digital storage for verified documents.
    \item free printer ink.
    \item louder speakers.
    \item paper envelopes.
  \end{enumerate}
\item \textit{[SSC CGL-pattern]} The Digital India programme was launched in the year:
  \begin{enumerate}[label=\Alph*.]
    \item 2015
    \item 1995
    \item 2005
    \item 2025
  \end{enumerate}
\item \textit{[JKSSB-pattern]} \textbf{(MSQ: select ALL correct options.)} Which are Indian digital/e-governance platforms?
  \begin{enumerate}[label=\Alph*.]
    \item UMANG
    \item DigiLocker
    \item UPI
    \item Paper telegram
  \end{enumerate}
\item \textit{[SSC CGL-pattern]} \textbf{(MSQ: select ALL correct options.)} Which statements about digital platforms are correct?
  \begin{enumerate}[label=\Alph*.]
    \item UPI enables instant bank payments by phone.
    \item DigiLocker holds verified digital documents.
    \item UMANG gathers many government services in one app.
    \item UPI is a paper money-order form.
  \end{enumerate}
\item \textit{[PYQ-adapted: JKSSB FAA 2024, QCRIAS item 672; SRC-16/SRC-19]} The UMANG app enables citizens to access:
  \begin{enumerate}[label=\Alph*.]
\item Virtual-reality games
   \item Streaming movies
   \item Government services and information
   \item Online fitness classes
  \end{enumerate}

\item \textit{[PYQ-adapted: JKSSB FAA 2024, QCRIAS item 671; SRC-16/SRC-19]} Smart Governance primarily involves:
  \begin{enumerate}[label=\Alph*.]
\item Making governance processes more complicated
   \item Integrating IT to improve service delivery and decisions
   \item Avoiding technology
   \item Reducing transparency
  \end{enumerate}

\item \textit{[Original]} A service that lets a citizen access services from central and state agencies through one mobile platform is an example of:
  \begin{enumerate}[label=\Alph*.]
\item Integrated e-governance delivery
   \item Disk defragmentation
   \item Local file indexing
   \item Spreadsheet recalculation
  \end{enumerate}

\item \textit{[Original]} A Government-to-Citizen service is abbreviated:
  \begin{enumerate}[label=\Alph*.]
\item G2C
   \item G2G
   \item G2B
   \item G2E
  \end{enumerate}

\item \textit{[Original, MSQ drill]} \textbf{(MSQ: select ALL correct options.)} Which are common e-governance interaction categories?
  \begin{enumerate}[label=\Alph*.]
\item G2C
   \item G2G
   \item G2B
   \item B2C means Government to Citizen
  \end{enumerate}

\end{enumerate}
\answersection{11}
\answersubheading{47}{Humanware and IT Governance}
\begin{enumerate}[label=\textbf{Q\arabic*.}, leftmargin=*]
\item \textbf{C.} Humanware (also called liveware or peopleware) is the people
  involved with a computer system. The physical parts (A) are hardware, the
  programs (B) are software, and firmware (D) is software stored in hardware.
\item \textbf{D.} Firmware is software stored in hardware, not a category of
  humanware. Users (A), operators (B), and system administrators (C) are all
  people, so they are humanware.
\item \textbf{B.} ICT stands for Information and Communication Technology. A,
  C, and D are not the correct expansions of the abbreviation.
\item \textbf{C.} The people who design, operate and use a system are its
  humanware. Hardware (A) is the machinery, software (B) is the programs, and
  firmware (D) is a type of stored software.
\item \textbf{A.} IT governance is the direction and control of an
  organisation's own IT resources (governance of IT); e-governance is the use
  of IT to govern and deliver public services (IT in governance). B reverses
  the two; C wrongly treats them as identical; D misclassifies a governance
  term as hardware.
\item \textbf{A.} E-governance is the use of ICT by government to deliver
  services and interact with citizens. Managing a company's internal IT (B) is
  IT governance, and repairing computers (C) or selling software (D) are
  operational or commercial activities, not e-governance.
\item \textbf{B.} G2C stands for Government to Citizen. A, C, and D are not
  the standard expansions of the term.
\item \textbf{C.} G2E means Government to Employee. A wrongly pairs G2B with
  Employee (it is Government to Business); B wrongly pairs G2G with Citizen
  (it is Government to Government); D wrongly pairs G2C with Business (it is
  Government to Citizen).
\item \textbf{A.} IT in governance increases transparency and accountability.
  More paperwork (B), slower processing (C), and a greater need for physical
  visits (D) are the opposite of what IT in governance achieves.
\item \textbf{C.} Statements I and III are correct: e-governance uses ICT to
  deliver services, and G2G means Government to Government. Statement II is
  wrong because IT governance is governance \emph{of} IT, not the use of IT to
  govern the state --- that is e-governance. A omits III; B includes the false
  II; D includes the false II as well.
\item \textbf{B.} An electronic record is information created, stored and
  retrieved in electronic form. A printed paper file (A) is a physical record,
  not an electronic one; a paper-storage device (C) is not a record; and a
  virus (D) is malware, not a record.
\item \textbf{A.} A scanned land document stored on a computer is an
  electronic record. A paper-only letter (B), a printed certificate (C), and a
  handwritten register (D) are all physical records.
\item \textbf{B.} Authentication verifies that a user or system is who or what
  it claims to be. Granting permission (A) is authorization; encryption (C)
  protects confidentiality; and backup (D) is data protection --- none of
  these is authentication.
\item \textbf{B.} Authentication proves identity (``who are you?''), while
  authorization grants or denies permission (``what may you do?''). A reverses
  the two; C and D wrongly merge the two distinct functions.
\item \textbf{A.} OTP stands for One-Time Password. B, C, and D are not the
  correct expansions of the abbreviation.
\item \textbf{B.} Two-factor authentication uses two \emph{different} factors,
  such as a password (something you know) and an OTP (something you have). Two
  passwords (A) are the same factor twice, so they are not two factors; C and
  D do not describe authentication factors.
\item \textbf{C.} A fingerprint is something you \emph{are}, a biometric
  factor. Something you know (A) is a password or PIN; something you have (B)
  is an OTP or token; D is not a real factor.
\item \textbf{A.} Applying for a certificate through an online government
  portal is a government service delivered digitally. Buying from a private
  shop (B) is e-commerce, not a public service; a personal email (C) and
  typing a letter (D) are ordinary personal computing tasks.
\item \textbf{B.} The system administrator manages user accounts and access
  rights. An end user (A) uses the system; a data-entry operator (C) feeds
  data; a programmer (D) writes software.
\item \textbf{B.} The operator runs the computer system and routine jobs. A
  mislabels the user as a programmer; C mislabels the programmer as a network
  manager; D mislabels the technician as an office user.

\item \textbf{A.} Humanware is the people component. B, C, and D name hardware items.
\item \textbf{A.} G2C means Government to Citizen. B, C, and D are not the recognised expansion.
\item \textbf{A.} Online certificates are e-governance. B, C, and D are manual chores.
\item \textbf{A, B, C.} Users, programmers, and admins are humanware. Cartridges (D) are consumables.
\item \textbf{A, B, C.} The humanware, e-governance, and training facts are correct. D is false: people and process matter too.
\item \textbf{A.} Humanware refers to people involved with computer systems. Firmware, shareware, and middleware are software terms.

\item \textbf{A.} IT governance aligns and oversees technology, risk, and accountability in support of organizational objectives. The other choices are too narrow or incorrect.

\item \textbf{A.} System administrators commonly manage users, permissions, and systems. The other roles do not generally control organization-wide access.

\item \textbf{A.} An end user uses the system for work or personal tasks. The other options describe specialized engineering or network functions.

\item \textbf{A, B, C, D.} Governance involves organizational leaders, IT teams, users, and relevant external oversight stakeholders.

\end{enumerate}
\answersubheading{57}{Digital Platforms and E-Governance}
\begin{enumerate}[label=\textbf{Q\arabic*.}, leftmargin=*]
\item \textbf{B.} e-NAM stands for National Agriculture Market. A (National
  Agricultural Marketing), C (New Agriculture Market), and D (National Agro
  Market) are not the correct full form.
\item \textbf{C.} NeGP stands for National e-Governance Plan. A uses
  ``Programme'', B uses ``New'' and ``Government'', and D uses
  ``Electronic'' instead of ``e-Governance Plan''.
\item \textbf{B.} NeGD stands for National e-Governance Division. A says
  ``Department'', C says ``Government'' and ``Directorate'', and D says
  ``Directorate''; none matches the official name.
\item \textbf{C.} G2C stands for Government to Citizen. A (Company), B
  (Community), and D (Country) are not the standard category.
\item \textbf{D.} e-NAM is the National Agriculture Market, an agricultural
  marketing portal. Online banking (A), telemedicine (B), and railway
  reservation (C) belong to other services.
\item \textbf{B.} e-NAM networks the existing Agriculture Produce Market
  Committee (APMC) mandis to create a unified national market. Stock
  exchanges (A), shopping malls (C), and supermarkets (D) are not the
  markets it links.
\item \textbf{C.} e-NAM was launched on 14 April 2016. The other years (2014,
  2015, 2017) do not match the launch date.
\item \textbf{D.} UMANG is a Government to Citizen (G2C) service. GeM (A) is
  G2B; e-Office (B) and PFMS (C) are G2G.
\item \textbf{B.} GeM is used by government to buy goods and services from
  businesses, so it is Government to Business (G2B). G2C (A) serves citizens;
  G2G (C) serves government bodies; G2E (D) serves government employees.
\item \textbf{C.} e-Office is used for paperless working between government
  departments, so it is Government to Government (G2G). G2C (A) serves
  citizens; G2B (B) serves businesses; G2E (D) serves employees.
\item \textbf{D.} G2S is not one of the standard e-Governance interaction
  categories. The four standard categories are G2C (A), G2B (B), G2G (C),
  and G2E.
\item \textbf{A.} Digital India was launched on 1 July 2015. The other dates
  (15 August 2014, 1 January 2016, 26 January 2015) are not the launch date.
\item \textbf{D.} ``Digital Signature Infrastructure'' is not one of the nine
  pillars. Broadband Highways (A), e-Kranti (B), and Information for All
  (C) are genuine pillars.
\item \textbf{C.} e-Kranti is Pillar 5 of Digital India and the upgraded
  National e-Governance Plan (NeGP 2.0). A describes a broadband scheme; B
  describes NeGD; D describes GeM.
\item \textbf{C.} The National e-Governance Division (NeGD), created in 2009,
  supports MeitY in the programme management and implementation of
  e-Governance projects. NeGP (A) is a plan; e-Kranti (B) is a pillar;
  NASSCOM (D) is an industry body, not a government division.
\item \textbf{A.} Statements I and II are correct: NeGP was approved in May
  2006, and e-Kranti is a pillar of Digital India. Statement III is false
  because NeGD is a division (an organisation), not a plan. B and C include
  false III; D includes all three.
\item \textbf{A.} The National Portal of India is india.gov.in. negd.gov.in
  (B) is the NeGD site; digitalindia.gov.in (C) is the Digital India
  programme site; gem.gov.in (D) is the Government e-Marketplace.
\item \textbf{B.} DigiLocker is a secure digital wallet for authentic
  documents, so this pairing is correct. A wrongly links e-NAM to
  procurement; C wrongly links UMANG to agricultural trading; D wrongly
  links GeM to document storage.
\item \textbf{B.} A digital signature provides authenticity and integrity of
  an electronic record. Encryption (A) provides confidentiality, not
  authenticity; compression (C) reduces size; backup (D) is for recovery.
\item \textbf{B.} e-NAM is online agricultural commodity trading (i)-b;
  MyGov is citizen engagement and participative governance (ii)-a; e-Office
  is paperless official file management (iii)-c. Option A wrongly pairs
  e-NAM with governance (a); C and D scramble the pairs.

\item \textbf{A.} UMANG delivers government services by mobile. B, C, and D are single-purpose non-government uses.
\item \textbf{A.} DigiLocker stores verified digital documents. B, C, and D are unrelated physical items.
\item \textbf{A.} Digital India launched in 2015. B, C, and D name other years.
\item \textbf{A, B, C.} UMANG, DigiLocker, and UPI are digital platforms. A paper telegram (D) is offline.
\item \textbf{A, B, C.} The UPI, DigiLocker, and UMANG facts are correct. D is false: UPI is digital, not paper.
\item \textbf{C.} UMANG is a government-service access platform. The other choices describe entertainment or fitness apps.

\item \textbf{B.} Smart Governance uses technology to improve efficiency, transparency, and service delivery. The other choices contradict those aims.

\item \textbf{A.} Integrated e-governance platforms provide access to services across agencies. The other options are local computer operations.

\item \textbf{A.} G2C means Government to Citizen. G2G is government-to-government, G2B government-to-business, and G2E government-to-employee.

\item \textbf{A, B, C.} G2C, G2G, and G2B are standard categories. B2C means Business to Consumer, not Government to Citizen.

\end{enumerate}